\documentclass[12pt]{article}
\title{Digest of The New York Times}
\author{The New York Times}

% \definecolor{dark-green}{rgb}{0,0.25,0}
% \usepackage[utf8]{inputenc}

\usepackage{fontspec,xunicode,xltxtra}
\setmainfont[Ligatures=Required,Mapping=tex-text]{Georgia}
% \setmainfont[Ligatures=Required,Mapping=tex-text]{Palatino Linotype}
\setsansfont[Ligatures=Required,Mapping=tex-text]{Myriad Pro}
\setmonofont[Ligatures=Required,Mapping=tex-text]{Courier New}
\usepackage{xeCJK}
\setCJKmainfont{宋体}
\setCJKmonofont{黑体}

% \usepackage{indentfirst}


\usepackage[dvips,dvipsnames,svgnames]{xcolor}
\definecolor{light-gray}{gray}{0.95}

%生成PDF的链接
\usepackage[dvipdfmx]{hyperref}
\hypersetup{
    bookmarks=true,         % show bookmarks bar?
    bookmarksopen=true,
    pdfpagemode=UseNone,    % options: UseNode, UseThumbs, UseOutlines, FullScreen
    pdfstartview=FitH,
    pdfborder=1,
    pdfhighlight=/P,
    pdfauthor={wuxch},
    unicode=true,           % non-Latin characters in Acrobat’s bookmarks
    colorlinks,             % false: boxed links; true: colored links
    linkcolor=blue,         % color of internal links
    citecolor=blue,        % color of links to bibliography
    filecolor=magenta,      % color of file links
    urlcolor=cyan           % color of external links
}
% 使用目录
\usepackage{makeidx}
% 使用图形
\usepackage{graphicx}


\usepackage{titletoc}
\dottedcontents{section}[1.5em]{}{2.3em}{1pc}

% % 页眉页脚的设定
% \usepackage{fancyhdr}
% \pagestyle{fancy}
% \lhead{\includegraphics[height=0.2in]{zte_logo}}

% \usepackage{footmisc}
\usepackage{pfnote}
\usepackage{fnpos}

% % 嵌入的代码显示
% \usepackage{listings}
% \lstset{language=C++, breaklines, extendedchars=false}
% \lstset{keywordstyle=\color{blue},
%         commentstyle=\color{SeaGreen},
%         stringstyle=\ttfamily,
%         showstringspaces=false,
%         backgroundcolor=\color{light-gray}}
% \lstset{basicstyle=\small}

\usepackage{booktabs}

\usepackage{tipa}
\usepackage{phonetic}
\usepackage[gen]{eurosym}

% \usepackage{CJK}
\usepackage{makeidx}
% \usepackage{glossary}

\usepackage{fancyhdr}

\usepackage{lastpage}
\cfoot{\textsf{\thepage\ of \pageref*{LastPage}}}

\pagestyle{fancy}
\fancyhead{}
\renewcommand{\headrulewidth}{0pt}
\pagenumbering{arabic}

\usepackage{graphicx}
\usepackage{picinpar}

\setlength{\parindent}{0in}
\setlength{\parskip}{2.5ex plus4mm minus3mm}
% \setlength{\parskip}{.3\baselineskip}

\setcounter{section}{0}

% bibtex
\usepackage[super]{natbib}
% \usepackage[globalcitecopy,labelstoglobalaux]{bibunits}
% \makeatletter
% \renewcommand{\@bibunitname}{\jobname.\the\@bibunitauxcnt}
% \makeatother

% \renewcommand{\refname}{缩略语}
% \renewcommand{\refname}{}

\makeatletter                   % Reference list option change
\renewcommand\@biblabel[1]{#1:  } % from [1] to 1: 
\makeatother                    %

\newcommand{\citenumfont}[1]{\textsf{#1}}

% \usepackage{tipa}
% \usepackage{phonetic}

\newcommand{\phonetic}[1]{{\fontspec[Colour=CC00CC]{Microsoft Sans Serif} #1}}


%%% Local Variables:
%%% TeX-master: "medical_trivia"
%%% End:


\makeindex
\begin{document}
\date{}
\thispagestyle{empty}
\renewcommand\contentsname{\textsf{Digest of The New York Times}}
{\footnotesize\textsf{\tableofcontents}}
\clearpage
\setcounter{page}{1}
\section{New Protests in Western China After Clashes}

\lettrine{S}{ome} women glared through their black veils at the paramilitary troops encircling them.
Others held identity cards of missing relatives in the air. Fists raised, tears in their eyes, they
demanded the release of sons and husbands seized by the police after Muslim Uighurs rioted in this
western regional capital days earlier.

And as the group of several hundred Uighur women beseeched journalists on a government-sponsored
tour here on Tuesday, they gave voice to broader concerns at the heart of the deadliest ethnic
violence to strike China in decades.

``They don't respect our lifestyle,'' said one woman, a 26-year-old who gave her name as Guli. ``We
want our dignity. We just want fairness, and we want equality.''

A wide variety of government policies here in the western desert region of Xinjiang, a lightly
populated area that covers about a sixth of China's total landmass, has for years led many of the
area's 10 million Uighurs to believe their culture and livelihoods were under assault by the Han
Chinese, the dominant ethnic group in China, according to local residents, foreign scholars and
recent studies of the area.

The policies include limits on religious practice, the phasing out of Uighur-language instruction in
schools and the reinforcement of better economic opportunities for the Han, from businesspeople to
migrant workers.

Uighurs are the largest ethnic group in Xinjiang, but Han migration, encouraged in part by
government incentives, is quickly changing the demographics here: census figures show that Han made
up 40 percent of the population in 2000, a huge leap over the 6 percent in 1949. Under the Chinese
Communist Party, Han have always held the power in Xinjiang. Wang Lequan, the party secretary of the
region, is a Han whose hard-line policies have inspired systems of control in other ethnic minority
regions of China, including Tibet.

``Fundamentally, the relationship between Uighur and Han is one of colonized to colonizer,'' said
Nicholas Bequelin, a China researcher at Human Rights Watch who has written about policies in
Xinjiang.

That dynamic may have laid the foundation for the riot on Sunday in which 156 people were killed and
more than 1,000 injured when angry Uighurs attacked Han civilians and battled with security forces
across the city. Government officials declined Tuesday to give an ethnic breakdown of the dead. The
riot began as a protest over government handling of a brawl between Uighur and Han factory workers
in southern China.

On Tuesday afternoon, thousands of Han Chinese armed with sticks, shovels, pipes and meat cleavers
tried marching to the Uighur quarter to exact revenge for those Han civilians who were killed on
Sunday. Paramilitary troops fired tear gas at the mob\cite{mob}, but not before the first wave got
into a brick-throwing battle with Uighurs perched on rooftops near Erdaoqiao Market, where the
rioting began on Sunday.

President Hu Jintao of China cut short a trip to Italy for the G8 summit and flew back to Beijing
early on Wednesday because of the unrest in Xinjiang, the official Xinhua news agency reported.

Many Han Chinese say the Uighurs, like China's 55 other ethnic minorities, actually enjoy generous
advantages under government policies. Uighur women, for example, can give birth to more than one
child without having to pay a fine, unlike the Han. Uighur students have extra points added to their
scores when taking the standardized tests that determine university placement.

But on issues that go to the heart of Uighur identity, the government takes a strict line, many
Uighurs say.

The vast majority of Uighurs are Sunni Muslims, but the practice of Islam is tightly circumscribed.
Government workers are not allowed to practice the religion. Imams cannot teach the Koran in
private, and study of Arabic is allowed only at designated government schools. Two of Islam's five
pillars -- the sacred fasting month of Ramadan and the pilgrimage to Mecca, the hajj -- are also
closely managed: students and government workers are compelled to eat during Ramadan, and passports
of Uighurs have been confiscated to force them to join official hajj tours.

Three years ago, in its annual report on international religious freedom, the State Department
singled out Xinjiang for criticism in a section on China: ``Officials in the Xinjiang Uighur
Autonomous Region tightly controlled religious activity, while elsewhere in the country, Muslims
enjoyed greater religious freedom,'' the report said.

On Tuesday, Abudurehepu, a religious leader in Xinjiang who supports the government, said at a news
conference here that ``our religious freedom is respected,'' noting that Xinjiang had more than
2,000 mosques.

He also said that ``the party and the government have been doing very well on ethnic policy, like
having Uighur kids going to Uighur-language schools.''

In fact, the government is phasing out the use of the Uighur language in schools. Many Uighur
parents know the importance of having their children learn Mandarin Chinese, but they are upset over
the disappearance of their native language from the education system. There are some bilingual
schools, but those generally relegate the Uighur language to a marginal role.

A 2009 Amnesty International report on threats to Uighur identity charts the recent history of the
erosion of the Uighur language in education, beginning with a policy in the 1990s that eliminated
Uighur as a language of instruction at the university level. Today, at Xinjiang University in
Urumqi, only Uighur poetry classes are taught in Uighur, the report says. In 2006, the government
began carrying out policies that make Chinese the main language of preschool instruction.

Since the central government adopted a ``develop the west'' campaign in the past decade, Xinjiang's
economy has grown quickly, and living standards on the whole have risen. But many Uighurs complain
about high unemployment and the growing income gap with Han Chinese, who control the largest
industries in Xinjiang: oil, agriculture and construction. They give many more contracts and jobs to
other Han.

``Uighurs feel cut out of this process,'' said a former resident of Kashgar, an oasis town near
China's western border where more than 200 protesters gathered on Monday.

The bingtuan, vast farms started by the military in the 1950s to employ demobilized troops, are
among Xinjiang's biggest moneymakers. But Mr.~Bequelin, the human rights researcher, said more than
90 percent of employees at bingtuan were Han.

Chinese officials deny that government policies contribute to ethnic unrest. They place blame for
the tensions on outside figures like the Dalai Lama, the Tibetan spiritual leader, or, in the case
of the latest Xinjiang riots, Rebiya Kadeer, a Uighur businesswoman and former political prisoner
who lives in Washington.

On Tuesday, as thousands of Han armed with makeshift weapons tried to attack the Uighur quarter, the
party secretary of Urumqi climbed atop a car and pleaded with them to go home.

\section{After Her Rise in China and Expulsion, a Uighur Becomes the Face of Her People}

\lettrine{A}{s} the global face of resistance to what she calls the worsening Chinese repression of
the Uighurs, Rebiya Kadeer is displaying the tenacity and sense of destiny that drove her improbable
climb inside China in decades past, from laundry girl to famed business mogul.

The Beijing government that hailed her as a model citizen in the 1990s, before imprisoning her for
stealing state secrets and sending her into exile in the United States in 2005, vilifies her as the
unseen hand behind protests that erupted Sunday in the Uighur homeland of western China. ``All the
difficulties in my life prepared me for the tough times we face now,'' said the woman, who is happy
to be called the ``Mother of the Uighurs,'' in an interview on Tuesday.

In a plain wool suit and a traditional Uighur cap topping waist-length pigtails, Ms.~Kadeer, 62,
veered from impish humor and warmth -- she leapt to pump the hand of a reporter who described
visiting her childhood town -- to intense, hand-waving condemnations of Chinese perfidy.

The walls of her small office in downtown Washington are covered with photographs of meetings with
President George W.~Bush and Laura Bush, and pictures of several of her 11 children, two of whom are
now in prison in China. They were sentenced to long terms after she came to the United States and
resumed work for Uighur rights.

The week's events have catapulted Ms.~Kadeer to a new level of global recognition, a prominence that
seems belied by the few modest rooms here where she and a few aides press their cause with
telephones, the Internet and passion.

This week, several office and personal phones rang incessantly, with reporters from around the world
seeking a word. Still, it became clear that the Uighurs, long downtrodden and little known in the
West, enjoy little of the glamour of their neighbors, the Tibetans. When Ms.~Kadeer led a march to
the Chinese Embassy on Tuesday, no more than several dozen supporters, mainly fellow exiles, showed
up.

If she was disappointed, she gave no sign. In the interview and in her autobiography, ``Dragon
Fighter,'' which came out this year, Ms.~Kadeer described her survival through famine, persecution
during Mao's Cultural Revolution and then -- as she threw herself into black-market trading of
cloth, underwear and other items -- the repeated seizure of her goods and money by corrupt or
overzealous officials.

She claims that she had, from the beginning, an irrepressible devotion to Uighur self-determination.
In her eyes, even her start in life brought an omen. Money and luck were running out in the mining
settlement where her father hoped to strike it rich, she wrote, in a story that may be too good to
investigate.

In accordance with tradition, her father went to bury the bloody birth linens. As he dug a hole, he
suddenly shouted, ``Gold!'' From that moment on, she wrote, her parents said, ``You don't belong to
us; you belong to the people.''

What is indisputable is that from early on she was a determined and shrewd businesswoman willing to
sell goods from a sack at the side of the road when necessary, buying and selling thousands of
sheepskins or logs when she saw the chance. As China's economy opened up in the 1980s, she expanded
into real estate and flourished. By the 1990s she was running trading companies all over Central
Asia, had built a famous women's bazaar and then a seven-story department store in Urumqi, the
capital of the region of Xinjiang, and ran a charity for Uighur women.

Her career had personal costs. In an unthinkable violation of Uighur custom, and angering her
relatives, she traveled for months at a time, leaving young children with a working husband or
relatives. ``Of course it was difficult for me as a woman to leave my children,'' she said. ``But I
found out that money is very important to the destiny of a nation, and I decided to find that
money.''

Five of her children are now in the United States, and have been working computers and phones night
and day this week, she said. Another five, including the two in prison, remain in China, and one
lives in Australia.

In the mid-1990s, as Chinese officials heralded her as an example of ethnic success and even made
her a member of the national legislature, she tried to work for change and never lost sight of her
political dream, Ms.~Kadeer said.

``I was sincere in my interactions with the Chinese government. I was hoping to solve the problems
of the Uighurs,'' she said. ``I still believe we can solve the problems.''

But she started speaking out about Uighur grievances and she kept ties with her husband, by then a
dissident living in the United States. In 1999 she was imprisoned.

Ms.~Kadeer dismisses Beijing's charge that she planned last Sunday's protests.

She is more than happy, however, to tell how she and the two organizations she heads, the Uighur
American Association and the World Uighur Congress, both of which receive financial support from the
National Endowment for Democracy, mobilized exile groups around the world to protest an episode in
Guangdong Province in late June.

Chinese officials say that two Uighur workers were killed by a small group of Han Chinese, who have
been detained; Ms.~Kadeer says, with evident sincerity, there is evidence that a mob killed up to 60
Uighurs while the police did nothing. But her version has not been independently verified, and
Chinese authorities accuse exiles of exaggerating the matter to incite anti-Chinese feelings.

The world congress, based in Munich, has just one paid staff member but is in touch with some 51
exile groups around the world. Ms.~Kadeer said that by June 30 she had called all of those groups to
encourage demonstrations outside Chinese embassies.

The rumors about mass killings in Guangdong were one trigger for Sunday's protests, but Ms.~Kadeer
challenged Beijing authorities to release the transcript of a call she made to a brother in Urumqi
on Saturday in which, she said, she urged him not to become involved in any demonstrations.

``Instead of blaming me, the Chinese government should start listening to the complaints of the
Uighur people and choose dialogue,'' Ms.~Kadeer said.

Her fame and force of personality have given the Uighurs a huge lift, but some exiles wonder about
her domination and future leadership.

``I've been looking for someone like me who can take over,'' she said on Tuesday. For now, she said,
``The people will not let me stop because my goal is their liberation.''

``Until I lose my consciousness, I'll stay on as the leader.''

\section{At a Factory, the Spark for China's Violence}

\lettrine{T}{he} first batch of Uighurs, 40 young men and women from the far western region of
Xinjiang, arrived at the Early Light Toy Factory here in May, bringing their buoyant\cite{buoyant}
music and speaking a language that was incomprehensible to their fellow Han Chinese workers.

``We exchanged cigarettes and smiled at one another, but we couldn't really communicate,'' said Gu
Yunku, a 29-year-old Han assembly line worker who had come to this southeastern city from northern
China. ``Still, they seemed shy and kind. There was something romantic about them.''

The mutual good will was fleeting.

By June, as the Uighur contingent rose to 800, all recruited from an impoverished rural county not
far from China's border with Tajikistan, disparaging chatter began to circulate. Taxi drivers traded
stories about the wild gazes\cite{gaze} and gruff\cite{gruff} manners of the Uighurs. Store owners
claimed that Uighur women were prone\cite{prone} to shoplifting. More ominously\cite{ominous}, tales
of sexually aggressive Uighur men began to spread among the factory's 16,000 Han workers.

Shortly before midnight on June 25, a few days after an anonymous Internet posting claimed that six
Uighur men had raped two Han women, the suspicions boiled over into bloodshed.

During a four-hour melee\cite{melee} in a walkway between factory dormitories, Han and Uighur
workers bludgeoned one another with fire extinguishers, paving stones and lengths of steel shorn
from bed frames.

By dawn, when the police finally intervened, two Uighur men had been fatally wounded and 120 other
people were injured, most of them Uighurs, according to the authorities.

``People were so vicious\cite{vicious}, they just kept beating the dead bodies,'' said one man who
witnessed the fighting, which he said involved more than a thousand workers.

Ten days later and 1,800 miles away, the clash in Shaoguan provoked a far greater spasm\cite{spasm}
of violence in Urumqi, the capital of the Xinjiang region. On July 5, a demonstration by Uighur
students protesting what they said was a lackluster investigation of the factory brawl gave way to a
murderous rampage\cite{rampage} against the city's Han residents, followed by killings carried out
by the Han.

In the end, at least 192 people died and more than 1,000 were injured, according to the government.
Of the dead, two-thirds were Han, the authorities said. Uighurs insist that the toll among their own
was far higher.

Shaoguan officials, who said that the rape allegations were untrue, contended that the violence at
the toy factory was used by ``outsiders'' to fan ethnic hatred and promote Xinjiang separatism.
``The issue between Han and Uighur people is like an issue between husband and wife,'' Chen Qihua,
vice director of the Shaoguan Foreign Affairs Office, said in an interview. ``We have our quarrels,
but in the end, we are like one family.''

Li Qiang, the executive director of China Labor Watch, an advocacy group based in New York that has
studied the Shaoguan toy factory, has a different view. He said the stress of low pay, long hours
and numbingly repetitive work exacerbated deeply held mistrust between the Han and the Muslim
Uighurs, a Turkic-speaking minority that has long resented Chinese rule.

``The government doesn't really understand these ethnic problems, and they certainly don't know how
to resolve them,'' Mr.~Li said.

In the government's version of events, the factory clash was the simple product of false rumors,
posted on the Internet by a disgruntled former worker who has since been arrested.

A few days later, the authorities added another wrinkle to the story, saying that the fight was
prompted by a ``misunderstanding'' after a 19-year-old female worker accidentally stumbled into a
dormitory room of Uighur men.

The woman, Huang Cuilian, told the state news media that she screamed and ran off when the men
stamped their feet in a threatening manner. When Ms.~Huang, accompanied by factory guards, returned
to confront the men, the standoff quickly escalated.

The Uighur workers have since been sequestered at an industrial park not far from the toy factory.
Officials refused to allow a reporter access to the workers, and a large contingent of police
officers blocked the hospital rooms where two dozen others were recovering from their wounds.

``They want to lead a peaceful life and not be bothered by the media,'' said Mr.~Chen, the Shaoguan
official. He said the government of Guangdong Province, where Shaoguan is located, and the factory
would provide them employment at a separate plant.

Officials at Early Light, a Hong Kong company that is the largest toy maker in the world, declined
to comment.

In the city of Kashgar, the ancient heart of Uighur civilization, the Shaoguan killings have
inflamed longstanding anger over the way China manages daily life in Xinjiang. Many Uighurs complain
about policies that encourage Han migration to the region and say the government suppresses Uighurs'
language and religion. When it comes to employment, they say coveted state jobs go to the Han; a
2008 report by a United States Congressional commission noted that government job Web sites in
Xinjiang set aside most teaching and civil service positions for non-Uighurs.

``If we weren't so poor, our children wouldn't have to take work so far from home,'' said Akhdar, a
67-year-old man who, like many others interviewed, refused to give his full name for fear of
reprisals from the authorities.

The Uighurs who work at the Shaoguan toy factory, all of them from Shufu County outside Kashgar, are
part of a growing wave of 1.5 million people who have migrated from Xinjiang to more prosperous
cities of coastal China. This year, more than 6,700 young men and women left Shufu County, according
to government figures, part of an ambitious jobs export program intended to relieve high youth
unemployment and provide low-cost workers to factories.

According to an article in the state-run Xinjiang Daily, ``70 percent of the laborers had signed up
for employment voluntarily.'' The article, published in May, did not explain what measures were used
to win over the remaining 30 percent.

But residents in and around Kashgar say the families of those who refuse to go are threatened with
fines that can equal up to six months of a villager's income. ``If asked, most people will go,
because no one can afford the penalty,'' said a man who gave his name only as Abdul, whose
18-year-old sister is being recruited for work at a factory in Guangzhou but has so far resisted.

Some families are particularly upset that recruitment drives are directed at young unmarried women,
saying that the time spent living in a Han city far away from home taints their marriage prospects.
Taheer, a 25-year-old bachelor who is seeking a wife, put it bluntly. ``I would not marry such a
girl because there's a chance she would not come back with her virginity,'' he said.

Still, a few Uighurs said they were thankful for factory jobs with wages as high as \$190 a month,
double the average income in Xinjiang. One man, a 54-year-old cotton farmer with two young
daughters, said he was ready to send them away if that was what the Communist Party wanted. ``We
would be happy to oblige,'' he said with a smile as his wife looked away.

Once they arrive in one of China's bustling manufacturing hubs, the Uighurs often find life
alienating. Mr.~Li of China Labor Watch said many workers were unprepared for the grueling work, the
cramped living conditions and what he described as verbal abuse from factory managers.

But the biggest challenge may be open hostility from Han co-workers, who like many Chinese hold
unapologetically negative views of Uighurs.

Many Han say they believe that Uighurs are given unfair advantages by the central government,
including a point system that gives Uighur students and other minorities a leg up on college
entrance exams.

Zhang Qiang, a 20-year-old Shaoguan resident, described Uighurs as ``barbarians'' and said they were
easily provoked to violence.

``All the men carry knives,'' he said after dropping off a job application at the toy factory, which
is eager to hire replacements for the hundreds of workers who quit in recent weeks.

Still, Mr.~Zhang acknowledged that his contact with Uighurs was superficial. When he was a student,
his vocational high school had a program for 100 Xinjiang students, although they were relegated to
separate classrooms and dorms.

If he had any curiosity about his Uighur classmates, it was quashed by a teacher who warned the Han
students to keep their distance. ``This is not prejudice,'' he said. ``It is just the nature of
their kind.''

\section{Senate Likely to Vote on Sotomayor in August}

\lettrine{J}{udge} Sonia Sotomayor completed her testimony before the Senate Judiciary Committee on
Thursday, acknowledging regret for having said that a ``wise Latina'' judge could decide a case
better than a white man while defending her role in a case involving New Haven firefighters in which
she was reversed by the Supreme Court.

But over her four days in the witness chair, Judge Sotomayor provided Republicans little
ammunition\cite{ammunition} with which to block the Senate from approving her elevation to the
Supreme Court, where she would become the nation's first Hispanic justice.

Senior Republican staff aides said in interviews they expected that at least one and perhaps as many
as three of the panel's seven Republicans might vote to approve the Sotomayor nomination and send it
to the full Senate, which is expected to confirm her in the first week in August.

Senator Patrick J.~Leahy, Democrat of Vermont and the committee chairman, has scheduled a vote for
next Tuesday, but Republicans on the panel have indicated they will ask the vote be delayed a week
until July 28.

The aides said they expected Senator Lindsey Graham of South Carolina, whose questioning of Judge
Sotomayor veer\cite{veer}ed between folksy\cite{folksy} support and wariness, to vote to approve the
nomination. They said he could be joined by Senators Orrin G.~Hatch of Utah, Charles E.~Grassley of
Iowa, John Cornyn of Texas or Tom Coburn of Oklahoma. Senators Hatch and Grassley, committee
veterans\cite{veteran}, have generally supported Democratic judicial nominees in the past. Senator
Coburn said he was ``mighty impressed'' with Judge Sotomayor after questioning her.

In 2005, John G.~Roberts Jr., who had been nominated to be chief justice by a Republican president,
George W.~Bush, received three votes from among the eight Democrats on the panel.

One of the aides, all of whom spoke on the condition of anonymity, said Republicans had for some
time given up any hope of derailing the Sotomayor nomination. But the aide said that committee
Republicans had taken satisfaction in two of Judge Sotomayor's concessions at the hearings, which
ended in the evening with panels of outside witnesses.

Not only did she back away from and express regret for her ``wise Latina'' comments, but the
official said Republicans were also pleased that she seemed to repudiate President Obama's
formulation that a judge needed to have empathy for those who came before the court.

Before she got up from the witness chair for good, Judge Sotomayor, who sits on the federal appeals
court in New York City, was confronted repeatedly by Republican questions about the apparent gulf
between her testimony and her past speeches in which she spoke about how personal experiences might
influence a judge's behavior. ``Look at my record over the last 17 years as a judge,'' she responded
regularly.

In eschewing\cite{eschew} her speeches, Judge Sotomayor gave answers that were reminiscent of
Clarence Thomas, who was challenged at his 1991 confirmation hearings about several speeches in
which he said that ``natural law,'' an unwritten moral code, sometimes trumped written law. Justice
Thomas said his speeches were philosophical musings, an answer Democrats dismissed but Republicans
accepted.

Senator Jeff Sessions of Alabama, the ranking Republican on the committee, remained Judge
Sotomayor's most relentless inquisitor\cite{inquisitor}. In one pointed colloquy\cite{colloquy},
Mr.~Sessions challenged her about being part of a three-judge panel that ruled against white and
Hispanic New Haven firefighters after the city threw out civil service exam results because not
enough members of minorities had achieved scores high enough for promotion.

Mr.~Sessions, who was denied a federal district court judgeship 23 years ago over accusations of
racial insensitivity, tried to turn the tables on Judge Sotomayor. ``Did you fail to show the
courage that the attorney general has asked us to show about race?'' he asked, referring to comments
by Attorney General Eric H.~Holder Jr.

``Sir, no, I did not show a lack of courage,'' she replied. ``It was a thorough, complete discussion
of the issues.''

Two of the firefighters who lost the case before Judge Sotomayor only to achieve victory last month
in a 5 to 4 Supreme Court ruling testified against her nomination in the afternoon.

The panel also heard from David Cone, a former major-league pitcher, who said Judge Sotomayor
``rescued baseball'' with her ruling in 1995 in which she barred team owners from unilaterally doing
away with the free agent and salary arbitration systems. In discussing what is perhaps Judge
Sotomayor's most celebrated case, Mr.~Cone said, ``All of us players, owners and fans are in her
debt.''

The National Rifle Association, after days of critical comments about Judge Sotomayor, formally
opposed her nomination on Thursday. The group said it did so because of concern over her commitment
to the Second Amendment, citing a case in which she took part involving the banning of a martial
arts weapon in New York State.

Officials of the gun group said they would not include a senator's vote on the nomination to rate
lawmakers on their fidelity to a pro-gun-rights agenda.

\section{Future Nominations Are at Stake in Hearings}

\lettrine{A}{s} the two parties skirmish over the Supreme Court nomination of Judge Sonia
Sotomayor, some of their rhetorical fire is aimed not at her but at the next justice President Obama
may get to pick.

After three days of testimony, Judge Sotomayor appeared to have made no major mistakes that would
jeopardize her confirmation in a Senate dominated by Democrats. So both sides are trying to use the
Judiciary Committee hearings to define the parameters of an acceptable nomination in case another
seat opens up during Mr.~Obama's presidency.

By forcing Judge Sotomayor to retreat from Mr.~Obama's desire for justices with ``empathy,''
Republicans have effectively set a new standard that future nominees will be pressed to meet. The
Republicans hope their aggressive questioning of Judge Sotomayor on race discrimination, gun control
and the death penalty will make it harder for Mr.~Obama to choose a more outspoken liberal in the
future.

Liberal activists, by contrast, hope the hearings demonstrate that a Democratic president has
nothing to fear from Republicans who have not rattled Judge Sotomayor. If she is confirmed by a
commanding vote that includes a number of Republicans, the activists argue, they will have given
Mr.~Obama more political running room next time to name a more full-throated champion of liberal
values.

``Democrats want to draw lines in one place in the sand, and Republicans want to draw lines in
another place in the sand,'' said Rachel Brand, a former assistant attorney general who helped
manage the confirmations of Chief Justice John G.~Roberts Jr.~and Justice Samuel A.~Alito Jr.

For all sorts of reasons, Mr.~Obama's next nomination fight could easily be a bigger confrontation
than this one. At the moment, the president still enjoys strong popular support and is replacing a
retiring liberal, Justice David H.~Souter, with another liberal. But if his next choice alters the
political balance or comes when his poll numbers are down, the battle could become much fiercer.

``It depends on who's left the court,'' said Fred McClure, a former White House official who worked
on the confirmations of three Republican-appointed justices. ``If the departing justice falls into
the middle-to-liberal wing of the court, then I think the president is free to nominate a similar
candidate. If on the other hand, the nominee is to replace the middle-to-right wing of the court,
then I think he's got to think differently.''

One lesson the White House has learned is to expect a fight on any nomination, fueling its desire
for a strong vote for Judge Sotomayor. White House officials say they hope to generate momentum for
the next one, though cautioning that each nomination is unique.

Several legal experts said Judge Sotomayor's testimony might make it harder for Mr.~Obama to name a
more liberal justice next time.

She repudiated the president's assertion that ``what is in a judge's heart'' should influence
rulings and rejected the liberal idea that the Constitution is a ``living'' document whose meaning
evolves with society. Instead, she said the Constitution was ``immutable'' and did not change except
by amendment. And she dismissed any role for foreign law in deciding cases, an influence some
liberal legal experts argue should be considered.

Louis Michael Seidman, a Georgetown University constitutional law professor, said Judge Sotomayor
adopted a ``fairy tale'' definition of judging that ignores the discretion\cite{discretion} they
have to decide hard cases where the legal materials do not dictate outcomes.

``She reinforced the official ideology, and it makes it that much harder for other judges later on
to talk to the American people as if they were adults about what courts actually do and what
constitutional law consists of,'' Mr.~Seidman said.

James R.~Copland, director of the Center for Legal Policy at the Manhattan Institute, said it was
significant that Judge Sotomayor had endorsed a ``traditionalist'' view of judging and distanced
herself from liberal theories. ``That limits somewhat his ability to go far afield from the
mainstream public understanding of the law,'' Mr.~Copland said of the president.

Manuel Miranda, chairman of the Third Branch Conference, a conservative advocacy group, said Judge
Sotomayor could make it easier for conservatives to attack future nominees. ``She is portraying
herself as someone who is bound by the rules that conservatives have been articulating for so many
years,'' Mr.~Miranda said. As a result, if Mr.~Obama picks someone with different views, Republicans
``now can say, 'You don't meet the Sotomayor test.' ''

But Nan Aron, president of the liberal Alliance for Justice, said that if Judge Sotomayor ended up
being confirmed by a significant margin, it would represent a ``strong vote of confidence in
President Obama's judicial picks in the future'' at both the appeals court and Supreme Court levels.

``Everyone has been saying get a good vote, and this administration will have a firm foundation and
a good record to move forward on pushing back against the ultraconservative hold that currently
exists on this court,'' Ms.~Aron said.

Either way, each side has taken the measure of the other and will emerge from this nomination with
ammunition for the next.

``There were some rigorous questions, but it wasn't Armageddon,'' said Kenneth M.~Duberstein, who
was Ronald Reagan's White House chief of staff and later managed the confirmations of Justices
Souter and Clarence Thomas. ``That's left for another day and another nomination.''

\section{Chinese Question Police Absence in Ethnic Riots}

\lettrine{A}{s} this shattered regional capital sorts through the corpses from China's deadliest
civil unrest in decades, another loss has become apparent: faith in the government's ability to
secure the peace and quell mass disturbances. In many neighborhoods, police officers remained absent
for hours as the carnage unfolded, witnesses say.

The bloodletting here on July 5, in which ethnic Uighurs pummeled\cite{pummel} and stabbed ethnic
Han to death, was just the latest episode in a nationwide upswing in large-scale street violence
that had already prompted concerned officials in Beijing to look for new ways to defuse such
outbursts. In all of the recent cases, not only were officials and security forces unable to contain
the violence, but average people clashed\cite{clash} with the police en masse -- a sign of the
profound distrust of local authority throughout much of China.

``In the last several years, the level of violence and speed with which these incidents can turn
violent has increased,'' said Murray Scot Tanner, an analyst of Chinese security. ``It raises a
very, very serious question: To what extent are the Chinese people afraid of their police anymore?''

In parts of the Uighur quarter and in poorer, mixed areas of south Urumqi, young Uighur men with
sticks, knives and stones went on a bloody rampage for about five hours while police officers
remained mostly absent, according to interviews with dozens of residents. In some areas where police
officers arrived but were outnumbered by rioters, the officers stood around or fled, witnesses said.

``Where were the police while people were being killed?'' said Cheng Wei, 41, a landscaper whose
neighbors, poor fruit vendors from Henan Province, lost a son in the riots. ``They were completely
useless.''

Large street protests that turn violent, and that officials and security forces have been powerless
to stop, have been on the rise in recent years, analysts say. The government usually avoids
reporting the number of protests or riots in China, but an article in January in Outlook Weekly, a
policy magazine published by Xinhua, the state news agency, said there were 90,000 such events in
2006, up from 60,000 in 2003.

The central government still can completely lock down areas when it anticipates protests, as it did
across the Tibetan plateau\cite{plateau} in the spring or for the 20th anniversary of the student
rallies at Tiananmen Square in June. But increasingly, security forces seem to have been caught
unaware.

The rampage\cite{rampage} by Uighurs on July 5 was followed for days by reprisal\cite{reprisal}
killings by Han vigilantes\cite{vigilante} who defied\cite{defy} police orders to refrain from
violence. At least 192 people were killed and 1,721 injured in all of the violence, most of them
Han, according to the government. Many Uighurs say the Uighur casualties\cite{casualty} have been
severely undercounted. The Han, who dominate China, are the majority in Urumqi, even though the
Uighurs, a Turkic people largely resentful of Chinese rule, are the biggest ethnic group in this
western region of Xinjiang.

In March 2008, rioters in Tibet openly defied police officers who, caught by surprise, largely
disappeared during the first 24 hours of violence. At least 19 people died.

Last month, tens of thousands of residents of Shishou, in Hubei Province, clashed with riot police
officers over the mysterious death of a hotel chef. A year earlier, in Weng'an County of Guizhou
Province, at least 30,000 people rioted over the handling of an inquiry into the death of a
17-year-old girl, torching police cars, the main police station and the government headquarters.

Frustration at legal injustice and Communist Party corruption is a common thread. The violence in
Xinjiang began as a peaceful protest on July 5, when Uighurs called for a proper inquiry into a
factory brawl\cite{brawl} in southern China that had left two Uighurs dead.

``The absence of an independent legal system is the party's biggest mistake, because when people
can't take their grievances\cite{grievance} to the courts, they take them to the streets,'' said
Nicholas Bequelin, an Asia researcher for Human Rights Watch.

So concerned are Chinese leaders over the rise in mass violence and the growing contempt for law
enforcement that they have taken new measures to ensure stability, with the 60th anniversary of the
establishment of the People's Republic coming up in October.

Vice President Xi Jinping, pegged\cite{peg} as the next leader of China, took charge of a committee
to ensure social stability. Separately, party officials and police officers down to the county level
have taken part in training for managing civil unrest. The drills include teaching them to disable
local Internet service during an outbreak and emphasizing that leaders take part in dialogue at the
front lines rather than resort to shows of force.

But party leaders and police officers in Urumqi failed to avert\cite{avert} disaster the night of
July 5 even though government officials say the police knew as early as 1 a.m. that day that Uighurs
were planning to hold a protest.

In the early evening of July 5, galvanized\cite{galvanize} by Internet messages, Uighurs began
gathering at People's Square in the city center, near the headquarters of the regional Communist
Party and government offices, to protest the handling of the earlier factory brawl. Police officers
quickly encircled the crowd, witnesses said.

A mile south, about 6 p.m., people also began gathering on the northern edge of the old Uighur
quarter, said Adam Grode, an American teacher who watched the scene from his 16th-floor apartment.
The crowd swelled to more than 1,000 people, including women and the elderly.

There were at first only a few traffic police officers standing around. But by 6:30 p.m., a line of
troops from the People's Armed Police, a paramilitary force, had formed to the north and was trying
to push the crowd down into the Uighur quarter. Some officers charged with batons\cite{baton}. The
crowd surged back against the troops, fists raised.

Another wave of troops arrived. They were better equipped, with body padding and riot shields,
Mr.~Grode said. Some had rifles slung across their bodies.

Young men began hurling stones and bricks as the police attacked with batons. People also threw
rocks at buses that had been halted. A full-fledged street battle erupted, though the police
officers at this point did not use their guns, Mr.~Grode said.

Just a few hundred yards south, at the busy Grand Bazaar area, there were few officers. The handful
there just stood by as rioters set upon any Han civilians they saw, witnesses said. One taxi driver,
who gave his name as Mr.~Han, said he was dragged from his car by Uighurs with knives while
policemen watched. He managed to escape.

After 8 p.m., rioters showed up in mixed neighborhoods about two miles southeast of the Uighur
quarter. Police officers did not arrive until after 1 a.m., witnesses said. These areas were among
the worst hit; witnesses said bodies were strewn all around Dawan North Road, for instance.

``The police arrived around 1:30 a.m., and they put down their riot shields to move bodies,'' said
Mr.~Cheng, the landscaper.

Earlier, at twilight\cite{twilight}, back in the northern half of the Uighur quarter, officers
sprinted through alleyways to beat down and handcuff Uighur men. By around 10 p.m., they had begun
opening fire with guns and tear gas rifles, Mr.~Grode said, adding that he heard occasional series
of single-shot gunfire. Another foreigner also said she heard gunfire after dark.

By 1 a.m., the rioting had ebbed, and police officers in the Uighur quarter were putting scores of
handcuffed men onto buses.

Han residents keep asking why security forces showed up so late in the southern neighborhoods, where
Han live close to Uighurs and are clearly vulnerable. Mr.~Tanner, the security analyst, said that 11
years after the Tiananmen Square protests, security forces were ordered to handle protests
cautiously, but that if rioting broke out, officers and paramilitary troops could use ``decisive
force'' as long as senior local officials had given approval. They are not supposed to let a riot
run its course, he said.

But security forces also make securing government buildings, financial centers and other strategic
points a top priority, Mr.~Tanner said. Indeed, a local reporter wrote that he saw many police
officers after 8 p.m. on Zhongshan Road, where government buildings are. This could help explain why
officers did not show up in the residential areas until much later.

At the most basic level, though, the policing failure appears rooted in the government's inability
to understand the Uighur-Han relationship. ``There's a severe failure of intelligence about society
and about social tensions,'' Mr.~Tanner said. ``In this case, what I think they were clearly
unprepared for is the level of organized intercommunal violence.''

Two days after the killings by the Uighurs, thousands of Han with sticks and knives clashed with
police officers as the Han tried storming the Uighur quarter. None of them trusted the government to
mete out proper punishment or to protect the Han.

A man who gave his name as Mr.~Li, waving a wooden chair leg, said, ``I'm here to safeguard
justice.''

\section{Why Japan's Cellphones Haven't Gone Global}

\lettrine{A}{t} first glance, Japanese cellphones are a gadget lover's dream: ready for Internet and
e-mail, they double as credit cards, boarding passes and even body-fat calculators.

But it is hard to find anyone in Chicago or London using a Japanese phone like a Panasonic, a Sharp
or an NEC. Despite years of dabbling\cite{dabble} in overseas markets, Japan's handset makers have
little presence beyond the country's shores.

``Japan is years ahead in any innovation. But it hasn't been able to get business out of it,'' said
Gerhard Fasol, president of the Tokyo-based IT consulting firm, Eurotechnology Japan.

The Japanese have a name for their problem: Gal\'apagos syndrome.

Japan's cellphones are like the endemic\cite{endemic} species that Darwin encountered on the
Gal\'apagos Islands -- fantastically evolved and divergent\cite{divergent} from their mainland
cousins -- explains Takeshi Natsuno, who teaches at Tokyo's Keio University.

This year, Mr.~Natsuno, who developed a popular wireless Internet service called i-Mode, assembled
some of the best minds in the field to debate how Japanese cellphones can go global.

``The most amazing thing about Japan is that even the average person out there will have a
superadvanced phone,'' said Mr.~Natsuno. ``So we're asking, can't Japan build on that advantage?''

The only Japanese handset maker with any meaningful global share is Sony Ericsson, and that company
is a London-based joint venture between a Japanese electronics maker and a Swedish
telecommunications firm.

And Sony Ericsson has been hit by big losses. Its market share was just 6.3 percent in the first
quarter of 2009, behind Nokia of Finland, Samsung Electronics and LG of South Korea, and Motorola of
Illinois.

Yet Japan's lack of global clout is all the more surprising because its cellphones set the pace in
almost every industry innovation: e-mail capabilities in 1999, camera phones in 2000,
third-generation networks in 2001, full music downloads in 2002, electronic payments in 2004 and
digital TV in 2005.

Japan has 100 million users of advanced third-generation smartphones, twice the number used in the
United States, a much larger market. Many Japanese rely on their phones, not a PC, for Internet
access.

Indeed, Japanese makers thought they had positioned themselves to dominate the age of digital data.
But Japanese cellphone makers were a little too clever. The industry turned increasingly inward. In
the 1990s, they set a standard for the second-generation network that was rejected everywhere else.
Carriers created fenced-in Web services, like i-Mode. Those mobile Web universes
fostered\cite{foster} huge e-commerce and content markets within Japan, but they have also increased
the country's isolation from the global market.

Then Japan quickly adopted a third-generation standard in 2001. The rest of the world dallied,
essentially making Japanese phones too advanced for most markets.

At the same time, the rapid growth of Japan's cellphone market in the late 1990s and early 2000s
gave Japanese companies little incentive to market overseas. But now the market is shrinking
significantly, hit by a recession and a graying economy; makers shipped 19 percent fewer handsets in
2008 and expect to ship even fewer in 2009. The industry remains fragmented, with eight cellphone
makers vying for part of a market that will be less than 30 million units this year.

Several Japanese companies are now considering a push into overseas markets, including NEC, which
pulled the plug on its money-losing international cellphone efforts in 2006. Panasonic, Sharp,
Toshiba and Fujitsu are said to be planning similar moves.

``Japanese cellphone makers need to either look overseas, or exit the business,'' said Kenshi
Tazaki, a managing vice president at the consulting firm Gartner Japan.

At a recent meeting of Mr.~Natsuno's group, 20 men and one woman crowded around a big conference
table in a skyscraper in central Tokyo, examining market data, delivering diatribes and frequently
shaking their heads.

The discussion then turned to the cellphones themselves. Despite their advanced hardware, handsets
here often have primitive, clunky interfaces, some participants said. Most handsets have no way to
easily synchronize data with PCs as the iPhone and other smartphones do.

Because each handset model is designed with a customized user interface, development is
time-consuming and expensive, said Tetsuzo Matsumoto, senior executive vice president at Softbank
Mobile, a leading carrier. ``Japan's phones are all 'handmade' from scratch,'' he said. ``That's
reaching the limit.''

Then there are the peculiarities\cite{peculiarity} of the Japanese market, like the almost universal
clamshell design, which is not as popular overseas. Recent hardware innovations, like solar-powered
batteries or waterproofing, have been incremental rather than groundbreaking.

The emphasis on hardware makes even the newest phones here surprisingly bulky. Some analysts say
cellphone carriers stifle innovation by demanding so many peripheral hardware functions for phones.

The Sharp 912SH for Softbank, for example, comes with an LCD screen that swivels 90 degrees, GPS
tracking, a bar-code reader, digital TV, credit card functions, video conferencing and a camera and
is unlocked by face recognition.

Meanwhile, Japanese developers are jealous of the runaway global popularity of the Apple iPhone and
App Store, which have pushed the American and European cellphone industry away from its obsession
with hardware specifications to software. ``This is the kind of phone I wanted to make,''
Mr.~Natsuno said, playing with his own iPhone 3G.

The conflict between Japan's advanced hardware and its primitive software has contributed to some
confusion over whether the Japanese find the iPhone cutting edge or boring. One analyst said they
just aren't used to handsets that connect to a computer.

The forum Mr.~Natsuno convened to address Gal\'apagos syndrome has come up with a series of
recommendations: Japan's handset makers must focus more on software and must be more aggressive in
hiring foreign talent, and the country's cellphone carriers must also set their sights overseas.

``It's not too late for Japan's cellphone industry to look overseas,'' said Tetsuro Tsusaka, a
telecom analyst at Barclays Capital Japan. ``Besides, most phones outside the Gal\'apagos are just
so basic.''

\section{Wikipedia May Be a Font of Facts, but It's a Desert for Photos}

\lettrine{S}{ome} may wonder: Could there be a bad picture of Halle Berry or George Clooney?

Just visit Wikipedia. There you'll find a fuzzy shot of Ms.~Berry from the mid-1980s, when she was
part of a U.S.O. tour with other Miss USA contestants. She is out of focus and wearing a
red-and-white baseball cap -- in short, she is barely recognizable. Mr.~Clooney, in his Wikipedia
entry, is shown in Chad wearing a khaki vest and a United Nations cap. Smiling, he is ruggedly
handsome in the company of two women who work for the United Nations; still this is hardly a glamour
shot.

Then there are big names like Howard Stern or Julius Erving who have no photograph at all on
Wikipedia.

At a time when celebrities typically employ a team of professionals to control their images,
Wikipedia is a place where chaos rules. Few high-quality photographs, particularly of celebrities,
make it onto this site. This is because the site runs only pictures with the most permissive
Creative Commons license, which allows anyone to use an image, for commercial purposes or not, as
long as the photographer is credited.

``Representatives or publicists will contact us'' horrified at the photographs on the site, said Jay
Walsh, a spokesman for the Wikimedia Foundation, which operates the Wikipedia encyclopedias in more
than 200 languages. ``They will say: 'I have this image. I want you to use this image.' But it is
not as simple as uploading a picture that is e-mailed to us.''

``In general,'' he added, ``we need them to know that giving us a photograph from Annie Leibovitz
won't work unless Annie Leibovitz is O.K. with it.''

Photographs are a glaring flaw in the Wikipedia model. Unlike the articles on the site, which in
theory are improved, fact checked, footnoted and generally enhanced over time, photographs are
static works created by individuals. A bad article can become a better article. A bad photograph
simply stays bad.

Wikipedians have tried to make up for this defect by organizing outings where groups of contributors
take high-quality photographs of buildings or objects. Likewise Wikipedia has tried to gain
permission from large photographic collections to use their material.

Last winter the German Federal Archives released 100,000 low-resolution digital copies under a
license so they could appear on Wikipedia. Recently a Wikipedia user, Derrick Coetzee, downloaded
more than 3,000 high-resolution photographs from the British National Portrait Gallery -- to serve,
in essence, as the head shots for important historical figures like Charlotte Brontë or Charles
Darwin.

The gallery threatened legal action against Mr.~Coetzee, saying that while the painted portraits may
be old and thus beyond copyright protection, the photographs are new and therefore copyrighted work.
The gallery is demanding a response by Monday from Mr.~Coetzee, who is being represented by the
Electronic Frontier Foundation. In an e-mail message on Friday a gallery spokeswoman, Eleanor
Macnair, wrote that ``contact has now been made'' with the Wikimedia Foundation and ``we remain
hopeful that a dialogue will be possible.''

But none of this has made much of an improvement in Wikipedia's photography. Any gallery of hideous
Wikipedia photographs would include the former N.B.A. star George Gervin, who is standing stiffly in
a suit in a shot that is cropped longer and thinner than would be typical even for a basketball
player. The unrestricted photograph came from the office of Senator John Cornyn of Texas, who has
been cut out of it.

As in Mr.~Gervin's case, the government is a prime source for public domain photographs. President
Obama, for example, looks composed and serious in the official portrait that sits on the upper
right-hand corner of his article.

But Wikipedia contributors also cull government collections for photographs of celebrities' meetings
with politicians, hoping to find something to post on the page.

The former home-run king Hank Aaron is shown in an out-of-context, oddly cropped photograph from a
1978 visit to the White House. Likewise the main photograph of Michael Jackson was taken from his
1984 visit with Ronald and Nancy Reagan.

Recent photographs on Wikipedia almost exclusively are the work of amateurs who don't mind giving
away their work. Amateur may be too kind a word; their photos tend to be the work of fans who happen
to have a camera. The opera singer Natalie Dessay is shown looking the wrong way at an autograph
signing; the actress Allison Janney appears in sunglasses at the Toronto Film Festival. The Coen
brothers, Joel and Ethan, are seen from middle distance at Cannes in 2001, with Ethan covering his
mouth, perhaps because he has just coughed.

Then there are the photographs taken from the stands, with the subject barely a fleck. Barry Bonds
is apparently the outfielder in the center of one photograph on his page; David Beckham can be
discerned with his hands on hips during a 1999 soccer match.

A few celebrities, like Pl\'acido Domingo and Oliver Stone, have had the foresight to provide their
own freely licensed photographs. And considering the money that stars spend to maintain their image,
it is surprising that more have not invested in high-quality, freely licensed photographs for
Wikipedia and other sites. Perhaps they don't recognize how popular Wikipedia is. In June, for
example, Ms.~Berry's article had more than 180,000 page views.

Also, it can be difficult to persuade a talent photographer to go along with that approach because
one free photograph can drive out all the others, said Jerry Avenaim, a celebrity photographer. He
is unusual in that he has contributed about a dozen low-resolution photographs to Wikipedia,
including a shot of the actor Mark Harmon, originally created for TV Guide.

In an interview Mr.~Avenaim still sounded torn about the idea of contributing his work. He said he
was trying to accomplish two goals: ``One, I really wanted to help the celebrities that I care about
to show them in the proper light they want to be shown,'' he said. ``Second, it is an interesting
marketing strategy for myself.''

He said that having his work on Wikipedia has increased his online visibility as reflected in
search-engine results and traffic to his Web site,, but that the costs are potentially high. ``This
is the lifeblood of my career,'' he said, noting that photographers may get paid very little for a
celebrity shot for a magazine. They make their money from resales of the image. And even a
low-resolution photograph that is available free -- say, his shot of Dr.~Phil -- becomes the default
photograph online and means there is no need to pay for another one of his shots.

That, ultimately, is the issue for photographers who might want to donate their work to Wikipedia,
but not the entire Internet.

``To me the problem is the Wikipedia rule of public use,'' Mr.~Avenaim said. ``If they truly wanted
to elevate the image on the site, they should allow photographers to maintain the copyright.''

\section{Across U.S., ESPN Aims to Be the Home Team}

\lettrine{N}{ot} content with being a sports colossus with broadcasts in 200 countries, ESPN is
taking aim at hometown sports coverage, threatening one of the last strongholds of local newspapers
and television stations.

ESPN has long dominated the coverage of national athletics, pumping out news and commentary on every
major sport (and some not-so-major ones) via an expanding network of cable channels, Web sites and
mobile services.

Now, after a promising test run in Chicago, ESPN is adding local offshoots to three more cities. On
Monday, ESPN, which is owned by the Walt Disney Company, plans to announce local Web sites in New
York, Los Angeles and Dallas -- in what executives say is only the ``first inning'' of their effort
to provide hyperlocal sports coverage in cities across the country.

``We look at this as a perfect example of offering a better product,'' said George W.~Bodenheimer,
co-chairman of Disney Media Networks.

Chicago news outlets, which have some experience competing directly with ESPN, acknowledge that the
threat is real.

``We are taking ESPN's marketing push seriously, and we are looking forward to the local sports turf
battle in the weeks and months ahead,'' said Bill Adee, The Chicago Tribune's editor for digital
media.

In less than three months, ESPN Chicago has become the city's top sports site, attracting about
590,000 unique visitors in June, according to data from comScore, an Internet measurement company.
Second place went to The Tribune's online sports section with 455,000 unique visitors.

ESPN Chicago does not seem to have cut into The Tribune's online sports audience as much as it has
slowed its growth, according to a review of the traffic data.

At The Los Angeles Times, which is about to face ESPN head-on, the associate editor, Randy Harvey,
said: ``It would be foolish to underestimate ESPN, but it comes down to resources. I don't see them
being able to replicate what we do.''

Once ESPN establishes itself in local markets, it plans to move deeper into local sports -- down to
the high school level and perhaps beyond -- by using social networking and other technology to
inform its journalism.

Just as CNN offers news from everyday people through its iReport service, ESPN could augment its
local coverage by tapping sports fans. Picture fathers uploading scores from their daughters' soccer
games.

``This initiative is perfect for exploring that opportunity,'' said Mr.~Bodenheimer, who also serves
as president of ESPN and ABC Sports.

The company would say only that its costs are ``minimal,'' using existing resources to expand; in
fact, ESPN says it is hiring only about 15 new people for the sites in Dallas, Los Angeles and New
York. The Dallas site will start in the fall, the other two early next year.

These moves have challenges, even for a company with pockets as deep as ESPN. Sports fans, loyal to
their local newspaper columnists, could balk at the continued march of ESPN across the sports news
landscape, and the media giant could easily run afoul of local quirks. Costs also may be difficult
to control, and maintaining editorial standards with content supplied by fans might be difficult.

Also, Fox Sports has successful cable networks in several regional markets, but its Web presence is
minimal and for the most part does not focus on high schools.

Still, ESPN's offshoots may hurt native news operations, possibly speeding the flow of fans from
print to the Web. Local papers' sports sections generally do not generate much advertising revenue,
but sports news helps keep subscribers loyal.

The New York Daily News and The Dallas Morning News had no comment over the weekend on ESPN's plans.

Sports also bring a global audience to the Web sites of local papers. ESPN wants those people, who
are probably also checking ESPN.com, to stay entirely within its arms.

ESPNChicago.com, introduced on April 13, is the model. The site covers the city's seven professional
teams and sports at nine area colleges, much of it already offered on ESPN's other platforms. Plans
are afoot to expand coverage to area high schools.

There is also a focus on pickup leisure sports. Organizing a neighborhood softball league? ESPN
Chicago has a tool to help. The site's coverage, from news wires and ESPN staff, is augmented by
original reporting from a local team of editors, columnists and bloggers -- some of them familiar to
sports fans from their previous work at The Chicago Sun-Times and the local NBC affiliate.

The site also offers a daily Chicago edition, three to six minutes long, of its flagship
``SportsCenter'' program.

``Huge'' is how Stacey Woelfel, chairman of the Radio Television News Directors Association,
described the potential threat to TV stations, in part because their Web presence in sports ``tends
to be fairly weak.''

ESPN, of course, is going after local ad dollars, as well as readers. Chicago Lincoln-Mercury
dealers, Hilton Hotels of Chicago, MillerCoors and Hawthorne Race Course have been among ESPN
Chicago's advertisers.

``A national sports brand that's dedicated to your local market? It hasn't really been seen
before,'' said Ray Elias, marketing director for StubHub, the ticket resale marketplace and an early
partner with ESPN Chicago.

Although local advertising in print, radio and television has plunged in recent years -- led by a
drop in spending from car dealerships -- online advertising is growing. The Kelsey Group, which
studies local media, values the local online ad market at about \$16 billion and expects it to
double by 2013.

John Kosner, ESPN's digital media chief, played down the competitive threat posed by his company's
plan. ``You've got outstanding newspapers and television stations that have been in these markets
forever and have strong followings,'' he said. ``Our presence will be additive.''

Mr.~Kosner even hinted at collaboration. ``It's not a foregone conclusion that we will be
independent'' in every market, he said. ESPN already has some content-sharing partnerships with
local newspapers.

ESPN, which broadcasts ``Monday Night Football,'' has looked for a way to go local for years, but
the growth potential was always outweighed by the costs of deploying local reporters and ad staffs.
Sites that simply aggregated what ESPN was already doing would be quickly dismissed by fans,
Mr.~Kosner said.

To address the costs, ESPN is building off its radio stations' Web sites. The company owns five
stations -- in Chicago, New York, Los Angeles, Dallas and Pittsburgh -- and counts more than 350
others as affiliates.

``We decided to take those sites, set up to service the station, and turn them on their ear to
service the consumer,'' said Jim Pastor, senior vice president of the ESPN Radio Station Group.

Localization will not pad ESPN's bottom line anytime soon. But it could turn into a significant
growth engine over time.

``Efforts like this might not sound so meaningful,'' said Michael Morris of UBS. ``But if this area
takes off, and ESPN's success record is astounding, it could really add up in a hurry.''

\section{In Britain, Judge Finds for Google in Libel Case}

\lettrine{G}{oogle} is not liable for defamatory\cite{defamatory} material that appears in its
search results, a British judge has ruled, a decision that lawyers call significant because of the
country's reputation as a haven for libel claimants\cite{claimant}.

The decision was in a case involving Metropolitan International Schools, a British company that
operates Internet-based training courses. The company wanted to sue Google over negative comments
posted by a reader of a technology news Web site; the comments appeared in text blurbs\cite{blurb}
with the results of Google searches related to Metropolitan.

David Eady, a High Court judge, ruled that Metropolitan could not sue Google, saying it was not a
publisher of the material. ``It has merely, by the provision of its search service, played the role
of a facilitator,'' he wrote.

Lawyers said the decision, the first of its kind in Britain, was consistent with court decisions in
several other European countries. In the United States, search engines are protected from liability
for the contents of the results they turn up. Several European countries have extended similar
protection.

Still, lawyers said Judge Eady's decision, published Friday, was significant because of
uncertainties over libel law in Britain. British courts are often seen as sympathetic to defamation
plaintiffs, and Web hosting services in Britain have been held responsible for the contents of sites
they host.

``It's significant because it's consistent with what we thought the law should mean,'' said Struan
Robertson, a technology lawyer at the firm Pinsent Masons in London. ``If the judge had ruled
otherwise, it would have been a terrible decision for search engines.''

Google said the decision ``reinforces the principle that search engines are not responsible for
content that is published on third-party Web sites.

``Justice Eady made clear if someone feels they have been defamed by material on a Web site then
they should address their complaint to the person who actually wrote and published the material, and
not a search engine, which simply provides a searchable index of content on the Internet,'' the
company's statement said.

Metropolitan International Schools could not be reached for comment.

\section{A Race to Provide Wi-Fi (but Not the Voice Part)}

\lettrine{T}{hough} it is not yet clear how extensive the market is for
in-flight Wi-Fi service, the competition to provide it is heating up. Row44, a California company
whose in-flight Wi-Fi technology uses satellite connections, is emerging as an aggressive competitor
to the market leader, Aircell, whose system uses ground-based stations.

Aircell said it had installed its system, called Gogo, on 460 airplanes, with technicians working
night shifts at airline maintenance centers across the country. The company hopes to have as many as
1,000 domestic airplanes outfitted by the end of the year.

Two airlines, Virgin America and AirTran, now have Gogo on their entire fleets. Other carriers
installing the Aircell system are Delta, American and United. Delta expects to have 90 percent of
its domestic mainline fleet outfitted by the end of September.

Meanwhile, John Guidon, the chief executive and a co-founder of Row44, said the company was ready to
begin installations, though it did not yet have a firm fleetwide contract. He estimated that Row44's
system would be installed on 500 to 1,000 airplanes worldwide by the end of 2010.

It would seem that Aircell is running so far ahead that competitors might be wary of entering the
market. But Row44 is banking on making its domestic entry to the market aboard Southwest Airlines,
which has been testing Row44 Wi-Fi on four of its more than 500 Boeing 737 aircraft.

Since June, after months of free trials on those planes, Southwest has been evaluating demand for
the service at ``various price points,'' said a spokeswoman, Whitney Eichinger.

Alaska Airlines has also been testing the Row44 system on a single aircraft.

Southwest is not expected to make any firm decisions on whether to install Wi-Fi fleetwide until
after the summer.

The big question about in-flight broadband service has always been whether enough passengers would
pay for it, especially during a poor economic period.

``Right now, as we run these trials we're seeing healthy single-digit numbers of paying customers
,'' Mr.~Guidon said. He said that number was acceptable for the initial phase of development.

``We feel that's a conservative number,'' Mr.~Guidon said. As more people acquire Wi-Fi enabled
personal devices, which are easier to use in the cramped space of coach seats than laptops, ``this
will become an essential service that customers will demand,'' he said.

Aircell told me it cost about \$100,000 to outfit an airplane with its land-based system. Mr.~Guidon
would not disclose how much the installation of Row44's satellite system costs for each plane, other
than to say, ``It's quite a lot more, let's put it that way.''

Prices for using the Aircell service have ranged as high as \$12.95, depending on the length of a
flight, though users of Wi-Fi enabled BlackBerrys and smartphones receive discounts. Last week,
Aircell began offering a 24-hour pass for \$12.95, covering multiple flights.

Mr.~Guidon anticipates that for Row44 service, airlines will charge about \$10 for long-haul flights
and \$5 for short-hauls.

Wi-Fi enabled devices, whether notebooks or smartphones, are capable of providing more than the
Internet and e-mail. One big question surrounding the proliferation of smartphones and other Wi-Fi
enabled devices on planes is whether they will soon lead to in-flight phone calls.

Both the Aircell and Row44 systems are technically capable of providing voice-over-Internet
protocol, called VoIP, service, which allows users to make phone calls over broadband Internet
connections.

Aircell and Row44 both block the service for voice calls, but there are Web forums devoted solely to
sharing tips on how to overcome the roadblocks.

Like many probably bad ideas, there seems to be a certain inevitability to having people engage in
phone calls aboard planes. In fact, some foreign carriers already allow it.

Earlier this year, for example, the big European budget airline Ryanair teamed up with a company
called OnAir to install capability for making calls with cellphones, Wi-Fi enabled smartphones and
BlackBerrys in flight on 20 airplanes.

``This is a political hot potato here in America, where there is concern over the social
implications of having people talking on phones, annoying other passengers,'' Mr.~Guidon said.
``Nevertheless, across the rest of the world this seems to be rolling out with no drama.''

He said that airlines would ultimately decide whether it should become available.

Until it does, ``we'll block voice over I.P.,'' he said. ``Nevertheless, there will continue to be a
cat-and-mouse game between people who try to figure out ways of getting voice over I.P. over Row44
service and our attempts to stop them. I think that's just what makes the world go around.''


\section{Earnings Climb Nearly 8\% at Yahoo}

\lettrine{C}{arol} Bartz, Yahoo's chief executive, is known for peppering\cite{pepper} her public
presentations with the odd expletive\cite{expletive} or two. But her conference call Tuesday to
discuss Yahoo's second-quarter financial results was entirely family-friendly. Then again, Ms.~Bartz
did not have much to swear about.

Yahoo's revenue declined 13 percent in the second quarter as advertisers continued to reduce
spending in the downturn. But cost-cutting initiatives helped to soften the blow, and profits
climbed nearly 8 percent, beating analysts' expectations.

Yahoo also said that it planned a new round of investment in products and a rebranding campaign,
which would lead to a drop in profit in the current quarter. Investors reacted by sending shares
down nearly 3 percent in after-hours trading.

``Over all, the long and protracted\cite{protract} turnaround process continues,'' said Jeffrey
Lindsay, an analyst with Sanford C.~Bernstein \& Company. ``We don't see a catalyst that is going to
turn things around soon.''

Yahoo continues to discuss a search and advertising partnership with Microsoft that would create a
more viable rival to Google. The talks have intensified recently, according to people briefed on
them.

Both companies have declined to discuss their talks publicly and, for the first time since Ms.~Bartz
became chief executive in January, analysts did not ask her about them.

But Ms.~Bartz recognized that Yahoo's search business would fare better with many more users, or
more ``scale,'' and she praised Bing, Microsoft's new search engine. ``I think Microsoft should be
getting kudos for Bing,'' Ms.~Bartz said in a conference call with investors. ``They've done a nice
job.''

Some analysts interpreted those comments as an indication that she appeared more willing to consider
a deal than ever before.

``She talked about scale being an issue,'' said Ross Sandler, an analyst with RBC Capital Markets.
``She talked about Microsoft having a pretty good search engine. That's all an admission that it
might make sense for these parties to come together.''

Yahoo reported net income of \$141 million, or 10 cents a share, compared with \$131 million, or 9
cents a share, a year earlier. Revenue dropped to \$1.57 billion, from \$1.8 billion a year earlier.

Net revenue, which excludes commissions paid to advertising partners, was \$1.14 billion, down 15
percent from \$1.35 billion a year earlier.

On average, Wall Street analysts had expected Yahoo to report net income of 8 cents a share on net
revenue of \$1.14 billion.

``Considering the economy I am pleased with our results,'' Ms.~Bartz said. ``Over all we are seeing
less fear in the marketplace.'' Ms.~Bartz said advertisers appeared ready to spend more, but she
added that it was too early to say how that would affect Yahoo.

Yahoo said search advertising revenue on its sites declined 15 percent, even as the number of
searches on Yahoo rose. By comparison, Google, whose business is largely driven by search
advertising, last week reported a 3 percent growth in revenue. Analysts said Yahoo did not give good
reasons for the disappointing results.

On Tuesday, Yahoo also introduced an overhauled home page, a major initiative intended to restore
its luster\cite{luster} with users, advertisers and investors.

The new home page allows users to preview Yahoo services like mail and Flickr, and also third-party
applications and sites, like Facebook, Gmail or BBC World News, in a section on the left side of the
page called My Favorites. As Internet users' attention is increasingly fragmented across a growing
number of sites and services, the company's decision to include third-party content on its home
represents an effort by Yahoo to establish itself as the ``center point of people's lives online,''
said Tapan Bhat, a senior vice president at Yahoo who is responsible for the company's home page.

With 114 million visitors in June in the United States alone, Yahoo.com remains the most visited
home page online. But Yahoo's efforts to cash in on its huge audience have faltered\cite{falter} in
recent years, as prices for premium brand ads have declined and as advertising networks have given
marketers the ability to reach large audiences across vast arrays of Web sites at lower cost.

The new design could help Yahoo. The preview windows for third-party applications will include ads
and should help advertisers reach large and specific groups of customers, Mr.~Bhat said.

\section{Challenge to Health Bill: Selling Reform}

\lettrine{W}{hat's} in it for me?

On the subject of health care reform, most Americans probably don't have a good answer to the
question. And that, obviously, is a problem for the White House and for Democratic leaders in
Congress.

Current bills would expand the number of insured -- but 90 percent of voters already have insurance.
Congressional leaders say the bills would cut costs. But experts are dubious. Instead, they point
out that covering the uninsured would cost billions.

So the typical person watching from afar is left to wonder: What will this project mean for me,
besides possibly higher taxes?

Barack Obama was able to rise from the Illinois State Senate to the presidency in large measure
because of his ability to explain complex issues and then to make a persuasive argument. He now has
a challenge worthy of his skills.

Our health care system is engineered, deliberately\cite{deliberate} or not, to resist change. The
people who pay for it -- you and I -- often don't realize that they're paying for it. Money comes
out of our paychecks, in withheld taxes and insurance premiums, before we ever see it. It then flows
to doctors, hospitals and drug makers without our realizing that it was our money to begin with.

The doctors, hospitals and drug makers use the money to treat us, and we of course do see those
treatments. If anything, we want more of them. They are supposed to make us healthy, and they appear
to be free. What's not to like?

The immediate task facing Mr.~Obama -- in his news conference on Wednesday night and beyond -- is to
explain that the health care system doesn't really work the way it seems to. He won't be able to put
it in such blunt terms. But he will need to explain how a typical household, one that has insurance
and thinks it always will, is being harmed.

The United States now devotes one-sixth of its economy to medicine. Divvy that up, and health care
will cost the typical household roughly \$15,000 this year, including the often-invisible
contributions by employers. That is almost twice as much as two decades ago (adjusting for
inflation). It's about \$6,500 more than in other rich countries, on average.

We may not be aware of this stealth \$6,500 health care tax, but if you take a moment to think, it
makes sense. Over the last 20 years, health costs have soared, and incomes have grown painfully
slowly. The two trends are directly connected: employers had to spend more money on benefits,
leaving less for raises.

In exchange for the \$6,500 tax, we receive many things. We get cutting-edge research and heroic
surgeries. But we also get fabulous amounts of waste -- bureaucratic and medical.

One thing we don't get is better health than other rich countries, whether it's Canada, France,
Japan or many others. In some categories, like emergency room care, this country seems to do better.
In others, like chronic-disease care, it seems to do worse. ``The fact that we spend all this money
and don't have better outcomes than other countries is a sign of how poorly we're doing,'' says
Dr.~Alan Garber of Stanford University. ``We should be doing way better.''

So far, no one has grabbed the mantle as the defender of the typical household -- the opponent of
spending that creates profits for drug companies and hospitals at no benefit to people's health and
at significant cost to their finances.

Republicans have actually come out against doing research into which procedures improve health. Blue
Dog Democrats oppose wasteful spending but until recently have not been specific. Liberals rely on
the wishful idea -- yet to be supported by evidence -- that more preventive care will reduce
spending. The American Medical Association, not surprisingly, endorses this notion of doing more
care in the name of less care.

Mr.~Obama says many of the right things. Yet the White House has not yet shown that it's willing to
fight the necessary fights. Remember: the \$6,500 tax benefits someone. And that someone has a
lobbyist. The lobbyist even has an argument about how he is acting in your interest.

These lobbyists, who include big names like Dick Armey and Richard Gephardt, have succeeded in
persuading Congress to write bills with a rather clever feature. They include some of the ideas that
would cut costs -- but defang them.

One proposal would pay doctors based on the quality of care, rather than quantity, but it's a pilot
project. Doctors who already provide good care may well opt in; doctors providing wasteful but
lucrative care surely will not. The bills would also finance research on which treatments are
effective. But Medicare officials would not be prevented from continuing to spend taxpayer money on
ineffective treatments.

In reaction, some people who should be natural supporters of reform have become critics. The Mayo
Clinic -- one of Mr.~Obama's favorite models of care -- says the legislation fails to ``help create
higher-quality, more affordable health care.''

On Thursday, Mr.~Obama will visit another example he likes to cite, the Cleveland Clinic. Its
successes capture what real reform would look like. Like Mayo, the Cleveland Clinic pays its doctors
a salary, rather than piecemeal, and delivers excellent results for relatively little money.

``I came here 30-some years ago,'' Delos Cosgrove, a heart surgeon who is the clinic's chief
executive, told me. ``And I have never received any additional pay for anything I did. It never made
a difference if I did five heart operations or four -- I got paid the same amount of money. So I had
no incentive to do any extra tests or anything.''

This is the crux of the issue, economists say: the current fee-for-service system needs to be
remade. The administration has made some progress, by proposing a powerful new Medicare overseer who
could force the program to pay for good results and stop paying for bad ones.

But even a strong Medicare plan won't be enough. Reform will need to attack the piecemeal system in
numerous ways. Among the most promising, which Mr.~Obama has resisted, is a limit on tax subsidies
for the costliest health insurance plans. This limit would give households and employers a reason to
become smarter shoppers.

Above all, reform can't revolve around politely asking the rest of the medical system to become more
like the Cleveland Clinic.

In recent weeks, polls have shown that a solid majority of Americans support the stated goals of
health reform. Most want the uninsured to be covered and want the option of a government-run
insurance plan. Yet the polls also show that people are worried about the package emerging from
Congress.

Maybe they have a point.

\section{A Defining Moment Nears for President}

\lettrine{S}{ix} months into his administration, President Obama is at a pivotal moment. He has
pushed through a \$787 billion economic stimulus package, bailed out Wall Street and, on Tuesday,
managed to beat the defense industry in the Senate, which voted to kill a high-profile fighter jet
program.

But the public, and lawmakers, are growing skittish over Mr.~Obama's next big plan, to remake the
American health care system. How he handles the issue over the next several weeks could shape the
rest of his presidency, shedding light on his political strength, his relationship with both parties
in Congress and the extent to which he is willing to bend in fighting for his agenda.

With some fellow Democrats balking over his insistence that both the House and the Senate pass
health legislation before the August recess, Mr.~Obama has a tough decision to make: Does he take a
hard line, demanding that lawmakers stick to his timetable -- and risk losing the support of
Republicans and moderate Democrats? Or does he signal flexibility, allowing lawmakers to take their
time -- and give opponents the chance to marshal their case against the bill?

``He's got to be careful that while he ratchets up the pressure, he doesn't bet his whole presidency
on whether this gets done before the August recess,'' said Kenneth M.~Duberstein, who orchestrated
President Ronald Reagan's first-term legislative strategy. ``He has a broad, broad agenda that he's
in a rush to enact, and if he's not careful he will be viewed as a steamroller who tries to get
things fast and not necessarily right.''

On Wednesday, Mr.~Obama will address the nation in a prime-time news conference. Rahm Emanuel, the
White House chief of staff, said in an interview that the president intended to use it as a
``six-month report card,'' to talk about ``how we rescued the economy from the worst recession'' and
the legislative agenda moving forward, including health care and energy legislation, which squeaked
through the House and faces a tough road in the Senate.

Polls show that Mr.~Obama is more popular than his own policies, a worrisome sign for a president
with such an ambitious agenda. Mickey Edwards, a former Republican congressman who is now vice
president of the Aspen Institute, said Mr.~Obama might be making a mistake in reading his election
as a mandate for dramatic change.

``A lot of people supported Obama because they wanted to repudiate the Bush administration,'' said
Mr.~Edwards, who backed Mr.~Obama for president. ``I was one of those people who supported him for
reasons other than the policies he is proposing. He seemed more thoughtful, more contemplative -- I
felt he had the right temperament to be president. But I think his health care proposal goes beyond
what the public at the moment is ready to accept.''

Mr.~Obama came into office promising a more bipartisan Washington tone, which he has so far been
unable to achieve. His actions in the coming weeks on health care may determine his long-term
relationship not only with Republicans but also with his fellow Democrats.

``I think this will be a major factor in defining his presidency,'' said Tom Daschle, the former
Senate Democratic leader, who remains a close adviser to the White House on health issues. ``Because
he's made it such an issue, and because he has invested so much personal time and effort, this will,
more than stimulus and more than anything he has done so far, be a measure of his clout and of his
success early on. And because it is early on, it will define his subsequent years.''

On the Republican side, one question is whether Mr.~Obama will succumb to the temptation to turn
health care into a partisan fight, even as he tries to court the opposing party. He is, after all,
still a popular new president confronting an unpopular Republican Party, and so it would be easy for
him to demonize Republicans as obstructionists who want to stand in the way of progress.

Senator Jim DeMint, Republican of South Carolina, gave Mr.~Obama an opening to do just that the
other day, and the president took it. Mr.~DeMint called health care a ``Waterloo moment'' that could
break Mr.~Obama. The president struck back, declaring, ``This isn't about me.'' But if Mr.~Obama
extends that line of attack to Republicans more broadly, and rams a bill through without their
support, any claim he may have to bipartisanship will quickly evaporate.

As for Democrats, Mr.~Obama faces a balance-of-power conundrum. He has said all along that he will
set out broad principles for a bill and leave the details to Congress. But now House Democrats in
the fiscally conservative Blue Dog Coalition, including seven who hold decisive votes on the Energy
and Commerce Committee, say they will not support the House bill without big changes.

One question for Mr.~Obama is whether to try to strong-arm them, and face a rebellion from some of
the very same conservative Democrats who helped put him in office. If he forces them to vote for a
bill their constituents do not like, on a timetable that feels too rushed for them, it could hurt
them at home. That could mean a bigger political problem for the White House: a resulting loss of
Democratic seats in the 2010 midterm elections.

Another question is how hard Mr.~Obama will push Congress as a whole to adopt his progressive
agenda, not only on health care but also on climate change and a variety of other issues.

The next few weeks, as the president tries to broker a health care deal, may well tell Americans
just how far he is willing to go.

\section{Obama Wins Crucial Round in Senate Vote on F-22}

\lettrine{W}{ith} some of his political capital on the line, President Obama won a crucial victory
on Tuesday when the Senate voted to strip out \$1.75 billion in financing for seven more F-22 jet
fighters from a military authorization bill.

The president had repeatedly threatened to veto the \$679.8 billion bill if it included any money
for the planes. The 58-to-40 vote clearly gives the Obama administration more leeway to overhaul
military spending.

The F-22, the world's most advanced fighter, has been a flashpoint in a battle over the
administration's push to shift more of the Pentagon's resources away from conventional warfare
projects, like the F-22, to provide more money for fighting insurgencies.

Senate aides said that some Democrats who otherwise might have voted for more planes sided with the
president out of concern that a loss could have hurt him in the fight for health care reform.

``The president really needed to win this vote,'' Senator Carl Levin, a Democrat from Michigan who
led the fight to cut financing for the plane, said after the vote.

Lockheed Martin, the prime contractor for the F-22, has estimated that work on the plane provides
25,000 jobs and indirectly supports about 70,000 others. But Robert M.~Gates, the defense secretary,
has said that the Pentagon needs to accelerate a new plane, the F-35, and that doing so would offset
the job losses.

About 1,000 suppliers in 44 states provide the jobs, which will gradually be phased out as some of
the 187 F-22s that have been ordered are completed.

About two-thirds of the jobs are in California, Texas, Georgia, Washington and Connecticut. Several
large unions who supported Mr.~Obama in his campaign for the presidency, back building more planes.

All four senators from California and Washington are Democrats, and they all voted in favor of
preserving the money for more planes.

The senators from Connecticut -- Senator Christopher J.~Dodd, a Democrat, and Senator Joseph
Lieberman, a former Democrat who is now an independent -- also voted to keep the money in the
budget, as did the four Republican senators from Georgia and Texas.

Senator John Kerry of Massachusetts, the Democratic presidential nominee in 2004, had long supported
the plane, partly because of jobs in Massachusetts, but he voted on Tuesday to strip out the money.
Senator Edward M.~Kennedy of Massachusetts, another senior Democrat who also had supported the
plane, is battling brain cancer and did not vote Tuesday.

Immediately after the vote, Mr.~Obama praised the Senate's decision, saying that any money spent on
the fighter was an ``inexcusable waste'' -- and that by following his lead the Senate had
demonstrated a commitment to changing Washington's ingrained habits.

He also received crucial support from his Republican rival in last year's presidential election,
Senator John McCain of Arizona, who co-sponsored an amendment with Mr.~Levin to remove the \$1.75
billion from the bill.

Mr.~McCain told reporters after the vote that the result was ``definitely attributable'' to the
strong push by the president and Mr.~Gates.

Mr.~McCain added that the vote ``really means there's a chance of us changing the way we do business
in Washington,'' particularly in terms of Pentagon contracting.

Mr.~Levin said that Mr.~Gates and the White House chief of staff, Rahm Emanuel, made phone calls to
influential senators to rally support.

``This was a very significant decision that the Senate made after a very, very tough battle,''
Mr.~Levin said.

Despite Mr.~Obama's veto threat, the Armed Services Committee had set the stage for Tuesday's
decision by voting 13 to 11 in late June to shift the \$1.75 billion from other parts of the
Pentagon's budget for 2010 to add the seven planes to the 187 that have been built or ordered.

The House has also voted to keep the plane alive by authorizing \$369 million to buy advanced parts
for 12 more F-22s. Ultimately, a conference committee will decide the next step.

And even though Tuesday's vote represented a test of strength for the White House, the issue could
potentially resurface in the Senate. Senator Daniel Inouye, a Democrat from Hawaii and the chairman
of the Senate Appropriations Committee, was among the Democrats who voted on Tuesday to preserve the
money for the planes, and his panel will be putting together a separate military appropriations bill
soon.

Senator Saxby Chambliss, a Republican from Georgia who led the fight to save the plane, said after
the vote that he was disappointed. But he added that he hoped a law banning exports of the plane
might be changed to allow it to be sold to allies like Japan. Lockheed assembles the planes in
Marietta, Ga.

Still, military analysts said supporters had made much more progress in saving the F-22 than most
experts had expected when Mr.~Gates announced plans in April to cancel it and other major weapons
systems.

Critics have long portrayed the F-22 as a cold war relic. The plane was designed in the late 1980s
and can perform tactical operations at higher altitudes than other fighters. It can cruise at
supersonic speeds without using telltale afterburners, and it has a stealthy skin that scatters
radar detection signals. Proponents see it as a form of insurance against possible wars with
countries like China.

But the F-22 has never been used in war, and the Pentagon's focus has shifted to simpler weapons
needed in Iraq and Afghanistan.

Air Force leaders recently agreed that they could make do with the F-22s already built or ordered,
instead of the 381 that the service had sought.

Mr.~Gates has said a new fighter, the F-35, is better designed to attack ground targets. The plane
will be used by the Navy, the Marine Corps and the Air Force, and the Pentagon plans to buy more
than 2,400 of them.

\section{Why China Will Keep Investing Abroad}

\lettrine{O}{vershadowed} by a sparkling economic report card for the second quarter, the leap in
China's official currency reserves last week is yet one more reason to bank on a transformative
surge in the country's investments abroad.

But the diplomatic storm whipped up by charges leveled by Beijing against four employees of Rio
Tinto shows just how bumpy the journey overseas may be if it triggers a backlash against Chinese
interests.

China has long been scouring the globe for energy and commodities to feed its thrumming economy.
What is new is the leadership's determination to increase outbound foreign direct investment, or
O.F.D.I., as it weans the economy off low-value, export-oriented manufacturing. The deal by Sinopec,
the largest Chinese oil refiner, to buy the Swiss oil explorer Addax for \$7.24 billion last month
was China's largest overseas acquisition yet.

The government is not throwing caution to the wind. Beijing blocked Bank of China's purchase of a
stake in La Compagnie Financière Edmond de Rothschild, the French merchant bank, and has responded
tepidly to a bid by Sichuan Tengzhong, a little-known machinery maker, for General Motors' Hummer
unit.

But the first half of 2009 may prove to be an inflection point for Chinese outbound foreign direct
investment, said Daniel Rosen, a visiting fellow at the Peterson Institute for International
Economics in Washington.

``Despite short-term anxieties, Chinese O.F.D.I. is poised to grow markedly in the medium and long
term, and the importance of these investments to Chinese firms is changing fundamentally as the
nation confronts the need to rebalance its growth model,'' Mr.~Rosen wrote in a paper whose
co-author was Thilo Hanemann.

With domestic economies of scale largely exhausted, companies will have a powerful incentive to move
abroad to upgrade their manufacturing and compete in more profitable areas like distribution, design
and branding. ``Made in China'' will increasingly give way to ``Made by China -- abroad,'' Mr.~Rosen
and Mr.~Hanemann argued.

On top of the commercial motives, greater capital outflows would make monetary policy easier to
conduct by reducing China's balance of payments surplus. Because the central bank buys most of the
foreign currency entering China to cap the yuan's exchange rate, every dollar that can be recycled
as outbound investment is one dollar less that must be bought and added to reserves.

Those reserves jumped \$177.9 billion in the second quarter to \$2.13 trillion, helping to generate
record credit and money supply growth that economists fear may be creating asset bubbles.

``We don't know what the right amount of domestic monetary creation is, but it's pretty safe to say
it's much, much lower than the increase in central bank reserves,'' said Michael Pettis, an
economics professor at Peking University.

Hence the significance of new government rules issued last week that will make it easier for
companies to finance outbound foreign direct investment.

``On top of China needing to invest -- they need to secure energy supplies, they need to acquire
technology -- there's also the element of recycling dollar revenues,'' said Qu Hongbin, chief China
economist at HSBC in Hong Kong.

China's outbound foreign direct investment has increased steadily this decade, and the outflow of
\$52 billion last year was a record. Still, the historical stock of the outbound direct investment
of \$170 billion is puny next to China's foreign exchange reserves and its inward foreign investment
stock of \$876 billion.

Mr.~Qu, however, said that outflows could reach \$100 billion to \$150 billion a year as soon as
2012. And Mr.~Rosen said in an e-mail message that sum would be ``entirely conceivable'' if there
were no political interference from Beijing.

Mr.~Rosen and Mr.~Hanemann identify four areas where progress, now blocked by the state's reluctance
to take its hands off the economy, is needed to promote outbound investment.

The government must completely pull back from companies' investment decisions. It must further
liberalize access to foreign exchange and O.F.D.I. funding. It must let all companies invest
overseas. And it needs a strategy that serves China's long-term interests -- defending pariah states
where China has invested hurts the image of the country's companies among consumers.

Similar reputational damage is where the detention by China of Stern Hu, an Australian and Rio's
chief negotiator in fractious iron ore pricing talks, may dim China's ambitions to expand abroad.
The perception that Beijing is exacting revenge on Rio for being outsmarted in the ore talks risks
hardening attitudes far beyond Australia.

After all, critics already complain that foreign companies in China cannot trade and invest on a
level playing field. So why should the West roll out the welcome mat for Chinese investors?

Mr.~Pettis of Peking University said he expected outbound foreign direct investment to keep growing
despite ``political noise'' generated by trophy transactions, but he acknowledged that Chinese
attitudes toward overseas investors would hardly encourage foreigners to welcome Chinese investment.

``It's too easy if you are opposed to a Chinese acquisition to point out the things that have
happened here in China,'' he said.

The ``callousness'' with which Beijing has blocked some inward investments raises questions about
its seriousness toward cross-border investments both ways, Mr.~Rosen and Mr.~Hanemann said. Yet the
implications of the Rio case are not clear-cut.

The charges laid at Rio's door reflected China's immaturity on matters of commercial regulation,
Mr.~Rosen wrote in his message. ``On the one hand, that immaturity is an impediment to cross-border
investment flows both ways. But on the other, and somewhat counterintuitively, it is one of the
reasons China is so attractive as an investor and as a place to invest.''

\section{China's Wide Reach in Africa}

\lettrine{A}{mong} Westerners, the economic partnership between China and Africa is often
overlooked. But in ``China Safari: On the Trail of Beijing's Expansion in Africa'' (Nation Books,
\$27.50), Serge Michel and Michel Beuret examine the roots of this relationship -- and argue that
China is engaged in a conquest of Africa that will have worldwide economic implications.

As French journalists, Mr.~Michel and Mr.~Beuret bring an acute awareness of their own country's
colonial history to the China-Africa story. Mr.~Michel, a former West Africa correspondent for Le
Monde, has also reported from Iran, Iraq and Afghanistan. Mr.~Beuret, the foreign editor of the
Swiss magazine L'Hebdo, has written extensively on human trafficking in China and Europe.

``China Safari'' is a fascinating, provocative work of firsthand reporting that illuminates an
important global economic story. The book also features a 16-page insert of color photographs shot
by Paolo Woods, who puts human faces on the book's sprawling story and highlights some of the stark
juxtapositions of African laborers and their Chinese bosses.

In 1976, under Mao Zedong, China completed the Tan-Zam railway, linking Zambia to the port of Dar es
Salaam in Tanzania. But, as the authors note, Chinese inroads into Africa really got a kick-start in
1995, when President Jiang Zemin made a speech urging Chinese business leaders to ``Go abroad!
Become world players!''

Chinese business interests in Africa have grown dramatically in recent years. The authors say that
bilateral trade between the regions quintupled, to \$55 billion, from 2000 to 2006, and that the
figure is expected to reach \$100 billion by 2010. Chinese business interests in Africa range from
oil, lumber, refining, agriculture, mining, textiles and banking to the construction of dams,
railroads, highways, bridges, airports and housing.

The authors contend that China's ambitions in Africa are grandly geopolitical as well as economic.
As Jacob Wood, a Shanghai-born housing developer based in Africa for more than 30 years, tells them:
``I'm going to be honest with you, China is using Africa to get where the United States is now, and
surpass it.''

According to one report cited by the authors, there are now about 750,000 Chinese living and working
in Africa, in countries including South Africa, Nigeria, Zambia, Sudan, Algeria, Congo, Zimbabwe,
Mozambique, Angola, Gabon, Guinea, Ethiopia, Cameroon, Egypt and Chad.

So far, China's ventures in Africa have produced decidedly bittersweet results.

On the one hand, the Chinese invest heavily in infrastructure projects and create hundreds of
thousands of jobs in return for being granted oil, mining and other concessions by African nations.
According to the authors, Chinese ``discretion and humility'' are a breath of fresh air against a
backdrop of ``colonial arrogance'' by France, Britain and other nations. In addition, the authors
say, ``Africa welcomes any competition that shakes up the Western, Lebanese and Indian business
monopolies.''

Many African leaders are enamored\cite{enamor} of the Chinese mix of authoritarianism and capitalism
in business affairs, an emphasis on efficiency and a lack of preaching\cite{preach} about human
rights, the authors say. Moreover, when the Chinese talk, they back up their words with concrete
actions.

``The Chinese build things, the Europeans don't,'' declares Claude Alphonse N'Silou, the minister of
construction and housing of the Congo Republic.

But there are increasing signs of trouble in the Sino-African economic paradise. The book says the
Chinese have been criticized for selling arms to Sudan and Zimbabwe in the past, and for showing
support for the Zimbabwean dictator Robert Mugabe. The authors also cite a United Nations report
that says that thousands of machetes used in the Rwandan genocide were imported from China.

And the everyday workplace is riven\cite{rive} with accusations of mistreatment of African laborers
by Chinese bosses, and accusations of virulent\cite{virulent} racism that seems to contradict the
Chinese image of discretion and humility.

An opposition leader in Zambia, meanwhile, complains that his country is allowing Zambia to become
``a province -- make that a district -- of China,'' adding, ``There's nothing any minister can do
when confronted by China's amoral scheming.'' The authors also cite instances of Chinese companies
recently abandoning infrastructure projects without public explanation.

Since ``China Safari'' went to press, there have been reports that China may be backing out of deals
in Guinea and Congo.

``CHINA SAFARI'' is hobbled by the fact that Mr.~Michel and Mr.~Beuret were often denied access to
important business people and projects. As a result, the authors occasionally resort to rambling
anecdotal narratives like one in which they trace the manufacture of plastic souvenir pyramids sold
at Egyptian archeological sites to -- where else? -- China.

Despite the recitations of socioeconomic and cultural disruptions and workplace abuses, the authors
reach what may be an overly sanguine\cite{sanguine} conclusion. They contend that ``China's arrival
has been a boon for a continent adrift,'' adding that the Chinese have ``given Africa a real sense
of worth, as much in the eyes of Africans themselves as in the eyes of foreigners.''

Still, it is not hard to join the authors in predicting that this joining of Chinese and African
interests will likely succeed to the chagrin of the rest of the business world.

\section{Unraveling how children become bilingual so easily}

\lettrine{T}{he} best time to learn a foreign language: Between birth and age 7. Missed that window?

New research is showing just how children's brains can become bilingual so easily, findings that
scientists hope eventually could help the rest of us learn a new language a bit easier.

``We think the magic that kids apply to this learning situation, some of the principles, can be
imported into learning programs for adults,'' says Dr.~Patricia Kuhl of the University of
Washington, who is part of an international team now trying to turn those lessons into more
teachable technology.

Each language uses a unique set of sounds. Scientists now know babies are born with the ability to
distinguish all of them, but that ability starts weakening even before they start talking, by the
first birthday.

Kuhl offers an example: Japanese doesn't distinguish between the ``L'' and ``R'' sounds of English
-- ``rake'' and ``lake'' would sound the same. Her team proved that a 7-month-old in Tokyo and a
7-month-old in Seattle respond equally well to those different sounds. But by 11 months, the
Japanese infant had lost a lot of that ability.

Time out -- how do you test a baby? By tracking eye gaze. Make a fun toy appear on one side or the
other whenever there's a particular sound. The baby quickly learns to look on that side whenever he
or she hears a brand-new but similar sound. Noninvasive brain scans document how the brain is
processing and imprinting language.

Mastering your dominant language gets in the way of learning a second, less familiar one, Kuhl's
research suggests. The brain tunes out sounds that don't fit.

``You're building a brain architecture that's a perfect fit for Japanese or English or French,''
whatever is native, Kuhl explains -- or, if you're a lucky baby, a brain with two sets of neural
circuits dedicated to two languages.

It's remarkable that babies being raised bilingual -- by simply speaking to them in two languages --
can learn both in the time it takes most babies to learn one. On average, monolingual and bilingual
babies start talking around age 1 and can say about 50 words by 18 months.

Italian researchers wondered why there wasn't a delay, and reported this month in the journal
Science that being bilingual seems to make the brain more flexible.

The researchers tested 44 12-month-olds to see how they recognized three-syllable patterns --
nonsense words, just to test sound learning. Sure enough, gaze-tracking showed the bilingual babies
learned two kinds of patterns at the same time -- like lo-ba-lo or lo-lo-ba -- while the
one-language babies learned only one, concluded Agnes Melinda Kovacs of Italy's International School
for Advanced Studies.

While new language learning is easiest by age 7, the ability markedly declines after
puberty\cite{puberty}.

``We're seeing the brain as more plastic and ready to create new circuits before than after
puberty,'' Kuhl says. As an adult, ``it's a totally different process. You won't learn it in the
same way. You won't become (as good as) a native speaker.''

Yet a soon-to-be-released survey from the Center for Applied Linguistics, a nonprofit organization
that researches language issues, shows U.S.~elementary schools cut back on foreign language
instruction over the last decade. About a quarter of public elementary schools were teaching foreign
languages in 1997, but just 15 percent last year, say preliminary results posted on the center's Web
site.

What might help people who missed their childhood window? Baby brains need personal interaction to
soak\cite{soak} in a new language -- TV or CDs alone don't work. So researchers are improving the
technology that adults tend to use for language learning, to make it more social and possibly tap
brain circuitry that tots would use.

Recall that Japanese ``L'' and ``R'' difficulty? Kuhl and scientists at Tokyo Denki University and
the University of Minnesota helped develop a computer language program that pictures people speaking
in ``motherese,'' the slow exaggeration of sounds that parents use with babies.

Japanese college students who'd had little exposure to spoken English underwent 12 sessions
listening to exaggerated ``Ls'' and ``Rs'' while watching the computerized instructor's face
pronounce English words. Brain scans -- a hair dryer-looking device called MEG, for
magnetoencephalography -- that measure millisecond-by-millisecond activity showed the students could
better distinguish between those alien English sounds. And they pronounced them better, too, the
team reported in the journal NeuroImage.

``It's our very first, preliminary crude attempt but the gains were phenomenal,'' says Kuhl.

But she'd rather see parents follow biology and expose youngsters early. If you speak a second
language, speak it at home. Or find a play group or caregiver where your child can hear another
language regularly.

``You'll be surprised,'' Kuhl says. ``They do seem to pick it up like sponges.''

\section{In China, a Rocky Ascent for Basketball}

\lettrine{W}{ith} 1.3 billion potential fans, China is increasingly seen as a financial promised
land for N.B.A. stars through endorsement deals, and the league itself has established a robust
organization here valued at \$2 billion.

But China's own professional league, the Chinese Basketball Association, has hardly enjoyed a smooth
ascendance alongside this country's basketball boom. American players and agents describe broken
contracts, unpaid wages, suspicions of game-fixing and rising resentment toward foreign players.
Several players have left China after failing to receive paychecks. Last month, the league announced
that it lost \$17 million last season, which ended in May.

Players and coaches in China's professional league said problems escalated\cite{escalate} last
season after the association loosened salary and court-time restrictions on foreign players, part of
an effort to heighten the game's appeal to China's growing N.B.A fan base and to bring in more
lucrative\cite{lucrative} sponsorship deals. The association also hoped the prowess\cite{prowess} of
imported players would help bolster China's basketball prospects for the 2012 Summer Olympics.

The efforts yielded conflicting results. TV ratings soared, and foreign players found starring roles
-- the top 15 scorers were non-Chinese, and players like Bonzi Wells and Donta\'e Jones -- who had
less than stellar N.B.A. careers -- frequently scored more than 40 points a game. At the same time,
the dominance of foreign players fueled frustration.

``Foreigners should play supporting roles, not dominate the game,'' said Zhang Xiong, director of
operations for the Chinese Basketball Association.

Li Xiaofeng, 20, a restaurant manager and C.B.A. fan, said: ``I don't like foreign players. They got
most of the chances to shoot and score. How about our own players? They don't have the chance to
bring their skill and talent into play.

``Our Chinese players' ability is limited by the current rule.''

Some Chinese state news media outlets went so far as to call imported players a
``malignant\cite{malignant} tumor.''

Meanwhile, China's most prominent homegrown player, Yao Ming, is an N.B.A. star. This month he
bought the financially troubled Shanghai Sharks, for whom he played five seasons before joining the
N.B.A.

Chinese players like Wang Yong of the Dongguan Leopards support the increased participation of
foreign players. ``Chinese and foreign players are a harmonious\cite{harmonious} blend,'' he said.
``I've learned a lot from them this season and feel I am a better player.''

Foreign players bridled\cite{bridle} at accusations that they were selfish, saying they were simply
following orders.

``The coaches tell you you're the main scorer,'' said Corsley Edwards, an American who played for
the Yunnan Bulls last season.

The dominance of international players is not the only systemic problem in the 18-team league.
Coaches, visiting players and their agents suspect that the outcome of some games is predetermined.

Players recounted locker-room lectures in which they were told to slack off on the court. On other
occasions, they said, the best players had to sit out particularly competitive games or were sent
home once their teams made the playoffs.

Gabe Muoneke, an American player who joined the Yunnan Bulls last season, said he was told by a
Chinese teammate that a game against the Shanghai Sharks in November was fixed.

``He said, 'Listen, my bookie told me we're going to win today, so don't worry,''' Muoneke said.

At the time, Shanghai and Yunnan were the two worst teams in the league, both without a victory. The
Bulls won, 107-97.

Muoneke said the incident confirmed what he and other players have long suspected: that game-fixing
is a problem for the Chinese league.

``It's common knowledge that Chinese teams bribe referees\cite{referee},'' he said.

Awvee Storey, a former N.B.A. player with the Liaoning Hunters, said he often sensed his Chinese
teammates were going through the motions. ``I felt a lot of times we were playing just to play and
not to win,'' he said.

The league and Storey's team denied allegations of game-fixing. The Yunnan Bulls contended that
their international players did not understand Chinese basketball.

``C.B.A. referees are not very good,'' said Wu Li, a team spokesman for the Liaoning Hunters. ``Lots
of people think referees make bad calls because they are being influenced by teams or coaches, but
we don't know of any proven cases.''

Wu went on to say that the foreign players' bad attitudes caused many problems.

When asked about game-fixing, Zhang of the C.B.A. said, ``There haven't been any problems like
that.''

Giovanni C.~Funiciello, an agent who has sent players to China for more than eight years, said that
although most games are played fairly, betting is a problem.

``Do I think some games to a degree are influenced?'' he said. ``Yeah, I would say so.''

The accusations have led to a spate\cite{spate} of articles in the state-controlled Chinese news
media about game-fixing and bribery, and the league vowed to crack down on such cheating.

Although the C.B.A. has denied the accusations, in November it announced harsher penalties for
official misconduct. Social interactions between referees and team officials are now prohibited.

``This season we will put a knife to the neck of any referee who is involved in match-fixing or
bribery,'' Liu Xiaonong, the commissioner of the C.B.A., said last year, according to China Daily.
``If a league is frequently linked to rumors of match-fixing, it means it has a big problem.''

Liu recently told Beijing Youth Daily that because of increased expenses from the C.B.A.'s expanded
schedule, ``the operation of the league and clubs are both in danger.''

The N.B.A. is undoubtedly watching these developments with interest. It has created an operation
here that involves marketing partnerships, merchandising and events. Companies including Bank of
China and ESPN invested \$253 million to acquire a stake in the N.B.A. endeavor. The N.B.A. is also
involved with building 12 basketball arenas in China.

``I'm not aware of cheating, but we're not involved with it,'' said David Stern, the commissioner of
the N.B.A. ``That's totally 100 percent under control of the C.B.A., and they've made it clear
they'd like to keep it that way.''

The Chinese league also faces allegations by its players that teams have reneged\cite{renege} on
contracts or failed to pay salaries.

Many former N.B.A. players were lured to China by six-figure salaries, in addition to free meals,
lodging and family visits. Local players, by contrast, earn about \$14,000 a season.

Faced with ballooning budgets and bleak championship prospects, some owners chose to throw in the
towel, leaving imports unpaid and abandoned.

The story is familiar to players like Muoneke and Edwards, who left the Yunnan Bulls over contract
disputes. Edwards played for three months, then the team's general manager told him he would not be
paid the rest of his salary. Filing a claim with the league got him nowhere, Edwards said.

His Chinese teammates face a similar plight\cite{plight}. In May, long after their season ended,
they sent a letter to the C.B.A. saying they had yet to receive 70 percent of their salaries and
other compensation. The matter remains unresolved.

\section{Clinton Speaks of Shielding Mideast From a Nuclear Iran}

\lettrine{S}{tiffening} the American line against Iran, Secretary of State Hillary Rodham Clinton
said Wednesday that the United States would consider extending a ``defense umbrella'' over the
Middle East if the country continued to defy international demands that it halt work that could lead
to nuclear weapons.

While such a defensive shield has long been assumed, administration officials in Washington
acknowledged Wednesday that no senior official had ever publicly discussed it. Some of the officials
said the timing of Mrs.~Clinton's remarks reflected a growing sense that President Obama needed to
signal to Tehran that its nuclear ambitions could be countered militarily, as well as
diplomatically.

It also signified increasing concern in Washington that other Middle East states -- notably Saudi
Arabia and Egypt -- might be tempted to pursue their own nuclear programs for fear Iran was growing
closer to realizing its presumed nuclear ambitions.

Mrs.~Clinton later clarified her comments on Iran, delivered in advance of a regional meeting here,
saying her warning that the United States might create such an umbrella shield did not represent any
backing away from the Obama administration's position that it must prevent Tehran from obtaining a
bomb capability. But her words suggested that the administration was developing a strategy should
all efforts at negotiation fail.

Her statement also came as Iran's internal divisions and crackdown on post-election protests have
complicated Mr.~Obama's pledge to ``engage'' Iran directly. Iranian officials have hinted that they
will present new proposals on the nuclear program, and American officials have said their offers to
negotiate stand.

Speaking during a televised town hall meeting in Bangkok, Mrs.~Clinton said, ``We want Iran to
calculate what I think is a fair assessment, that if the U.S.~extends a defense umbrella over the
region, if we do even more to support the military capacity of those in the gulf, it's unlikely that
Iran will be any stronger or safer, because they won't be able to intimidate and dominate, as they
apparently believe they can, once they have a nuclear weapon.''

Asked about Mrs.~Clinton's comments, Sir Nigel Sheinwald, the British ambassador to the United
States, said, ``I don't think it should be read as an acceptance of an Iranian nuclear weapon'' but
rather as a statement intended to ``reassure our partners in the gulf.''

A senior White House official said he believed that Mrs.~Clinton was speaking for herself and that
she was, as she insisted, restating existing policy.

Mrs.~Clinton's invocation of a defense umbrella is reminiscent of the so-called nuclear umbrella
that Washington extends to its Asian allies: implicitly, the promise of an American reprisal if they
are attacked by nuclear weapons. But she did not use the term nuclear, and a senior State Department
official cautioned that her remarks should not be interpreted to mean that.

After meeting the foreign ministers of China, Russia, Japan and South Korea, Mrs.~Clinton also said
that the United States would not offer new incentives to North Korea to return to negotiations. She
said all of the other nations that had engaged in talks with North Korea in the past five years were
united in demanding that North Korea undertake a ``complete and irreversible denuclearization''
before receiving any economic or political incentives from them.

She did not detail the steps that would be part of such a process, though she confirmed that they
could include the disabling of the Yongbyon nuclear complex. Last year, North Korea began to
dismantle that complex, where it runs a nuclear reactor and reprocess fuel rods to recover
plutonium, but it vowed in June to restart production there.

The United States has had an uncharacteristically visible presence at this gathering of the
Association of Southeast Asian Nations, or Asean. It signed a friendship treaty with Asean's 10
members and called on one country, Myanmar, to release the imprisoned pro-democracy leader Daw Aung
San Suu Kyi.

Israel's minister of intelligence and atomic energy, Dan Meridor, told Israeli Army radio: ``I was
not thrilled to hear the American statement from yesterday that they will protect their allies with
a nuclear umbrella, as if they have already come to terms with a nuclear Iran. I think that's a
mistake.''

Mrs.~Clinton said she was trying to make even starker the choice Iran faced if it did not agree to
abandon its program.

The administration has talked about bolstering the military capacity of Iran's neighbors in the
Persian Gulf so they could better meet the threat of a heavily armed Iran. It has also defended the
proposed missile defense system in Eastern Europe as a potential shield against Iran.

``It faces the prospect, if it pursues nuclear weapons, of sparking an arms race in the region,''
Mrs.~Clinton said. ``That should affect the calculation of what Iran intends to do, and what it
believes is in its national security interest.''

On North Korea, Mrs.~Clinton tried to project a united front, saying that China, Russia, Japan and
South Korea had pledged to carry out the United Nations sanctions adopted in June against the North
after its recent nuclear and missile tests.

Mrs.~Clinton also reiterated\cite{reiterate} concerns that North Korea might be transferring nuclear
technology to Myanmar, which American officials refer to by its former name, Burma. She is to
deliver a statement on North Korea on Thursday. In an excerpt provided to reporters, the tone
remained unyielding, but the United States pledged to give North Korea ``significant economic and
energy assistance'' if it undertook a verifiable denuclearization.

At the ministers meeting, Mrs.~Clinton demanded that Myanmar release Mrs.~Aung San Suu Kyi, who
could face years in prison on charges that she violated her house arrest. ``It's so critical that
she be released from this persecution that she has been under,'' she said later at a news
conference. ``If she were released, that would open up opportunities, at least for my country, to
expand our relationship with Burma, including investments in Burma.''

Mr.~Obama extended a ban on American investments in Myanmar in May, but an official said the
president could rescind\cite{rescind} it.

\section{China Fears Ethnic Strife Could Agitate\cite{agitate} Uighur Oasis}

\lettrine{A}{li} the tour guide seemed nice enough and his English flowed with grammatical
perfection -- a useful attribute in a city where most people speak a Turkic language that sounds
nothing like Chinese.

``Sure, I will take you wherever you want to go, but first I have to call my friend and see if he
will drive us,'' Ali said, turning away. After a quick exchange, he hung up the phone and politely
announced that his friend was actually a government minder who would soon be arriving to guide the
would-be clients away from any potential trouble.

The destination his ``friend'' had in mind? The airport, where the reporters, subject to a ban on
foreign journalists, would be escorted\cite{escort} onto the next flight out of town.

``Sorry,'' Ali said as the journalists prepared to flee in a taxi. ``But if I didn't make that call,
I would get in big trouble.''

Kashgar, the ancient Silk Road oasis and backpacker lure, has been besieged by fear since ethnic
rioting about two weeks ago claimed at least 197 lives in Urumqi, the capital of this northwestern
expanse known as the Xinjiang Uighur Autonomous Region. Although the two cities are separated by
about 700 miles of punishing desert and snow-draped mountains, the authorities are especially
anxious about potential unrest in Kashgar, a city of 3.4 million that is 90 percent Uighur, a Muslim
minority that has long had a mercurial\cite{mercurial} relationship with the Han Chinese who govern
Xinjiang.

The authorities have good reason to be skittish\cite{skittish}. Last August, at least 16 military
police officers were killed in an attack here, unnerving the government just as
dignitaries\cite{dignitary} and athletes were arriving in Beijing for the 2008 Olympics. The police
called it a terrorist strike by two Uighur men armed with explosives and machetes, though some
witnesses later challenged that version of events.

In the early 1990s, Kashgar was also the scene of bombings and demonstrations; at least 21 people
were killed and thousands were arrested during one particular army crackdown. The city has long been
a crucible\cite{crucible} for Uighur self-determination, even if nationalist aspirations were never
the same after a Chinese warlord vanquished the newborn East Turkestan Republic, a short-lived
nation that called Kashgar its capital for a few months in 1933.

Although it is rapidly being bulldozed in the name of modernization, Old Kashgar and its ancient
dusty warrens remain the heart of Uighur culture and a beguiling\cite{beguile} draw for tourists. To
China's leadership, however, the city is also an incubator\cite{incubator} for those seeking to
create a Uighur homeland by the borders of Pakistan, Afghanistan and a handful of other
predominantly Muslim countries whose names end with ``stan.''

This time around, Kashgar has been relatively quiet. During the turmoil in Urumqi, a crowd of 200
people tried to protest outside the city's Id Kah Mosque, the largest in China, and were quickly
dispersed by the police, according to Xinhua, the state news agency.

But while journalists in Urumqi can roam with relative freedom, the few foreign reporters who made
it to Kashgar were promptly hustled\cite{hustle} out of town.

``The situation may look calm now, but it could change at any second,'' a local government official
told Mark MacKinnon, a writer for The Globe and Mail, a Canadian newspaper, as he and his colleagues
were sent to the airport.

The uncertainty and sense of isolation have been only magnified by the blocking of access to the
Internet and shutdown of text messaging and international phone service that has severed
communications in Kashgar and the entire region. The blackout has been especially challenging for
export companies, banks, factory owners and academics, and some of them say they have been told that
Internet and phone service will be curtailed\cite{curtail} until at least October, when China
celebrates the 60th anniversary of the Communist revolution.

``I'm expecting a group of Swiss tourists next week, but I have no way of knowing whether they're
still coming,'' said one beleaguered tour operator.

Like Urumqi, which has been flooded with soldiers since July 5, Kashgar is patrolled by young men in
military camouflage, many of whom ride through the city day and night, their green army trucks
draped with ostensibly calming slogans like ``National Separatists Are Our Enemy.''

But the government's most effective weapon against potential trouble is largely unseen: the
neighborhood committees made up of appointed Uighur cadres and citizens who, driven by fear or
ambition, are ready to do the government's bidding.

``You have to be careful because informers are everywhere,'' said Ismail, a secondary school teacher
who used only one name for his own safety. He said his brother had been detained after publicly
criticizing plans to tear down the old mud-and-straw homes that, until recently, flanked\cite{flank}
Kashgar's historic mosque. ``I would not trust anyone if I were you,'' he said.

His words were not hyperbole\cite{hyperbole}. By late last week, hotel clerks, tour guides and taxi
drivers had been instructed to be on the lookout for pesky\cite{pesky} foreign journalists. A woman
employed by a state-owned tourism company told of a meeting during which her boss warned that people
caught assisting reporters would lose their jobs -- as would members of their immediate families.

The campaign appeared to be extremely effective. When his passengers asked to be taken to a rural
county known for its unemployed and disaffected residents, one Uighur driver called the police and
then warned other drivers against helping the passengers escape.

After several close encounters with the authorities, the foreigners made it to the well-irrigated
countryside that forms a lush buffer between Kashgar and the vast Taklimakan Desert stretching 590
miles to the east. In one town, a group of old men hacking at the soil spoke rapturously about the
pace of modernization that had made farming, and their lives, much easier.

``We have electricity, fertilizer and motorbikes now,'' one of the men said.

Later, after some gentle prodding, the farmers allowed that life was not without difficulties. One
man, pointing to a row of unfinished brick houses, said local officials had demolished the
villagers' old homes and promised that the government would pay for the construction of new ones.

``The homes they're building are half as large, and now we have to pay half their cost,'' he said as
his neighbors nodded with disgust. ``We don't have that kind of money.''

The men continued on for a while, speaking animatedly as the tour guide's face registered a
kaleidoscope of troubled expressions. Their ranting done, the guide, a graduate student best left
unidentified, paused before declining to render their words into English.

``I'm sorry,'' he said. ``But it's better for everyone if I just pretend I didn't hear that.''

\section{Concerns on Health Care Reflect Conflicting Goals}

\lettrine{A}{s} Democratic Congressional leaders try to round up the votes to remake the health care
system, they face a range of concerns about the cost and scope of the legislation among centrist
lawmakers in each party whose support is vital to a deal.

Representative Charlie Melancon, Democrat of Louisiana, suspects that a proposed government-run
health insurance plan will be another unaffordable ``entitlement program'' that puts a stranglehold
on the federal budget and on future generations.

Senator Ben Nelson, Democrat of Nebraska, expressed fears that the health care legislation could
inadvertently cause some Americans to lose their existing coverage. And he wants to be sure that any
bill includes new ``incentives for healthy lifestyles.''

Representative Mike Ross, Democrat of Arkansas, wants to reduce or eliminate the disparities in what
Medicare pays for health services in rural areas versus urban centers, a gap that he says has forced
hospitals to close and doctors to move away.

Senator Susan Collins, Republican of Maine, worries about requiring employers to provide insurance
to workers. Without categorically rejecting the idea, she said: ``My inclination is to oppose an
employer mandate. Employers want to provide health insurance and do so if they can afford it.''

Many of the centrists said they shared the same concerns: that the legislation proposed so far is
too expensive; does not sufficiently reduce health care costs over the long term; and would raise
taxes too much, or in ways they oppose.

If there is one thing centrists in the House and the Senate agree on, it is that they are being
pushed way too fast to act on a hugely complex bill with an astronomical price tag of roughly \$1
trillion over 10 years, prompting a loud chorus of demands to slow down.

But aside from taking more time, and missing President Obama's deadline of having each chamber pass
a health care bill before the August recess -- a goal that seems impossible to meet -- getting to
``yes'' will require an artful set of compromises.

In many cases, the concerns of centrist lawmakers reflect conflicting goals: to expand health
coverage to nearly all Americans while reducing the growth of health spending, cutting the cost of
the bill and minimizing new taxes.

In the House, a group of seven Democrats who are members of the fiscally conservative Blue Dog
Coalition have stalled the proceedings of the Energy and Commerce Committee, the last of three House
panels needed to approve the bill.

The House speaker, Nancy Pelosi, said Wednesday that she shared some of the group's concerns but
also said, ``I have no question that we have the votes on the floor of the House to pass this
legislation'' -- a remark that some took as a threat to barrel ahead.

``Why are we rushing it?'' asked Mr.~Melancon, a Blue Dog on the Energy and Commerce Committee.
``Let's get it right.''

And he effectively dared House leaders to press ahead without fully addressing his group's concerns,
noting that the Blue Dogs potentially had enough votes in Energy and Commerce to stall the bill and
that leaders would have to resort to pressure tactics.

``We're going to need some orthopedists around here to take care of the broken bones and twisted
arms,'' he said.

The best hope for a bipartisan compromise seems to lie with a small group of lawmakers in the Senate
Finance Committee, which continued to meet on Wednesday.

The committee chairman, Senator Max Baucus, Democrat of Montana and a centrist with a long history
of collaborating with Republicans, has resisted pressure from Senate leaders and the White House
while insisting that a compromise is within reach.

Republicans on that committee, including Senator Charles E.~Grassley of Iowa, are insisting that any
new taxes to help pay for the legislation come from within the health care arena. They have
expressed no interest in the income surtax, on high-earners, proposed by the House that would raise
\$544 billion over 10 years and instead are exploring a menu of smaller items including fees on
private insurers.

Last week, a bipartisan group of senators, including two on the Finance Committee, sent a letter to
the Senate leadership pleading for more time. A number of other senators have said they want the
panel to continue working carefully toward a deal.

Ms.~Collins, who joined Democrats in supporting the economic stimulus measure this year, said she
had several big concerns about the House health care bill and a separate measure approved last week
on a party-line vote by the Senate Committee on Health, Education, Labor and Pensions.

She said that she was ``very skeptical'' of the new government health insurance plan that Democrats
want to create, and that she was ``not a proponent of an employer mandate,'' which would require
employers to provide health benefits to workers or pay a fee to the government.

Instead, she said, ``I would like to encourage more employer coverage by giving generous tax credits
to smaller businesses.''

As for reaching a deal, Ms.~Collins said, ``we have a long way to go.''

The Senate Finance Committee has been working on a compromise idea that would create private
nonprofit health insurance cooperatives instead of a government-run insurance plan to compete with
private insurers in an effort to reduce costs.

Mr.~Ross, of Arkansas, said the House Blue Dogs were open to considering the co-op idea and any
number of other proposals that might reduce the cost of the legislation while still expanding
insurance coverage to most Americans.

At a meeting at the White House on Tuesday, the Blue Dogs said, they won agreement from the Energy
and Commerce Committee chairman, Representative Henry Waxman of California, to try to curtail
Medicare spending by giving the executive branch new power to set annual payment rates for doctors,
hospitals and other health care providers, based on recommendations from an independent advisory
council.

Rates are now set by statutory formulas, and Congress is besieged by lobbyists pleading for bigger
increases each year.

But Mr.~Ross said that addressing long-term costs was just one area of concern for the Blue Dogs.
They also want to reduce or eliminate the disparity in Medicare payments that result in lower rates
for health providers in rural areas.

Mr.~Ross said his hometown, Prescott, now has three doctors -- down from six -- and that private
insurers in Arkansas pay 30 percent more than Medicare, making doctors reluctant to accept new
Medicare patients.

\section{Daring to Dream of a Resurgent AOL}

\lettrine{S}{hortly} after Tim Armstrong took over as chief executive of AOL, he asked to see the
list of business deals that were being negotiated. He saw 900 of them.

It was too many by far. ``If you looked through the deal sheet, would you have been able to see the
strategy of the company?'' he asked. ``I had a hard time.''

The deals were small and incremental. At best, he said, ``you would have thought it was a small- to
medium-size Internet company.''

Mr.~Armstrong wants AOL to think big again. Three months after leaving a senior job as Google's
president of advertising sales, he is formulating his ambitious recovery plan for AOL. He wants to
make AOL the biggest creator of premium\cite{premium} content on the Web and the largest seller of
online display advertising.

Mr.~Armstrong plans to outline his five-point strategy on Friday for the company at an all-hands
meeting under a large tent on the its half-empty campus near Dulles International Airport outside
Washington. Beyond talking about business lines, however, Mr.~Armstrong's primary challenge is to
address what he calls AOL's ``crisis of confidence.'' He wants the weary and beaten-down company to
grow again.

``AOL has a choice to make,'' he said. ``We either lose slowly or win quickly. We are choosing to
win quickly.''

Nine and a half years after Steve Case combined the company with Time Warner, AOL suffers from
myriad\cite{myriad} problems. It has long since lost the mantle of king of the Internet to
Mr.~Armstrong's former employer, Google. It has suffered through wrenching waves of mass layoffs,
management turmoil\cite{turmoil} and constant bickering\cite{bickering} with its corporate parent.
Time Warner plans to shed the unit by year end. Meanwhile, AOL struggles with the prospect of fading
into irreversible irrelevance, a collection of tired brands for a shrinking core of customers
hanging on mainly because they are too lazy to change their AOL.com e-mail addresses.

This year, AOL is expected to post revenue of about \$3.2 billion, down 38 percent in two years. A
majority of that revenue is advertising, but AOL's 6.2 million remaining customers for its dial-up
Internet service are highly profitable and the most avid readers of its content. About 200,000 of
them cancel service every month.

Mr.~Armstrong and a core group of managers brought from Google are trying ``to change the DNA of the
company,'' in the words of Jeffrey A.~Levick, a longtime aide to Mr.~Armstrong who is now the
president of AOL's advertising unit.

``People put numbers on the board and I say, 'You are missing a lot of zeros and a lot of commas,'
'' he said. ``I don't think people have talked that way around here since the days of Steve Case.''
Investors just see zeros when they think about the potential value of an independent AOL.

``Expectations from myself and Wall Street for AOL are still dire,'' said Richard Greenfield, an
analyst with Pali Capital. Still, he said the choice of Mr.~Armstrong to run the company ``is far
better than we would have expected.'' And if Mr.~Armstrong leads even a modest turnaround in AOL, it
``could surprise people and lead to substantial upside.''

The market value of AOL after the spinoff to Time Warner shareholders will depend on how much debt
Time Warner saddles\cite{saddle} it with. Mr.~Greenfield said AOL may be worth \$2 billion to \$3
billion, far less than its \$20 billion valuation in 2005, when Google invested \$1 billion in it --
a deal Mr.~Armstrong helped negotiate.

Mr.~Armstrong has long thought big. Mr.~Levick worked at an advertising agency in Chicago when he
met Mr.~Armstrong, who was selling an early form of advertising for a tiny search engine named
Google.

``Here was a man standing at a whiteboard drawing the picture of all advertising all coming in
through one place, Google,'' Mr.~Levick said. ``I don't think even anyone saw how big this was going
to be, but Tim's plan was for this to be bigger than anyone's wildest imagination, even the people
at Google.''

Despite having the wealth that comes with being one of Google's early employees, Mr.~Armstrong, 38,
said he was lured\cite{lure} to AOL create a new kind of media company. ``One of the biggest
challenges in the media business is also one of the biggest opportunities,'' he said. ``If you tried
to recreate AOL's assets, it would be incredibly expensive.''

Before Mr.~Armstrong can move forward with his strategy, he must stabilize AOL's ranks. In the last
three years, the company has had three chief executives and five heads of ad sales.

Mr.~Armstrong's regular Tuesday product reviews include the people who build and operate each
product as well as their bosses. The all-hands meeting on Friday, which will be Webcast to all
employees, will be the third such event he has held.

Mr.~Armstrong's five-point plan was also determined by a collaborative process in a two-day meeting
in New York. The top 100 employees sorted through three dozen current and potential business lines
for the company.

Mr.~Armstrong asked simply, ``What can we win?'' Among the hot debates were what AOL's role should
be in social networking and in search.

Eventually, the assembled employees voted on their top five ideas. Separately, Mr.~Armstrong wrote
his top five on a blackboard, turning it so the audience could see it only after the vote. The only
difference: Mr.~Armstrong wanted to include AOL's Truveo video search company in the top priorities.
But he deferred to the group and assigned Truveo instead to a new unit called AOL Ventures, where he
is putting noncore businesses, like the Bebo social network, that might eventually be sold.

Mr.~Armstrong's plan is to compete directly with Yahoo, Microsoft and Google to become the dominant
network for display ads. Mr.~Armstrong says the company's technology, with the data it has on
millions of consumers accumulated over nearly 25 years, will give it an edge\cite{edge}.

``Nobody owns the display space today,'' he said.

Where Mr.~Armstrong will find much less competition is content, which most Internet companies find
too expensive to produce. AOL already operates more than 70 specialized blogs, including Boombox (on
hip-hop music), WalletPop (on personal finance) and Paw Nation (on pets).

Mr.~Armstrong wants to cover more topics in more countries with much more video.

The combination of specialized content and display advertising, he said, should make AOL appealing
to large consumer products companies with big marketing budgets like Procter \& Gamble.

``If you ask P.\& G.~what companies have the products that make you feel most comfortable, with the
best content and the best targeting, AOL is already on the list today,'' Mr.~Armstrong said. ``Our
aim is to move AOL to the top of the list.''

\section{Amazon's Expanding With Deal for Zappos}

\lettrine{A}{mazon}.com is trying the shoe retailer Zappos.com on for size.

The companies announced on Wednesday that Amazon was acquiring Zappos, based in Henderson, Nev., for
10 million shares of Amazon stock, worth nearly \$900 million at its current level.

Amazon also said it would give Zappos employees an additional \$40 million in cash and stock.

The acquisition is the largest in Amazon's history.

Online shoe sales hit \$4.3 billion last year, up from \$3.3 billion in 2007, according to Forrester
Research.

Zappos, a private company founded 10 years ago, has won fans with perks like free shipping and
personalized service. It is the largest player in that market.

Amazon has tried without much success to burnish its shoe offerings in the face of competition from
Zappos. In 2007, it introduced a separate site, Endless.com, to sell shoes and handbags. But while
Zappos received 4.5 million visitors in June, Endless.com got 777,000, according to comScore.

Amazon also sells shoes on its main site and allows outside retailers to sell them there as well.

``I don't think Amazon has gotten much traction'' with Endless, said Sucharita Mulpuru, an analyst
at Forrester. ``People will go to Amazon to purchase a lot of their core media products and consumer
electronics. But it is not thought of in the same light when shopping for apparel, accessories and
footwear.''

Zappos, Ms.~Mulpuru said, ``is one of the few sites out there that people think of first when they
think of footwear.''

The companies said that the management team of Zappos would remain intact and that Zappos would
continue to operate as an independent division of Amazon from its offices outside Las Vegas. Zappos,
whose backers include the venture capital firm Sequoia Capital, also said it planned to maintain its
shipping facility in Louisville, Ky.

The companies had held conversations about a deal over the years, but they took on a serious tone a
couple of months ago, a person briefed on the matter said. Amazon initially sought to pay cash, but
Zappos asked for an all-stock deal, this person said.

The extra cash and restricted stock for employees is meant to keep them on board and preserve the
company's culture, the person said. The deal is expected to close in the fall.

On the Zappos blog on Wednesday, Tony Hsieh, Zappos' chief executive, said that the decision to
accept Amazon's offer was a difficult one. This was in part because Amazon focuses on low prices and
convenience, while Zappos has carved out a favorable reputation among consumers for personalized
service, free overnight shipping and a policy allowing buyers to return any pair of shoes free.

``We plan to continue to run Zappos the way we have always run Zappos -- continuing to do what we
believe is best for our brand, our culture and our business,'' Mr.~Hsieh wrote.

Jeffrey P.~Bezos, Amazon's chief executive, said in a video clip about the deal that Zappos ``has a
customer obsession, which is so easy for me to admire.''

He added, ``I get all weak-kneed when I see a customer-obsessed company.''

Zappos appears to engender friendly feelings even among some of its smaller competitors.

Korey Buzzell, who runs the independent site Shoe-Store.net, said Zappos had been an amicable
competitor in the past, sending customers to his site when it could not fulfill their orders.

``They have sent a lot of business our way over the years, and in return I myself refer people to
Zappos if I can't help them,'' he said. ``Amazon is different. I don't use Amazon that much for
shoes, in fact I try to avoid them. I hope Zappos doesn't change with an acquisition like this.''

Amazon was advised in the deal by Lazard, and Zappos by Morgan Stanley.

\section{Mom-and-Pop Operators Turn to Social Media}

\lettrine{T}{hree} weeks after Curtis Kimball opened his cr\`eme br\^{u}l\'ee cart in San Francisco,
he noticed a stranger among the friends in line for his desserts. How had the man discovered the
cart? He had read about it on Twitter.

For Mr.~Kimball, who conceded that he ``hadn't really understood the purpose of Twitter,'' the
beauty of digital word-of-mouth marketing was immediately clear. He signed up for an account and has
more than 5,400 followers who wait for him to post the current location of his
itinerant\cite{itinerant} cart and list the flavors of the day, like lavender\cite{lavender} and
orange creamsicle.

``I would love to say that I just had a really good idea and strategy, but Twitter has been pretty
essential to my success,'' he said. He has quit his day job as a carpenter to keep up with the
demand.

Much has been made of how big companies like Dell, Starbucks and Comcast use Twitter to promote
their products and answer customers' questions. But today, small businesses outnumber the big ones
on the free microblogging service, and in many ways, Twitter is an even more useful tool for them.

For many mom-and-pop shops with no ad budget, Twitter has become their sole means of marketing. It
is far easier to set up and update a Twitter account than to maintain a Web page. And because
small-business owners tend to work at the cash register, not in a cubicle\cite{cubicle} in the
marketing department, Twitter's intimacy\cite{intimacy} suits them well.

``We think of these social media tools as being in the realm of the sophisticated, multiplatform
marketers like Coca-Cola and McDonald's, but a lot of these supersmall businesses are gravitating
toward them because they are accessible, free and very simple,'' said Greg Sterling, an analyst who
studies the Internet's influence on shopping and local businesses.

Small businesses typically get more than half of their customers through word of mouth, he said, and
Twitter is the digital manifestation of that. Twitter users broadcast messages of up to 140
characters in length, and the culture of the service encourages people to spread news to friends in
their own network.

Umi, a sushi restaurant in San Francisco, sometimes gets five new customers a night who learned
about it on Twitter, said Shamus Booth, a co-owner.

He twitters about the fresh fish of the night -- ``The O-Toro (bluefin tuna belly) tonight is some
of the most rich and buttery tuna I've had,'' he recently wrote -- and offers free seaweed salads to
people who mention Twitter.

Twitter is not just for businesses that want to lure customers with mouth-watering descriptions of
food. For Cynthia Sutton-Stolle, the co-owner of Silver Barn Antiques in tiny Columbus, Tex.,
Twitter has been a way to find both suppliers and customers nationwide.

Since she joined Twitter in February, she has connected with people making lamps and candles that
she subsequently ordered for her shop and has sold a few thousand dollars of merchandise to people
outside Columbus, including to a woman in New Jersey shopping for graduation gifts.

``We don't even have our Web site done, and we weren't even trying to start an e-commerce
business,'' Ms.~Sutton-Stolle said. ``Twitter has been a real valuable tool because it's made us
national instead of a little-bitty store in a little-bitty town.''

Scott Seaman of Blowing Rock, N.C., also uses Twitter to expand his customer base beyond his town of
about 1,500 residents. Mr.~Seaman is a partner at Christopher's Wine and Cheese shop and owns a bed
and breakfast in town. He sets up searches on TweetDeck, a Web application that helps people manage
their Twitter messages, to start conversations with people talking about his town or the mountain
nearby. One person he met on Twitter booked a room at his inn, and a woman in Dallas ordered sake
from his shop.

The extra traffic has come despite his rarely pitching his own businesses on Twitter. ``To me,
that's a turn-off,'' he said. Instead of marketing to customers, small-business owners should use
the same persona they have offline, he advised. ``Be the small shopkeeper down the street that
everyone knows by name.''

Chris Mann, the owner of Woodhouse Day Spa in Cincinnati, twitters about discounts for massages and
manicures\cite{manicure} every Tuesday. Twitter beats e-mail promotions because he can send tweets
from his phone in a meeting and ``every single business sends out an e-mail,'' he said.

Even if a shop's customers are not on Twitter, the service can be useful for entrepreneurs, said
Becky McCray, who runs a liquor store and cattle ranch in Oklahoma and publishes a blog called Small
Biz Survival.

In towns like hers, with only 5,000 people, small-business owners can feel isolated, she said. But
on Twitter, she has learned business tax tips from an accountant, marketing tips from a consultant
in Tennessee and start-up tips from the founder of several tech companies.

Anamitra Banerji, who manages commercial products at Twitter, said that when he joined the company
from Yahoo in March, ``I thought this was a place where large businesses were. What I'm finding more
and more, to my surprise every single day, is business of all kinds.''

Twitter, which does not yet make money, is now concentrating on teaching businesses how they can
join and use it, Mr.~Banerji said, and the company plans to publish case studies. He is also
developing products that Twitter can sell to businesses of all sizes this year, including features
to verify businesses' accounts and analyze traffic to their Twitter profiles.

According to Mr.~Banerji, small-business owners like Twitter because they can talk directly to
customers in a way that they were able to do only in person before. ``We're finding the emotional
distance between businesses and their customers is shortening quite a bit,'' he said.

\section{Chinese Worker Commits Suicide Over Missing iPhone Prototype}

\lettrine{C}{hinese} worker Sun Danyong was responsible for handling the prototypes of one of the
world's hottest products -- the iPhone. When one of the gadgets went missing and his company began
investigating him, he jumped off his apartment building and killed himself.

The death -- which involves allegations\cite{allegation} that security guards roughed up the worker
-- prompted Apple Inc.~on Wednesday to issue a terse\cite{terse} statement, insisting that all the
company's contractors must treat workers with respect and dignity.

The 25-year-old Sun started his new job last year after earning a degree in business management. He
moved from his native Yunnan province -- a poor region sharing a border with Myanmar -- to the
southern boomtown of Shenzhen, home to the sprawling\cite{sprawl} factory complex run by Foxconn
Technology Group. The Taiwanese manufacturer has long been one of Apple's key suppliers.

Sun's job involved shipping iPhone prototypes to Apple. It was an extremely sensitive position for a
company like Apple, known for shrouding its new product launches in secrecy and suspense -- a
strategy that's consistently helped whip\cite{whip} up the just-can't-wait-to-buy-it feeling among
consumers worldwide.

Although Apple and Foxconn confirmed Sun's suicide, they would not provide much information about
the circumstances of his death. Many details have been reported by the state-run Southern Metropolis
Daily, one of the region's most popular and aggressive newspapers.

The paper's account, which hasn't been disputed by the companies, said: Sun reported on July 13 that
he was missing one of the 16 fourth-generation iPhones in his possession. Foxconn security guards
searched his apartment, detained him and beat him. In the early morning of July 16, a distraught Sun
jumped from the 12th floor of his apartment building.

Jill Tan, an Apple spokeswoman in Hong Kong, issued only a brief statement about the incident.

``We are saddened by the tragic loss of this young employee, and we are awaiting results of the
investigations into his death,'' Tan said. ``We require our suppliers to treat all workers with
dignity and respect.''

Apple's success amid\cite{amid} the recession is due in part to the iPhone. More than 5.2 million of
the devices were sold in the third quarter this year -- seven times the number sold at the same time
last year, the company said Tuesday when announcing its earnings. The sales spike was related to a
newly released version of the device, it said.

The Cupertino, Calif.-based company's earnings jumped 15 percent in the third quarter -- growth
largely propelled by laptop and iPhone sales.

One of Apple's most important manufacturing partners has long been Foxconn, owned by Taiwan's Hon
Hai Precision Industry Co.~-- the world's biggest contract manufacturer of electronics. The
corporate behemoth\cite{behemoth} has also produced computers for Hewlett-Packard Co., PlayStation
game consoles for Sony Corp.~and mobile phones for Nokia Corp.

Foxconn executive Li Jinming said in a statement that Sun's death showed that the company needed to
do a better job helping its employees with psychological pressures.

``Sun Danyong graduated from a good school. He joined the company in 2008. He had an extremely
bright future. The group and I feel deep pain and regret when a young person dies like this.''

The company also said its security chief has been suspended and turned over to the police.

The police declined to respond to questions from The Associated Press.

Foxconn's security chief, Gu Qinming, was quoted by the Southern Metropolis Daily as saying he never
hit Sun. Gu reportedly said that after three security personnel searched Sun's apartment and did not
find the phone, the employee followed orders to go to Gu's office on July 15.

``I got a bit agitated. I pointed my finger at him and said that he was trying to shift the blame,''
Gu was quoted as saying.

He added, ``I was a little angry and I pulled his right shoulder once to get him to tell me what
happened. It (the beating) couldn't have happened,'' the paper reported.

\section{Clinton Trades Jibes\cite{jibe} With North Korea}

\lettrine{T}{he} United States and North Korea fell into an acrimonious\cite{acrimonious} exchange
on Thursday, with the North Korean government ridiculing Secretary of State Hillary Rodham Clinton
as a ``schoolgirl'' and a ``pensioner,'' two days after she referred to its leaders as ``unruly
children.''

At a meeting of Southeast Asian nations here, the war of words competed for attention with
Mrs.~Clinton's campaign to marshal worldwide pressure on the North Koreans to
dismantle\cite{dismantle} their nuclear weapons program.

On Thursday, the Foreign Ministry in North Korea issued a statement criticizing remarks Mrs.~Clinton
made this week to ABC News, in which she said the best response to North Korea's behavior would be
to ignore it, as one would a child clamoring for attention.

``We cannot but regard Mrs.~Clinton as a funny lady, as she likes to utter such
rhetoric\cite{rhetoric}, unaware of the elementary etiquette\cite{etiquette} in the international
community,'' the North Korean statement said. ``Sometimes she looks like a primary schoolgirl and
sometimes a pensioner going shopping.''

North Korea said it would defend its sovereignty\cite{sovereignty} against the United States, which
it accuses of aiming nuclear weapons at it.

The ill will surfaced vividly during a meeting of the Association of Southeast Asian Nations, or
Asean, at this Thai resort when the North Korean delegation turned up on a hotel podium\cite{podium}
to deliver a scheduled statement.

The officials were told that Mrs.~Clinton was due to speak soon, though she was running late, as she
has often this week. After huddling, they stalked away and held a news conference nearby, at which
they reiterated that North Korea would never return to multiparty talks with South Korea, Japan,
Russia, China and the United States, talks that are aimed at curbing its nuclear ambitions.

For her part, Mrs.~Clinton said she was encouraged by the international support for pressuring North
Korea. Even Myanmar, she said, responded to requests by China and other countries to track a North
Korean freighter this month that American officials suspected was carrying illicit cargo.

``The international community's response to North Korea's actions has been unequivocal and nearly
unanimous, leading to a new consensus,'' Mrs.~Clinton said at a news conference, during which she
read a lengthy statement restating the American policy on North Korea.

She said there was a commitment to carry out the sanctions called for in a United Nations Security
Council resolution adopted in June after North Korea's recent arms tests. Among its measures, the
resolution bans weapons shipments to North Korea and seeks to squeeze the sources of financing for
its nuclear and missile programs.

Mrs.~Clinton singled out China, an influential neighbor, for asking officials in Myanmar, formerly
Burma, to help in dealing with the North Korean freighter, the Kang Nam 1, which was steaming toward
Myanmar. The vessel eventually turned around on its own, and she called China's pressure a
``proximate cause.''

She said she would discuss further steps in pressing North Korea with senior Chinese officials
during consultations with China next week in Washington. Mrs.~Clinton is leading the strategic and
economic dialogue jointly with Treasury Secretary Timothy F.~Geithner.

The United States has talked to China and other countries about a package of incentives, including
economic and energy aid, which could be offered to North Korea in return for dismantling its nuclear
program.

But even before Thursday's vitriolic\cite{vitriolic} statements from North Korea, American officials
said they were more focused for now on inflicting pain on North Korea than on luring it back to the
bargaining table.

``We are not interested in half measures,'' Mrs.~Clinton said. ``We have no desire to pursue
protracted negotiations that will only lead us right back to where we have already been.''

Mrs.~Clinton said the North Koreans had been intransigent in their public statements during the
conference. Other senior American officials said the tone of the North's statements was openly
hostile.

Still, Mrs.~Clinton may have contributed to the chilly atmosphere in her remarks just before the
meeting. ``Maybe it's the mother in me,'' she said to ABC News, ``the experience I've had with small
children and teenagers and people who are demanding attention: don't give it to them.''

\section{Microsoft Revenue Falls 17\%}

\lettrine{M}{icrosoft}, the once-swaggering giant of the personal computer industry, has been
humbled, both by the recession and by problems of its own making.

On Thursday, the world's largest software company reported its worst fiscal year since it initially
sold stock to the public in 1986. Year-over-year revenue and full-year sales of Microsoft's flagship
Windows software dropped for the first time.

``Clearly, Microsoft is not immune to the economic downturn,'' said Brendan Barnicle, a software
analyst with Pacific Crest Securities.

Many prominent companies tied to the PC industry have watched about one-fourth of their revenue
vanish as business customers in particular have scaled back new PC purchases. Microsoft makes more
money from versions of Windows tied to business computers than it does on cheaper machines aimed at
consumers. If businesses start buying again, Microsoft should benefit from higher overall sales and
rising profits.

But the economy's pinch is not solely to blame for Microsoft's problems. The company's Windows Vista
software, hailed in 2007 as the most significant product in the company's history, has failed to
attract businesses in any meaningful way because of problems with compatibility and speed.

According to a study by Forrester Research, 86 percent of corporate PCs continue to rely on the
eight-year-old Windows XP.

Microsoft's Windows profits have also fallen because of rising interest in the cheap, compact
laptops known as netbooks, which rely on the lower-priced Windows XP instead of Vista.

Earlier this month, Google revealed plans to sell a direct competitor to Windows for the netbook
market, and major PC makers have backed Google in what is most likely a bid to put even more pricing
pressure on Microsoft.

The big question haunting Microsoft is when sales of its old-line products like Windows and Office
will come back and fuel the company's aggressive strikes into areas like search, mobile device
software and online business applications.

In its fourth quarter, Microsoft's net income fell to \$3.05 billion, or 34 cents a share, for the
period ended June 30. The figures represent a 29 percent drop in net income from the \$4.30 billion,
or 46 cents a share, that Microsoft reported in the period a year ago.

Excluding charges tied to legal matters, layoffs and investments, Microsoft earned 36 cents a share
in the quarter, meeting the forecast of analysts surveyed by Thomson Reuters.

Microsoft's 17 percent drop in quarterly revenue was more troubling to Wall Street, which pushed the
company's shares down more than 7 percent, to \$23.70, in after-hours trading. The company took in
\$13.10 billion in the quarter, missing analysts' estimate of revenue by \$1.27 billion.

For the full year, Microsoft's revenue declined 3 percent, to \$58.44 billion, while its net income
fell 18 percent, to \$14.57 billion.

On a positive note, Microsoft echoed recent comments from Intel, saying that there had been an
increase in PC sales and healthier spending in Asia and the United States.

``There are some signs that we have at least seen the worst,'' said Christopher P.~Liddell, the
chief financial officer at Microsoft, in a conference call.

In numerous public appearances, Microsoft's chief executive, Steven A.~Ballmer, has warned that
people should not expect a major bounce-back in technology spending when the economy recovers.
Rather, he suggested that a new, low bar had been set and that companies needed to adjust to such a
reality.

Microsoft's recent missteps have resulted in a more cautious public stance. In a recent speech,
Mr.~Ballmer was equivocal about whether the company's new Windows 7 software, to be introduced in
October, would prompt a surge in PC and software sales. ``Maybe it will. Maybe it won't,'' he said.

``They want to keep expectations tame\cite{tame} in front of Windows 7,'' said Israel Hernandez,
director of software research at Barclays Capital. ``They want to keep things in check until they
have better visibility on the economy.''

Mr.~Ballmer's most boisterous\cite{boisterous} recent comments have revolved around Microsoft's
staying power. He has vowed that the company will keep pouring money into research and development
until it gets search, mobile software and virtualization software right. Microsoft spends close to
\$10 billion a year on research -- more than any other technology company.

Analysts contend that Microsoft has survived the downturn well enough, posting profits at a time
when losses were more expected. They say the company seems destined to capitalize on large-scale PC
upgrades, as companies finally give into the pressure of buying new machines.

``They are at a great stage right now,'' said Mr.~Barnicle.

In the recession, Microsoft has laid off thousands of people and sought to cut far deeper than usual
into its cost structure. ``In my mind, we are a stronger company than we were a year ago,''
Mr.~Liddell said.

If the economy recovers as hoped next year, Microsoft appears poised\cite{poised} to find its
``mojo,'' as Mr.~Ballmer describes it. The company will have revamped its traditional products and
will be releasing a new wave of online business software.

Microsoft already appears confident enough to revisit a search and advertising combination with
Yahoo, which Microsoft failed to acquire outright after extensive discussions last year. Talks about
a partnership have intensified in recent weeks, and some people at both companies say they believe a
deal could be near.

\section{Pi\"ech, VW's Chairman, Is Seen as Master Tactician\cite{tactician}}

\lettrine{T}{he} ouster\cite{ouster} of Porsche's chief executive, Wendelin Wiedeking, underlines
the tactical skills of Ferdinand Pi\"ech, Volkswagen's chairman, whose ability to outmaneuver rivals
and allies has seen him steadily increase his influence at both companies during a 46-year career in
the auto industry.

It also marks a crucial step toward adding a 10th brand to an automotive empire built by Mr.~Pi\"ech
that includes low-end, mass-market and luxury carmakers as well as trucks, among them the VW,
Bentley, Bugatti, Skoda, Audi and Scania brands.

Mr.~Pi\"ech's career might then eclipse that of his illustrious\cite{illustrious} grandfather,
Ferdinand Porsche, who founded Porsche and designed the original VW Beetle.

Mr.~Pi\"ech has often been underestimated because he tends to keep employees and business partners
in the dark, leaving them to guess his next move, say people who have worked for him.

While he was the development chief at Audi, he had different teams work on the aerodynamics of the
Audi 100, using wind tunnels in Hamburg, Stuttgart, Wolfsburg and Turin to complete parts of the
development. He had the test results sent to Ingolstadt, where Audi is headquartered.

``I was in the middle of it all, putting together the pieces of the puzzle,'' he said in his
autobiography.

Gaining control of Porsche would also mark a return of Mr.~Pi\"ech's management influence there,
where he spent nine years before quitting as board member for research and production in 1971.

At the time, Porsche's owning families, his relatives, decided they did not want a family member to
lead the Stuttgart-based maker of the iconic 911 sports car.

Although Mr.~Pi\"ech's aspirations to lead Porsche were dashed in the 1970s, he went on to hold
senior positions at Audi and Mercedes-Benz before rising to become Volkswagen's chief executive in
1993, and eventually chairman in 2002.

The soft-spoken Mr.~Pi\"ech is unfazed by tension, and often encourages rivalries, people who have
worked for him say.

``It's not possible to take a company to the top by focusing on the highest level of harmony,'' he
said in his autobiography.

Because the Pi\"ech clan controls almost half the votes in Porsche, and the owning families need 75
percent approval for major decisions, Mr.~Pi\"ech's side of the family retains a veto even at
Porsche.

Mr.~Pi\"ech also demonstrated Volkswagen's control over Porsche in 2008 by helping to block a
cooperation deal that would have given Porsche access to VW technologies. Mr.~Pi\"ech later voted in
favor of the cooperation, but only after infuriating his cousin, Wolfgang Porsche, the Porsche
chairman.

Porsche is heavily dependent on Volkswagen, which delivers or assembles components for a third of
all Porsche cars, including the body and engine for the four-door Panamera and Cayenne luxury
models.

Mr.~Pi\"ech is widely admired for his intimate knowledge of engineering. He helped to revive
Volkswagen by developing a modular construction system that reduced costs by allowing several
different cars to share the same underpinnings.

But Mr.~Pi\"ech clashed\cite{clash} with shareholders and some of his own managers for putting the
brakes on cost-cutters and betting instead on quality improvements and technical innovations in the
hope that higher profit would follow.

Mr.~Pi\"ech has relied on the support of labor leaders to keep his career alive at VW, sometimes
even at the expense of executives he had appointed.

Wolfgang Bernhard, a former executive of Daimler and Chrysler who once headed the VW brand, and
Bernd Pischetsrieder, once the VW chief executive, both left the company after Mr.~Pi\"ech backed
union leaders who sought to limit cost-cutting measures at VW.

Under Mr.~Pi\"ech's watch, VW acquired the Bentley, Bugatti and Lamborghini brands in a single year.

He also pushed the VW brand into the upscale market by developing the VW Phaeton and the
unprofitable Bugatti Veyron, a move that put him on a collision course with Mr.~Wiedeking, who
thought Mr.~Pi\"ech was squandering VW's money on ``vanity projects.''

\section{Porsche Chief Was Biggest Casualty\cite{casualty} of His Battle for VW}

\lettrine{C}{omfortably} in control of his own company and angling for the much larger Volkswagen,
Porsche's chief executive, Wendelin Wiedeking, described the process before him in November 2007 as
an intense chess match.

``Even now, none of the players involved knows how many moves will still be required and how long it
will take until the game is finally over,'' Mr.~Wiedeking said during an interview at the time.

But Mr.~Wiedeking's metaphor proved far too playful for what transpired, as financial markets and
feuding families took their toll on his attack plan. He had started a war, and ended up as its
biggest casualty.

On Thursday, Porsche's supervisory board fired Mr.~Wiedeking and paved the way for a merger with
Volkswagen, the company Porsche had dreamed of owning.

The departure of Mr.~Wiedeking, 56, marked a tumultuous\cite{tumultuous} end for the best-paid chief
executive in Germany, a hyper-aggressive manager who ruffled feathers for years but nonetheless
showed considerable industrial prowess.

In an ironic twist, Mr.~Wiedeking also anticipated a core problem of car companies worldwide -- a
lack of economies of scale -- but his solution of Porsche taking over Volkswagen became his undoing.

He was constantly in the news in Germany, not the least because he earned \texteuro77.4 million, or
\$110 million, last year, a princely sum in a country that was skeptical of high pay packages even
during good times. But his contract guaranteed him 0.9 percent of the company's profit, and in the
past few years it was greater than anyone could have anticipated. The provision was granted in the
1990s, when Porsche's future was in doubt.

Porsche said Thursday that Mr.~Wiedeking would walk away with a \texteuro50 million severance
package, only days after rumors suggested that the figure might be four times that.

Mindful of the austere times, Mr.~Wiedeking announced that he would put \texteuro25 million into a
charitable trust that would promote ``socially fair development'' at Porsche facilities. He also
planned to donate \texteuro500,000 each to three organizations that look after infirm journalists.

A mechanical engineer by training, he joined Porsche in 1983 at age 31. He left for an auto-parts
maker in 1988 but returned to Porsche as head of production in 1991 and began whipping the nearly
bankrupt company back into shape. At one point he made a show of destroying a bin for extra parts
along its assembly line to make his point about the need for leaner, Japanese-style production
processes.

He became chief executive a year later and transformed the company into one of the world's most
profitable automakers, with annual production of around 100,000 vehicles. He dropped the
money-losing 928 and 968 models, overhauled the iconic 911 and developed two new models, the Boxster
convertible and the Cayenne sport utility vehicle.

Mr.~Wiedeking's prot\'eg\'es rose as well. Michael Macht, who will now succeed Mr.~Wiedeking, ran
the production lines and was the first head of Porsche Consulting, an enterprise that began by
evangelizing Porsche suppliers on the need to cut costs.

In 2005, long before the industry's current consolidation drive, Mr.~Wiedeking had huddled in
Salzburg, with members of the two families that control the sports car maker -- the Porsches and the
Pi\"echs -- and floated the idea of buying VW, on the theory that only a large company could afford
the huge investments in new technologies that environmental regulations and fuel-saving imperatives
would demand.

Initially shocked at his audacity, the families eventually bought in.

But Mr.~Wiedeking's strategy looked rash rather than bold after financial market chaos last year
transformed the landscape for heavily indebted companies.

Though the automaker that Mr.~Wiedeking had so painstakingly rebuilt responded flexibly to the sharp
drop-off in sales, a \texteuro9 billion debt load proved unbearable. Ferdinand Pi\"ech, a member of
Porsche's founding family, a board member, and also chairman of Volkswagen, pounced on the
opportunity to reverse the terms of Porsche's audacious bid. But he insisted that Mr.~Wiedeking
would have to go, and that Porsche would have to bring in an outside investor. Along with his close
confidant and chief financial officer, Holger H\"arter, the once mighty Mr.~Wiedeking departed.

Mr.~Wiedeking drew lengthy and thunderous applause before he said his farewell to about 5,000
Porsche workers Thursday.

``We should not look back, no matter what happens,'' he said. ``We should now work constructively on
our future.''

\section{Amazon Reports Lower-Than-Expected Sales}

\lettrine{A}{mazon.com} posted weaker-than-expected earnings on Thursday, punctuated by a steep
decline in its flagship business of selling media products like books, music and DVDs.

Amazon, based in Seattle, said its net profit fell 10 percent, to \$142 million, or 32 cents a
share, in its second quarter, from \$158 million, or 37 cents a share, in the quarter a year
earlier. Revenue climbed 14 percent, to \$4.65 billion, coming in slightly below Wall Street's
expectations.

Analysts polled by Thomson Reuters on average had expected \$4.67 billion in revenue and earnings of
32 cents a share. A charge for a \$51 million legal settlement to Toys ``R'' Us depressed the
company's profit.

Shares of Amazon fell more than \$6, or nearly 7 percent, in after-hours trading. It more than
doubled in the last eight months as investors applauded the company's ability to navigate the
recession.

``People had looked at their recent performance and assumed that Amazon was, relatively speaking,
exempt from the current downturn,'' said Jeffrey Lindsay, an analyst at Sanford C.~Bernstein.
Amazon's quarter ``was by any measure a good performance, but expectations had gotten ahead of
themselves,'' he said.

For some analysts, the most worrisome indicator was a sharp falloff in Amazon's media sales in North
America. The company's original business of selling books, CDs and DVDs showed no growth in the
second quarter over the prior year, after posting 15.6 percent growth in all of 2008. Worldwide
media sales grew only 1 percent, while worldwide sales of electronics and general
merchandise\cite{merchandise} grew 35 percent.

In a conference call with analysts, Thomas J.~Szkutak, Amazon's chief financial officer, attributed
the falloff in media to an industrywide slowdown in video game and console sales and to an
unfavorable comparison with the same quarter last year, when popular products like Wii Fit were
introduced.

But analysts seemed to worry that the emergence of digital alternatives in Amazon's Video on Demand,
MP3 and Kindle stores -- which generally have lower prices than their physical counterparts -- might
be weakening one of the core drivers of Amazon's business.

``I think people are worried that the Kindle and their other digital efforts may be
cannibalizing\cite{cannibalize} their real-world sales,'' Mr.~Lindsay said.

Nevertheless, Amazon has been accelerating its shift from physical to digital media formats. This
month the company lowered the price of the Kindle reading device to \$299 from \$359 and began
selling the Kindle DX, which has a larger screen, for \$489. Amazon's Web site says the DX is sold
out, with a two- to three-week wait.

On Wednesday, Amazon announced plans to buy the shoe retailer Zappos.com for stock currently worth
more than \$800 million, the largest purchase in its 15-year history. Amazon executives said
Thursday that the acquisition was intended to help the company establish a separate outpost in the
rapidly growing business of selling shoes and apparel\cite{apparel} online.

Mr.~Szkutak said Zappos had roughly \$635 million in revenue last year and made a small profit. He
said Amazon intended to let the company function as a ``standalone business'' and was not
considering layoffs or other cost-saving measures once the deal was completed. ``This is not about
synergies\cite{synergy}; this is about growing in categories that we think are very interesting,''
he said.

Analysts generally applauded the purchase of Zappos. Marianne Wolk, an analyst at the Susquehanna
Financial Group, said buying Zappos ``eliminates a competitor in the market, and a competitor that
could have been even more powerful if it expanded into more categories.''

Amazon issued only a broad forecast for its current quarter, saying that it expected revenue to rise
in a range of 11 percent to 23 percent from a year earlier, and that operating income could fall as
much as 22 percent or increase as much as 36 percent.

Amazon's chief executive, Jeffrey P.~Bezos, did not participate in a conference call with analysts;
the company said he was traveling. But earlier in the day he did address one matter: the deletion of
the George Orwell books ``1984'' and ``Animal Farm'' from the Kindle libraries of customers who had
bought them from a vendor that lacked the rights to sell them. Angry customers said this was
equivalent to the company's breaking into their houses to take back a conventional book.

In a posting on an Amazon bulletin board, Mr.~Bezos wrote an apology, saying the move was ``stupid,
thoughtless, and painfully out of line with our principles.''

\section{Obama Wades Into a Volatile Racial Issue}

\lettrine{A}{mericans} got a rare glimpse Wednesday night of the perspective that a black president
can bring to a racially charged situation.

In response to a question at his prime-time news conference about the arrest of Henry Louis Gates
Jr., the black Harvard professor, in his own home over the weekend, Mr.~Obama declared that the
Cambridge, Mass., police had ``acted stupidly.''

Mr.~Obama's response was his most animated\cite{animate} performance of the hour-long news
conference, and represented an extraordinary plunge by a president into a local law-enforcement
dispute. And it opened a window into a world from which Mr.~Obama is now largely shielded,
suggesting the incident had struck a raw nerve with the president.

In the public spotlight, Mr.~Obama has sought to transcend\cite{transcend}, if not avoid, the issue
of race. As a candidate, he tried to confine his racial references to the difficulty of catching a
cab in New York, although he was forced to confront it directly during the Pennsylvania primary when
his pastor\cite{pastor}, Jeremiah Wright, became an issue. And last week, at the 100th convention of
the NAACP in New York, he spoke in uncharacteristically personal terms about his rise to power as a
black man, while warning black Americans not to make excuses for their failure to achieve.

Wednesday night's press conference seemed to be a different deal as the president leaped into a
highly charged controversy that has ignited passions across talk radio, the blogosphere and the
old-fashioned water cooler.

But in fact, racial profiling was a major issue for Mr.~Obama when he was in the Illinois
legislature. He was the chief sponsor of a bill, which became law, that requires police to record
the race, age and gender of all drivers they stop for traffic violations and for those records to be
analyzed for evidence of racial profiling.

And so the substance of his response was not as surprising as the fact that a president so quickly
joined the fray\cite{fray}.

The police were called to Professor Gates's house on July 16 after a report of a
burglary\cite{burglary} in progress. Professor Gates said he told the police that he lived in the
house and that he was jimmying open a damaged front door. Still, the police report said he was
arrested for ``loud and tumultuous\cite{tumultuous} behavior in a public space.'' He was held in
police custody\cite{custody} for four hours. Professor Gates said he was the victim of racial
profiling and has demanded an apology, but the police officer involved, who is white, has said he
has nothing to apologize for.

On Tuesday, disorderly conduct charges against Professor Gates were dropped, and the city of
Cambridge, its police department, the Middlesex County district attorney's office and Professor
Gates issued a joint statement calling the incident ``regrettable and unfortunate.''

Mr.~Obama, asked on Wednesday what the incident said about race relations in America, noted up front
that Professor Gates is a friend and that his comments might be biased\cite{biased}. He said
``words'' had been exchanged and added:

``Now, I don't know, not having been there and not seeing all the facts, what role race played in
that, but I think it's fair to say, number one, any of us would be pretty angry; number two, that
the Cambridge police acted stupidly in arresting somebody when there was already proof that they
were in their own home; and, number three, what I think we know, separate and apart from this
incident, is that there is a long history in this country of African-Americans and Latinos being
stopped by law enforcement disproportionately. And that's just a fact.'' He added later that the
incident was ``a sign of how race remains a factor in this society.''

He also used biting humor, grinning broadly as he imagined being in Professor Gates's seemingly
preposterous\cite{preposterous} circumstance of being arrested after trying to get into his own
home.

``Here, I'd get shot,'' Mr.~Obama said, referring to his new address of 1600 Pennsylvania Avenue.

The statement was a bit of political jujitsu\cite{jujitsu} that acknowledged the intense security
that surrounds any president while letting sink in the image of what would happen to anyone, let
alone a black man, who might seem to be breaking into the White House.

Mr.~Obama's response lit up the blogosphere immediately after the press conference. The debate
developing overnight was whether Mr.~Obama had gone too far in his response.

On nytimes.com, one commenter had this to say: ``Shouldn't we all be outraged that a black man can
somehow be arrested attempting to break into his own home in 2009? Obama is right on the money on
this one. Do we have to be black to get this?''

Another said: ``I don't think the arrest was racially motivated. I believe it was a contest of egos.
Mr.~Gates lost that battle since he was beyond the ivy walls. I agree with the President that the
arrest was stupid, but I think there was plenty of blame to go around.''

It could not be determined how well Mr.~Obama knows Professor Gates. But the professor, a widely
respected expert in the field of race relations, had very kind words for Mr.~Obama's
pivotal\cite{pivotal} speech on race relations after the Wright affair threatened to sink his
candidacy.

``I think it was brilliant,'' Professor Gates said of the speech in an interview with Tavis Smiley
at the time. ``It is a great speech about race, and race relations, particularly between black
people and white people at the beginning of the 21st Century.''

\section{Obama Criticizes Arrest of Harvard Professor}

\lettrine{P}{resident} Obama bluntly\cite{bluntly} accused the police of acting ``stupidly'' in
arresting the Harvard scholar Henry Louis Gates Jr.~last week after an officer had established that
Mr.~Gates had not broken into his own home in Cambridge, Mass.

Mr.~Obama stopped short of accusing the police department of racial profiling, as Mr.~Gates has
done.

But at a prime-time White House news conference that was otherwise largely devoted to health care,
Mr.~Obama weighed in on the Gates case and suggested that the police should never have arrested him.
He added that African-Americans and Hispanics in the United States have long been familiar with
racial profiling by law enforcement.

``There's a long history in this country of African-Americans being stopped disproportionately by
the police,'' Mr.~Obama said. ``It's a sign of how race remains a factor in this society.''

For Mr.~Obama, who has rarely talked about race, this was an unusual instance of openly addressing
an issue of deep concern to African-Americans.

The Gates arrest came up in the final question of the hourlong news conference in which Mr.~Obama
largely stuck to familiar talking points about why the country needs to overhaul its health care
system right now. Lynn Sweet of The Chicago Sun-Times asked him about the case and what it said
about race relations in America.

Mr.~Obama paused, then said, ``Well, I should say at the outset that Skip Gates is a friend, so I
may be a little biased here.''

Then he made his only joke of the evening, as he speculated about what would happen if he were seen
trying to force the door of his own home? ``I guess this is my house now,'' he said, ``so it
probably wouldn't happen.'' Then, after a beat, he added, ``Let's say my old house in Chicago. Here,
I'd get shot.''

The president then became serious, taking up a chronology\cite{chronology} of the events last week
after the police received a report of a possible break-in at the home of Mr.~Gates, a leading
authority on African-American history.

``The police are doing what they should,'' he said. ``There's a call. They go investigate. What
happens?

``My understanding is that Professor Gates then shows his I.D. to show that this is his house, and
at that point he gets arrested for disorderly conduct.''.

``I don't know, not having been there and not seeing all the facts, what role race played in that,''
Mr.~Obama continued. ``But I think it's fair to say, No.~1, any of us would be pretty angry; No.~2,
that the Cambridge police acted stupidly in arresting somebody when there was already proof that
they were in their own home; and No.~3, what I think we know separate and apart from this incident
is there is a long history in this country of African-Americans and Latinos being stopped by police
disproportionately. That's just a fact.''

The Cambridge police dropped disorderly conduct charges against Professor Gates on Tuesday, but he
says he still wants a personal apology. Professor Gates has said he thinks it was because he is
black that the officer, Sgt. James Crowley, had not at first believed he lived in the upscale home.

Many of Professor Gates's friends say they believe he may still file charges against the Cambridge
Police Department.

On Wednesday the union representing Sergeant Crowley said it was standing behind the officer. The
union, the Cambridge Police Superior Officers Association, said in a statement that Sergeant Crowley
was a ``highly respected veteran supervisor'' who had its ``full and unqualified support.''

Sergeant Crowley told The Associated Press that he had followed proper procedures and would not
apologize.

The union's statement said, ``His actions at the scene of this matter were consistent with his
training, with the informed policies and practices of the department, and with applicable legal
standards.'' The police department did not respond to requests for comment late last night.

\section{Chinese-American Children Sent to Live With Kin Abroad Face a Tough Return}

\lettrine{G}{ordon}, 3, would not look his parents in the eyes, and refused to call them Mom and
Dad. He erupted in tantrums\cite{tantrum} and sometimes cried nonstop for half an hour.

``We did not know why,'' said his mother, Winnie Liu, recalling the desperation that sent them to a
neurologist to check Gordon for autism\cite{autism}, and to a hospital that referred them to
Butterflies, a mental health program for very young children on the Lower East Side of Manhattan.

Finally they learned the reason for their child's distress -- and the reason social service agencies
that help families from China are facing a sharp rise in such developmental problems.

Like thousands of other Chinese immigrants responding to financial and cultural pressures, Ms.~Liu
and her husband, Tim Fang, had sent Gordon to live with his grandparents, thousands of miles away in
Fujian Province, a few months after his birth in New York. Working long hours in the restaurant
business, they had not brought him back to the United States until he was old enough to attend
all-day public preschool.

And now he saw them as strangers who had stolen him away to a strange land.

``The children that have that experience come back with tremendous needs,'' said Nina Piros,
director of early childhood programs at University Settlement, a nonprofit agency that estimates
that 400 of the 1,000 children served by its Butterflies program are returnees from China. ``They
come here and they're totally traumatized.''

Some act out in frightening and confusing ways, she said, banging their heads on walls, refusing to
speak, or wandering aimlessly in the classroom. These signs of extreme trauma have often been
misunderstood as symptoms of autism. But they are the marks of the emotional dislocations these
young children have endured.

Less severely affected youngsters are helped through supportive workshops for their teachers and
parents. But about two dozen in the Butterflies program need the kind of intensive therapy that
eventually helped Gordon and his parents bond, said Andrea D.~Bennett, director of Butterflies,
which was started three years ago with money from the City Council.

The phenomenon of American-born children who spend their infancy in China has been known for years
to social workers, who say it is widespread and worrying. About 8,000 Chinese-born women gave birth
in New York last year, so the number of children at risk is substantial, according to the
Chinese-American Planning Council, a social service agency that hopes to get a grant to educate
parents about the pitfalls of the practice and help them find alternatives.

But no one tracks the numbers, and the issue has only recently seized the attention of
early-childhood researchers like Yvonne Bohr, a clinical psychologist at York University in Toronto,
who calls such children ``satellite babies.''

Their repeatedly disrupted attachments to family members ``could potentially add up to a mental
health crisis for some immigrant communities,'' Dr.~Bohr wrote in an article in May in The Infant
Mental Health Journal. She cited classic research like the work of Anna Freud, who found that young
children evacuated during the London blitz were so damaged by separation from their parents that
they would have been better off at home, in danger of falling bombs.

Dr.~Bohr, who is undertaking a longitudinal study of families with satellite babies, cautions that
the older research was shaped by Western values and expectations. Chinese parents, including
university-educated professionals she has studied, are often influenced by cultural traditions: an
emphasis on self-sacrifice for the good of the family, a belief that grandparents are the best
caretakers, and a desire to ground children in their heritage.

Sending babies back to grandparents is also done in some South Asian communities, she said.

But Amanda Peck, a spokeswoman for University Settlement, which has been serving newcomers to the
Lower East Side since 1886, said that while family separations are a feature of migration in many
ethnic groups, the satellite-baby phenomenon seems rare outside the Chinese community.

Some children are better able to adapt, whether because of natural resilience, more supportive
parenting or the age at which disruptions occurred. Even in severe cases like Gordon's, the
Butterflies program has had success in overcoming the worst consequences of separation with
therapeutic play and support for parent and child, said Victoria Chiu, its bilingual therapist.

But for many children, new separations are in store even after they return to the United States. In
one typical case, parents migrated to work in a Chinese restaurant in South Carolina, taking a
school-age child along, but leaving a baby in China and a 3-year-old with grandparents in New York.

``The 3-year-old, he wouldn't even smile,'' Ms.~Chiu said. ``When he sat in circle time, his whole
little body was just slumped.''

Gordon, now 7, keeps up with his second-grade classmates and has learned to control his temper, said
his parents, who own Wild Ginger, a restaurant on Broome Street. In imperfect but fluent English,
his mother recounted the hard climb to that happy resolution, and revisited the scene of major
turning points: a tiny playroom under the eaves of the old settlement house, where a dollhouse and a
big plush dog played a role in healing her son.

Dressed as a superhero, Gordon would often rescue the dog from a pretend fire in the dollhouse,
saving him from ``the bad guys,'' as Ms.~Chiu and his mother played along.

``I was the bad lady,'' Gordon's mother, 31, recalled ruefully. ``Then the play changed, and he
tried to save Mom from the bad guy.''

The therapist explained: ``He was trying to find mastery over things he had no control over. We
started introducing scenarios to help him develop trust in his parents' authority over his life.''

Ms.~Liu, who was 17 when she immigrated to New York on a green card sponsored by her father, pressed
a hand to her heart. ``This wonderful therapist, this program, help us read the child's mind,'' she
said. ``Now he hug me, and he say 'Mommy' sometimes.''

Still, Gordon remains more withdrawn than typical 7-year-olds. Ms.~Liu said she struggles with guilt
and regret.

``I advise all Chinese families, do not send your kids away, no matter how hard, because that loss
cannot be made up,'' she said. ``Money is not so important. Nothing can make up for the sensation of
love between parents and children.''

The shuttling of babies first caught public attention in New York a decade ago, when women workers
from Fujian Province, deep in debt to the ``snakeheads'' who had smuggled them into the country, had
little choice but to send their infants back to their extended families.

Typically, such children returned at school age. Their tough adjustment to the change in language,
customs and parental discipline was generally likened to the problems of other immigrant children,
who must often cope with long-delayed reunions after being left behind for years.

Now, however, because of the expansion of free full-day preschool in recent years, satellite babies
return and start classes as young as 2 years, 9 months.

Their parents, including many lawful permanent residents and citizens like Gordon's mother, assume
that the children will adjust more easily because they are so young. But early childhood is the
crucial time for learning to form attachments and feel empathy, and serious disruptions carry
lifelong consequences, psychologists say, including higher rates of depression and dysfunction.

Many families are unaware of the potential psychological damage, said Hong Shing Lee, chief
operating officer of the Asian-American Federation of New York.

That was the case for the family of Alisa Chen, now 4. Alisa was 6 months old when her mother, Qiao
Yuni Chen, a waitress unable to afford day care, took her to her grandmother in China. When
Mrs.~Chen returned more than a year later to visit -- and to leave Alisa's baby sister, Angie -- she
was heartbroken by Alisa's rejection. Only in the last two weeks of a three-month stay was Alisa
willing to sleep at her mother's side.

Alisa started preschool at University Settlement in August, only a week after arriving in New York;
two months later, teachers referred her to Butterflies.

``She seemed kind of lost, not picking up English, withdrawing from her peers,'' Ms.~Chiu recalled.
``She seemed anxious that her mom wouldn't pick her up.'' Another problem was the mother's
expectations: The only toy in their home was a letter board more appropriate for a 6-year-old than
for a child turning 4.

Mrs.~Chen, whose husband is now in the Army in South Carolina, threw herself into becoming a more
supportive parent, Ms.~Chiu said. Though she spoke little English, she phonetically memorized songs
like ``Itsy Bitsy Spider.'' At a 99-cent store, the therapist helped her pick playthings that would
allow her daughter to express herself.

The payoff was obvious when the preschooler returned from a class trip to the Bronx Zoo one recent
afternoon. Pigtails bouncing, her smile electric with joy, Alisa threw herself into her mother's
arms. Ms.~Chiu beamed.

Next month, Alisa's little sister arrives from China to begin Head Start.

\section{Ruling May Bolster\cite{bolster} Claim to I.B.M. Fortune}

\lettrine{T}{he} Maine Supreme Court on Thursday upheld an adult same-sex adoption by a descendant
of the founder of I.B.M., with a share of the multimillion-dollar fortune at stake.

The ruling may bolster the standing of the adoptee as she pursues her claim in Connecticut to what
she says is her share of the family fortune.

But the case may have fairly limited implications for other adult adoptees, said Michael P.~Koskoff,
a lawyer in Bridgeport, Conn., who represents the adoptee, Patricia Ann Spado. Ms.~Spado was adopted
by Olive F.~Watson, a granddaughter of the I.B.M. founder, Thomas J.~Watson Sr., in 1991. They
became involved in 1979.

The purpose of adult adoptions was often to establish financial security and inheritance rights for
same-sex partners, long before any states had legalized same-sex marriage.

``Now that same-sex couples can get married, adoption would not be the method of choice to establish
a legally binding relationship,'' Mr.~Koskoff said.

Maine is one of the few states that allowed adult adoptions when Ms.~Watson, who was 43 and owned a
home in Maine, adopted Ms.~Spado, then 44.

They broke up less than a year after the adoption. But Ms.~Spado has claimed she is a legal
grandchild of Mr.~Watson and a beneficiary of his trusts because she is Ms.~Watson's daughter.

In 2005, two trustees of the Watson trusts challenged the adoption. They said that Ms.~Spado had
fraudulently claimed to be living in Maine and that the adoption violated state policy prohibiting
adoptions involving same-sex couples. A probate court ruled in 2008 in favor of the trustees,
annulling the adoption.

But on Thursday, the State Supreme Court vacated that order, saying that there was ``insufficient
evidence of fraud'' and that the trustees had not proved a violation of policy. Moreover, the court
noted, Ms.~Watson herself opposed the annulment.

\section{For Public, Obama Didn't Fill in Health Blanks}

\lettrine{A}{s} Craig Brown watched President Obama's news conference on Wednesday night on his
TiVo-equipped television, he kept hitting the pause button so he could throw questions at the image
frozen on the screen.

How much will this health care plan really cost, he asked. How can we cover nearly everybody without
higher taxes or debt? Who is going to decide which treatments are allowed? Why cannot they just get
rid of the waste without changing the whole system?

Like many in the country, Mr.~Brown, a 36-year-old father of four who lives in an Atlanta suburb,
has grown increasingly anxious about Washington's efforts to reconfigure health care and what it may
mean for his middle-class family. Although he and his wife, Judith, supported John McCain in the
presidential race, they find Mr.~Obama an earnest\cite{earnest} and compelling
pitchman\cite{pitchman}. But they remain frustrated by the lack of available detail about his plan's
contours\cite{contour} and cost.

They say they feel they are being asked to buy on spec from a government they do not trust. And they
have lots of questions.

``The bottom line is there are so many unknowns,'' said Ms.~Brown, 35, who works part time at her
church and cares for her young children. ``What we do know is there is going to be more government
control, and with more control you're going to have fewer choices. It's an innate\cite{innate} part
of being American to have those choices.''

A similar unease was apparent in three other living rooms where families gathered to watch the news
conference. An affluent small-business owner from near Chicago, a middle-class manager from Denver,
and an uninsured worker from Cleveland each expressed skepticism that change would improve their
lots.

Although she may well benefit from Mr.~Obama's plan to subsidize health insurance for the working
poor, Rowena Ventura, the uninsured worker from Cleveland, wondered whether she could afford it.
``I'm worried because they're talking about forcing people to buy insurance,'' said Ms.~Ventura, a
registered Democrat and part-time health care worker. ``You just can't ask any more of me. You just
can't.''

Ms.~Ventura, 44, who also attends community college, has moved her ailing mother into the living
room of the house she shares with her disabled husband. She said she recently discovered a
lump\cite{lump} on her left foot but cannot afford to see a doctor about it. Yet she is
cynical\cite{cynical} about Mr.~Obama's prescription.

``You see,'' she said, gesturing at Mr.~Obama on the television, ``he's saying he wants to continue
private insurance, but then he says they're part of the problem. Well, which is it? It's just
ridiculous.''

Dean Raschke, a McCain voter who owns two Chicago-area businesses, one providing roadside assistance
and the other making debit cards, said he worried that Washington would end up taxing the health
benefits he provides to his 50 employees. He said he also feared that Congress would raise his
income taxes to pay for the plan, although his earnings are well below the \$1-million-a-year
threshold now being considered.

``I have very conflicted emotions because I do want to help people who aren't as fortunate as we
are,'' said Mr.~Raschke, 38, watching the news conference with his wife, Jill. ``But I have a big
issue about what this health care plan would do to small businesses like mine that already have a
health care plan. I'm afraid that people could be unintentionally harmed.''

Recent polls have detected a modest slippage in public support for the kinds of changes being
considered in Congress, and in Mr.~Obama's handling of health care. The president has made the case
for his plan at scripted events each day this week, including at a town-hall-style meeting in
Cleveland on Thursday.

Mr.~Obama acknowledged the spectrum\cite{spectrum} of concerns during Wednesday's news conference.

``I understand that people are feeling uncertain about this,'' he said. ``They feel anxious, partly
because we've just become so cynical about what government can accomplish.'' He said he understood
that people might prefer the devil they know.

But the president's expression of empathy\cite{empathy} provided scant comfort to the Browns. They
still did not feel they were getting straight talk, as when Mr.~Obama responded to a question about
what Americans would have to sacrifice.

``He said they're going to have to give up paying for things they don't need, and that was an
awesome answer for a politician,'' Mr.~Brown said sardonically. ``You mean I don't have to give up
anything I already have?''

The Browns are Jamaican immigrants who met in college in Florida. Mr.~Brown gained citizenship in
1999; his wife expects to do so next year. The family is insured through his job at a family-owned
trophy\cite{trophy} shop, where he earns about \$38,000 a year.

Mr.~Brown said he realized that his escalating\cite{escalate} insurance premiums, which have doubled
since 2006, had suppressed his wages. He noted that he and his wife were still struggling to pay off
\$3,000 in uncovered medical expenses from the birth of their youngest child.

But the Browns said Mr.~Obama and the Democrats had not convinced them of the need for radical
change. They said the notion of establishing a new government health plan to compete against private
insurers seemed un-American. They questioned the wisdom and fairness of taxing the rich. And they
said individuals should bear more responsibility for staying healthy.

``I know the system is not perfect, but I'm not completely convinced it's broken,'' Mr.~Brown said.
``And even if it's broken, I'm not sure the government is the solution.''

Unlike the Browns, Liz Wessen, 32, a manager for a market research firm in Denver, supported
Mr.~Obama in November. But that good will does not negate her nervousness about the money being
spent in Washington.

``My only concern is that this comes on the heels of the stimulus package,'' Ms.~Wessen said from
her home in the Highlands neighborhood. ``Where is this money supposed to be coming from? I'm not
sure if this is the best time to fix another enormous problem.''

Watching the president, she said she was pleased to hear that the Democrats wanted to prevent
insurers from denying coverage to those with pre-existing conditions and to allow those who change
jobs to hold on to their coverage. But she said she wanted more specifics and wished that Mr.~Obama
would dictate terms to Congress rather than merely prod lawmakers to act.

``I think the press conference was more convincing people of his motives than it was to actually
explain the program,'' Ms.~Wessen said. ``I expected it to be more.''

\section{Studios Bring Another Dimension to Comic-Con}

``\lettrine{H}{ow} many of you have ever wanted to go to another planet?'' shouted the filmmaker
James Cameron, addressing 6,500 fans here on Thursday.

With those words, the future came to Comic-Con -- and it was wearing funny glasses.

For the first time, the giant annual convention for fans of comics, movies, superheroes, science
fiction and video games that has become a major event on Hollywood's calendar -- now in its 40th
year -- devoted almost an entire day to presenting films in 3-D, culminating with a 25-minute
glimpse of ``Avatar,'' the eagerly anticipated science-fiction thriller due in December from
Mr.~Cameron and 20th Century Fox.

``I'm not going to waste your time, because I know you're here to see stuff,'' said Mr.~Cameron, who
spoke only a few words before showing pieces of his film.

The scenes portrayed the adventures of a human experiencing a new world, the planet Pandora, through
the body of an alien. And they were delivered with a startling verisimilitude that seemed to plant
flying insects and waving ferns on the heads of people in the next row -- and had Comic-Con fans
roaring with approval.

The ``Avatar'' presentation cleared a crucial test for Fox, which has invested more than \$200
million in the movie, the most conspicuous example of a recent spate of expensive 3-D efforts from
the major studios that will roll out over the next few years.

Mr.~Cameron's devotion to next-generation 3-D methods has raised enormous expectations around a film
made with technology that some people predict could have an impact on movies comparable to the
introduction of sound.

It was a measure of the movie's importance that Tom Rothman, co-chief of Fox Filmed Entertainment,
personally introduced Mr.~Cameron.

``Moments like these are rare in the life of a movie company,'' Mr.~Rothman said. Mr.~Cameron's last
feature film was the box-office juggernaut ``Titanic,'' well over a decade ago.

Actually, it had been Disney executives who led the push to use 3-D promotions at Comic-Con, where
projection technology had previously been unequal to such large-scale presentations.

No small problem was the distribution of Dolby 3-D glasses by the thousands, with repeated pleas
that they eventually be returned, though many were left broken on the floor.

Even as the first trailers were screened, hours before Mr.~Cameron's presentation, doubts lingered
about whether the big 3-D day would work.

``Some of the 3-D effects aren't going to look absolutely perfect,'' said Patton Oswalt, the
comedian and voice actor who moderated a presentation of Disney's forthcoming ``A Christmas Carol,''
``Alice in Wonderland'' and ``Tron: Legacy.'' Mr.~Oswalt asked attendees not to blog about effects
being ``quite subpar.''

But the fans ``oohed'' and ``aaahed'' over the displays of technical virtuosity, starting with
motion-capture film from Robert Zemeckis's ``A Christmas Carol,'' featuring Jim Carrey.

Sony Pictures followed with its ``Cloudy with a Chance of Meatballs,'' and New Line weighed in with
a 3-D sequel to the horror film ``Final Destination.''

But if technology ruled the day, it took flesh-and-blood movie stars to bring down the house. A
surprise appearance by Johnny Depp -- who showed up beside the director Tim Burton in support of his
Mad Hatter's role in ``Alice'' -- caused a near riot, with fans rushing the stage for pictures.

A couple of hours later, the cast of Summit Entertainment's ``The Twilight Saga: New Moon'' provided
the day's emotional pop. Teenage girls started lining up on the street to see snippets from the
second movie in the series. Many wore cardboard Burger King crowns emblazoned with ``Team Edward,''
referring to Edward, the nonbloodsucking vampire at the franchise's center.

Media attention on the Summit Entertainment franchise is so enormous that the studio opted to hold a
news conference at the nearby Hilton Hotel on Thursday in place of the normal press panels after the
big Hall H presentation. (``It is for the personal safety of our stars,'' said a Summit staff
member.)

The ``New Moon'' trio -- Robert Pattinson, Kristen Stewart and Taylor Lautner -- took the stage in
the Hilton ballroom looking as if the intense scrutiny was starting to wear on them. A
disheveled-looking Ms.~Stewart nervously took questions in front of about 30 television and video
cameras. She described the greatly anticipated film as ``seriously emotional'' and coming from an
``empty, completely dark place.''

In a lighter moment, a television news reporter asked the young stars what questions they were tired
of being asked, and Mr.~Lautner said, ``Fans ask me to growl for them and I really don't enjoy doing
that.'' Ms.~Stewart said people could ask her anything they liked. But a Summit executive cut off a
question about off-screen chemistry between cast members.

The trade-show floor and the streets surrounding the convention center presented the usual array of
Comic-Con curiosities. Two men in ``Star Wars'' Stormtrooper uniforms mingled with women dressed as
Alice in Wonderland. Vampires and grim reapers were in abundance. The cable network A\&E staged a
mock funeral to promote ``Chris Angel Mindfreak'' -- complete with coffins carried by men wearing
black lipstick -- in the middle of a main walkway. Women on roller skates zigzagged through the
crowds, apparently to promote Drew Barrymore's roller-derby movie, ``Whip It.''

Only a couple of months ago, those who run the convention said they doubted that 3-D could be
properly displayed in the huge hall where Hollywood studios have become accustomed to introducing
prospective blockbusters like ``Watchmen'' and ``X-Men Origins: Wolverine. ``We're about to find
out,'' David Glanzer, the convention's director of marketing, said just before the day's session
began.

Two years ago, Paramount Pictures screened a piece of ``Beowulf'' in 3-D here. But that occurred
only in a small, preconvention preview.

This time, 3-D was the main event. Asked if he intended ever again to make a conventional,
two-dimensional live action film, Mr.~Zemeckis could only muster a maybe.

``Never say never,'' he said.

But, he added: ``Right now, I'm dedicated to sending these forms into the world. I love the control
these art forms give the director.''

\section{Asia Trip Propels Clinton Back Into Limelight}

\lettrine{A}{t} every stop during her visit to Southeast Asia this week, Secretary of State Hillary
Rodham Clinton delivered the same message: ``The United States is back.'' She was talking about
America's role in a region where its visibility had dimmed during the Bush years, but she might as
well have been talking about herself.

After being sidelined at home in recent weeks by a broken elbow, Mrs.~Clinton thrust herself back
into the limelight, making headlines with unexpected statements about Iran, North Korea and Myanmar
-- at least one of which her aides felt obliged to clarify. And she reaffirmed that overseas, at
least, she is not just the nation's chief diplomat but one of its most reliable celebrities,
trailing perhaps only President Obama in star power.

In India, Mrs.~Clinton was a ubiquitous\cite{ubiquitous} presence on the three main news channels,
giving interviews, sipping coffee with tycoons in Mumbai, marveling over climate-friendly buildings
and dining at a New Delhi restaurant where one of the menus is named after her husband, former
President Bill Clinton.

In Bangkok, she bantered\cite{banter} with two Thai journalists in a televised town hall meeting,
which, like previous public events with Mrs.~Clinton in Turkey, Japan and South Korea, ended up
being a cross between a Council on Foreign Relations panel and an episode of ``Dr.~Phil.''

What, the secretary of state was asked, does she talk about with her husband? ``We talk about our
dog, who got sick and had to go to the vet,'' she said.

And what does she make of the beleaguered American economy? Mrs.~Clinton invoked the name of Ben
S.~Bernanke, the Federal Reserve chairman, and his recent Congressional testimony before lamenting
that unemployment was still high in many parts of the country. ``We're not out of the woods,'' she
said, in a guarded tone. ``I guess that's the best way to say it.''

Mrs.~Clinton batted away suggestions that she had been marginalized. That, she said, was a
canard\cite{canard} propagated by the news media, which jumped on the fact that her recovery from
the elbow injury forced her to cancel two foreign trips, including one with Mr.~Obama to Russia.

``I'm not with the president on the trip and all of a sudden everybody goes, 'Oh, where is she?
She's gone, disappeared,' '' Mrs.~Clinton said in a tone of mock horror, as the audience giggled.

Certainly, she seemed determined to continue her breakneck schedule back home. On Friday, after
flying from Thailand, she scheduled a full day of meetings, including one with the Iraqi prime
minister, Nuri Kamal al-Maliki. On Sunday, Mrs.~Clinton is to appear live on ``Meet the Press.'' And
next week, she will be co-chairwoman of a high-level meeting with Chinese officials, along with
Treasury Secretary Timothy F.~Geithner.

Mrs.~Clinton stayed in the news, even from unlikely places like Phuket, a lush\cite{lush} beach
resort that played host this week to a meeting of the Association of Southeast Asian Nations.

She warned about nuclear links between North Korea and Myanmar. And she offered a rare glimpse into
what the United States might do if Iran did not respond to its diplomatic overture. Her comments
about the possibility of extending a ``defense umbrella'' over the Middle East ricocheted around
foreign-policy circles.

Mrs.~Clinton's aides did their best to convince reporters traveling with her that she was not
signaling any change in administration policy. But when a secretary of state muses\cite{muse}
publicly about one of the world's most sensitive diplomatic problems, the story quickly takes on a
life of its own.

In India, Mrs.~Clinton encountered a different problem: a savvy environment minister, Jairam Ramesh,
who seized on a visit by her to a ``green building'' outside New Delhi to repeat India's rejection
of American demands that it accept binding cuts in its greenhouse gas emissions.

Fortunately for her, most of the audiences seemed to care less about greenhouse gases and security
umbrellas than about hearing dish from her about her life, her ambitions and, particularly, her
relationship with Mr.~Obama.

On the inevitable question of how she felt working for the man she had fought in the primary
campaign, she said, ``The president is the president; I tried to be the president and was not
successful.''

Nor did she bite at questions about whether she would ever run for the White House again. ``That's
not anything I'm at all thinking about,'' she said. Mrs.~Clinton did offer some morsels\cite{morsel}
about Mr.~Obama's courtship of her to be the secretary. She said she resisted his initial job offer
and gave him names of people she thought would be suitable candidates. Mr.~Obama told her he needed
her, however, and his persistence wore her down.

``He gave me an enormous amount of authority as secretary of state, and really everything I asked
for so that I could the job that he wanted me to do,'' she said, ``and I was running out of
excuses.''

\section{1 Plus 1: Shanghai Tweaks Child Rules}

\lettrine{T}{hirty} years after China began enforcing a one-child policy, this city is actively
encouraging young couples to have a second child.

City officials are planning to visit homes, slip leaflets under doors and offer emotional counseling
and financial incentives.

The world's most populous country, it seems, wants to be a little bigger.

The new push, which aims to tackle growing worries about the country's shrinking work force and
aging population, is the most public effort yet to counter a policy that is considered both a
tremendous success and a terrible failure. While it has kept population growth under control, it has
also led to forced abortions.

China is not doing away with the one-child policy, which still largely applies to urban residents,
but is allowing more exceptions to the rule. Shanghai, with 20 million residents one of China's
biggest cities, is leading the effort.

``We advocate eligible couples to have two kids because it can help reduce the proportion of the
aging people and alleviate\cite{alleviate} a work force shortage in the future,'' Xie Lingli,
director of the Shanghai Population and Family Planning Commission, was quoted as saying in the
Friday issue of China Daily, the country's largest English-language newspaper.

Already, ethnic minorities are allowed to have more than one child, and many rural residents can
have a second child if the first child is a girl, because many families still follow tradition and
favor male heirs.

But over the decades, China's urban residents have faced greater restrictions, enforced with a
variety of incentives and fines.

Now, those restrictions are being lifted. A couple made up of two parents who have no siblings
themselves can now have a second child.

That is the demographic Shanghai is appealing to in the hope of rebalancing a population that is
aging fast, with about 22 percent of its residents already over the age of 60.

The country as a whole is facing a similar problem. In 2006, about 8 percent of China's population
was over the age of 65. But that proportion is expected to triple to about 322 million people, or
nearly a quarter of the country's population, by 2050, according to the United Nations.

\section{Store From Middle America Opens an Outlet\cite{outlet} in Midtown}

\lettrine{J}{. C}. Penney, the middle American department store, is suddenly in the middle of
Manhattan. This week, the value-priced chain quietly opened its first store in the city, in Midtown
-- a mere falafel cart or two away from its rival Macy's. The grand opening is July 31.

``We always knew we wanted a presence in Manhattan,'' said Michael J.~Boylson, executive vice
president and chief marketing officer for J.~C. Penney, which is based in Plano, Tex. ``It's only
been within the last four to five years it was even financially feasible.''

Feasibility may not mean success, of course. Near term, Penney's, like retailers of all stripes, is
still trying to claw its way out of the recession. And the horizontal sprawl of the typical Penney's
store does not translate readily to Manhattan real estate.

``It's not so easy to hit a home run,'' said Marshal Cohen, chief industry analyst for NPD Group.

Yet like the ing\'enue in the musical ``42nd Street,'' nearly every big-box retailer dreams of
someday coming to Manhattan. This year, Costco will open its first Manhattan store in East Harlem.
Next spring, the upscale outlet chain Nordstrom Rack will open its first Manhattan store in Union
Square. Home Depot and Best Buy have opened several stores in the city in the last decade.

The daunting\cite{daunt} prospect of adapting the big-box business to the Big Apple has limited
Target to the city's outer boroughs, while Wal-Mart has not ventured in from the suburbs.

The retailers that do set up shop in Manhattan come for the exposure, the cachet and the pipeline to
some of the nation's wealthiest consumers, even in a recession.

Although Penney's has stores in the Bronx, Queens and Staten Island, Manhattan represents its own
geographic and psychographic space.

``It is the center of media in the United States,'' said Mr.~Boylson, who added that being in
Midtown would enable his company to introduce itself to Manhattan residents (and news media) who may
have ``an outdated perception'' of Penney's. The company has been using the tag line ``We're
stepping up our style'' in its advertising, and highlighting its clothing by the likes of Charlotte
Ronson and Nicole Miller.

Penney's, which will occupy two lower levels of the Manhattan Mall, on the Avenue of the Americas
near 33rd Street, also contends it can make money in Manhattan.

``We think now, more than ever, the value that we bring to the table is going to play very, very
well,'' Mr.~Boylson said.

Most of the borough's boutiques and department stores cater to upscale consumers, he said, so
Penney's prices ``are going to fill a void.'' Executives have said the chain tends to do well in
ethnically diverse markets and is among the most popular clothing stores for Hispanic consumers.

By opening in a major traffic hub and aiming at households earning about \$65,000 a year, J.~C.
Penney hopes to attract new customers and take its national retail battle with Macy's to its rival's
home turf. Advertisements for the new store zing Macy's with barbs like, ``the real miracle is now
on 33rd Street,'' referring of course to the film ``Miracle on 34th Street,'' about a Macy's Santa
Claus.

Penney's is tackling Manhattan even as June comparable sales fell 8.2 percent from June 2008. The
heady days of 2006 when the company was reporting robust sales increases seem a dim memory. The
relatively good news, however, is that while Penney's is still reporting declines with other
retailers, the double-digit drops of last year have not been repeated since January.

At 153,000 square feet, Penney's Manhattan store (which includes a Sephora cosmetics shop) is about
the same size as its other locations. But the ceilings are significantly lower -- eight feet high
instead of the usual 11. To create the illusion of height, the company installed white panels that
float slightly below a black ceiling. Also, directional signs that hang from the ceiling are
daintier than usual.

Like all out-of-towners, retailers go to great lengths to make it in Manhattan. They overhaul their
typical big-box store designs, rethink how they manage inventory, hire more employees and even offer
different merchandise. And they must contend with local concerns -- the survival of nearby small
businesses, traffic patterns, noise and labor practices.

Even after years of planning, they still must do a lot of on-the-job learning. Home Depot, for
instance, learned from customers that, in a borough\cite{borough} of tiny kitchens, diminutive
ranges, not usually stocked by the chain, are a must.

``That was one of those baptisms\cite{baptism} by fire,'' said Steve Lieberman, a regional
merchandising manager for Atlanta-based Home Depot. And the chain's popular free do-it-yourself
workshops had to be tailored to Manhattan residents. Instead of clinics on lawn care, the Manhattan
Home Depot offers seminars on how to pot a plant. Another class that has been popular in Manhattan
is how to install your own toilet.

``We sell a boatload of toilet seats and toilets'' in Manhattan, Mr.~Lieberman said. And delivery is
much more important than in the suburbs; customers cannot lug\cite{lug} a toilet home on the F
Train. (Penney's is offering same-day delivery in Manhattan for \$15 for up to 70 pounds of
merchandise.)

Then there is the special character of New Yorkers. Executives have estimated that a city customer's
``wait tolerance'' will be no more than 90 seconds.

To hurry things along in its Manhattan shoe department -- which will offer more shoes (some 41,000
pairs) than any other Penney's -- the retailer hired shoe runners. When a customer wants to try on a
pair, a sales associate will use an electronic hand scanner to send the shoe's size and type to the
runners in the stockrooms. The runners will fetch the shoes and then place them on dumbwaiters that
zoom up to the sales floors.

Other New-York-minute amenities\cite{amenity} include an electronic queuing system that directs
customers to cash registers and estimates wait times -- almost as if the checkout line were a Walt
Disney World ride.

A machine called the FindMore has a touch screen to enable shoppers to view every item in the store
and on Penney's Web site. Even if the store is out of an item, customers can still buy it using the
machine.

Among the trickier challenges Manhattan poses to retailers is figuring out how to distribute a
big-box store model over multiple floors.

James Damian, senior vice president of the enterprise design group for Minneapolis-based Best Buy,
said chains had to determine whether their brand could survive vertically.

``Most big-box retailers fail because they've never kicked the tires on that,'' he said.

One consideration is how well each floor will accommodate throngs of shoppers. ``The traffic is
ferocious,'' Mr.~Damian said.

That and the scarcity of storage space in Manhattan means big-box chains typically must house their
inventory outside the borough, and use small trucks to bring it into their stores. Mr.~Boylson said
Penney's was storing inventory in a distribution center in Manchester, Conn., and expected to be
``running smaller trucks to New York almost around the clock.''

The chains also have to increase the number of employees, cash registers and fitting room stalls.

The Manhattan Penney's has 500 employees, making it the company's largest store staff. It has 100
cash registers; most Penney's have 58. The fitting rooms have twice the number of stalls, 61, of the
average Penney's.

Elbow room is tight, even for dummies. The Manhattan store has more mannequins\cite{mannequin} than
any other J.~C. Penney outlet, to fill the myriad\cite{myriad} display cases near the escalators and
entrances.

The company is also using technology to help manage potential crowds: if an hour goes by and the
bathrooms have not been cleaned, a buzzer automatically goes off, alerting store staff.

Heavy customer traffic is among the reasons the company likes Manhattan Mall. Consumers can enter
from the street, or from the 34th Street Herald Square subway station and the 33rd Street PATH train
station.

Company executives declined to say whether the recession had enabled them to get a break on the
rent, but Mr.~Boylson said the mall owner, Vornado Realty Trust, allowed Penney's to design the
space as it desired.

The Manhattan store merchandise\cite{merchandise} is essentially the same as that in other Penney's,
though during a recent tour, Pete Sadler, a district manager, said the store planned to take market
share from Macy's by offering more plus- and petite-size clothing.

A lot of merchandise is already on sale, like 40 percent off cultured pearl jewelry, and 50 percent
off clothing brands like St. John's Bay.

Walking the aisles, Mr.~Sadler squeezed \$3.99 pillows and stacks of white bath towels, also \$3.99.

``We think the people in Manhattan are going to love these prices.''

\section{New Factory Seen as Boon\cite{boon} for Upstate}

\lettrine{A}{ decade} ago, it was 200 acres of pine trees near Saratoga Lake. A decade from now, it
could be the booming heart of a new high-tech corridor, filled with some of the most sophisticated
labs and plants in the country.

After years of haggling and setbacks, officials finally broke ground Friday on a \$4.2 billion plant
that will manufacture advanced microprocessors and would be the most advanced facility of its kind
in the world. The plant, to be built on a 222-acre site, will be bigger than the Jacob K.~Javits
Convention Center in Manhattan.

Elected officials are counting on the expected completion of the factory in 2011 to provide a major
boost to the ailing upstate region, where for years they have promised, and mostly failed, to
deliver an economic resurgence.

The plant, which will manufacture chips for California-based Advanced Micro Devices and other chip
design firms, will generate 1,400 new manufacturing jobs, officials say.

With the new factory as an anchor, they say, they can attract more investment to the area, building
a cluster of businesses and academic centers that could ultimately rival Route 128 outside of Boston
or North Carolina's Research Triangle.

``You've got everything in place here to make the whole upstate region a leader in chip fab
manufacture and high technology generally,'' said Charles V.~Wait, president of Adirondack Trust
Corporation, a bank that helped finance early efforts to bring tech companies to the area. ``When I
grew up, you couldn't find a job here in Saratoga, and everyone left after college.''

The plant will receive \$1.37 billion in state aid over 15 years, making it one of the most heavily
subsidized projects in state history. About half that money is capital aid that was appropriated in
this year's budget, a considerable gamble at a time of severe economic crisis.

``I was concerned about the value of taxpayer dollars versus the number of jobs that would be
created,'' said Gov. David A.~Paterson, speaking at a news conference after the groundbreaking.
``But when I started to see that we would move to global leadership in semiconductor development, I
then realized that these are where the jobs of the next decade can come from.''

The new factory is also a gamble for Advanced Micro Devices, the world's second-largest producer of
computer chips.

Since January, the company has posted hundreds of millions of dollars in losses, and it expects
sales to continue declining. In March, it partnered with the Advanced Technology Investment Company,
created by the Abu Dhabi government, to spin off GlobalFoundries, which will build and own the New
York plant. The Abu Dhabi backing will provide a hedge against any threats to the project posed by
the continued recession in the United States.

While chip design is a big business in the United States, most chips are manufactured in China,
Taiwan and Singapore. The factory here, some experts say, will help make the country globally
competitive in chip manufacture for the first time in years.

``The grants and tax support level the playing field on a global basis,'' said Douglas A.~Grose, the
chief executive of Global Foundries. ``We looked all over the world for the next site. Every country
that wanted a facility like this stepped forward.''

For state economic development officials, the groundbreaking was the culmination of more than a
decade of planning that began during the administration of Gov. George E.~Pataki.

Before state officials could entice an advanced chip factory to the region, they had to build an
entire infrastructure to support it. That included not only a network of new roads and utility
upgrades, but the creation of local academic centers, including an institute for nanotechnological
engineering at the State University at Albany.

``There were times when we had some doubt along the way,'' said Mr.~Grose.

Even after Advanced Micro Devices settled on the project site, in 2006, there were hurdles,
including zoning issues, the company's financial difficulties and the need for more state financing.
To smooth the way for timely construction, Mr.~Paterson -- once a skeptic about the project --
jumped in to negotiate a labor agreement under which the factory would be built mostly by New York
workers.

At Friday's ceremony, a mound of dirt was piled in an air-conditioned tent with seven shovels ready
for the groundbreaking ceremony, one for each of the speakers. Champagne awaited in an adjoining
tent. The lineup of elected officials included Mr.~Paterson, Mr.~Pataki, United States Senator
Charles E.~Schumer and Joseph L.~Bruno, the former State Senate leader who retired last year and who
regards the factory as part of his legacy.

Mr.~Bruno said he was so eager to see the project through final approval that he had put off
retiring.

``I knew it was going to happen,'' he said.

\section{Microsoft Offers Europe a Choice of Browsers}

\lettrine{M}{icrosoft} has offered to give European computer users a choice of Web browsers, an
unprecedented move that the company, the world's largest software maker, hopes will avert another
antitrust conviction and big fine.

Microsoft said users of new computers running Windows 7 in Europe would get a ballot\cite{ballot}
screen that would let them easily download other browsers from the Internet and turn off Microsoft's
default Internet Explorer.

If European officials approve the proposal and drop their antitrust case against Microsoft, the
version of Windows 7 with the ballot screen could be available in Europe shortly after the official
release of the software on Oct.~22, the company's general counsel, Bradford L.~Smith, said in an
interview.

The European Commission, the European Union's executive body, said in a statement released in
Brussels that it ``welcomes this proposal and will now investigate the practical effectiveness in
terms of ensuring genuine consumer choice.''

The offer by Microsoft, which a decade earlier had fought attempts to open Windows to another
competing browser, Netscape, represents an unprecedented concession in its latest antitrust clash
with European regulators.

Mr.~Smith said the concessions had the potential to settle the latest European antitrust case.
``This is an important step forward and follows extensive and lengthy discussion between Microsoft
and the commission. But until the commission approves something, nothing is really done. Given the
state of the economy, we believed it was very important to move forward quickly,'' he said.

``This would mark a big step forward in addressing a decade of legal issues and would be good news
for European consumers and our partners in the industry,'' the company said in its statement.

The ballot\cite{ballot} screen would not be an option for buyers of PCs using Windows 7 software in
the United States or in other markets outside of the European Union. But Microsoft said it also
planned to release confidential technical information worldwide that would help competitors design
software that worked more smoothly with Windows, Windows Server, Exchange, Office and SharePoint.

Its rivals were taking a wait-and-see attitude. ``It is, of course, a good development that
Microsoft will make changes to allow users to choose their own default Web browser,'' said John
Lilly, the chief executive of Mozilla, which makes Firefox, the second-most-used browser after
Internet Explorer. But Mr.~Lilly said he reserved judgment on the proposal until seeing details,
including which browsers would be able to participate in the ballot and under what terms, and
whether Microsoft would use automatic updates to offer a browser choice to existing Windows users.
Google, which makes the Chrome browser, had no comment.

John Briggs, a partner at Axinn, Veltrop \& Harkrider, a Washington law firm specializing in
antitrust litigation, called the development a breakthrough for Microsoft and the commission.

``I think it is very significant that Microsoft is willing to provide access to other browsers,'' he
said. ``It's also significant that the commission and Microsoft are negotiating without having
proceedings go all the way.''

Microsoft and European regulators had clashed for almost a decade in the first antitrust case, as
the software maker tried to rebuff charges of abusing its dominance in operating systems to aid its
media players and servers. The legal wrangling\cite{wrangle} eventually cost Microsoft more than \$1
billion in fines, and reams of negative publicity.

Microsoft abandoned the case in October 2007 after losing an important appeal before the European
Court of First Instance in Brussels, which upheld the commission's 2004 judgment.

The commission's browser case -- which has drawn support from Microsoft's rivals like Google, Sun
Microsystems, Nokia and others -- appeared to be on a similar, acrimonious track.

Mr.~Briggs said Microsoft had learned a costly lesson in the first antitrust case. ``They got their
lunch handed to them in the first case and I think they probably realized it may be easier to give
up a little than to lose a lot again.''

\section{NBC Benches a Heavy Hitter for Fall}

\lettrine{D}{espite} noise about a 12-month schedule, summer retains most of its minor-league
trappings for network television: new shows, like ``Survivor,'' ``American Idol'' and ``Dancing With
the Stars'' have often been introduced between June and August, and if they prove they have the
goods, they get a call up to the big leagues in September.

That has not happened in a few summers, mainly because nothing the networks have introduced has
resembled a young Mickey Mantle.

But there is one current summer show that is putting up some power numbers this year. NBC's
``America's Got Talent,'' a somewhat old-fashioned variety show competition featuring singers,
dancers, jugglers and ventriloquists, is the undisputed ratings champ of the summer. It dominates
every hour it is on, posting viewer totals -- more than 13 million -- that would rank it among the
top shows at any time of the year. Tuesday's edition, with 13.7 million, had more viewers than any
nonsports show since the ``Idol'' finale last May.

So surely the time has come for NBC to elevate its summer hit to television's top rung of
competition. Yet NBC remains hesitant to take the show to the big time of regular season television
competition.

``We program year-round,'' said Ben Silverman, the co-chairman of NBC Entertainment, arguing for
keeping the show as a summer powerhouse.

The decision puzzles Brad Adgate, senior vice president of Horizon Media, who noted that NBC has had
a long struggle in prime time with no certain hits on the immediate horizon.

``Every network could use another hit show,'' Mr.~Adgate said. ``And this is certainly a proven
hit.''

NBC is already facing some questions about its fall schedule. A planned new drama, ``Parenthood,''
had to be delayed, forcing the network's only backup, a series called ``Mercy,'' onto the air sooner
than expected.

NBC's plan, Mr.~Silverman said, is to push ``America's Got Talent'' through until the second week of
September when its two-night finale will provide an initial lead-in for NBC's biggest programming
gamble of the fall: the new Jay Leno prime-time hour, scheduled to run every weeknight at 10.

``We can keep looking at the number of hours'' for the show, Mr.~Silverman said, suggesting the
series could run deeper into the fall in the future.

Paul Telegdy, the executive in charge of reality programs for NBC, underscored the argument, citing
the opportunity to promote NBC's coming series while competitors are pulling in only a fraction of
the number of viewers. ``We're really quite happy ruling the roost in the summer,'' Mr.~Telegdy
said.

Fox's talent competition series, ``So You Think You Can Dance,'' has already been ticketed for a new
fall address. In head-to-head contests this summer, ``America's Got Talent'' has crushed the dance
series, both in total viewers and among the young adult viewers that many advertisers prefer.

There are other issues with moving the show into the teeth of television's regular season. The first
is the complicated logistics involved in setting up a season.

``The audition process is arduous,'' Mr.~Telegdy said. ``It takes about eight months.'' That would
mean a plan would have to be in place near the start of the calendar year to get the show ready for
a fall run.

But the major factor in NBC's decision may be Simon Cowell, the show's creator and executive
producer, who also happens to be the biggest star reality television has created, thanks to his
judging job on ``Idol.''

Mr.~Cowell's value to ``Idol'' is so great that his new contract with Fox will almost surely shatter
records. Though he does not appear on ``America's Got Talent`` (his deal with Fox makes his on-air
work exclusive to that network), he is so closely associated with it that Mr.~Silverman said that he
would not consider running the show from January on, when ``Idol'' is on the air.

This is not a strict contractual limitation, Mr.~Silverman conceded, but ``there are a lot of
relationships involved,'' he said.

Mr.~Cowell would welcome moving the show to the fall. ``I've said for some time that we'll end up
with an audience of 20 million viewers,'' Mr.~Cowell said in a telephone interview (from a yacht in
the south of France, where he is vacationing). ``It's a good show; it can work any time of year.''
He'd particularly like to see it compete with ABC's ``Dancing with the Stars'' in the fall.

``We crushed that show with the 'Britain's Got Talent' in the U.K. this year,'' he said.

Mr.~Cowell said the next few weeks should tell a lot and perhaps force NBC's hand. The show will end
its audition phase after this week, moving on to live shows involving the top 40 finalists.

``We'll see how the live shows go,'' Mr.~Cowell said. ``If we can get the live shows up to 14
million or 15 million viewers, then I think we'd have a serious conversation with NBC.''

\section{Gates Faults U.S. on Data Privacy and Immigration}

\lettrine{I}{n} a far-ranging speech on Friday, Bill Gates criticized the American government's
policy on immigration and data privacy, predicted giant leaps in technology in the near future and
explained why he had to shut down his Facebook page.

``Over the next decade, the entire way we interact'' with computers will change, Mr.~Gates, the
chairman of Microsoft, told hundreds of government officials and information technology executives
in New Delhi. Mr.~Gates spoke of cellphones that would recognize people around them or be used to
test for diseases, computers equipped with voice recognition and an Internet that was used for much
more than Web pages.

While the recession has been a ``big deal,'' it has not slowed innovation, he said, in part because
countries like India and companies like Microsoft are investing in education and research for a new
generation of computer scientists.

Microsoft is angling to work on India's national identity card project, Mr.~Gates said, and he will
be meeting with Nandan Nilekani, the minister in charge. Like Mr.~Gates, Mr.~Nilekani stopped
running the technology company he helped to start, Infosys, after expanding it into one of the
biggest players in the business. He is now tasked with providing identity cards for India's 1.2
billion citizens starting in 2011. Right now in India, many records like births, deaths,
immunizations and driving violations are kept on paper in local offices.

Mr.~Gates was also critical of the United States government's unwillingness to adopt a national
identity card, or allow some businesses, like health care, to centralize data-keeping on
individuals.

``It has always come back to the idea that 'The computer knows too much about you,' '' he said.

The United States ``got off to a bad start'' when it comes to using computers to keep data about its
citizens, he said. Doctors are not allowed to share records about an individual patient, and virtual
doctor visits are banned, he said, which ``wastes a lot of money.'' The United States ``had better
come up with a better model'' for health care, he said.

He was also critical of Congress's stance on immigration, and said he would like to see immigration
exceptions for ``smart people.'' Canadian laws are more favorable, he said, because they allow
immigrants to work if they are offered a high-paying job. Microsoft has created ``a lot of jobs in
Canada for that reason,'' he said.

Asked whether he ever ``unplugs,'' abandoning e-mail messages, computers and his cellphone entirely,
Mr.~Gates laughed and said ``I'm not a 24-hour technology person.'' He said he read a lot ``and
sometimes not on a screen.'' He added that he was not big on text messaging. ``All these tools of
technology let us waste our time if we're not careful,'' he said.

Mr.~Gates said the buzzwords ``social networking'' applied to something that had been around for a
long time -- a way to communicate with numerous people at the same time.

He acknowledged that he once had a Facebook page, but every day ``ten thousand people tried to be my
friend.'' He said he spent too much time trying to decide ``Do I know them? Don't I know them?''
Ultimately, he said, ``I had to give it up.''

\section{Obama Shifts Tone on Gates After Mulling Scale of Debate}

\lettrine{P}{resident} Obama tried Friday to defuse a volatile national debate over the arrest of a
black Harvard University professor as he acknowledged that his own comments had inflamed tensions
and insisted he had not meant to malign the arresting officer.

Mr.~Obama placed calls to both the professor, Henry Louis Gates Jr., and the man who arrested him,
Sgt. James Crowley, two days after saying the police had ``acted stupidly'' last week in hauling
Professor Gates from his home in handcuffs. Mr.~Obama said he still considered the arrest ``an
overreaction,'' but added that ``Professor Gates probably overreacted as well.''

``I obviously helped to contribute ratcheting it up,'' the president said in an appearance in the
White House briefing room. ``I want to make clear that in my choice of words, I think I
unfortunately gave an impression that I was maligning the Cambridge Police Department or Sergeant
Crowley specifically, and I could have calibrated those words differently.''

Mr.~Obama's unusual personal intervention and public statement came just four hours after the White
House said he had no more to say on the matter. But after talking with Michelle Obama and some of
his closest friends amid unrelenting publicity, his advisers said, the president reversed course in
hopes of quashing a dispute that had set off strong reactions and made it harder for the White House
to focus attention on his efforts to pass health care legislation.

The Gates case has become the first significant racial controversy Mr.~Obama has confronted since
being sworn in as the nation's first African-American president. The improvisational handling of it
underscored the delicate challenges for a leader who has tried to govern by crossing old lines and
emphasizing commonalities over differences.

Advisers said both his sharp statement, which was made at Wednesday night's news conference, and his
toned-down remarks on Friday reflected strains of his experiences. He was personally outraged by the
arrest and wanted to speak bluntly about it, aides said. And they said he was distressed that his
words proved polarizing and contrary to his instincts for conciliation.

Whether he succeeded in tamping down the emotions of the case remained to be seen. In their
telephone conversation, Mr.~Obama said, Sergeant Crowley suggested that he and Professor Gates come
to the White House to share a beer with the president. Mr.~Obama then conveyed that idea in his
phone call with Professor Gates.

Professor Gates said in an e-mail message afterward that he was ``pleased to accept his invitation''
to come to the White House and meet Sergeant Crowley. ``After all, I first made the offer to meet
with Sgt. Crowley myself, last Monday,'' he wrote. ``I told the president that my entire career as
an educator has been devoted to racial healing and improved race relations in this country. I am
determined that this be a teaching moment.''

Sergeant Crowley made no public comments after his conversation with the president. He has denied
doing anything wrong and has declined to apologize to Professor Gates.

The episode stemmed from a misunderstanding when Professor Gates returned to his Cambridge home on
July 16 and found his door stuck. A woman reported seeing someone trying to break into the house and
the police responded. Although the arresting police officer became aware that Professor Gates was in
his own home, the police said he was belligerent and arrested him for disorderly conduct. The charge
was later dropped.

Mr.~Obama defended his decision to weigh in. ``The fact that this has become such a big issue I
think is indicative of the fact that, you know, race is still a troubling aspect of our society,''
he said. ``Whether I were black or white,'' he said, commenting ``is part of my portfolio.''

Mr.~Obama first discussed with aides how to address the arrest during a meeting before his Wednesday
news conference. Aides said Mr.~Obama, a Harvard-trained lawyer, zeroed in on the fact that the
arrest came after police confirmed that Professor Gates was in his own home.

But his use of the word ``stupidly'' at the news conference that evening generated angry responses
from Cambridge police, and some of his aides privately rued the word choice. Mr.~Obama, who said he
was surprised at the response, discussed the issue over dinner with friends at his home in Chicago
on Thursday during a quick trip there for a fund-raiser, according to people close to the family. On
Friday morning, they said, he also talked it through with Mrs.~Obama.

By then, the controversy had dominated White House staff meetings. Robert Gibbs, the press
secretary, had told reporters at 10 a.m. that Mr.~Obama had nothing more to say. Some advisers had
concluded the furor would not dissipate unless Mr.~Obama made another statement, while others were
wary of him revisiting the episode and particularly did not want him to apologize, they said.

During the morning, police union members held a news conference in Cambridge calling on Mr.~Obama to
apologize for demeaning Sergeant Crowley and suggesting it was Professor Gates who had made it a
racial incident.

``The facts of this case suggest that the president used the right adjective but directed it to the
wrong party,'' said Sgt. Dennis O'Connor, president of the Cambridge Police Superior Officers
Association.

Sgt. Leon Lashley, an African-American officer at the Gates house that day, separately told The
Associated Press that he supported Sergeant Crowley's actions ``100 percent.''

The police event contributed to what one White House aide called a ``critical mass,'' but aides said
it was not the deciding factor, noting that Mr.~Obama had not watched. Shortly after noon, Mr.~Obama
called his senior adviser, David Axelrod. ``I'm going to call Sergeant Crowley and then I think I
ought to step into the press room and address it,'' Mr.~Axelrod said he said.

The president dictated some thoughts intended to avoid directly blaming either the professor or the
officer, and speechwriters had less than two hours to craft remarks. Mr.~Obama called Sergeant
Crowley about 2:15 p.m. and they spoke for five minutes. He went to the briefing room to make his
statement, then called Professor Gates about 3:15 p.m.

Mr.~Obama said the issue was making it harder for him to focus attention on health care. ``I don't
know if you've noticed, but nobody has been paying much attention to health care,'' he said.

He did not apologize but softened his language. ``I continue to believe, based on what I have heard,
that there was an overreaction in pulling Professor Gates out of his home to the station,'' he said.
``I also continue to believe, based on what I heard, that Professor Gates probably overreacted as
well.''

Mr.~Obama described Sergeant Crowley as an ``outstanding police officer and a good man'' who has ``a
fine track record on racial sensitivity.'' But he said the incident showed that ``because of the
difficulties of the past, you know, African-Americans are sensitive to these issues.''

John Payton, president of the NAACP Legal Defense Fund, said that unlike white presidents who could
dance around racial issues, Mr.~Obama had to be direct. ``That's the whole difference. Bush could
punt. Obama can't punt,'' he said. ``This issue resonates with him.''

Christopher Edley Jr., a former adviser to President Bill Clinton on race issues and now law school
dean at the University of California, Berkeley, said the episode dispelled the ``rosy hopefulness''
stemming from Mr.~Obama's election ``in case anybody needed more evidence that we're not beyond
race.''

\section{A Dispute Over Obama's Birth Lives on in the Media}

\lettrine{T}{he} conspiracy theorists who have claimed for more than a year that President Obama is
not a United States citizen have found receptive ears among some mainstream media figures in recent
weeks.

Despite ample evidence to the contrary, the country's most popular talk radio host, Rush Limbaugh,
told his listeners on Tuesday that Mr.~Obama ``has yet to have to prove that he's a citizen.'' Lou
Dobbs of CNN said that Mr.~Obama should do more to dispel the claims. Larry King, also of CNN, asked
guests about it. Chris Matthews debated it with guests on MSNBC, and ``NBC Nightly News'' even did a
segment debunking the theory.

Cable news is often stretched for news in the summer, but the birth certificate case has been fueled
by the combustible combination of luck, compelling video, media-savvy doubters -- and an outlandish
topic.

The theory that Mr.~Obama was born in Kenya, his father's homeland, first took root among some
staunchly conservative elements. In response, the Obama campaign scanned the candidate's
``certification of live birth'' from Hawaii's department of health and published it on the Internet.
Numerous third parties have examined and confirmed the birth certificate and concluded that it met
the requirements for citizenship.

``This smear was thoroughly debunked during the election,'' said Eric Burns, the president of Media
Matters for America, a liberal media monitoring organization.

But advocates of the issue, who are sometimes termed birthers, called attention to themselves
through frequent calls to talk radio shows, e-mail messages to news organizations and a videotaped
question at a Congressional town hall. Since mid-July, Mr.~Dobbs has discussed it repeatedly on his
radio show and on TV, emphasizing that he believes Mr.~Obama is a citizen, while wondering in
particular why Mr.~Obama has not provided a fuller copy of his birth certificate so that ``all of
this nonsense goes away.'' Mr.~Dobbs's coverage has been criticized by Media Matters and other
organizations.

In an e-mail message to Mr.~Dobbs's producers on Thursday, Jon Klein, the president of CNN's
domestic networks, wrote that ``it seems this story is dead -- because anyone who still is not
convinced doesn't really have a legitimate beef.''

In an interview Friday, Mr.~Klein said the e-mail message should not be interpreted as an order to
stop debating the subject. He defended Mr.~Dobbs's broadcast, saying that ``what we do here all the
time is dig into the truth about all kinds of controversial issues.''

The claims about Mr.~Obama's citizenship have percolated on the Internet and on right-leaning radio
talk shows for months. But from more mainstream outlets, ``the coverage has been minuscule'' until
this month, said Philip J.~Berg, a lawyer who has challenged Mr.~Obama's citizenship in a number of
lawsuits, unsuccessfully to date.

Mr.~Berg added: ``If I could file a lawsuit against the national media, including The New York
Times, I would do that, because I think you've all done an injustice by not covering Obama.''

Phil Griffin, the president of MSNBC, said the week's recurrent cable news coverage was set off by
two back-to-back events: the introduction of a bill that would require presidential candidates to
provide a copy of their original birth certificate (it has nine sponsors in the House and one in the
Senate), and the video of a town hall held by a Republican representative that turned into a hearing
on Mr.~Obama's citizenship.

At the town hall in late June, the representative, Michael N.~Castle of Delaware, was angrily
questioned by a woman who said of Obama, ``he is not an American citizen, he is a citizen of
Kenya.'' Members of the crowd applauded, and Mr.~Castle was heckled when he said Mr.~Obama was
indeed an American citizen.

It took more than a week for the incident to receive widespread attention but since then, ``there's
been a vast uptick,'' in mainstream coverage, Mr.~Berg said. One notable exception is the Fox News
Channel, which has hardly addressed the controversy this week.

A number of MSNBC shows, mostly in mocking tones, have covered the birthers' claims in recent days.
On ``Hardball'' Thursday evening, the radio talk show host G.~Gordon Liddy asserted that Mr.~Obama
may be an ``illegal alien,'' leading the host Mr.~Matthews to hand Mr.~Liddy a photocopy of the
birth certificate.

Mr.~Griffin said the claims were legitimate to cover ``in that there's a segment of our population
that believes this and keeps bringing it up.''

``It's racist,'' Mr.~Griffin said. ``Just call it for what it is.''

A lot of journalists ``live with this issue; we get e-mails, we get asked about it,'' Brian
Williams, the anchor of the ``NBC Nightly News,'' said on Wednesday's broadcast.

He showed an excerpt from the YouTube video in a segment about the birthers that he said were
``spreading lies'' about Mr.~Obama.

``So many conspiracy theories, so little time,'' he concluded.

\section{Amazon Faces a Fight Over Its E-Books}

\lettrine{L}{ast} week, Jeffrey P.~Bezos, chief executive of Amazon, offered an apparently heartfelt
and anguished\cite{anguish} mea culpa\cite{meaculpa} to customers whose digital editions of George
Orwell's ``1984'' were remotely deleted from their Kindle reading devices.

Though copies of the books were sold by a bookseller that did not have legal rights to the novel,
Mr.~Bezos wrote on a company forum that Amazon's `` 'solution' to the problem was stupid,
thoughtless and painfully out of line with our principles.''

An apology was not enough for many people.

A growing number of civil libertarians and customer advocates wants Amazon to fundamentally alter
its method for selling Kindle books, lest it be forced to one day change or recall books, perhaps by
a judge ruling in a defamation\cite{defamation} case -- or by a government deciding a particular
work is politically damaging or embarrassing.

``As long as Amazon maintains control of the device it will have this ability to remove books and
that means they will be tempted to use it or they will be forced to it,'' said Holmes Wilson,
campaigns manager of the Free Software Foundation.

The foundation, based in Boston, is soliciting\cite{solicit} signatures from librarians, publishers
and major authors and public intellectuals. This week it plans to present a petition\cite{petition}
to Amazon asking it to give up control over the books people load on their Kindles, and to
reconsider its use of the software called digital rights management, or D.R.M. The software allows
the company to maintain strict control over the copies of electronic books on its reader and also
prevents other companies from selling material for the device.

Two years after Amazon first introduced the Kindle and lighted a fire under the e-books market,
there is increasing awareness of how traditional libraries of paper and ink differ from those made
of bits and bytes. The D.R.M. in Amazon's Kindle books, backed up by license agreements with
copyright holders, prevents customers from copying or reselling Kindle books -- the legal right of
``first sale'' that is guaranteed to owners of regular books.

D.R.M. has created a new dynamic between consumers and the vendors of digital media like books and
movies. People do not so much own, but rent this media. And the rental agreement can be breached by
the manufacturer at any time, sometime with little or no notice.

People are also worried that the very architecture of network-connected devices like the Kindle,
TiVo or iPod give tech companies unprecedented control over digital media and by extension, the free
exchange of ideas.

Once upon a time, retailers sold customers a product and then walked away after the transaction.
Today's specialized devices often keep an umbilical\cite{umbilical} cord to their vendor, loading
updates and offering convenient ways to make purchases. These devices also limit the extent to which
people can load independent software and customize their experiences.

Such tethered\cite{tether} systems provide significant advantages to the consumer. Companies can
keep their own records of what people buy and restore the content if it is inadvertently lost.
Device software can be kept up to date, and vendors can track what people buy and make personalized
recommendations for new material they might like.

Randal C.~Picker, a law professor at the University of Chicago, says he thinks Amazon was right to
delete the improperly sold versions of ``1984'' and argues such systems can also allow companies to
better enforce copyright laws. He notes that the harm to the Orwell book buyers was minimal, since
their money was refunded after copies were deleted from their Kindles.

``Because copyright infringement\cite{infringement} was poor and lax in the offline world, it should
also be that way in the online world? I don't understand that logic,'' Mr.~Picker said. ``The whole
point of moving online is that it creates new opportunities.''

But critics say that any device capable of interfering with how its owner uses media is potentially
dangerous. ``I worry that systems like these tethered appliances are gifts to regulators,'' said
Jonathan Zittrain, a professor at Harvard Law School and author of the book, ``The Future of the
Internet -- and How to Stop It.'' Mr.~Zittrain predicts that governments in some parts of the world
will want to use it ``like a line item veto for content,'' removing objectionable sentences or
chapters in some books.

``It could happen first in jurisdictions\cite{jurisdiction} like the United Kingdom, where there
isn't as rich a First Amendment tradition and where libel\cite{libel} suits happen much more
frequently,'' he said.

Whether or not people are bothered by these possibilities may in part be a function of their age, as
a new generation grows up with an implicit understanding of the rules around these networked devices
and learns to live with them.

``I'd like to live in a perfect world where I own this content and can do whatever I want with it,''
said Justin Gawronski, a high school student whose copy of ``1984'' was erased by Amazon, but who
recently declined when a lawyer asked him to join a class-action lawsuit over the
incident. Mr.~Gawronski said, ``This is probably going to happen again and we just have to learn to
live with it.''

\section{Final Frontier for Wireless Hard to Break Through}

\lettrine{W}{ith} a vast, unserved population, Africa used to be a no-lose proposition for mobile
phone operators looking for new customers. But after a wave of investment, the wireless industry's
final frontier is showing signs of age.

In Egypt, France T\'el\'ecom has been sparring\cite{spar} for two years with its junior partner,
Orascom, for control of MobiNil, the leading Egyptian mobile operator and France T\'el\'ecom's
largest African holding.

In Johannesburg, Bharti Airtel, the largest mobile operator in India, is in the second year of
on-again, off-again negotiations to merge with MTN, the largest mobile African operator, with 91
million customers in 21 countries. The companies have set a deadline of Friday to reach an agreement
or cut off talks.

And last Monday, Vivendi, the French media company that owns a majority stake in SFR, broke off
talks with Zain, the Bahraini group that is looking to sell its 16 sub-Saharan African phone
businesses serving 40 million people.

Clinching\cite{clinch} deals for African mobile phone business, once a relatively low-risk
proposition for growth-starved mobile operators, is not as easy as it was. That is especially true
in sub-Saharan Africa, where rising competition and the cost of providing service to remote areas
has eaten into profits.

``What we are seeing is an evolution in the valuation of African mobile phone companies,'' said
Kristoff Puelinckx, managing partner in Dubai at Delta Partners, an industry consultant that advises
Middle Eastern and African operators. ``Operators are investing heavily in their networks, but it is
not as clear as it was that the financial returns will be there in the end.''

The more sanguine\cite{sanguine} view of those returns has contributed to the intense haggling over
acquisition prices that has dragged out the Bharti-MTN merger in Johannesburg and the sale of Zain's
African operations, the second-largest on the continent after MTN. The hard bargaining is a change
from five years ago, when most African markets were underserved and the use of mobile devices was a
fraction of what it was in Europe and North America.

Investors have long since filled the vacuum, and many sub-Saharan markets have up to seven operators
competing for consumers who spend less than \$5 a month on service. In Botswana, cellphone
penetration exceeds 80 percent, and in South Africa, it has topped 100 percent.

But the growth is slowing in some markets. In Kenya, where competition has led to lower prices,
there has been little increase in cellphone use, said Andre Wills, an analyst at Africa Analysis, a
business consultancy in Johannesburg.

``We are still bullish\cite{bullish} on the African telecom sector,'' Mr.~Wills said. ``But some
markets are showing signs of maturity. The outlook isn't as rosy\cite{rosy} as it was.''

The challenge of making a profit in sub-Saharan Africa is one reason Zain, which also has 25 million
customers in six Middle Eastern countries, may be looking to sell, said Mr.~Puelinckx.

Zain's biggest investors, the Kuwaiti sovereign wealth fund and the Kharafi family of Kuwait, which
together own 40 percent of Zain, are hoping to profit from a sale of Zain's African businesses,
Mr.~Puelinckx said.

``But their perceptions of the value of Zain's African businesses are far beyond what the market is
willing to pay at the moment,'' he added.

Antoine Abou Khalil, a Zain spokesman, said the company declined to comment on the sales talks.
Representatives for Orascom, France T\'el\'ecom and Vivendi also declined to comment.

Zain's sub-Saharan businesses generated 55 percent of its \$7.4 billion in 2008 sales, but only 7.3
percent of its \$1.2 billion profit, according to its financial report.

In the first quarter of this year, seven of Zain's sub-Saharan businesses reported a loss as the
company struggled to make money in very low-cost markets, said Nicholas Jotischky, an analyst at
Informa Telecoms \& Media in London.

``It's increasingly an issue as operators expand into rural areas, where most of the remaining
unpenetrated areas lie,'' Mr.~Jotischky said.

MobiNil, the Egyptian operator that started in 1998, has been consistently profitable and even sells
advanced services that are rare in Africa, like wireless broadband. That is why the two
shareholders, Orascom, owned by the Sawiris family, and France T\'el\'ecom are fighting to push each
other out.

The fight began in 2007, when Orascom objected to spending plans by MobiNil, triggering a series of
countervailing buyout offers that both sides have rejected. France T\'el\'ecom, which owns 72.1
percent of MobiNil, has offered to pay \texteuro528 million, or \$750 million, to Orascom for its
27.9 percent stake in the holding company.

But France T\'el\'ecom's attempt to gain control has foundered on compensation to minority
shareholders of MobiNil's listed operating company, ECMS, based in Cairo. The largest is Orascom,
with 20 percent. France T\'el\'ecom has offered to pay \texteuro1.47 billion to the shareholders,
representing an additional \texteuro600 million for Orascom.

But Orascom, and the Egyptian stock market regulator, have rejected the offer as too low. In March,
the International Arbitration Court in Geneva, which was asked by Orascom to decide the dispute,
sided with France T\'el\'ecom.

An adviser to France T\'el\'ecom who was not authorized to speak publicly described the situation as
a stalemate.

Frank Oehler, the head of business development for new markets including Africa at Nokia Siemens
Networks, a maker of wireless network gear, said negotiations like those taking place in Egypt are a
signs of the new wave of foreign investment changing the landscape of the African mobile industry.

At the same time, Mr.~Oehler said, all operators are under new pressure to justify the purchase of
new network equipment, given the slim profit margins that are a reality in sub-Saharan Africa.

``Investors have woken up to the costs of doing business in Africa,'' Mr.~Oehler said. ``The
pressure to control costs, which was always there, is now much greater than it was.''

\section{Ericsson Wins Auction\cite{auction} for Nortel Assets}

\lettrine{N}{ortel} Networks announced during the weekend that Ericsson of Sweden had won an auction
of the bankrupt Canadian telecommunications company's wireless technologies unit with an offer of
\$1.13 billion, thwarting\cite{thwart} a bid by the rival Nokia Siemens Networks.

Last month, Nortel signed a deal to sell the profitable division to Nokia Siemens for \$650 million.
A Canadian government-owned bank also agreed to provide \$300 million in financing for the
transaction. But other companies were invited to submit bids to a court overseeing Nortel's
bankruptcy protection process.

The final days of the auction had produced considerable controversy in Canada, where the sale was
front-page news. The maker of the BlackBerry, Research In Motion, which replaced Nortel as the
leading Canadian technology company, said earlier this week that it had been effectively, and
unfairly, shut out of the bidding. R.I.M. said that it had been prepared to offer about \$1.1
billion.

MatlinPatterson Global Advisors, a private equity firm and Nortel creditor based in New York, also
condemned the apparent piecemeal liquidation of Nortel and entered a bid in an attempt to keep the
company intact\cite{intact}. That breakup accelerated last week when Nortel announced a deal to sell
a division that makes communications systems for governments and large corporations to Avaya for
\$475 million.

After R.I.M. called on the Canadian government to investigate the sale of what was widely seen as
Nortel's most valuable operation, some Canadians argued that because Nortel had benefited
extensively from government assistance, those assets should remain under Canadian control.

It was not apparent if enough of the wireless unit's assets remain in Canada to trigger a mandatory
review of the sale under Canada's foreign investment laws. The Canadian government, citing privacy
laws, does not comment on that process. The sale must be approved by courts in Canada and the United
States.

In February, Verizon Wireless picked Ericsson and Alcatel-Lucent to provide equipment to increase
the speed of its data networks beginning next year. While Nortel has long been a major equipment
supplier to Verizon, that company was shut out.

Nortel's profitable wireless business, which provides a steady stream of revenues from contracts and
royalty payments from patents, is based on a declining technology known as CDMA. Once dominate in
North America, and largely ignored by Ericsson, that technology is being replaced by LTE. Nortel has
developed its own LTE systems which Ericsson, which is the world's leading producer of wireless
network infrastructure, will now be able to exploit that or at least prevent their use by
competitors.

Ericsson will rename the Nortel operation Ericsson CDMA and said that it would offer jobs to at
least 2,500 of its workers, about a third of whom are based in Canada.

\section{For Mozilla and Google, Group Hugs Get Tricky}

\lettrine{B}{oxes} lined the cubicles and hallways in the offices of Mozilla on a recent afternoon,
and its chief executive, John Lilly, seemed a bit disoriented as he looked for a place to sit.
Mozilla, which makes the Firefox Web browser, had just moved from one end of this city to the other,
mainly to gain more space for its growing work force.

Yet it was hard not to read symbolism into the move. Mozilla's old offices were next door to
Google's sprawling headquarters. For several years, Google has been Mozilla's biggest ally and
patron\cite{patron}. But in September, it also became Mozilla's competitor when it unveiled its own
Web browser, Chrome.

So it seemed only natural for Mozilla to move out from under Google's shadow.

``We've learned how to compete with Microsoft and Apple,'' says Mr.~Lilly, a soft-spoken, earnest
38-year-old. ``Google is a giant, of course, and competing with them means we are competing with
another giant, which is a little tiring.''

Those big companies weren't giving much thought to browsers when Firefox was released in 2004, and
neither were most ordinary Web users. A browser was just a window onto the Web, and people used
whatever was already installed on a computer. Usually that meant Microsoft's Internet Explorer.

Since then, Firefox has captured nearly a quarter of the browser market by focusing on speed,
security and innovation. Its success is all the more remarkable because it was built and marketed by
a far-flung community of programmers, testers and fans -- mostly volunteers -- coordinated by a
nonprofit foundation. It is a shining example of the potential of open-source software, which anyone
can modify and improve, and its ascent is one of Silicon Valley's most unusual success stories.

In short, Mozilla showed the world that browsers matter. Now the challenge is to keep proving that
Mozilla matters.

The rise of Firefox unleashed a new wave of innovation and competition among browser makers.
Microsoft and Apple, which makes the Safari browser, have narrowed the gap with recent upgrades.
That makes it less likely that people will take the trouble to seek out and install Firefox.

At the same time, the Web has been migrating from PCs to powerful mobile phones like the iPhone.
Firefox won't have a mobile version ready until later this year.

And then there is Google. After introducing Chrome, a lightning-fast browser designed to run
increasingly complex Web applications, Google upped the ante. This month it said it would put Chrome
at the center of a new operating system -- the software that handles the most basic functions of a
PC.

``Google, Apple and Microsoft can all throw a lot of resources toward improving their browsers.
Mozilla, not so much,'' says Rob Enderle, principal analyst at the Enderle Group. ``When it was them
against Microsoft, it wasn't such a big problem. Now that there are other alternatives, it becomes
harder for them to retain relevance.''

Despite Mozilla's mighty and increasingly competitive rivals, the spread of Firefox has continued
unabated. Nearly 300 million people around the world use it, making Firefox not only the most
successful open-source consumer product, but also one of the most successful software programs ever.

To a large extent, that success sprang from a disparate community that coalesced around Firefox and
was harnessed by Mitchell Baker, Mr.~Lilly's predecessor. Ms.~Baker, whom Mr.~Lilly calls the
``conscience'' of Mozilla, remains its chairwoman and is actively involved in managing it.

Ms.~Baker, 52, seems to embody Mozilla's anticorporate ethos\cite{ethos}. Unlike the clean-cut
Mr.~Lilly, Ms.~Baker has a decidedly counterculture look. Her hair, dyed a reddish color, is closely
cropped on one side only, and she is prone to wearing sandals with hiking socks.

She organized a recent meeting of nonprofit groups at Mozilla that ended with what she called a
``psychic group hug'' -- not a literal embrace, but a chance for everyone there to describe in one
word how they were feeling.

For Mozilla and its millions of fans, Firefox is not just cool software but also a cause: to ensure
that no company, whether Microsoft, Google or anyone else, can tilt the Web to its advantage by
tweaking its browser to favor its products or applications. Microsoft appeared ready to use that
tactic after its browser vanquished Netscape's.

Ms.~Baker envisioned Firefox as a counterweight to that possibility, and a way to make browsing
safer. In the years after Netscape's demise, Microsoft essentially stopped improving Internet
Explorer, and the browser quickly became vulnerable to security threats, an explosion of pop-up ads
and other annoyances.

Firefox was faster, safer and blocked pop-ups. It also offered some compelling innovations, like
tabs, which allowed users to have multiple pages open inside a single browser window. Word of its
virtues spread quickly, first in the tech world, then through a rapidly expanding fan base.

Today the model remains the same. Only a small fraction of the people involved in building Firefox
are paid employees at Mozilla, which has about 250 workers. An additional 1,000 or so programmers
contributed code for the most recent Firefox release. There are also tens of thousands of other
volunteers who help test and promote Firefox, write add-ons and help translate it into more than 70
languages.

``We succeeded because more people got engaged, helped us build a better product and helped us get
the product into the hands of people,'' Ms.~Baker says. ``We succeeded because of the mission.''

That community's fervor was heard in a recent weekly conference call for engineers who build
Firefox. The call, moderated but open to anyone, had 30 participants, some in Mozilla offices and
others scattered worldwide.

The conversation bounced from one technical topic to another: user interfaces, bug tracking,
security. Then someone interrupted to announce that in just a few hours, one million people had
downloaded the new Firefox 3.5. The group erupted in cheers. (By the next afternoon, the number had
topped six million.)

``Mozilla is about a community coming together and saying it can compete with the largest software
company in the world,'' says Sandeep Krishnamurthy, director of the business administration program
at the University of Washington at Bothell, who recently wrote a paper on Firefox's success. ``There
really is nothing like it.''

Across the globe, Vineel Reddy, 21, an engineering student in Hyderabad, India, basks in the
satisfaction of having contributed to that success. Drawn to the Firefox mission, but not
particularly good at programming, he decided to work on promoting Firefox. He rallied some friends
who were skilled at video editing, rented a camera and produced a flashy clip that compares Firefox
to a slick sports car.

The video has been seen more than 33,000 times, and Mr.~Reddy says he gets daily e-mail messages
thanking him. ``This is the best experience I have had,'' he says.

Mr.~Lilly says it was that kind of dedication from volunteers that enticed him to move from a career
in venture capital and to join Mozilla in 2005.

``As Americans we've lost the sort of civic engagement, the participation in making the world what
we want and what we think it should be,'' says Mr.~Lilly, who became chief executive last year.
``That, as a mission and as a product ethos, resonated\cite{resonate} with me.''

MR. LILLY readily acknowledges that Google's entrance into the browser market rocked the Mozilla
boat. ``Life was simpler before they did this,'' he says.

That said, Chrome's release does not signal a return to the browser wars of the 1990s, when
Microsoft poured resources into crushing the upstart Netscape.

For starters, Mozilla and Google have long had an agreement that makes Google the standard home page
when people start Firefox, and sends them to Google when they type something into the search box at
the top of the browser.

Google pays Mozilla hefty fees in return.

The deal accounted for 88 percent of Mozilla's \$75 million in revenue in 2007, according to its
most recent tax filings, and it was recently renewed through 2011. (The gusher of income from Google
prompted the nonprofit Mozilla Foundation to set up a taxpaying subsidiary, the Mozilla Corporation,
in 2005.)

The deal has helped Google gain market share in search. But Google has also been among the biggest
beneficiaries of Firefox's success in other ways.

Google's fortunes are tied to the advancement of the Web. As browsers become faster, more
standardized, more secure and more capable of running complex Web applications, Google's services,
like search, Gmail, maps and office software, become easier to use and more popular, and Google
earns more money. Firefox's innovations have helped with this.

``Mozilla has done an amazing job,'' says Sundar Pichai, a vice president of engineering at Google
who heads the development of Chrome.

Mr.~Pichai says that because of Mozilla's vital role, the company thought long and hard about the
impact Chrome would have on Firefox. Google eventually came to believe that it could help spur even
more innovation in browsers by building its own, he says. And it made Chrome open-source, so any
advances it makes could be adopted by others, including Mozilla.

``We were all very clear that if the outcome was that somehow Mozilla lost share to Google, and
everything else remained the same, internally, we would have been seen as having failed,''
Mr.~Pichai says.

So far, Chrome doesn't appear to have hurt Firefox. Chrome has grabbed just under 2 percent of the
browser market, according to Net Applications, a company that tracks browser use. During the same
period, Firefox's share has kept growing, to 22.5 percent from 19.5 percent. Microsoft's has
continued to decline, to 66 percent from 72, though it argues that most of that loss has been on
computers that don't readily support Internet Explorer, like those from Apple.

In many ways, Google and Mozilla are fighting the same battle, albeit not with the same objectives.
They both contend that the Web should be open and based on common standards -- Mozilla because it is
its mission, Google because it is good for its business.

``Most days we are aligned with them,'' Mr.~Lilly says. ``Their focus on the open Web is pretty
amazing.''

But he says financial pressures may someday push Google to start using Chrome to favor its own
services. That danger, he says, ``clarifies for me how important it is for independent organizations
like Mozilla to exist.''

Few at Mozilla fear that Google will pull its financial support. And even if it did, Mozilla says
that it could find another company willing to cut a similar deal. The real challenge for Firefox is
that over time, as Google and others invest in developing and marketing their browsers, the mission
that drove Firefox's success will lose urgency.

``Mozilla performed a really good service, but you have to wonder what their relevance is going to
be going forward,'' says Matt Rosoff, an analyst at Directions on Microsoft, an independent firm
that tracks the company. ``They keep Microsoft honest. But if Google is pushing innovation in its
own browser, it can play that role.''

Mr.~Lilly acknowledges that with the intense competition, it will be harder for Firefox to stand
out. Firefox and its supporters will have to work harder to find areas of browser technology where
they want to focus their efforts, he says, and let Google, Apple and Microsoft push innovation in
other areas.

But he says the renewed competition is a testament to the success of Firefox and of the Mozilla
mission.

``This is the world we wanted, and the world we made,'' he says. ``We wanted a world where people --
normal human beings -- could make meaningful choices about the browser technology that they use.
That's what we have today.''

\section{Files Vanished, Young Chinese Lose the Future}

\lettrine{F}{or} much of his education, Xue Longlong was silently accompanied from grade to grade,
school to school, by a sealed Manila envelope stamped top secret. Stuffed inside were grades, test
results, evaluations by fellow students and teachers, his Communist Party application and -- most
important for his job prospects -- proof of his 2006 college degree.

Everyone in China who has been to high school has such a file. The files are irreplaceable histories
of achievement and failure, the starting point for potential employers, government officials and
others judging an individual's worth. Often keys to the future, they are locked tight in government,
school or workplace cabinets to eliminate any chance they might vanish.

But two years ago, Mr.~Xue's file did vanish. So did the files of at least 10 others, all 2006
college graduates with exemplary\cite{exemplary} records, all from poor families living near this
gritty north-central town on the wide banks of the Yellow River.

With the Manila folders went their futures, they say.

Local officials said the files were lost when state workers moved them from the first to the second
floor of a government building. But the graduates say they believe officials stole the files and
sold them to underachievers seeking new identities and better job prospects -- a claim bolstered by
a string of similar cases across China.

Today, Mr.~Xue, who had hoped to work at a state-owned oil company, sells real estate door to door,
a step up from past jobs passing out leaflets and serving drinks at an Internet cafe. Wang Yong, who
aspired to be a teacher or a bank officer, works odd jobs. Wang Jindong, who had a shot at a job at
a state chemical firm, is a construction day laborer, earning less than \$10 a day.

``If you don't have it, just forget it!'' Wang Jindong, now 27, said of his file. ``No matter how
capable you are, they will not hire you. Their first reaction is that you are a crook\cite{crook}.''

Perhaps no group here is more vilified\cite{vilify} and mistrusted than China's local officials, who
shoulder much of the blame for corruption within the Communist Party. The party constantly vows to
rein them in; in October, President Hu Jintao said a clean party was ``a matter of life and death.''

Critics contend that China's one-party system breeds graft that only democratic reforms can check.
But China's leaders say the solution is not grass-roots checks on power, but smarter oversight and
crime-fighting.

Public policy specialists say China is shifting its emphasis from headline-grabbing corruption cases
to more systematic\cite{systematic} ways to hold officials accountable. The government opened an
anticorruption hot line last month to encourage whistle-blowers. A few localities require that
officials disclose their family assets to the party.

But in Wubu, a struggling town of 80,000 banked by steep hills and coal mines, citizens say that
local officials answer to no one, and that anyone who dares challenge them is punished.

``When the central government talks about the economy and development, it sounds so great,'' said
Mr.~Wang, the day laborer. ``But at the local level, corrupt officials make all their money off of
local people.''

Student files are a proven moneymaker for corrupt state workers. Four years ago, teachers in Jilin
Province were caught selling two students' files for \$2,500 and \$3,600; the police suspected that
they intended to sell a dozen more. In May, the former head of a township government in Hunan
Province admitted that he had paid more than \$7,000 to steal the identity of a classmate of his
daughter, so his daughter could attend college using the classmate's records.

While not quite as important as in Communist China's early days, when it was a powerful tool of
social control, the file, called a dangan, is an absolute requirement for state employment and a
means to bolster\cite{bolster} a candidate's chances for some private-sector jobs, labor experts
say. Because documents are collected over several years and signed by many people, they are
virtually impossible to replicate.

So in September 2007, when one Wubu graduate sought work at a local bank and discovered that his
file was gone, word spread fast. For the next two years, his parents and a group of other parents in
similar straits said, they sought help at every level of the bureaucracy.

The government's answer, they said, was to reject any inquiry, place the graduates' parents under
police surveillance and repeatedly detain them. Last February, they said, five parents trying to
petition the national government were locked in an unofficial jail in Beijing for nine days.

``We are so exhausted,'' said one tearful mother, Song Heping. ``Our nerves are about to snap from
this torture. The officials who were responsible not only have not been punished, they have been
promoted.''

Wubu officials did not respond to repeated inquiries. One Chinese television journalist said they
told him they had resolved the matter simply by creating new folders. But families say the folders
held nothing but brief, error-riddled r\'esum\'es that employers reflexively reject as fake.

The parents are uniformly poor: one father drives a three-wheel taxi, earning just 15 cents per
passenger.

Mr.~Xue's parents sacrificed even more than most, in the belief that education would lead their
children out of poverty. They earn just \$450 a year growing dates, and live near a dirt mountain
path, drinking well water and cooking over a wood fire.

Mr.~Xue, the oldest child, wore secondhand clothes and skipped meals throughout high school. When he
won admission to a university in Xian, 400 miles away, his parents borrowed to cover the \$1,500 in
annual expenses. Initially, it seemed the bet would pay off: he said he had had a chance to work at
an oil company with a monthly salary of \$735.

But the job evaporated\cite{evaporate} with his dangan. ``It was a catastrophe\cite{catastrophe},''
he said. Now he earns a base salary of \$90 a month as a door-to-door salesman and lives in a tiny,
dingy room in a Xian slum.

The woman he hoped to marry left him because her parents said he would never have a stable job. His
mother suffered a nervous breakdown, and the family debt ballooned. his father, Xue Ruzhan, said he
owed more than \$10,000 -- more than twice what his property is worth.

``What is the point of continuing to live?'' the father said. ``Sometimes I want to commit suicide.
These corrupt officials destroyed all our hopes.''

Including, it seems, the hopes of Longlong's younger sister, Xiaomei, an 11th grader who once
thought she would follow him to a university degree.

No more. ``I want to quit,'' she said during a school lunch break. ``My brother graduated from
college. What good did it do him?''

\section{China Steel Executive Killed as Workers and Police Clash}

\lettrine{C}{hina}'s state-run press confirmed Monday that a riot broke out at a steel mill in north
China Friday evening, leaving the executive of another steel mill dead.

The report, in the English-language China Daily, provided few details on the mayhem\cite{mayhem},
but a report on Saturday by a Hong Kong-based group, the Information Center for Human Rights and
Democracy, which broke the story on the riot, said that at least 30,000 workers were involved and
that about 100 people were wounded.

The riot, at the Tonghua Iron and Steel Works in Jilin Province in northern China, broke out after a
visiting steel executive from a related company threatened mass layoffs at the Tonghua steel mills
as part of a major restructuring of the state-owned company, China Daily said.

The riot followed a pattern of massive demonstrations that have taken place in various parts of the
country over the past few years, many involving citizens outraged over government corruption or
threatened with layoffs or orders to relocate.

The China Daily report said Chen Guojun, the steel executive who was beaten to death, had threatened
3,000 Tonghua steelworkers with layoffs, which he had said could take place within three days. He
also had signaled that larger jobs cuts were likely at the struggling steel mill.

The report said the rioters blocked the police, ambulances and government officials from reaching
Mr.~Chen before he died.

Mr.~Chen was general manager at Jianlong Steel, a large, privately owned company based in Beijing.
Jianlong had acquired a large stake in Tonghua, a state-owned company with as many as 50,000
employees, and Jianlong was working to restructure it.

The China Daily report quoted a police officer identified only as Wang as saying, ``Chen
disillusioned workers and provoked them by saying most of them would be laid off in three days.''

The officer also said that Mr.~Chen's warning that 30,000 jobs would eventually be cut to 5,000
``infuriated the crowd.''

China's steel industry, the world's largest, is just beginning to recover from a sharp downturn that
took place in the second half of last year, forcing many smaller mills to halt production.

The government has been trying to restructure and consolidate the steel industry, which is seen as
inefficient and plagued by corruption.

\section{IPhone Maker in China Is Under Fire After a Suicide}

\lettrine{W}{hen} a closely guarded prototype of a new Apple iPhone went missing at a huge factory
here two weeks ago, an internal investigation focused on a shy, 25-year-old employee named Sun
Danyong.

Mr.~Sun, a college graduate working in the logistics department, denied stealing the iPhone. But he
later complained to friends that he had been beaten and humiliated by the factory's security team.
On the night he was questioned, he sent an anguished\cite{anguish} text message to his girlfriend.

``Dear, I'm sorry. Go back home tomorrow,'' he wrote, according to a message she later posted
online. ``I ran into some problems. Don't tell my family. Don't contact me. I'm begging you for the
first time. Please do it! I'm sorry.''

Soon after, in the early-morning hours of July 16, Mr.~Sun apparently jumped to his death from the
12th floor of an apartment building in what his employer, Foxconn Technology, says was a suicide.

Apple and Foxconn, one of the world's biggest manufacturers of consumer electronics and a major
Apple supplier, issued statements last week expressing sorrow for the death. Foxconn said it
suspended one security officer, pending a police investigation, and that the company was now
considering counseling services for its employees.

The Apple statement said: ``We are saddened by the tragic loss of this young employee, and we are
awaiting results of the investigations into his death. We require that our suppliers treat all
workers with dignity and respect.'' The company would not comment further.

The local police bureau declined to answer questions about the case. But reports of the apparent
suicide have set off a firestorm of criticism of Foxconn's treatment of Mr.~Sun, labor conditions at
its factories and the pressures Apple places on suppliers to abide by the culture of secrecy that
surrounds its development of new products.

The case also underscores the challenges that global companies face in trying to safeguard their
designs and intellectual property in the hotly contested smartphone market, particularly here in the
southern Chinese city of Shenzhen, an electronics manufacturing center known for piracy and
counterfeiting\cite{counterfeit}.

Apple's popular iPhone is already widely imitated and counterfeited in China. And there are regular
rumors on Chinese Web sites about new Apple prototypes leaking out of Chinese factories.

``When you outsource to a third party, you lose some control,'' says Dane Chamorro, general manager
in China at Control Risks, a global consulting firm. ``And if you're outsourcing to China, it's
going to be even more challenging. There's going to be a bounty\cite{bounty} on every design.''

Labor rights groups say the worker's death should compel Apple to improve conditions at its supplier
factories in China and prevent worker abuse.

Foxconn, part of Taiwan's Hon Hai group, has also been sharply criticized because of suspicions
about unduly harsh treatment of the worker.

Foxconn, which produces electronics for some of the world's best-known brands, like Sony and
Hewlett-Packard, operates a cluster of sprawling factories in southern China. One of its Shenzhen
campuses has nearly 300,000 workers.

But some labor rights activists say the company treats employees harshly, routinely violating labor
laws.

In an e-mail message on Thursday, China Labor Watch, which monitors Chinese factories and is based
in New York, blamed Mr.~Sun's death on ``Foxconn's inhumane and militant management system, which
lacks fundamental respect for human rights.'' The group said it published an in-depth study of
Foxconn last year, detailing its abuses.

James Lee, general manager of China operations at Foxconn, defended the company's labor practices in
a lengthy interview on Friday, and also said the company would strive to improve management of its
facilities.

``It's very difficult for the company to defend itself against such charges,'' Mr.~Lee said of
complaints from labor rights groups. ``You're welcome to look at how employees are treated here.''

A reporter toured two of the company's campuses in Shenzhen on Friday, including the one where
Mr.~Sun worked. The campuses were so large they contained retail stores, banks, post offices and
high-rise dormitories with outdoor swimming pools.

The reporter was not allowed to see manufacturing lines because the company said it had to protect
trade secrets.

Outside the gates of one campus, most workers interviewed independently of the company said they
were well treated. One of about 15 workers questioned admitted to being forced to work overtime
above the legal limit.

In his interview, Mr.~Lee, the Foxconn manager, said the company also had a duty to protect the
intellectual property of its customers, and that it was honestly seeking answers to what happened to
the product.

Foxconn said it still does not know what happened to the missing iPhone. The company said Mr.~Sun
was given 16 prototypes on July 9 or 10 to deliver to research and development, and failed to report
one missing until three days later.

The company says his explanation for the missing phone did not seem credible and that he had had
problems before.

``Several times he had some products missing, then he got them back,'' Mr.~Lee said. ``We don't know
who took the product, but it was at his stop.''

In an interview with Southern Metropolis Daily newspaper last week, the security officer suspended
by Foxconn denied beating Mr.~Sun, saying only that he ``became a little angry'' and grabbed
Mr.~Sun's right shoulder.

Even so, the company paid compensation to Mr.~Sun's family. It declined to say how much, but
Mr.~Sun's brother cited a figure of 300,000 renminbi, or more than \$44,000, and said Mr.~Sun's
girlfriend was also given an Apple laptop computer.

Mr.~Sun's brother doubts he stole the prototype.

``He was honest and modest. He would never steal anything,'' said Sun Danxiong, 28, his brother.

Mr.~Sun grew up in a small, impoverished village in southwest Yunnan province and ranked first in
his high school, his family says. He graduated from the Harbin Institute of Technology, one of the
nation's top schools, before joining Foxconn about a year ago.

On Thursday, with his son Danxiong standing nearby, holding a box with Sun Danyong's ashes, the
father, Sun Yangdong, said Foxconn had treated the family well. But he said he was still in shock
that his son could leap from a building because he was so gentle and tender.

Soon after, a security guard, who was joined by two men wearing Foxconn shirts, threatened to ``beat
up'' a journalist's translator if she persisted in asking the family questions. Foxconn officials
later said the guard was not on their staff and might have been with the police bureau.

Back in Yunnan, Mr.~Sun said that on the night of his brother's death, he had e-mailed friends,
angry about Foxconn's questioning of him. In one message, Mr.~Sun said he was locked up and beaten.
``A Fortune 500 company even has these things,'' he wrote.

On Sunday, Danxiong said some of his brother's friends told him Mr.~Sun killed himself out of anger
at Foxconn. His brother said: ``They told me he was extremely angry at Foxconn; they humiliated him
and he wanted to resist the company, and planned to do something big.''

\section{`Nightline' Is Thriving in Late Night}

\lettrine{I}{n} the heavyweight battle being waged in late-night television, an awful lot of rounds
lately have been going to a third man in the ring, the one not slugging it out with punch lines.

While most of the attention, starting in June, has focused on the fight between the new host of
NBC's ``Tonight Show,'' Conan O'Brien, and his CBS rival, David Letterman, ABC's venerable news
program ``Nightline'' has managed mostly to fly under the radar -- in an upward
trajectory\cite{trajectory}.

Viewership for the news show is up 14 percent in the last six weeks compared with the same week a
year ago, and, in the most recent two weeks, the program has frequently grabbed the most viewers of
the three shows.

`` 'Nightline' is a rare example of the broadcast networks' counterprogramming; it almost never
happens outside of prime time,'' said Brad Adgate, senior vice president for research at Horizon
Media.

Ted Koppel, the longtime, widely celebrated anchor, left ``Nightline'' in November 2005, and
predictions of the show's imminent\cite{imminent} demise\cite{demise} were
prevalent\cite{prevalent}.

Instead the show thrived, by revamping its format to report three topics a night instead of one and,
most critically, by going back to producing editions live each weeknight.

``We do feel good,'' said James Goldston, the executive producer. ``This is kind of the
culmination\cite{culmination} of a process that has gone on for a long time, since we relaunched the
show.''

But, as Mr.~Goldston acknowledged, ``Nightline'' was never really a contender for most-watched
status in late night until Jay Leno left ``Tonight.'' ``I think there's another big shift happening
in late night,'' Mr.~Goldston said. David Westin, the president of ABC News, said, ``When Jay left,
some people were freed up to consider other alternatives.''

Everyone who works on ``Nightline'' knows it has been a marked show inside some quarters of ABC. A
desire to have an entertainment show, and the big profits that one could bring, has led to some
suggestions that ``Nightline'' might have to give up its 11:35 p.m. time slot. (These rumblings were
loud when ABC pursued Mr.~Letterman in 2002, and recurred, more quietly, recently when the network
recently weighed moving its own late-night star, Jimmy Kimmel, up a half-hour.)

In interviews, Mr.~Westin, Mr.~Goldston and Cynthia McFadden, one of the three ``Nightline''
anchors, all acknowledged an awareness that the program's performance remained under constant
appraisal. As Ms.~McFadden put it, ``We know we have to earn it every day.''

Mr.~Westin said, ``There are no guarantees, but we feel very good about where we are. The network is
very pleased with what they are getting with 'Nightline.' ''

(One factor not often mentioned is that ``Nightline'' has the advantage of being rated for just a
half-hour, while the other two shows run for a full hour at a time of night when many viewers are
drifting away to bed as the minutes roll by. The ABC program would win less often if its ratings
were compared to the first, much higher-rated half-hours of the entertainment shows.)

However the audience is parsed, ``Nightline'' is a real player in this high-stakes competition --
and it has every reason to want to get that story out. One way the show has been doing that lately
is with lots of Michael Jackson coverage.

The program offered reports related to Mr.~Jackson's death on 12 of 16 nights at one point, 10 times
as its lead story.

But as Mr.~Goldston pointed out, ``We've won on nights when we covered health care and Afghanistan
as well.'' (One health care night was big because of a potent guest, President Obama.)

``I think it's been helped by a heightened interest in news this summer from notable deaths,
confirmation hearings, the economy, health care,'' Mr.~Adgate said, ``but it's unusual it would
increase versus 2008 since it was such an interesting political year.''

Mr.~Goldston defended the extent of the Jackson coverage as justified by the intense interest in
what he called ``a cultural touchstone.'' One reason the show was able to pursue the Jackson story
so often was the connection another of its anchors, Martin Bashir, had with Mr.~Jackson. Mr.~Bashir
conducted one of the most memorable interviews with the singer.

Mr.~Goldston called Mr.~Bashir ``a formidable weapon'' during the Jackson coverage.

Ms.~McFadden said the Jackson story was ``of interest and importance'' to the public. And though
each addressed the Jackson issue separately, Mr.~Goldston, Mr.~Westin and Ms.~McFadden all made a
point of saying they would not apologize in any way for covering the Jackson story as much as
``Nightline'' did.

``It's wrong-headed and myopic to think our success is Michael Jackson-driven,'' Ms.~McFadden said.
``The notion that we're only benefiting because of Michael Jackson is laughable. We won't be No.~1
for long if that's the case.''

What has really made the recent success possible, Mr.~Goldston said, is a combination of increased
news interest from viewers, the opportunities afforded by the change at NBC and technological
changes in newsgathering.

``Everyone who works here is capable of shooting stories themselves and editing them,'' Mr.~Goldston
said. Ms.~McFadden said that she and one producer went to Afghanistan and did all the camerawork.

Mr.~Goldston said ``Nightline'' also ``got onto Twitter early on'' and used the service to get out
the message of what it was covering. Now more than 800,000 people follow the show on Twitter.

But mainly, he said, the show is succeeding because of ``a relentless focus on the news.'' Coming on
after local newscasts has always been to the show's advantage, though, before the transition to the
new format, ``Nightline'' was retaining only about 39 percent of the local news audience,
Mr.~Goldston said. Now the figure is up to 54 percent.

That is good enough for the show to be profitable, although ABC declined to offer any specific
figures.

``It's a profitable show,'' Mr.~Goldston said. ``We're not at the point of a Leno on the 'Tonight
Show' -- not even close. But it does make money -- and we can make more.''

\section{Clown of `Morning Joe' Gets a Show of His Own}

\lettrine{W}{illie} Geist, who starts a predawn shift on MSNBC on Monday, likes to joke that his
viewers will be insomniacs\cite{insomniac}, drug dealers and drunks.

But more people are awake at 5:30 a.m. than he might think. His new program, ``Way Too Early,'' is
opening another front in the fierce competition between cable news channels.

Mr.~Geist is a member of the ensemble on the three-hour ``Morning Joe,'' the laid-back morning show
with a concentration of fans in the Northeast Corridor from Washington to New York. On that show,
when he is not asking questions of guests like the White House press secretary, Robert Gibbs, he
plays the humorist with a news sensibility.

``I don't think anybody's quite accurately branded me. I'm not sure I could do it myself,'' he said.

It will become easier on Monday when the 34-year-old becomes host of his own TV show for the first
time.

``Way Too Early'' is envisioned as ``a pregame show to 'Morning Joe,' '' said Chris Licht, the
executive producer of both programs. It will be live at 5:30, while competitors like CNN, HLN
(formerly Headline News) and Fox News are running taped shows.

Although Mr.~Geist jokingly calls the time period ``obscene,'' Nielsen estimates that almost 17
percent of the American households with televisions, 19 million total, are tuned in at that time.

As a lead-in for ``Morning Joe,'' the show will try to build the early-morning audience for MSNBC.
From 5:30 to 9 a.m., the channel is currently ranked in fourth place among cable news channels for
the 25- to 54-year-old group, the most desirable to advertisers.

Mr.~Geist, the son of the CBS correspondent Bill Geist, is the rare producer to leap to an on-camera
position. He started in TV as an editor and a producer at CNN Sports Illustrated, the largely
forgotten cable sports news venture that was closed in 2002, and continued to work for CNN's sports
unit. Having harbored an interest in reporting on-camera, he said he decided to give it a try one
day while producing a segment in Dearborn, Mich.

``I realized that at some point, I just had to force their hand,'' he said.

In 2004 he moved to Fox Sports Net, where he helped write and produce a sports commentary show. The
program ended after just nine months, but nevertheless caught the eye of Rick Kaplan, who was the
president of MSNBC, a unit of NBC Universal, at the time.

When Mr.~Kaplan hired the conservative commentator Tucker Carlson for an 11 p.m. talk show, he asked
Mr.~Geist and three laid-off colleagues to imagine what the show could look like. Their ideas stuck,
and before long Mr.~Geist had found another opportunity for on-camera practice: he would dash out of
the control room at 11:55 p.m. for a pop-culture-detritus segment called ``The Cutting Room Floor.''

Mr.~Geist outlived Mr.~Carlson's show; he was sought by Joe Scarborough, the leader of the ``Morning
Joe'' troupe, as a sidekick character and he quickly grew into a full co-host.

On ``Morning Joe,'' where he will continue to appear, he manages to mix straight news and humor
every morning. On Thursday, he presented an over-the-top tribute to Gidget, the Chihuahua who had
starred in Taco Bell commercials and died last week at the age of 15. At the end of the three-hour
show, he tried to say ``good night sweet, sweet Gidget,'' prompting gales of laughter from the
show's co-hosts.

Mr.~Geist is perhaps best known for such segments, deemed ``News You Can't Use'' on ``Morning Joe,''
which good-naturedly mock the conventions of TV news. ``I always like overplaying something,'' he
said in an interview.

But the 5:30 a.m. show, he said, will be ``a news show at its core.'' Because of the early hour he
will be calling guests at home; the first guest will be Mr.~Gibbs.

Phil Griffin, the president of MSNBC, said he would like ``Way Too Early'' to be a fast-paced
``SportsCenter of news,'' something for ``an aggressive businessperson who is on the treadmill, who
wants a download of information.''

The before-dawn hour may be a testing ground for the type of show Mr.~Geist wants to do in the
future. But he rebuffed speculation about a shot at prime time. Mr.~Griffin said only that ``I want
to give Willie a platform.''

\section{Site Wins Fashion Fans by Letting Them Design}

\lettrine{T}{he} fashion magazines Vogue, InStyle and Lucky may rule the newsstand racks. But
online, they are also-rans, overlooked by the fashion-conscious in favor of Polyvore, an upstart Web
site far from Fifth Avenue.

Polyvore is a user-generated fashion magazine filled with user-generated ads. The people who go to
it play fashion editor and create collages featuring pictures of clothes, accessories and models
from across the Web. Readers view the collages, which the site calls ``sets,'' and if they click on
a dress or necklace, they are taken to the Web site that sells it.

Founded by three ex-Yahoo engineers, Polyvore has been focused on getting people to visit the site.
It seems to be working. Polyvore had more than 835,000 unique visitors in June, almost 25 percent
more than the traffic to Style.com, run by Vogue, and InStyle.com, according to Compete, a Web
analytics firm. It is also far bigger than the Web sites of Lucky and Harper's Bazaar. While other
fashion magazine sites have been struggling to hold an online audience, Polyvore has tripled its
traffic in the last year.

Now it is shifting its focus to making money. It runs ads, like the magazine sites do. But it also
earns a commission when users click on or purchase clothes from certain e-commerce sites, though
only about a quarter of outbound clicks make money for Polyvore. The company is now trying to forge
relationships with clothing and accessories sites in return for uploading their product catalogues
to the site.

``To compete with the Vogues and InStyles, it's not just about transacting and maximizing the
dollars per page view,'' said Peter Fenton, a partner at Benchmark Capital, which invested \$2.5
million in Polyvore. ``There's this aspirational side and entertainment side, which none of the
sites up until now have done a good job at tapping into.''

At the same time, Polyvore is giving stale e-commerce sites a much-needed jolt of inspiration.
``Online retail started around digital cameras. Now, sites are using the same engines to sell
shirts, but that's not the way they should be sold,'' said Pasha Sadri, co-founder and chief
executive of Polyvore. ``Clothing is much more visual.'' And unlike cameras or books, clothes and
accessories are bought as part of a whole outfit, though many sites still show pieces of clothing
individually.

Mr.~Sadri, a corduroy-wearer who is not particularly fashion-forward himself, conceived the idea for
Polyvore in 2007, while working at Yahoo. He was the software engineer behind Yahoo Pipes, a tool
that allows people to pull together content from across the Web, similar to what Polyvore does.

When Polyvore users surf the Web, they use a tool called the Clipper, downloaded and saved on the
toolbar, to select images and save them to Polyvore, where anyone can use them in a collage.

Brands and e-commerce sites can also upload their items to Polyvore, though today 95 percent of the
images come from users. Polyvore attributes images with a hyperlink to the original site. (It gets
about five requests a week, usually from photographers or painters, to remove images.)

To create a set, users drag and drop images and manipulate them. Polyvore also offers fonts for
text, and audio clips from Amazon. Its 928,000 registered users create 28,000 new sets a day. A set
inspired by Blake Lively, the ``Gossip Girl'' actress, includes a picture of her with a leather
jacket similar to the one she is wearing, available on the Zadig \& Voltaire Web site for \$840. It
also features accessories for the outfit, like Coach boots, Ray-Ban aviator sunglasses and a
T-Mobile Sidekick.

On a blog called Boutonnieres \& Bow Ties, brides-to-be submit photos of their wedding dresses.
Samantha Shih, owner of a menswear company called 9Tailor, uses Polyvore to create the groom's
outfit, complete with tie and cuff links, and uploads it to her blog.

To increase revenue, Polyvore is engaging fashion companies to sponsor content on the site. Tory
Burch promoted its summer collection, which was inspired by Venice, by running a contest on Polyvore
that asked users to create Venetian-themed sets using the new Tory Burch pieces.

Lori's Shoes, a boutique in Chicago and online, recently noticed a surge of traffic coming from
Polyvore, which John Coyle Steinbrunner, the company's creative director, had not heard of.

Now Lori's Shoes wants to advertise on Polyvore. ``It just gets us a better customer than banner ads
or cost-per-click ads, they're just instantly more loyal,'' Mr.~Steinbrunner said. ``As opposed to
just being a customer, they get to help determine what the look and aesthetic of the shoe is.''

Some online retailers, including Charlotte Russe and Torrid, license Polyvore's technology to use on
their own sites. Sucharita Mulpuru, an e-commerce analyst at Forrester Research, says she thinks
licensing is the most promising use for Polyvore and competitors like Kaboodle.

``On the retail sites themselves, people are looking to be inspired and are down the funnel of the
purchase process,'' she said. But while using Polyvore's site, ``people are in a different mind-set,
they're looking to browse, they might not be looking to buy.''

Polyvore also plans to sell data on customer preferences it compiles on the site. It could
potentially tell a retailer that a type of shoe is more popular in Manhattan than Los Angeles, so it
would know where to stock the shoe. Or designers could upload images of new items before deciding to
produce them to get input from fashion-savvy users.

It could also give buyers information about trends in real-time, faster than monthly magazines, said
Jess Lee, Polyvore's product manager. This fall, for example, watch for recent trends bubbling up on
the site: exposed zippers, fingerless gloves and butterfly prints.

\section{Recession Shadowing Chicago Bid for Games}

On a recent afternoon, Mayor Richard M.~Daley delivered his annual speech on the condition of the
city he has run for 20 years. Revenues may fall \$250 million short. Some city workers must take 15
unpaid days this year, including Mr.~Daley. More than 400 workers were laid off that very afternoon,
after talks with two unions collapsed.

In the same address, Mr.~Daley pressed forward with the city's efforts to host the 2016 Summer
Olympics, which carry an expected \$3.3 billion price tag. A decision by the International Olympics
Committee is due in October, and Chicago is considered a favorite among the four finalists.

Polls here suggest broad support for bringing the Olympic Games to the city. But increasingly, the
economic downturn is taking a central role in the local debate over the bid as more residents raise
concerns that Chicago taxpayers, already struggling, could be left paying the bills despite
assertions from organizers that no city dollars will be needed.

``How do we know?'' a resident, Douglas Brown, demanded of leaders of the Olympics bid during a
recent neighborhood meeting on the South Side.

``We can't take your word for it,'' Mr.~Brown said, adding, ``When do we get our guarantees to make
us sleep at night?''

At the same time, Mr.~Daley and other supporters of the Games argue that the Olympics would be a
force -- perhaps the force -- to lift Chicago from this financial gloom, with seven years of new
construction, jobs and tourism.

Asked about the difficulties of lobbying for an Olympics bid during a recession, Lori Healey, the
president of Chicago 2016, the bid committee here, said: ``I think it makes it easy. People are
hungry for jobs and opportunities.''

Earlier events that placed Chicago on an international stage, Ms.~Healey said, also came during
periods of financial gloom: a World's Fair in 1893 and again in 1933.

To hear Ms.~Healey and other bid leaders tell it, there is no downside. If the International
Olympics Committee were to choose Chicago over Madrid, Rio de Janeiro and Tokyo on Oct.~2, advocates
predict the Games would not only break even but would also make money (as have, they say, earlier
Olympics in the United States), generate more than \$22 billion in indirect economic impact on the
city and create \$1 billion in new tax revenue. Many of the sites needed for the events would not
require construction because they already exist.

Organizers say private financial support is mounting, with \$60 million raised so far for the bid,
and no city dollars are expected to be needed for either the bid or the Games.

On that last point, however, residents of Chicago seem skeptical. They have heard promises before.

This spring, a \$1.15 billion deal to privatize the city's parking meter system turned into a fiasco
after City Hall's inspector general called it a dubious financial deal and after motorists said they
poured money into fancy new meters that, in turn, spat out error messages. A few years ago,
Millennium Park, a downtown centerpiece, opened behind schedule and millions of dollars over budget.

``You all are projecting we're going to make a lot of money,'' a resident, Robin Kaufman, told
Olympics planners at a neighborhood meeting, one in a series intended to shore up support. ``But the
bankers were projecting they were going to make a lot of money. Bernie Madoff was predicting he was
going to make a lot of money.''

Ms.~Kaufman lifted a sign that read, ``No Blank Checks.''

At a high school auditorium on the West Side, where the bid leaders showed glossy Olympics
schematics and stood beside toned former Olympians, Stephanie Patton asked, ``Why should we trust
you?''

Even without a city having been picked for the 2016 event, Chicago agreed this summer to spend \$86
million on land for an Olympics village. Bid leaders say private developers will ultimately foot the
bill for a project that is to create permanent housing regardless of the Olympics.

Bid leaders say the Games themselves are expected to make \$450 million in profits. In case of a
shortfall, though, a ``safety net'' package for the Games will include an expected \$1 billion in
private insurance, according to bid leaders, as well as pledges from the State of Illinois of \$250
million and the City of Chicago of \$500 million.

Then last month, Mr.~Daley indicated that he would sign a host contract required by international
Olympics officials, including a standard blanket provision offering the city's backing -- one that
Chicago leaders had earlier said they hoped to modify.

The notion that the city could then be responsible beyond \$500 million ``set off some alarm bells''
for aldermen, one of them, Joe Moore, said. Mr.~Moore added that the City Council was seeking more
details before it signed off on it. And a group here opposed to the event, No Games Chicago, said
the prospect seemed to have suddenly stirred an outpouring of interest from people who had been
silent on the Olympics.

``The levee broke around that issue,'' said a founder of the group, Bob Quellos.

What strikes some residents as particularly puzzling is the bid committee's refrain that, as a
private nonprofit entity, it is separate from city government and public money. Technically, that
may be, but skeptics note the committee and City Hall share goals and often seem intricately
intertwined; Ms.~Healey, for instance, stepped down as the mayor's chief of staff to lead the
committee.

Crucial to maintaining residents' support for the Games, polls suggest, is convincing them that
their dollars will not be spent. In February, a poll by The Chicago Tribune found that 64 percent of
residents of Chicago and its suburbs favored having an Olympics but that 75 percent were against the
use of tax money to cover shortfalls.

At one of the community forums, a city official likened the efforts of those planning the Olympics
Games to Jean Baptiste Pointe DuSable, who is believed to have been the first non-Indian settler of
Chicago, and Daniel H.~Burnham, whose 1909 Plan of Chicago largely designed the city with the long
strip of lakefront parkland that defines it.

Then the bid leaders brought out their big guns for these crowds: a videotape of Barack Obama, then
a candidate for president, in Chicago, his hometown, smiling and shaking hands with Mr.~Daley during
a rally last summer for the bid and declaring his wish that an Olympics reveal Chicago as ``not just
a city that works, but a city that inspires.''

Patrick G.~Ryan, the founder of the Aon Corporation and the chairman of the bid committee, told one
crowd that the thought that private money being raised for the Games might otherwise pay for
schools, garbage collection or city workers' salaries was wrong.

``To say we could spend the money either hosting the Games or on something else is really a false
choice,'' Mr.~Ryan said. ``This money only comes in if we win.''

\section{Internship Is Secured, but a Seat? That's Tricky}

Susan Cho's internship at the Kings County district attorney's office this summer has presented some
different challenges than her previous stint working alongside prosecutors in Rochester. ``There's
sometimes trouble finding a seat,'' Ms.~Cho, 26, who is soon to be a second-year student at
Villanova University School of Law, said of her spell at 350 Jay Street in Downtown Brooklyn.

Last summer, 135 law students toiled as unpaid interns with the Brooklyn district attorney. This
year there are 195 -- a number that Carol L.~Moran, the director of legal hiring, believes is a
record.

The bumper crop of interns, whose placements end July 31, is the result of an increase in both
applications and the percentage of students who accepted offers, most likely a fallout of the
recession, as many corporate law firms have cut back on their summer programs. In addition to the
intensified competition to score a full-time -- and paid -- job after graduation, the intern glut
has put a premium on office chairs, desk space and meaty assignments.

``It's much harder for them to find stuff for us to do,'' said one intern who spoke on the condition
of anonymity for fear of alienating his bosses. ``Definitely some people feel they haven't done
anything.''

Ms.~Moran, the hiring director, said the Brooklyn office received about 500 applications for this
summer's positions, up from 400 in 2008. ``We made close to 200 offers last year, and ended up with
135,'' she said. ``This year we made slightly more offers -- close to 250. I was looking for 150
interns, and ended up with 195.''

Jerry Schmetterer, a spokesman for the district attorney's office, said the glut ``appears to be
economy related,'' saying that fewer students this year asked to defer the decision to accept an
unpaid internship -- which many did in the past while waiting to see if they had snared lucrative
employment.

The Queens district attorney has also seen an uptick in interns, with 83 (from 60 last year); but
Manhattan (about 50) and the Bronx (20) have their usual number this summer. A spokesman for the
Staten Island district attorney said the office had seen more interest, but kept hiring stable, at
27, to ensure proper supervision and space.

The popular internships come amid a brutal legal job market: according to the National Association
for Legal Career Professionals, the employment rate for 2008 graduates -- 90 percent as of February
2009 -- represented the first decrease since 2003.

``We know law firms hired smaller groups of summer associates,'' said James G.~Leipold, the group's
executive director. ``There's a lot of competition for a smaller number of opportunities.''

Meg Reuter, assistant dean for career planning at New York Law School, said that while students are
considering the public sector as an alternative to scarce corporate work, such jobs are not immune
to economic conditions.

``The federal government can deficit spend, so they can hire even when they don't have the tax
revenue, but state and local government cannot,'' she said.

In Brooklyn, Ms.~Moran said that although about 75 percent of the lawyers on Jay Street had once
been interns, only a few of this summer's cohort would get job offers. She said the office was now
hiring about 7 to 8 percent of applicants, down from 10 to 12 percent in previous years.

Those who supervise interns, and some of the interns themselves, said there has been enough work to
go around. ``I do research, draft motions, call witnesses,'' said Venice Choi, 22, who is working in
the homicide bureau.

But other interns pass the hours doing crossword puzzles or playing games on the computer; some have
taken to observing random court proceedings to fill their days.

Surviving an unpaid summer in New York is another challenge. Some interns received grants from their
schools -- typically about \$5,000 -- intended to support students with an interest in public
interest law. Others are living with relatives or working at paid jobs in May and August to make
ends meet.

``I'm receiving credit,'' said Stephanie Mishler, an intern in the homicide bureau who will start
her second year at Seton Hall University School of Law this fall. ``It'll pay off in the long run.''

The swell of interns has helped in at least one area: On Wednesday evening, the rackets division
took on the green zone trial bureau in a softball game at Prospect Park, and each team fielded a
brace of students.

Before the game, John Rudikoff, an assistant district attorney, said that one of the rackets
interns, Evan Gidez, 24, soon to be a third-year student at Brooklyn Law School, had already proved
his worth on the field.

``In the first game Evan made an amazing catch; he's an excellent ballplayer,'' Mr.~Rudikoff said.
``We didn't know that, but it would have given us even more reason to hire him.''

\section{With Jobs Scarce in Japan, Women Clamber\cite{clamber} to Become Professional Flirts}

\lettrine{T}{he} women who pour drinks in Japan's sleek\cite{sleek} gentlemen's clubs were once
shunned\cite{shun} because their duties were considered immodest: lavishing\cite{lavish} adoring
(albeit\cite{albeit} nonsexual) attention on men for a hefty\cite{hefty} fee.

But with that line of work, called hostessing, among the most lucrative\cite{lucrative} jobs
available to women and with the country neck-deep in a recession, hostess positions are increasingly
coveted\cite{covet}, and hostesses themselves are gaining respectability and even acclaim. Japan's
worst recession since World War II is changing mores.

``More women from a diversity of backgrounds are looking for hostess work,'' said Kentaro Miura, who
helps manage seven clubs in Kabuki-cho, Tokyo's glittering red-light district. ``There is less
resistance to becoming a hostess. In fact, it's seen as a glamorous job.''

But behind this trend is a less-than-glamorous reality. Employment opportunities for young women,
especially those with no college education, are often limited to low-paying, dead-end jobs or temp
positions.

Even before the economic downturn, almost 70 percent of women ages 20 to 24 worked jobs with few
benefits and little job security, according to a government labor survey. The situation has worsened
in the recession.

For that reason, a growing number of Japanese women seem to believe that work as a hostess, which
can easily pay \$100,000 a year, and as much as \$300,000 for the biggest stars, makes economic
sense.

Even part-time hostesses and those at the low end of the pay scale earn at least \$20 an hour,
almost twice the rate of most temp positions.

In a 2009 survey of 1,154 high school girls, by the Culture Studies Institute in Tokyo, hostessing
ranked No.~12 out of the 40 most popular professions, ahead of public servant (18) and nurse (22).

``It's only when you're young that you can earn money just by drinking with men,'' said Mari Hamada,
17.

Many of the cabaret\cite{cabaret} clubs, or kyabakura, are swank\cite{swank} establishments of dark
wood and plush cushions, where waiters in bow ties and hostesses in evening gowns flit\cite{flit}
about guests sipping fantastically expensive wine.

Some hostesses work to pay their way through college or toward a vocational degree, or to save up to
start their own businesses.

Hostessing does not involve prostitution, though religious and women's groups point out that
hostesses can be pressured into having sex with clients, and that hostessing can be an entry point
into Japan's sprawling underground sex industry.

Hostesses say that those are rare occurrences, and that exhaustion from a life of partying is a more
common hazards in their profession.

Young women are drawn nonetheless to Cinderella stories like that of Eri Momoka, a single mother who
became a hostess and worked her way out of penury\cite{penury} to start a TV career and her own line
of clothing and accessories.

``I often get fan mail from young girls in elementary school who say they want to be like me,'' said
Ms.~Momoka, 27, interviewed in her trademark seven-inch heels. ``To a little girl, a hostess is like
a modern-day princess.''

Even one member of the Japanese Parliament, Kazumi Ota, was a hostess. That revelation once would
have ignited a huge scandal, but it has not. She will run for re-election on the leading opposition
party ticket, the Democratic Party of Japan, in the national election next month, and the ticket is
expected to unseat the ruling party.

It is unclear how many hostesses work in Japan. In Tokyo alone, about 13,000 establishments offer
late-night entertainment by hostesses (and some male hosts), including members-only clubs frequented
by politicians and company executives, as well as cheaper cabaret clubs.

Hostesses tend to drinks, offer attentive conversation and accompany men on dates off premises, but
do not generally have sex for money. (Men who seek that can go to prostitutes, though prostitution
is illegal.)

Hostesses are often ranked according to popularity among clients, with the No.~1 of each club
assuming the status of a star.

Mineri Hayashi has made it to the top of her club, Celux, six years after coming to Tokyo from
northern Japan. One recent evening, she readied herself for an elaborate birthday event her club was
throwing in her honor.

Outside the club, bigger-than-life posters of Ms.~Hayashi adorned the street. At the club, a dozen
men put up balloons and lined up Champagne bottles.

The club's clientele is diverse, including workaday salarymen, business owners and other men
unwinding after work.

Celux hopes to make more than \$60,000 on Ms.~Hayashi's birthday party, which will be attended by
scores of regulars.

``Life has been fun, and I want to keep on having fun,'' Ms.~Hayashi said, placing a
tiara\cite{tiara} in her hair. She talks of plans to retire next year and travel abroad.

Her 17-year-old sister, who also wants to be a hostess, may succeed her. Ms.~Hayashi is supportive.
``I just want her to be happy,'' she said.

Popular culture is also fueling hostessing's popularity. TV sitcoms are starting to
depict\cite{depict} cabaret hostesses, or kyabakura, as women building successful careers. Hostesses
are also writing best-selling books, be they on money management or the art of conversation.

A magazine that features hostess fashion has become wildly popular with women outside the trade, who
mimic\cite{mimic} the heavily made-up eyes and big, coiffed hair.

But Serina Hoshino, 24, another Tokyo hostess, is exhausted from the late nights and heavy drinking.

Slumped in her chair at the M.A.C. hair salon, she talked about endless after-hours dates with
clients. Stumbling back home at dawn, she sleeps the rest of the day. On her days off, she hardly
leaves her apartment.

Her reward is about \$16,000 a month, almost 10 times the salary of most women her age.

``It's nice to be independent, but it's very stressful,'' Ms.~Hoshino said, speaking through a cloud
of hair spray and cigarette smoke.

In recent months, clubs have also started to feel the squeeze of the bad economy. Hostess wages are
starting to fall to as little as \$16 an hour. Still, that rate remains above many daytime jobs
here.

So, the young women keep coming. The Kabuki-cho district is lined with dark-suited scouts recruiting
women. One club recruiter said some women turn up to interviews with their mothers in tow, which
never would have happened when the job was less respectable.

``Women are being laid off from daytime jobs and so look for work with us,'' said Hana Nakagawa, who
runs a placement agency for higher-end clubs in Tokyo.

She gets about 40 inquiries a week from women looking for hostess jobs, twice as many as before the
downturn.

Atsushi Miura, an expert on the issue, says hostessing will be popular among Japanese women as long
as other well-paying jobs are scarce.

``Some people still say hostesses are wasting their life away,'' he said. ``But rather than
criticizing them, Japan should create more jobs for young women.''

\section{To Create Jobs, Tennessee Looks to New Deal Model}

\lettrine{C}{ritics} elsewhere may be questioning how many jobs the stimulus program has created,
but here in central Tennessee, hundreds of workers are again drawing paychecks after many months out
of work, thanks to a novel use of federal stimulus money by state officials. Here in one of
Tennessee's hardest-hit areas, some workers were cutting down pine trees with chainsaws and clearing
undergrowth on a recent morning, just past the auto parts factory that laid them off last year when
it moved to Mexico. Others were taking applications for unemployment benefits at the very center
where they themselves had applied not long ago. A few were making turnovers at the Armstrong Pie
Company (``The South's Finest Since 1946'').

The state decided to spend some of its money to try to reduce unemployment by up to 40 percent here
in Perry County, a rural county of 7,600 people, 90 miles southwest of Nashville where the
unemployment rate had risen to above 25 percent after its biggest plant, the auto parts factory,
closed.

Rather than waiting for big projects to be planned and awarded to construction companies, or for tax
cuts to trickle through the economy, state officials hit upon a New Deal model of trying to put
people directly to work as quickly as possible.

They are using welfare money from the stimulus package to subsidize\cite{subsidize} 300 new jobs
across Perry County, with employers ranging from the state Transportation Department to the
milkshake place near the high school.

As a result, the June unemployment rate, which does not yet include all the new jobs, dropped to
22.1 percent.

``If I could have done a W.P.A. out there, I would have done a W.P.A. out there,'' said Gov. Phil
Bredesen of Tennessee, a Democrat, referring to the Works Progress Administration, which employed
millions during the Great Depression.

``I really think the president is trying to do the right thing with the stimulus,'' Mr.~Bredesen
said, ``but so much of that stuff is kind of stratospheric\cite{stratospheric}. When you've got 27
percent unemployment, that is a full-fledged depression down in Perry County, and let's just see if
we can't figure out how to do something that's just much more on the ground and direct, that
actually gets people jobs.''

Tennessee is planning to pay for most of the new jobs, which it expects will cost \$3 million to \$5
million, with part of its share of \$5 billion that was included in the stimulus for the Temporary
Assistance for Needy Families program, the main cash welfare program for families with children. The
state did not wait for the federal paperwork to clear before putting residents of Perry County back
to work.

Other states are still drawing up plans for spending the welfare money, which is typically used for
items like cash grants for families and job training. Some are likely to use part of it to subsidize
employment, as Tennessee is doing, but it is hard to imagine many other places where the creation of
so few jobs could have such an immediate and outsize impact as it did in this bucolic\cite{bucolic}
county.

A stimulus job came just in time for Frank Smith, 41, whose family was facing
eviction\cite{eviction} after he lost his job as a long-haul truck driver. Then he landed a job with
the Transportation Department.

``The day I came from my interview here, I was sitting in the court up here where I was being
evicted,'' Mr.~Smith said after a sweaty morning clearing trees under a hot sun to make room for new
electric poles. ``Luckily I'm still in the same place. There's a lot of people that were totally
displaced.''

Scott and Allison Kimble married after meeting on the assembly line at the Fisher \& Company auto
parts plant. When the factory closed last year and relocated to Mexico, the Kimbles, along with many
of their friends and neighbors, found themselves out of work. Now Mr.~Kimble has a stimulus job
working for the Transportation Department, and Ms.~Kimble has one in what has become a growth
industry, taking telephone applications for unemployment benefits.

``I know what they feel like,'' she said between calls. ``I've been in their position.''

Michael B.~Smith, 53, who drove a forklift at the plant for 31 years, now drives a Caterpillar to
clear land for a developer. Robert Mackin, 55, who lost his job, his health insurance and his home,
now has a job with the Transportation Department, a rental home, health insurance and an added
benefit: the state employee discount when his daughter goes to a state college.

``With a degree, she can always go somewhere,'' Mr.~Mackin said.

The impact has been enormous, all across the county. Even the look of the place is changing,
following the old W.P.A. model. In addition to the jobs for adults, there are 150 summer jobs for
young people, some of whom have been working with resident artists to paint murals depicting local
history on the buildings along Main Street in Linden, the county seat.

Over all, two-thirds of the new jobs are in private sector businesses, which are
reimbursed\cite{reimburse} by the state for the salaries of eligible\cite{eligible} stimulus
workers. Some, in retail, might be hard to sustain when the stimulus money runs out in September
2010. Other businesses say the free labor will help them expand, hopefully enough to keep a bigger
work force.

The Commodore Hotel Linden, a newly restored 1939 hotel that has brought new life to downtown, has
seen an increase in its bookings since it has expanded its staff thanks to the stimulus. And the
Armstrong Pie Company expects to be able to keep on the new bakery assistants and drivers it hired
with stimulus money, saying the new workers have helped the company triple its pie production and
expand its reach through central Tennessee.

The county mayor, John Carroll, has been working to lure new industry to the area. Walking through
the cavernous\cite{cavernous}, empty Fisher plant, Mr.~Carroll pointed to a forgotten display case
filled with dozens of awards for safety and manufacturing excellence. ``What we can offer,'' he
said, ``is a great work force.''

Mr.~Kimble said the new jobs had given him and his wife paychecks, health insurance and a reason to
get up each morning. But he said he hoped that a big, long-term employer would move in soon.

``This job here is not a permanent fix,'' he said. ``We still need some kind of industry to look and
come into Perry County. But for right now we've got hope, and when you've got hope, you've got a
way.''

\section{Business-Class Bargains for Overseas Travelers}

\lettrine{F}{rom} the airlines' point of view, these are tough times for premium\cite{premium}
overseas travel. In May, according to the International Air Transport Association, global demand for
business- and first-class seats was down over 26 percent from May 2008. It was the 12th
consecutive\cite{consecutive} month of year-on-year declines in demand for those premium seats,
which at one time were a gold mine for airlines.

But a crisis for airlines has also created opportunities for business travelers bound overseas and
determined to make the trip in productive comfort (lots of space to work in and those cushy lie-flat
beds to rest in, for example) -- without buying a ticket that can cost as much as a good used car.

With a degree of flexibility in travel plans, you can now fly overseas in business class for a
fraction of what it used to cost. Airlines started introducing spot fare sales for international
premium seats last fall when demand fell in a deteriorating economy. Now, as conditions have
worsened, airlines seem to be settling into a basic restructuring of international premium-class
fares that resembles the way airlines have long priced leisure fares in the back of the plane.

``I don't know even what counts as a sale anymore,'' said Joe Brancatelli, who publishes
Joesentme.com, a subscription Web site for business travelers. ``Airlines are now yield-managing
fares up front like they do in the back. It's looking like a permanent sale environment.''

Airlines adopted yield-management strategies in the 1980s to sell coach seats efficiently. Yield
management assumes that airline seats are perishable and can be marked at
fluctuating\cite{fluctuate} prices based on calculations for demand, starting months in advance and
continuing until the time the plane takes off.

Because of reductions in corporate spending, business travelers have begun behaving more and more
like leisure travelers, eschewing immediate convenience for lower prices -- even in the
international premium niche.

I can't begin to summarize the current discount fare environment. But here are a few current and
arbitrarily chosen examples. Contemplating a business trip to Shanghai? With a little flexibility,
you can travel in luxury and sleep in a lie-flat bed. Air Canada has a sale in which a first-class
ticket between Los Angeles and Shanghai is available for less than \$3,500. A year ago, a traveler
could have easily paid over \$15,000 on various airlines for that ticket.

Business class from New York to London? For the British Airways swanky Club World cabin, that will
be about \$2,544 round trip for late summer travel, with an advance purchase. The walk-up fare for
Club World on that route used to be about \$11,000 and is still about \$7,500.

New York to Amsterdam? On the mostly business-class planes operated by OpenSkies, a British Airways
subsidiary, the summer fare is about \$1,300, round trip.

As always, you need to check individual airline Web sites for the fine print on advance purchase
restrictions and fees.

Or, given the current free-for-all in high-end fares, you may even want to consult with a real live
travel agent. That's because finding the right premium fares at the right time has become
cumbersome\cite{cumbersome} even for business travelers who used to book travel themselves in a less
frantic environment.

``We all know the airlines are reducing capacity, but that's going to take time,'' said Fran Kramer,
who specializes in international bookings for DePrez Travel, a big agency in Rochester with clients
all over the country. ``While they do that and figure out what works,'' she added, ``there are
certainly a lot of opportunities'' for finding major fare discounts on international premium travel.

Ms.~Kramer said that airlines, desperate for whatever revenue they can get, sometimes reach out to
agents with last-minute offers for premium seats that are unsold, even at promotional fares, just
before a flight. ``Yesterday, we got a call 24 hours in advance offering my client the option to go
to Europe in business class for \$500 extra,'' she said.

Airlines can't continue losing money indefinitely, but they're still flying those fancy cabins
overseas. So big discounts on premium travel will continue until supply can be brought into some
profitable relationship with demand. Even then, those stunning \$11,000 round-trip business-class
fares between, say, New York and London may be a thing of the past.

Airline executives have been looking ahead with reduced expectations.

\section{In First Class, but Treated Otherwise}

\lettrine{I}{ HATE} flying. But I like to see different places. So I usually just try to make the
best of the situation.

I'm an entrepreneur, so I do watch my travel budget and I generally fly coach or business class.
Like a lot of people, I fantasize about traveling first class. Just the name alone -- first class --
conjures up a fantasy of perfection.

I am so incredibly na\"ive.

My husband and I had to go to Mumbai for business. Although we were flying into London, with a
connection to Mumbai, we thought we should take the opportunity to get to Dubai, too. We knew that
some of the Hindi-language film people, or Bollywood types, were in Dubai so it would be smart for
us to do the same thing.

When we landed in London, we were very out of sorts and very hungry. But we had to make arrangements
to go to Dubai and then run to another terminal to catch our connection.

As we were checking in, the agent overheard my husband and me talking about the meeting we had with
the musical artist 50 Cent. He was a huge fan.

I made a joke that I would give him some 50 Cent gear that I had in my carry-on if he would upgrade
our tickets.

The next thing I knew we had two first-class boarding passes.

When we got on the plane we were ushered into first class, and I was excited.

But I don't think the attendant was too thrilled. Our ticket stubs showed that we were ``upgrades,''
not people who actually paid for the amenities\cite{amenity}.

The attendants were paying a lot of attention to everyone, bringing extra pillows, fluffing blankets
and doing everything in their power to make these other travelers comfortable. People were getting
refills for their wine glasses. I couldn't get a glass of water.

My husband and I were starved. I'm not a picky\cite{picky} eater, and I figured I couldn't go wrong
with ordering risotto when it was time to get some food. Unfortunately, the risotto came smothered
in cheese. I can't eat cheese; I get sick.

I told the attendant, and she looked at me like I was crazy. I couldn't get another meal. So I asked
for a snack, only to be told that snacks were served only in business or coach class.

At this point, I really didn't know what to do. So I just sat there like a good little passenger.

Once we were set for landing, the attendant did come up to me to chat. I thought she was going to
apologize or try to make nice. Instead, she said she wanted to know why the airline was trying to be
so nice to us by upgrading our seating arrangements. Apparently, the airline never gives upgrades,
so it must have done something really horrible to us if it gave us an upgrade.

I was no longer going to be the good little passenger. So I said the first thing that came into my
head. I told her my best friend was married to the airline's chairman, and we always get upgraded.

It was a complete lie, but it was worth seeing the look of fright on the attendant's face. I thought
she was going to faint. She scurried\cite{scurry} away to talk to another attendant, who was her
boss.

The next thing I know, my husband and I are surrounded by attendants offering us wine, food, fluffy
pillows and blankets. It was almost comical, especially since the flight was set to land in about 10
minutes.

I still want to believe that first class is, well, first class. So I would take an upgrade again.
Maybe next time, my first-class fantasy will come true.

\section{Friendly Skies, Revisited}

\lettrine{T}{here} was a time when airlines mimicked the hotel business, sending out their own
version of concierges in airport lobbies. But lately, the airport terminal has seemed more like a
test kitchen of ways to eliminate all human contact with customers.

Labor, after all, has traditionally been the airlines' largest expense, and the frontline airport
staff is especially endangered by innovations like online check-in and self-service kiosks. Since
2001, American carriers have shed 154,000 full-time jobs, according to the Air Transport
Association.

But some airlines, stung by a wave of negative publicity about their treatment of customers, are
reversing course. One is Delta Air Lines, now the world's largest carrier after its merger with
Northwest. It is reviving a service program known as the ``Red Coats,'' for the crimson-jacketed
airport service people who assist passengers in need.

Delta has already brought the agents back to Kennedy Airport, and last month it added more than 100
at its biggest airport hub, Atlanta. As of Aug.~1, about 500 of these agents will be in circulation
at more than a dozen airports, including terminals in Washington, Los Angeles, Boston, Cincinnati,
Memphis, Newark and New York.

The Red Coats program, which dates back to the 1960s, ``was one of our better-known services,'' Gil
West, Delta's senior vice president of airport customer service, said in an interview. The decision
to disband the program in 2005, while the airline reorganized under Chapter 11 bankruptcy, was
viewed as a mistake by some employees, he said. He said the new agents were being drawn from within
Delta's ranks but that they would be selected through a ``highly competitive'' process and would be
paid more than other airport agents. The new Red Coats will be able to use hand-held devices to
print boarding passes and issue vouchers to airport clubs.

Other airlines have recently taken steps to improve their relations with customers. The chief
executive of AMR's American Airlines, Gerard J.~Arpey, said recently that the company was giving out
bonuses of \$100 to \$200 to customer service workers who met customer satisfaction goals in major
categories, ranging from baggage handling to ``gate interactions and boarding experience.''

At the same time, however, airlines are finding more ways to reduce contact with customers at
airports. Nearly half of all domestic travelers check in online or via a self-service kiosk at
airports, according to surveys by Airport Interviewing and Research Inc., an industry consultant in
White Plains. Checking in via mobile phone is another variation that is catching on in the United
States.

Some foreign airlines have gone further. Ryanair is simply eliminating airport counter agents, and
by next year will require all customers to check in online. Lufthansa now has ``self-boarding''
turnstiles at some of its gates at Frankfurt and Munich airports. Lufthansa also plans to test a new
type of kiosk that will print out meal vouchers and other documents for travelers delayed at major
connecting hubs. But Lufthansa also took note when customers complained about difficulties in using
an earlier version of its self-service kiosks, and the airline went back to the drawing board. ``You
have to be careful to pay attention to what is valued by your customer,'' said Jan-Wilhelm
Breithaupt, director of global passenger services for Lufthansa.

However welcome these advances are from an efficiency standpoint, consumer advocates point out they
can be a source of frustration when flights are canceled. ``Kiosks cannot solve complex travel
issues or emergency situations,'' said Kate Hanni, director of FlyersRights.org, which is pushing
for passage of a federal airline passengers' bill of rights that has recently gained momentum in
Congress.

Airports, aware their reputation is at stake as well, are also taking steps to improve their image
with consumers. ``The airport more than ever is becoming a caregiver,'' said Ira Weinstein,
president of Airport Interviewing and Research. ``If something goes wrong, the airport gets a black
eye.'' Several larger airports have initiated new programs in the last few years. Houston, for
example, started something called ``Houston-friendly'' that rewards airport agents for good service.

And a ``customer care'' program run by the Port Authority of New York and New Jersey, the operator
of New York's airports, has recently stepped up training to help its workers deal with increasingly
frustrated fliers. Among other things, these aides are instructed in the ``seven elements of
customer care,'' including ``smile, be friendly and approachable'' and ``initiate positive
contact.''

But there are limits to what airport employees can do. They do not have access to airline computers
and are not empowered to handle requests like overnight accommodations for stranded passengers or an
upgrade to first class on an overbooked flight.

So Delta's decision to revive the Red Coat has been the topic of some lively discussion in
business-traveler message boards on Web sites like flyertalk.com, where some recent postings focused
on what one person described as the ``big question -- will they be able to make 'battlefield
decisions' employing logic'' and cut through the red tape?

Mr.~West of Delta said that they would. He gave as an example the case of a business traveler flying
from New York to Latin America who missed his connection in Atlanta.

``The Red Coat would meet the inbound flight at the gate, they'd rebook the person on the next
available flight and scan the person's boarding pass right there and issue a new one,'' and if no
flight was available until the next day, the agent could dispense\cite{dispense}
vouchers\cite{voucher} or offer other assistance, all without having to go to a supervisor for
permission -- a common source of annoyance to frequent fliers who often have to wait for a
resolution.

``They are able to make those decisions on a case-by-case basis,'' he said.

\section{Silverman Quits NBC Universal; Cable Executive Succeeds Him}

\lettrine{T}{he} deal maker is gone. Now it's the cable guy's turn.

Ben Silverman, who called himself a ``rock-star television executive'' and tried to turn around
NBC's struggling prime-time schedule with reality shows and advertising-driven programs like ``The
Celebrity Apprentice'' and ``Knight Rider,'' announced on Monday that he would leave as co-chairman
of NBC Universal Entertainment to start a new venture with the media and Internet mogul\cite{mogul}
Barry Diller.

At the same time, NBC said that it would promote Jeff Gaspin, a longtime executive who oversees
NBC's lineup of cable channels, like USA, Bravo, CNBC and MSNBC; its syndication\cite{syndication}
business; and the Spanish-language network Telemundo, to chairman of NBC Universal Television
Entertainment. Mr.~Gaspin will run the entertainment operation of the NBC network, including its
prime-time schedule.

By awarding Mr.~Gaspin supervision of the flagship broadcast channel as well as the cable channels,
NBC Universal is acknowledging a shift that viewers and advertisers already recognize: that the
distinction between broadcast and cable is evaporating\cite{evaporate}.

``The distinction is definitely waning,'' said David Joyce, an analyst with Miller Tabak \& Company.
``It makes NBC seem more like a cable network company.''

NBC, with its long tradition as the ``Must See TV'' network, may now be considered just another part
of Mr.~Gaspin's large portfolio. The scripted shows on USA now rival those on NBC, especially in the
summertime months, when NBC's only hit is the reality show ``America's Got Talent.'' Last Thursday,
when NBC struggled to draw three million viewers to its sitcom repeats, a new episode of USA's most
popular show, ``Burn Notice,'' attracted 6.8 million viewers.

``USA has been putting on better original dramas than most broadcast networks,'' said Steve
Sternberg, a prominent TV analyst and a former executive vice president at the media agency MAGNA
Global.

Jeff Zucker, the president of NBC Universal, has said that NBC is now primarily a cable company. But
he tried to dismiss the notion that Mr.~Gaspin's arrival was a takeover by the cable side. ``Jeff
has a lot of experience on the broadcast side,'' he said. Many of the primary programming decisions
will be made by Angela Bromstad, who is president of prime-time programming.

In an interview, Mr.~Gaspin acknowledged the change in television fortunes. ``You can't dismiss the
fact that cable is just much more profitable than broadcast television. I don't see that changing.
Cable is a much superior business model.'' Without tipping his hand about particular shows, he added
that he would need to re-evaluate NBC's schedule if the network were to attract higher advertising
rates.

``Clearly we are going to have to re-message NBC's brand,'' he said. ``We have to make shows that
are successful enough to charge a premium to advertisers. It's the only way to do well in network
television.''

The transition comes at a trying time for both NBC Universal and its parent company, General
Electric, which has been hit hard by the recession and is paring\cite{paring} back selectively in a
strategy to rely more on its strong industrial businesses.

NBC has often been mentioned by analysts as a business outside the company's core of expertise that
could be a candidate for sale.

The unit's financial performance has lagged\cite{lag} recently. In the second quarter, NBC
Universal's revenue fell 8 percent, to \$3.57 billion, and profits fell 41 percent, to \$539
million, weighed down by weakness in advertising and a disappointing performance by the movie unit.

``Silverman was brought in to add firepower to the creative side of NBC Universal, and clearly
they've been disappointed,'' said Steven Winoker, an analyst at Bernstein Research. ``But I do not
think this is a signal that G.~E. is reducing its commitment to NBC.''

Mr.~Silverman ran Reveille, an independent production company behind hits like ``The Office,'' ``The
Biggest Loser'' and ``Ugly Betty'' when he was hired by NBC two years ago, but with the network
stuck in fourth place, his departure was not a surprise.

Even two months ago, Mr.~Silverman insisted in an interview that he was ``a happy worker at NBC''
but the network didn't renew his contract, clearly signaling that he had no future there without
formally firing him. Mr.~Silverman said that two weeks ago he began serious talks about a new role
with Mr.~Diller, who had been his partner in two previous ventures.

From the start, Mr.~Silverman was an awkward fit for NBC, which went after him aggressively to help
turn around their prime-time lineup. Under an unusual arrangement, Mr.~Silverman was allowed to keep
his stake in Reveille, although not a hand in its day-to-day operations, even as the company was
selling programs to NBC, an indication of just how prized Mr.~Silverman had been. (The company was
later sold to Shine, owned by Elisabeth Murdoch.)

But the high hopes for his tenure were quickly tempered by programming misses and personal missteps.
He saved the reality franchise ``The Apprentice,'' expanded ``The Biggest Loser'' and supported edgy
comedy like ``Parks and Recreation,'' but many of his choices, like ``American Gladiators'' and
``Knight Rider'' were failures that seemed out of place on the network known for quality dramas like
``ER'' and ``The West Wing.''

Mr.~Silverman also alienated\cite{alienate} people in Hollywood by insulting other executives in an
interview and became probably the first G.E. executive to throw a party featuring women in bikinis
and white tigers in cages. (Even Mr.~Silverman's departure was out of step with the G.E. corporate
culture. Twelve minutes before NBC sent out the press release, the news appeared on the Twitter
account of Mr.~Silverman's close friend, the ``American Idol'' host Ryan Seacrest.)

As his shows struggled, Mr.~Silverman's duties were curtailed\cite{curtail} and he spent much of his
time making deals with advertisers for sponsorships and product placements -- the part of the
business he promoted as the future of the industry. ``Ben did terrific things for us'' in the areas
of marketing and advertiser integration, Mr.~Zucker said, adding that Mr.~Silverman, ``really didn't
have the time'' to turn around NBC's ratings.

Mr.~Silverman said on Monday that he would remain at NBC through September to support the transition
to Mr.~Gaspin, and to back the start of the new fall TV series for the NBC network.

``I am invested in what happens in the fall and they are invested in my helping with the
transition,'' he said.

A senior executive at a large production studio who requested anonymity because his company has
deals with NBC said that programming, not advertising, still ruled in the television business.

``When you're the head of the network, you live and die by the shows you develop and schedule,'' he
said. ``You can talk all you want about disintermediation and bringing in advertisers, but at the
end of the day the measure of success is whether you create shows that people want to watch.''

\section{A Debate Rages in Iceland: Independence vs. I.M.F. Cash}

\lettrine{J}{ust} months after an epic\cite{epic} banking collapse forced Iceland into the arms of
the International Monetary Fund, this island nation is locked in a fierce debate over how to pay off
its creditors without ceding\cite{cede} too much of its vaunted independence.

The balance Iceland strikes between bowing to the policy demands of the global financial community
and satisfying the desires of its increasingly resentful population of 300,000 will be closely
watched as I.M.F. programs in beaten-down economies from Latvia and Ukraine to Hungary and Romania
enter a crucial phase.

''When you impose austerity\cite{austerity}, it becomes very painful and comes at a cost,'' said
Simon Johnson, a former I.M.F. economist who now teaches at the Massachusetts Institute of
Technology. But many Icelanders are blaming the I.M.F. and in this case, he says, that is not
warranted.

``Iceland is a rich country that behaved recklessly\cite{recklessly} and helped destabilize the
world financial system,'' Mr.~Johnson said. ``They will have to take their medicine.''

While those in Iceland's left-leaning government will not put it so bluntly\cite{bluntly}, that is
broadly the case they are making.

The first country to throw its government out of office as a result of the global financial crisis,
Icelanders could see the government that replaced it topple\cite{topple} too, leaving the country
rudderless -- unless it wins approval for a deal to repay Britain and the Netherlands the \$5.7
billion loan it used to compensate foreign depositors for losses in Icelandic banks.

A vote on the measure in the country's Parliament is scheduled for next week. But even Iceland's own
government is riven.

``This is an attack on our sovereignty\cite{sovereignty},'' said Ogmundur Jonasson, the country's
health minister. ``It reminds me of old colonial times. Gordon Brown had no harsh words for the
United States when Lehman Brothers went down and billions of pounds went to the U.S.~That was
friendship -- this is 'Take the little guy and nail him to the wall.' '' To not pass the bill, the
government says (most of it anyway), would lead to the I.M.F. and other outside lenders withdrawing
funds, further jeopardizing\cite{jeopardize} the country's fragile condition.

But detractors\cite{detractor} say passing it would increase Iceland's debt burden to 200 percent of
gross domestic product, making it one of the most leveraged nations in the world. Ultimately, they
say, it could drive Iceland to default.

At the crux\cite{crux} of this debate is the Icesave, or ``Iceslave,'' as it is called here. Icesave
accounts were a top-of-the-market gambit\cite{gambit} by Landsbanki, the most aggressive of the
failed Icelandic banks, to raise cash by extending its branch network from tiny Reykjavik to the
high streets of London. The reaction to the agreement to make good on the accounts encapsulates all
the swelling anger that Icelanders now bear toward bankers, foreign creditors and I.M.F. technocrats
-- not necessarily in that order.

Lilja Mosesdottir is an economist and a back-bench member of Parliament in the governing Left Green
party. But if she were to vote now, she says, she would vote against the government bill.
Ms.~Mosesdottir, new to politics, swept into power this winter when the conservative party was
overturned by the ``pots and pans revolution.''

``It is like after a war and you are the loser,'' she said, taking a quick coffee break from
back-room negotiations over the deal. ``This is an agreement that will lead to a sovereign default,
and we don't want that to happen.''

Whether or not she is right about default, the war analogy is apt. Iceland has lost billions, and
others are now dictating the terms of its recovery.

The resentment\cite{resentment} felt here is rooted in a belief that Iceland's core virtue of
flinty\cite{flinty} self-reliance has been defiled by its bankers and foreign creditors. It is a
sentiment\cite{sentiment} that stretches far into the country's history and culture -- from the
Nordic sagas to the quest for autonomy of Bjartur of Summerhouses, the impoverished sheep farmer in
Halldor Laxness's ``Independent People,'' the country's best-known modern literary work.

As the rhetoric\cite{rhetoric} escalates\cite{escalate}, Iceland's finance minister, Steingrimur
J.~Sigfusson, a lifelong leftist, finds himself in the awkward position of defending the Icesave
plan as well as the severe economic restrictions that the country has been forced to endure to
qualify for more money from the I.M.F. and other Nordic lenders. Such measures include sharp cuts in
health spending and higher gas prices. Higher interest rates have pushed unemployment to about 8
percent, from 1 percent, in little more than a year.

Mr.~Sigfusson scoffs\cite{scoff} at any notion of default and argues that the deal to repay
creditors was the best that could have been achieved. With a term of 15 years, a low interest rate
and a seven-year grace period, the deal is flexible enough to allow Iceland to repay it, he says,
especially if the economy recovers and the government is successful in selling Landsbanki's foreign
assets.

``This is the greatest tragedy of all, but it has to be done,'' he said, looking gaunt from the
hours of parliamentary arm-twisting that now consume his days.

As to the widely held belief that it is the I.M.F. and not the government that is dictating policy,
Mr.~Sigfusson acknowledges that he is in close contact with the I.M.F.'s representative here.

He points to frequent disagreements, especially over the fund's recommendation that the government
maintain high interest rates as well as capital controls -- a prescription he describes as similar
to wearing a belt and suspenders at the same time. But he emphasizes that it is Iceland, not the
I.M.F., that has the final word.

``This is a trial not just for us, but the I.M.F., too,'' he said. ``They have a lot at stake here
as they must show that they are flexible enough to adapt their program to a developed Nordic welfare
state.''

Known to many here as ``the governor of Iceland,'' Franek Rozwadowski, the I.M.F. representative,
argues that this designation is inaccurate. As part of its program, Iceland must turn a deficit that
is now 13 percent of G.D.P. to a surplus by 2013.

``It would be more accurate to call the relationship a collaboration in which Iceland has engaged
the fund to help design its recovery program,'' he said.

On Aug.~3, the I.M.F. is expected to discuss whether to disburse a second tranche of its \$2.1
billion loan to Iceland (about a quarter has been disbursed so far). Mr.~Rozwadowski says Iceland is
on target with steps to balance its budget, and he hails Mr.~Sigfusson for political courage.

Such niceties are thin gruel for many Icelanders whose personal debts have skyrocketed in the wake
of the precipitous fall in Iceland's currency, the krona.

Gunnar Sigurdsson, a theater director, says his car loan -- which was tied to a basket of Swiss
francs and Japanese yen -- has doubled since the crisis began; his mortgage payments have jumped
over 35 percent.

Personal bankruptcy is inevitable, he says, and he is now trying to make a ``Roger and Me''-type
documentary -- training his camera on Iceland's top politicians, bankers and, if he is lucky,
Dominique Strauss-Kahn, the head of the I.M.F. ``I have had enough of this stupidity,'' he said. ``I
just want answers.''

\section{Texas Hospital Flexing\cite{flex} Muscle in Health Fight}

\lettrine{O}{ne} of the largest sources of campaign contributions to Senate Democrats during this
year's health care debate is a physician-owned hospital in one of the country's poorest regions that
has sought to soften measures that could choke\cite{choke} its rapid growth.

The Democratic Senatorial Campaign Committee collected nearly \$500,000 at a reception here on March
30, mostly from physicians and others affiliated\cite{affiliate} with Doctors Hospital at
Renaissance, financial disclosure records show.

The event was held at the home of a prominent McAllen developer, Alonzo Cantu, a hospital founder,
investor and board member who has raised prodigious\cite{prodigious} sums from the Rio Grande Valley
for an array of Democrats.

Another event at Mr.~Cantu's home, in September 2007, brought in at least \$800,000 for the
committee's House counterpart, the Democratic Congressional Campaign Committee, according to
disclosure reports. The House speaker, Nancy Pelosi, was in attendance and cut a ribbon at the
hospital's new women's center while in town.

The hospital, which is in Edinburg, adjacent to McAllen, is working both sides of the aisle. Its
political action committee, Border Health PAC, split \$120,000 last year among House and Senate
candidates, including Republicans.

Although Congressional negotiations over health care legislation are continuing, Doctors Hospital
seems to be getting much of what it wants. Thus far, physician-owned hospitals have been insulated
from some of the most onerous\cite{onerous} potential restrictions in the health care legislation
moving through Congress.

Representative Pete Stark, a California Democrat who wants to clamp\cite{clamp} down on
physician-owned hospitals, said their formidable\cite{formidable} lobbying had helped eliminate his
proposal to limit physician ownership to 40 percent at any hospital.

``Particularly led by these guys in Texas, these guys who have been raising tons of money for
contributions,'' Mr.~Stark said in an interview. ``I am sure that some of my colleagues have been
willing to hear them out.''

The flood of dollars from just a single interest of modest size illustrates just how high the stakes
are in this year's health care negotiations, and how even sideline issues generate intense lobbying.

The gleaming\cite{gleam}, well-equipped Doctors Hospital at Renaissance\cite{renaissance}, which has
expanded to 503 beds from 30 in six years, has become a footnote in the health care debate. It was
featured unflatteringly in a June article in The New Yorker about geographic
disparities\cite{disparity} in health care spending, a story that President Obama has cited
repeatedly in speeches and meetings.

The article, which is sharply disputed\cite{dispute} by hospital officials, posited\cite{posit} that
physician ownership provided ``an unholy temptation to overorder'' tests and procedures because
doctors earn not only their fees but also a share of the hospital's profits. At Doctors Hospital,
where 353 of its 452 owners are physicians, net revenue amounted to \$64 million in 2008.

Leading members of Congress have long been concerned about the potential for conflicts of interest,
lapses\cite{lapse} in patient safety and cherry-picking of patients with the best insurance at the
country's 230 physician-owned hospitals. Past efforts to restrict ownership by doctors have
stalled\cite{stall}, but language to that effect is currently included in health care legislation in
Congress, though in ways that are unlikely to hurt Doctors Hospital.

The campaign committees, because they are controlled by each party's Congressional leadership, are
viewed as effective vehicles for gaining influence at high levels. But top Democrats, including
Senator Robert Menendez of New Jersey, chairman of the Senate committee, said that Mr.~Cantu had
been a longtime patron\cite{patron} and that the Texas contributions would have no special impact.

Mr.~Cantu, a son of Mexican immigrants, said the primary goal of the fund-raising was to
lure\cite{lure} politicians to the border so they could see firsthand the region's needs for jobs,
education, infrastructure and health care.

``We've always been ignored by the Republicans because we're the blue section of the state, and the
Democrats have taken us for granted,'' he said. ``What do politicians listen to? Typically, it's
money and votes.''

But he and other Doctors Hospital officials acknowledged they also have a particular interest in
protecting the hospital, one of the largest employers in the region. On at least three occasions in
the last two months, delegations from McAllen have made the rounds of Capitol Hill to discuss their
concerns, and blunt\cite{blunt} the effects of the New Yorker article.

They have argued they are being unfairly grouped with boutique specialty hospitals that do not have
emergency departments and that cater to privately insured patients. Eighty-eight percent of Doctors
Hospital patients are either on public insurance or uninsured, 750 babies are delivered there a
month, and no one is turned away because of inability to pay, they said.

Physician ownership, they added, has meant major investments in the latest equipment and good
staffing ratios for nurses. Appealing to local pride, the hospital markets itself as the first in
the area to offer services like PET scans, robotic surgery and breast imaging, which once required
trips to Houston or San Antonio.

Representative Rub\'en Hinojosa, a Democrat whose district includes the hospital, said he set up
some of the Capitol Hill meetings. Mr.~Hinojosa said he never mentioned the campaign contributions
and could not judge whether the money had opened doors.

But, he said, the hospital's officials ``are listened to,'' adding, ``Every chairman of a committee
that has influence over them had meetings with them.''

The Senate Finance Committee has yet to release its final draft, but bills passed by two House
committees would prevent the opening of new physician-owned hospitals by disqualifying them from
receiving Medicare reimbursements. Existing facilities like Doctors Hospital would be grandfathered
in.

One key provision would limit a hospital's ownership by doctors to the level in place at the time of
enactment. That is a change from previous language in House bills to restrict physician ownership to
40 percent. It would have forced Doctors Hospital, where physicians have an 82-percent stake, to be
sold or required some of its owners to divest.

Mr.~Stark said the change had been made ``at a higher pay grade than mine.'' The physician-owned
hospitals, he said, ``just thought they could buy their way out of it, and it's a sad commentary on
the Congress.''

Speaker Pelosi supported the tougher provisions in previous bills but aides said she was willing to
accept a compromise to win the new restrictions.

Aides to Senator Max Baucus. Democrat of Montana, the chairman of the Finance Committee, and to
Senator Charles E.~Grassley of Iowa, the ranking Republican, said the Senate would never have
adopted the 40 percent limit and that it was only fair to allow existing hospitals to maintain their
ownership structures.

Another provision in the House bill would allow physician-owned hospitals to expand to up to double
their size, with government approval, but only in fast-growing areas with heavy Medicaid populations
and high bed occupancy rates. Doctors Hospital officials said they thought they would qualify.

The officials said they could not judge what influence they might have had in shaping the
legislation. ``I can't say we had a direct impact,'' said Susan S.~Turley, the chief financial
officer. ``I think every physician-owned hospital was talking to their congressmen.''

Whether or not there is cause and effect, some of the hospital's beneficiaries have passionately
supported its position. In November, the South Texas doctors raised nearly \$60,000 for
Representative Sheila Jackson Lee, a Democrat from Houston. Last week, Ms.~Lee, who has a
physician-owned hospital in her district, inserted comments in The Congressional Record against
restricting the hospitals' growth. She did not respond to requests for comment for this article.

Disclosure reports show that about 220 donors from the McAllen area contributed to the Senate
Democratic committee for the March event. Receipts totaled at least \$425,000, about 2 percent of
the group's collections for the first half of the year.

``We're not na\"ive,'' said Dr.~Ambrosio Hernandez, a pediatric surgeon who gave \$5,000. ``We
understand that politics plays a role in everything.''

There is a wide gulf between rich and poor in McAllen, where gated neighborhoods built by Mr.~Cantu
coexist with former shantytowns that only recently gained paved roads and indoor plumbing.

The largest gift to the Senate Democrats -- the maximum allowable \$30,400 -- came from Mr.~Cantu,
whose construction company built the hospital. Other major donations included \$26,783 from
Dr.~Lawrence Gelman, the hospital's chief executive, and gifts of \$10,000 to \$15,000 from three
hospital board members, six doctors and two top executives at a bank where Mr.~Cantu is chairman and
that has made loans to hospital investors.

Dr.~Gelman also attended a fund-raiser related to health care on Monday in Washington for the
National Republican Senatorial Committee, but because he has already contributed the maximum of
\$30,400, he did not have to pay the \$5,000 ticket price.

Mr.~Cantu, who began his working life as a migrant grape-picker, has long been generous to
Democratic and philanthropic\cite{philanthropic} causes. Disclosure filings show he and his family
members have donated more than a quarter of a million dollars to federal campaigns since 1995. He
said he raises money from doctors, vendors, employees and friends, but has never reimbursed anyone
for a donation, which would violate federal election law.

``Just for the record,'' he said, ``we're not that stupid.''

Mr.~Cantu helped bundle hundreds of thousands of dollars in contributions for Hillary Rodham
Clinton's 2008 presidential campaign. He was the only private citizen invited to accompany
Ms.~Clinton, now secretary of state, to El Salvador in June for the inauguration\cite{inauguration}
of its new president.

He gave \$2,300 to Mr.~Obama last October, after Ms.~Clinton had withdrawn from the race. But like
others here, he is not pleased about the president's depiction of health care in McAllen.

``What's so upsetting,'' he said, ``is that to make his case he threw McAllen under the bus.''

\section{House Democrats End Impasse on Health Bill}

\lettrine{E}{fforts} to pass sweeping health care legislation took a big step forward on Wednesday
as House Democratic leaders reached an agreement with fiscally conservative party members that would
cut the bill's cost and exempt many small businesses from having to provide health benefits to
workers.

The agreement, brokered by aides to President Obama, overcame a 10-day impasse and would allow a
pivotal\cite{pivotal} House committee to resume work on the bill, with an expectation that the panel
could approve it later this week.

Under the deal, the Democratic leaders promised to defer a vote by the full House until September,
so lawmakers could test public sentiment on the measure, which could fundamentally restructure
one-sixth of the nation's economy.

Elements of the agreement reflect priorities shared by centrist members of both parties who have
been trying for months to forge a compromise in the Senate Finance Committee.

Under the House agreement, the federal government would still establish a public insurance plan to
compete with private insurers, but the public plan would not use Medicare fee schedules to pay
doctors and hospitals, as envisioned in the original House bill. Instead, the public plan would
negotiate rates with health care providers, as private insurers often do.

Representative Mike Ross of Arkansas, the chief negotiator for the fiscally conservative Democrats
known as Blue Dogs, said the changes were ``a huge win for us.''

Blue Dogs said the change reduced the likelihood that the government plan would compete unfairly
with private insurers by forcing doctors and hospitals to accept below-market rates. As part of the
deal, states could, in addition, set up nonprofit cooperatives to offer coverage to individuals,
families and small businesses. The Senate Finance Committee is coalescing\cite{coalesce} around a
similar idea, as an alternative to a government-run insurance plan, but it was not in the original
House bill.

``We have cut the cost of the bill substantially,'' Mr.~Ross said. ``We have delayed a floor vote
until September. We have protected small business.''

But some of the concessions to Blue Dogs set off a revolt\cite{revolt} among members of the
Congressional Progressive Caucus\cite{caucus}, who said they feared that the public insurance plan
was being weakened.

``We do not support this,'' said Representative Lynn Woolsey, Democrat of California, co-chairwoman
of the progressive caucus. ``It's a nonstarter.''

In the Senate, negotiations in the Finance Committee were also moving forward -- to such a degree
that Senate Republican leaders became worried that a deal might be near. One Republican negotiator,
Senator Michael B.~Enzi of Wyoming, tried to quash that idea, saying the group was ``nowhere near a
deal.''

``I don't see a way that we can finish before the recess,'' Mr.~Enzi said -- just 24 hours after the
Senate Democratic leader, Harry Reid of Nevada, had said he expected the committee to approve a bill
next week.

The Energy and Commerce Committee, the third of three House panels to take up the legislation, was
to resume work on Thursday.

As he traveled to Raleigh, N.C., and Bristol, Va., Mr.~Obama welcomed the agreement. ``I'm
especially grateful that so many members, including some Blue Dogs on the Energy and Commerce
Committee, are working so hard to find common ground,'' he said. ``Those efforts are extraordinarily
constructive in strengthening this legislation and bringing down its cost.''

The president's travels on Wednesday were part of a White House sales pitch, as officials have been
fretting\cite{fret} that the script has been hijacked by critics who say that some of his proposals
could hurt Americans who already have health insurance.

``These folks need to stop scaring everybody,'' Mr.~Obama said.

After marathon negotiations, four of seven Blue Dog holdouts on the Energy and Commerce Committee
reached the agreement with the chairman of the panel, Representative Henry A.~Waxman, Democrat of
California.

The Congressional Budget Office said the bill would ultimately provide coverage to most people who
are uninsured. Mr.~Ross said the Blue Dogs had won concessions that should bring the 10-year cost
below \$1 trillion, a goal shared by the Senate Finance Committee.

Senator Max Baucus, the Montana Democrat who is chairman of the committee, said the latest estimates
from the budget office indicated that his bill would cost ``under \$900 billion.''

The House negotiators, including Mr.~Ross and aides to Mr.~Waxman, said they had agreed to make
these changes in the bill:

Most employers would still be required to provide health insurance to workers or pay a new federal
tax, but more small businesses could qualify for the exemption, which would be available to
businesses with annual payrolls of \$500,000 or less, compared with a threshold of \$250,000 in the
original House bill. The maximum tax rate, 8 percent of wages, would apply to employers with
payrolls exceeding \$750,000, rather than the original threshold of \$400,000.

Medicaid would be expanded, as under the original bill, but states would pay a small share of the
additional costs, perhaps 7 percent. The federal government would have paid all the additional cost
under the original bill.

People with low or moderate incomes could still get federal subsidies to help them buy insurance,
but they might have to spend slightly more of their own income -- a maximum of 12 percent, rather
than 11 percent.

Representative Earl Pomeroy, a Blue Dog from North Dakota, voted against the bill in the Ways and
Means Committee on July 17 but said Wednesday that he could support the revised bill.

He said he was pleased that the bill would exempt more of the smallest employers, and that the
public plan would not use Medicare payment rates, which are low in North Dakota.

Mr.~Ross gave few details about the new provisions to control costs, other than to say that the
government would spend less on subsidies and Medicare would save money by linking hospital payments
to the quality of care provided, under a new system of ``value-based purchasing.''

\section{New Poll Finds Growing Unease on Health Plan}

\lettrine{P}{resident} Obama's ability to shape the debate on health care appears to be
eroding\cite{erode} as opponents aggressively portray his overhaul plan as a government takeover
that could limit Americans' ability to choose their doctors and course of treatment, according to
the latest New York Times/CBS News poll.

Americans are concerned that revamping\cite{revamp} the health care system would reduce the quality
of their care, increase their out-of-pocket health costs and tax bills, and limit their options in
choosing doctors, treatments and tests, the poll found. The percentage who describe health care
costs as a serious threat to the American economy -- a central argument made by Mr.~Obama -- has
dropped over the past month.

Mr.~Obama continues to benefit from strong support for the basic goal of revamping the health care
system, and he is seen as far more likely than Congressional Republicans to have the best ideas to
accomplish that. But reflecting a problem that has hindered\cite{hinder} efforts to bring major
changes to health care for decades, Americans expressed considerable unease about what the end
result would mean for them individually.

``We need to fix health care,'' Mary Bevering, a Democrat from Fort Madison, Iowa, said in a
follow-up interview, ``but if the government creates the system, I'm afraid the quality of care will
go down and costs will go up: We will pay more taxes.''

``It's going to come down to regulation,'' Ms.~Bevering said. ``What also worries me is whether we
will be told what physician we can have.''

The poll was taken at a moment of extreme fluidity\cite{fluidity}, both in terms of the complicated
negotiations in the House and the Senate as lawmakers and the administration sort out the substance
and politics of competing proposals, and in the efforts by both sides to define the stakes of the
health care debate for the public.

With Congress now almost certain to recess until after Labor Day without floor votes on any specific
plan, a vigorous\cite{vigorous} advertising and grass-roots effort to shift public opinion is likely
in the next month or two. The poll offers hope to both sides.

The changes in the public's attitude over the past month, even if not huge, suggest one reason
Mr.~Obama sought so hard to get Congress to vote on some version of an overhaul before heading home.

Opponents of the proposed health care overhaul have already spent \$9 million on television
advertisements raising concerns about it, said Evan Tracey, the chief operating officer of Campaign
Media Analysis Group, which tracks political advertising. The advertisements are financed by the
Republican National Committee and aimed at constituents of wavering lawmakers. The committee is also
running radio spots.

Officials said the advertising would accelerate as the legislators returned home for the summer. The
advertisements present the overhaul as a risky experiment, or a government takeover of health care
that would prevent people from choosing their own doctors.

Mr.~Obama is making an intense effort to rebut those claims. On Wednesday, he flew to Raleigh, N.C.,
for a town-hall-style meeting to address the kinds of public concerns reflected in the poll results.

``First of all,'' Mr.~Obama said, ``nobody is talking about some government takeover of health care.
I'm tired of hearing that. I have been as clear as I can be. Under the reform I've proposed, if you
like your doctor, you keep your doctor; if you like your health care plan, you keep your health care
plan. These folks need to stop scaring everybody, you know?''

Mr.~Obama sought in particular to reassure people who already have health insurance and whom the
overhaul plans under consideration in Congress would benefit by preventing insurers from dropping
them or diluting\cite{dilute} their coverage if they become ill, while also bringing rapidly rising
costs under control. And he sought to stoke a sense of urgency for getting a bill signed this year.

``If we do nothing, I can almost guarantee you your premiums will double over the next 10 years,
because that's what they did over the last 10 years,'' Mr.~Obama said. ``It will eat into the
possibility of you getting a raise on your job because your employer is going to be looking and
saying, 'I can't afford to give you a raise because my health care costs just went up 10, 20, 30
percent.' ''

The national poll was conducted by telephone starting on Friday and ending on Tuesday. It involved
1,050 adults, and has a margin of sampling error of plus or minus three percentage points.

Mr.~Obama's job approval rating has dropped 10 points, to 58 percent, from a high point in April.

And despite his efforts -- in speeches, news conferences, town-hall-style meetings and other forums
-- to address public misgivings, 69 percent of respondents in the poll said they were concerned that
the quality of their own care would decline if the government created a program that covers
everyone.

Still, Mr.~Obama remains the dominant figure in the debate, both because he continues to enjoy
relatively high levels of public support even after seeing his approval ratings dip, and because
there appears to be a strong desire to get something done: 49 percent said they supported
fundamental changes, and 33 percent said the health care system needed to be completely rebuilt.

The poll found 66 percent of respondents were concerned that they might eventually lose their
insurance if the government did not create a new health care system, and 80 percent said they were
concerned that the percentage of Americans without health care would continue to rise if Congress
did not act.

By 55 percent to 26 percent, respondents said Mr.~Obama had better ideas about how to change health
care than Republicans in Congress did.

There is overwhelming support for a bipartisan\cite{bipartisan} agreement on health care, and here
again, Mr.~Obama appears in the stronger position: 59 percent said that he was making an effort to
work with Congressional Republicans, while just 33 percent said Republicans were trying to work with
him on the issue.

Over all, the poll portrays a nation torn\cite{torn} by conflicting impulses and confusion.

In one finding, 75 percent of respondents said they were concerned that the cost of their own health
care would eventually go up if the government did not create a system of providing health care for
all Americans. But in another finding, 77 percent said they were concerned that the cost of health
care would go up if the government did create such a system.

\section{Airlines Follow Passengers Onto Social Media Sites}

\lettrine{W}{hen} United Airlines baggage\cite{baggage} handlers in Chicago damaged Dave Carroll's
\$3,500 guitar last year, he tried to get his restitution\cite{restitution} from the airline the
old-fashioned way. But after months and months of phone calls and faxes with various
customer-service representatives, the airline refused to accept his claim. So Mr.~Carroll, a
professional country music singer from Canada, channeled his frustration into a song and a video,
which he posted on YouTube.

Since it first appeared July 6, ``United Breaks Guitars'' has been viewed more than 4.4 million
times. United has scrambled\cite{scramble} to respond, also in nontraditional ways. Its first
comments came not in a press release, but via Twitter.

As the global recession and the consequent drop in air traffic force carriers to cut back on
services and consider charging extra for everything from checked bags to onboard meals, the
Internet, and social media Web sites in particular, is giving once-faceless travelers a global --
and instantaneous -- platform from which to air their grievances.

Confronted with these middle-seat Davids-turned-Goliaths, airlines, like other consumer-oriented
businesses, are racing to find their own ways to use social media channels like Twitter, Facebook
and YouTube to forge deeper relationships with passengers -- before things go wrong.

``Thanks to social computing, travelers' tales are no longer confined\cite{confine} to what they
tell to their coworkers and neighbors,'' said Henry H.~Harteveldt, a vice president and airline and
travel industry analyst at Forrester Research in San Francisco, adding that roughly 60 percent of
travelers in Europe and North America engage in some form of social networking online. ``They are
out there in public for the whole world to see.''

Airlines in the United States have been the quickest to embrace social media as a low-cost public
relations and marketing tool, in particular to spread the word about fare sales or to make
announcements about new routes or services. Carriers like Southwest Airlines, JetBlue and Alaska
Airlines are among the most active users, each with online ``followings'' in the tens, if not
hundreds of thousands of people. Even some airports, like Hartsfield-Jackson International in
Atlanta and Logan International in Boston, are using dedicated channels on Facebook and YouTube to
provide travelers with information like how to use the airport train system or to give updates on
construction projects or changes to rental-car facilities.

``With social media, the communication goes both ways,'' said Steven Lott, a spokesman for the
International Air Transport Association, based in Geneva. ``If there is any kind of operational
problem, someone can send pictures, a video on YouTube or tweet about it and the word gets out in
seconds. At the same time, airlines that use social media effectively can come up with a public
relations response just as quickly, using it to dispel rumors or myths and get out accurate
information.''

Southwest Airlines, which has more than 380,000 Twitter followers and another 70,000 on Facebook,
leveraged these online audiences after a football-sized hole appeared in the fuselage of one of its
planes during a flight from Nashville to Baltimore on July 13. Within an hour of the plane's
emergency landing, in Charleston, West Virginia, passengers from the flight were tweeting about it
and posting images of the hole, said Christi Day, a member of Southwest's emerging media staff.

``We closely monitored the conversations that were developing,'' said Ms.~Day. ``As soon as we felt
we were ready, we got out as much factual information as we could.'' That included an official
statement that Southwest's entire fleet would be inspected immediately and that passengers onboard
the affected flight would get refunds.

Last December, Oliver Beale, a passenger traveling on Virgin Atlantic from Mumbai to London, was so
put off by his in-flight meal that he wrote a blistering letter of complaint to the airline's owner,
Richard Branson, complete with photos of the food. Copies of the missive, in which Mr.~Beale
described his ``culinary journey from hell,'' spread rapidly via e-mail and the story was quickly
picked up by bloggers and traditional media from Sydney to Southampton.

The airline's public-relations team quickly got out the message on Facebook that Mr.~Branson had
personally phoned Mr.~Beale and invited him to help select food and wine for future flights.

``Our response was so fast, some people even accused us of setting the whole thing up as a publicity
stunt,'' said Paul Charles, a Virgin Atlantic spokesman. ``We took that as a compliment.''

Others have had less positive experiences with social media. In March, someone set up a Twitter
profile purporting to be an official Ryanair channel. Several tweets made abusive remarks about the
airline's customers, describing them as ``mostly stupid'' and ``a bit dense.''

Ryanair complained to Twitter and had the rogue\cite{rogue} feed shut down. But the airline also
discontinued several official feeds set up by its marketing department.

For low-cost carriers like Ryanair, analysts recognized there might be disincentives. ``For most
airlines that have a large following on Twitter or Facebook, it's pretty clear that they have a few
people dedicated full-time to it,'' said David Beckerman, a vice president at OAG , an aviation
market research firm. JetBlue, for example, which has more than 960,000 Twitter followers, has 35
staff members dedicated to updating its feed.

Still, he said, ``it's weird that airlines aren't using social media more,'' particularly in Asia,
where large numbers of travelers are equipped with Web-enabled cellphones. Only a handful of Asian
carriers, among them Malaysia Airlines and the low-cost Air Asia, use social media actively, he
said.

In the case of Mr.~Carroll and his guitar, United's apology -- which first went out to its 18,600
Twitter followers about 24 hours after the video appeared -- has failed to attract a fraction of the
attention than did the singer's original protest. The airline has not responded on its own corporate
YouTube channel, although viewers have posted more than a dozen comments related to the song. ``This
company needs someone online for damage control (no pun intended),'' wrote one.

Robin Urbanski, a United spokeswoman, said in an e-mailed statement that the airline should have
responded much sooner and that Mr.~Carroll's video would be used for corporate training purposes.

``United's experience is the latest and best example of how social computing is affecting the
business reputations of travel companies,'' said Mr.~Harteveldt of Forrester. ``Every mistake as
well as every success is going to be Facebooked and Twittered.''

\section{Senators Seek a Ban on Texting and Driving}

\lettrine{S}{tates} that do not ban texting by drivers could forfeit\cite{forfeit} hundreds of
millions of dollars in federal highway funds under legislation introduced Wednesday in the Senate.

Under the measure, states would have two years to outlaw the sending of text and e-mail messages by
drivers or lose 25 percent of their highway money each year until the money was depleted.

``Studies show this is far more dangerous than talking on a phone while driving or driving while
drunk, which is astounding,'' said Senator Charles E.~Schumer, Democrat of New York, one of four
Democratic senators to introduce the proposal.

Mr.~Schumer said the authors were responding to recent studies that have begun to quantify the risks
of texting while driving. One study released this week from the Virginia Tech Transportation
Institute found that truck drivers face a 23 times greater risk of crash or near crash when texting
than when not doing so.

Research from the University of Utah, which used a driving simulator to study the ability of
motorists to multitask, found an eightfold greater risk of crashing when texting. By comparison,
Utah researchers showed that drivers using a cellphone to talk face a four times greater risk of
crashing, about equal to someone with a 0.08 blood alcohol level, generally the legal limit for
intoxication.

Currently, texting while driving is banned in 14 states, including Alaska, California and New
Jersey, as well as the District of Columbia. The legislature in New York recently passed such a
measure and sent it to the governor for a signature.

Regulation of the roadways generally happens at the state level. But the federal government has
exerted pressure on the states based on the threat of withholding federal highway funds, as Congress
did in 1984 to pressure states to raise the minimum drinking age to 21 years.

Mr.~Schumer said that the legislation was essentially based on the drinking age law.

The Governors Highway Safety Association, a group that represents state highway safety agencies in
every state, opposes texting while driving but does not support the proposed legislation.

``We oppose sanctioning\cite{sanction} states since there is not yet a proven effective method for
enforcing a texting or cellphone ban,'' an association spokesman, Jonathan Adkins, said.

Safety advocates said that such concerns about enforcement were raised about seat belt laws but
argued that the value of such laws -- even if they could not be enforced all the time -- created
awareness about the issue and set societal guidelines for the behavior.

The other sponsors of the Senate measure include Mary L.~Landrieu of Louisiana, Robert Menendez of
New Jersey and Kay R.~Hagan of North Carolina.

\section{Microsoft and Yahoo Are Linked Up. Now What?}

\lettrine{T}{he} bumpy\cite{bumpy}, marathon mating dance between Microsoft and Yahoo finally
concluded on Wednesday, when the two companies announced a partnership in Internet search and
advertising to take on the industry powerhouse Google.

But there was plenty of skepticism about whether the new partners could make a serious
dent\cite{dent} in Google's dominance.

Even with the deal, the Microsoft-Yahoo search operation will be dwarfed by Google -- with a 28
percent market share in the United States, versus 65 percent -- and will face an uphill struggle to
try to wean people away from Google's simple white search page.

If Yahoo and Microsoft cannot persuade people to switch, they will not build the larger audience
that will bring in more revenue from ads tied to searches.

``This battle is won or lost as the user sits at the keyboard,'' said Peter S.~Fader, a professor at
the Wharton School of the University of Pennsylvania and co-director of its Interactive Media
Initiative. ``Google is winning for good, consumer-friendly reasons. You can't just buy that.''

The Microsoft-Yahoo pact represents a pragmatic\cite{pragmatic} division of duties between the two
companies, instead of the blockbuster deal Microsoft, No.~3 in the search market, was shooting for
last year when it bid \$47.5 billion to buy Yahoo, No.~2 in search.

That hostile\cite{hostile} offer was ultimately withdrawn by Microsoft, and its collapse led to lots
of soul-searching at Yahoo and the replacement of its co-founder Jerry Yang with an outsider, Carol
Bartz, in the chief executive role.

Under the pact, Microsoft will provide the underlying search technology on Yahoo's popular Web
sites. The deal will give a lift to Microsoft's search engine, which it recently overhauled and
renamed Bing. Its search ads will have broader reach and become more lucrative.

Bing, which tries to put search results in better context than rivals, has won praise and favorable
reviews, after Microsoft spent years falling farther and farther behind Google in search.

For Yahoo, the move furthers the strategy under Ms.~Bartz to focus the company on its strengths as a
publisher of Web media sites in areas like finance and sports, as a marketer and leader in online
display advertising.

``This deal allows Yahoo to invest in what we should be investing in for the future -- audience
properties, display advertising and the mobile Internet experience,'' Ms.~Bartz said in an interview
on Wednesday. ``Our vision is to be the center of people's lives online.''

The terms of the 10-year agreement give Microsoft access to Yahoo's search technologies. Yahoo will
receive a lucrative 88 percent of the search-generated ad revenue from its own sites for the first
five years of the deal, much higher than is standard in the industry.

After the takeover bid failed, the companies renewed talks about a partnership last summer. The
talks included discussion of a large upfront payment from Microsoft.

But when Ms.~Bartz joined Yahoo at the start of this year, the interest on the Yahoo side shifted.
Ms.~Bartz was more interested in steady revenue to ensure the longer-term financial health of Yahoo
instead of a big payment, she said in a conference call Wednesday.

Shares of Yahoo fell 12 percent, to \$15.14, after the deal was announced, apparently reflecting
investors' disappointment in the lack of a payment. Shares of Microsoft rose slightly.

``It feels kind of like a stab in the chest,'' said Darren Chervitz, the co-manager of the Jacob
Internet Fund, which owns about 100,000 shares of Yahoo. ``It certainly feels like Yahoo is giving
away their strong and hard-fought share of the search market for really a modest price.''

Now, Yahoo's financial fate will be inextricably linked with Microsoft for years. ``My sense is that
Yahoo will regret making this move,'' Mr.~Chervitz said.

If the deal is completed next year as planned, and after the partnership is fully in place in three
years, Yahoo estimates that its operating income will increase by \$500 million a year, based on the
anticipated higher search traffic and ad revenue, and a substantial drop in its investment in
technology development.

Steven A.~Ballmer, Microsoft's chief executive, said in an interview that Ms.~Bartz had driven a
hard bargain. ``Look,'' he said, ``she got 88 percent of the revenue and none of the cost.''

Still, Mr.~Ballmer added that he won something he badly wanted as well: ``I got an opportunity to
swing for the fences in search.''

The crucial issue for the partnership, analysts say, is its ability to stop and reverse the advance
of Google, which has steadily gained in search and search ad share in the last few years while Yahoo
and Microsoft have kept on fading. Reversing the trend, they say, would give the partnership
newfound credibility with advertisers and publishers.

``If Microsoft and Yahoo are 30 percent and growing in search, then the dynamics of the market can
shift,'' said David B.~Yoffie, a professor at the Harvard Business School.

Mr.~Fader of Wharton said he was not sure the partners would be able to shake up the business.
``Microsoft is buying some market share, but there is no evidence they are going to change the game
in any fundamental way,'' he said. ``What the Microsoft-Yahoo partnership needs is real
breakthroughs to deliver disruptive\cite{disruptive} innovation in search.''

Advertisers and Web publishers say they will be looking to the combination to improve its search
technology and service and put more pricing pressure on Google, which has turned the small text ads
that appear next to search results into a multibillion-dollar business.

``It could be a more competitive marketplace, but that's not certain,'' said Bob Liodice, president
of the Association of National Advertisers, a trade group. ``Google still holds two-thirds of the
market.''

Branding was one important consideration in the deal. Yahoo will still control the look of the
search features on its sites and will determine how search technology may be tailored differently
for, say, entertainment and finance sites. But Yahoo's search will include a logo saying ``Powered
by Bing.''

And Yahoo will be able to tap into records of searches for its own purposes, like monitoring the
online behavior of anonymous users to more efficiently place online display advertisements.

Throughout a conference call and the later interviews, Ms.~Bartz and Mr.~Ballmer emphasized that
combining the No.~2 and No.~3 companies in search would not harm competition but enhance it. Google
was rarely mentioned by name, but it was the subtext of the conversation and the deal itself.

Ms.~Bartz pointedly said the partnership would ``put choice back in the hands of consumers,
advertisers and publishers,'' all of whom, she said, were ``increasingly concerned'' about the
rising power of Google.

Microsoft and Yahoo said they expected resistance from Google. But Microsoft's general counsel,
Bradford L.~Smith, said he looked forward to explaining the details of the planned partnership to
antitrust officials in Washington and Brussels.

``There is a compelling case that this is going to increase competition,'' Mr.~Smith said.

The Microsoft-Yahoo stance, legal analysts noted, amounts to the assertion that Google is so
dominant in Internet search and search advertising that the best way to foster competition and
innovation is with a duopoly -- with the Microsoft-Yahoo partnership creating a credible rival to
Google.

``That is a pretty strong story, and I suspect the one they will present to the Justice
Department,'' said Andrew I.~Gavil, a law professor at Howard University.

\section{Chinese Workers Say Illness Is Real, Not Hysteria\cite{hysteria}}

\lettrine{T}{ian} Lihua was just beginning her morning shift when she felt a wave of
nausea\cite{nausea}, then numbness in her limbs and finally dizziness that gave way to
unconsciousness. In the days that followed, more than 1,200 fellow employees at the textile mill
where Ms.~Tian works would be felled by these and other symptoms, including
convulsions\cite{convulsion}, breathing difficulties, vomiting\cite{vomit} and temporary
paralysis\cite{paralysis}.

``When I finally came to, I could hear the doctors talking but I couldn't open my eyes,'' she said
weakly from a hospital bed last month. ``They said I had a reaction to unknown substances.''

Ms.~Tian and scores of other workers say the ``unknown substances'' came from a factory across the
street that produces aniline\cite{aniline}, a highly toxic chemical used in the manufacture of
polyurethane, rubber, herbicides and dyes.

As soon as the Jilin Connell Chemical Plant started production this spring, local hospitals began
receiving stricken workers from the acrylic\cite{acrylic} yarn factory 100 yards downwind from
Connell's exhaust stacks. On some days, doctors were overwhelmed and patients were put two to a bed.

A clear case of chemical contamination\cite{contamination}? Not so, say Chinese health officials who
contend that the episode is a communal\cite{communal} outbreak of psychogenic illness, also called
mass hysteria. The blurry vision, muscle spasms and pounding headaches, according to a government
report issued in May, were simply psychological reactions to a feared chemical exposure.

During a four-day visit, a team of public health experts from Beijing talked to doctors, looked at
blood tests and then advised bedridden workers to ``get a hold of their emotions,'' according to
patients and their families.

Western medical experts say fear of poisoning can lead people to describe symptoms that exist mainly
in their minds. But outbreaks of psychogenic illnesses on the scale of what has been reported in
Jilin are rare, they say.

The official diagnosis has done little to ease anxieties in Jilin, an industrial city in northeast
China where verdant\cite{verdant} low-rise mountains form a backdrop to a thicket of smokestacks.
More than two months since the health complaints began, at least two dozen people remain
hospitalized, and many others insist that they are suffering from toxic poisoning. Local residents
say the ``mass hysteria'' verdict\cite{verdict} is an attempt to cover up
malfeasance\cite{malfeasance}.

``How could a psychological illness cause so much pain and misery?'' said Zhang Fusheng, a
29-year-old textile worker, gasping as he lay tethered to an oxygen line in the hospital, his limbs
seized up and his eyes darting back and forth. ``My only wish is to get better so I can go back to
work and take care of my family.''

In May more than 1,000 residents blocked railroad tracks in the city for hours to draw attention to
the sick workers. Their ire intensified after the State Administration of Work Safety posted a
statement on its Web site describing the problem as a ``chemical leak'' and advising other companies
to learn from Connell Chemical's mistake. After a few hours, however, the statement had been
removed.

``We are simply laboratory mice in Connell's chemical experiment,'' said Xie Shaofeng, 34, a textile
worker whose wife remains hospitalized.

The episode comes at a time of rising environmental degradation\cite{degradation} in China brought
on by decades of heady growth and lax pollution controls. Although many people here have long lived
with sullied\cite{sully} air and water, they are increasingly aware of the toll that they take on
human health and are demanding greater restrictions on noxious industries.

Fear of contamination was heightened last fall after the government acknowledged that thousands of
children had been made ill by milk adulterated with melamine, an ingredient used in the manufacture
of plastics.

The Ministry of Health in Beijing declined to provide details of their findings in Jilin, but
according to local officials, investigators found no evidence of organ damage that would point to
chemical exposure. They added that those claiming to be sick had been in different parts of the
sprawling textile factory and offered inconsistent descriptions of the odor of what they said caused
their symptoms.

Although they say those who fell ill in Jilin could have been poisoned, psychogenic experts outside
China say it is also possible for some to have been affected by toxic fumes\cite{fume} while others
exhibited psychosomatic illnesses set off by real poisonings.

Robert E.~Bartholomew, a sociologist at the International University College of Technology in
Malaysia, said the government's handling of the episode, including the ban on reports in the news
media, might be fueling paranoia\cite{paranoia}. ``The best way to handle psychogenic illness is to
be open and transparent, which tends to dissipate\cite{dissipate} concerns,'' said Mr.~Bartholomew,
a co-author of ``Outbreak! The Encyclopedia of Extraordinary Social Behavior.'' If it is indeed a
case of mass hysteria, he said, it would be the largest instance of workplace psychogenic illness on
record.

The episode is not Jilin's first experience with the perils of aniline. In 2005, an explosion at
another factory that produced the volatile substance killed eight people and sent 100 tons of deadly
benzene and nitrobenzene into the Songhua River, tainting drinking water for millions of people
downstream.

Public anxiety was high even before the new \$125 million aniline plant opened in early April.
During a test run last September, two security guards standing in front of the textile plant were
overcome by fumes. Connell paid them compensation, although it is unclear what adjustments were made
to the manufacturing process and, more important, the venting of its airborne byproducts, a mix of
carbon monoxide, hydrogen sulfide and nitrogen oxide.

Shortly after production began, Li Hongwei, a 34-year-old Connell worker, collapsed and died on the
job. Although rumors suggested that he had been poisoned, factory officials insist that Mr.~Li
succumbed to a heart attack. His family, which received a compensation package that included a job
for his wife and a monthly \$200 stipend for his mother and son, declined to talk to reporters.

After Mr.~Li's death, the authorities forced Connell to halt production for a month. But in early
June, not long after it resumed, Wang Shulin, a 38-year-old technician at the textile mill, went
into convulsions\cite{convulsion} while on the job. He was sent to the hospital but died just as
doctors were administering a CT scan, according to co-workers. The cause of death was a brain
hemorrhage.

Factory officials insist that Mr.~Wang's death had nothing to do with chemical exposure.

Such assurances have done little to quiet fears that Connell continues to taint the air. Li
Jingfeng, 35, an electrician at an ethanol plant that abuts the aniline plant, said chemical
detectors at his factory had gone off five or six times in the last month, forcing workers to
evacuate. ``Everyone is nervous about what's coming out of that place,'' he said.

Those who continue to insist that they were poisoned have placed local officials in a difficult
position. Some patients have been sent to other cities for treatment; those who refuse to leave
local hospitals say doctors have been given orders to stop their medication. To get the skittish
back to work, factory officials have added an incentive of \$20 to \$30 to monthly salaries that
range from \$120 to \$200.

In interviews, a half-dozen of those still hospitalized in Jilin said they had not been given a
diagnosis nor were they allowed to see their medical records. One of them, Deng Yanli, 30, who is
troubled by convulsions and constant dizziness, showed a receipt for 10 medications that included
vitamin injections, pills to combat nausea and other treatments commonly given to stroke victims.
She said doctors at Jihua Hospital stopped administering the drugs in early June in an effort to get
her to leave.

The hospital director referred questions to the Jilin City Health Bureau, which issued a statement
saying, ``We have done our best diagnosing and treating these patients.''

Officials at Connell, which has resumed full production, say they are eager to move past the
episode. Although privately owned, the plant has a complicated corporate structure that includes
investors from Hong Kong and a number of local government officials. The aniline plant and the
neighboring textile mill are partly owned by one another, and Connell, according to a company Web
site, also runs a pharmaceutical concern that supplies Jilin City hospitals with 90 percent of their
intravenous drugs.

Cementing the company's prominence is its president, Song Zhiping, a representative to the National
People's Congress, China's legislative body.

Connell executives said Ms.~Song was traveling during a reporter's visit to their offices, but Xu
Zhongjie, vice chairman for corporate governance, said Ms.~Song felt wounded by the allegations
against her company, which he described as preposterous\cite{preposterous}. ``I come here every day,
and do I look sick?'' he asked with a broad smile. ``If we were spreading poison, the government
wouldn't allow us to continue production, and I have faith in the government.''

\section{Chinese Officials Harass\cite{harass} Group That Fights Hepatitis\cite{hepatitis}
  Bias\cite{bias}}

\lettrine{I}{n} the realm of potential threats to China's stability, an organization that advocates
on behalf of people infected with hepatitis B would seem to be low risk.

But on Wednesday, the group's director, Lu Jun, found himself squaring off against four security
officials who were trying to cart away stacks of literature they claimed had been printed without
official permission.

In the end, Mr.~Lu scored a partial victory. After eight hours looking through drawers and
photographing volunteers, the inspectors walked off with 90 pamphlets\cite{pamphlet}, but Mr.~Lu
prevented them from delving into the group's computer files. ``I fear this is not the end of it,''
he said Thursday.

The raid on Mr.~Lu's organization, the Yi Ren Ping Center, comes at a precarious time for China's
nongovernmental organizations, many of which operate in a kind of legal gray zone. Two weeks ago,
officials used a bureaucratic infraction as the reason to shut down the country's pre-eminent legal
rights center, Gongmeng, or Open Constitution Initiative. The closing followed a separate disbarment
of 53 lawyers known for taking on civil rights and corruption cases. Just before dawn on Wednesday,
the founder of Gongmeng, Xu Zhiyong, was taken into police custody, and he has not been heard from
since.

``The permissible space in which civil society groups can operate was already small, but right now
that circle is getting smaller and smaller,'' said Sharon Hom, the executive director of Human
Rights Watch, which is based in New York. ``If an organization is creating an independent voice,
putting together a newsletter or organizing people in any way, it's going to feel the full
brunt\cite{brunt} of the authorities.''

Although it is unclear exactly why the government is tightening its grip on such organizations,
legal experts and rights activists generally agree that it may be related to the celebrations, three
months from now, of the 60th anniversary of China's Communist revolution. A similar clampdown took
place in the months before the 2008 Summer Olympics, when security officials in Beijing stepped up
the harassment of dissidents and encouraged thousands of migrant workers to return to the
countryside.

``It's basically a foolish attempt to make the year as peaceful and uneventful as possible,'' said
Jiang Tianyong, a lawyer who was among those blocked from renewing their licenses.

Another explanation, Mr.~Jiang and others say, is that some powerful segments of China's leadership
feel threatened by the rise of independent entities working to advance causes like labor rights or
clean water, or in the case of the Yi Ren Ping Center, protection for people with hepatitis B.

There is widespread trepidation\cite{trepidation} over hepatitis B in China, a fear that has been
intensified by an explosion in advertising for medical testing services and sham\cite{sham} cures.
Even though it is preventable with a vaccine -- and most of those infected will not become ill --
state-owned companies, medical schools and food-processing plants have come to believe that it is
sensible policy to bar the infected.

Under Chinese law, carriers of hepatitis B cannot work as teachers, elevator operators, barbers or
supermarket cashiers. In a recent survey of 113 colleges and universities, conducted by the Yi Ren
Ping Center, 94 acknowledged that infected applicants, required to take blood tests, would be
summarily rejected.

Many of the 120 million carriers in China got the virus in the 1970s and 1980s, when a single
contaminated syringe was sometimes used to inoculate hundreds of people at a time against diseases.
The second-biggest group of carriers, about 40 percent of the total, according to the government,
got the virus from their mothers during childbirth.

An online bulletin board maintained by Mr.~Lu's group is a heart-rending clearinghouse for stories
of people fired from jobs, or students denied college educations, after mandatory blood tests
revealed their statuses. There are also scores of tales about the ashamed and the distraught who
killed themselves.

``People are so afraid of this virus, they don't act responsibly,'' said Wang Li, an engineer who
just graduated from a prestigious\cite{prestigious} Beijing university and saw two job offers
evaporate this year when blood tests showed that he had the virus. ``The only thing they told me
was, 'You are not suitable for work.' ''

Founded in 2006 by Mr.~Lu, who is also infected, the Yi Ren Ping Center provides up-to-date medical
information and tries to arrange legal help for those it considers wrongly dismissed from jobs. It
also encourages its 300,000 members to press for antidiscrimination laws. Last summer the center was
forced to move its Web site to an overseas server after it mysteriously vanished from the Internet.

Although his organization does not seek to challenge the government's authority, Mr.~Lu recognizes
that its mission can stir discomfort among the powerful and mighty. ``After all, it is these people
who are maintaining the status quo of discrimination,'' he said in his office on Thursday. ``And of
course, according to the government, there is no such thing as discrimination in China. There are
only misunderstandings.''

\section{China Sentences Women Held After Protest Attempt}

\lettrine{A}{ woman} and her two sisters, who came to Beijing from southern China during the
Paralympics\cite{Paralympics} in September to protest property seizures but were arrested, have been
sentenced to one year of detention\cite{detention} for vandalism\cite{vandalism}, the woman said in
a telephone interview on Thursday.

The woman, Huang Liuhong, had already been held for nearly a year in a hotel -- known as a black
jail -- while awaiting trial.

Ms.~Huang was released from the hotel, the Liuzhou Cement Factory Rest House, on July 17. She said
she expected to remain under a form of house arrest for one year in her hometown, under police
surveillance.

The case is one of several that starkly illustrate how the Summer Olympics and the Paralympics in
Beijing last year failed to expand freedom of speech in China, despite assertions by the
international organizers of those games that the events would push the Chinese government toward
more democratic policies.

Ms.~Huang traveled with 10 others from the town of Liuzhou in Guangxi Province to Beijing last
September to protest four cases of property seizure involving local officials.

But after being interviewed by an American journalist, they were seized by plainclothes police
officers who had followed them from Guangxi. Ms.~Huang, two older sisters and their 79-year-old
mother, all of whom had traveled to Beijing, were arrested.

The mother was soon released, but Ms.~Huang and her infant son were kept for 314 days in a hotel in
Liuzhou. Her two sisters were held in a detention center.

Ms.~Huang said she and her sisters were not put on trial until June 19. The judge found them guilty
of vandalism and sentenced them to one year in prison followed by two years of
probation\cite{probation}. The two older sisters were released July 14, Ms.~Huang said, each having
served almost the entire sentence under the verdict because of their time in the detention center.

Others from the group of 11 that came to Beijing were also put on trial.

After her release from the so-called black jail, Ms.~Huang was given a week to visit her husband in
the southern city of Shenzhen, where she is now. But she will have to return to Liuzhou and spend a
year under a version of house arrest, during which she can live in her home and move around Liuzhou
but not leave the town, she said.

``They said someone will be watching me,'' she said.

The court in Liuzhou had no immediate comment on Thursday; copies of court documents provided by
Ms.~Huang showed that she had been found guilty of vandalism.

Ms.~Huang said the police accused her of breaking the window of a police van after she was detained
in Beijing. She said she never damaged the van.

The three sisters plan to appeal the verdict, Ms.~Huang said, but such judgments are rarely
reversed.

\section{Motorola Posts a Profit, Aided by Cost-Cutting}

\lettrine{M}{otorola} beat Wall Street expectations Thursday by eking out a small profit in its
second quarter, but the struggling telecommunications company is pinning its hopes for revival on
new smartphones priced far below most current smartphones like the iPhone.

Sanjay K.~Jha, the company's co-chief executive in charge of its wireless handset unit, said in an
interview that the phones will be designed to work with data plans with lower monthly fees than the
\$30 to \$60 common with current smartphones. He foresees also selling the phones in poorer
countries with developing economies.

``In vast portions of emerging markets, the way to get access to the Internet will be through
wireless data connectivity,'' he said.

Motorola, which is selling half the number of handsets it sold a year ago, said it was developing
its new line of phones to run on Google's Android operating system. It has two models set for
introduction this fall, with more coming next year, the company said.

In the mean time, the company is slashing its costs to match its rapidly falling revenue by
canceling products, closing facilities, and firing people. It has eliminated 8,000 of its 60,000
jobs so far this year.

As a result of this cost-cutting, Motorola produced net income of \$26 million, or a penny a share,
compared with \$4 million, or break-even on a per-share basis, a year ago.

Excluding a few onetime items, Motorola lost 1 cent a share, compared with a loss of 4 cents a share
predicted by analysts.

Revenue was \$5.5 billion, down 31 percent from \$8.08 billion a year ago. That is less than the
\$5.6 billion in revenue Wall Street expected. Shares of Motorola rose 9.4 percent Thursday, to
\$7.19.

In the wireless division, revenue fell 45 percent, to \$1.8 billion, in the second quarter. The unit
posted an operating loss of \$253 million, compared with an operating loss of \$346 million a year
ago.

Motorola's other two business lines also showed sharply declining sales, but they were profitable.

Analysts, however, said that the fate of the company rested largely with the success of its new
handsets, which will compete with at least a dozen other Android phones by the end of the year, not
to mention the iPhone, the Palm Pre and various BlackBerry models.

``They have a lot hanging on two devices,'' said Mark Sue, an analyst with RBC Capital Markets.

Mr.~Sue said he expected Verizon and T-Mobile to be the first carriers to offer Motorola's new
phones. Mr.~Jha declined to name the carriers, but he said Motorolahad firm orders from two
companies.

The two handsets, Mr.~Jha said, will have the features now typical of smartphones like the iPhone.
One will be designed to focus on surfing the Internet and the other will specialize in connecting to
social networks. Both will be offered with an extra-fee service from Motorola that will help phone
owners communicate with each other and social networks.

The lower-cost phones to be introduced next year will have fewer features, smaller screens and
cameras with lower resolutions. Moreover, they will be designed to limit the customer's use of
multimedia features that consume a lot of wireless capacity. That will allow carriers to offer
lower-priced data plans.

\section{As Sales Drop Sharply, Kodak Misses Estimates}

\lettrine{E}{astman} Kodak posted disappointing second-quarter results on Thursday as weak demand
hurt its film and digital photography businesses, sending its shares down 12 percent.

Kodak said sales fell sharply as a global recession limited consumer spending on travel and other
activities that spur the usage of cameras and snapshot printers.

Kodak's traditional film business continued to shrink, with sales falling 30 percent in the quarter,
hurt in part by uncertainty over Hollywood labor talks, which have since been resolved.

Kodak posted a second-quarter loss of \$189 million, or 70 cents a share, compared with a profit of
\$495 million, or \$1.72, a year earlier.

Excluding special items, the loss was 43 cents a share. Analysts, on average, expected a loss of 36
cents a share, according to Reuters Estimates.

Kodak said its gross profit margin fell to 18.5 percent, from 23.6 percent, as sales declined and
intellectual property licensing royalties decreased.

Revenue fell 29 percent, to \$1.77 billion, from \$2.49 billion. Analysts expected \$1.83 billion,
according to Reuters Estimates.

Analysts shared shareholders' lack of enthusiasm over the results, with several calling attention to
Kodak's ability to generate cash. Kodak had negative cash flow of about \$158 million in the
quarter.

Kodak has spent the last six years shifting its focus from cameras and traditional film to digital
printing systems and long-life photographic ink. It has said it needs partners to further develop
some of its cutting-edge technologies, like CMOS image sensors and superthin OLED screens, but tight
credit markets have muted plans by other companies to invest in Kodak's commercial printing systems
and services. Its shares fell 37 cents to \$2.91 on the New York Stock Exchange.

\section{In a Quebecer's Heart, Pepsi Occupies a Special Place}

\lettrine{I}{n} 1967 Charles de Gaulle created a media and diplomatic storm after he fueled Quebec's
separatism by shouting ``Vive le Quebec libre!'' from a balcony at Montreal's city hall.

Last month the former French president's endorsement of a Quebec nation reappears in a television
commercial selling Pepsi to the province's French-speaking majority.

The ad, which marks the soft drink maker's 75th anniversary in Canada, is in keeping with Pepsi's
unusual history. Since the 1970s, Pepsi's Quebec bottling unit has gone its own way, using marketing
that appeals, in various ways and with different degrees of subtlety, to the nationalistic sense of
the province's French-speaking population.

Quebec is one of the developed world's few markets where the positions of Pepsi and Coca-Cola are
reversed. According to ACNielsen MarketTrack, Pepsi's main brand commands 29.9 percent of the retail
soft drink market in Quebec, based on volume, compared with 12.3 percent for Coke.

And all of the brands owned by Pepsi control just over 60 percent of the province's soft drink
business -- a dominance that even led to a common slur\cite{slur}. The Canadian Oxford dictionary
defines ``pepsi'' as derogatory term for a French Canadian, ``from the perceived Qu\'eb\'ecois
preference for Pepsi-Cola.''

Until the 1970s, Pepsi followed the approach still used by most advertisers for the Quebec market.
Television spots developed for English Canada were double-shot with French-speaking actors
substituted if there were speaking roles. For other ads, it was simply a matter of changing onscreen
text and substituting French voiceovers.

Perhaps there is no connection -- certainly no one at Pepsi will draw the link -- but Pepsi's
intense focus on the Quebec market coincided with the political changes of the 1970s which
eventually brought the separatist Parti Qu\'eb\'ecois to power and to many English-speaking
residents leaving the province. While Pepsi was being promoted by Michael Jackson in the rest of the
world 25 years ago, the company hired Claude Meunier, a moderately successful member of a Quebec
comedy duo, to perform tiny sketch comedy pieces as Pepsi ads. Using broad, often slapstick humor,
Mr.~Meunier spoofed a variety of stereotypical Quebecers, including a hockey player who seemed to
have taken one too many pucks to the head.

Mr.~Meunier's characters often spoke slang-laden joual, a dialect that was once mostly heard in
working-class, French neighborhoods of Montreal. The ads turned Mr.~Meunier into a big celebrity in
Quebec and Pepsi became the province's top soft drink. (Unusually, Mr.~Meunier still appears in some
of Pepsi's Quebec ads, although since 2003 he has mainly promoted Diet Pepsi.) In recent years, the
Quebec message has been more overt in ads. One mocks Quebec stereotypes from a peculiar provincial
law that sets June 30 as the expiration date for most residential leases (and thus transforms July 1
into a chaotic, provincewide moving day) to a Quebec City so cold that an actor must consume a
bottle of Pepsi as if it were a Popsicle.

English Canadians don't escape notice. One ad features a man that Pepsi describes simply as ``a
tourist'' but who dresses unfashionably and acts in a manner that many Quebecers associate with
English Canadians. He orders a Coke at a Quebec City restaurant, provoking evident hostility among
the French-speaking patrons. Its tagline translates roughly as: ``Here, it's Pepsi.''

The cameo appearance of Charles de Gaulle is not the only portion of the fast-paced and largely
humorous 75th anniversary commercial (created by BBDO Montreal), which might offend some. One
sequence lumps images of Queen Elizabeth II, Canada's head of state (or more precisely her waving
hand and her hat), and Pope John Paul II with Bonhomme, the snowman mascot of Quebec City's winter
carnival.

While acknowledging that Pepsi's Quebec ads sometimes push limits in the interests of humor, Sylvain
Charbonneau, vice president for Pepsi's Canadian bottling group and its director general for Quebec,
said that during his 20 years with the company in Quebec, he had never known it to be second-guessed
by the Canadian head office, which is in suburban Toronto.

``It's easier for us to have credibility with 60-point share, but that wasn't the case 25 years
ago,'' he said in a boardroom at a recently upgraded bottling plant in an industrial area of
Montreal. ``They listen to us, give a chance to take risks, make mistakes. It was important for them
to have someone who understands the market.''

Coke has done some Quebec-specific advertising, although none for the last two years. In advance of
next year's Winter Olympics, which will be held in Canada, David M.~Moran, a spokesman for Coca-Cola
Canada of Toronto, said that the company would feature Quebec athletes in forthcoming commercials.
He said the company had signed ``a major Quebec celebrity'' to act as a spokeswoman, although he
declined to identify her.

While tailoring its message to French-speaking Quebecers has been successful for Quebec,
Mr.~Charbonneau acknowledges that the approach may not be appropriate for many companies.

``It's more costly for a population of seven million to do a unique advertising campaign year after
year,'' he said. ``It costs you more per capita than it would for the U.S.~with 300 million people.
But it pays out for us and Quebecers don't mind laughing about themselves.''

\section{Ads Follow Web Users, and Get Deeply Personal}

\lettrine{F}{or} all the concern and uproar over online privacy, marketers and data companies have
always known much more about consumers' offline lives, like income, credit score, home ownership,
even what car they drive and whether they have a hunting license. Recently, some of these companies
have started connecting this mountain of information to consumers' browsers.

The result is a sea change in the way consumers encounter the Web. Not only will people see
customized advertising, they will see different versions of Web sites from other consumers and even
receive different discount offers while shopping -- all based on information from their offline
history. Two women in adjoining offices could go to the same cosmetic site, but one might see a
\$300 Missoni perfume, the other the house-brand lipstick on sale for \$2.

The technology that makes the connection is nothing new -- it is a tiny piece of code called a
cookie that is placed on a hard drive. But the information it holds is. And it is all done
invisibly.

``Now, you're traveling the Internet with a cookie that indicates you're this type of consumer: age
group X, income level, urban versus rural, presence of children in the household,'' said Trey
Barrett, a product leader at Acxiom, one of the companies offering this linking to marketers.

Advertisers and marketers say this specificity is useful, taking out the guesswork involved in
online-only profiling, and showing products to the people most likely to be interested. Retailers
including Gap, Sephora and Victoria's Secret are using this tactic.

But consumer advocates say such unseen tracking is troubling. On the old Internet, nobody knew you
were a dog. On the new targeted Internet, they now know what kind of dog you are, your favorite
leash\cite{leash} color, the last time you had fleas and the date you were neutered.

``The industry's love affair with persistent cookies has made it virtually impossible for users to
go online without being tracked and profiled,'' said Marc Rotenberg, executive director of the
Electronic Privacy Information Center, in an e-mail message.

While Congress has been holding hearings on online privacy lately, the sessions have focused on
online behavioral targeting. The industry has argued that no government intervention is needed, an
argument that the Federal Trade Commission has so far accepted.

Consumers can avoid cookie-based tracking by deleting cookies from their computers or setting their
browsers not to accept cookies. But few do, and privacy advocates say it is easy for companies to
add cookies without users noticing.

For decades, data companies like Experian and Acxiom have compiled reams of information on every
American: Acxiom estimates it has 1,500 pieces of data on every American, based on information from
warranty cards, bridal and birth registries, magazine subscriptions, public records and even dog
registrations with the American Kennel Club.

Patrick Williams, the publisher of the personal finance magazine Worth, recently asked Acxiom to
find the names and addresses of 10,000 Americans from each of 11 cities who had houses worth more
than \$1 million, net worth of over \$2 million, lived within a few miles of other rich people and
subscribed to business publications.

``They are the scariest data research company around -- they know far too much,'' said Mr.~Williams,
who said he was very happy with the amount of information it gave him.

Companies like Acxiom and a competitor, Datran Media, make the connection between online and offline
data when a person registers on a Web site or clicks through on an e-mail message from a marketer.

Datran's cookies include 50 to 100 pieces of information. Both companies say cookie data is
anonymous and generalized. Datran and Acxiom then sell advertising on Web sites like NBC.com,
Facebook and Yahoo to companies that use their data.

For marketers, all this data is a boon. Beltone New England, a hearing-aid company, asked
Datran to find people online who were 65 and over, owned a house, were head of a household, made
more than \$35,000 a year and lived in New England so it could show them ads. Datran also tested the
same ads with a wider group of people.

``What was surprising was we found the majority of responders turned out to be women 35 to 40 who
had elderly parents at home,'' said Perry Ebel, Beltone's director of marketing and business
development. He said he was changing his offline marketing to include that group.

By using real-world data online, marketers can customize messages even further -- showing different
products to people with different shopping habits, whether it is in ads, an e-mail message or in
semipersonalized Web pages.

Rodale, which publishes books along with magazines like Men's Health and Prevention, uses Acxiom
data to help determine which promotional e-mails to send to which customers. Offers aimed at women
might be accompanied by an e-mail message offering a Father's Day subscription to Men's Health for
him and a free book on losing belly fat for her. Young men might get another offer -- a book on sex
positions. Some marketers are using offline data more subtly -- for example, showing a budget
shopper a discount offer and a regular shopper a full-price section.

``The people who buy less frequently and are most price-conscious may get a better deal than someone
who buys more frequently, who would buy anyway,'' said Christopher S.~Marriott, global managing
director of Acxiom Digital, a division of Acxiom.

Of course, shoppers would have little reason to think their experience or their ads are being
personalized based on their home value or Volvo ownership.

``It is a little Big Brother-ish,'' said Betsy Coggswell, 49, a social worker in Fullerton, Calif.,
who shops online regularly. Still, she said, she wasn't shocked. ``Every time you put out
information about yourself -- people have got to understand -- it's going to be collected by
somebody.''

Some online companies avoid matching online and offline profiles. In 2000, DoubleClick abandoned
plans to connect online and offline data after a huge outcry. Google, which later acquired
DoubleClick, has been conducting studies that connect the two areas, but it does not currently
collect or serve ads based on such personal information without user permission, Sandra Heikkinen, a
Google spokeswoman, said.

While Acxiom, Datran and some of their partners address their use of tracking in their privacy
policies, such policies have become worthless, Mr.~Rotenberg said. ``Real transparency means that
the user gets access to the information, not to a policy about the information,'' he said.

Paul M.~Schwartz, a law professor and privacy expert at the law school of the University of
California, Berkeley, said the unwitting participation by consumers makes online marketing different
from offline.

``Interactive media really gets into this creepy Orwellian thing, where it's a record of our
thoughts on the way to decision-making,'' he said. ``We're like the data-input clerks now for the
industry.''

\section{With WiMax, Walking on the Wireless Side in Baltimore}

\lettrine{B}{altimore}, the home of the hapless Orioles and a favorite backdrop for so-called
realistic TV crime dramas, just happens to be one of the most wired cities in the country. It
ascended to the throne after Clearwire introduced WiMax service last winter, giving the city a
preview of what the company is slowly building throughout the country.

Calling the city ``wired'' is not quite accurate because, HBO series aside, WiMax has little to do
with copper; it delivers the Internet through radio signals broadcast from cellphone towers. The
service is much faster, though, than what many cellphone networks currently provide with their 3G
networks. Indeed, WiMax delivers speeds much faster than many DSL circuits, rivaling many cable
modems. I often clock downloads at 6 megabits per second (equivalent to basic cable service in many
areas) and uploads at faster than 1 megabit per second.

Sprint, one of the owners of Clearwire, calls WiMax a ``4G network'' to distinguish it from the 3G
networks that connect smartphones and offer speeds from 0.4 to 1 megabit per second. To make matters
a bit more confusing, Clearwire sells ``pre-WiMax'' service in 47 cities at speeds that top out at 2
megabits per second. In an effort to clarify the matter, Clearwire is using the brand name ``Clear''
to apply to full WiMax service.

For the last six months, I've used a full WiMax/4G equipped netbook to test the service around
Baltimore. The Acer Aspire One with a Sprint U300 WiMax card I used is an ideal companion for
sending and receiving e-mail messages. It's small enough to take almost everywhere but it's large
enough to act like a PC -- a PC that's always connected to a very big Wi-Fi hotspot.

Adding WiMax to a laptop may make it easier to read e-mail messages often, but the real value of the
bandwidth becomes apparent when the PC does something more than just send text. VoIP software like
Skype turns it into a video phone, a browser pointed at Hulu acts like a television that can fetch
shows on command and there's also GPS service for finding directions. It's a smartphone with a
normal keyboard and a very open software marketplace. All of the PC software built for the desktop
also works with the small package.

The U300 attaches to any laptop, but Netbook manufacturers are offering more and more machines with
built-in cards. Lenovo, for instance, will add the WiMax/Wi-Fi Link 5150 card to some of their
laptops for an additional \$30. Prices vary depending on the model.

WiMax is also one of the first wireless services that's being actively marketed to people sitting on
their couches at home. Xohm, the brand name originally given to Baltimore's WiMax service (it's now
being merged into Clearwire and the Xohm brand will disappear), sells a home base station meant to
compete directly with DSL or cable for \$35 a month, and sometimes there is a short-term discount. A
base station and a laptop card together cost \$50 a month. The service options are getting
complicated and the prices vary in different communities. Comcast, another investor in Clearwire, is
starting a wireless service in Portland, Ore., under its own brand name. It will charge \$50 a month
for service within Portland that is promised to deliver downloads at a rate of 4 megabits per
second.

Sprint also sells a card that offers both 4G service where WiMax is available and 3G service where
Sprint's cellphone network is all that's there. It's a better choice for anyone who travels outside
of the cities where WiMax is appearing. Sprint charges \$80 for the modem and \$80 a month for the
service.

The WiMax coverage in Baltimore is good but far from comprehensive. The signal blankets downtown and
many of the neighborhoods, but it stops just a few blocks from my house. The cellphone tower is on
one side of the hill and my house is on the other. Clearwire says that it is slowly expanding
coverage but I've seen little change in the map over the last few months.

WiMax can also suffer from the same problems that affect all wireless services. Rain and snow absorb
the signal, reducing the quality of the service during storms, an effect the industry calls ``rain
fade.'' Trees and other plants are filled with water and can cause the same problems even when the
sun is shining. Thick walls are also a challenge. Being closer to the tower is always better for
service. All of these effects work together, so it's no surprise that the maps of WiMax service show
that early deployment is concentrated in the densest part of the city where trees are rare. That's
where the most people will find the best reception.

But wireless also comes with advantages. I've averaged about one visit from the phone or cable
company every year or so because the copper wires coming to my house need their care. Both services
require internal wiring that must either be fished through the walls or the baseboard. Wireless
service to the home avoids these problems.

The Xohm/Clearwire base station can sit anywhere in the house and it can even be moved, but putting
it on a higher floor near a window improves service. If the house isn't near the tower, it may even
help to put it on the closest side of the house where there's no dense foliage in the path. Using
the WiMax laptop card alone also works under the same conditions.

The more I used the laptop while traveling around the city, the more aware I became of the time it
took Windows to start. While the Internet service may be available everywhere, it took several
painful minutes for Windows to boot up. To make matters worse, the U300 card needs its own minute or
so to look for a signal.

Some netbook manufacturers are experimenting with adding a simpler operating system that can start
much faster than Microsoft Windows. These machines can boot up in less than 10 seconds but only by
loading a lightweight operating system that offers a few basic services like a Web browser or Skype.
The full version of Windows is still available if you need it. Needless to say, offering such a
start-up time will change the utility of these microlaptops considerably.

Speeding the start-up time will be crucial if the netbooks want to compete with smartphones for
casual use by people on the go. The bandwidth is ready to supply full-size applications that augment
reality with an endless heavy stream of data. Scott Richardson, chief strategy officer of Clearwire,
told me: ``I did a demonstration in Portland with some computer industry guys. I was driving down a
road going 60 miles per hour and I got a 14 megabits-per-second download.''

Portland came online in January of this year and Clearwire just announced that it was selling
service in Atlanta and Las Vegas a few weeks ago. This isn't the end. Clearwire says it's on to
Chicago, Dallas-Fort Worth, Honolulu, Seattle and Charlotte, N.C. They're all fine cities, even if
they don't all have the same level of baseball and the same fertile source of inspiration for
narrative crime dramas.

\section{When Auto Plants Close, Only White Elephants Remain}

\lettrine{T}{he} sheer size of the sites has inspired grand visions for redevelopment -- a \$1
billion football stadium, a huge Hollywood movie studio, even the world's largest indoor tennis
complex.

But for the communities saddled with a huge, empty auto plant, the reality is dismal\cite{dismal}.

Abandoned car factories, sprawling over hundreds of acres, often stand vacant for years awaiting
demolition, environmental cleanup and a willing developer.

Since 2004, General Motors, Ford and Chrysler have closed 22 major auto plants in the United States.
Only eight of those have found buyers. And in the wake of the G.M. and Chrysler bankruptcies,
another 16 plants will be shut by 2011.

The most optimistic redevelopment proposals, like a football stadium for the Atlanta Falcons or a
movie studio in the small Michigan town of Wixom, a Detroit suburb, are long shots at best.

``The plants, whether they're still standing or reoccupied, are always going to be a haunting
reminder of what we were, what we've gone through, and where we still need to go,'' said
Representative Thaddeus McCotter, Republican of Michigan, whose district includes an old Ford plant
in Wixom and a G.M. plant that will soon close.

Even sites in attractive locations are hard to sell in the weak economy. With so many companies
being squeezed financially, there is a glut\cite{glut} of available commercial real estate.

``Even if you only go back three or four years, it was easier than today,'' said Phil Horlock, head
of Ford's land development division.

The loss to the local community when a plant closes goes well beyond jobs. Tax revenue evaporates
and related businesses vanish.

Industry analysts estimate that each job in a plant helps create another five to seven jobs.

``Some of those are direct suppliers, but then there are places that workers spend money, like
grocery stores, restaurants and day care,'' said Kristin Dziczek of the Center for Automotive
Research in Ann Arbor, Mich.

Ford's 4.7-million-square-foot Wixom factory, which closed in 2007, was the company's largest
assembly plant in the United States. More than six million cars were built there over 50 years.

At its peak in the late 1980s, the factory employed nearly 4,000. Now it's an empty shell of rusting
corrugated metal, surrounded by desolate parking lots and a barbed wire fence.

The plant fronts an Interstate highway, and stretches almost a mile along Wixom Road. It once
provided 40 percent of the town's property taxes, but now accounts for less than 15 percent.

The town has about 13,000 residents, and relied heavily on the paychecks of plant workers.

``When it closed, a lot of businesses around us closed too,'' said Moe Leon, owner of the Bullseye
Sports Bar and Grill on Wixom Road. ``We're fighting night and day to stay above water.''

Last year, a team of executives from Warner Brothers toured the plant as a potential site for a new
studio. Other developers have floated proposals for a hotel, an ice hockey arena and a business
district devoted to green technology.

Ford officials said several ideas were under consideration, but there was no timetable for a sale.

Mr.~Horlock said Ford had sold five large factories in the last five years. He said the company
generally did not raze a plant, or begin to clean up any toxic wastes, until a deal was sealed.

The auto company sold one assembly plant, in Lorain, Ohio, to a developer who leases out space to
small industrial firms. Honda, a competitor of Ford's with two assembly plants in Ohio, recently
started storing excess inventory of cars and minivans in the plant's parking lots. But most closed
plants languish.

After G.M. closed a factory in Doraville, Ga., last year, there was an initial rush of interest from
developers. One of them suggested the site be used for a new stadium for the National Football
League's Falcons, but the team's owners have so far shown little enthusiasm for that idea.

Luke Howe, an assistant to Doraville's mayor, said other proposals for the 165-acre site range
``from the ridiculous to the sublime,'' and no developers had come forward with adequate financing.

``We knew it would be a lengthy process,'' he said. ``It just happened to fall during one of the
worst economic times.''

But some factory sites have remained vacant through good times as well -- a sign of just how
difficult it is to create a second life for them.

G.M. demolished most of its giant Buick City manufacturing complex in Flint, Mich., a decade ago.
The 200-acre property appeared to finally have a future as a transportation hub for long-haul trucks
and rail cars, but G.M. could not complete a deal before filing for bankruptcy on June. 1.

The Flint site, along with dozens of other factories and properties, is part of the old G.M. that is
still in bankruptcy and will most likely be sold, eventually, in the liquidation process. The
company's best assets were transferred to a new corporate entity that now operates as General
Motors. Auto companies sell their old plants for a song, compared to what they put into them. The
government in Oklahoma County, which covers much of Oklahoma City, for example, paid \$55 million
last year for a four-million-square-foot plant that will be leased to the United States Air Force to
expand a nearby base. Just eight years ago, G.M. invested \$700 million in the plant to modernize
it.

Environmental problems can also hamper a sale. While every plant has different pollution issues,
G.M.'s restructuring chief, Albert A.~Koch, estimated in bankruptcy court that the company's
liabilities for all its closed sites was \$530 million.

One old G.M. plant that has found a new life is in Linden, N.J.~The company shut down the factory in
2005, demolished it two years ago, and sold it to a developer for a mixed-use project called Legacy
Square.

The 104-acre site will be anchored by a Super Wal-Mart store, and ultimately will pay more taxes to
the city than G.M. did.

But the 2,400 permanent jobs will fall far short of the 6,000 workers that G.M. once employed in the
city.

``The United Auto Workers had very high-paying jobs there, and those will not be comparable to the
jobs that are going to be in the shopping mall,'' said Linden's mayor, Richard J.~Gerbounka.

But, he added, ``We're better off than having 104 acres of vacant property sitting there.''

The worst of the empty auto plants is located, perhaps fittingly, on the downtrodden east side of
Detroit.

The 3.5-million-square-foot factory has been crumbling since the Packard Motor Car Company closed
its doors more than 50 years ago. Trees grow on the plant's roof, and chunks of concrete regularly
fall from the bridge that connects two of its buildings.

Trespassers often explore its rotted interior, and photographs and videos of the ruins are easily
found on the Internet.

Vandals have set fires several times this year in the piles of wooden pallets, tires and garbage
that litter the complex. It is not unusual to see clouds of thick smoke pouring from the building on
a summer evening.

\section{A Primary for Senator? Get Serious}

\lettrine{T}{he} Empire State is a splendid nickname for New York, maybe perfect. Some leading
political lights act as if the state is their empire, not part of a republic in which supreme power
is supposed to rest with its citizens.

Democratic Party elders have worked overtime to keep Kirsten E.~Gillibrand, New York's appointed
junior senator, from facing any serious threat in a 2010 primary. With a mighty assist from the
Obama White House, they have tried to frighten off anyone who could give her a tough run. They seem
quite good at being scary. A Long Island congressman, Steve Israel, got the message and dropped all
thoughts of a race. So did Scott M.~Stringer, the Manhattan borough president.

A few people may yet run. But the only one with reasonably broad name recognition is Representative
Carolyn B.~Maloney of the Upper East Side.

Ms.~Maloney has made a big deal of saying that she will announce her candidacy. A formal declaration
of sorts was expected early this week. Nothing happened. Those who pay attention to these things
have been left to wonder if an announcement is still to come, or if waiting for Ms.~Maloney ranks
with waiting for Godot. Or maybe Guffman.

Whatever the congresswoman decides, the Empire State's Democratic chieftains could not be blunter
about their druthers: They don't want Ms.~Gillibrand to break into so much as a mild sweat in a
primary.

Politicians are always prattling on about the glories of democracy. They boast about how they got
rid of those nasty smoke-filled back rooms, where sebaceous party hacks used to select the
candidates, and replaced them with primaries and with openness.

Truth is, politicians hate primaries. Those who are in command do, anyway. They go out of their way
to knock off anyone who opposes their preferred candidates, apparently in the belief that matters of
such weight cannot be left to lesser mortals, sometimes known as voters.

In this regard, New York Republicans and Democrats are equal offenders. The state Republicans under
Gov. George E.~Pataki worked hard to shut out any challenge to their designated favorites in the
state's presidential primaries: Bob Dole in 1996 and George W.~Bush in 2000. Party bosses would have
gotten their way, too, had the courts not intervened and said, uh-uh.

With Ms.~Gillibrand, the Democratic desire to eliminate democratic competition is especially
striking because not a single voter has had a chance to pass judgment on whether she belongs in the
United States Senate.

It's not just that she was appointed. The man who named her, Gov. David A.~Paterson, was not elected
to his own office. It's a double whammy, intensified for some New Yorkers by Ms.~Gillibrand's -- how
to put it -- demonstrated flexibility on certain core issues.

``What's the point of a Democratic primary if we don't use that process?'' said Liz J.~Abzug, a
political consultant who supports Ms.~Maloney and has opposed those trying to keep the congresswoman
out of the running.

WHEN sheer political muscle doesn't bump off would-be contenders, pettifoggers can be counted on to
limit ballot access, even with incumbents.

We just saw an example of this in the Democratic primary race for public advocate. City Councilman
Bill de Blasio had collected enough voter signatures to choke a stable of horses. Yet the city's
Board of Elections knocked him off the ballot because his papers contained a meaningless typo. After
a few days, common sense prevailed; the board let Mr.~de Blasio back in the race.

But it has stood by its decision to shut out City Councilman Alan J.~Gerson, who is seeking
re-election in Lower Manhattan. Mr.~Gerson also had some sloppy paperwork. He plans to ask a state
court next week to restore him to the Democratic primary ballot.

His lawyer is Lawrence A.~Mandelker, who did the Republican leaders' handiwork back when they tried
to lock unfavored candidates out of those presidential primaries. He has changed, Mr.~Mandelker
said. His conversion hardly compares to that of St. Paul on the road to Damascus. Still, he said,
working to keep qualified people off ballots ``becomes more tawdry the older you get.''

But then, the words ``tawdry'' and ``politics'' tend to find each other in the Empire State, whose
motto is ``Excelsior,'' often translated from the Latin as ``Ever Upward.'' Perhaps ``Exasperate''
would be more on the money.

\section{Obama and Sister to Share a Town}

\lettrine{T}{he} far-flung Obama clan\cite{clan} is coming closer together. President Obama's
half-sister, Maya Soetoro-Ng, is moving her family here from Hawaii and will spend the next several
months living in the nation's capital, White House officials say.

The extended stay means that the siblings will live in the same city, at least for a while, for the
first time in years.

Ms.~Soetoro-Ng, who has given up her job as a high school teacher in Hawaii, has a book contract and
a new baby. Her husband, Konrad Ng, a professor at the University of Hawaii, will become the
scholar-in-residence at the Smithsonian Institution's Asian Pacific American Program here next
month.

Museum officials say Mr.~Ng will stay in Washington through December, when he will return with his
family to Hawaii. There, he will complete his yearlong term as resident scholar. The family is not
expected to move into the White House, though they are very close to the Obamas.

``Maya made many trips to Chicago and elsewhere throughout the campaign, and she and her brother
have a great relationship,'' said Alan S.~King, a lawyer in Chicago who is a friend of Mr.~Obama.

Mr.~King said he believed that the president would be ``very happy to have Maya and her family close
by.''

Katie McCormick Lelyveld, a spokeswoman for Michelle Obama, said Ms.~Soetoro-Ng was unavailable for
comment on Thursday.

In a recent interview, Ms.~Soetoro-Ng, 38, described the first lady as a ``big sister'' figure in
her life. And Ms.~Soetoro-Ng, who campaigned on Mr.~Obama's behalf before the presidential election,
has often spoken warmly about her relationship with her older brother, which she says has remained
strong even though they have often lived far apart.

The siblings, who share the same mother, Stanley Ann Dunham Soetoro, spent several years together in
Indonesia and in Hawaii before Ms.~Dunham decided to return to Indonesia with Ms.~Soetoro-Ng while
Mr.~Obama remained in Hawaii with his grandparents. (Ms.~Soetoro-Ng's father was Indonesian;
Mr.~Obama's father was Kenyan.)

Still, Ms.~Soetoro-Ng has said, they always stayed connected.

Growing up, she said, Mr.~Obama introduced her to music, literature and new neighborhoods. He helped
her decide which dates she should keep and which she should ditch. He spoke at her wedding and
helped her grieve when her father died.

As adults, they often celebrated Christmas in Hawaii, playing Scrabble -- Mr.~Obama
trounced\cite{trounce} her in an intense match the month before the Iowa caucuses -- and savoring
time with their families.

And after their mother died of ovarian cancer in 1995, they stood together near the water's edge in
Hawaii and scattered her ashes in the Pacific.

``He's someone who really enlarged my perspective,'' Ms.~Soetoro-Ng told the Canadian television
network CTV in November.

``He took me to colleges around the country to help me make a good decision about where to go to
school,'' she said. ``He took me with him to go work in the South Side of Chicago and to help
communities and to do voter registration and showed me that it was important to be active and to be
an activist.''

Since the presidential election, Ms.~Soetoro-Ng has spoken at several universities in California,
emphasizing the role that educators can play in promoting diversity and conflict resolution.

In April, she landed a contract from Candlewick Press of Cambridge, Mass., to write a children's
book, ``Ladder to the Moon,'' that is inspired by her mother and her 4-year-old daughter, Suhaila.
She is also working on a book about peace education and conflict resolution in high schools.

School officials at La Pietra Hawaii School for Girls, where Ms.~Soetoro-Ng taught history, said she
left her job when her contract ended at the end of the school year. And she has already spent some
time here this summer.

In June, she served on the President's Commission on White House Fellowships, which selects talented
young people to spend a year working for senior White House staff members and other officials.

\section{House Bucks President on Spending for Military}

\lettrine{T}{he} House on Thursday approved President Obama's plan to kill the F-22 fighter jet. But
Democratic leaders bucked White House veto threats on other programs, and they heatedly rejected a
Republican effort to strip more than 550 earmarked expenditures from the \$636 billion military
bill.

Mr.~Obama and other political leaders had hailed\cite{hail} last week's vote in the Senate to cancel
the F-22 as a sign of their progress in changing military spending practices.

But in sometimes tense exchanges on the House floor on Thursday, two Republicans, Representatives
Jeff Flake of Arizona and John Campbell of California, sought to embarrass some other
representatives and to suggest that little had changed in how Congress pushes pet military projects.

While the House voted 269 to 165 to approve an amendment that stripped out money for building more
F-22s, it overwhelmingly rejected efforts by Mr.~Flake and Mr.~Campbell to cut up to \$2.7 billion
in earmarks, including money that lawmakers had inserted on behalf of specific companies on 553
smaller projects.

The bill also included more than \$1 billion to continue work on larger projects the administration
wants to kill, like a new presidential helicopter, and nearly \$1.2 billion for combat planes that
the Pentagon did not request.

But in an echo of the Republican accusations of pork-barrel spending in the huge economic stimulus
package this year, much of the debate on the military bill focused on earmarks like \$1.5 million
for a gunshot-detection system and \$8 million for a project to upgrade torpedo capabilities.

Mr.~Flake said \$200 million of the earmarks were requested by companies that had been clients of
the PMA Group, a lobbying firm that is under investigation over its ties to several congressmen,
including Representative John P.~Murtha, the influential Pennsylvania Democrat who was in charge of
the bill.

Mr.~Flake, a longtime critic of earmarks, said some of the companies had performed poorly on
government work. He also said many of the companies had raised substantial amounts of campaign money
for Mr.~Murtha and others legislators, including Republicans, who sponsored their projects.

``That simply doesn't look right,'' Mr.~Flake said, adding that the House ``ought to want to have a
higher standard.''

Mr.~Murtha, who put together the bill as the chairman of the House military appropriations
subcommittee, defended the practice and lashed out at Mr.~Flake.

Mr.~Murtha said he had long sought earmarks as a way to create jobs in Pennsylvania. He said the
state had sent a disproportionate number of soldiers to Iraq and Afghanistan compared with the level
of military contracts it had typically received.

He said that he had an obligation to bring work to the small businesses in his area and that some of
the earmarks, like \$2.2 million for body-armor research, grew out of his visits with troops in
Afghanistan and with wounded troops in hospitals.

``I'm not sure how often Mr.~Flake goes to the hospital,'' Mr.~Murtha said.

Mr.~Murtha, who was visibly angry, said several times that the bill included language instructing
the Pentagon to seek competitive bids on the earmarked projects, and other lawmakers noted that
Congress has the final say on spending decisions.

But Mr.~Flake said that in the past, most companies that had sought the earmarks ended up with the
contracts.

The overall bill, which would set military appropriations for 2010, passed by 400 to 30. The Senate
will take up its version later this year, and the two bills will need to be reconciled in
conference.

Mr.~Obama had repeatedly threatened to veto any bill that included more money for the F-22, the
world's most advanced fighter, as the Pentagon seeks to shift more from high-tech weaponry to
simpler systems the troops can use now.

After the Senate voted 58 to 40 last week to cancel the F-22, Mr.~Murtha decided to give up the
fight for it as well. He offered the amendment on Thursday that removed \$369 million for parts,
which had been meant to be a down payment toward 12 more planes.

The administration has also received support from Congress to scale back a \$160 billion Army
modernization plan and other programs.

But the White House warned this week that the president's advisers would consider recommending a
veto if the House went ahead with plans -- as it did Thursday -- to try to save the troubled effort
to create a new presidential helicopter and to finance development of an alternative engine for
another new fighter plane, the F-35.

After Thursday's vote, Tommy Vietor, a White House spokesman, said the president had asked for a
review of the House bill, including the earmarks.

Mr.~Vietor said that over all, Congressional spending bills included 25 percent fewer earmarks this
year than last year.

Still, he said, ``We knew from the start that changing the old ways of Washington would not happen
overnight.''

\section{`Clunkers' Auto Rebate\cite{rebate} Plan So Popular That It's Broke}

\lettrine{N}{ew-car} shoppers appear to have already snapped up all the \$1 billion that Congress
appropriated for the ``cash for clunkers'' program, leading the Transportation Department to tell
auto dealers Thursday night to stop offering the rebates.

But a White House official said the program had not been suspended, creating confusion about its
status. The program offers \$3,500 to \$4,500 for people who trade in an old car for a new one with
higher fuel economy.

In a statement issued Thursday evening, Robert Gibbs, the White House press secretary, said: ``We
are working tonight to assess the situation facing what is obviously an incredibly popular program.
Auto dealers and consumers should have confidence that all valid CARS transactions that have taken
place to date will be honored.''

The program, formally known as the Car Allowance Rebate System, was scheduled to be offered until
Nov. 1, or as long as the money was available. But the program was so successful that it has
exhausted all the money allocated within the first week. Dealers have submitted applications on
behalf of consumers seeking rebates on about a quarter-million vehicles.

The National Automobile Dealers Association surveyed its members in recent days and warned the
Transportation Department on Thursday that it had a very large backlog of applications, said Bailey
Wood, a spokesman for the association.

Late in the day, the group said the Transportation Department had responded by telling it to stop
taking applications at midnight. The government and the dealers were concerned that buyers would
close trade-in deals to buy new cars assuming they had a big rebate coming only to discover later
that money was not available.

The dealers' group said late Thursday night that it had not heard about the White House policy
reversing the decision. Mr.~Wood said that his group would ask Congress and the White House to add
money to the plan.

Transportation Secretary Ray LaHood has already been making calls to members of Congress, telling
them about the situation. The Michigan delegation was planning a meeting Friday morning to discuss
the situation, a Congressional aide said.

On Thursday evening, the government Web site describing the program, www.cars.gov, still showed a
chart shaped like a fuel gauge that indicated \$779 million was available for trade-ins of cars and
light trucks. Earlier Thursday, the Transportation Department issued a news release that said that
applications for fewer than 23,000 vehicles had been submitted as of Wednesday, with a rebate value
of just under \$100 million.

The Transportation Department had begun accepting applications for the rebates on Monday, when rules
putting the program in place took effect. But car dealers had been accumulating the applications
since July 1, when Congress put the law into effect.

The program had two goals: aiding the ailing car industry and improving fuel economy of the vehicles
on the road.

Cars submitted under the program were to be junked\cite{junk}. They had to be less than 25 years old
and have a fuel economy, as rated by the window sticker, of 18 miles a gallon or less.

The size of the rebate depended on the fuel economy of the replacement vehicle. Consumers were also
supposed to receive the scrap value of their trade-ins.

From the dealers' point of view, the program was a resounding success.

``Two hundred and fifty thousand vehicles in four weeks?'' Mr.~Wood said. ``One word comes out of my
mouth: Wow.''

As word spread unofficially on Thursday night, car dealers were suddenly unsure of what to tell
would-be buyers.

A Ford dealership in Paramus, N.J., did not know of the apparent suspension until a reporter called
seeking comment.

Other dealers said they had no idea what the status of the program was, or whether the deals that
they had already signed would be honored by the government. Some said they were notified by e-mail
message by fellow dealers.

The dealers' association, however, had been warning that the money would go quickly.

Under the program, a buyer who picked a car with a mileage improvement of more than four miles per
gallon but less than 10 were eligible for \$3,500; a buyer whose new vehicle was rated 10 miles per
gallon or better than the old one was eligible for \$4,500.

Until the cash-for-clunkers program began, the auto industry had been on track for annual sales of
about 10 million units, down from the peak of about 16 million units a year.

\section{In House, Freshman Democrats Make a Stand}

It took Representative Dan Maffei of New York two tries, \$4 million and the retirement of a
Republican incumbent\cite{incumbent} to win his House seat last year. After all that, he wants to
avoid becoming a one-term wonder because of an unpopular vote on health care.

``We can't afford to make enemies,'' said Mr.~Maffei, a former Congressional aide and a Democratic
member of the freshman class that played a role this week in slowing House consideration of a health
insurance overhaul, upending plans by Speaker Nancy Pelosi to meet President Obama's goal of
approving a bill before the August recess.

Of 35 first-term Democrats, Mr.~Maffei, who represents the Syracuse area, and 25 others occupy
suburban, small-town and rural seats they took out of Republican hands, seats Republicans are eager
to take back. As a result, junior Democrats want to be sure the emerging health care plan is one
they can embrace, particularly after they have already had to cast a difficult vote on
climate-change legislation.

``As long as I feel comfortable that I had the opportunity to address issues that I have raised, I
am willing to take the tough vote,'' said Representative Frank Kratovil Jr., Democrat of Maryland.
``What I am not willing to do is follow the lead simply to follow the lead.''

As the House prepared to leave town until after Labor Day, the health bill was taking on some of the
ideological hue of House freshmen, many of whom represent districts in Southern and Western states
that were previously out of reach of Democrats, far from the urban centers that have long been the
party's base.

But the changes in the health care legislation have already provoked a backlash from the party's
progressive wing, setting up an August struggle not only with Republicans but also among Democrats
for the overall direction of the health plan.

Working somewhat as a bloc, and also through other groups like the fiscally conservative Blue Dogs,
the freshmen helped not only to postpone a floor vote but also to spur concessions on cost and
regional disparities. In addition, in response to freshman demands, party leaders are talking about
raising the threshold for any surtaxes that could hit small businesses.

``On issues where we agree, we are not afraid to speak out together and flex our muscles,'' said
Representative Jared Polis, Democrat of Colorado and an author of a letter to Ms.~Pelosi challenging
the proposed surtax. ``And one of those issues in the context of the health care debate was speaking
out against tax increases we saw as too expensive.''

Typically, Congressional freshmen are more seen than heard. But the political precariousness and the
sheer size of the freshman class -- along with the sophomores elected in 2006, they represent the
difference between a Democratic majority and minority -- give them more sway than usual.

``They are influential,'' said Representative Steny H.~Hoyer of Maryland, the majority leader, who
helps tend to the care and feeding of the freshmen. ``One of the concerns of the freshmen was to
give them time, so they are getting time.''

``It gives them an opportunity to explain to their constituents their views of the bill, and
secondly, they have more time to figure out what changes, if any, they need to effect,'' said
Mr.~Hoyer, who is urging the freshmen to take the offensive on health care and emphasize the
potential benefits of the overhaul.

In an indication of the importance attached to the class, Speaker Pelosi meets weekly with the
freshmen for a Wednesday breakfast. Representative Chris Van Hollen of Maryland, the chairman of the
Democratic Congressional Campaign Committee, is sending them home for August with a briefing package
that offers recommendations on how to sell the health care plan.

To hear Republicans talk, the freshmen are going to need each recommendation.

``I think it's safe to say that over the August recess, as more Americans learn more about their
plan, they're likely to have a very, very hot summer,'' Representative John A.~Boehner of Ohio, the
House Republican leader, said Thursday.

Republican strategists are planning a monthlong assault on Democrats deemed to be vulnerable, and
many freshmen are at the top of the list.

Not only is the first re-election campaign considered the best opportunity to beat an incumbent, but
freshmen from formerly Republican districts in states like Virginia, Ohio and Pennsylvania have
already been under the gun over the climate-change legislation. Thirteen of the 35 districts now
held by Democratic freshmen were carried by the Republican candidate for president in 2008.

Republicans say they do not believe that any modifications to the health plan will be enough to
protect Democrats.

``If Democrats think that their situation is improving on health care, then they haven't been
reading the polls,'' said Ken Spain, a spokesman for the National Republican Congressional
Committee.

Mr.~Spain called the emerging bill ``political malpractice for a Democrat in even the safest of
seats.''

But some of the freshmen say they are willing to take the heat as long as they can have some say in
shaping the legislation.

``We wanted to tackle the problems that neither party has had the guts to face for the last few
decades, and we are ready to take them on even if that is unpopular,'' said Representative Tom
Perriello, a Virginia freshman who narrowly defeated a Republican incumbent. ``But only if we feel
what we are doing is actually solving the problem and not just patching it up.''

Mr.~Perriello said August could amount to a monthlong town-hall-style meeting to allow Democrats to
explore health care with residents of their districts. He and other freshmen said that as recent
survivors of competitive races, they had a good sense of how to gauge public opinion and act
accordingly.

``We all just finished very tough races,'' said Representative Jim Himes of Connecticut. ``Very few
of us are in safe Democratic districts. So we listen very, very hard.''

\section{Wal-Mart Contests Use of Name by Canadian Union}

\lettrine{W}{al-Mart}'s feelings about unions in its stores are well known. In Quebec, the company
does not want a union even breathing its name.

Wal-Mart Canada has asked a Quebec court to stop a union Web site from using the company's name,
color schemes, a variation of its logo, a parody of its slogan and even photographs of people
wearing blue vests.

Wal-Mart says that it is seeking the injunction only to protect its trademarks. But the United Food
and Commercial Workers Canada, which represents employees at the retailer's only North American
store with a union contract, said on Wednesday that the company was trying to thwart its organizing
efforts.

``This injunction request is an over-the-top assault on freedom of speech and on our ability to
effectively communicate with Wal-Mart workers,'' Wayne Hanley, the union's president for Canada said
in a statement. ``Wal-Mart's response to the success of www.walmartworkerscanada.ca is just another
outrageous\cite{outrageous} example of how the largest retailer in the history of the world will use
its bottomless legal budget to manipulate the collective bargaining process and do just about
anything to discourage its associates from joining the union.''

Andrew Pelletier, a spokesman for Wal-Mart Canada, said a recent redesign of the union's Web site
``in a way we feel infringes\cite{infringe} on our trademark'' led the company to seek an
injunction.

``If you don't take steps to protect your trademark, the trademark becomes vulnerable,''
Mr.~Pelletier said from Mississauga, Ontario, where the Canadian unit is based. ``Anyone who knows
Wal-Mart knows we are an advocate of free speech and open communication.''

Derek Johnstone, a spokesman for the union, said the Web site was set up in 2003. It contents are
directed at Wal-Mart employees, he said, and are a mixture of negative articles about the company's
labor practices and advice about legal rights under Canadian labor laws.

Next to a large union logo, the Web site features a variation of Wal-Mart's current symbol in which
a circle of stick people holding hands replaces a ring of sticks. Next to a photograph of three
people wearing blue vests is the slogan ``Get Respect. Live Better!,'' a play on Wal-Mart's ``Save
Money. Live Better.''

On Wednesday, the union added banners reading ``Under the Threat of Censorship'' to the Web site.

Ariel Katz, a professor of intellectual property and Internet law at the University of Toronto, said
that Canadian courts had often sided with companies in trademark and copyright lawsuits against
unions.

``It's a very serious issue, the extent to which trademark law can be used to silence speech,'' he
said. ``The basic question of trademark is whether people who view the Web site would think it is
sponsored by Wal-Mart.''

In one prominent case, the Federal Court of Canada ruled in 1996 that the Canadian Auto Workers
union had violated Michelin's copyright with fliers and posters depicting Bibendum, the Michelin
tire man, about to stomp\cite{stomp} on a workers.

But Professor Katz said a recent Supreme Court of Canada decision about photocopying may strengthen
the union's case against Wal-Mart.

While both the Canadian branches of the union and Wal-Mart are based in Ontario, the company has
filed its complaint with the Quebec Superior Court. Mr.~Pelletier said that was partly because the
Web site was in French as well as English ``and in our experience Quebec courts have the most
expedient processes in resolving lawsuits of this nature.''

Quebec is also the province where the union has been most successful.

In addition to the Wal-Mart store near Montreal, which received a court-imposed contract this
spring, the union has organized at another Wal-Mart store and a Wal-Mart tire center. A court is
also hearing a union challenge to Wal-Mart's decision to close a store in the province after a
successful union drive.

This week, the United Food and Commercial Workers added its own variation of the disputed Michelin
ad: an animated video on their Web site in which a man and a woman are crushed by a giant boot
labeled ``Walmart's Army of Lawyers.''

\section{China to Try Suspects Held After Riots}

\lettrine{C}{hina} will begin trials in the next few weeks for suspects it accuses of playing a role
in the deadly riots that shook the capital of the Xinjiang region in early July, state media outlets
reported Friday.

The English-language China Daily newspaper said officials were organizing special tribunals to weigh
the fate of ``a small number'' of the 1,400 people who have been detained, most of them Uighurs, a
Turkic-speaking Muslim minority whom security forces have blamed for much of the killing.

Earlier this week, the authorities arrested an additional 253 suspects, many through tips provided
by residents of Urumqi, the regional capital where the violence took place. On Thursday, the
authorities published the photographs of an additional 15 people, all but one of them Uighur, who
they say had a hand in the unrest. Those who provide information leading to an arrest can collect as
much as \$7,350 in reward money.

``The police urged the suspects to turn themselves in,'' China Daily wrote, quoting an unidentified
law enforcement official. ``Those who do so within 10 days will be dealt with leniently, while
others will be punished severely.''

In the days after the riots, the head of the Communist Party in Urumqi was blunt about what awaits
those convicted of the most serious offenses. ``To those who have committed crimes with cruel means,
we will execute them,'' said the official, Li Zhi.

The riots, the worst outbreak of ethnic strife in China's recent history, began July 5 after
protests over the deaths of Uighur factory workers in another part of China turned into a murderous
rampage\cite{rampage}. The violence, which lasted three days, claimed 197 lives, most of them Han
Chinese beaten to death on the streets, according to the government. The Han are the dominant ethnic
group in China.

Uighur advocates overseas, however, insist that the official death toll undercounts the number of
Uighurs killed by the paramilitary police and during revenge attacks by the Han that followed the
initial rioting.

China has accused outsiders of instigating the unrest, heaping most of the blame on Rebiya Kadeer,
the 62-year-old leader of the World Uighur Congress, which advocates self-determination for China's
Uighurs. They say Ms.~Kadeer, a businesswoman who spent years in a Chinese jail before going into
exile, organized the killings from her home in Washington.

In recent weeks Ms.~Kadeer has been on an aggressive campaign to convince the world that her people
are the primary victims of the rioting. During a visit to Japan on Wednesday, she told reporters
that 10,000 people had disappeared overnight in the days following the unrest. ``Where did they
go?'' she asked. ``Were they all killed or sent somewhere? The Chinese government should disclose
what happened to them.''

Her claims have infuriated China, with one official in Xinjiang describing her remarks as
``completely fabricated\cite{fabricate}.'' Ms.~Kadeer says she cannot reveal the source of her
information because to do so would endanger those who provided it.

If the trials that followed the 2008 riots in Tibet are any guide, the court hearings in Xinjiang
will be swift. According to China Daily, the accused will be appointed lawyers who have ``received
special training,'' as have the judges who will preside over the cases. Each trial will be heard by
a panel of three or seven judges, and the majority opinion will prevail\cite{prevail}.

Human rights groups, however, say they have little confidence the tribunals will be fair. They
expect the proceedings to be closed to the public, as are most trials in China, and they note that
the defendants will not have lawyers of their own choosing.

``Without independent legal counsel, you don't have any clue as to what evidence has been collected
and through what means,'' said Renee Xia, international director of Chinese Human Rights Defenders,
which is based in Hong Kong. ``Were they tortured or coerced to confess? Trials can be speedy, but
it doesn't mean they will be fair.''

\section{China Could Use Some Honest Talk About Race}

\lettrine{W}{hen} the city of Detroit erupted in some of the worst rioting in American history over
a five-day period in July 1967, the Johnson administration responded by naming a high-level
commission to investigate the incident and more generally to weigh in on the troubled issue of race
relations in the United States.

The panel, known as the Kerner Commission, undertook to plumb\cite{plumb} three key questions:
``What happened? Why did it happen? What can be done to prevent it from happening again?'' And in a
simple but powerful phrase that helped define the era, it concluded that ``Our nation is moving
toward two societies, one black, one white -- separate and unequal.''

The Kerner Commission did not introduce the concept of minority civil rights in the United States.
That movement began to gain critical mass in the 1950s, through direct citizen action by people like
Rosa Parks, who refused to surrender her bus seat to a white person in Montgomery, Alabama, and was
arrested and tried for her defiance of racism, sparking a 381-day boycott of public transportation
by blacks in the city.

What the Kerner Commission did, rather, was signal recognition at the highest levels of American
society that the United States had major racial problems, along with civil rights deficiencies that
seriously marred our democracy. And recent events in Cambridge, Massachusetts, where Henry Louis
Gates, Jr., the country's most prominent black academic, was suspected of burglary and arrested in
his own house, demonstrate that questions of civil rights in America still preoccupy us.

This is the second year in a row of severe turmoil\cite{turmoil} in western China, following the
uprising that swept Tibetan areas in March of 2008. The events of recent weeks in China's Xinjiang
region, where were nearly 200 people died during unrest and a dozen members of the predominantly
Muslim Uighur minority were killed by police (according to official figures), demonstrate if nothing
else how China desperately awaits its own civil rights moment.

The Kerner Commission's famous old questions would be a good place to start: What exactly happened
and why? And an open and honest Chinese conversation about race, ethnicity, religion and identity is
long overdue and would go a long way toward healing papered-over divisions that run deep in this
society.

The response of the system here so far, alas, has shown no such willingness. The official media,
operating in their mouthpiece of power mode, have rushed to certain conclusions about the events,
namely that the trouble was instigated by ``splittists,'' and that sinister\cite{sinister} foreign
forces were at work behind the rioting.

Openness and transparency about the events of Urumqi would be welcome but by themselves would only
constitute a first step, no more. China has made great, and often insufficiently acknowledged
strides away from totalitarianism in the last generation, but one area where the rigidities of the
past linger on is in the politics of ethnicity.

China clings to the fiction that areas where ethnic minorities have historically predominated,
places like Xinjiang and Tibet, with distinctive languages and cultures and
lingering\cite{lingering} memories of self-rule, are ``autonomous regions.'' This, even as these
areas are governed by local party leaderships appointed by Beijing and heavily dominated by members
of the country's Han majority. This, also, as Beijing floods these areas with Han economic migrants,
for the purpose of settling and securing China's rough western frontier, raising local living
standards and to assimilate the local people into the ways of the Han.

Although this effort lacks in candor\cite{candor} and transparency, not to mention the possibility
of meaningful input from or consent by the locals, it would be wrong to conclude it is entirely
undertaken out of bad faith. The materialists who rule China seem to genuinely believe that economic
development is the answer to almost every question, and their favorite statistic relating to
Xinjiang is the doubling of the region's economy between 2002 and 2008.

At best, this statistic is misleading, though. Most of the economic growth in Xinjiang is related to
the expansion of the petroleum sector, which is overwhelmingly dominated by Han. Indeed the unrest
there seems fueled in part by a sense of among Uighurs that they are losing ground economically to
the Han in their own homeland.

I interviewed a Uighur barber in Urumqi two years ago who complained that the newcomers form their
own social and business networks and often enjoy government support of one kind or another. This
man, who had been trained in petrochemical engineering in Russia, said he had been unable to find a
job in that booming sector. Han, he said, hire Han.

A new study, published in the China Quarterly by Brenda L.~Schuster, reveals other gaps in the
economic statistics. ``In life expectancy, infant mortality, maternal mortality and morbidity,
Uyghur people are much worse off than Han,'' the report reads. It then speaks of how ``group
specific psychological stress and the socio-economic and demographic changes of the past 60 years
could be major factors.''

Many African-Americans, particularly in urban areas, where health indicators persistently lag behind
those of the general population, even at similar income levels, would readily recognize such
stresses. China, meanwhile, clings to the old Maoist-era fable\cite{fable} of the country as one big
happy ethnic family, even as it labors hard in Xinjiang to discourage Islamic worship and otherwise
dilute\cite{dilute} Uighur culture.

Two years of violence may not yet make a trend, but this myth has just become a lot harder to
sustain, even among China's Han majority, who may yet come to appreciate that respect for
differences rather than forced assimilation is the better recipe for harmony.

\section{Now on YouTube, Local News}

\lettrine{W}{ith} its ability to collect articles and sell advertisements against them, Google has
already become a huge force in the news business -- and the scourge\cite{scourge} of many
newspapers. Now its subsidiary YouTube wants to do the same thing to local television.

YouTube, which already boasts of being ``the biggest news platform in the world,'' has created a
News Near You feature that senses a user's location and serves up a list of relevant videos. In
time, it could essentially engineer a local newscast on the fly. It is already distributing hometown
video from dozens of sources, and it wants to add thousands more.

YouTube says it is helping TV stations and its other partners by creating a new -- but so far not
fiscally significant -- source of revenue.

But news media companies may have reasons to be wary. Few TV stations have figured out how replicate
profits on the Internet. YouTube can easily act as another competitor.

So for now, most of the YouTube videos near you come from nontraditional sources: radio stations,
newspapers, colleges and, in the case of a fledgling\cite{fledgling} San Francisco outfit called
VidSF, three friends who despise the local TV diet of fires and homicides.

``It really levels the playing field,'' said Kieran Farr, a founder of VidSF who covers the city's
culture and uploads his segments to YouTube.

News Near You, started in the spring, is only part of YouTube's push into news video. This summer,
the company invited the more than 25,000 news sources listed on Google News to become video
suppliers. The site is also promoting videos from ABC News, The Associated Press, Reuters and other
outlets.

This year, it began featuring breaking news videos -- including ones submitted by citizens in Iran,
where protests are being captured by cellphone users -- on its home page.

So far, the localized videos are no replacement for a print or TV diet of news. On Sunday, visitors
near Baltimore saw a news report about a teen assistance program; in Chicago they saw a WGN-TV
feature about street performers; and in Los Angeles, they saw a review of an electric motorcycle
produced by The Los Angeles Times. Producers often count the views in hundreds, not thousands.

To date, nearly 200 news outlets have signed up with YouTube to post news packages and split the
revenue from the advertisements that appear with them. In addition, Google searches now show YouTube
videos alongside news articles, helping the videos reach a wider audience.

YouTube's sheer breadth -- it is visited by 100 million Americans each month -- makes it a powerful
force for promotion, as well as a potential threat to entrenched media companies. And those
companies already have more than enough to worry about: much of the local media marketplace has
collapsed in recent years as classified ads have moved online, automakers have
curtailed\cite{curtail} ad spending and news and entertainment options have proliferated.

YouTube, meanwhile, is still trying to turn a profit nearly three years after it was acquired by
Google. Because copyright concerns prevent it from placing ads on amateur videos, it has striven to
sign up professional partners to seed the site with ad-friendly content. News is one obvious option.

``Google can only gain by splitting revenue with people who have feet on the street in local
markets,'' said Terry Heaton, a senior vice president at AR\&D, a company that advises locally
focused media organizations.

Google said in June that it was pleased with YouTube's trajectory and indicated that it expected the
site to be profitable in the not-too-distant future, but did not specify when.

While YouTube can gain by adding local video, it remains to be seen whether established news outlets
will benefit. Google's scraping\cite{scraping} of print headlines and links has led some to assign
blame to the company for the financial struggles of newspapers. The chief executive of Dow Jones
recently called Google a ``digital vampire'' that was ``sucking the blood'' from newspapers by
harvesting their free articles.

What YouTube is doing is somewhat different. It is not sending digital spiders around the Web to
collect videos automatically; instead, it is asking news outlets to sign up as partners and
promising a wider audience for their material.

YouTube's push to organize local news video began in earnest in the spring when the News Near You
module was introduced. The module uses the Internet address of a visitor's computer to determine the
user's location and whether any partners are located within a 100-mile radius. If so, seven days of
local videos are displayed.

But in many places, namely urban markets, 100 miles can hardly be counted as a local area; Steve
Grove, the head of news and politics for YouTube, said, ``we'll get a smaller radius as we bring on
more partners.''

Mr.~Grove said about 5 percent of users who see the News Near You module watch at least one local
news video, a rate that YouTube sees as encouraging.

``The relevancy factor kicks into gear,'' he said. ``Suddenly these videos actually matter to you
because they're about your neighborhood or where you live.''

YouTube's arms-wide-open approach forces stations to judge whether YouTube is a friend or a
foe\cite{foe}, echoing a question that newspapers have grappled with for years. (Some have deemed
Google a ``frenemy.'')

``For stations, there won't be significant revenue from YouTube in the short term,'' said Andy
Plesser, the editor of the online video site Beet.tv. That is partly because few of the videos reach
an audience of millions, he said.

He also suggested that the local news feature faces resistance from station owners. ``Many simply
don't see the value of being on YouTube,'' he said.

In New York, the cable news channel NY1, for instance, said it was concentrating on drawing visitors
to its own site, rather than sharing videos elsewhere.

``It's an old conundrum\cite{conundrum},'' Marc Nathanson, an executive producer of NY1.com, said.
``When you have valuable content, do you make people come to you for it? Or do you put it out there
to enhance the brand and hope that users find their way back to you for more?''

He guessed that NY1 would offer some videos on sites like YouTube in the future.

Meanwhile, a new breed of local news broadcasters -- including ones without broadcasting licenses,
the traditional barrier to entry in local markets -- is emerging online.

``Radio stations, newspapers, universities, advocacy organizations, churches and other local groups,
and individuals are becoming news producers,'' Mr.~Plesser said.

Mr.~Heaton said stations should treat YouTube as a marketing machine for their old-fashioned
television newscasts.

``As long as Google wants to pay for the bandwidth'' to host videos, he said, ``let them.''

Mr.~Grove said YouTube had not met ``a ton of resistance'' from news outlets and the main reason
there are so few participants is that the initiative is new.

Mr.~Grove has stumbled across some local members, including The Dallas Morning News and The
Cincinnati Enquirer, that he had not even known were on the site. (Any Internet user can create a
YouTube account and start uploading videos free.) The New York Times created a YouTube channel in
2006 and has posted more than 1,300 videos on it, essentially highlighting some of the content on
its own Web site.

National news, too, is being curated by the site. The ``top headlines'' section of the site collects
videos from prominent TV sources and crossreferences them with the trends from Google News, then
produces the equivalent of a 30-minute newscast for users to watch, either in full or from one
segment to another.

In the future, more of the News Near You could come from people who do not report the news for a
living. As the protests in Iran continue to demonstrate, citizens are able to provide much of the
spot video from breaking news, even though they may lack the objectivity of professionals.

The new iPhone includes a video recording capability with a ``send to YouTube'' button, suddenly
making it simple and fast to upload clips.

Rachel Sterne, the founder and chief executive of the citizen journalism site GroundReport, said the
feature ``trains laymen to be reporters.'' And YouTube says it is developing tools to automatically
spotlight those citizen videos as they come in.

\section{Out of Search Business, Yahoo Shifts Its Focus}

\lettrine{C}{arol} A.~Bartz, chief executive of Yahoo, has been hobbled\cite{hobble}.

Three weeks ago, doctors gave her a new left knee, made of titanium and plastic. As a result, she is
limping around Yahoo's headquarters here, occasionally standing to hold a chair and stretch her leg
while a bottle of Percocet sits at the ready on her desk.

Yahoo also underwent invasive surgery recently, selling its search business to Microsoft for an
initial 88 percent share of search revenue in a 10-year deal. Since that move last week, Yahoo too
has limped along, with disappointed investors asking whether Ms.~Bartz could have won more favorable
terms and with Yahoo shares tumbling more than 15 percent.

In a lengthy interview on Friday before departing for a vacation, Ms.~Bartz said she sold the search
business because Yahoo could no longer continue to match the level of investment Google and
Microsoft were making in searching, one of the Web's most lucrative\cite{lucrative} and
technologically complex businesses.

While reducing the marketing and infrastructure costs associated with search, the deal will also
provide money that Yahoo can use to bolster\cite{bolster} other businesses. Ms.~Bartz plans to
invest the money in Yahoo's display ad, content and mobile services technology.

``My first reaction when I got here was that I wouldn't even do a search deal,'' she said, ``until I
looked at our expense structure and our actual options and looked at what our prime job was, which
is to grow audience.''

Yahoo will lose some of its most talented engineers to Microsoft and as many as 400 employees
through layoffs. The deal also undercuts years of investment around search technology. By selling
the technological crown jewels, the company may lose some of its high-tech credibility among
employees and others in Silicon Valley, as well as among customers.

But the core of Yahoo remains intact, according to Ms.~Bartz. ``We haven't
eviscerated\cite{eviscerate} the company,'' she said.

Given Yahoo's battering by investors, Ms.~Bartz lamented that Yahoo and Microsoft failed to explain
the relationship better to Wall Street. She blamed herself for a comment she made several weeks ago,
at an industry conference, that Microsoft would have to pay Yahoo ``boatloads of cash'' to win its
search business. That statement, she said, helped to solidify an expectation that Yahoo would
receive \$1 billion or more upfront as part of the deal.

``I made a mistake. I was never interested in doing it for upfront money. That doesn't help me
operate a business,'' Ms.~Bartz said. She noted that such a payment would have had significant tax
consequences while contributing only \$3 million in annual interest to Yahoo's bottom line.

Wall Street had become enamored\cite{enamor} with the previous, more immediately lucrative proposals
between Microsoft and Yahoo. Last year, trying to improve its search business and better compete
with its archrival Google, Microsoft offered \$46 billion to buy all of Yahoo. Analysts estimate
that the new deal -- involving what many people saw as Yahoo's most important asset -- is worth only
around \$4 billion to \$5 billion.

``It's rather like getting a Picasso and saying, 'You know, the canvas costs \$200, the paint cost
\$300, so we'll sell it to you for \$500,' '' said Jeffrey Lindsay, an analyst at Sanford
C.~Bernstein. ``I've never seen investors so angry.''

Microsoft's chief, Steven A.~Ballmer, has gone to unusual lengths to defend the deal on Yahoo's
behalf. Last week, he told a group of investors gathered at Microsoft's headquarters that it was
``sort of unbelievable'' that Yahoo got to keep such a high percentage of the ad sales while
spending nothing on the underlying infrastructure to run the search operation.

And Ms.~Bartz argues that the more generous offers came during far different economic times. In
addition, she said, analysts have underestimated the two and half years it will take to integrate
the two companies' search and advertising systems. The companies also face the daunting\cite{daunt}
process of getting approval from regulatory agencies that are now closely scrutinizing deals on
antitrust grounds.

Ms.~Bartz said that she avoided Mr.~Ballmer's immediate advances after taking the helm of Yahoo.
``He called my first day,'' she said. ``I told him: 'Go away. I haven't even found the bathroom.' ''

When Ms.~Bartz finally began negotiations, Mr.~Ballmer presented two options: a large upfront
payment or a higher percentage of revenue tied to ads sold on Yahoo Web sites. ``He wasn't going to
pay twice,'' Ms.~Bartz said.

Ultimately, Mr.~Ballmer, who was unavailable to comment, and Ms.~Bartz ended up on the phone
debating minute details, like how fast things like search results and ads would appear, at one point
haggling over a margin of 100 milliseconds.

Ms.~Bartz also fought to play down Microsoft's Bing brand on the Yahoo search results page, with
just a bit of text instead of a logo, and she pushed Mr.~Ballmer hard to make sure that Yahoo
retained the right to sell its own display ads to big advertisers like Procter \& Gamble and
Unilever.

``I had to get Steve through the, 'we get to sell part,' '' she said. Meanwhile, Ms.~Bartz has
plenty of other work to do at Yahoo. She has picked up on the dysfunction throughout the company
during her eight months on the job.

For one, Yahoo has not poured enough money into creating a top-class technology infrastructure
capable of delivering a variety of services to consumers. ``We are actually behind in investing,''
she said. ``We should have invested more.''

Ms.~Bartz also complained about a lack of balance among the different product groups. Motivated
employees in the company's sports section have turned it into a must-visit Web destination by
breaking news stories and providing witty commentary. Other sections have languished.

This is often a result of what Ms.~Bartz sees as insufficient communication and barriers among the
company's various American and international properties over sharing videos and other content. ``We
had a fight for 18 months here on what a video button should look like,'' Ms.~Bartz said.

Spending money to fix these problems made more sense to Ms.~Bartz than hemorrhaging cash in a bid to
keep pace with Google and Microsoft.

Microsoft has backed its new search engine, Bing, with billions of dollars in infrastructure and
marketing investments. So far, the search engine has been well received. While encouraged by some of
Bing's innovations, Ms.~Bartz figures that Google will most likely match them soon. ``That is the
arms race I don't want to be in,'' she said.

But some investors fear that Yahoo may be trading one of the most profitable businesses of the 21st
century -- search -- for a deflating, hypercompetitive media business that appeared to peak in the
20th century.

``I don't think the surviving pieces are very interesting,'' said Michael Lippert, manager of the
Baron iOpportunity fund, which acquired Yahoo shares earlier this year on the anticipation of a
lucrative search partnership. Mr.~Lippert was disappointed by the terms of the deal but did not make
out too badly; his fund also owns Microsoft stock, which rose slightly last week.

``Ballmer got himself a fantastic deal,'' he said.

\section{Newspaper Apologizes for Inventing a Scandal}

\lettrine{``I}{t's} a Scandal,'' an exceptionally large front-page headline in The Saint John
Telegraph-Journal proclaimed last month above an article about what the prime minister of Canada
did, or did not, do with a communion\cite{communion} wafer\cite{wafer}.

Last week, the wafer issue resurfaced on the front page of the newspaper. But this time as more of
an embarrassment. In a lengthy correction, the New Brunswick newspaper said that ``there was no
credible support'' for the article's claim that Prime Minister Stephen Harper had pocketed the wafer
during the funeral of a former governor-general or that a Roman Catholic priest was demanding an
explanation from his office.

Unusually, the newspaper apologized not just to Mr.~Harper but also to the two reporters whose names
appeared on the article, adding that the errors were introduced ``in the editing process'' and were
included ``without the knowledge of the reporters and without any credible support for them.''

The fallout did not stop there. Shortly after the article appeared, the publisher, a member of New
Brunswick's most powerful family whose holdings include The Telegraph-Journal, and the newspaper's
editor were removed from their jobs.

In early July, Mr.~Harper, who is a Protestant, attended the funeral of Rom\'eo LeBlanc, a former
governor-general, at a Roman Catholic Church in New Brunswick. In about three seconds of television
news footage he can be seen receiving a wafer, or ``host,'' from an archbishop\cite{archbishop}.
Instead of immediately placing it in his mouth, as is customary for Catholics, Mr.~Harper dropped
the hand holding the wafer to his side. What happened to the wafer after that was not captured on
camera.

The Telegraph-Journal, however, wrote several days after the funeral that ``the prime minister
slipped the thin wafer that Catholics call 'the host' into his jacket pocket.'' Msgr. Brian
Henneberry, the vicar general and chancellor of the Diocese of Saint John, was quoted as saying that
if Mr.~Harper had pocketed the wafer, ``it's worse than a faux pas, it's a scandal from the Catholic
point of view.''

The article made it clear, however, that he was not at the funeral, had not seen the television
footage and did not have any direct knowledge of the wafer's fate. (Monsignor Henneberry did not
respond to requests for comment.)

The article was picked up by the national news media and followed Mr.~Harper to the G-8 summit
meeting last month in Italy.

At a news conference there, Mr.~Harper insisted that he ate the wafer, adding, ``this is a low
moment in journalism.''

Annie Smith, a spokeswoman for the Irving-owned media holding company Brunswick News, said that no
one at the newspaper would comment on the article or the retraction. She confirmed that James
C.~Irving, the publisher, and Shawna Richer, the editor and a former reporter for The Globe and Mail
who once contributed sports articles to The New York Times, had been removed from their positions.
Ms.~Smith declined to say, however, whether either was still an employee.

Mr.~Irving could not be reached. Ms.~Richer did not respond to requests for an interview.

The newspaper's retraction, however, did not end the debate. Some commentators in Atlantic Canada
argue that the apology raises almost as many questions as it answers. And the affair has again drawn
attention to the Irving family's extraordinary control over not just New Brunswick's news media (it
owns all of the daily papers in the province and controls the weekly market) but also its political
and economic power (the Irvings' businesses include oil refineries, forestry, retailing, paper pulp
production and food processing.)

That combination that has prompted supporters of Mr.~Harper, a Conservative, to suggest in blogs and
on Twitter that the Irvings planted the story to improve the fortunes of the opposition Liberal
Party. At the same time, some Liberal bloggers claim that the retraction was a bid by the family to
curry favor with Mr.~Harper. Neither side has offered much in the way of collaboration.

Among those baffled by it all is Parker Donham, a prominent journalist in Atlantic Canada who now
works as a communications consultant. On his blog, Contrarian, Mr.~Donham said that because the
article was entirely about Mr.~Harper's ``nontransubstantiation'' of the wafer, it is difficult to
determine what was added during editing.

``So which part of the story is true, and which fabricated?'' Mr.~Donham wrote. ``That's the central
problem with the Irvings' highly selective account of the events: we can't tell.''

Mr.~Harper, at least, now appears satisfied.

``We appreciate and accept the apology from The Telegraph-Journal,'' Kory Teneycke, his spokesman,
wrote in an e-mail message on Friday.

\section{New York Post Sends In the Clown}

\lettrine{T}{he} New York Post made the phrase ``media circus'' literal when it hired a clown to
roam the corridors of the New York State Capitol.

Sending in the clown, a playful attempt by The Post to satirize\cite{satirize} the political power
struggle that paralyzed the government for more than a month, was part of several days of
circus-themed coverage in June that offered readers some of the paper's more creative headline
writing and photo illustrations. (Think clown faces superimposed on senators' bodies and headlines
that used words like ``bozo'' liberally.)

But the clown, Connie Adams, was real. Ms.~Adams, a retired New York City high school physical
education teacher who lives in Yonkers, typically performs at children's birthday parties and
corporate picnics. In an interview, Ms.~Adams said that she almost told The Post no when she was
asked to travel to Albany for a day so a photographer could take pictures of her standing next to
politicians.

She was reluctant to cancel her standing Wednesday night gig\cite{gig} at a local Burger King and
make the two-and-a-half-hour drive to the capital. But she said she was grateful for her 15 minutes
of fame, which earned her a mention on National Public Radio, as well as numerous color photos in
The Post and a video on its Web site.

``I'm the envy of all the clowns,'' Ms.~Adams said with a laugh, adding that people had written to
her through her Web site, sunnytheclown.net, asking her to send them autographed copies of the June
11 edition of The Post, which featured a picture of her being ordered out of the Senate lobby by a
sergeant at arms.

She said the publicity had not generated any new business for her. But she is recognized more
frequently. On a recent trip to the grocery store, she said she caught a stranger's eye. She had
just finished a gig and was still wearing her clown attire.

``I don't generally like to go out when I'm in clown,'' she said. ``You have to stay in character,
and it's really hard. But some guy turns to me and says: 'Sunny the Clown! I just saw you in the
paper!' ''

But she was not responsible for the Post op-ed that appeared under the byline ``Sunny the Clown.''

``Oh, I didn't write that,'' she said, sounding somewhat disappointed. The writing was left to a
reporter in New York, she said, who read a draft of the article to Sunny before it went to print.

``I thought it was pretty funny,'' she said. ``My friends did, too. And I had to say, 'It wasn't me,
O.K.?' ''

\section{With Enough Soldiers, the Army Is Looking for a Few Good Officers}

\lettrine{A}{fter} three straight years of growth -- helped in part by a sagging job market -- the
Army is starting to find itself a bit bottom heavy, swollen with young recruits but short on
officers to lead them. Starting Monday, it is hoping a fresh promotional effort can reverse that.

For the first time in its history, the Army is introducing an advertising campaign to recruit
officers. The ads in many ways resemble the force's mainstream recruitment effort -- camouflaged
soldiers carrying big guns and standing at attention, with patriotic music as the soundtrack -- but
have been tweaked to appeal to achievement-oriented college graduates who could qualify for one of
its officer training programs.

For example, two of the TV commercials could be mistaken for ads from I.B.M. or Accenture until the
Army's signature music chimes in about halfway through. Another shows pictures of famous generals
like George Washington, Douglas MacArthur and Colin Powell, while a voice-over says, ``Officers in
the U.S.~Army can rise to any challenge. Can you?''

The goals are to attract ambitious young Americans who might normally consider the Army beneath
their career objectives and give the Army a jolt of much-needed creative leadership.

``It's a tough environment out there,'' said Lt.~Gen. Benjamin C.~Freakley, head of the Army
Accessions Command, which oversees recruiting. ``It's no longer where the enemy lines up on one side
of the field and the coalition lines up on the other side and the referee blows the whistle. It's a
very complicated battlefield to figure out, and there are no referees.

``It is a different era, and it requires a different kind of thinker,'' he said.

The Army's growing pains are reflected in its recruitment statistics. The force has expanded by more
than 50,000 troops since 2005, reaching 544,000 at the end of 2008. And last month, the defense
secretary, Robert M.~Gates, announced his intention to add 22,000 troops in the next year. The
expansion of the Army has come in response to wars in Iraq and Afghanistan that taxed the
scaled-back force that emerged from the 1990s.

But some of that growth can be attributed to relaxed standards. Only 83 percent of new recruits in
2008 held a high school diploma, missing the Army's goal of 90 percent for the third
consecutive\cite{consecutive} year. A growing percentage of new recruits during that time have also
been scoring in the lowest acceptable range on the Army's vocational aptitude test.

So while the force is meeting or exceeding its goals in terms of pure numbers, it is falling short
in its search for men and women qualified to lead their peers.

``The economic downturn has made it really easy for the Army to recruit in recent years, and that
was a big surprise to people who thought a big shooting war would discourage people from signing up
in the first place,'' said Loren B.~Thompson, a defense analyst with the Lexington Institute, a
research organization. ``But because of changes in threats and shifts in domestic demographics, the
military is looking for a different type of officer today than what it was seeking in the past, and
that requires special recruiting efforts.''

Four TV commercials will provide the public face of the campaign, all of which begin running Monday.
The two that slowly reveal themselves to be Army ads tell the story of high-ranking corporate
executives who were once Army officers: Joseph DePinto, chief executive of 7-Eleven; and Otto
Padron, a senior vice president at Univision.

``We learned through research that when the `Army Strong' music came on in that first second of the
commercial, those achievement-oriented students would reject it because they would think it was not
for them,'' said George Dewey, executive creative director at McCann Erickson, the Army's creative
agency. ``When these kids find out that the Army has produced these superstars, it's an `aha' moment
for them.''

That is also the thinking behind an ad that features George Washington and Buzz Aldrin. The final
ad, which more closely resembles a typical Army commercial, features a current officer explaining
his decision to join the R.O.T.C.

The TV ads will mostly run during the same programs as the mainstream recruitment commercials, like
``CSI: Miami,'' ``Law and Order'' and sports programming. Web ads will run on ESPN.com and other
major sites, but also destinations like Stack.com, a site for achievement-oriented athletes.

McCann Erickson is part of McCann Worldgroup, a collection of creative, media, digital and public
relations agencies, each of which has a hand in the wide-ranging campaign. Worldgroup, which is part
of the Interpublic Group of Companies, took over the Army's advertising account in 2006, when it
introduced the ``Army Strong'' tag line, replacing ``Army of One.''

The online centerpiece of the campaign is a microsite, GoArmy.com/officer, where potential
candidates can answer a series of questions that determines which of the four paths to becoming an
officer is right for them. Those paths are Army R.O.T.C., the Military Academy at West Point, direct
commission and officer candidate school.

The site also contains videos, each about five minutes long, explaining more about each of the
paths. Potential candidates can use the site to download information and make contact with a
recruiter. The site was created by MRM, the digital arm of McCann Worldgroup.

The Army is also working with Major League Baseball to produce a program called ``Leaders of the
Diamond,'' a series of interviews with all-stars about leadership and dedication that will appear on
MLB.com. It is also hosting panel discussions at universities around the country this fall at which
students can talk to officers directly.

The goal of the campaign is not just to recruit officers now, said General Freakley, but also to
begin doing a better job of marketing the officer ``product'' to young Americans.

``If you think about it as brand or product management, we have this product within our brand that
gets no notoriety\cite{notoriety},'' he said. ``For those who just graduated college, now is the
time to become aware they can come to officer candidate school. We think the timing is right to get
the notion out.''

\section{Two Sides Take Health Debate Outside Capita}

\lettrine{W}{ith} Republicans mobilizing against the proposed health care overhaul, President Obama,
Congressional Democrats and leading advocacy groups are laying the groundwork for an August
offensive against the insurance industry as part of a coordinated campaign to sell the public on the
need for reform.

The effort will feature town-hall-style meetings by lawmakers and the president, including a swing
through Western states by Mr.~Obama, grass-roots lobbying efforts and a blitz of expensive
television advertising. It is intended to drive home the message that revamping\cite{revamp} the
health care system will protect consumers by ending unpopular insurance industry practices, like
refusing patients with pre-existing conditions.

``I think what we want to communicate is that this is going to give people who have insurance a
degree of security and stability, the protection that they don't have today against the sort of
mercurial\cite{mercurial} judgments of insurance bureaucrats,'' said David Axelrod, a senior adviser
to Mr.~Obama, adding, ``Our job is to help folks understand how this will help them.''

Revamping health care is the president's top legislative priority, and people on all sides of the
debate agree that August, when lawmakers leave Washington to take the pulse of constituents, will be
crucial to shaping public opinion. With Republicans making headway by casting the legislation as a
costly government takeover, Democrats have decided they must answer the question on the minds of
those now insured: ``What's in it for me?''

That has led to a campaign of increasingly harsh\cite{harsh} rhetoric\cite{rhetoric} against the
insurance industry, which says it favors an overhaul but is working to defeat Mr.~Obama's call for a
government-run insurance plan to compete against the private sector. On Friday, Speaker Nancy
Pelosi, Democrat of California, promised a ``drumbeat across America'' to counter what she termed a
``shock and awe\cite{awe}, carpet-bombing by the health insurance industry to perpetuate the status
quo.''

The tough talk, however, has risks. The industry trade group, America's Health Insurance Plans, is
urging members to confront Democrats at public meetings, and the rising tensions could make it
difficult for the president to keep insurers at the negotiating table.

The drumbeat will begin Monday, when Kathleen Sebelius, the health and human services secretary,
travels to Hartford to talk about what the White House now calls ``health insurance reform.''
Senator Christopher J.~Dodd of Connecticut, who disclosed Friday that he has prostate cancer and
pointedly reminded Americans that he was fortunate to have health coverage, will be among several
Democratic lawmakers present.

Also Monday, Senator Michael Bennet, Democrat of Colorado, will appear with doctors, nurses and
administrators at St. Joseph Hospital in Denver to discuss ``how insurance company procedures are
burdening our physicians, nurses and patients,'' a spokeswoman said. Throughout the recess,
Democratic lawmakers will hold similar events, coordinated with advertising by allied groups.

``We understand the future of health reform could hinge\cite{hinge} on how the conversation with the
American people goes in the next six weeks,'' said Representative Chris Van Hollen, Democrat of
Maryland and assistant to the speaker, who is coordinating the House effort.

Republicans understand that and will also be campaigning hard.

In the Senate, Republicans will meet this week to coordinate strategy, but some plans are already in
motion for public meetings and a blizzard of radio and television appearances. Senators Tom Coburn
of Oklahoma, a family practice doctor, and John Barrasso of Wyoming, an orthopedic surgeon, will
take their ``Senate Doctors Show,'' an Internet program, on the road to argue that the Democratic
plan will not improve care or control costs.

In the House, Representative Mike Pence of Indiana, chairman of the Republican Conference,
distributed a packet to colleagues on Friday urging them to argue that the Democrats' plan would
include ``more than \$800 billion in new tax hikes'' and ``harmful cuts'' to Medicare that would
``result in millions of seniors losing their health coverage.''

The Democrats are getting a lift from a little-known group of former Obama campaign operatives
called Unity '09 that has held weekly strategy meetings, away from the White House, to bring
together administration officials, labor unions, health advocacy groups and other backers of the
legislation. Mr.~Axelrod said he had attended as an ``infrequent visitor.''

The current message is an eight-point list of ``Health Insurance Consumer Protections'' the White
House Web site promises will ``bring you and your family peace of mind.'' Mr.~Obama picked up on the
theme last week, promising members of AARP that he would ``reform the insurance companies so they
can't take advantage of you.''

The hard line is a departure for the White House, which began its overhaul campaign by trying to win
over constituency\cite{constituency} groups -- doctors, pharmaceutical companies, hospitals and
insurers among them. Insurers played a leading role in killing a health care overhaul in the Clinton
administration, but Karen Ignagni, president of America's Health Insurance Plans, said at a White
House meeting in March that the industry would cooperate this time.

In an interview Friday, Ms.~Ignagni noted that the industry had endorsed many of the
administration's proposed changes, including ending the practice of refusing coverage for
pre-existing conditions, and said it would work with lawmakers to develop a bill that did not
include a public plan.

``The rhetoric that we are hearing is reminiscent of '93, '94, but we're on the 2009 playbook,'' she
said, adding, ``The inconvenient fact is that we support those reforms.''

Republicans say the new rhetoric reflects desperation among Democrats to find a message that will
resonate\cite{resonate}.

Broadly speaking, Mr.~Obama's goal is to extend coverage to the nation's 47 million uninsured while
also slowing the growth of health care spending. The president wanted the House and Senate to pass
legislation before the August recess, but the House broke on Friday without doing so and the Senate
will not complete a bill before recessing after this week.

Polls show that the public is growing uneasy; a New York Times/CBS News survey last week found that
while Mr.~Obama still has strong support for revamping health care, Americans are concerned that an
overhaul would reduce the quality of care, increase out-of-pocket costs and tax bills and limit
their options in choosing doctors.

``August is going to be a critical period for closing the deal,'' said Ron Pollack, director of the
advocacy group Families USA. The group is now in a strange-bedfellows advertising partnership with
the pharmaceutical\cite{pharmaceutical} industry. Their latest advertisement, a three-week, \$4
million campaign, updates the ``Harry and Louise'' advertisements that helped kill health
legislation in the 1990s.

In the new advertisement, called ``Get the Job Done,'' Louise tells Harry that Americans need good
coverage. To which Harry adds, ``Even if they have a pre-existing condition.''

\section{Hundreds More Chinese Are Rounded Up}

\lettrine{C}{hinese} security officials have detained an additional 319 people whom they suspect of
taking part in the unrest in the far western region of Xinjiang, the official media reported Sunday.

The arrests, according to the Xinhua news agency, came days after 253 other people were taken into
custody, most of them Uighurs, a Turkic-speaking, predominately Sunni Muslim minority that the
Chinese government has blamed for much of the killing last month in the city of Urumqi.

``Every suspect, without any exception, will be arrested and punished according to the law to root
out any hidden danger,'' Xinhua said. It added that many of those detained in recent days had been
turned in by people outraged by the violence, which claimed 197 lives, nearly two-thirds of the
victims Han Chinese. ``It aroused the feelings of local residents, and they had more initiative in
exposing those suspects and provided a huge amount of clues,'' Xinhua said.

The police have not released a tally of those being held, but that figure exceeds 2,000 if the
numbers published by the state news media in recent weeks are added up. Rebiya Kadeer, the leader of
the World Uighur Congress, an exile group, has claimed that 10,000 people have disappeared since the
July 5 riots. Chinese officials have rejected that figure as ``completely fabricated.''

Whatever the true number, the authorities have said that only a small number of those in detention
will stand trial this month. Those found guilty are likely to face harsh punishment, possibly
execution.

The government says the rioting was instigated and orchestrated from abroad by those seeking to
disrupt Xinjiang's ``ethnic harmony.'' Uighur exile groups, however, say the strife was set off by
China's repressive governance of Xinjiang and, on July 5, a heavy-handed police response to peaceful
demonstrations in Urumqi.

On Sunday, Xinhua also published what it said was an interview with the imam of one of Urumqi's
biggest mosques\cite{mosque}, who offered an account of how three Uighur men tried to
incite\cite{incite} congregants to kill in the name of Islam during afternoon prayers on July 13, a
week after the riots had subsided.

The imam, Abdushukur Rehmutura, said one of the men stood up among the congregants, took out a green
banner and shouted: ``Holy war! Holy war! Go with us.'' According to the Xinhua article, two other
men then revealed a bag of 20 knives, and when no one would join them, they ran into the street.
Minutes later, after ignoring police orders to drop their weapons, according to the account, the men
were shot and killed.

\section{Mobile Gives Manga a Lift}

\lettrine{E}{iko} Moori, a 24-year-old office worker, felt the pull of manga one day when she saw an
ad on her mobile phone screen. Because she grew up on a steady diet of manga, a form of Japanese
cartoon, the offer to view ``Mischievous Kiss,'' a comic love story she was devoted to as a
teenager, for just a few dimes an episode was irresistible.

She took the plunge: For a total of ¥1,600, or about \$17, she viewed the entire series, over 70
episodes in all.

``It brought back memories,'' she said, adding that she had once owned the whole series, which she
kept in her room. ``I sold all of them because they were crowding the space. But for just a few
clicks on my mobile, I was able to read them all again.''

For a variety of reasons, the mobile version has manga booming again. In the year ending in March,
Japanese manga publishers raked in ¥32.9 billion in revenue, up 43 percent from the previous year
and from next to nothing in 2003, when manga first became available by cellphone, according to
Impress R\&D, a research company in Tokyo, which published the data at the annual Tokyo
International Book Fair in July.

The soaring numbers are generating hope -- and some controversy -- in a venerable industry that has
yielded internationally popular titles like ``Dragon Ball,'' published by Shueisha, and ``Sailor
Moon,'' created by Kodansha.

With revenue for manga in magazine and book form having inched down for more than a decade -- to
¥448 billion in 2008 from the peak of ¥586 billion in 1995, according to the Research Institute for
Publication in Tokyo -- some analysts are calling the cellphone ``the savior'' of the form and a
necessary kick in the pants for a sagging industry.

``At present, cellphone comics is the only one roaring ahead in sales,'' said Shinichi Yoshizawa,
director of digital media business development at Kodansha, the largest Japanese publisher of books,
manga and magazines.

Publishers would not reveal publicly what percentage of their revenue came from mobile manga, but
privately, many said about 10 percent, with mobile sales fast gaining on paper sales. So far, it is
unclear whether the mobile sales will make up for lost revenue.

But readers are taking to the mobile version, largely because the user experience has been good and
the pay platform convenient.

Ms.~Moori, the office worker, said she liked the mobile experience because ``on a whim, I can decide
to read them on the train.'' Others cited the graphics, saying manga looked especially good lit up
against the dark screen of a cellphone.

Yusuke Nakabayashi, a media consultant at Nomura Research Institute, a unit of Nomura Securities,
said the mobile infrastructure in Japan was crucial. Third-generation networks, which currently
offer the best high-speed Internet experience, cover nearly 100 percent of Japan, and the phones
made here have crisp, high-resolution displays and software designed for viewing manga. The software
allows readers to zoom in and out and move sideways to reflect the dynamism of the sequences, he
said.

A Japanese company called Celsys creates the reader software, called book-surfing, which is
downloaded from the mobile carrier's Web portal and installed in the phone. For payment, the major
cellphone service providers, like NTT DoCoMo and KDDI, add charges for the mangas to subscribers'
phone bills.

But perhaps the most crucial driver behind the cellphone comic boom is its attractiveness to women.
You can tell just by looking at the best-selling titles on mobile comic sites -- all involve
romance.

``It's a bit hard commenting publicly on this, but the most popular comics on the mobile are
adult-oriented ones for women,'' including love stories with sexually explicit content,
Mr.~Nakabayashi said. Translation: Women who do not want to be seen reading these titles in public
places like the train helped create the market for manga on the cellphone, which accords them
privacy in ways that magazines and books do not.

``Cellphone is a very private medium, much more so than a comic magazine,'' said Yutaka Tashiro,
director of rights group at Shueisha, another leading Japanese manga publisher. The romance- and
sex-filled manga -- known as shojo manga -- are read by women in their 20s, Mr.~Tashiro said. In
fact, about 70 percent of cellphone manga readers are women, he said.

So while the explosion in cellphone manga is not a simple story of migration from print to digital,
most experts agree that the future of manga lies more on the mobile than on paper.

``The idea of printing the magazine, loading them on the track and delivering them to distribution
centers, that whole model is on the decline,'' said Noboru Rokuda, professor at Kyoto Seika
University, which has a manga faculty, and a longtime artist himself. ``I like to keep the paper
manga tradition going, but there is an inexorable transition away from paper and into the digital
medium.''

If the migration is slow, it is the publishers themselves standing in its way, critics say.
Publishers of leading titles like Kodansha, Shueisha and Shogakukan -- who are also leading book
publishers -- have in most cases released content for mobile use that had already been published in
magazines and books.

The longtime practice has been to serialize new titles in the weekly or monthly comic magazines
first. If they prove popular, they are rolled into a single title. If the momentum continues, they
might be turned into TV animation series and into video games.

``They are afraid of destroying the old business model,'' said Mr.~Nakabayashi, of Nomura Research
Institute.

Analysts and fans say that all is not perfect with mobile manga. On the small screen, some of the
artists' intent is lost in translation. That is one reason that the mainstay manga readership -- men
-- have mostly stayed with the paper version. On the magazine pages, comic strips come in all shapes
and sizes, big, small and rectangular and sometimes come in a two-page spread, allowing for more
variety.

``Manga strips are placed in sequence on the pages, creating a special visual effect,'' said
Mr.~Rokuda, the manga professor. ``It is one rare medium that allows readers to have a bird's-eye
view on several strips and flow back and forth in time. That's lost when each strip is displayed on
the screen one by one.''

\section{Nissan Unveils Electric Car, to Be Available Next Year}

\lettrine{N}{issan}, Japan's third-largest automaker, took the wraps off its much-awaited electric
car on Sunday, naming the hatchback the Leaf.

Nissan and its French partner, Renault, have been aggressive proponents of pure electric vehicles
with no emissions, announcing plans to mass-market the clean but expensive cars globally in 2012.

Nissan will begin selling the first Leaf cars in the United States, Japan and Europe toward the end
of 2010, adding two more models soon after. It expects production to start with about 200,000 units
a year.

Coordinating the car's introduction with the inauguration of Nissan's global headquarters in
Yokohama, south of Tokyo, the chief executive, Carlos Ghosn, drove up to the stage in a sky-blue
Leaf prototype. Nissan is returning to the port city of Yokohama, where it was founded in 1933,
after being based in Tokyo's high-end Ginza shopping district for 41 years.

Hurt by sliding vehicle sales worldwide since the financial crisis hit last year, Nissan has
suspended its goals set under a midterm business plan, with the exception of its aggressive push
into electric cars.

With oil prices topping \$60 even in a recession and environmental regulations tightening all the
time, Mr.~Ghosn said he was optimistic about electric vehicles entering the mainstream, expecting
them to represent one in 10 new cars globally by 2020.

Other automakers like Toyota and Volkswagen have announced plans to start producing electric cars,
but they say it could take decades for the vehicles to spread because of their high cost, limited
driving range and long charging times with the current battery technology.

Nissan did not announce pricing for the five-seater Leaf, but Mr.~Ghosn said the price, without the
expensive lithium-ion battery that Nissan is considering leasing, would be within the range of a
comparable gasoline-engine car.

The Leaf has a top speed of more than 76 miles an hour and a cruising range of at least 100 miles --
a distance that covers the needs of 80 percent of drivers, Nissan said.

The carmaker has lagged its domestic rivals Toyota and Honda in offering hybrid cars, which have
taken the Japanese market by storm this year with the help of government stimulus steps.

For years, Mr.~Ghosn has played down the importance of hybrid cars -- both for consumers hoping to
save money at the pump and for manufacturers looking to make profit. But in what could signal a
shift in strategy, Mr.~Ghosn said Nissan would reserve the option of offering hybrid vehicles if the
market demanded it.

\section{Trendy Japanese Flock\cite{flock} to Hybrids}

\lettrine{S}{hoppers} who stop into Toyota dealerships here and across Japan generally have just one
thing on their minds: buying a Prius.

After years of indifference toward hybrids -- in part because many gas-only models that got
excellent mileage were already available in showrooms -- drivers in Japan are lining up for the
Prius, the Honda Insight and other models powered by gasoline and electricity.

The enthusiasm for hybrids is also a result of new models arriving in showrooms. Toyota introduced
the third-generation Prius in May; in June, it was the best-selling car in Japan, elbowing aside the
Insight, which in March became the first hybrid to hold the top sales spot.

Honda has been selling about 10,000 Insights a month, twice as many as had been expected in the
company's internal forecasts. There has been a five-month waiting list for the Prius in Japan,
similar to the situation in California last year when gas prices swept past \$4 a gallon.

``Ordinarily, small cars don't have much status, but people don't care about changing from a bigger
vehicle to a smaller vehicle if it's a Prius these days,'' said Toshiyuki Yokoyama, the general
manager at Kanagawa Toyota Motor Sales, which runs about 50 dealerships in and around Yokohama.
``Suddenly, everyone wants the car as soon as possible.''

Japanese consumers are only now embracing hybrids because, as in Europe, there have been plenty of
cheaper fuel-efficient alternatives. But the sales of new, reasonably priced models have jumped in
the past year or so because of government subsidies and a blizzard of media coverage that has fueled
Japan's fad-conscious consumers.

The question is whether the recent surge in sales of ``eco cars,'' as the Japanese call hybrids and
alternative-fuel vehicles, is a lasting trend or one pumped up by short-term subsidies and high
gasoline prices.

In the United States, where Japanese automakers had trouble keeping up with demand for hybrids in
2008, sales have slowed (along with the rest of the market). Toyota suspended work on a factory in
Mississippi devoted to building the Prius.

Japanese are eager consumers of trendy products -- everything from the latest handbags to video
games -- and sales often fall as fast as they rise.

Automakers in Japan say, however, that sales of hybrids will remain strong even if government
subsidies of up to \$4,000 are phased out. Though down from the peaks reached in 2008, gas prices
are still relatively high, and in a weak economy consumers are looking to spend less on gasoline.

From an early age, Japanese are taught to be more efficient because their country has few natural
resources. To reduce the use of air-conditioners, businessmen a few years ago were encouraged to do
without neckties. Appliance companies eagerly promote waterless washing machines.

But it is price that matters most, which is why Toyota and Honda used less costly technology to cut
the cost of their hybrids below three million yen, or \$30,000, the rough benchmark for a luxury car
in Japan.

Price was definitely an issue for Eiji Suga, a storekeeper in Kawagoe, a city one hour from Tokyo,
who reserved a white Prius in April.

Mr.~Suga's 13-year old Honda Accord, which has 43,000 miles on it, had very little resale value. But
because of the car's age, he qualified for a 250,000 yen (about \$2,500) government subsidy if he
scrapped the car and bought a more fuel-efficient replacement.

He will receive another 150,000 yen in tax breaks for buying an eco car. The price of the new Prius
was about 300,000 yen lower than its predecessor, too.

``Together with the tax breaks and other government benefits, it was easy to choose,'' said
Mr.~Suga, who added that seven of his friends in the neighborhood were also buying the new Prius.

Nissan and Mitsubishi have also announced plans to leapfrog hybrids with new electric vehicles.

But the durability of the hybrids on the market may make it hard for automakers in Japan to persuade
consumers to buy electric cars.

Tadashi Ishida, a retired electronics company executive, bought his Prius in 2006 to replace a
Toyota Mark II, taking advantage of a 190,000 yen government rebate. Mr.~Ishida, who lives in
Yokohama, has been particularly impressed with the car's fuel economy.

He uses just a quarter-tank of gas to get to his family's summer house about 75 miles away, about
half of what it took to get there in the Mark II. He said that hybrids were no passing fad.

``Average Japanese people do not consider the hybrid car a status symbol,'' Mr.~Ishida said. ``Our
Prius just fits our needs in terms of fuel consumption, daily family use and easy driving.''

\section{America Hears a Gaffe\cite{gaffe}, Russia Sees a Plot}

\lettrine{A}{fter} Vice President Joseph R.~Biden Jr.~gave an interview to The Wall Street Journal
portraying Russia as a limping and humbled nation, many in Washington responded last week with a
helpless shrug: There's crazy Joe, they said, the guy who once told a wheelchair-bound state senator
to stand up for a round of applause.

But in Russia, they weren't shrugging. Within hours, a top Kremlin aide had released a
barbed\cite{barbed} statement comparing Mr.~Biden to Dick Cheney. Commentators announced Mr.~Biden's
emergence as Washington's new ``gray cardinal\cite{cardinal}'' -- the figure who, from the shadows,
makes all the decisions that matter. Others said Washington's mask had been torn off, revealing
Mr.~Obama's ``reset'' as at best insubstantial and at worst duplicitous.

American officials spent several days trying to convince their Russian counterparts that Mr.~Biden's
words were, for lack of a better label, a gaffe. Russia's highest officials have kept silent on the
matter, but their initial responses were skeptical.

``Biden has said this in such a way that the whole world heard it,'' said Alexei K.~Pushkov, who is
the anchor of the current events show ``Post-Scriptum.'' ``And then there are secret, furtive calls
in the night, dragging Russian officials from their supper. They want to say this is not true. But
somehow everybody still thinks it is.''

Among the reasons for their skepticism: In today's Russia, politicians just don't run off at the
mouth. Not so long ago, Russian public life was a symphony of embarrassing episodes. Remember when
Boris Yeltsin confused Norway with Sweden, suggested that Germany and Japan had nuclear arsenals,
and toppled over while saluting an honor guard in Uzbekistan?

That all ended with the presidency of Vladimir V.~Putin. Mr.~Putin, now Russia's prime minister,
occasionally departs from statesmanlike language, as when he threatened to hang the Georgian
president by his testicles or offered a French reporter an especially thorough circumcision. But
coming from Mr.~Putin, these statements are expressions of Russian might, something like a political
philosophy -- never, ever mistakes.

For anyone subordinate to the president to allow themselves that freedom is inconceivable, said
Vladimir V.~Pozner, the host of a talk show on state television.

``If it's not the No.~1 man or woman, clearly that person has been instructed to say what he or she
said,'' Mr.~Pozner said. ``It's psychologically very difficult for a Russian to believe otherwise.
If you write in The New York Times whatever you write, I'm sure Mr.~Putin will say, 'Of course. It
was ordered.' ''

It will also be hard to convince the Kremlin that the comments don't indicate a deeper drama.
Russians have spent months searching for clues to Mr.~Obama's true intentions; when Mr.~Obama killed
a fly during a television interview shortly before traveling to Moscow, for example, several
analysts here interpreted it as a message to Russia.

Mr.~Biden has now supplied evidence for two plotlines -- a deep rift within the administration, or a
``sophisticated game,'' said Andrei V.~Ryabov, a political analyst at Moscow's Carnegie Center. This
ambiguity, he said, plays into the conviction of Mr.~Putin and his team that real events take place
far from view, among a handful of powerful individuals, and that public politics are ``no more than
puppetry, decoration in the theater.''

``Nothing accidental can happen in this system,'' Mr.~Ryabov said. ``Everything has a hidden
meaning.'' Even accidental words from officials are likely to be read closely; as a Russian proverb
has it, ``What a sober man has on his mind, a drunk puts on his tongue.''

Mr.~Pushkov was among those who put little credence in Mr.~Obama's overtures, and to him,
Mr.~Biden's words offer a far more honest assessment of American policy. He says he reads in them a
split in Washington between cold war heavyweights and a president too weak to bring them to heel.

``It's not just a question of schools of thought,'' he said, dryly, but something far more serious.
Schools of thought, he added, are something to be ``exercised on a veranda with a cup of coffee on a
summer evening.''

Of course, every warming of the relationship between Moscow and Washington has been a tenuous
process, punctuated by false starts and furious backpedaling.

In 1974, after signing on to the idea of ``peaceful coexistence,'' Leonid Brezhnev seems to have
been called on the carpet by a Central Committee concerned about ceding ground to the United States;
he went on to repudiate two key agreements with the Americans. Jimmy Carter, under a drumbeat of
criticism for caving in to Russia, halted ratification of the second strategic arms limitation
treaty after the Soviets invaded Afghanistan in 1979; he explained that the invasion had changed his
view of Moscow's intentions.

This thaw seemed tentative, too, even before Mr.~Biden's words. The coming months could bring
renewed fighting in Georgia, or another gas crisis with Ukraine, or a deadlock on the renegotiation
of the Strategic Arms Reduction Treaty.

``At this point these were just words -- unfortunate words, reckless words, but still, it was just
words, not of the president but of the vice president,'' said Dimitri K.~Simes, the president of the
Nixon Center. ``The question is what is going to happen next.''

\section{Confusion Surrounds `Clunkers' Program}

\lettrine{W}{hen} the White House and Congress decided late last week to give their ``cash for
clunkers'' program a reprieve, they did little to relieve harried and befuddled\cite{befuddle}
automobile dealers and customers trying to sort it out.

``There is absolute frustration across the board,'' said Alex Kurkin, a Miami lawyer representing
several car dealers. ``They're not really confident about any deals, and no one can give them advice
about what they should be telling their customers.''

The program allocated \$1 billion for rebates of up to \$4,500 to customers trading in clunkers for
more fuel-efficient new vehicles. It was set to expire in November, but on Thursday, the federal
Transportation Department said it had been so popular that the money was gone in days.

Officials advised dealers to stop offering the rebates, and the program appeared suspended. But
White House officials said it had not been. On Friday, the House authorized \$2 billion more to
replenish the program. But the Senate will not vote until next week, so until then, the status of
the program is unclear.

While many in Washington hailed the program as a success for the auto industry, it has left some
dealers anxious over whether there will be enough money to repay them for rebates they have already
given.

Barry Magnus, the general manager at DCH Paramus Honda in New Jersey, said he had sold more than two
dozen new cars under the program but was still owed more than \$80,000. The government had said it
would take 10 days to reimburse\cite{reimburse} dealers, and that was before the program ran out of
money. The government Web site where dealers were supposed to register their deals
crashed, and many dealers have not been able to update their information.

The program requires that the clunkers be put out of service for good. The engines must be
destroyed. So, after a certain point, any decision to trade in an old car is irreversible.

Nick Clites, who is in charge of used cars at DCH, took a reporter through the process as he prepped
a 1988 BMW, with 214,000 miles on the odometer, for its death. He drained the oil, then
donned\cite{don} a silky blue protective suit, goggles and gloves and poured a sodium silicate
solution into the engine. He revved\cite{rev} the car, and within seconds, the solution hardened,
the engine seized up and was kaput.

Wreckers were called to pick up the car, remove the spoiled engine and crush the car. They have to
certify that the cars have been crushed and rendered inoperable before the government will process
reimbursement claims.

When the program appeared to run out of money, dealers became even more anxious to get the clunkers
off their lots quickly.

``Oh my God, what a mess today,'' Sally Ann Maggio, a co-owner of Hackensack Auto Wreckers, said
Friday. ``We have been overwhelmed with phone calls from the dealerships. They're hitting the panic
button.''

Ms.~Maggio said she did not think much of the program to begin with because she made her profit by
reselling engines from junked cars. But the scrap is worth something.

Dan Cohen, a retired lawyer on Long Island, made a deal to trade in his 1988 Jeep for a new Toyota
in early July, before the program went into effect (and when showrooms were quiet). He qualified for
a \$4,500 dealer rebate and said he signed a contract that also gave him \$500 as scrap value. When
he returned Friday to complete the deal, the showroom was ``crawling with customers,'' he said --
and the dealer cut the scrap fee to \$100.

So the deal is off. ``I may not get the \$4,500,'' Mr.~Cohen said, ``if I go somewhere else and the
federal allotment is used up already.'' He said that if he could not find a comparable deal, he
might sue. (Calls to the dealer for comment did not go through.)

Mr.~Kurkin, the Miami lawyer, said many dealers were adding clauses to their sales contracts, saying
that if the rebates did not come through, the customer would have to pay the difference.

``If a dealer doesn't have a separate document addressing this possibility, the dealer will likely
have to eat it,'' he said. ``I certainly see a lot of litigation over this.''

\section{A Smart Shower May Even Know Your Song}

\lettrine{I}{n} these digital times, even the venerable\cite{venerable} bathroom faucet\cite{faucet}
is getting an update. Companies like Moen and Hansgrohe are now replacing faucet handles in tubs and
showers with touch screens and other electronic controls that are smart enough to store all the
details of time and temperature of a favorite shower or bath -- and produce them at a touch.

With the RainBrain shower from Hansgrohe, the mix of hot and cold water, its volume and flow, and
even music from a favorite MP3 playlist can be retrieved by tapping a small screen, said Nicolas
Grohe, director of marketing and product development for the American branch of the company in
Alpharetta, Ga. (Hansgrohe is based in Schiltach, Germany.) The RainBrain will be offered in the
United States beginning in March, he said, and cost about \$4,500.

Even amid a deluge of water, users can easily tell the system to regulate the showerheads, which
pour water from overhead and both sides, as well as a handheld sprayer, Mr.~Grohe said. ``The dials
and controls are big,'' he said. ``It's difficult to miss them.''

A bar marked with a plus sign, for instance, makes the water hotter, a minus cools it. If users
prefers a quick alternative, there's also a rotary dial.

The menu button on the touch screen offers a program called Scottish shower that varies the water
temperature, automatically switching from hot to cold. ``You can tell it how hot and cold you want
it,'' Mr.~Grohe said, and to set how long each cycle lasts.

Another Hansgrohe program, Cascade, turns different showerheads on and off in rhythmic, pulsing
jets. The company has posted a simulationof the control panel on its Web site.

Another electronic system, ioDigital, was introduced in May by Moen, a company based in North
Olmsted, Ohio, that is known for its faucets. The system is available for a standard single-head
shower (about \$1,200), for Roman or whirlpool tubs (about \$1,300) and for a shower with multiple
heads, called a vertical spa (about \$2,500), said Mark Knurek, senior product manager. The price
includes the showerheads, water spouts for tubs, and related fixtures, as well as the electronic
controls.

The control panel for the ioDigital shower or bath can be mounted in the shower, on the rim of a
bathtub or outside the shower on a wall, he said. A cable connects the interface to the valve
assembly that regulates the flow of hot and cold water. The assembly can be placed in the wall
behind the shower, for example, or farther away, so long as the location is an interior wall
insulated from the outside.

``We've put the assemblies into crawl spaces, in vanities and even in basements,'' Mr.~Knurek said.

In May, Debbie Brown, a chemist who lives with her husband and two children in Bath, Ohio, decided
to use ioDigital to retrofit a shower and whirlpool tub in the master bathroom. She liked the option
of electronic controls to replace the old, defective faucet handles, in part because there would be
minimal damage to the bathroom.

``We didn't have to tear up the ceramics,'' she said.

Most of the work was done on the wall near the bathroom, she said, where the plumber placed the
valve assembly.

The electronic control panel for the bath is on the wall above the tub, as is the shower control. A
fiberglass specialist filled in the holes in the tub where the faucet handles once were. ``You can't
tell anything was there,'' she said.

Installing the new shower controls was straightforward, too, she said. Mrs.~Brown saved a few pieces
of tile from the original installation that could be used to cover up the holes from the faucet
handles.

Mrs.~Brown likes the new, streamlined look of the bath and shower. ``And you don't have all of the
handles to clean,'' she said.

She also enjoys being able to program the water temperature. Each family member has chosen a
different setting, saved with a number from 1 to 4.

``I'm No.~1,'' she said. ``I set the temperature I want, hit the button and the water is right where
I want it all of the time, instead of never quite right.''

A remote control (\$135) can be used to start the shower from another room, or to top off water in
the bath.

There is, of course, one big drawback to electronic controls: when there's no power, there's no
shower. Unlike the humble faucets, those touch screens

and panels won't work without electricity.

\section{10 Years Ago, an Omen\cite{omen} No One Saw}

\lettrine{T}{en} years ago Sunday, on an island off the coast of America, something impossibly
glamorous\cite{glamorous} took place. Partygoers took boats from Manhattan to the home of the Statue
of Liberty to plop\cite{plop} pashalike on pillows and blankets and munch\cite{munch} on lamb chops
while Macy Gray sang, their faces illuminated with multicolored Chinese lanterns and fireworks
curated and narrated by George Plimpton.

``That was a silver-flanged fleur-de-lis,'' said the voice, highly recognizable but
disembodied\cite{disembody} by darkness.

This was the Talk magazine launch party on Aug.~2, 1999 -- simply called The Party at the time --
and it seemed as if a new era of media fabulousness had been christened\cite{christen}. The Hearst
Corporation and Miramax, owned by Disney, decided to finance a new general interest magazine led by
Tina Brown, fresh off her triumphs at Vanity Fair and The New Yorker, that would lead the national
conversation.

It was all kicked off with the kind of lavish\cite{lavish} party that would seem unthinkable in the
current context.

Sponsored liquor\cite{liquor} flowed, women teetered about on heels in deep grass, and the A-list
guest list -- Mr.~Kissinger, please meet Miss um Ms.~uh meet Madonna -- was a testament to the power
of the synergized word. Content was king and Ms.~Brown was its queen.

``Now you're not exactly the tired masses, the huddled masses, but then again, I'm an immigrant who
toiled here on the Concorde,'' she said to the crowd after being introduced by Queen Latifah. ``But
I just want to say, here's to Lady Liberty tonight.''

Too bad nobody saw the sharks circling in the harbor. Rather than the culmination\cite{culmination}
of a century of press power, the Talk party was the end of an era, a literal fin de si\`ecle. Flush
with cash from the go-go '90s and engorged by spending from the dot-com era, mainstream media
companies seemed poised on the brink\cite{brink} of something extraordinary. But that brink ended up
being a cliff.

``It seems like that happened in the 18th century,'' said Ms.~Brown by phone last Friday.

Magazines are on pace to book little more than half of the advertising pages that the industry did
10 years ago, and dozens of longtime titles have disappeared. The last big magazine introduction --
Portfolio at Cond\'e Nast Publications -- flamed out this spring after two years at a cost of more
than \$80 million. Now even Cond\'e Nast Publications, the world headquarters of printed luxury, has
brought in the bean counters from McKinsey with an eye toward further cuts. There may never be
another large magazine launch ever, and certainly not one that was accompanied by the fanfare of
Talk.

I'm still ashamed to admit that I wasn't one of the lucky 1,000 people invited to the party -- old
prerogatives\cite{prerogative} die hard -- so I was trapped on shore, covering it secondhand with a
nose pressed up against the glass. But it is worth thinking about how this future, or lack of one,
arrived so unforeseen.

Ten years ago, journalists, long the salarymen of the publishing economy, began gorging\cite{gorge}
on big contracts and options from digital start-ups like shrimp\cite{shrimp} at a free buffet. With
coveted\cite{covet} writers commanding \$5 for every typed word into magazines that were stuffed to
the brim with advertising, there was a fizziness, some would say recklessness, in the air. The
industry was drunk on its own prerogatives, working a party that seemed as if it would never end.

Peter Kaplan, the former editor of The New York Observer, attended the party and oversaw coverage of
the event.

``Tina, for all the excellence of her antenna, was scratching the air, and like many of us, was
unable to pull in the new signal,'' he said. ``She failed to see that it was probably already over
and that there was something slightly hollow about that event.''

Most of us who covered media did not fully understand the implications of the new technology that
could publish and distribute information at zero marginal cost. The Web was viewed as a niche, as a
way to supplement and enhance the printed product, certainly not a threat that would make many of
those publications obsolete.

``Most of the talk at the Talk party was about the party itself,'' said Kurt Andersen, a novelist,
radio host and founder of Spy magazine. ``It was weird and interesting because you were sort of
wandering around in the dark out there and bumping into people. There was a meta quality to the
thing, a self-consciousness, that in retrospect\cite{retrospect} was probably telling.''

At least Ms.~Brown did not compose a rap ode to the new magazine. That fell to Mr.~Big, Ron Galotti,
the former Cond\'e Nast publisher who managed to get a flock of advertisers to buy the
hype\cite{hype} and commit to the first four issues of a magazine they had never seen. After Talk
closed, Mr.~Big quit Manhattan media and moved to a farm in Vermont. Maybe he knew something we
didn't. (He did not return a call.)

``It was the end of something extraordinary, but none of us knew it at the time,'' Ms.~Brown says
now. ``What followed was a very turbulent odyssey, not just for me, but for all of us. There has
been a volcanic\cite{volcanic} realignment that none of us foresaw.''

After Talk closed early in 2002, Ms.~Brown hosted a television show on CNBC and wrote a book about
Princess Diana. Demonstrating a nimbleness that has characterized her entire career, she is now
running The Daily Beast, a scrappy but promising digital media site owned by Barry Diller's
InterActiveCorp.

She pays her writers, increasingly an exception these days, but there are no huge contracts or boat
rides to sylvan lawns full of impossibly famous people. The Daily Beast has 1.5 million unique
visitors a month, according to Quantcast, and has kicked up some notice, but its opening party of a
hundred or so took place at the very much land-locked Pop Burger on Ninth Avenue in Manhattan. There
was no Sarah Jessica Parker, no Robert De Niro and no Hugh Grant in attendance. There were a lot of
bloggers instead.

Modern media success is enabled by brutal\cite{brutal} cost control and using hard, fast numbers to
convince advertisers they will get a return for their spending. Once stalwart magazines like
BusinessWeek are up for grabs and entire formerly lucrative categories have been wiped out. The
magazine canard of associative glamour, of selling aspiration by the bucket-load with page after
page of pricy merchandise, is all but dead but for a few exceptions.

Ms.~Brown once wrote the book ``Life as a Party'' and on that night, it was.

``I was aware it was a historic night,'' Ms.~Brown said. ``We were on a boat and I was with Natasha
Richardson. We were talking and laughing, looking at the lights of the twin towers. And then a big
wave came over the side of the boat and soaked us both. Now Natasha is gone, the towers are gone.
It's very, very sad, but I am very excited by this new world we are heading into.''

\section{Staying Tough in Crackdown on Immigrants}

\lettrine{A}{fter} early pledges by President Obama that he would moderate the Bush administration's
tough policy on immigration enforcement, his administration is pursuing an aggressive strategy for
an illegal-immigration crackdown that relies significantly on programs started by his predecessor.

A recent blitz\cite{blitz} of measures has antagonized\cite{antagonize} immigrant groups and many of
Mr.~Obama's Hispanic supporters, who have opened a national campaign against them, including small
street protests in New York and Los Angeles last week.

The administration recently undertook audits of employee paperwork at hundreds of businesses,
expanded a program to verify worker immigration status that has been widely criticized as flawed,
bolstered a program of cooperation between federal and local law enforcement agencies, and rejected
proposals for legally binding rules governing conditions in immigration detention centers.

``We are expanding enforcement, but I think in the right way,'' Janet Napolitano, the homeland
security secretary, said in an interview.

Ms.~Napolitano and other administration officials argue that no-nonsense immigration enforcement is
necessary to persuade American voters to accept legislation that would give legal status to millions
of illegal immigrants, a measure they say Mr.~Obama still hopes to advance late this year or early
next.

That approach brings Mr.~Obama around to the position that his Republican rival, Senator John McCain
of Arizona, espoused\cite{espouse} during last year's presidential campaign, a stance Mr.~Obama
rejected then as too hard on Latino and immigrant communities. (Mr.~McCain did not respond to
requests for comment.) Now the enforcement strategy has opened a political rift with some immigrant
advocacy and Hispanic groups whose voters were crucial to the Obama victory.

``Our feelings are mixed at best,'' said Clarissa Martinez De Castro, immigration director of the
National Council of La Raza, which has joined in the criticism, aimed primarily at Ms.~Napolitano.
``We understand the need for sensible enforcement, but that does not mean expanding programs that
often led to civil rights violations.''

Under Ms.~Napolitano, immigration authorities have backed away from the Bush administration's
frequent mass factory roundups of illegal immigrant workers. But federal criminal prosecutions for
immigration violations have actually increased this year, according to a study by the Transactional
Records Access Clearinghouse, a nonpartisan group that analyzes government data. In April, there
were 9,037 immigration cases in the federal courts, an increase of 32 percent over April 2008, the
group found.

Ms.~Napolitano said in the interview that she would not call off immigration raids entirely as some
Hispanic lawmakers have suggested. ``We will continue to enforce the law and to look for effective
ways to do it,'' she said. ``That's my job.''

Ms.~Napolitano, who as governor of Arizona sparred with Republican legislators seeking tougher steps
against illegal immigration, said she was looking for ways to make enforcement programs inherited
from President George W.~Bush less heavy-handed. She also wants to put the enforcement focus on
illegal-immigrant gang members and convicts\cite{convict} and on employers who routinely hire
illegal immigrants so as to exploit them.

Immigration authorities have started audits of employees' hiring documents at more than 600
businesses nationwide. If an employer shows a pattern of hiring immigrants whose documents cannot be
verified, a criminal investigation could follow, Ms.~Napolitano said.

She has also expanded a federal program, known as E-Verify, that allows employers to verify
electronically the identity information of new hires. Immigrant and business groups have sued to try
to stop the program, saying the databases it relies on are riddled with inaccuracies that could lead
to American citizens' being denied jobs.

But officials of the Homeland Security Department say technological improvements have enhanced the
speed and accuracy of E-Verify. With 137,000 employers now enrolled, only 0.3 percent of 6.4 million
queries they have made so far in the 2009 fiscal year have resulted in denials that later proved
incorrect, the officials say. That, opponents note, still means false denials for more than 19,000
people.

In addition, Ms.~Napolitano has expanded a program that runs immigration checks on every person
booked into local jails in some cities. And she recently announced the expansion of another program,
known as 287(g) for the provision of the statute authorizing it, that allows for cooperation between
federal immigration agents and state and local police agencies.

In extending 287(g), federal officials also drew up a new agreement, which all of some 66 localities
currently participating have been asked to sign, that is intended to enhance federal oversight and
clarify the priority on deporting those immigrants who are criminal fugitives or are already behind
bars.

But advocates for immigrants were dismayed that the new agreement did not include strong protections
against ethnic profiling. They were surprised, they say, that Ms.~Napolitano did not terminate the
cooperation agreement with the sheriff of Maricopa County, Ariz., Joe Arpaio, who calls himself the
``toughest sheriff in America.'' Latino groups in Arizona have accused Mr.~Arpaio of using the
program to harass Hispanic residents.

``If they reform the 287(g) program and Arpaio doesn't change, it won't be reform,'' said Frank
Sharry, executive director of America's Voice, a national immigrant advocacy group.

Ms.~Napolitano said it would be up to Mr.~Arpaio, like other current participants, to decide whether
to sign and abide by the new cooperation agreement. Separately, the Justice Department has opened a
civil rights investigation of Mr.~Arpaio's practices.

The Obama administration has received support for its immigration position from a leading Democrat,
Senator Charles E.~Schumer of New York, the chairman of the Judiciary subcommittee on immigration,
who will be writing an immigration overhaul bill later this year.

In preparation for what is likely to be a furious debate, Mr.~Schumer has called on Democrats to
show that they are serious about immigration enforcement and is even asking them to stop using the
term ``undocumented'' to refer to immigrants who are here illegally.

Democrats have to ``convince the American people there will not be new waves of illegal immigrants''
after an overhaul passes, Mr.~Schumer said in an interview.

Republicans who oppose any legalization of the status of illegal immigrants say they remain
unimpressed by the new enforcement measures.

``After 20 years of broken promises, it takes a lot more than token gestures,'' said Representative
Brian P.~Bilbray, a California Republican who heads an immigration caucus in the House.

Michael A.~Olivas, a professor of immigration law at the University of Houston, said Hispanic
advocates were irked\cite{irk} by the enforcement measures because they had seen scant\cite{scant}
sign that the administration was also moving deliberately toward an overhaul bill.

``We literally have the worst of all worlds,'' Professor Olivas said.

\section{Obama Renews Vow of No Middle-Class Tax Increase}

\lettrine{T}{he} White House tried Monday to douse\cite{douse} speculation that it might raise taxes
on the middle class in violation of President Obama's campaign promise, just a day after two of his
top economic advisers left the door open to such a move to rein in spiraling\cite{spiral} deficits.

Mr.~Obama told his economic team in a meeting at the White House that he intended to stand by his
promise not to increase taxes on families making less than \$250,000, aides said. He then sent his
spokesman out to repeat that message in front of the television cameras.

``The president made a commitment in the campaign. He's clear about that commitment, and he's going
to keep it,'' said Robert Gibbs, the White House press secretary.

The renewal of the promise came a day after Treasury Secretary Timothy F.~Geithner and Lawrence
H.~Summers, director of the National Economic Council, both refused to rule out tax increases on the
middle class while discussing ways to pare\cite{pare} the deficit. The two were speaking on separate
Sunday morning talk shows, venues where administration officials are usually well prepared on the
official line before appearing.

``It's never a good idea to absolutely rule things out no matter what,'' Mr.~Summers said on ``Face
the Nation'' on CBS. Mr.~Geithner, on ABC's ``This Week With George Stephanopoulos,'' said, ``We
can't make these judgments yet about exactly what it's going to take'' to tame the deficit.

Conservative critics interpreted those comments as laying the groundwork for trying to wriggle out
of Mr.~Obama's campaign pledge.

``Obama should fire Geithner and Summers,'' said Grover Norquist, president of Americans for Tax
Reform, a group that opposes tax increases. They ``went on national television and implied the
president lied his way into office and that he is open to raising taxes.''

The developments come at a time when the White House and Congressional Democrats, trying to figure
out how to pay for expanding health care coverage, are considering proposals to increase taxes on
the wealthiest Americans. Some critics from the left have suggested that Mr.~Obama should not limit
tax increases to the rich so that a broader cross section of Americans would be invested in the new
health care system, as they are in Social Security and Medicare.

But the White House is trying to fend\cite{fend} off attacks portraying Mr.~Obama as a tax-and-spend
liberal. Mr.~Gibbs said that he had read the transcripts from Sunday's shows ``a few times'' to
study what had been said and that the president had made a point of reminding Mr.~Geithner and
Mr.~Summers of his position, but was not scolding them.

``We talked about it as an issue,'' Mr.~Gibbs said, but added, ``This wasn't a, you know, like
'school is in' type of thing.''

Mr.~Gibbs seemed exasperated\cite{exasperated} at repeated questions on the matter at his daily
briefing.

``If you don't trust what I'm going to tell you, I don't know why we do this,'' he said finally.

Asked why Mr.~Geithner and Mr.~Summers did not repeat the president's campaign promise, he said,
``They left it to me.''

\section{Rivalry\cite{rivalry} Between Apple and Palm Intensifies Over Access to iTunes}

\lettrine{T}{he} Palm Pre has a large touch screen, slide-out keyboard and fast Web browsing. Palm
also likes to point out that another selling point is the smartphone's ability to link to iTunes,
Apple's music and media store.

Trouble is, Apple wants to make sure the iPhone is the only cellphone that can do that. So it
changed its software to block the Pre's access to iTunes.

Now Palm is calling foul and is trying to rally the consumer electronics industry to its defense.
Palm says that Apple, which allows only its own devices to connect directly with iTunes software, is
misusing the standards put in place to foster interoperability between computers and devices using a
U.S.B. connection.

Palm has filed a complaint with the U.S.B. Implementers Forum, an industry group established by
technology companies that developed the technology that links computers to other electronic devices,
claiming Apple is restraining trade. Predicting the outcome of this particular filing is tricky,
said Mike Abramsky, an analyst with the investment firm RBC Capital Markets.

``There isn't much precedent for this case,'' he said. ``It's breaking new ground. In my mind,
ultimately the users are the arbiters\cite{arbiter} in the outcome of these situations.''

The forum declined to comment on when it would respond to Palm's filing.

Currently, Palm, which once again has the Pre synching with iTunes, accomplishes the feat\cite{feat}
by duping\cite{dupe} the iTunes software into recognizing the phone as an Apple music player,
allowing it to synchronize and transfer files between the phone and a personal computer. The tactic
has raised some eyebrows among those who think that Palm's approach itself is a breach of the
standards set by the U.S.B. governing board.

But Palm says its strategy of masquerading\cite{masquerade} as an iPod is acceptable because it is
in response to Apple's restriction. ``We think we are consistent with our compliance,'' said Douglas
B.~Luftman, an associate general counsel for Palm. ``We're not trying to appear to be anything we're
not -- except for interoperability purposes with iTunes.''

Other devices not made by Apple work with iTunes, but generally require an additional download or
intermediary step to work. Palm's strategy, which allows an effortless connection between the
smartphone and iTunes, is a first for the industry.

If anything, Palm's sparring with Apple underscores ``how critical the iTunes ecosystem has become
for consumers,'' said Michael Gartenberg, vice president for strategy and analysis at Interpret, a
market research firm based in Los Angeles and New York. ``ITunes is the center of gravity where
consumers keep their content.''

Palm is not the only company clashing with Apple over its hardware and software policies. Google and
Apple are tussling\cite{tussle} over Apple's refusal to allow certain Google applications on the
phone. On Friday, the Federal Communications Commission said it was looking into those actions.

For Palm and Apple, the iTunes battle is part of a bigger rivalry. The Palm Pre, introduced in June
to glowing reviews, has been lauded as a worthy competitor to the iPhone. The competition is also
set against Palm's luring some of Apple's senior executives. Jonathan J.~Rubinstein, Palm's chief
executive, once oversaw iPod development at Apple, and Mike Bell, senior vice president for product
development, is another Apple alumnus.

``It's possible Apple is worried that customers might start buying Palm phones instead of iPhones or
iPods,'' said J.~Gerry Purdy, the chief analyst of mobile and wireless at the research firm Frost \&
Sullivan.

Mr.~Gartenberg, the analyst at Interpret, agreed. ``Apple understands that seamless relationship
between the iPod, the iPhone and iTunes,'' he said. ``It's a big driver behind why consumers are
buying their devices, and they're going to try and protect that.''

Tom Neumayr, a spokesman for Apple, said, ``As we've said before, newer versions of iTunes software
may no longer provide synching functionality with unsupported digital media players.''

The stakes are especially high for Palm, which once dominated the smartphone market with its Treo
handsets. The company, which has been steadily losing market share since 2007, is counting on the
Pre to reverse its fortunes. Paul Coster, an analyst with J.~P. Morgan who follows the company,
estimated that Palm shipped close to 180,000 devices in its first two weeks and could reach as many
as 2.5 million in the fiscal year ending in May 2010.

For Pre owners, accessing iTunes is important, said Katie Mitic, senior vice president for product
marketing at Palm. ``We appreciate that many, many of our consumers use iTunes as their music or
photo management system,'' she said. ``We just want to make it as easy as possible for them to use
it on the Pre.''

Of course, the Pre also has software restrictions on synching with the Zune, the portable media
player that is Microsoft's answer to the iPod, but since it has few users, Palm is not fighting for
access there. ``The Zune has a much smaller market share,'' Mr.~Gartenberg said. ``There are fewer
people using it to control their media, so there are very few other vendor devices who are saying
they want to hook into the Zune software.''

Palm's tiff with Apple also has some people wondering why Apple even bothers. Mr.~Purdy, the
analyst, described Apple's decision to thwart Palm's working with iTunes as ``shooting themselves in
the foot.''

``All this means is that Palm's customers would be able to purchase music through iTunes,''
Mr.~Purdy said. ``This would offer an opportunity for increased revenue for Apple.''

Tim Wu, a professor at Columbia who specializes in telecommunications law, copyright and
international trade, said, ``There's something very unseemly about what Apple is doing.'' He added,
``It's very counter to the ideals of openness, which is a concept Apple pioneered in computing.'' In
2007, Steven P.~Jobs, Apple's chief executive, issued a call to the music industry for openness,
titled ``Thoughts on Music.''

As for Palm, Mr.~Wu said, ``It sounds like an uphill battle, in terms of trying to stop Apple from
doing this.''

But Palm may have a shot. ``The history suggests that openness wins,'' said Mr.~Wu, citing examples
like AT\&T's attempts to restrict the devices attached to its phone lines and Apple's early attempts
to sell printers that worked only with Macs.

It was not immediately clear how long Palm planned to battle future software patches from Apple that
disable the smartphone's compatibility with iTunes. The company did not say whether it would pursue
lawsuits or enlist government aid on restraint of trade grounds.

``This is a classic technology cat-and-mouse game,'' Mr.~Gartenberg said. ``It often comes down to
which side tires first.''

\section{New President Tries to Get Toyota Back on the Road}

\lettrine{G}{eneral} Motors, Ford and Chrysler are not the only ones working through wrenching
restructurings. Toyota is, as well, though with a much lower profile.

The results have yet to show, and in fact, no one at Toyota expects 2009 to be anything but dismal.

Its new president, Akio Toyoda, has moved quickly since taking charge in June, when he declared his
dismay at the company's financial crisis.

``Like everyone in the company, I am extremely frustrated'' about the automaker's decline,
Mr.~Toyoda said at his first news conference as president. ``So we must start again from the ground
up.''

His first target was Toyota's top management. Already, 40 percent has retired or been reassigned.

Four of Toyota's five executive vice presidents, the group that now leads the company under
Mr.~Toyoda, are new to their jobs. Only one, Takeshi Uchiyamada, Toyota's product development chief,
is a holdover from the team that surrounded Mr.~Toyoda's predecessor, Katsuaki Watanabe.

Further, the four newcomers are each in charge of a global region on top of duties within the
company. (Mr.~Uchiyamada is the only one with a single focus.)

Now Mr.~Toyoda is headed for the United States. On Wednesday, he will address a management seminar
sponsored by the Center for Automotive Research, near Traverse City, Mich., an annual gathering of
the clan in the auto industry.

Toyota dealers and employees, many of whom know Mr.~Toyoda from the years he worked in the United
States, are eager for a shot of adrenalin\cite{adrenalin} from their new boss, the grandson of
Toyota's founder.

``He knows the business,'' said P.~Compton Cramer Jr., the owner of Cramer Toyota of Venice, Fla.,
who first met Mr.~Toyoda a decade ago.

``He's going to bring a new kind of energy to the company,'' added Josephine Cooper, Toyota's group
vice president of government and industry affairs, who had been Mr.~Toyoda's hostess in Washington
before his promotion.

Mr.~Toyoda's biggest problem, executives and analysts say, is transforming a lineup dotted with
mundane\cite{mundane}-looking vehicles that buyers bought because of their reliability, not because
of their style appeal.

It is a problem shared with Detroit, where executives are striving to inject their automobiles with
excitement, too. Unlike G.M., Ford and Chrysler, however, Toyota cannot reach back to its roots and
tap into a DNA of sexy sheet metal.

And it faces a difficult balancing act: become too aggressive, and Toyota risks chasing off the
practical-minded consumers who have been its base for a half-century. Some have already been scared
off.

Pamela Templeton, an owner of Fort Myers Toyota in Florida, said Toyota was caught ``a little
flat-footed'' when gas prices hit record levels last year. Despite a longtime focus on small cars,
the company bet heavily this decade on trucks like the Tundra and sport utility vehicles like the FJ
Cruiser.

So did Ms.~Templeton, who spent \$13 million to open a separate 165,000-square-foot showroom just
for light trucks. The market's demise forced her to eliminate customer pampering like a massage
chair and cappuccino machine. ``We had to make hard choices, and I'd rather keep people than
fluff\cite{fluff},'' she said.

Only two years ago, Toyota was riding the crest of a deliberate growth strategy\cite{stray} in which
it doubled in size and became the world's biggest car company. Although most carmakers have
suffered, Toyota's rapid decline shocked many at the automaker and the experts who follow it.

Yoshimi Inaba, named by Mr.~Toyoda as Toyota's new head of North American operations, faces heavy
pressure to turn things around here. ``Without North America, we are not likely to come back to
global proficiency,'' Mr.~Inaba told reporters in Detroit last month.

Toyota has an advantage over its Detroit rivals in its streamlined processes for engineering and
building its cars, said James P.~Womack, an author and expert on company efficiency. But he said the
company, famous for caution and self-reflection, strayed when growth, not quality, became its main
concern.

``They sort of quit being paranoid\cite{paranoid},'' Mr.~Womack said. ``I'm neutral to buy on Toyota
over the long term. It's going to take them a while to get back on their feet.''

Ms.~Cooper faces a unique concern in Washington, where the Obama administration spent \$65 billion
this year to shepherd G.M. and Chrysler through bankruptcy and to pay for their restructurings.

She said she was constantly being asked how her company, which spent years trying to emphasize its
investments in the United States, viewed the White House support for its rivals.

``My job is to make sure there is not more attention paid to Toyota because of this,'' Ms.~Cooper
said.

But Toyota cannot avoid scrutiny\cite{scrutiny} as it decides what to do with its half of New United
Motor Manufacturing Inc., or Nummi, a 25-year-old joint venture with G.M. in Fremont, Calif.

G.M. has pulled out of the deal and classified the plant as one of the assets of the old G.M., which
are being liquidated in bankruptcy.

Toyota said last week that it might dissolve its half of the venture, which was the brainchild of
Mr.~Toyoda's father, Shoichiro Toyoda, and paved the company's way into the American car market in
1984.

That has set up a confrontation with the United Automobile Workers union, which represents workers
at the plant. Sergio Santos, the local union president, argues that the plant's experienced work
force is among the country's best, consistently ranking ahead of other Toyota factories on
productivity surveys.

``We have Super Bowl trophies to back up what we do,'' Mr.~Santos said.

Union leaders are pressing lawmakers in California to provide incentives so Toyota can keep the
plant open, while U.A.W. leaders have urged members to lobby Congress for worker protections.

Even if the venture dissolves, Mr.~Santos said he hoped the plant could be sold back to Toyota or
another new owner like Roger Penske, the former racecar driver and entrepreneur who recently
purchased Saturn.

Senior executives in Japan will make the final decision on Nummi, Ms.~Cooper said.

Closing the plant would be a turning point for Toyota, representing both the end of its first
chapter in the United States and an acknowledgment of its harsh financial reality, Mr.~Womack said.

He said the downturn had presented Toyota with the chance to stop emulating Detroit and to detour to
a more realistic path.

Ms.~Templeton, the Florida dealer, agreed. ``Instead of trying to be the manufacturer for everyone,
they need to get back to being the best,'' she said.

\section{Spurring\cite{spur} Sales, Car Rebate Plan Is Left Up in Air}

\lettrine{T}{he} fate of the ``cash for clunkers'' program remained uncertain on Monday even as
sales figures from automakers demonstrated that people had flocked to dealers to trade in old gas
guzzlers.

The White House urged the Senate to add \$2 billion to the program, as the House voted to do last
Friday before leaving for its August recess. Still, dealers around the country stopped promising the
rebates to car shoppers on Monday, because of uncertainty about how much of the \$1 billion
initially allocated had been used up, or when or whether more money would be available. The Senate
begins its recess this Friday.

The short-term tonic\cite{tonic} of the first billion dollars was evident, though, in sales figures
that automakers reported Monday. New-vehicle sales rose last month to the highest level in nearly a
year, and in the final week of July, cars and trucks were rolling off dealers' lots at almost the
same rate they had before the recession began.

Dealers estimated that they sold a quarter-million cars with the rebate money.

And the Transportation Department reported that the average gas mileage of the vehicles being bought
was significantly higher than required to qualify for a rebate of \$3,500 to \$4,500. Of 120,000
rebate applications processed so far, the department said the average gas mileage of cars being
bought was 28.3 miles per gallon, for S.U.V.'s, 21.9 miles per gallon, and for trucks, 16.3 miles
per gallon.

``The statistics are much better than anybody dreamed they would be,'' said Senator Dianne
Feinstein, Democrat of California. The actual mileage gain so far, she said, was not due to the
details of the law but ``the good judgment of the American people.''

Senator Feinstein, along with Senator Susan M.~Collins, Republican of Maine, was the author of an
early version of a ``cash for clunkers'' bill that would have required bigger improvements in
mileage. Their decision to express support for extending the current version of the program, at a
news conference late Monday afternoon, was an important signal to other senators concerned about
whether the program was doing enough for the environment.

Senator Charles E.~Schumer of New York, an early backer of the Feinstein-Collins approach, also
voiced support for an extension. ``If it ain't broke, don't fix it,'' he said. ``It's working in
every way. It's working as stimulus, it's working to help families, it's working to improve
mileage.''

The White House spokesman, Robert Gibbs, said at his daily briefing that President Obama would lobby
for the extension of the program on Tuesday when the 60 Democratic senators came to the White House
for lunch.

``It's good for dealers and auto manufacturers,'' Mr.~Gibbs said. ``It's good for our energy
security and our environment.''

But the prospects for renewal were still uncertain.

Jim Manley, a spokesman for the Senate majority leader, Harry Reid, said of the statements by the
three senators, ``This is encouraging news, but in the Senate, where any one senator can stall
things, we still need to reach agreement with the Republicans so we can get this through the
chamber.''

The program squeaked through the Senate the first time it came up, passing 60-36, the minimum needed
under Senate rules. In contrast, the House voted 298 to 119 to provide the initial \$1 billion for
the program. On Friday the vote to add \$2 billion more was 316 to 109.

The Senate Republican leader, Mitch McConnell of Kentucky, said Monday on the Senate floor that the
Obama administration had botched the execution of the program by miscalculating how popular it would
be. This was a reason to be more deliberate in acting on a health care overhaul, he said.

Mr.~McConnell did not, however, speak out against adding \$2 billion to the clunker program, though
his spokesman, Don Stewart, said some Republicans had ``real objections to it,'' because it was
borrowing the \$2 billion from one loan program -- at the Energy Department -- and funneling it to
another.

The program, formally known as the Car Allowance Rebate System, was meant to last until Nov. 1 or as
long as the money lasted. But the program basically ran out of money within the first week as
dealers submitted applications to the Transportation Department for a quarter-million cars they had
sold in July.

Auto dealers have been uncertain about how to handle new applications for the program since Thursday
night, when the administration sent mixed signals about the future of the program.

In Houston on Monday, one dealer, George DeMontrond, said, ``We are not going to do anymore until we
get further information.''

Noting the two potential levels of rebates, he said, ``at \$3,500 to \$4,500 a car, that is the
difference between making a little money and losing a lot. We're not prepared to take that risk
until we fully understand what's going to happen next.''

In Detroit, the automakers acknowledged sales were certain to drop off again when the trade-in
program ended -- whether immediately or later this year -- and they conceded that there was a long
way to go before the industry could climb out of its disastrous slump.

At the Ford Motor Company, sales rose 2.3 percent in July, the company's first year-over-year sales
increase in nearly two years and the first for any of the six largest automakers since last August.
Ford's compact sedan, the Focus, was the top-selling vehicle bought by those who turned in a clunker
to be scrapped, the government said.

``We're seeing people that we haven't seen in years and years and years, coming back to where they
bought the car 10 years ago,'' said Patrick Dazzo, the general manager of Joe Rizza Ford Lincoln
Mercury in North Riverside, Ill. ``I had every person in my dealership working nonstop. It was
great. Very encouraging.''

Three other automakers -- Hyundai, Kia and Subaru -- also reported increases. Detroit's other car
companies, Chrysler and General Motors, still sold fewer vehicles than they did a year ago, but the
declines of 9 percent and 19 percent, respectively, were much smaller than in recent months.

No automaker revealed how many of its sales were to people who traded in a clunker, except Hyundai,
which said the program accounted for 22 percent of sales. Ford and others said they did not know
which sales were part of the program but could tell that it was a significant factor based on the
types of vehicles sold, the fact that sales suddenly jumped after July 24 and anecdotal reports from
dealers.

Even so, sales are down 32 percent compared with the first seven months of 2008.

\section{Findings May Explain Gap in Cancer Survival}

\lettrine{S}{cientists} say they have made a discovery that may help explain the racial gap in
cancer survival, providing clues to why white patients often outlive blacks even when they have what
appear to be the same cancers.

The insights come from research at the University of Maryland into throat cancer and squamous-cell
cancers of the head and neck, which have been increasing sharply in recent years, apparently because
of the human papillomavirus -- the same sexually transmitted virus that causes cervical cancer and
is the target of a vaccine for girls.

The virus can also be spread through oral sex, causing cancer of the throat and tonsils, or
oropharyngeal cancer.

The new research builds on earlier work suggesting that throat cancer tumors caused by the virus
behave very differently from other throat cancers, and actually respond better to treatment. And the
new research suggests that whites are more likely than blacks to have tumors linked to the virus,
which may explain the poor outcomes of African-Americans with HPV-negative tumors.

University of Maryland researchers did the study after finding that their white patients with throat
cancer were surviving 70 months on average, compared with 25 months for their black patients, even
though all were treated at the same hospital.

The racial disparity in survival for oropharyngeal cancers explained most of the gap between blacks
and whites for all head and neck cancers, the researchers said.

``We were shocked to see this in our own institution, where more than half of the patients we treat
are African-American,'' said Dr.~Kevin J.~Cullen, director of the Greenebaum Cancer Center at
University of Maryland and senior author of the new study, in the September issue of Cancer
Prevention Research. Around the same time, the Maryland researchers were analyzing specimens of head
and neck tumors gathered from participants in a treatment trial called the TAX 324 study, to
determine how many tumors were linked to HPV.

The results were striking: the TAX 324 patients whose tumors were caused by the virus responded much
better to treatment with chemotherapy and radiation. And they were also overwhelmingly white.

While about one-half of the white patients' throat tumors were HPV-positive, only one of the black
patients had a tumor caused by the virus, Dr.~Cullen said.

``There was no difference in the survival between black and white patients in the TAX 324 trials if
you subtracted out the HPV-positive patients,'' he said.

The racial gap has often been explained as a result of late diagnosis among African-Americans, lack
of access to care and less aggressive treatment, but experts said that in the case of oropharyngeal
cancer, there appeared to be distinct biological differences between the tumors.

This suggests that the racial gap in survival for this particular cancer may trace back to social
and cultural differences between blacks and whites, including different sexual practices, experts
said.

At a briefing for reporters, leading cancer experts called the new report a landmark paper that
would transform the treatment of oropharyngeal cancers and challenge doctors to develop new
treatment options for patients with HPV-negative tumors.

Dr.~Otis Brawley, medical director of the American Cancer Society, who wrote an editorial
accompanying the report, said that changing sexual practices were increasing rates of head and neck
cancers, and perhaps others as well.

``There is a huge public health message here,'' he said.

\section{Doctors in Cuba Start Over in the U.S.}

\lettrine{I}{n} 1991, Carlos Domínguez, a family doctor in one of Havana's poorest neighborhoods,
bought a boat for 12,000 pesos -- the equivalent of saving his entire paycheck for three years -- to
escape the government that had trained him to be an international doctor.

The boat was old and needed to be outfitted with the transmission from a 1952 Ford, one of the many
American cars that still cruise the streets of Havana. The mechanic warned him there was no reverse
gear. The boat could only go forward.

``Perfect,'' Dr.~Domínguez, now 46, said he replied. ``I don't plan on coming back. From now on, I'm
just going forward.''

And so, armed with his grandfather's World War II compass, he left Cuba and made his way to Miami,
rowing the last seven hours after the gasoline ran out. He was 28 years old and ready to resume his
life as a doctor.

But first he needed to pass four exams given only in English, and then put in several years of
training as a hospital resident.

Dr.~Domínguez, who had been taught Russian in his military school in Cuba, knew no English. Still,
he passed one exam before failing the second by a few points. Already married and saddled with
family responsibilities, he put away his medical school books, and signed up for a program to become
a nurse in one year. Since 2001 he has worked as a hospice admissions nurse, a job that allows him
to work with patients while avoiding the hurdles that doctors have to overcome to practice medicine
in the United States.

While the rest of the country is suffering from a shortage of primary care physicians, Miami is
awash\cite{awash} with Cuban doctors who have defected in recent years. By some estimates, 6,000
medical professionals, many of them physicians, have left Cuba in the last six years.

Cuban doctors have been fleeing to South Florida since Fidel Castro seized power in 1959, but the
pace intensified after 2006, when the Department of Homeland Security began a program that allowed
Cuban medical personnel ``who study or work in a third country under the direction of the Cuban
government'' to travel to the United States legally. The program has effectively turned a crowning
achievement of Cuba's foreign policy on its head.

In the 50 years since the revolution, Cuba has sent more than 185,000 health professionals on
medical missions to at least 103 countries. About 31,000, most of them doctors, are in Venezuela,
where they work in exchange for cheap oil and other trade benefits for the Cuban government.

And more are in the pipeline. Cuba's official news agency reported that more than 25,000 health
professionals graduated this year, ``the largest graduation ever.''

But many doctors on the island are now vying to be tapped for an international mission, in part
because they know that no matter where they are sent, they will be one step closer to a visa to the
United States.

The missions have earned Cuba much recognition, goodwill and bargaining power. President Obama told
reporters at the end of a recent hemispheric meeting in Trinidad that he found it ``interesting'' to
learn from Latin American leaders ``about the thousands of doctors from Cuba that are dispersed all
throughout the region, and upon which many of these countries heavily depend.''

Yet for many Cuban doctors, who earn the equivalent of \$25 a month, the lure\cite{lure} of a life
of freedom and opportunities in the United States is too strong to resist. And so these children of
the revolution, educated by a Communist regime to reject capitalism and embrace socialism, have
ended up in Miami, often tending to elderly Cubans who fled the island before the doctors were born.

Ana Carbonell, chief of staff for Representative Lincoln Diaz-Balart, Republican of Florida, said
more than 2,000 Cubans had already settled in the United States under the parole program.

``It brings to our community highly qualified professionals at a time of great need,'' Ms.~Carbonell
said. ``They work alongside U.S.-trained doctors, and they enhance any practice or wherever they
work.''

Many have been able to obtain licenses and practice medicine. Others have chosen to settle for
careers in the medical field but not as physicians, and some work in fields that have nothing to do
with medicine.

``I know neurosurgeons who are working in warehouses or factories or as gas attendants,'' said Julio
C\'esar Alfonso, 40, who graduated from medical school in Cuba in 1992 and works as a clinic manager
in Miami. ``But I know many more who are working as nurses, medical assistants and technicians.''

Lianete P\'erez, 37, works as a medical assistant in the office of a pediatrician in Miami. A former
anesthesiologist, Dr.~P\'erez longed to leave Cuba, arrived in 2002 and is studying to take the
medical exams later this year. Unlike other doctors who resent having their skills tested years
after medical school, she said she welcomed the chance to go back to the books.

``There are enormous differences between medicine in Cuba and in the United States,'' she said. ``I
can't tell you that Cuban doctors are not well trained, but I can tell you that the books we used
were edited in 1962, and for me, coming here was like starting all over again.''

Aside from old books, Cuban medical students and doctors must contend with a lack of modern
equipment and, often, of drugs and diagnostic tools taken for granted in developed countries. But
many expatriate doctors say their dealings with patients in Cuba were more humane and less rushed
than they are in the United States.

Foreign doctors trained in languages other than English face immense challenges getting a license to
practice in the United States. Not only must they relearn their profession in English, but many,
like Dr.~Domínguez, must also work to support themselves and their families. Cuban doctors, in
particular, tend to be older by the time they arrive in the United States, sometimes too old to
dedicate years to studying for exams and finding and completing a residency program.

Finding a residency program is almost impossible, they say, in part because most hospitals seek
young doctors who can work long hours for little pay. A 40-year-old doctor is often deemed too old
to be accepted in a residency program. Dr.~J.~Donald Temple, an associate professor of medicine at
the University of Miami, said foreign-trained doctors faced the widespread perception that medical
schools abroad are not as rigorous as medical schools in this country.

``They are not going to be accepted as readily as American-trained doctors are for some of the more
competitive residency programs,'' said Dr.~Temple, who runs a training program for Latin American
and Caribbean doctors to become leaders in the medical field in their own countries. ``Many programs
would not even consider them.''

Their loss is Miami's gain. Just about any hospital in Miami can boast of having highly trained
Cuban doctors working as nurses or medical assistants or even lab technicians.

Carlos P\'erez-Sedano, 40, is one of them. He said he was sent to Ghana in September 2005. A year
later, just two months after the parole program was announced, he applied for a United States visa.
By December 2006, he was in Miami.

He now works as a case manager at Kendall Regional Hospital and is studying to take the exams for
his Florida medical license. He says that half of the case managers in his department are Cuban
doctors, and that at least 50 Cuban doctors work in other areas of the hospital, not all of them as
physicians. Hospital officials at Kendall Regional say they keep no record of the nationality of
their employees.

One out of every four doctors in the United States was trained overseas, said an American Medical
Association official who deals with international doctors, but who is not authorized to speak to
reporters and declined to give her name.

``We feel that international medical graduates are an integral part of our work force, and we don't
see that changing any time soon,'' she said. ``We rely very heavily on them because they take
positions that are in underserved areas and go into primary care. Without them our country would be
in dire straits, especially minority populations and the poor.''

This is the path that Dr.~Juan A.~Bereao took when he came from Cuba 11 years ago. A surgeon in
Cuba, he now works as a house doctor in several South Florida hospitals, including Kendall. For
three years before that, he was a resident in internal medical at Bronx-Lebanon Hospital in New York
City, where most of his patients were poor and from ethnic or racial minorities.

Dr.~Domínguez, the hospice nurse, said he was mostly at peace with the fact that no one calls him
``Doctor'' anymore, except some old patients from Cuba he sometimes run into in Miami. He earns more
than \$100,000 a year, sends his two children to private schools and vacations with his family in
France and Spain. But he said a part of him would always pine for the profession -- and the title.

``I've had to get used to think as a nurse, but it's difficult,'' he said. ``Deep down, I'm still a
doctor.''

\section{Divorce, It Seems, Can Make You!}

\lettrine{M}{arried} people tend to be healthier than single people. But what happens when a
marriage ends?

New research shows that when married people become single again, whether by divorce or a spouse's
death, they experience much more than an emotional loss. Often they suffer a decline in physical
health from which they never fully recover, even if they remarry.

And in terms of health, it's not better to have married and lost than never to have married at all.
Middle-age people who never married have fewer chronic health problems than those who were divorced
or widowed.

The findings, from a national study of 8,652 men and women in their 50s and early 60s, suggest that
the physical stress of marital loss continues long after the emotional wounds have healed. While
this does not mean that people should stay married at all costs, it does show that marital history
is an important indicator of health, and that the newly single need to be especially vigilant about
stress management and exercise, even if they remarry.

``When your spouse is getting sick and about to die or your marriage is getting bad and about to
die, your stress levels go up,'' said Linda Waite, a sociology professor at the University of
Chicago and an author of the study, which appears in the September issue of The Journal of Health
and Social Behavior. ``You're not sleeping well, your diet gets worse, you can't exercise, you can't
see your friends. It's a whole package of awful events.''

The health benefits of marriage, documented by a wealth of research, appear to stem from several
factors. Married people tend to be better off financially and can share in a spouse's employer
health benefits. And wives, in particular, act as gatekeepers for a husband's health, scheduling
appointments and noticing changes that may signal a health problem. Spouses can offer logistical
support, like taking care of children while a partner exercises or shuttling a partner to and from
the doctor's office.

But in the latest study, researchers sought to gauge the health effects of divorce, widowhood and
remarriage in a large cohort of people over time.

Among the 8,652 people studied, more than half were still married to their first spouse. About 40
percent had been divorced or widowed; about half of that group were remarried by the time of the
study. About 4 percent had never married.

Over all, men and women who had experienced divorce or the death of a spouse reported about 20
percent more chronic health problems like heart disease, diabetes and cancer, compared with those
who had been continuously married. Previously married people were also more likely to have mobility
problems, like difficulty climbing stairs or walking a meaningful distance.

While remarrying led to some improvement in health, the study showed that most married people who
became single never fully recovered from the physical declines associated with marital loss.
Compared with those who had been continuously married, people in second marriages had 12 percent
more chronic health problems and 19 percent more mobility problems. A second marriage did appear to
heal emotional wounds: remarried people had only slightly more depressive symptoms than those
continuously married.

The study does not prove that the loss of a marriage causes health problems, only that the two are
associated. It may be that people who don't exercise, eat poorly and can't manage stress are also
more likely to divorce. Still, researchers note that because the effect is seen in both divorced and
widowed people, the data strongly suggest a causal relationship.

One reason may be changes at the cellular level during times of high stress. In an Ohio State
University study, scientists analyzed blood samples of people undergoing the stress of caring for a
loved one with Alzheimer's disease. The research focused on telomeres, which insulate and protect
the ends of chromosomes; with aging, telomeres shorten and the activity of a related enzyme also
declines.

Compared with a control group, the Alzheimer's caregivers showed telomere patterns associated with a
four- to eight-year shortening of life span. Dr.~Waite said the stress of divorce or widowhood might
take a similar toll, leading to chronic health and mobility problems.

None of this suggests that spouses should stay in a bad marriage for the sake of health. Marital
troubles can lead to physical ones, too.

In a series of experiments, scientists at Ohio State studied the relationship between marital strife
and immune response, as measured by the time it takes for a wound to heal. The researchers recruited
married couples who submitted to a small suction device that left eight tiny blisters on the arm.
The couples then engaged in different types of discussions -- sometimes positive and supportive, at
other times focused on a topic of conflict.

After a marital conflict, the wounds took a full day longer to heal. Among couples who exhibited
high levels of hostility, the wound healing took two days longer than with those who showed less
animosity.

``I would argue that if you can't fix a marriage you're better off out of it,'' said Janice
Kiecolt-Glaser, an Ohio State scientist who is an author of much of the research. ``With a divorce
you're disrupting your life, but a long-term acrimonious marriage also is very bad.''

\section{Senate Is Expected to Extend `Clunker' Trade-Ins}

\lettrine{T}{he} secretary of transportation and the Senate majority leader both expressed
confidence Tuesday that the Senate would revive the ``cash for clunkers'' program by voting to
infuse it with another \$2 billion, as the House did last week.

The Senate minority leader said the Senate would hold a vote, although he did not predict the
result.

The Senate is scheduled to go home on Friday and simply getting the bill to the floor before then
requires a consensus\cite{consensus} among the senators. If the bill were amended\cite{amend}, it
would effectively kill the program until Labor Day, since the House has already begun its August
recess and would be unlikely to return to approve the Senate version.

But administration officials and Senate Democratic leaders said the Senate would approve the House
version, giving more Americans a chance to get rebates by trading in gas guzzlers for more
fuel-efficient vehicles.

``There obviously is a real pent-up demand in America,'' said Ray LaHood, the transportation
secretary. ``People love to buy cars, and we've given them the incentive to do that. I think the
last thing that any politician wants to do is cut off the opportunity for somebody who's going to be
able to get a rebate from the government to buy a new automobile.''

Harry Reid, Democrat of Nevada and the Senate majority leader, said, ``We'll pass cash for
clunkers.'' And Mitch McConnell, Republican of Kentucky, who is the minority leader, said there
would be a vote, but did not say how it would come out.

Even Senator Jim DeMint, Republican of South Carolina, a vehement\cite{vehement} opponent of the
program, said he would not block a vote.

But it was not clear Tuesday evening when the Senate might actually take up the issue. It was
scheduled to take up the confirmation of Judge Sonia Sotomayor for a Supreme Court seat on Wednesday
morning.

Mr.~LaHood said that the clunkers program was still operating but would run out of money at the end
of the week without a fresh infusion. Dealers around the country, though, said they had stopped
making deals under the program because they were not confident that there would be any money left
from the original \$1 billion allocated by Congress. Dealers have mostly been advancing the amount
of the rebate, \$3,500 to \$4,500, to their customers and waiting for repayment from Washington.

The dealers said they believed that about 250,000 cars had been sold under the clunkers program, but
data entry has lagged sales.

Mr.~LaHood said on Tuesday that 157,000 transactions had been entered in the computer system, with a
rebate value of \$664 million. The predominant transaction so far has been to trade in a pickup,
S.U.V. or van for a sedan\cite{sedan}, he said, with the result that the new vehicles go 61 percent
further on a gallon than the old vehicles did.

Some senators had argued that the terms of the program were too lenient, and that the government
should specify tighter mileage standards. But the actual fuel economy improvement between the
vehicles junked and their replacements were far better than the law required.

Another concern is when to end the program. Congress had expected the first billion dollars to last
for months, but it was consumed in about one month. Senator Ben Nelson, Democrat of Nebraska, a
fiscal conservative, said: ``I don't see how you ever shut it down. You just keep feeding the kitty
and just keep putting more money into it unless you have a way to end it.''

\section{Giant Particle\cite{particle} Collider Struggles}

\lettrine{T}{he} biggest, most expensive physics machine in the world is riddled\cite{riddle} with
thousands of bad electrical connections.

Many of the magnets meant to whiz high-energy subatomic particles around a 17-mile underground
racetrack have mysteriously lost their ability to operate at high energies.

Some physicists are deserting the European project, at least temporarily, to work at a smaller,
rival machine across the ocean.

After 15 years and \$9 billion, and a showy ``switch-on'' ceremony last September, the Large Hadron
Collider, the giant particle accelerator outside Geneva, has to yet collide any particles at all.

But soon?

This week, scientists and engineers at the European Center for Nuclear Research, or CERN, are to
announce how and when their machine will start running this winter.

That will be a Champagne moment. But scientists say it could be years, if ever, before the collider
runs at full strength, stretching out the time it should take to achieve the collider's main goals,
like producing a particle known as the Higgs boson thought to be responsible for imbuing other
elementary particles with mass, or identifying the dark matter that astronomers say makes up 25
percent of the cosmos.

The energy shortfall could also limit the collider's ability to test more exotic ideas, like the
existence of extra dimensions beyond the three of space and one of time that characterize life.

``The fact is, it's likely to take a while to get the results we really want,'' said Lisa Randall, a
Harvard physicist who is an architect of the extra-dimension theory.

The collider was built to accelerate protons to energies of seven trillion electron volts and smash
them together in search of particles and forces that reigned earlier than the first trillionth of a
second of time, but the machine could run as low as four trillion electron volts for its first year.
Upgrades would come a year or two later.

Physicists on both sides of the Atlantic say they are confident that the European machine will
produce groundbreaking science -- eventually -- and quickly catch up to an American rival, even at
the lower energy. All big accelerators have gone through painful beginnings.

``These are baby problems,'' said Peter Limon, a physicist at the Fermi National Accelerator
Laboratory in Batavia, Ill., who helped build the collider.

But some physicists admit to being impatient. ``I've waited 15 years,'' said Nima Arkani-Hamed, a
leading particle theorist at the Institute for Advanced Study in Princeton. ``I want it to get up
running. We can't tolerate another disaster. It has to run smoothly from now.''

The delays are hardest on younger scientists, who may need data to complete a thesis or work toward
tenure. Slowing a recent physics brain drain from the United States to Europe, some have gone to
work at Fermilab, where the rival Tevatron accelerator has been smashing together protons and
antiprotons for the last decade.

Colliders get their oomph from Einstein's equivalence of mass and energy, both expressed in the
currency of electron volts. The CERN collider was designed to investigate what happens at energies
and distances where the reigning theory, known as the Standard Model, breaks down and gives nonsense
answers.

The collider's own prodigious energies are in some way its worst enemy. At full strength, the energy
stored in its superconducting magnets would equal that of an Airbus A380 flying at 450 miles an
hour, and the proton beam itself could pierce 100 feet of solid copper.

In order to carry enough current, the collider's magnets are cooled by liquid helium to a
temperature of 1.9 degrees above absolute zero, at which point the niobium-titanium cables in them
lose all electrical resistance and become superconducting.

Any perturbation, however, such as a bad soldering job on a splice, can cause resistance and heat
the cable and cause it to lose its superconductivity in what physicists call a ``quench.'' Which is
what happened on Sept.~19, when the junction between two magnets vaporized in a shower of sparks,
soot and liberated helium.

Technicians have spent most of the last year cleaning up and inspecting thousands of splices in the
collider. About 5,000 will have to be redone, Steve Myers, head of CERN's accelerator division, said
in an interview.

The exploding splices have diverted engineers' attention from the mystery of the underperforming
magnets. Before the superconducting magnets are installed, engineers ``train'' each one by ramping
up its electrical current until the magnet fails, or ``quenches.'' Thus the magnet gradually grows
comfortable with higher and higher current.

All of the magnets for the collider were trained to an energy above seven trillion electron volts
before being installed, Dr.~Myers said, but when engineers tried to take one of the rings' eight
sectors to a higher energy last year, some magnets unexpectedly failed.

In an e-mail exchange, Lucio Rossi, head of magnets for CERN, said that 49 magnets had lost their
training in the sectors tested and that it was impossible to estimate how many in the entire
collider had gone bad. He said the magnets in question had all met specifications and that the
problem might stem from having sat outside for a year before they could be installed.

Retraining magnets is costly and time consuming, experts say, and it might not be worth the wait to
get all the way to the original target energy. ``It looks like we can get to 6.5 relatively
easily,'' Dr.~Myers said, but seven trillion electron volts would require ``a lot of training.''

Many physicists say they would be perfectly happy if the collider never got above five trillion
electron volts. If that were the case, said Joe Lykken, a Fermilab theorist who is on one of the
CERN collider teams, ``It's not the end of the world. I am not pessimistic at all.''

For the immediate future, however, physicists are not even going to get that. Dr.~Myers said he
thought the splices as they are could handle 4 trillion electron volts.

``We could be doing physics at the end of November,'' he said in July, before new vacuum leaks
pushed the schedule back a few additional weeks.

``It's not the design energy of the machine, but it's 4 times higher than the Tevatron,'' he said.

Pauline Gagnon, an Indiana University physicist who works at CERN, said she would happily take that
energy level. ``The public pays for this,'' she said in an e-mail message, ``and we need to start
delivering.''

\section{Korean Police Raid Occupied Factory}

\lettrine{F}{iring} water cannons, the police raided Ssangyong Motor Company on Tuesday to reclaim
the plant, which had been occupied for more than two months by hundreds of workers resisting
layoffs.

After a day of battles, with workers firing nuts and bolts with slingshots and police helicopters
showering workers from above with tear-gas-laced water, the police gained control of all but one of
the factory buildings, where some 500 workers continued to make a last stand.

The workers also hurled fire bombs from rooftops. At least 23 police officers were injured, the
police said, while the union said ``many'' workers were inside the factory were also injured.

Dark smoke billowed from the factory in Pyeongtaek, 45 miles south of Seoul. Workers burned tires
when thousands of riot police officers moved in behind plastic shields. Police helicopters hovered
overhead, throwing water laced with tear gas on the workers.

After hours of fighting, the police occupied several buildings but retreated from a paint shop,
where at least 520 workers made their last stand. The building is stocked with
inflammable\cite{inflammable} material.

Ssangyong, the country's fifth-largest carmaker, filed for bankruptcy in January amid falling sales
and mounting debts.

The company has also had an acrimonious\cite{acrimonious} relationship with its Chinese owner, the
Shanghai Automotive Industry Corporation, which bought a controlling share of the business in 2004.
Shanghai Auto, China's largest carmaker, was seeking to expand overseas at a time when South Korea
was beginning to embrace China as an economic partner.

As the company's fortunes waned\cite{wane}, management laid off hundreds of workers in 2006, leading
to cross-national recriminations and a strike in which workers took over the plant and locked out
management for nearly two months.

In December, union members held several Shanghai Auto officials hostage for seven hours, accusing
them of absconding\cite{abscond} with proprietary technology.

In April, the company announced plans to shed 36 percent of its workforce. About 900 employees
barricaded themselves inside the paint shop and other buildings on May 22.

As the police siege on the factory persisted, many workers broke ranks and deserted the factory. But
the protesters rejected the management's latest offer to reduce the number of layoffs. They said in
a statement that rather than being divided, they would ``die together.''

Ssangyong threatened to seek liquidation if the workers prolong the occupation.

\section{Opposition Woos\cite{woo} Japan's Voters With Costly Vows}

\lettrine{I}{n} a recent YouTube video posted by Japan's governing party, a smooth talker with an
uncanny resemblance to the country's main opposition leader, woos his date with sweet promises: a
life without worries about child care costs or retirement, if only she will marry him.

The wide-eyed woman asks how he will pay for such a carefree lifestyle.

``Just choose me,'' the suitor snaps. ``I'll sort out the details after we're married.''

The not-so-subtle message from the governing Liberal Democratic Party is that the opposition
Democratic Party of Japan is offering voters much the same kinds of promises. The voters have seemed
ready to be wooed.

The opposition party is riding a surge in popularity before elections for the lower house of
Parliament on Aug.~30. Its platform promises more cash in hand for Japanese struggling through a
deep recession: cash allowances to families with children, free tuition and more social security --
all without adding to the country's growing debt.

The Liberals have called the plan nonsensical and sure to bankrupt Japan. Instead, they promise to
stimulate the economy by investing in strategic industries, like solar power and ``green'' cars.
Growth will then quickly trickle down to the common people, the Liberals say.

The opposition ``focuses on giving out money with no strategy for economic growth,'' Prime Minister
Taro Aso said of the Democrats. ``We can no longer afford to pass our debt onto our children and
grandchildren.''

The two parties are, in essence, pushing two versions of how Japan, the world's second-largest
economy, should work: as it has for the last half-century, by propping up its corporations, or by
distributing wealth to the people.

If voters embrace the Democrats, Japan's economy could shift from a long reliance on exports for
growth, a model that has been largely discredited in the recent collapse of exports in the global
downturn. Instead, a Democratic Party victory could result in a redistribution of wealth to ordinary
Japanese to spur consumer spending.

``We will not pursue a traditional growth model -- that is, rely on exports to grow our economy,''
Katsuya Okada, secretary general of the Democratic Party of Japan, told reporters on Monday. ``First
and foremost, we need to improve peoples' livelihoods.''

But the party's smorgasbord of policies -- which include yearly allowances of 312,000 yen (\$3,297)
per child, free tuition at public high schools and an end to highway tolls and gasoline tax -- will
not come cheap.

The Democrats say the program will cost 7.1 trillion yen to start next April, and an additional 16.8
trillion yen in 2011. That comes to more than 3 percent of Japan's 510-trillion-yen economy.

With its public debt already surpassing 180 percent of gross domestic product, Japan cannot afford
such spending. By some measures, its debt-to-G.D.P. ratio is second only to Zimbabwe's.

But the party's proposals have proved attractive. The Kumano family in central Tokyo, for instance,
would receive 624,000 yen (\$6,593) in child benefits a year for its two boys, now 8 and 9, until
they are about 16. The father, Eiichi Kumano, who runs a child care business, says the money would
be a relief.

``The boys are getting older, and we have to start thinking about college,'' Mr.~Kumano said. ``In
the short term, we're thankful for the cash. But I do worry about how the government will pay for
all this in the long run.''

The Democrats say they intend to finance their program by eliminating the governing party's wasteful
public spending, a legacy of decades of pork-barrel politics.

``We intend to rebuild the entire budget from scratch,'' said Masaharu Nakagawa, the Democrats'
shadow finance minister. ``We will eliminate all that's wasteful.'' Economists agree that the
opposition party would have a freer hand in uprooting vested interests and writing a more efficient
budget. The party has, until now, been virtually ignored by corporate Japan. In 2007, Keidanren, the
country's main business federation, contributed almost 2.9 billion yen to the governing Liberals,
and just 80 million yen to the Democrats.

``A growth strategy that focuses on ordinary citizens could be better than current policy, which is
constrained by longstanding business interests,'' said Ryutaro Kono, chief economist for Japan at
BNP Paribas.

The opposition is also exploring raising money by eliminating generous tax breaks and subsidies that
have long supported weak Japanese companies at the expense of the wider economy. About two-thirds of
Japanese companies, and half of large companies, do not pay corporate taxes because of obscure tax
subsides for research, investment incentives or because they lose money, according to the
Organization for Economic Cooperation and Development.

By changing corporate tax policies and tapping idle government funds, the opposition party hopes to
avoid raising Japan's 5 percent consumption tax. The Liberals insist it will soon be necessary to
raise that tax to meet government expenditures.

Businesses say the Democrats' proposals would hobble Japanese companies in international
competition. Executives are already stepping up the pressure on the opposition to ease its tough
stance on corporations, which includes much tighter emissions goals for greenhouse gases.

``The economy will stall, or even worse -- backslide,'' Tetsuya Nomura, president of the
construction company Shimizu Corporation and the leader of the industry's main lobbying group, told
reporters. ``None of us can put food on the table,'' he said.

Some analysts question the ability of the opposition party, which has never been in office, to
navigate Japan's complex economy. Formed a decade ago, the party is a mix of lawmakers from a
spectrum of political views, including socialists and free-market advocates. Its agenda is likely to
face fierce resistance from Japan's powerful bureaucrats, whom the Democrats have vilified as
keepers of the status quo.

Economists note that neither party has a viable plan for long-term growth. With the elections for
Parliament's more powerful lower chamber looming, neither side has addressed some of the more
painful changes that Japan must make to build a more productive economy, including bringing more
competition to domestic industries, revising rigid labor laws and increasing efficiency in financial
markets.

\section{ESPN Limits Social Networking}

\lettrine{S}{oon} after ESPN issued 12 guidelines to its employees about social networking on
Tuesday, Ric Bucher, one of its N.B.A. writers and analysts, wrote on Twitter, ``The hammer just
came down, tweeps: ESPN memo prohibiting tweeting info unless it serves ESPN.''

He then added, ``My guess is I can still tweet about my vacation/car shopping, etc.''

The guidelines are more detailed than Bucher described them. But they restrict the freedom that ESPN
employees might previously have enjoyed.

``We've been in the social networking space for a long time, and will continue to be there,'' said
Chris LaPlaca, an ESPN spokesman. ``But we want to be smarter about how we do it.'' He said that
Bucher's ``interpretation of the policy is mistaken.''

The guidelines say that on-air talent, reporters and writers are prohibited from having
sports-related blogs or Web sites and that they will need a supervisor's approval to discuss sports
on any social networking sites. They will also be restricted from discussing internal policies or
detailing how stories are ``reported, written, edited or produced.''

The guide that Bucher focused on reads, ``The first and only priority is to serve ESPN-sanctioned
efforts, including sports news, information and content.'' Violating the new guidelines could lead
to suspension or dismissal.

ESPN will soon deploy modules to simultaneously link ESPN employees' tweets and other social
networking feeds to various ESPN's Web sites. ``The key phrase is write it once, publish it
everywhere,'' LaPlaca said. He said the policy is ``not meant to diminish our having conversations
with fans, but to do them so that if you're not on Facebook, you can still get the information.''

\section{Hot Story to Has-Been: Tracking News via Cyberspace}

\lettrine{L}{ike} a lot of new ideas, Media Cloud started with a long-running argument among
friends. Ethan Zuckerman and a handful of his colleagues at the Berkman Center for Internet and
Society at Harvard Law School found themselves in endless disputes about the mainstream media and
newer digital variations. Who sets the agenda? How is public debate shaped? What topics are covered
or ignored?

Anecdotes favoring one side or another were as plentiful as pop-ups, but a comprehensive and
reliable database that could track the daily rhythm of the news cycle over time and was available
for public use didn't exist. So Mr.~Zuckerman and others at Berkman decided to create one.

The result is Media Cloud, a system that tracks hundreds of newspapers and thousands of Web sites
and blogs, and archives the information in a searchable form. The database, at mediacloud.org, will
eventually enable researchers to search for key people, places and events -- from Michael Jackson to
the Iranian elections -- and find out precisely when, where and how frequently they are covered,
said Mr.~Zuckerman, whose official title is senior researcher, though he acknowledges that a more
accurate label would be computer geek and international development specialist. (At the moment only
a small sample of Media Cloud's tools are on the public Web site.)

The findings, which can be graphed or mapped, can demonstrate the evolution of a report and
variations in coverage. Users get to ``do the fun part, which is analyzing the data,'' Mr.~Zuckerman
said, ``while we do the hard part of this, which is collecting it.'' Eventually users will be able
to compare the top 10 news events covered by Fox News, The Atlanta Journal-Constitution and the BBC,
for example, or chart the terms that appear most frequently in The New York Times, compared with
leading blogs, or create a world map showing which countries receive the most media attention, or
follow the path of a particular report to see if it dominates the news or dies out.

For the past decade or so, many researchers have used link analysis to figure out how information
spreads, said Yochai Benkler, a Harvard Law School professor at Berkman who has been involved in
creating Media Cloud. You could identify which Web sites were linked to most frequently and infer
whose sites were most influential. But researchers have pretty much squeezed all that they can from
that approach, Mr.~Benkler said. Although Media Cloud is still in its early stages, it is among
``the next generation of tools that actually look at what people are saying,'' he said, adding, it
is ``a better microscope.''

There are other kinds of media trackers. Cornell University researchers, for example, have developed
MemeTracker, which maps the daily news cycle by grabbing repeated quotations from one million online
sources. (A meme is anything -- an idea, a phrase -- that spreads by imitation from one person to
another.)

Its graphs, which can be viewed at memetracker.org, display the reports that are competing against
one another for attention on a given day, as well as those that have staying power or quickly
disappear. A recent paper on MemeTracker's experience during the presidential campaign was hailed by
experts as a landmark piece of work.

Pew Research Center's Project for Excellence in Journalism offers a news-coverage index, which is
laboriously compiled by having 14 people sample leading reports produced by 55 outlets. Media Cloud
is much less exact, Mr.~Zuckerman said, but it can automatically scan hundreds, and eventually
thousands, of sources.

Amy S.~Mitchell, deputy director of Pew's journalism project, said Media Cloud ``offers the public a
great opportunity to play around with looking at a wide swath of media at more of a surface level.''
But, she added, it cannot really capture the nuances of the news agenda of the news media. ``There
are certain things that computer algorithms cannot do that individuals can,'' she said.

Since every method has virtues and drawbacks, she added, ``I think there is tremendous value in
having both approaches.''

What Media Cloud offers that no one else does is a tool anyone can use to answer all sorts of
questions about the media landscape. One topic that Mr.~Benkler and Mr.~Zuckerman have long been
debating is whether the Internet has helped open up the public sphere to more voices, or whether it
just serves as an echo chamber, simply repeating information and views that the mainstream media
already circulate.

Mr.~Benkler is using Media Cloud to test his theory that digital media is widening the circle of
voices somewhat. Sites that he characterizes as ``one link out'' from the most visible (like The
Huffington Post, Talking Points Memo and Instapundit) are entering into the conversation, he argues.

Who has the power to place an idea on the national agenda is another question that Mr.~Benkler said
Media Cloud could help answer. For instance, how is the conversation about the recession and the
financial crash shaped? Using some of the database's more specialized tools, Mr.~Benkler
investigated who first floated the idea for a temporary takeover of the financial system by the
government, as was done in Sweden in the 1990s.

Paul Krugman, a columnist for The New York Times, first raised the idea in September 2008, but it
was a cluster of influential economic and political bloggers like Brad DeLong and Matthew Yglesias
who kept the idea alive. After the subject disappeared for a couple of months, the bloggers then
resurrected it early this year as Washington began discussing the details of a bank rescue plan. To
Mr.~Benkler, this preliminary evidence suggests that the network of public media has given a voice
to some people who in the past may have had useful ideas but were, as a practical matter, unable to
inject them into the national conversation.

``If you're actually trying to map where an idea starts and how it moves through the public sphere,
you need a database like we're developing, with time-stamped data,'' Mr.~Benkler said, explaining
that services like Google and Lexis/Nexis are not as comprehensive or do not provide that level of
detailed information. Media Cloud also enables more fine-grained analyses by examining language and
context. ``How does rhetoric change over time, and what's the role of the Internet and the
mainstream media in that?'' Mr.~Zuckerman asked.

Some of his colleagues, for example, have been tracking the frequency of the words bailout and
stimulus to pinpoint when one term overtook the other. Media Cloud mapped the results to show how
the term bailout, used constantly in the news in the fall, eventually gave way to the word stimulus
after President Obama took office. The results were graphed, illustrating precisely when the two
lines crossed -- where, as Mr.~Zuckerman would say, ``one meme took over from the other.''

Media Cloud's founders have put out an open call on their Web site for research ideas. The system
provides a platform, Mr.~Zuckerman said, on which others can build using their own kinds of tracking
software. The point is to start with facts rather than impressions.

As Mr.~Zuckerman noted, a lot of anecdotes don't necessarily add up to the truth.

\section{In Release of Journalists, Both Clintons Had Key Roles}

\lettrine{F}{ormer} President Bill Clinton left North Korea on Wednesday morning after a dramatic
20-hour visit, in which he won the freedom of two American journalists, opened a diplomatic channel
to North Korea's reclusive government and dined with the North's ailing leader, Kim Jong-il.

Mr.~Clinton departed from Pyongyang, the capital, around 8:30 a.m. local time, along with the
journalists, Laura Ling, 32, and Euna Lee, 36, on a private jet bound for Los Angeles, according to
a statement from the former president's office.

The North Korean government, which in June sentenced the women to 12 years of hard labor for
illegally entering North Korean territory, announced hours earlier that it had pardoned the women
after Mr.~Clinton apologized to Mr.~Kim for their actions, according to the North Korean state
media.

President Obama contacted the families of the women on Tuesday evening, according to administration
officials, but the White House said it would withhold public comment until the former president
landed on American soil.

Mr.~Clinton's mission to Pyongyang was the most visible by an American in nearly a decade. It came
at a time when the United States' relationship with North Korea had become especially chilled, after
North Korea's test of its second nuclear device in May and a series of missile launchings.

It ended a harrowing ordeal for the two women, who were stopped on March 17 by soldiers near North
Korea's border with China while researching a report about women and human trafficking. They faced
years of imprisonment in the gulaglike confines of a North Korean prison camp.

And it catapulted Mr.~Clinton back on to the global stage, on behalf of a president who defeated his
wife, Hillary Rodham Clinton, in a bitter primary campaign last year, and who later asked her to be
his secretary of state.

Mrs.~Clinton was deeply involved in the case, too. She proposed sending various people to Pyongyang
-- including Mr.~Clinton's vice president, Al Gore -- to lobby for the release of the women, before
Mr.~Clinton emerged as the preferred choice of the North Koreans, people briefed on the talks said.

About 10 days ago, these people said, Mr.~Gore, who co-founded Current TV, the San Francisco-based
media company that employs Ms.~Ling and Ms.~Lee, called Mr.~Clinton to ask him to undertake the
trip. Mr.~Clinton agreed, as long as the Obama administration did not object.

The riveting\cite{rivet} tableau\cite{tableau} of a former president, jetting into a diplomatic
crisis while his wife was embarking on a tour of Africa in her role as the nation's chief diplomat,
underscored the unique and enduring role of the Clintons, even in the Obama era.

On Wednesday in Nairobi, Kenya, Mrs.~Clinton said, ``I'm very happy and relieved to have these two
young women, Laura Ling and Euna Lee, on their way home to their families.''

``I spoke to my husband on the airplane and everything went well; we are extremely excited they will
be reunited,'' she said, adding, ``it was just a good day to be able to see this happen.'' She said
she would have more to say on the matter later, when the journalists and her husband landed in
America.

Mr.~Clinton's trip to Pyongyang came just two weeks after North Korea issued a harsh personal attack
on Mrs.~Clinton, in response to comments she made comparing its nuclear test and missile launchings
to the behavior of an attention-seeking teenager.

The North Korean Foreign Ministry objected to her ``vulgar remarks'' and called her ``a funny lady''
who was neither intelligent nor diplomatic. ``Sometimes she looks like a primary-school girl and
sometimes a pensioner going shopping,'' a spokesman said.

The episode evidently did not stop consideration of sending her husband as an envoy\cite{envoy}. But
the initiative was cloaked\cite{cloak} in secrecy and came after weeks of back-channel talks between
the United States and North Korea through its United Nations mission. In addition to Mr.~Gore, the
White House's list of potential candidates included Gov. Bill Richardson of New Mexico.

North Korea signaled its desire to have Mr.~Clinton act as a special envoy in conversations with
Ms.~Ling and Ms.~Lee, who relayed that message to their families in the middle of July, according to
a senior administration official. The message was passed to Mr.~Gore, who contacted the White House,
which then explored whether such a mission would be successful.

Mr.~Obama did not speak directly with Mr.~Clinton before the mission. But his national security
adviser, Gen. James L.~Jones, contacted the former president to sound him out. The senior official
said the administration did ``due\cite{due} diligence'' with the North Koreans to ensure that if
Mr.~Clinton went, he would return with the journalists. He also denied that Mr.~Clinton apologized
as a condition of obtaining the pardons from the government.

As president, Mr.~Clinton had sent Mr.~Kim a letter of condolence\cite{condolence} on the death of
his father, Kim Il-sung, according to a former official. For Mr.~Kim, the former official said,
freeing the women was a ``reciprocal humanitarian gesture.''

Mr.~Kim is believed to have suffered a stroke last year. American officials said they thought his
declining health had set off a succession struggle, complicating the Obama administration's dealings
with the North.

The families of the American journalists issued a statement saying they were ``overjoyed'' by news
of the pardon and thanked Mr.~Obama and Mrs.~Clinton. ``We especially want to thank President Bill
Clinton for taking on such an arduous\cite{arduous} mission and Vice President Al Gore for his
tireless efforts to bring Laura and Euna home,'' the statement said.

Current TV said in a statement that it too was ``overjoyed'' and that the hearts of its employees
went out to Ms.~Ling and Ms.~Lee for ``persevering through this horrible experience.''

The Obama administration said Ms.~Ling and Ms.~Lee were in good health.

Administration officials said Mr.~Clinton went to North Korea as a private citizen, did not carry a
message from Mr.~Obama for Mr.~Kim and had the authority to negotiate only for the women's release.

``This was 100 percent about the journalists,'' said a senior administration official, who spoke on
the condition of anonymity because of the delicacy of the matter. ``We knew Kim Jong-il would
probably seek a meeting with Clinton. But that's not what this visit was about.''

Still, North Korea, clearly seeing a propaganda opportunity at home and a rare chance for a measure
of favorable publicity abroad, welcomed Mr.~Clinton with the fanfare of a state visit. It broadcast
a group portrait, as well as photos of Mr.~Kim gesturing and talking to Mr.~Clinton; of the former
president accepting flowers from a North Korean girl; and of Mr.~Clinton, seated across a
negotiating table from Mr.~Kim, each flanked by their aides. Among those greeting Mr.~Clinton at the
airport was Kim Kye-gwan, North Korea's chief nuclear negotiator.

Among those accompanying Mr.~Clinton was David Straub, a former director of the Korea desk at the
State Department, who had held talks with the North Koreans through what is known as the ``New York
connection.''

Also on hand was John Podesta, an informal adviser to the Obama administration who served as
Mr.~Clinton's chief of staff in the final years of his presidency, when the former president yearned
to travel to North Korea to clinch a deal that would have curbed its nuclear program.

That visit never happened -- partly because the White House concluded that a deal was not assured --
and President George W.~Bush put the brakes on direct talks with North Korea, setting the stage for
eight years of largely fruitless efforts to stop the North's nuclear ambitions.

Given Mr.~Clinton's stature and his long interest in the North Korean nuclear issue, experts said it
was likely that his discussions in North Korea ranged well beyond obtaining the release of Ms.~Ling
and Ms.~Lee.

``It would be someplace between surprising and shocking if there wasn't some substantive discussion
between the former president, who is deeply knowledgeable about the nuclear issue, and Kim
Jong-il,'' said Robert L.~Gallucci, who negotiated with North Korea in the Clinton administration.

Mr.~Clinton has sought to find the right place in the Obama era, eager to play a role without
stepping on the toes of the new president or certainly of his secretary of state.

The last time the two spoke, said the White House press secretary, Robert Gibbs, was in March, when
Mr.~Obama invited Mr.~Clinton to a ceremony in Washington for signing legislation expanding the
AmeriCorps program created by Mr.~Clinton.

In interviews last spring, Mr.~Clinton said that he would be happy to do anything Mr.~Obama asked
him to do, but that ``I try to stay out of their way.''

Mr.~Clinton's mission may be less of an issue for Mr.~Obama than for Mrs.~Clinton. The same day he
landed in North Korea, she arrived in Kenya, kicking off an 11-day journey through Africa -- a visit
now largely eclipsed by her husband's travels.

\section{China Tries Campaigner On '08 Quake}

\lettrine{H}{uman} rights advocates are calling on the Chinese government to cancel the criminal
trials of two men who pushed for official investigations into the causes of widespread school
collapses during the devastating 2008 earthquake in Sichuan Province.

The trial of one man, Huang Qi, began Wednesday but adjourned\cite{adjourn} without a verdict.
Mr.~Huang, a well-known blogger and civil rights campaigner, is accused of possessing state secrets,
which carries a sentence of five years to life. The second defendant, Tan Zuoren, a writer and also
a prominent rights advocate, faces a potential five-year sentence for subversion and is to go on
trial Aug.~12. The charges are broad ones the Chinese government often uses to silence people who
publicly challenge the government.

``These trials are not about a reasonable application of the law, but about silencing government
critics whose work has considerable public benefit and sympathy,'' Sophie Richardson, the Asia
advocacy director at Human Rights Watch, an advocacy group based in New York, said in a written
statement released Tuesday. ``The government is likely seeking to squelch\cite{squelch} those who
cause it embarrassment, but in the process it is undermining domestic and international confidence
in its ability to cope in a transparent way with natural disasters.''

Mr.~Huang's wife, Zeng Li, said in a telephone interview that her husband's trial began at 10:30
a.m. and continued for three hours. It was unclear Wednesday night when the trial, closed to the
public, would resume.

The May 2008 earthquake was the most devastating natural disaster in China in decades, killing
nearly 69,000 people and leaving about 18,000 missing, all presumed dead, according to official
estimates. Initial reports from official news media said about 7,000 schoolrooms collapsed, and as
many as 10,000 schoolchildren may have perished.

Many of the schools collapsed even though buildings next to them remained standing, which grieving
parents and advocates attributed to shoddy\cite{shoddy} construction and corruption. Officials in
Sichuan blamed the earthquake itself, not bad construction.

Mr.~Huang, who tried to help five bereaved fathers push for investigations and compensation, was
detained by the police in June 2008 and formally charged the next month. Mr.~Huang ran the Tianwang
Human Rights Center, an informal advocacy group, in Chengdu, the provincial capital of Sichuan, and
discussed rights issues on his Web site, www.64tianwang.com.

Mr.~Tan was detained in March and charged the next month. The charge of subversion is related to his
attempts to investigate construction of the schools and to organize an event to observe the 20th
anniversary of the killing of civilians by the government during the Tiananmen Square protests of
June 1989, according to Human Rights Watch.

\section{Streets Empty in Plague-Stricken Town}

\lettrine{P}{ublic} buses were ordered off the roads of a remote Chinese town to control the
possible spread of pneumonic plague that has killed three people and seemed poised to claim a fourth
victim, residents and the authorities said Wednesday.

One patient hospitalized with the highly infectious lung disease was ``near death'' while another
person was in serious condition, according to a notice from the health bureau of the northwestern
province of Qinghai, where the stricken town of Ziketan is located.

Seven others infected were ``basically stable,'' the notice said.

The police have set up checkpoints around Ziketan, a farming town of 10,000 people in an ethnically
Tibetan area, sealing it off to prevent the spread of the disease that can kill in as little as 24
hours if left untreated. Residents have reported that some people have tried to flee the town on
foot, though it is unclear whether they made it past the checkpoints.

The authorities have set up a cordon\cite{cordon} with a 28-kilometer, or 17-mile, radius around
Ziketan, more than 480 kilometers west of Beijing, residents said. Calls to the town and provincial
public security bureaus rang unanswered.

Public buses were pulled off the streets, and few shops were open, according to a food seller in
Ziketan, surnamed Han. Residents have been told to disinfect their homes and shops, he said.

``Yesterday afternoon there were police patrolling on the streets, advising the shops to close,''
said Mr.~Han, who gave only his surname because of the sensitivity of the issue. ``I took a stroll
out of my shop earlier this morning and found that only around 20 shops are open in town. There are
few people or vehicles on the streets.''

China has had previous cases of plague, a disease that circulates mainly among small animals.
Experts have said most cases in China's northwest occur when hunters are contaminated while skinning
infected animals. Pneumonic plague is the least common and most deadly form of the disease. It can
be directly spread between humans because the bacteria is airborne. But if treated early with
antibiotics, it is curable.

\section{A Small-Business Guide to Intellectual Property}

\lettrine{T}{he} two most precious resources for any small-business owner are time and money. That's
why when the subject of intellectual property comes up, many owners run in the other direction. They
see images of expensive lawyers and use that as an excuse to ignore the topic, reasoning that it is
a problem for big companies to worry about.

The trouble is, with the rise of competition through the Internet and on the global market,
understanding intellectual property is more critical than ever for small-business owners. Let's
explore some of the common fallacies:

1. For small-business owners, it's not worth the time or effort to secure intellectual property
rights.

Daniel Lubetzky, chief executive of New York City-based Kind Snacks, had high hopes when he and his
company attended the Natural Products Expo West in Anaheim, Calif., in March. And who could blame
him, since his Kind Plus bars had been named the best new product at the Natural Products Expo East
last October?

But it didn't take long before Mr.~Lubetzky knew something had gone wrong: He kept hearing how one
of his competitors had copied the packaging, look and feel of his bars.

Fortunately for Mr.~Lubetzky, he had secured crucial components of intellectual property like
trademarks, trade dress (the look and feel of a product) and Web addresses after founding his
company. Unlike a patent, which can cost up to \$25,000 to secure, trademarks and Web addresses can
be obtained relatively cheaply and without the aid of a lawyer.

With the legal documentation to back up his intellectual property rights, Mr.~Lubetzky sent the
offending company a cease-and-desist letter, which achieved the desired result. ``Too many
entrepreneurs forget there is more to I.P. than just patents,'' said Mr.~Lubetzky, who happens to be
a lawyer.

2. Once I get a trademark, my brand is safe.

It may be. But consider what happened to Tracey Deschaine, who runs a restaurant called Dixie Picnic
in Ocean City, N.J.

When Ms.~Deschaine opened her business in 2006, she secured trademarks on her business name and logo
and on the name of her signature item, ``upcakes,'' which are upside-down frosted cupcakes. The
problem, she says, was that even though she had obtained the trademarks, someone monitoring the
activity on the United States Patent and Trademark Office's Web site had spotted her application and
secured upcakes.com as the Web address, or U.R.L., before she could.

``I had no idea that even though I have a trademark, someone else could just go register the
U.R.L.,'' she said. ``I wish I had planned ahead and bought the site before I did that.''

3. Having a patent gives me the right to produce something.

This is a very fundamental misunderstanding. Actually, what a patent does is give you the right to
prevent someone else from producing what your patent covers. ``Having a strong I.P. position helps
ensure that other people pay you for your innovation like they would a toll on a road,'' Mr.~Kocher
said.

But even if you do have a patent, there's no guarantee that someone won't try to get around it.
There's also no guarantee that you will win if you fight that person. But if you have your I.P.
ducks in a row and a commitment to do whatever you can to defend those rights, you do have a
fighting chance -- even in a fight against a much larger company.

Consider the example of Cryptography Research, a 20-employee technology firm in San Francisco that
specializes in data security. Beginning in 2004, the company made the decision to pursue litigation
against the credit card giant Visa, which Cryptography asserted was infringing on its patents
covering smart cards. To pursue the case against Visa, however, Cryptography's founder, Paul Kocher,
knew he needed a serious war chest in addition to his patent portfolio.

That's why he decided to sell off another piece of his business, patents covering technology that
protects Blu-ray discs from piracy, to Macrovision, which is now known as Rovi, in 2007 for \$45
million. ``All of a sudden we became a formidable opponent for someone who thought we couldn't
fight,'' Mr.~Kocher said. In the end, the gamble paid off, as the two companies settled out of
court, with Visa's agreeing to license the technology from Cryptography.

4. If I have a patent or trademark in the United States, I don't need to worry about the rest of the
world.

It depends on your business model. Intellectual property rights, which also include country-specific
U.R.L.'s, need to be obtained country by country, some of which protect them better than others. The
cost can vary, too.

In Japan, for example, it is notoriously expensive to acquire patents. In addition, the annual fees
required to maintain the patents there are often prohibitively expensive for small businesses, said
Gary Johnson, chief executive of Blue Spark Technologies, a manufacturer based in West Lake, Ohio,
that makes small, flexible batteries used in things like radio frequency identification tags.

``What we have done is to develop a strategy to go after I.P. protection in a limited number of
countries that we think we are most likely to sell or manufacture in, like the U.S.~and China,'' he
said. ``A lot of the choice comes down to what your business plan tells you.'' To decide what your
international I.P. strategy should be, consult a lawyer and conduct some cost-benefit analysis to
see if expanding your I.P. rights makes sense.

5. People who collect patents but don't actually make anything are ``patent trolls,'' parasites who
can make money only by filing lawsuits against real businesses.

The term ``patent troll'' was coined in the wake of the epic lawsuit fought between NTP, a small
holding company, and Research in Motion, which makes the hugely popular BlackBerry. The focal point
of the dispute was a patent for wireless e-mail delivery held by NTP -- something that R.I.M.
eventually would pay millions of dollars to license. But what most people remember about the story
is the lawsuits and the notion that NTP was somehow in the wrong for trying to enforce its patent,
mostly because it didn't make any products itself.

But consider that many inventors never set out to build a company, only to partner with someone who
would bring their products to life. Thomas Edison, for instance, received more than 1,000 patents --
many of which he licensed to other companies. ``He created what we might consider the first
innovation factory,'' says Mark Blaxill a co-founder of 3LP Advisors, an intellectual property
consulting company based in Boston.

A more recent example is Trident Design, a company founded by an inventor, Chris Hawker, which
patented and then licensed the design for the PowerSquid. Like Edison, Mr.~Hawker's company invents
products, builds an intellectual-property wall around them and then licenses them to other
companies.

The PowerSquid is now manufactured by a division of Phillips Electronics and sold by a spinoff of
Trident called Flexity. ``Our entire business model is leveraging our I.P.,'' Mr.~Hawker said.

\section{Bill Clinton and Journalists in Emotional Return to U.S.}

\lettrine{F}{ormer} President Bill Clinton arrived in the United States Wednesday morning after a
dramatic 20-hour visit to North Korea, in which he won the freedom of two American journalists,
opened a diplomatic channel to North Korea's reclusive government and dined with the North's ailing
leader, Kim Jong-il.

The private plane carrying Mr.~Clinton and the journalists, Laura Ling, 32, and Euna Lee, 36, landed
at 5:50 a.m. Pacific Daylight Time at Bob Hope Airport in Burbank, just outside Los Angeles.

The two women stepped off the plane in jeans and sweaters, rushing down the stairs to be reunited
with their families, who clustered around them. Ms.~Lee, in tears, embraced her husband, Michael
Saldate, and knelt to hug her 4-year-old daughter, Hana. Ms.~Ling kissed her husband, Iain Clayton.
Mr.~Clinton stepped off the plane a few moments later, embracing former Vice President Al Gore, the
founder of the media company that employs the journalists.

``Thirty hours ago, Euna Lee and I were prisoners in North Korea,'' Ms.~Ling said in brief remarks
to reporters, blinking back tears. ``We feared that at any moment we could be prisoners in a hard
labor camp. Then suddenly we were told that we were going to a meeting.

``We were taken to a location and when we walked through the doors, we saw standing before us
President Bill Clinton,'' she said, recounting the final moments of her ordeal\cite{ordeal}. ``We
were shocked, but we knew instantly in our hearts that the nightmare of our lives was finally coming
to an end. And now we stand here home and free.''

Mr.~Gore then spoke. ``President Obama and countless members of his administration have been deeply
involved,'' in the effort to bring the women home, he said. ``To everybody who has played a part in
this,'' he said, ``we are so grateful.''

The North Korean government, which in June sentenced the women to 12 years of hard labor for
illegally entering North Korean territory, announced hours before the jet's departure from North
Korea that it had pardoned the women after Mr.~Clinton apologized to Mr.~Kim for their actions,
according to the North Korean state media.

Mr.~Clinton's wife, Secretary of State Hillary Rodham Clinton, said Wednesday that the
administration was ``extremely excited'' that the women would be reunited with their families. But
she denied that her husband had apologized.

President Obama, who contacted the families of the women on Tuesday evening, said that he, too, was
``extraordinarily relieved'' at the journalists' return.

``I want to thank President Bill Clinton -- I had a chance to talk to him -- for the extraordinary
humanitarian effort that resulted in the release of the two journalists,'' Mr.~Obama said outside
the White House on Wednesday morning.

Mr.~Clinton's mission to Pyongyang was the most visible by an American in nearly a decade. It came
at a time when the United States' relationship with North Korea had become especially chilled, after
North Korea's test of its second nuclear device in May and a series of missile launchings.

It ended a harrowing ordeal for the two women, who were stopped on March 17 by soldiers near North
Korea's border with China while researching a report about women and human trafficking. They faced
years of imprisonment in the gulag-like confines of a North Korean prison camp.

And it catapulted Mr.~Clinton back on to the global stage, on behalf of a president who defeated
Mrs.~Clinton in a bitter primary campaign last year, and who later asked her to be his secretary of
state.

Mrs.~Clinton was deeply involved in the case, too. She proposed sending various people to Pyongyang
-- including Mr.~Gore -- to lobby for the release of the women, before Mr.~Clinton emerged as the
preferred choice of the North Koreans, people briefed on the talks said.

About 10 days ago, these people said, Mr.~Gore called Mr.~Clinton to ask him to undertake the trip.
Mr.~Clinton agreed, as long as the Obama administration did not object.

In an interview Wednesday with NBC's ``Today'' show in Nairobi, Kenya, Mrs.~Clinton said the final
request to Mr.~Clinton had come from the White House.

``When the message came to us from the young women themselves, to their families, to former Vice
President Gore, and then to the administration that sending my husband would be the best way to
ensure their release, of course we took that very seriously, discussed it,'' she said, according to
a transcript. ``The White House reached out, as they said, to my husband, to ask him if he would be
willing to do that.''

The riveting\cite{rivet} tableau\cite{tableau}, of a former president jetting into a diplomatic
crisis while his wife was embarking on a tour of Africa in her role as the nation's chief diplomat,
underscored the unique and enduring role of the Clintons, even in the Obama era.

At a news conference in Nairobi on Wednesday, Mrs.~Clinton said the case of the captured Americans
had been handled separately from America's dispute with North Korea over its nuclear program.

``We have been working hard on the release of the two journalists, and we have always considered
that a separate issue,'' she said. The future of the United States' relationship with Pyongyang, she
continued, was ``really up to them.''

Mr.~Clinton's trip to Pyongyang came just two weeks after North Korea issued a harsh personal attack
on Mrs.~Clinton, in response to comments she made comparing its nuclear test and missile launchings
to the behavior of an attention-seeking teenager.

The North Korean Foreign Ministry objected to her ``vulgar remarks'' and called her ``a funny lady''
who was neither intelligent nor diplomatic. ``Sometimes she looks like a primary-school girl and
sometimes a pensioner going shopping,'' a spokesman said.

The episode evidently did not stop consideration of sending her husband as an envoy\cite{envoy}. But
the initiative was cloaked in secrecy and came after weeks of back-channel talks between the United
States and North Korea through its United Nations mission. In addition to Mr.~Gore, the White
House's list of potential candidates included Gov. Bill Richardson of New Mexico.

North Korea signaled its desire to have Mr.~Clinton act as a special envoy in conversations with
Ms.~Ling and Ms.~Lee, who relayed that message to their families in the middle of July, according to
a senior administration official. The message was passed to Mr.~Gore, who contacted the White House,
which then explored whether such a mission would be successful.

Mr.~Obama did not speak directly with Mr.~Clinton before the mission. But his national security
adviser, Gen. James L.~Jones, contacted the former president to sound him out. The senior official
said the administration did ``due diligence'' with the North Koreans to ensure that if Mr.~Clinton
went, he would return with the journalists.

As president, Mr.~Clinton had sent Mr.~Kim a letter of condolence on the death of his father, Kim
Il-sung, according to a former official. For Mr.~Kim, the former official said, freeing the women
was a ``reciprocal humanitarian gesture.''

Mr.~Kim is believed to have suffered a stroke last year. American officials said they thought his
declining health had set off a succession struggle, complicating the Obama administration's dealings
with the North.

The families of the American journalists issued a statement saying they were ``overjoyed'' by news
of the pardon and thanked Mr.~Obama and Mrs.~Clinton. ``We especially want to thank President Bill
Clinton for taking on such an arduous mission and Vice President Al Gore for his tireless efforts to
bring Laura and Euna home,'' the statement said.

Current TV said in a statement that it was also ``overjoyed'' and that the hearts of its employees
went out to Ms.~Ling and Ms.~Lee for ``persevering through this horrible experience.''

Mr.~Clinton and the journalists traveled on a private jet owned by Stephen Bing, a real estate heir
and a major Democratic Party contributor, who donated the use of the jet as a favor, said Andy
Meyers, the chief executive officer of Shangri-la Industries, which was founded by Mr.~Bing.

Administration officials said Mr.~Clinton went to North Korea as a private citizen, did not carry a
message from Mr.~Obama for Mr.~Kim and had the authority to negotiate only for the women's release.

``This was 100 percent about the journalists,'' said a senior administration official, who spoke on
the condition of anonymity because of the delicacy of the matter. ``We knew Kim Jong-il would
probably seek a meeting with Clinton. But that's not what this visit was about.''

Still, North Korea, clearly seeing a propaganda\cite{propaganda} opportunity at home and a rare
chance for a measure of favorable publicity abroad, welcomed Mr.~Clinton with the
fanfare\cite{fanfare} of a state visit. It broadcast a group portrait, as well as photos of Mr.~Kim
gesturing and talking to Mr.~Clinton; of the former president accepting flowers from a North Korean
girl; and of Mr.~Clinton, seated across a negotiating table from Mr.~Kim, each flanked by aides.
Among those greeting Mr.~Clinton at the airport was Kim Kye-gwan, North Korea's chief nuclear
negotiator.

Among those accompanying Mr.~Clinton was David Straub, a former director of the Korea desk at the
State Department, who had held talks with the North Koreans through what is known as the ``New York
connection.''

Also on hand was John Podesta, an informal adviser to the Obama administration who served as
Mr.~Clinton's chief of staff in the final years of his presidency, when the former president yearned
to travel to North Korea to clinch a deal that would have curbed its nuclear program.

That visit never happened -- partly because the White House concluded that a deal was not assured --
and former President George W.~Bush put the brakes on direct talks with North Korea, setting the
stage for eight years of largely fruitless efforts to stop the North's nuclear ambitions.

Given Mr.~Clinton's stature and his long interest in the North Korean nuclear issue, experts said it
was likely that his discussions in North Korea ranged well beyond obtaining the release of Ms.~Ling
and Ms.~Lee.

``It would be someplace between surprising and shocking if there wasn't some substantive discussion
between the former president, who is deeply knowledgeable about the nuclear issue, and Kim
Jong-il,'' said Robert L.~Gallucci, who negotiated with North Korea in the Clinton administration.

Mr.~Clinton has sought to find the right place in the Obama era, eager to play a role without
stepping on the toes of the new president or, certainly, the secretary of state.

The last time the two had spoken, said the White House press secretary, Robert Gibbs, was in March,
when Mr.~Obama invited Mr.~Clinton to a ceremony in Washington for signing legislation expanding the
AmeriCorps program created by Mr.~Clinton.

In interviews last spring, Mr.~Clinton said that he would be happy to do anything Mr.~Obama asked
him to do, but that ``I try to stay out of their way.''

Mr.~Clinton's mission may be less of an issue for Mr.~Obama than for Mrs.~Clinton. The same day he
landed in North Korea, she arrived in Kenya, beginning an 11-day journey through Africa -- a visit
now largely eclipsed by her husband's travels.

\section{Clinton and Gore, Together Again}

\lettrine{I}{t} was a gripping moment: Bill Clinton, the former president, and Al Gore, his vice
president, sharing a lengthy embrace as Mr.~Clinton delivered two journalists who worked for Mr.
Gore back to American soil from captivity in North Korea.

The tableau at the airport in Burbank, Calif., on Wednesday morning was a visible reminder of how
circumstances had conspired over the past few days to bring the reigning -- and sometimes warring --
names of the Democratic Party together around the drama of Mr.~Clinton's trip to North Korea. There
on the tarmac\cite{tarmac} were the two dominant Democrats of the 1990s, Mr.~Gore and Mr.~Clinton,
who had played a critical role negotiating the end of an international crisis on behalf of President
Obama and Hillary Rodham Clinton, the secretary of state and Mr.~Obama's rival for the nomination in
2008.

Mr.~Clinton and Mr.~Gore did not linger\cite{linger} together in public on Wednesday. They have had
relatively little contact after their unhappy parting nearly nine years ago, according to associates
of the two men. Still, over eight years, they have both built post-White House identities and
reputations that seem to have left them at peace with their shared pasts and, it would seem, with
each other.

``A special thanks to President Clinton,'' Mr.~Gore said, as Mr.~Clinton gazed at him, his head
nodding. ``My partner and friend. So grateful.''

Mike Feldman, a Democratic consultant who has worked for both men, said: ``Look, they've both gone
off and done different things. But when you work side by side as they did, that's a pretty strong
bond and that doesn't go away. People say they don't spend a lot of time together -- well, they live
in different cities.''

For just a moment -- even with extra heft and grayer hair -- it was easy to recall the happier and
era-defining time when Mr.~Clinton and Mr.~Gore captured the attention of the nation, campaigning in
a bus caravan in 1992 as the Democratic candidates for president and vice president.

On the campaign trail, and later in the White House, the two men and their spouses presented
themselves as friends: Sharing a bus and jokes and warm embraces, Mr.~Gore introducing Mr.~Clinton
at every stop on the way, feeding off each other's riffs and energies. Still, associates of both men
said, that relationship was always more pragmatic\cite{pragmatic} than personal, a recognition by
two very ambitious men, whose similar backgrounds almost ensured a rivalry, that they had stumbled
across a chemistry that could deliver them the White House.

The tensions that seem to invariably sour relationships between presidents and vice presidents by
the end emerged in the Clinton White House, building to the moment when Mr.~Gore made clear to
associates his revulsion after learning that Mr.~Clinton had had an affair with an intern.

Mr.~Gore confronted Mr.~Clinton in the Oval Office after Mr.~Gore ran for president in 2000 and lost
the electoral college vote to George W.~Bush on a ruling by the Supreme Court. In the course of a
brutal one-hour exchange, aides to both men said, Mr.~Gore blamed Mr.~Clinton for his loss --
suggesting he had been dragged down by the burden of Mr.~Clinton's impeachment proceedings -- and
Mr.~Clinton angrily attacked Mr.~Gore for distancing himself from Mr.~Clinton and his record.

Mr.~Gore avoided appearing with Mr.~Clinton in the first years after his defeat, and was viewed in
Clinton circles as implicitly supporting Mr.~Obama in his competition last year with Mrs.~Clinton --
though one associate of Mr.~Clinton said the former president was appreciative in the end that
Mr.~Gore had sat it out.

But much of that had passed by the time Mr.~Gore called Mr.~Clinton 10 days ago and asked for his
help.

And while they have crossed paths from time to time, there has clearly been nothing like the moment
on Wednesday when Mr.~Gore -- spotting Mr.~Clinton at the top of the steps coming off the aircraft
-- walked forward, clapping his hands loudly, leading the audience in applause. They shook hands
first, and than hugged for a full five seconds, Mr.~Clinton patting his hand along Mr.~Gore's back,
lingering, before Mr.~Clinton went on to greet the waiting family members.

\section{After Clinton Trip, U.S. Studies Signals From N. Korea}

\lettrine{A}{ day} after former President Bill Clinton's flight into North Korea to win the freedom
of two American journalists, the Obama administration moved Wednesday to send a stern\cite{stern}
message to the North Korean government: nothing has changed.

While President Obama celebrated the emotional reunion of the journalists, Laura Ling and Euna Lee,
with their families, he said Mr.~Clinton's ``humanitarian mission'' did not ease the demands of the
United States and many allies that North Korea alter its behavior if it wants to escape its
isolation.

But even as the administration took pains through its statements to assure allies abroad and voters
at home that it would not be lured na\"ively into yet another round of fruitless talks with the
North Korean government, officials were busily calculating whether the psychology in the North had
changed in a way that might create opportunities for resolving the long confrontation over its
nuclear ambitions.

They were also evaluating preliminary impressions brought back by Mr.~Clinton and the other
Americans on the trip about North Korea's reclusive leader, Kim Jong-il, and any diplomatic signals
Mr.~Kim might have been trying to convey.

Administration officials said the White House had no plans to change its strategy for negotiating
with North Korea, which involves imposing strict sanctions on the government and inviting it back to
multiparty negotiations only if it agrees to give up its nuclear weapons program.

``We have said to the North Koreans that there's a path to better relations,'' Mr.~Obama said in an
interview with MSNBC. ``We just want to make sure the government of North Korea is operating within
the basic rules of the international community that they know is expected of them.''

Mr.~Obama's message, repeated by Secretary of State Hillary Rodham Clinton, illustrated the
determination of the White House not to allow Mr.~Kim to turn Mr.~Clinton's visit to his advantage
in talks with the United States and other countries.

Still, the events of the past two weeks show that relations between the United States and the North
are, if nothing else, rich in diplomatic theater, with a new American administration working to find
the right tone and strategy for dealing with a country that has confounded multiple presidents,
including Mr.~Clinton.

Officials said the administration's restatement of its policy did not mean that it was blind to the
opportunities that could flow out of the mission's happy outcome.

``How this impacts the psychology of the North Koreans, no one can tell,'' said a senior
administration official, who spoke on the condition of anonymity because of the delicacy of the
issue.

The White House began to debrief Mr.~Clinton later on Wednesday, after he had short phone
conversations with Mr.~Obama and Mrs.~Clinton, who is traveling in Africa. After delivering Ms.~Ling
and Ms.~Lee to their families at an airport in Burbank, Calif., Mr.~Clinton flew to his home in
Chappaqua, N.Y.

Among other issues, administration officials said, there is no established protocol for a
face-to-face debriefing of a busy former president about intelligence and diplomatic matters. Do you
summon\cite{summon} him to the White House, or do you send a delegation to visit him?

Mr.~Clinton has not spoken publicly about the trip, which was cloaked in secrecy until just before
he arrived in North Korea.

He flew across the Pacific and back on a Boeing 737 that belongs to Steve Bing, a film producer and
friend of the former president. Mr.~Bing has been a major donor to Democratic candidates.

For all the talk of the North's isolation, the United States and North Korea hold regular
back-channel talks through North Korea's United Nations mission. The so-called New York connection
laid the groundwork for Mr.~Clinton's visit and will probably play a part in what comes next.

The release of the journalists could break an impasse between the United States and North Korea
because, officials and outside experts said, there was virtually no prospect of reviving nuclear
negotiations as long as the women appeared to be headed to a North Korean prison camp.

At the same time, Mr.~Kim is likely to expect some kind of payback for the pardons, political
analysts said.

How the United States responds to Mr.~Kim's expectations may determine whether Mr.~Clinton's visit
leads to a genuine opening between the countries.

``They're going to expect us to take our foot off the pedal,'' said Victor Cha, who negotiated with
North Korea in the Bush administration. ``If China calls and says North Korea is willing to come
back to six-party talks in Beijing if we relax the sanctions, this becomes a tough call.''

In a sign of the administration's sensitivity, a senior official said the United States had briefed
officials at the highest levels of the Chinese and Russian governments last weekend about the nature
of Mr.~Clinton's mission, to reassure them that the White House did not intend the visit to lead to
bilateral talks with North Korea.

One American official noted that Mr.~Clinton, in his talks with Mr.~Kim, broached the issue of
Japanese and South Koreans abducted by North Korea -- politically emotive matters for those
countries, in their dealings with the North.

Some of the administration's calculations will be influenced by Mr.~Clinton's observations about
Mr.~Kim's health and mood. He suffered a stroke last year, and his diminished health led to
speculation that the North was caught up in a struggle to succeed him.

American officials are also eagerly interviewing David Straub, a former director of the State
Department's Korea desk, who accompanied Mr.~Clinton. Mr.~Straub speaks Korean and took part in
back-channel communications with North Koreans when he was in the government.

As the administration deliberates over North Korea, there are reports that the North is suspected of
spreading its technology. The Sydney Morning Herald in Australia reported last week that Myanmar was
building a nuclear reactor and plutonium extraction plant, with help from North Korea.

The acting State Department spokesman, Robert A.~Wood, said Tuesday, ``We're obviously looking into
these types of reports.'' Mrs.~Clinton said recently that there was concern about possible transfers
of nuclear technology from North Korea to Myanmar.

\section{Intimate\cite{intimate} Homecoming Is Played Out in Public}

\lettrine{F}{or} months, Laura Ling's uncle sent her scores of unanswered e-mail messages,
suggesting that she practice meditation while held in confinement in North Korea. For weeks at a
time, Ms.~Ling's family in California would hear nothing, plunged into dark worries about her health
and future yet unable to do anything about her well-being.

And so, after Ms.~Ling's return to the United States on Wednesday as part of an extraordinary
diplomatic mission, her mother brought her soup.

``Laura loves this soup,'' said Mary Ling, standing in front of her daughter's house in North
Hollywood with a bouquet of flowers and some nourishment. ``It's a special Chinese herbal soup.''

Exhausted, emotionally drained and clearly moved to be reunited with their families after a long
plane trip home with former President Bill Clinton, the American journalists Laura Ling, 32, and
Euna Lee, 36, began the rest of their lives in Los Angeles on Wednesday, their loved ones elated
and, for now, on the other side of the cameras.

It was an intimate homecoming remarkably played out in a public setting that intertwined politics,
the Clinton family, journalism and the secrecy of a spy novel.

Both women wept as they descended the steps of the private airplane they rode with Mr.~Clinton from
North Korea, Ms.~Lee bending over for a second as if overcome with joy, Ms.~Ling pumping two fists
in the air. Both were met at the bottom by their husbands, whom they grabbed and clung to tightly,
and Ms.~Lee bent over to gaze at her 4-year-old daughter, Hana, whose little arms she grasped.

Ms.~Ling had called her sister Lisa from the airplane, and Lisa urged her to remember to mention her
surprise at seeing Mr.~Clinton in North Korea when she hit the tarmac in Burbank early Wednesday
morning. And then the sisters got to wrap their arms around each other. ``I can't ever describe what
my whole body went through,'' Lisa Ling said.

Laura Ling said she and Ms.~Lee had sensed the support for them while in prison, even without
constant communication.

``We could feel your love all the way in North Korea,'' Ms.~Ling said. ``It is what sustained us in
our darkest hours. We are very grateful we were granted amnesty by the government of North Korea.''

Iain Clayton, Ms.~Ling's husband, said that ``it was the best feeling of my life'' to see his wife
descend from the plane to his arms. ``I looked in her eyes,'' Mr.~Clayton said, ``and saw she was
seeing her family at the bottom of the steps, and I saw the joy and relief in her.''

The two women were arrested on March 17 near the North Korean border with China while reporting on
human trafficking for Current TV, a San Francisco-based media company co-founded by former Vice
President Al Gore. In June, the women were sentenced to 12 years of hard labor for illegally
entering the country, but they were freed Tuesday after Mr.~Clinton negotiated their release.

During the women's captivity, supporters held vigils\cite{vigil} in San Francisco, Washington and
other cities. Lisa Ling, a former co-host of ``The View'' and correspondent for ``The Oprah Winfrey
Show,'' played a pivotal role in drawing news media attention to their plight.

Lisa Ling said her sister and Ms.~Lee had been held in isolation for most of their detention,
without any knowledge of each other's well-being. North Korean authorities permitted the women to
call their families on a few occasions. At one point, Ms.~Ling told her family on the phone to write
to Ms.~Lee ``and tell her I'm thinking about her and I love her.''

And during a July call, Ms.~Ling and Ms.~Lee told their families that the North Koreans had told
them that they would be willing to grant them amnesty if ``an envoy in the person of President
Clinton would agree to come to Pyongyang and seek their release,'' according to a senior Obama
administration official who briefed reporters.

The proposal was then reported to Mr.~Gore, who passed it on to Mr.~Clinton.

After their daybreak return to Los Angeles, both women sequestered themselves inside their homes on
Wednesday, enjoying their reunions with their families. Attached to a wall outside Ms.~Lee's home
was a sign that read ``Welcome Home Euna'' placed by a neighbor. Los Angeles police officers stood
outside keeping the news media at bay.

Ms.~Ling's family was the more forthcoming. ``She's been in isolation for three and a half months,''
Lisa Ling said about her sister. ``She's trying to decompress. She hasn't been able to speak freely
in months. Before she wants to speak her story, she wants to make sure of her state of mind. She is
lethargic\cite{lethargic} and tired now.''

Dr.~Jerry Wang, a brother of Ms.~Ling's mother, said he sent e-mail to Ms.~Ling throughout her
imprisonment, though he knew she could not respond.

``Even though it's been hard, it's been a very good experience for our family,'' Dr.~Wang said. ``It
brought the family closer in a lot of ways just to be together in our thoughts and prayers for
Laura.''

Ms.~Lee made a call to the journalism unit at Current TV around 11 a.m. and employees gathered
around a speakerphone to hear her. She told the employees that she was grateful to be home, and
thanked them for the letters that reached her in prison, via the Swedish ambassador, which she said
she had read over and over.

Questions will inevitably arise about whether the women will return to their journalism careers. The
trip to China was the first overseas assignment for Ms.~Lee, who speaks Korean. Ms.~Ling sometimes
worked 12-hour days and shuttled frequently between the San Francisco headquarters of Current TV and
Los Angeles. Dr.~Wang said he was sure his niece would continue with journalism because it ``has
always been her passion.''

For now, there is food. ``She is really anxious to have fresh fruit and food,'' Lisa Ling said of
her sister. ``She told us there were rocks in her rice,'' while in detention, she said. ``There will
be a sushi dinner sometime really soon.''

\section{For Today's Graduate, Just One Word: Statistics}

\lettrine{A}{t} Harvard, Carrie Grimes majored in anthropology\cite{anthropology} and
archaeology\cite{archaeology} and ventured to places like Honduras, where she studied Mayan
settlement patterns by mapping where artifacts were found. But she was drawn to what she calls ``all
the computer and math stuff'' that was part of the job.

``People think of field archaeology as Indiana Jones, but much of what you really do is data
analysis,'' she said.

Now Ms.~Grimes does a different kind of digging. She works at Google, where she uses statistical
analysis of mounds of data to come up with ways to improve its search engine.

Ms.~Grimes is an Internet-age statistician, one of many who are changing the image of the profession
as a place for dronish number nerds. They are finding themselves increasingly in demand -- and even
cool.

``I keep saying that the sexy job in the next 10 years will be statisticians,'' said Hal Varian,
chief economist at Google. ``And I'm not kidding.''

The rising stature of statisticians, who can earn \$125,000 at top companies in their first year
after getting a doctorate, is a byproduct of the recent explosion of digital data. In field after
field, computing and the Web are creating new realms of data to explore -- sensor signals,
surveillance tapes, social network chatter, public records and more. And the digital data surge only
promises to accelerate, rising fivefold by 2012, according to a projection by IDC, a research firm.

Yet data is merely the raw material of knowledge. ``We're rapidly entering a world where everything
can be monitored and measured,'' said Erik Brynjolfsson, an economist and director of the
Massachusetts Institute of Technology's Center for Digital Business. ``But the big problem is going
to be the ability of humans to use, analyze and make sense of the data.''

The new breed of statisticians tackle that problem. They use powerful computers and sophisticated
mathematical models to hunt for meaningful patterns and insights in vast troves\cite{trove} of data.
The applications are as diverse as improving Internet search and online advertising,
culling\cite{cull} gene sequencing information for cancer research and analyzing sensor and location
data to optimize the handling of food shipments.

Even the recently ended Netflix contest, which offered \$1 million to anyone who could significantly
improve the company's movie recommendation system, was a battle waged with the weapons of modern
statistics.

Though at the fore, statisticians are only a small part of an army of experts using modern
statistical techniques for data analysis. Computing and numerical skills, experts say, matter far
more than degrees. So the new data sleuths\cite{sleuth} come from backgrounds like economics,
computer science and mathematics.

They are certainly welcomed in the White House these days. ``Robust, unbiased data are the first
step toward addressing our long-term economic needs and key policy priorities,'' Peter R.~Orszag,
director of the Office of Management and Budget, declared in a speech in May. Later that day,
Mr.~Orszag confessed in a blog entry that his talk on the importance of statistics was a subject
``near to my (admittedly wonkish) heart.''

I.B.M., seeing an opportunity in data-hunting services, created a Business Analytics and
Optimization Services group in April. The unit will tap the expertise of the more than 200
mathematicians, statisticians and other data analysts in its research labs -- but that number is not
enough. I.B.M. plans to retrain or hire 4,000 more analysts across the company.

In another sign of the growing interest in the field, an estimated 6,400 people are attending the
statistics profession's annual conference in Washington this week, up from around 5,400 in recent
years, according to the American Statistical Association. The attendees, men and women, young and
graying, looked much like any other crowd of tourists in the nation's capital. But their rapt
exchanges were filled with talk of randomization, parameters, regressions and data clusters. The
data surge is elevating a profession that traditionally tackled less visible and less lucrative
work, like figuring out life expectancy rates for insurance companies.

Ms.~Grimes, 32, got her doctorate in statistics from Stanford in 2003 and joined Google later that
year. She is now one of many statisticians in a group of 250 data analysts. She uses statistical
modeling to help improve the company's search technology.

For example, Ms.~Grimes worked on an algorithm to fine-tune Google's crawler software, which roams
the Web to constantly update its search index. The model increased the chances that the crawler
would scan frequently updated Web pages and make fewer trips to more static ones.

The goal, Ms.~Grimes explained, is to make tiny gains in the efficiency of computer and network use.
``Even an improvement of a percent or two can be huge, when you do things over the millions and
billions of times we do things at Google,'' she said.

It is the size of the data sets on the Web that opens new worlds of discovery. Traditionally, social
sciences tracked people's behavior by interviewing or surveying them. ``But the Web provides this
amazing resource for observing how millions of people interact,'' said Jon Kleinberg, a computer
scientist and social networking researcher at Cornell.

For example, in research just published, Mr.~Kleinberg and two colleagues followed the flow of ideas
across cyberspace. They tracked 1.6 million news sites and blogs during the 2008 presidential
campaign, using algorithms that scanned for phrases associated with news topics like ``lipstick on a
pig.''

The Cornell researchers found that, generally, the traditional media leads and the blogs follow,
typically by 2.5 hours. But a handful of blogs were quickest to quotes that later gained wide
attention.

The rich lode\cite{lode} of Web data, experts warn, has its perils\cite{peril}. Its sheer volume can
easily overwhelm statistical models. Statisticians also caution that strong correlations of data do
not necessarily prove a cause-and-effect link.

For example, in the late 1940s, before there was a polio\cite{polio} vaccine\cite{vaccine}, public
health experts in America noted that polio cases increased in step with the consumption of ice cream
and soft drinks, according to David Alan Grier, a historian and statistician at George Washington
University. Eliminating such treats was even recommended as part of an anti-polio diet. It turned
out that polio outbreaks were most common in the hot months of summer, when people naturally ate
more ice cream, showing only an association, Mr.~Grier said.

If the data explosion magnifies longstanding issues in statistics, it also opens up new frontiers.

``The key is to let computers do what they are good at, which is trawling these massive data sets
for something that is mathematically odd,'' said Daniel Gruhl, an I.B.M. researcher whose recent
work includes mining medical data to improve treatment. ``And that makes it easier for humans to do
what they are good at -- explain those anomalies.''

\section{U.S. to Reform Policy on Detention\cite{detention} for Immigrants}

\lettrine{T}{he} Obama administration intends to announce an ambitious plan on Thursday to overhaul
the much-criticized way the nation detains immigration violators, trying to transform it from a
patchwork of jail and prison cells to what its new chief called a ``truly civil detention system.''

Details are sketchy\cite{sketchy}, and even the first steps will take months or years to complete.
They include reviewing the federal government's contracts with more than 350 local jails and private
prisons, with an eye toward consolidating many detainees in places more suitable for noncriminals
facing deportation -- some possibly in centers built and run by the government.

The plan aims to establish more centralized authority over the system, which holds about 400,000
immigration detainees over the course of a year, and more direct oversight of detention centers that
have come under fire for mistreatment of detainees and substandard -- sometimes fatal -- medical
care.

One move starts immediately: the government will stop sending families to the T.~Don Hutto
Residential Center, a former state prison near Austin, Tex., that drew an American Civil Liberties
Union lawsuit and scathing news coverage for putting young children behind razor wire.

``We're trying to move away from 'one size fits all,' '' John Morton, who heads the Immigration and
Customs Enforcement agency as assistant secretary of homeland security, said in an interview on
Wednesday. Detention on a large scale must continue, he said, ``but it needs to be done thoughtfully
and humanely.''

Hutto, a 512-bed center run for profit by the Corrections Corporation of America under a \$2.8
million-a-month federal contract, was presented as a centerpiece of the Bush administration's tough
approach to immigration enforcement when it opened in 2006. The decision to stop sending families
there -- and to set aside plans for three new family detention centers -- is the Obama
administration's clearest departure from its predecessor's immigration enforcement policies.

So far, the new administration has embraced many of those policies, expanding a program to verify
worker immigration status that has been widely criticized, bolstering\cite{bolster}
 partnerships
between federal immigration agents and local police departments, and rejecting a petition for
legally binding rules on conditions in immigration detention.

But Mr.~Morton, a career prosecutor, said he was taking a new philosophical approach to detention --
that the system's purpose was to remove immigration violators from the country, not imprison them,
and that under the government's civil authority, detention is aimed at those who pose a serious risk
of flight or danger to the community.

Janet Napolitano, the secretary of homeland security, said last week that she expected the number of
detainees to stay the same or grow slightly. But Mr.~Morton added that the immigration agency would
consider alternative ways to assure that those who face deportation -- and are not dangerous -- do
not flee.

Reviewing and redesigning all facilities, programs and standards will be the task of a new Office of
Detention Policy and Planning, he said. Dora Schriro, special adviser to Ms.~Napolitano, will become
the director, assisted by two experts on detention management and medical care. The agency will also
form two advisory boards of community groups and immigrant advocates, one focusing on detention
policies and practices, the other on detainee health care.

Mr.~Morton said he would appoint 23 detention managers to work in the 23 largest detention centers,
including several run by private companies, to ensure that problems are promptly fixed. He is
reorganizing the agency's inspection unit into three regional operations, renaming it the Office of
Detention Oversight, and making its agents responsible for investigating detainee grievances as well
as conducting routine and random checks.

``A lot of this exists already,'' he said. ``A lot of it is making it work better'' while
Dr.~Schriro's office redesigns the detention system, which he called ``disjointed'' and ``very much
dependent on excess capacity in the criminal justice system.''

Asked if his vision could include building new civil detention centers, he said yes. The current
32,000-bed network costs \$2.4 billion a year, but the agency is not ready to calculate the cost of
a revamped system.

Vanita Gupta, an American Civil Liberties Union lawyer who led the lawsuit against the Hutto center,
was jubilant over the decision to stop sending families there, but cautious about the other
measures.

``The ending of family detention at Hutto is welcome news and long overdue,'' she said in an e-mail
message. ``However, without independently enforceable standards, a reduction in beds, or basic due
process before people are locked up, it is hard to see how the government's proposed overhaul of the
immigration detention system is anything other than a reorganization or renaming of what was in
place before.''

Ms.~Gupta said the changes at Hutto since 2006 illustrated the importance of enforceable rules.
Before the A.C.L.U. lawsuit was settled in 2007, some children under 10 stayed as long as a year,
mainly confined to family cells with open toilets, with only one hour of schooling a day. Children
told of being threatened by guards with separation from their parents, many of them asylum-seekers
from around the world.

Only through judicial enforcement of the settlement, she said, have children been granted such
liberties as wearing pajamas at night and taking crayons into family cells. The settlement also
required the agency to honor agency standards that had been ignored, like timely reviews of the
decision to detain a family at all. Some families have been deported, but others were released or
are now awaiting asylum decisions in housing run by nonprofit social service agencies.

That kind of stepped-up triage could be part of the more civil detention system envisioned by
Mr.~Morton and Dr.~Schriro, who has been reviewing the detention system for months and is expected
to report her recommendations soon.

But the Hutto case also points to the limits of their approach, advocates say. Under the settlement,
parents and children accused of immigration violations were detained when possible at the country's
only other family detention center, an 84-bed former nursing home in Leesport, Pa., called the Berks
Family Shelter Care Facility. The number detained at Hutto has dropped sharply, to 127 individuals
from as many as 450.

Advocates noted that Berks, though eclipsed by the criticism of Hutto -- the subject of protest
vigils, a New Yorker article and a documentary -- also has a history of problems, like guards who
disciplined children by sending them across the parking lot to a juvenile detention center, and
families' being held for two years.

The Hutto legal settlement expires Aug.~29. In the most recent monitoring report last month,
Magistrate Judge Andrew W.~Austin wrote: ``Although the use of this facility to hold families is not
a violation of the settlement agreement, it seems fundamentally wrong to house children and their
noncriminal parents this way. We can do better.''

Mr.~Morton, a career prosecutor, seemed to agree. Hutto will be converted into an immigration jail
for women, he said, adding: ``I'm not ruling out the possibility of detaining families. But Berks is
the better facility for that. Hutto is not the long-term answer.''

\section{On the Fairway, New Lessons in Saving Water}

\lettrine{S}{ix} years ago, when Georgia's state government rewrote its rules for water use during
droughts, it cut no slack for an obvious culprit\cite{culprit}: golf courses.

With emerald\cite{emerald} fairways that glistened\cite{glisten} even in the most
blistering\cite{blistering} conditions, they were a tempting target.

Yet golf course managers were indignant\cite{indignant}. They argued that they were reining in water
use in dozens of ways, like planting native grasses and auditing sprinkler spray patterns. Instead
of being penalized, they said, they should be emulated\cite{emulate}.

It took a while, but from the South to the arid\cite{arid} West, their wish is coming true. Mindful
that global warming could provoke more and longer dry spells, state governments are increasingly
consulting golf courses on water strategies.

In Georgia, golf course managers have emerged as go-to gurus on water conservation for both
industries and nonprofit groups.

Marriott International is applying lessons learned at its golf course here to its resort properties
in other states. Habitat for Humanity is landscaping front yards with drought-tolerant plants
recommended by golf superintendents.

``Look, if you want to learn how to irrigate\cite{irrigate}, these are the guys to ask,'' said
Garith Grinnell, who recently retired from the Southern Nevada Water Authority.

Such accolades\cite{accolade} are a turnabout for a business that is often faulted for harming the
environment through excess use of water and pesticides\cite{pesticide}.

In Georgia, the shift in perspective came about largely because of a crippling drought that peaked
in 2007. By that year, 97 percent of the clubs that belonged to the Georgia Golf Course
Superintendents Association had voluntarily adopted what are viewed as best-management practices for
water use, reducing consumption, they estimated, by 25 percent in just three years.

Lake Lanier, Atlanta's main source of water, had meanwhile dropped to record low levels, exposing
muddy bottom not seen in half a century. It dawned on state and local water managers that golf
courses might have some useful know-how.

Golfing grounds managers ``are great technical assistance to me,'' said Kathy Nguyen, president of
the Georgia Water Wise Council, a state association of water professionals that encourages
conservation. ``I can call them up and talk to them about different technologies.'' (Georgia's
drought eased significantly this year.)

Ms.~Nguyen has relied on golf superintendents in drafting guidelines for homeowners like letting
grass grow longer, fixing leaks in hoses as promptly as possible and keeping lawn mower blades razor
sharp. (Grass cut by duller blades is more frayed and requires more water to stay healthy.)

The golf industry still draws strong criticism from environmentalists. Turf is, after all, the
thirstiest of plants. The average American golf course drinks up some 50 million gallons of water a
year -- comparable to the yearly usage of 1,400 people. In the West, the figures are higher.

Yet that reality, coupled with rising water prices, is what led to strides like irrigating golf
courses with ``gray water,'' or nonindustrial wastewater that is recycled for other purposes.

Tom Bancroft, chief scientist with the National Audubon Society, says that for all the progress golf
has made, it remains a deeply problematic industry. Many courses ``use fertilizers that can run off
into fresh water, and many use pesticides in lawn and grass,'' Mr.~Bancroft said. (Audubon
International, a separate group, works with golf courses to encourage wildlife preservation.)

Mark Esoda, superintendent of the Atlanta Country Club in suburban Marietta, where initiation fees
are \$85,000, acknowledges that practices among the nation's courses range from indifferent to
conscientious. But Mr.~Esoda maintains that he and other superintendents have a lot to teach
municipalities about watering their ball fields and homeowners about tending to their yards.

Zipping around the course on pine-shaded paths, Mr.~Esoda stopped abruptly near the seventh hole. He
gestured toward a patch of newly laid turf of zoysia, a warm-weather shade grass native to Southeast
Asia and Australia. On shady parts of the course it is replacing fescue, a genus of cold-weather
shade grass that can live through the winter and thus requires five months' more watering and
mowing.

Mr.~Esoda said he had also installed affordable monitors that prevent automatic sprinklers from
activating during or right after a rain.

And when isolated dry spots appear on the greens, he said, he sends staff members out with watering
cans rather than turning on the sprinkler system.

Finally, Mr.~Esoda has made an aesthetic adjustment after years of savoring the green glow of a
perfect lawn. ``Crispy around the edges is O.K.,'' he confided.

Water is just one area where golf courses and environmentalists may find a rapprochement, said
Anthony L.~Williams, director of grounds at Marriott's Stone Mountain public courses just outside
Atlanta.

As metropolitan areas sprawl outward, golf courses may be the only large-scale green space for miles
around, offering crucial potential habitat for migrating birds and other wildlife.

Mr.~Williams, who has a degree in local horticulture, has been letting native grasses take over his
lawns. Off the fairways he does not even bother to mow, and on the greens he is maintaining grass at
one-sixteenth of an inch higher than typical courses. It makes playing slow, he allows, but
``consistent.'' He has also replaced all the flowering annuals with perennials, which generally
require less water, choosing those that are attractive to native wildlife.

Since he took charge of the two courses in 2005, Mr.~Williams has cut water consumption by 45
percent, he said, and witnessed the return of some wildlife species like the red-tailed hawk.

The changes have come with a price, like the occasional large brown spot on the fairway. But
Mr.~Williams says the golfers do not mind.

``I just stand out there on the greens and explain, 'We are doing this so your grandchildren can
come out here and play,' '' he said. ``People understand that.''

\section{China Faces Delicate Task of Reining\cite{rein} in Bank Lending}

\lettrine{W}{hen} China announced three weeks ago that its economy had grown by 7.1 percent in the
first half of this year, this country appeared to be a lone bright spot during the global recession.

But many economists now worry that too much of China's growth was fueled by aggressive,
state-directed lending that could eventually result in a soaring number of bad loans and mounting
government debt.

While banks in the United States and Europe are still reluctant to make loans because of fears they
will not get their money back, Chinese banks issued a record 7.4 trillion yuan, or \$1.1 trillion,
in loans during the first six months of this year, mostly to big state-owned companies and
government infrastructure projects.

Here are just a few examples of the largess: In March, the city of Guangzhou was given a loan of
81.3 billion yuan to improve road transportation; in May, China's Aviation Industry Corporation said
it would receive 100 billion yuan to help the company export high-technology equipment; and a few
weeks ago, the China National Nuclear Corporation received loan approval for nearly 100 billion yuan
to advance the development of nuclear power.

``They opted for a very quick fix,'' said Stephen Roach, an economist and chairman of Morgan Stanley
Asia. ``Surging investment, fueled by the most rapid bank lending in history, accounted for nearly
90 percent of China's G.D.P. growth in the first half of this year. And that is worrisome.''

Mr.~Roach said China's growth remained too heavily weighted toward investment, rather than
consumption, creating unhealthy, imbalanced growth.

Analysts say that in China, new loans have grown this year at nearly three times the pace of a year
ago and that some of those loans may have been funneled into the resurgent Chinese stock and
property markets, creating the risk of new asset bubbles. The Shanghai stock market is up about 84
percent this year.

A similar lending binge in the 1990s led to an explosion of bad debt that left the biggest
state-owned Chinese banks nearly insolvent after the Asian financial crisis, until they were bailed
out by the government in a series of moves that ended in 2004.

This unprecedented move to stimulate the economy through bank lending now appears far more ambitious
than the \$586 billion economic stimulus package that Beijing announced last December.

In the United States, political leaders debate the extent to which the government ought to intervene
in the financial market. In China, there is little debate. The government simply orders banks to
lend -- and it happens almost instantly.

Some economists, including Mr.~Roach, fear that these moves could prove unsustainable and hinder the
long-term growth prospects for China.

But not everyone is so bearish. Some economists say that China's growth this year -- coming in at
7.9 percent in the second quarter -- is a positive sign for the global economy and a stunning
turnaround after the slowdown last year. They also say that Beijing is adept at changing course and
finding new policy tools for dealing with threats to growth. While China's growth rates have never
been exact and debate about the precise pace of expansion has persisted, no other big economy has
grown faster during the past three decades.

But if China's growth falters\cite{falter} later this year or early next year, some experts believe,
the global recession could deepen and lengthen. After all, many economists had forecast that much of
global growth over the next decade would come from China.

There is one thing economists agree on: The pace of lending in China must slow significantly in the
second half of this year, if Beijing is to avert a crisis.

Nicholas Lardy, an economist and China specialist at the Peterson Institute for International
Economics in Washington, contends\cite{contend} that China's investments in infrastructure have been
more prudent\cite{prudent} than in the 1990s and that lending will slow significantly in the second
half of this year. Still, he cautions: ``If that does not happen, I will become much more
pessimistic\cite{pessimistic}.''

Chinese regulators have already begun to acknowledge the risks of overly aggressive lending.

Last month, the Chinese central bank called for stricter supervision of bank loans because some
stimulus spending appeared to be directed toward wasteful government projects. The China Banking
Regulatory Commission also recently took steps to lower risk by ordering Chinese banks to raise
their bad-loan reserve ratios by the end of the year.

``Rapid expansion of bank loans in the first half year boosted the country's economic growth,'' Liu
Mingkang, the country's top banking regulator said in a speech in late July. ``But it also increased
the possibilities of financial risks.''

Still, Beijing seems unlikely to tighten monetary policy sharply. The government needs to keep the
economy on pace for 8 percent growth this year, quickly enough to create jobs and prevent social
instability, particularly before celebrations this October on the 60th anniversary of the founding
of the People's Republic of China.

At a high-level Politburo meeting last month, President Hu Jintao called for the country to maintain
a ``relatively loose'' monetary policy and signaled that the direction of macroeconomic policy
``should be maintained.''

The question is whether economic planners can strike a balance in encouraging growth. Restricting
lending could stall the recovery and make it difficult for Beijing to meet its growth target.
Allowing the lending boom to continue could sow the seeds of financial disaster.

``They've got to ensure there's enough loans for the real economy,'' said Wang Tao, an economist at
UBS Securities in Beijing. ``But they also have to slow down the loan growth to prevent an asset
bubble or future increase in nonperforming loans.''

In the first half of this year, China's bank loans were up more than 300 percent from a year ago, to
more than 7 trillion yuan, about equivalent to 25 percent of China's gross domestic product in 2008.

Chinese banks were able to do this because they have strong balance sheets relative to Western
banks. Bad loans were cleared off the books years ago and initial public stock offerings in recent
years pumped billions into many of China's big banks. The banks had also avoided many of the toxic
assets American and European banks held.

Through June, China's banking regulators said nonperforming loans at China's commercial banks
accounted for just 1.74 percent of outstanding loans -- down slightly from the beginning of the
year.

But many experts say that could change in the coming years, as these new loans come due. They
suspect that government pressure on banks to make loans this year will eventually result in a
spike\cite{spike} in loan defaults.

``If you tell bankers, lend as much as you like -- lend, lend, lend, and don't worry about the risk,
never in history has there not been a large increase in misallocated capital,'' said Michael Pettis,
a professor of finance at Peking University. ``It's never happened.''

Indeed, Fitch Ratings, the credit rating agency, warned in May that Chinese banks were facing new
risks, particularly because this year loan growth skyrocketed at a time when corporate profits were
weakening. ``Ordinarily, falling corporate earnings are met with tightened lending,'' the rating
agency said in its report. ``But in China precisely the reverse is evident, illustrating that
despite years of reform Chinese banks still retain an important policy function in upholding local
enterprises.''

Andy Xie, an economist based in Shanghai, said that one of the side effects of this huge lending
spree\cite{spree} was speculation\cite{speculation} in commodities, stocks and property -- all of
which could saddle the banks and local governments with huge debts, once prices fell.

``This could be a national debt issue,'' Mr.~Xie, a former Morgan Stanley economist, said. ``All of
this is government money.''

Recently, big state-owned property developers, flush with loans from state-owned banks, have made
record bids at state land auctions. And the rise in the Shanghai stock market is a feverish
turnaround following the sell-off last year.

No one has precise statistics or even clear proof that loans are flowing into the stock market, but
most analysts seem to believe that to be the case. Besides, they say, there was tremendous evidence
that in 2006 and 2007 many state-owned companies, even power producers, had set up independent units
to invest in stocks.

Among economists, though, an even bigger worry is that China is not doing more to rebalance its
economy, away from investment and exports and toward consumption. The country, they say, needs a new
growth model.

With exports down more than 20 percent from a year ago, China appears to be stimulating growth with
loans to state-owned companies and government works. But how long can that last? Mr.~Roach, the
Morgan Stanley economist, said not for long. China needs to stimulate its consumers.

\section{China Investigates Top Nuclear Official}

\lettrine{T}{he} top official of China's civilian and military nuclear power programs is being
investigated for ``grave violations of discipline,'' a phrase often used in corruption inquiries,
the Chinese Communist Party's disciplinary committee has announced.

The official, Kang Rixin, is the general manager and Communist Party secretary of China National
Nuclear Corporation, a vast holding company that is spearheading plans to increase the nation's
capacity to generate nuclear power at least sixfold in the next decade.

Mr.~Kang, 56, also is a member of the Communist Party Central Committee, the party's senior ruling
body, and sits on the same party disciplinary committee that is investigating him.

The announcement, issued late Wednesday, gave no indication of the accusations against Mr.~Kang. The
Beijing-based business journal Caijing and China Business News Daily quoted unidentified sources as
saying the inquiry centered on about \$263 million in company funds that apparently were invested
and lost in the stock market.

The publications also stated that investigators were looking into suspected bidding irregularities
in nuclear power plant contracts.

The inquiry is the latest in a series of investigations, involving high-ranking Chinese officials,
that some experts believe are supported by President Hu Jintao, who has often called for a
nationwide crackdown on corruption.

It was unclear whether Mr.~Kang had been suspended from his corporate or party posts. A press
official at the nuclear corporation hung up in the midst of a request to comment on the report.

With more than 100 sub-units, 14 scientific institutes and 280,000 employees, China National Nuclear
is among a select group of state-owned companies whose business is monitored at the highest levels
of the national government. Its businesses cover everything from uranium mining to reactor design
and construction to the building of nuclear weapons.

The company has increasingly focused on nuclear power, and turned a profit of \$701 million last
year. China has 11 nuclear power plants, with another 24 on the drawing boards or under way,
according to the World Nuclear Association, an industry group.

Mr.~Kang took over the reins at China Nuclear with a pledge to transform it from a stodgy
state-owned enterprise to a modern profit-oriented company. The company's Hong Kong subsidiary, CNNC
international Limited, recently bought a Canadian uranium mining firm, Western Prospector Group Ltd.

Mr.~Kang was named one of China's top business managers in 2007.

Several members of China's powerful elite have fallen to corruption charges in the past year, led by
Xu Zongheng, the politically connected mayor of Shenzhen, who lost his job in mid-June. Last month
saw the conviction of Chen Tonghai, the former chairman of the Chinese oil company Sinopec, on
charges of taking some \$29 million in bribes from 1999 to mid-2007.

\section{In First Return to Japan Court, Jurors\cite{juror} Convict and Sentence}

\lettrine{J}{apan}'s first jury trial in more than a half-century concluded Thursday as a panel of
three judges and six lay jurors convicted a Tokyo man of murder and then sentenced him to 15 years
in prison.

The jury in Tokyo District Court convicted Katsuyoshi Fujii, 72, of fatally stabbing a 66-year-old
female neighbor with a survival knife.

Mr.~Fujii had pleaded guilty, saying he attacked the woman after she knocked down bottles of water
he had placed around his house to ward off cats, according to local news reports.

In the most dramatic change to Japan's criminal justice system since the end of World War II, juries
are now empowered to decide guilt and sentencing. And in marked contrast to most Western judicial
systems, jurors may directly question witnesses, defendants and victims.

Although reticent\cite{reticent} at first, all six lay jurors asked questions during the trial,
which lasted just four days. Criminal trials in Japan sometimes drag on for years, and the lack of
transparency in the judicial process has often been frustrating to the public.

The verdict closed a week of intensive coverage by Japanese media, with up-to-the-minute reports
from the courthouse about details such as the questions being asked and the jurors' facial
expressions as they looked at bloody photos from the scene.

Juries in criminal trials were discontinued in Japan in 1943, as World War II intensified. A 2004
law reinstituted the jury system, beginning in 2009.

But opinion polls have shown the Japanese public to be highly skeptical of the jury system,
primarily because of deep cultural aversions\cite{aversion}, including a reluctance to express one's
opinions in public, to argue with colleagues and to question authority. Many Japanese have also
expressed reluctance to participate in cases involving capital punishment.

To educate and win over the public, Japanese courts have held hundreds of mock\cite{mock} trials in
recent years.

The jurors, known as ``saiban-in,'' are selected from the election rolls and must be at least 20
years old. They also must have at least a junior-high level education. Politicians and lawyers are
excluded.

Six lay jurors are teamed with three professional judges, with each of the nine having one vote.
Even if all three judges vote that a defendant is guilty, a not-guilty ruling by at least five of
the jurors will prevail\cite{prevail}. But even if all six jurors vote guilty, the ruling will not
stand unless at least one of the judges agrees.

Some legal analysts are concerned that the judges will dominate the jury deliberations and perhaps
skew the verdicts. Others see them as providing valuable guidance to uncertain citizen-jurors.

``The notion that judges should serve as educators for lay judges was supported by the long-standing
belief among Japanese elites that the average Japanese citizen lacks the political maturity to
participate in governance,'' the legal scholar Ingram Weber wrote in the East Asia Law Review. ``It
was this notion that was used to torpedo\cite{torpedo} previous attempts to introduce lay
participation in criminal justice.''

The new system, Mr.~Weber said, has been ``designed to democratize the criminal legal process'' and
serve as a way of ``injecting a measure of common sense and public values into court decisions by
giving lay persons a determining power over fact-finding and sentencing.''

According to people at the court cited Thursday by Kyodo News, four of the six jurors were women --
an office worker, 50; a piano teacher, 51; a nutritionist, 41; and a contract employee, 38. The two
men were an office worker, 43, and a part-time worker, 61.

After the trial, the part-timer was quoted by Kyodo as saying, ``I feel I have carried out an
important duty.''

\section{Sotomayor Confirmed by Senate, 68-31}

\lettrine{V}{oting} largely along party lines, the Senate on Thursday confirmed Judge Sonia
Sotomayor as the 111th justice of the Supreme Court. She will be the first Hispanic and the third
woman to serve on the court.

Chief Justice John G.~Roberts Jr.~was expected to administer the oath\cite{oath} of office to Judge
Sotomayor, 55, in the next few days, with a formal ceremony likely in September. She succeeds
Justice David H.~Souter, who retired in June.

Democrats celebrated the successful nomination and relatively smooth confirmation process as a
bright spot in a summer when they have been buffeted by several challenges, including rocky progress
on their attempts to overhaul the nation's health care system, President Obama's falling approval
ratings, the climbing unemployment rate and other lingering economic problems.

Shortly after the vote, President Obama said he was ``deeply gratified'' and confident that Judge
Sotomayor would become an outstanding justice. The ideals of ``justice, equality, opportunity'' that
guide the high court are the very ones that made the judge's ``uniquely American story'' possible in
the first place, the president said.

Judge Sotomayor's confirmation was never in much doubt, given Democrats' numerical advantage in the
Senate. But the final vote -- 68 to 31 -- represented a partisan divide. No Democrat voted against
her, while all but 9 of the chamber's 40 Republicans did so. Senator Edward M.~Kennedy, Democrat of
Massachusetts, is ailing\cite{ail} and did not vote.

During three days of debate on the Senate floor, Republicans labeled Judge Sotomayor a liberal
judicial activist, decrying several of her speeches about diversity and the nature of judgments, as
well as her votes in cases involving Second Amendment rights, property rights and a
reverse-discrimination claim brought by white firefighters in New Haven.

``Judge Sotomayor is certainly a fine person with an impressive story and a distinguished
background,'' the Senate minority leader, Mitch McConnell, Republican of Kentucky, said this week.
``But a judge must be able to check his or her personal or political agenda at the courtroom door
and do justice evenhandedly, as the judicial oath requires. This is the most fundamental test. It is
a test that Judge Sotomayor does not pass.''

But Democrats portrayed Judge Sotomayor as a mainstream and qualified judge whose life -- rising
from a childhood in a Bronx housing project to the Ivy League and now the Supreme Court -- is a
classic American success story. And they called her judicial record moderate and mainstream.

``Judge Sotomayor's career and judicial record demonstrates that she has always followed the rule of
law,'' Senator Patrick J.~Leahy, Democrat of Vermont and chairman of the Senate Judiciary Committee,
said Thursday. ``Attempts at distorting that record by suggesting that her ethnicity or heritage
will be the driving force in her decisions as a justice of the Supreme Court are demeaning to women
and all communities of color.''

From the moment Mr.~Obama chose her in May, many political strategists warned Republicans that
opposing the first Latina nominated to the Supreme Court would jeopardize the party in future
elections. In the waning days of the debate, some Democrats sought to portray Republican opposition
as a grave insult to Latinos.

``Republicans will pay a price for saying 'no' to this judge,'' Senator Robert Menendez, Democrat of
New Jersey, said in Spanish at a news conference Wednesday.

And in July, the National Rifle Association, which historically has stayed out of judicial
nomination fights, came out against Justice Sotomayor and said it would include senators'
confirmation vote in its legislative scorecard on gun-rights issues for the 2010 election -- a
pointed threat to Democrats from conservative-leaning states.

But attempts to appeal to interest-group politics in the confirmation process largely
faltered\cite{falter}.

The final vote was ``a triumph of party unity over some of the interest group politics that you
would have expected to play a bigger role,'' said Curt Levey, executive director of the conservative
Committee for Justice, which opposed Judge Sotomayor's confirmation.

Many Republicans took pains to emphasize that their vote against Judge Sotomayor did not mean they
were anti-Latino. They praised her credentials and her biography, saying they were troubled only by
what they said was her judicial philosophy.

Before announcing his opposition to her nomination, Senator John McCain of Arizona, last year's
Republican presidential nominee who has been sympathetic to calls by Latinos and others for
reforming the nation's immigration laws, first described her as an ``immensely qualified candidate''
with an ``inspiring and compelling'' life story. And he dwelled on his support for Miguel Estrada,
an appeals-court nominee of President George W.~Bush whom Democrats blocked from a vote even though
``millions of Latinos would have taken great pride in his confirmation,'' Mr.~McCain said.

Many other Republicans echoed Mr.~McCain's approach in explaining their votes. On Thursday, for
example, Senator Orrin G.~Hatch, Republican of Utah, spoke at length about the ``unfair and
disgraceful'' treatment of Mr.~Estrada, while criticizing Judge Sotomayor's record.

``I wish President Obama had chosen a Hispanic nominee whom all senators could support,'' Mr.~Hatch
said.

Juan Hern\'andez, who served as Hispanic outreach coordinator for Mr.~McCain's presidential
campaign, said most Republicans had not done enough to persuade Hispanics that they were welcome in
the party.

``It's not good enough to give two or three lines about Hispanics and then say, 'No, I'm not going
to vote for Sotomayor,' `` he said. ``We're just losing Hispanics left and right. It's amazing, in
the Republican Party -- we're doing it to ourselves.''

But Manuel A.~Miranda, chairman of the Third Branch Conference, a coalition of conservatives who
opposed the Sotomayor nomination, said Hispanics were ideologically diverse and would understand
that Republican opposition to a particular liberal-leaning judge did not mean they were hostile to
Hispanics -- especially since her confirmation hearing was civil, he said.

``Hispanics are not going to be offended by the opposition because Republicans didn't torment her,''
Mr.~Miranda said. ``Republicans can take this vote because they treated her well.''

For many Hispanic voters, the symbolism of the first Latina joining the Supreme Court -- and the
memory of who opposed her -- could be all that lingers, said Janet Murguía, president of the
National Council of La Raza, an Hispanic advocacy group.

``This is a singularly definitive historic moment,'' she said. ``So it is a vote, I think, that will
matter to the Latino community and will be remembered by the Latino community.''

A few people have argued that Justice Benjamin Cardozo, who served in the 1930s, should be
considered technically the first Hispanic on the Supreme Court. He was a Sephardic Jew whose family
believed its ancestors came from Portugal.

But the term ``Hispanic'' was uncommon during that era, and it usually means people from the
Americas with a Spanish-language heritage. On Thursday, the Hispanic National Bar Association hailed
her as the first Hispanic justice

What also remains to be seen is whether Democratic senators -- especially those from
conservative-leaning states and those who have received high ratings from the National Rifle
Association in the past -- will pay a political price for voting to confirm Judge Sotomayor despite
the group's opposition.

Andrew Arulanandam, an N.R.A. spokesman, declined to comment about the vote, but he did say it was
too early to know how much weight his group would give to the Sotomayor vote when putting together
its scores and endorsements for the 2010 election cycle.

Still, despite the seeming impotence of the gun-rights group's ability to intervene in the
nomination fight, Mr.~Miranda said he believed the threat of lower ratings might have had led more
Republicans to vote against Judge Sotomayor, noting that many had cited her alleged lack of support
for Second Amendment rights in explaining their votes.

``That was a seismic\cite{seismic} shift,'' Mr.~Miranda said.

Matthew Dowd, a former political adviser to Mr.~Bush who had warned Republicans to be civil,
disagreed. He said the Supreme Court confirmation process had simply become increasingly polarized
along party lines, regardless of a nominee's qualifications or the stance of groups like the
National Rifle Association.

``My view is that gun rights had nothing to do with it,'' he said. ``Supreme Court nominations have
become dodgeball games, with Democrats lining up on one side and Republicans lining up on our
side.''

\section{Journalists Entered N. Korea, Sister Says}

\lettrine{O}{ne} of the two journalists held captive for months in North Korea conceded to relatives
after her release on Wednesday that she had ``very briefly'' crossed into the Communist country, her
sister said Thursday.

Lisa Ling, the television personality, revealed that her sister -- Laura Ling, 32, a reporter with
Current TV -- planned to write an editorial explaining the events that led North Korean authorities
to detain her and a fellow American journalist, Euna Lee, 36. But Ms.~Ling, speaking by telephone on
CNN Thursday evening in her first interview since the release, revealed that the two women did
apparently cross into North Korea from China as they were researching a story on human trafficking
on March 17. The women were subsequently convicted and sentenced to 12 years hard labor.

``She did say that they touched North Korean territory very, very briefly,'' Ms.~Ling said of her
sister, adding later, ``She said that it was maybe 30 seconds, and everything just got sort of
chaotic. It's a very powerful story and she does want to share it.''

Ms.~Ling, a former co-host of The View, said that she would leave it to her sister to divulge the
remaining details for her captivity.

The women's return to the United States on Wednesday was set in motion by former president Bill
Clinton, who traveled to the reclusive country earlier this week to negotiate their release with the
North Korean leader, Kim Jong-il, who granted them a pardon. But until Thursday, it was unclear
whether the women had actually crossed the border into North Korea or had simply been captured on
the Chinese side.

Either way, what followed were more than two months of harrowing confinement, Ms.~Ling said on
Thursday. She said that her sister had spent much of her time in isolation, with only three small
meals a day that consisted primarily of rice with rocks mixed in, small vegetables, and fragments of
fried fish ``which she developed a reaction to.'' Ms.~Ling was confined to a small cell, where she
would walk in circles for exercise, read books, and occasionally bathe.

``Bathing was a little bit difficult because they didn't have hot water,'' Ms.~Ling said, ``So she
would fill up buckets and she would say, 'O.K., on Saturday I'm going to wash my hair.'''

Among Ms.~Ling's only human contact, her sister said, was with the guards who were assigned to her
cell.

``I think she won a lot of her captors over,'' Ms.~Ling said. ``She had some really lovely things to
say about the people who were watching her. She had two guards in her room at all times. And even
though they couldn't speak together, they developed a sort of strange kinship.''

In her short time home, Ms.~Ling has been racked by a fear of being alone. On Wednesday, exhausted
from her ordeal, Ms.~Ling wanted to take a brief nap but pleaded for her sister to be there when she
woke up, her sister said. She added that her sister's colleague, Ms.~Lee, has also been trying to
adjust to her new life at home with her family.

``I hear that Euna's 4-year old daughter does not want her to leave her sight,'' Ms.~Ling said.
``She keeps following her from room to room.''

\section{Migrants to China's West Bask\cite{bask} in Prosperity}

\lettrine{T}{hey} marched through the streets of Beijing, Shanghai and countless small towns
propelled by patriotic cheers and thumping drums. It was 1956, and Mao Zedong was calling on China's
youth to ``open up the west,'' the vast borderland known as Xinjiang that for centuries had defied
subjugation.

After a monthlong journey by train and open-air truck, thousands arrived at this Gobi Desert army
outpost to find that the factory jobs, hot baths and telephones in every house were nothing but
empty promises to lure them to a faraway land.

``We lived in holes in the ground, and all we did night and day was hard labor,'' recalled Han
Zuxue, a sun-creased 72-year-old who was a teenager when he left his home in eastern Henan Province.
``At first we cried every day but over time we forgot our sadness.''

More than five decades of toil later, men and women like Mr.~Han have helped transform Shihezi into
a tree-shaded, bustling oasis whose canned tomatoes, fiery\cite{fiery} grain alcohol and enormous
cotton yields are famous throughout China.

This city of 650,000 is a showcase of the Xinjiang Production and Construction Corps, a uniquely
Chinese conglomerate\cite{conglomerate} of farms and factories that were created by decommissioned
Red Army soldiers at the end of the civil war.

``Put your weapons aside and pick up the tools of construction,'' one popular slogan went. ``Develop
Xinjiang, defend the nation's borders and protect social stability.''

With a total population of 2.6 million, 95 percent of it ethnic Han Chinese, Shihezi and a string of
other settlements created by the military are stable strongholds in a region whose majority non-Han
populace has often been unhappy under Beijing's rule. Last month, that discontent showed itself
during vicious ethnic rioting that claimed 197 lives in Urumqi, the regional capital, which is a
two-hour drive away.

The government says that most of the dead were Han Chinese bludgeoned by mobs of Uighurs, Muslims of
Turkish ancestry whose presence in Xinjiang has been steadily diluted by migration from China's
densely populated east.

``Ever since we arrived they've resented us and had no appreciation for how we've improved this
place,'' said He Zhenjie, 76, who has spent his adult life leveling sand dunes, planting trees and
digging irrigation ditches. ``But we're here to stay. The Uighurs will never wrest Xinjiang away.''

Even if many Uighurs view the settlers as nothing more than Chinese colonists, many Chinese consider
the bingtuan, meaning soldier corps, a major success. In one fell swoop Mao deployed 200,000 idle
soldiers to help develop and occupy a resource-rich, politically strategic region bordering India,
Mongolia and the Soviet Union, a onetime ally turned menace\cite{menace}.

Shihezi and other bingtuan settlements quickly became self-sufficient, a relief to a government
lacking resources, and its ``reclamation warriors'' worked without pay those first few years,
steadily turning thousands of acres of inhospitable scrubland into some of the country's most
fertile terrain.

With an annual output of goods and services of \$7 billion, the settlements run by the bingtuan
include five cities, 180 farming communities and 1,000 companies. They also report directly to
Beijing and run their own courts, colleges and newspapers.

``During peaceful times, they are a force for development, but if anything urgent happens, they will
step out and maintain social stability and combat the separatists,'' said Li Sheng, a researcher at
the Chinese Academy of Social Sciences and a former bingtuan member who writes about the region's
history.

In those early years, the ranks of the bingtuan were fortified by petty criminals, former prisoners
of war, prostitutes and intellectuals, all sent west for ``re-education.'' During the mid-1950s,
40,000 young women were lured to Xinjiang with promises of the good life: they arrived to discover
their main purpose was to relieve the loneliness of the male pioneers and cement the region's Han
presence through their progeny.

Demographics have always been a tactical\cite{tactical} element of the campaign to pacify the
region. In 1949, when the Communists declared the establishment of the People's Republic of China,
there were just 300,000 Han Chinese in Xinjiang. Today, the number of Han has grown to 7.5 million,
just over 40 percent of the region's population. The percentage of Uighurs has fallen to 45 percent,
or about 8.3 million.

Their grievances have multiplied even as Xinjiang has grown more prosperous, thanks in part to its
huge reserves of natural gas, oil and minerals. Many Uighurs complain about the repression of their
Islamic faith, official policies that marginalize their language and a lack of job opportunities,
especially at government bureaus and inside the bingtuan.

During a recent visit to Shihezi, armed paramilitary policemen stopped every car and bus entering
the city. But only Uighurs were made to step out of vehicles for identification checks and searches.

Neatly laid out on a grid, its sidewalks graced by apple trees and elms, the city is populated by
the sturdy and defiantly proud who think of Xinjiang as China's version of Manifest Destiny, the
doctrine undergirding the westward expansion of the United States in the 19th century. But just
beneath the self-satisfaction runs a deep vein of bitterness, especially among those who arrived in
the 1950s and 1960s.

``I thought I was going to be a nurse, but I ended up sweeping the streets and cleaning toilets,''
said Yue Caiying, who moved here in 1963, and, like many of those with an education, was forced to
set aside personal ambition.

Lu Yiping, an author who spent five years interviewing women trucked into Xinjiang from Hunan
Province, tells of girls lured with promises of Russian-language classes and textile-mill jobs. In
an interview published online, he told the story of arriving women greeted by Wang Zhen, the
famously hard-line general who helped tame the region. ``Comrades, you must prepare to bury your
bones in Xinjiang,'' he quoted Mr.~Wang as telling the women.

Still, for many early settlers, Xinjiang offered an escape from the deprivation that stalked many
rural areas between 1959 and 1962, when Mao's disastrous attempt to start up China's
industrialization led to famine that killed millions.

Early settlers like Ma Xianwu, who arrived here in 1951 and helped dig the first thatch-covered pits
that served as shelter, offer a typical mix of conflicted emotions. He expressed wonder at the city
he had helped create, but also sorrow over the hardship he and others had endured.

``People would lose ears and toes to frostbite,'' said Mr.~Ma, who is 94 and nearly toothless.

But any sense of bitterness has faded. ``We were serving the motherland,'' he said, waving off the
adulation of a visitor. ``The glory belongs to the party. I'm just one drop of water in the ocean.''

\section{Lawyer's Detention Shakes China's Rights Movement }

\lettrine{C}{hina}'s nascent\cite{nascent} legal rights movement, already reeling\cite{reel} from a
crackdown on crusading\cite{crusade} lawyers, the kidnapping of defense witnesses and the shuttering
of a prominent legal clinic, has been shaken by the detention of a widely respected rights defender
who has been incommunicado since the police led him away from his apartment 12 days ago.

The man who was detained, Xu Zhiyong, 36, a soft-spoken and politically shrewd\cite{shrewd} legal
scholar who has made a name representing migrant workers, death row inmates\cite{inmate} and the
parents of babies poisoned by tainted milk, is accused of tax evasion\cite{evasion}. The charge is
almost universally seen here as a cover for his true offense: angering the Communist Party
leadership through his advocacy of the rule of law.

If convicted, he could face up to seven years in prison.

``We're all shocked by his detention, because Xu Zhiyong has always tried to avoid taking on radical
and politically sensitive cases,'' said Teng Biao, a colleague. ``His only interest is fighting for
the rights of the vulnerable and trying to enhance China's legal system.''

Mr.~Teng helped Mr.~Xu establish the Open Constitution Initiative, a six-year-old nonprofit legal
center that the authorities closed last month, charging that it was improperly registered and that
it failed to pay taxes.

Mr.~Xu is not the first rights advocate in China to face the wrath\cite{wrath} of the authorities in
recent years. Gao Zhisheng, a vocal lawyer, vanished into police custody six months ago, and Chen
Guangcheng, a blind lawyer, was beaten and then jailed after exposing abuses in China's
birth-control program.

Although rights lawyers and grass-roots social organizations have always been tightly controlled
here, the pressure has intensified in recent weeks. More than 20 lawyers known for taking on
politically tinged\cite{tinge} cases were effectively disbarred\cite{disbar}, and the police raided
a group that works to ease discrimination against people with Hepatitis B.

Last week, China's justice minister gave a speech saying that lawyers should above all obey the
Communist Party and help foster a harmonious society. To improve discipline, the minister said, all
law firms in the country would be sent party liaisons\cite{liaison} to ``guide their work.''

But given Mr.~Xu's international stature and reputation for working within the law, legal scholars
both in China and abroad say his prosecution suggests a new level of repression.

``What makes his detention particularly disturbing is that he's a special figure in so many ways,''
said Paul Gewirtz, director of the China Law Center at Yale Law School, which helped Mr.~Xu
establish his legal center, known here by its Chinese name, Gongmeng. ``He's at the forefront of
advancing the rule of law, which is something everyone agrees China needs for its ongoing
development.''

After 30 years of reform, China's legal system is at a critical juncture\cite{juncture}. Law schools
continue to pump out thousands of graduates each year, and the courts, even if imperfect, have
increasingly become a forum for resolving disputes. Late last month the Supreme People's Court
announced reforms intended to markedly reduce executions.

But as lawyers here discover, there are limits to China's embrace of judicial reform.

The Constitution, which includes guarantees of free speech and human rights, is unenforceable in
court. Judges routinely ignore evidence, making determinations based on political considerations.
And when it comes to vaguely defined offenses like ``subversion of state power'' or the invoking of
``state secrets'' laws, even the best-trained lawyers are powerless to defend the accused.

He Weifang, a law professor and legal adviser to Gongmeng, said conservative forces in the Communist
Party were increasingly wary of lawyers, who they suspect are ultimately seeking to challenge
one-party rule. Their greatest fear, Mr.~He said, is that advocacy lawyers and civil society
organizations could one day lead a pro-democracy movement among the poor and
disenfranchised\cite{disenfranchise} citizens they represent.

``What the authorities don't appreciate, though, is that lawyers are leading these people to the
courts, where their complaints can be resolved by rule of law,'' he said. ``People like Xu Zhiyong
can only help the government solve some of the problems it faces.''

According to Gongmeng, Mr.~Xu is being held at the Beijing No.~1 Detention Center, although public
security officials have not confirmed that he is in their custody\cite{custody}. Peng Jian, a lawyer
who is advising Gongmeng, said the authorities had imposed a \$208,000 penalty for nonpayment of
taxes due on donations from Yale.

A day after the raid on Gongmeng's office, Mr.~Xu held a news conference to say that the accusations
were baseless. He described the attack on his research center as a battle between corrupt officials
and society's most vulnerable citizens. ``We believe conscience will surely triumph over the evil
forces,'' he said.

A week later, police officers came to his door and led him away. Another employee of the research
center, Zhuang Lu, was also taken away the same day.

Soon after graduating from Peking University law school, Mr.~Xu became immersed\cite{immersed} in
the case of a 27-year-old graphic artist who was beaten to death in 2003 in police custody in the
southern city of Guangzhou. The artist, Sun Zhigang, had been arrested under vagrancy\cite{vagrancy}
laws that allowed the police to detain people for traveling outside their registered hometowns
without a permit.

Mr.~Xu led a campaign to end the practice, which gained widespread media attention. A few months
later, the State Council abolished the system.

That same year Mr.~Xu rose to the defense of a muckraking editor jailed in Guangzhou after his
newspaper, Southern Metropolis, ran a series of articles about Mr.~Sun's death. The editor, Cheng
Yizhong, said Mr.~Xu helped rally lawyers and journalists, leading to his release five months later.
``Only Xu had the courage to take on my case,'' he said.

More recently, he tried to build a case against black jails, the illegal holding cells that some
officials use to silence persistent critics. Last year, friends say, he was roughed up several times
while gathering evidence from petitioners who had come to Beijing to press their grievances to the
central government.

Raised in a Christian home in Henan Province, Mr.~Xu was fond of noting his birth in a county called
Minquan, which translates as ``civil rights.'' In an interview last year with The Economic Observer,
a Chinese weekly, he said this had a profound impact on his social consciousness.

``I strive to be a worthy Chinese citizen, a member of the group of people who promote the progress
of the nation,'' he said. ``I want to make people believe in ideals and justice, and help them see
the hope of change.''

\section{Online Market Flourishes in China }

\lettrine{I}{n} the months leading up to his college graduation in June, Yang Fugang spent most of
his days away from campus, managing an online store that sells cosmetics, shampoo and other goods he
often buys from local factories.

Today, his store on Taobao.com -- China's fast-growing online shopping bazaar -- has 14 employees,
two warehouses and piles of cash.

``I never thought I could do this well,'' said Mr.~Yang, 23, who earned \$75,000 last year. ``I
started out selling yoga mats and now I'm selling a lot of makeup and cosmetics. The profit margins
are higher.''

Taobao fever has swept Mr.~Yang's school, Yiwu Industrial and Commercial College, where
administrators say a quarter of its 8,800 students now operate a Taobao shop, often from a dorm
room.

Across China, millions of others -- recent college graduates, shopkeepers and retirees -- are also
using Taobao to sell clothes, mobile phones, toys and just about anything else they can find at
neighborhood stores and wholesale markets or even smuggle out of factories.

Internet analysts say this booming marketplace -- reminiscent of the early days of eBay, when
Americans started emptying their attics\cite{attic} for online auctions -- has turned Taobao into
China's newest Internet darling.

Though just six years old, Taobao (Chinese for ``to search for treasure'') already has 120 million
registered users and 300 million product listings. Its merchants produced nearly \$15 billion in
sales last year.

The company claims that sales through its Web site are already larger than any Chinese retailer.
And, Internet analysts say, sales on its site this year will surpass Amazon.com's expected sales of
about \$19 billion.

``This is the next big segment for China's Internet,'' said Jason Brueschke, an Internet analyst at
Citigroup in Hong Kong. ``It's their Amazon and eBay combined.''

Like eBay, Taobao does not sell anything itself; it simply matches buyers and sellers. It has a firm
foothold in China because many parts of the country still have poor transportation and some local
authorities favor their own government-owned outlets, making the retailing system inefficient.

The global recession also left once-booming factories overflowing with goods the rest of the world
does not seem to want.

The so-called Taobao addicts are helping to pick up the slack in a sluggish economy. ``I can't live
without Taobao,'' said Zhang Kangni, a graduate student in Shanghai. ``First, it's cheaper. I found
a dress at a store in Shanghai. It's a Hong Kong brand that sells for \$175. I found it on Taobao
for \$33.''

But skeptics\cite{skeptic} ask: Can Taobao actually make a profit and emerge as a true Web
powerhouse?

The company is not publicly traded and therefore does not disclose financial information, but
listings are free on Taobao and the company makes no money from online transactions. Almost all
Taobao's \$200 million in revenue comes from advertising, which the company says covers virtually
all its operational costs.

The company has been criticized, however, for contributing to a flourishing trade in
counterfeit\cite{counterfeit} goods. Taobao brushes aside such criticism, saying it has a new
program that is effectively cracking down on counterfeits.

Company executives also say Taobao is poised\cite{poised} to earn huge profits, but that their first
priority is creating an online community.

``Our vision for Taobao is to build a consumer's paradise, where people can shop online and have
fun,'' Jonathan Lu, Taobao's president, said. ``If you make the company better and better, profits
will naturally follow.''

His confidence in Taobao's future comes from the company's lineage\cite{lineage}. It is a division
of the Alibaba Group, which was founded by Jack Ma. In the past decade, Mr.~Ma has created an
Internet conglomerate\cite{conglomerate} with strong financial backing from Yahoo, Goldman Sachs and
the Softbank Group of Japan. Yahoo owns about 40 percent of Alibaba.

Alibaba.com -- the conglomerate's flagship Web site -- connects small businesses from around the
world with Chinese exporters. Taobao.com does something similar for consumers who want to sell to
other consumers.

When Taobao was founded in 2003, it appeared to have no chance. EBay and its Chinese partner,
EachNet, controlled 90 percent of China's online shopping. But Mr.~Ma, a former English teacher,
quickly undermined eBay's fee-based service by offering free listings on Taobao, essentially giving
away ads to anyone who wanted to sell.

At the time, eBay executives ridiculed the strategy, with many repeating that ``free is not a
business model.''

But almost immediately, the site took off, and in 2006, eBay pulled out of China, citing dwindling
market share and large losses. Today, it is Taobao that commands 80 percent of China's e-commerce
market, according to iResearch.

``Taobao is dominant,'' said Richard Ji, an Internet analyst at Morgan Stanley in Hong Kong.
``They're like an online Wal-Mart.'' Mr.~Ji says Taobao is a threat not only to traditional
retailers but also to big Chinese Internet companies, like Baidu, a leading search engine, because
they are competing with Taobao for many of the same advertisers.

Taobao has thrived, Internet analysts say, because people do not need much capital to start online
stores. This year, Taobao says its site could help create half a million new jobs, mostly among
young people opening new online stores.

Bao Yifen, a 23-year-old recent college graduate, opened her clothing shop with a \$5,000 investment
in 2007. Today, her Taobao store has sales of about \$4,000 a month.

``Three times a week I go to the wholesale market,'' Ms.~Bao said. ``It's a huge market. About 70 to
80 percent of the stuff is factory leftovers. There are even some brands, but they just cut the
labels off.''

Items smuggled into China from Hong Kong, Europe or the United States are also sold on Taobao,
evading high import duties and enabling sellers to profit by undercutting the prices of merchandise
in regular stores. An Apple MacBook Air that sells for \$2,225 in Beijing, for instance, costs just
\$1,508 in Hong Kong, a difference of 33 percent.

Counterfeit goods are also readily available, even though Taobao claims to have removed two million
``fake branded goods'' from the site.

Nevertheless, many Taobao sellers acknowledge dealing in illegal goods.

``I work in an O.E.M. factory that produces laptops and electronic devices for Sony,'' said one such
seller, who identified himself Mr.~Feng, referring to an original equipment manufacturer that
produces goods for global companies. ``We have Sony's core technology and exactly the same raw
materials and components, so we set up our own store selling netbooks and laptops on Taobao.''

A spokesman for Sony, Takashi Uehara, said the company had no comment but was looking into the
matter.

Here in Yiwu, which claims to be the site of the world's biggest wholesale market, Taobao has
started to change the look of Yiwu Industrial and Commercial College.

The school's vice dean, Jia Shaohua, points out an area designated as a start-up site for students
seeking to get rich. He points to students taking orders by computer, packaging products, sorting
inventory and taking photos of the items for display online, then adds, ``Around the school now,
there is a whole Taobao industrial chain.''

Every afternoon, even this summer, when the school should be relatively empty, one can hear the
ripping sounds of tape being wrapped around boxes in a building that could pass for a United Parcel
Service shipping terminal.

``The students don't need a lot of money,'' Mr.~Jia said. ``They just get orders and go find the
items at local factories.''

Mr.~Yang, the cosmetics seller, has become a campus hero. He operates his own warehouses a few miles
from the school, in the basements of a pair of residential buildings.

Standing in his crowded warehouse, near boxes of Neutrogena sun block, hairpins, toothbrushes and a
wide assortment of cosmetics, Mr.~Yang says business could not be better.

``Soon, I'll reach \$150,000 a month in sales,'' he said, flashing a big grin.

\section{In a Digital Future, Textbooks Are History }

\lettrine{A}{t} Empire High School in Vail, Ariz., students use computers provided by the school to
get their lessons, do their homework and hear podcasts of their teachers' science lectures.

Down the road, at Cienega High School, students who own laptops can register for ``digital
sections'' of several English, history and science classes. And throughout the district, a Beyond
Textbooks initiative encourages teachers to create -- and share -- lessons that incorporate their
own PowerPoint presentations, along with videos and research materials they find by sifting through
reliable Internet sites.

Textbooks have not gone the way of the scroll yet, but many educators say that it will not be long
before they are replaced by digital versions -- or supplanted altogether by lessons assembled from
the wealth of free courseware, educational games, videos and projects on the Web.

``Kids are wired differently these days,'' said Sheryl R.~Abshire, chief technology officer for the
Calcasieu Parish school system in Lake Charles, La. ``They're digitally nimble. They multitask,
transpose and extrapolate. And they think of knowledge as infinite.

``They don't engage with textbooks that are finite, linear and rote,'' Dr.~Abshire continued.
``Teachers need digital resources to find those documents, those blogs, those wikis that get them
beyond the plain vanilla curriculum in the textbooks.''

In California, Gov. Arnold Schwarzenegger this summer announced an initiative that would replace
some high school science and math texts with free, ``open source'' digital versions.

With California in dire straits, the governor hopes free textbooks could save hundreds of millions
of dollars a year.

And given that students already get so much information from the Internet, iPods and Twitter feeds,
he said, digital texts could save them from lugging around ``antiquated, heavy, expensive
textbooks.''

The initiative, the first such statewide effort, has attracted widespread attention, since
California, together with Texas, dominates the nation's textbook market.

Many superintendents are enthusiastic.

``In five years, I think the majority of students will be using digital textbooks,'' said William
M.~Habermehl, superintendent of the 500,000-student Orange County schools. ``They can be better than
traditional textbooks.''

Schools that do not make the switch, Mr.~Habermehl said, could lose their
constituency\cite{constituency}.

``We're still in a brick-and-mortar, 30-students-to-1-teacher paradigm,'' Mr.~Habermehl said, ``but
we need to get out of that framework to having 200 or 300 kids taking courses online, at night,
24/7, whenever they want.''

``I don't believe that charters and vouchers are the threat to schools in Orange County,'' he said.
``What's a threat is the digital world -- that someone's going to put together brilliant \$200
courses in French, in geometry by the best teachers in the world.''

But the digital future is not quite on the horizon in most classrooms. For one thing, there is still
a large digital divide. Not every student has access to a computer, a Kindle electronic reader
device or a smartphone, and few districts are wealthy enough to provide them. So digital textbooks
could widen the gap between rich and poor.

``A large portion of our kids don't have computers at home, and it would be way too costly to print
out the digital textbooks,'' said Tim Ward, assistant superintendent for instruction in California's
24,000-student Chaffey Joint Union High School District, where almost half the students are from
low-income families.

Many educators expect that digital textbooks and online courses will start small, perhaps for those
who want to study a subject they cannot fit into their school schedule or for those who need a few
more credits to graduate.

Although California education authorities are reviewing 20 open-source high school math and science
texts to make sure they meet California's exacting academic standards in time for use this fall --
and will announce this week which ones meet state standards -- quick adoption is unlikely.

``I want our teachers to have the best materials available, and with digital textbooks, we could see
the best lessons taught by the most dynamic teachers,'' said John A.~Roach, superintendent of the
Carlsbad, Calif., schools. ``But they're not going to replace paper texts right away.''

Whenever it comes, the online onslaught -- and the competition from open-source materials -- poses a
real threat to traditional textbook publishers.

Pearson, the nation's largest one, submitted four texts in California, all of them already available
online, as free supplements to their texts.

``We believe that the world is going digital, but the jury's still out on how this will evolve,''
said Wendy Spiegel, a Pearson spokeswoman. ``We're agnostic, so we'll provide digital, we'll provide
print, and we'll see what our customers want.''

Most of the digital texts submitted for review in California came from a nonprofit group, CK-12
Foundation, that develops free ``flexbooks'' that can be customized to meet state standards, and
added to by teachers. Its physics flexbook, a Web-based, open-content compilation, was introduced in
Virginia in March.

``The good part of our flexbooks is that they can be anything you want,'' said Neeru Khosla, a
founder of the group. ``You can use them online, you can download them onto a disk, you can print
them, you can customize them, you can embed video. When people get over the mind-set issue, they'll
see that there's no reason to pay \$100 a pop for a textbook, when you can have the content you want
free.''

The move to open-source materials is well under way in higher education -- and may be accelerated by
President Obama's proposal to invest in creating free online courses as part of his push to improve
community colleges.

Around the world, hundreds of universities, including M.I.T. and King Fahd University of Petroleum
and Minerals in Saudi Arabia, now use and share open-source courses. Connexions, a Rice University
nonprofit organization devoted to open-source learning, submitted an algebra text to California.

But given the economy, many educators and technology experts agree that the K-12 digital revolution
may be further off.

``There's a lot of stalled purchasing and decision making right now,'' said Mark Schneiderman,
director of federal education policy at the Software \& Information Industry Association. ``But it's
going to happen.''

For all the attention to the California initiative, digital textbooks are only the start of the
revolution in educational technology.

``We should be bracing ourselves for way more interactive, way more engaging videos, activities and
games,'' said Marina Leight of the Center for Digital Education, which promotes digital education
through surveys, publications and meetings.

Vail's Beyond Textbooks effort has moved in that direction. In an Empire High School history class
on elections, for example, students created their own political parties, campaign Web sites and
videos.

``Students learn the same concepts, but in a different way,'' said Matt Donaldson, Empire's
principal.

``We've mapped out our state standards,'' Mr.~Donaldson said, ``and our teachers have identified
whatever resources they feel best covers them, whether it's a project they created themselves or an
interesting site on the Internet. What they don't do, generally, is take chapters from textbooks.'' 

\section{Sotomayor Sworn In as Supreme Court Justice }

\lettrine{S}{onia} Sotomayor took the judicial oath on Saturday, becoming the first Hispanic and the
third woman to serve on the Supreme Court.

At just past 11 a.m., Chief Justice John G.~Roberts Jr.~administered a pair of oaths to her in two
private ceremonies at the Supreme Court building, completing her ascent to a life-tenured position
as the nation's 111th justice -- the first to be nominated by a Democratic president since 1994.

In the first ceremony, she took the standard oath affirmed by all federal employees, swearing to
``support and defend the Constitution of the United States against all enemies, foreign and
domestic.'' Only the chief justice, Justice Anthony M.~Kennedy, Justice Sotomayor's immediate family
and a court photographer attended this ceremony. Her mother, Celina Sotomayor, held a Bible for the
ritual.

They then walked to the court's East Conference Room for the judicial oath, joined by several dozen
friends and family members, where Justice Sotomayor swore to ``administer justice without respect to
persons, and do equal right to the poor and to the rich, and that I will faithfully and impartially
discharge and perform all the duties incumbent upon me'' under the Constitution and laws of the
United States.

With that, Chief Justice Roberts -- who smoothly recited the words of the oath, which he appeared to
be reading from a paper as a precaution against the kind of stumble that prompted a do-over
swearing-in after President Obama's inauguration -- shook her hand with a smile, congratulated her
and welcomed her to the court, saying she could now ``begin her work as an associate justice without
delay.''

Justice Sotomayor hugged her mother and then shook the hands of several other people in the room.
She made no statement.

Other witnesses included President Obama's White House counsel, Gregory B.~Craig; Cynthia Hogan,
assistant White House counsel; Judge Robert A.~Katzmann of the United States Court of Appeals for
the Second Circuit; and Representative Nydia M.~Vel\'azquez, Democrat of New York, who is the
chairwoman of the Congressional Hispanic Caucus and who, like Justice Sotomayor, is Puerto Rican.

The Supreme Court appointment of Justice Sotomayor, 55, capped a life story that began in a Bronx
housing project, where she was raised by her widowed mother. She attended two Ivy League
universities, worked as a prosecutor in the office of the Manhattan district attorney, and then left
for a stint at a corporate law firm.

Throughout the 1980s, she became increasingly engaged in civic life, joining the board of the Puerto
Rican Legal Defense and Education Fund and winning appointments to the New York City Campaign
Finance Board and the board of the State of New York Mortgage Agency.

Her rise brought her to the attention of Senator Daniel Patrick Moynihan, Democrat of New York, who
recommended that President George Bush appoint her to the federal bench, which she joined in 1992.
President Bill Clinton later elevated her to the Court of Appeals for the Second Circuit. President
Obama nominated her in May to succeed Justice David H.~Souter.

Saturday's ceremonies ended a 10-week confirmation battle in which Republicans called her a
``judicial activist,'' citing her speeches on topics like foreign law and the role of diversity in
the judiciary, including a now-famous comment about the superior judgment of a ``wise Latina.''

They also criticized her votes in cases involving Second Amendment rights, property rights and a
racial discrimination lawsuit brought by white firefighters in New Haven.

But Democrats insisted that she was a highly qualified, mainstream and moderate jurist. In the end,
Democrats' numerical advantage in the Senate virtually ensured her success. The Senate on Thursday
voted 68 to 31 to confirm her. No Democrat voted against her, while all but 9 of the chamber's 40
Republicans did so.

Justice Sotomayor can now prepare for her first case -- a campaign-finance dispute left over from
the last term -- but she still has several rituals to attend.

President Obama, who did not attend Saturday's events, will be at a White House reception on
Wednesday in Justice Sotomayor's honor.

And on Sept.~8, the Supreme Court will hold a special invitation-only investiture ceremony, a
special sitting at which she will take her seat for the first time, probably followed by an
appearance alongside Chief Justice Roberts for a traditional picture on the front steps of the
court.

\section{Coffee Can Wait. The Day's First Stop Is Online. }

\lettrine{K}{arl} and Dorsey Gude of East Lansing, Mich., can remember simpler mornings, not too
long ago. They sat together and chatted as they ate breakfast. They read the newspaper and competed
only with the television for the attention of their two teenage sons.

That was so last century. Today, Mr.~Gude wakes at around 6 a.m. to check his work e-mail and his
Facebook and Twitter accounts. The two boys, Cole and Erik, start each morning with text messages,
video games and Facebook.

The new routine quickly became a source of conflict in the family, with Ms.~Gude complaining that
technology was eating into family time. But ultimately even she partially succumbed, cracking open
her laptop after breakfast.

``Things that I thought were unacceptable a few years ago are now commonplace in my house,'' she
said, ``like all four of us starting the day on four computers in four separate rooms.''

Technology has shaken up plenty of life's routines, but for many people it has completely altered
the once predictable rituals at the start of the day.

This is morning in America in the Internet age. After six to eight hours of network deprivation --
also known as sleep -- people are increasingly waking up and lunging for cellphones and laptops,
sometimes even before swinging their legs to the floor and tending to more biologically urgent
activities.

``It used to be you woke up, went to the bathroom, maybe brushed your teeth and picked up the
newspaper,'' said Naomi S.~Baron, a professor of linguistics at American University, who has written
about technology's push into everyday life. ``But what we do first now has changed dramatically.
I'll be the first to admit: the first thing I do is check my e-mail.''

The Gudes' sons sleep with their phones next to their beds, so they start the day with text messages
in place of alarm clocks. Mr.~Gude, an instructor at Michigan State University, sends texts to his
two sons to wake up.

``We use texting as an in-house intercom,'' he said. ``I could just walk upstairs, but they always
answer their texts.'' The Gudes recently began shutting their devices down on weekends to account
for the decrease in family time.

In other households, the impulse to go online before getting out the door adds an extra layer of
chaos to the already discombobulating morning scramble.

Weekday mornings have long been frenetic, disjointed affairs. Now families that used to fight over
the shower or the newspaper tussle over access to the lone household computer -- or about whether
they should be using gadgets at all, instead of communicating with one another.

``They used to have blankies; now they have phones, which even have their own umbilical cord right
to the charger,'' said Liz Perle, a mother in San Francisco who laments the early-morning technology
immersion of her two teenage children. ``If their beds were far from the power outlets, they would
probably sleep on the floor.''

The surge of early risers is reflected in online and wireless traffic patterns. Internet companies
that used to watch traffic levels rise only when people booted up at work now see the uptick much
earlier.

Arbor Networks, a Boston company that analyzes Internet use, says that Web traffic in the United
States gradually declines from midnight to around 6 a.m. on the East Coast and then gets a huge
morning caffeine jolt. ``It's a rocket ship that takes off at 7 a.m,'' said Craig Labovitz, Arbor's
chief scientist.

Akamai, which helps sites like Facebook and Amazon keep up with visitor demand, says traffic takes
off even earlier, at around 6 a.m. on the East Coast. Verizon Wireless reported the number of text
messages sent between 7 and 10 a.m. jumped by 50 percent in July, compared with a year earlier.

Both adults and children have good reasons to wake up and log on. Mom and Dad might need to catch up
on e-mail from colleagues in different time zones. Children check text messages and Facebook posts
from friends with different bedtimes -- and sometime forget their chores in the process.

In May, Gabrielle Glaser of Montclair, N.J., bought her 14-year-old daughter, Moriah, an Apple
laptop for her birthday. In the weeks after, Moriah missed the school bus three times and went from
walking the family Labradoodle for 20 minutes each morning to only briefly letting the dog outside.

Moriah concedes that she neglected the bus and dog, and blames Facebook, where the possibility that
crucial updates from friends might be waiting draws her online as soon as she wakes. ``I have some
friends that are up early and chatting,'' she said. ``There is definitely a pull to check it.''

Some families have tried to set limits on Internet use in the mornings. James Steyer, founder of
Common Sense Media, a nonprofit that deals with children and entertainment, wakes every morning at 6
and spends the next hour on his BlackBerry, managing e-mail from contacts in different parts of the
world.

But when he meets his wife, Liz, and their four children, ages 5 to 16, at the breakfast table, no
laptops or phones are allowed.

Mr.~Steyer says he and his sons feel the temptation of technology early. Kirk, 14, often runs
through much of his daily one-hour allotment of video-game time in the morning.

Even Jesse, 5, has started asking each morning if he can play games on his father's iPhone. And
Mr.~Steyer said he constantly feels the tug of waiting messages on his BlackBerry, even during
morning hours that are reserved for family time.

``You have to resist the impulse. You have to switch from work mode to parenting mode,'' Mr.~Steyer
said. ``But meeting my own standard is tough.''

\section{As Health Care Debate Rages, Obama Takes to the Stump\cite{stump} }

\lettrine{P}{resident} Obama tried on Tuesday to defuse fears about his plan to overhaul the
nation's health care system, an issue at the center of one of the fiercest public-policy debates in
decades, telling a friendly audience in New Hampshire that a lot of misinformation is being spread.

``If you like your health-care plan, you can keep your health-care plan,'' the president told a
gathering in Portsmouth, N.H. ``You will not be waiting in any lines. This is not about putting the
government in charge of your health insurance.''

``For all the chatter and the yelling and the shouting and the noise, what you need to know is this:
If you don't have health insurance, you will finally have quality, affordable options, once we pass
reform. If you do have health insurance, we will make sure that no insurance company, or a
government bureaucrat, gets between you and the care that you need.''

The president went on, to applause. ``And we will do this without adding to our deficit over the
next decade, largely by cutting out the waste and insurance company giveaways in Medicare that
aren't making any of our seniors healthier,'' he said.

The president said some Republicans had been helpful, and that he hoped to achieve health care
reform on a bipartisan basis, ``but the most important thing is getting it done for the American
people.''

The New Hampshire event was part of a campaign to fight questionable but potentially damaging
charges that the president's vision would inevitably lead to ``socialized medicine,'' ``rationed
care'' and even forced euthanasia for the elderly.

The tone of the national debate over the future of health care has become increasingly emotional,
even bitter, as reflected in comments by lawmakers across the political specrtum.

In introducing a Web site to defend the president's proposals, White House officials were tacitly
acknowledging a difficult reality: they are suddenly at risk of losing control of the public debate
over a signature issue for Mr.~Obama and are now playing defense in a way they have not since last
year's campaign.

Senator Bernard Sanders of Vermont, an independent who is one of the most liberal members in either
House, said on Tuesday that ``the Republicans are the party of do-nothingism, and because of them it
is very hard to move forward.''

But Mr.~Sanders said in an interview on MSNBC that ``frankly, the Democrats have not handled this as
clearly and effectively as they might have.''

A different perspective was offered by Representative Peter King, a Long Island Republican far to
the right of Mr.~Sanders. Mr.~King said it was quite understandable that many Americans are not
enthusiastic about ``the radical type of reform that President Obama's talking about.''

The health care system should be changed ``incrementally'' rather than by major surgery, Mr.~King
said in another interview on MSNBC. The congressman said he thought the White House had made a
tactical\cite{tactical} error in its approach on health care. ``It may not be perfect,'' he said,
conceding that Americans ``may have problems with it.''

``But it's not the rabid-type issue that had to be solved by Aug.~1 of this year, the way President
Obama was saying,'' Mr.~King said.

And Senator Arlen Specter of Pennsylvania, a Republican-turned-Democrat, faced a crowd of emotional
constituents in Lebanon, Pa. One participant drew loud applause when he said illegal immigrants
should not be covered (Mr.~Specter agreed), and another complained that the legislation was as
complicated ``as a Russian novel.''

Mr.~Specter said the federal government has a ``social compact'' with the American people to ``take
care of people who need some help.'' He pledged again not to vote for something that would add to
the budget deficit, and he tried to reassure the crowd that people who are happy with their present
health insurance do not have to worry about losing it.

``So far, no bill has passed the Congress,'' Mr.~Specter noted. ``In the House of Representatives,
five committees have passed bills, but the House has not passed a bill. In the Senate, we're still
working on a bill, trying to get bipartisanship.''

``I know the American people are sick and tired of Republicans and Democrats fighting, and the
American people would like to see some bipartisanship and coming together in the public interest,''
the senator said, in an appeal for calm that was not entirely successful.

President Obama, speaking at a summit of North American leaders in Mexico on Monday, sounded an
optimistic note, predicting that ``the American people are going to be glad that we acted to change
an unsustainable system so that more people have coverage.''

But aides to Mr.~Obama said the rapidly escalating\cite{escalate} threat to his health care plans
had led him to order them to come up with a crisper message.

And Democratic Party officials enlisted in the fight by the White House acknowledged in interviews
that the growing intensity of the opposition to the president's health care plans -- within the last
week likened on talk radio to something out of Hitler's Germany, lampooned by protesters at
Congressional town-hall-style meetings and vilified in television commercials -- had caught them off
guard and forced them to begin an August counteroffensive.

In the process, the administration has had a harder time getting across the themes it wanted to
strike in this period: that the current system is unsustainable and that Mr.~Obama's plan holds
concrete benefits for people who already have health insurance as well as for those who do not.

``We all had a good sense that some of this was going to take place,'' said Brad Woodhouse, the
communications director for the Democratic National Committee. ``To be fair, I think we were
probably a little surprised -- just a little -- at the use of swastikas and the comparisons to Adolf
Hitler and the Third Reich that even Rush Limbaugh has fanned the flames on. And we were a little
surprised at the mob mentality.'' (Mr.~Woodhouse's use of the phrase ``mob mentality'' was itself
part of the Democratic effort to paint opponents speaking out against the plan as part of an unruly
but organized effort.)

For some of Mr.~Obama's supporters, the newly galvanized opposition to his proposed policies
provided a troubling flashback to the successful effort to stop President Bill Clinton's similarly
ambitious plans 16 years ago -- a fight Mr.~Obama's aides had studied carefully to avoid making the
same fatal mistakes.

White House officials say such fears are unwarranted, arguing that the conservative protests are
getting outsize coverage on cable news. ``Don't associate loud with effective,'' Rahm Emanuel, the
White House chief of staff, said in an interview, adding that he detected no anxiety from supportive
lawmakers in politically vulnerable districts. ``What is coming across is a lot of noise and a lot
of heat without a lot of light.''

And White House officials say their August counteroffensive is a break from the Clinton approach,
which is now viewed as having failed to adequately address critics.

Mr.~Obama will take the lead this week as he continues a series of public meetings to counter the
opposition, events White House officials hope will offer a high-profile opportunity to confront and
rebut critics.

As part of the effort, Speaker Nancy Pelosi and Representative Steny H.~Hoyer, the House majority
leader, wrote an opinion article in USA Today on Monday calling conservative protests at
Congressional town-hall-style meetings ``un-American'' for ``drowning out opposing views.'' (That
prompted a swift rebuke from the House minority leader, Representative John A.~Boehner, among other
Republicans.)

New television commercials disputing the conservative attacks are in the works, Mr.~Woodhouse said,
and allied members of Congress have been sent home for the August break with a set of poll-tested
talking points intended to shift the focus to the administration's advertised benefits of the plan
from the scary situations opponents have laid out.

``There's a whole set of rumors that the old playbook would tell you not to do anything about
because you draw attention,'' said Dan Pfeiffer, the White House deputy communications director.
``The lesson we've learned is you ignore these rumors at your peril, and the right answer is to take
them head on in as big a way as possible.''

It was only weeks ago that Mr.~Obama was pressing both the House and the Senate to complete work on
their versions of health legislation before recessing for the summer, a goal that went unmet as
divisions erupted among Democrats as well as between the two parties.

After getting much of what he wanted on high-risk initiatives like the economic stimulus package and
bailouts of banks and auto companies, Mr.~Obama had yet to face the full force of conservative
opposition to his policies. Some supporters now wonder whether his earlier glide path left him
unprepared for the sudden surge of opposition from conservative groups, which have found a rallying
point on health care.

``The expectation was that things have gotten so bad in the last 16 years that there would be
consensus\cite{consensus} on the need to act this time,'' said Howard Paster, who was Mr.~Clinton's
chief lobbyist in 1993. ``That was a mistake, that assumption.''

Mr.~Obama's team won early, high marks for diverging from the Clinton approach, specifically by
emphasizing the need to control costs and improve coverage for those who are already insured instead
of making the same moral-duty argument Mr.~Clinton had about the need to cover the uninsured.

Yet once Congress started filling in the details this summer and its analysts began pricing the
House and Senate packages, the estimates of the government's cost caused sticker shock again.

And once again that drew taxpayers' attention to the main reason for those costs: covering the
uninsured, through more Medicaid spending and subsidies for people to buy insurance and small
businesses to provide it.

That helped conservatives who had been struggling to gain traction\cite{traction} on health care to
speak to a constituency that has managed to gain significant anti-Obama attention this year, the
fiscally hawkish ``tea party'' activists opposed to the president's spending. They have dismissed
Mr.~Obama's promises that his plan will be fully paid for through offsetting spending cuts or
increased taxes, and have cast the plan as a costly takeover of health care by the government.

``I think the combination of spending a trillion dollars that we don't have and another rushed
process really triggered this,'' said Matt Kibbe, the president of the conservative group
FreedomWorks. ``People started paying attention.''

\section{China's Incinerators Loom as a Global Hazard }

\lettrine{I}{n} this sprawling metropolis in southeastern China stand two hulking brown buildings
erected by a private company, the Longgang trash incinerators. They can be smelled a mile away and
pour out so much dark smoke and hazardous chemicals that hundreds of local residents recently staged
an all-day sit-in, demanding that the incinerators be cleaner and that a planned third
incinerator\cite{incinerator} not be built nearby.

After surpassing the United States as the world's largest producer of household garbage, China has
embarked on a vast program to build incinerators as landfills run out of space. But these
incinerators have become a growing source of toxic emissions, from dioxin to mercury, that can
damage the body's nervous system.

And these pollutants, particularly long-lasting substances like dioxin and mercury, are dangerous
not only in China, a growing body of atmospheric research based on satellite observations suggests.
They float on air currents across the Pacific to American shores.

Chinese incinerators can be better. At the other end of Shenzhen from Longgang, no smoke is visible
from the towering smokestack of the Baoan incinerator, built by a company owned by the municipal
government. Government tests show that it emits virtually no dioxin and other pollutants.

But the Baoan incinerator cost 10 times as much as the Longgang incinerators, per ton of
trash-burning capacity.

The difference between the Baoan and Longgang incinerators lies at the center of a growing
controversy in China. Incinerators are being built to wildly different standards across the country
and even across cities like Shenzhen. For years Chinese government regulators have discussed the
need to impose tighter limits on emissions. But they have done nothing because of a bureaucratic
turf war, a Chinese government official and Chinese incineration experts said.

The Chinese government is struggling to cope with the rapidly rising mountains of trash generated as
the world's most populated country has raced from poverty to rampant consumerism. Beijing officials
warned in June that all of the city's landfills would run out of space within five years.

The governments of several cities with especially affluent, well-educated citizens, including
Beijing and Shanghai, are setting pollution standards as strict as Europe's. Despite those
standards, protests against planned incinerators broke out this spring in Beijing and Shanghai as
well as Shenzhen.

Increasingly outspoken residents in big cities are deeply distrustful that incinerators will be
built and operated to international standards. ``It's hard to say whether this standard will be
reached -- maybe the incinerator is designed to reach this benchmark, but how do we know it will be
properly operated?'' said Zhao Yong, a computer server engineer who has become a neighborhood
activist in Beijing against plans for an incinerator there.

Yet far dirtier incinerators continue to be built in inland cities where residents have shown little
awareness of pollution.

Studies at the University of Washington and the Argonne National Laboratory in Argonne, Ill., have
estimated that a sixth of the mercury now falling on North American lakes comes from Asia,
particularly China, mainly from coal-fired plants and smelters but also from incinerators. Pollution
from incinerators also tends to be high in toxic metals like cadmium.

Incinerators play the most important role in emissions of dioxin. Little research has been done on
dioxin crossing the Pacific. But analyses of similar chemicals have shown that they can travel very
long distances.

A 2005 report from the World Bank warned that if China built incinerators rapidly and did not limit
their emissions, worldwide atmospheric levels of dioxin could double. China has since slowed its
construction of incinerators and limited their emissions somewhat, but the World Bank has yet to do
a follow-up report.

Airborne dioxin is not the only problem from incinerators. The ash left over after combustion is
laced with dioxin and other pollutants. Zhong Rigang, the chief engineer at the Baoan incinerator
here, said that his operation sent its ash to a special landfill designed to cope with toxic waste.
But an academic paper last year by Nie Yongfeng, a Tsinghua University professor and government
adviser who sees a need for more incinerators, said that most municipal landfills for toxic waste
lacked room for the ash, so the ash was dumped.

Trash incinerators have two advantages that have prompted Japan and much of Europe to embrace them:
they occupy much less real estate than landfills, and the heat from burning trash can be used to
generate electricity. The Baoan incinerator generates enough power to light 40,000 households.

And landfills have their own environmental hazards. Decay in landfills also releases large
quantities of methane, a powerful global warming gas, said Robert McIlvaine, president of McIlvaine
Company, an energy consulting firm that calculates the relative costs of addressing disparate
environmental hazards. Methane from landfills is a far bigger problem in China than toxic pollutants
from incinerators, particularly modern incinerators like those in Baoan, he said.

China's national regulations still allow incinerators to emit 10 times as much dioxin as
incinerators in the European Union; American standards are similar to those in Europe. Tightening of
China's national standards has been stuck for three years in a bureaucratic war between the
environment ministry and the main economic planning agency, the National Development and Reform
Commission, said a Beijing official who insisted on anonymity because he was not authorized to
discuss the subject publicly.

The agencies agree that tighter standards on dioxin emissions are needed. They disagree on whether
the environment ministry should have the power to stop incinerator projects that do not meet tighter
standards, the official said, adding that the planning agency wants to retain the power to decide
which projects go ahead.

Yan Jianhua, the director of the solid waste treatment expert group in Zhejiang province, a center
of incinerator equipment manufacturing in China, defended the industry's record on dioxin, saying
that households that burn their trash outdoors emit far more dioxin.

``Open burning is a bigger problem according to our research,'' Professor Yan said, adding that what
China really needs is better trash collection so that garbage can be disposed of more reliably.

Critics and admirers of incinerators alike call for more recycling and reduced use of packaging as
ways to reduce the daily volume of municipal garbage. Even when not recycled, sorted trash is easier
for incinerators to burn cleanly, because the temperature in the furnace can be adjusted more
precisely to minimize the formation of dioxin.

Yet the Chinese public has shown little enthusiasm for recycling. As Mr.~Zhong, the engineer at the
Baoan incinerator, put it, ``No one really cares.''

\section{G.M. Hopes New Lineup Brings Back Excitement }

\lettrine{G}{eneral} Motors is focusing on smaller cars and more fuel-efficient vehicles to lure
back consumers who had given up on the struggling auto giant.

On Tuesday, G.M. showed off a lineup that was decidedly leaner and greener than what the company
offered in the months leading to its bankruptcy filing on June 1.

And just as G.M. has shed jobs, plants and operations, the company has plans to shrink its portfolio
of gas-guzzling trucks and big cars in favor of smaller vehicles.

``Everything we are doing is consistent with the viability plan turned into the government,'' said
Thomas G.~Stephens, G.M.'s vice chairman for global product development. ``And everything we do has
to be best-in-class in fuel economy.''

The government has committed about \$50 billion this year to rescue G.M. from financial collapse.
In return, American taxpayers hold a 60 percent stake in the company.

But while the Obama administration helped clean up G.M.'s balance sheet, the company's management
team is faced with the tall task of bringing some excitement back to dealer showrooms.

G.M.'s sales in the United States have plunged 37 percent this year, and its market share has fallen
just below 20 percent, a far cry from the company's glory days in the 1960s, when it held a 50
percent share of the American market and dominated competitors. As part of its effort to reverse
that trend and return to profitability, G.M. also began selling new cars and trucks on the auction
Web site eBay on Tuesday.

G.M. executives found out just how far the company had fallen on Monday, when more than 100
consumers -- some of them vocal critics of the company -- were invited to test its latest products.

The event, held at G.M.'s proving grounds outside Detroit, was as humbling for some executives as
the months of debate over the company's federal bailout.

``I learned a lot,'' said Fritz Henderson, who became G.M.'s chief executive in April. ``The
awareness of our products was painfully, almost pitifully low.''

Mr.~Henderson has vowed to change both the perception of G.M.'s vehicles and make them more
competitive with models from other auto companies.

``People are looking at us with a healthy degree of skepticism,'' Mr.~Henderson said.

The company's core brands -- Chevrolet, Cadillac, Buick and GMC -- will be broadened to appeal to
younger buyers who value fuel economy and technology more than size and horsepower.

The biggest change is coming at Buick. Once the bastion of large, four-door sedans marketed to
consumers older than 60, the brand is getting a decidedly youthful makeover.

Buick's largest car, the Lucerne, will be discontinued. Within three years, the brand will add a
compact car, a midsize sedan and two small crossover vehicles that combine the attributes of a car
and a sport utility vehicle.

Chevrolet will also get two new cars -- the tiny Spark and the compact Cruze -- and the GMC brand
may add a small crossover that mimics the boxy styling of vehicles produced by Toyota, Nissan and
Hyundai.

G.M.'s strategy is similar to that employed by its largest domestic rival, the Ford Motor Company --
the only American automaker that has not needed federal money to survive.

Ford is introducing a series of passenger cars that it hopes will ease its long dependence on pickup
trucks and S.U.V.'s for profits.

Now that G.M. has come through bankruptcy with far less debt and a tighter cost structure, analysts
said the company must move quickly to recast its image and products.

``G.M. has no choice at this point,'' said Ron Harbour, head of the automotive division of the
consulting firm Oliver Wyman. ``They have got to convince people that their cars are better and
their brands are cool.''

The company's restructuring plan calls for cutting the number of models in G.M. dealerships to 34,
from 48, by next year. Several of the models will disappear when the Saturn, Pontiac, Hummer and
Saab brands are closed or sold.

G.M. executives conceded that the company could no longer afford any laggards in its new lineup.

``We have to do great cars no matter what segment we enter,'' said Mr.~Stephens, who succeeded
Robert A.~Lutz as G.M.'s product czar.

G.M.'s new board is pressing management to start increasing sales again after many years of decline.
That will prove difficult in the current economic environment, in which vehicle sales have sunk to
their lowest point in more than 25 years.

But the company has some promising entries already entering the market.

The Chevrolet Equinox crossover vehicle has been a hot seller, as has the Camaro muscle car. At
Cadillac, a sharply restyled version of the SRX is winning raves from some consumers, even as its
angular appearance polarizes others.

``You might call the design aggressive, where I call it distinctive,'' Mr.~Stephens said. ``Its
looks are purely Cadillac. When you drive up, everybody knows what it is.''

Design will play a critical role in G.M.'s comeback, particularly when it begins introducing smaller
Buicks that could compete with the company's Chevy products.

Mr.~Stephens said that Buick would be focusing on owners of some foreign luxury brands like Honda's
Acura division. The new Buicks will be packed with new technology and upscale materials, which G.M.
hopes will allow the brand to command higher prices.

But G.M. has run into trouble in the past when its brands overlapped and too many of its models
looked alike.

``Ford has already started down this road, and the competition is hardly letting up,'' said Joseph
Phillippi, a principal in the firm AutoTrends Consulting. ``G.M. has obviously taken a lot of costs
out. But lower costs aren't going to help them sell cars that people aren't crazy about.''

\section{No.~1 Fugitive\cite{fugitive} of China Fears Death if Repatriated }

\lettrine{T}{he} lawyer for China's most wanted man says that his admission of guilt to smuggling
will not end his decade-long bid to stay in Canada.

The lawyer, David Matas, said Sunday that Lai Changxing's fight to stay in Canada was based on a
fear of torture and possible execution in China, not on his guilt or innocence.

Mr.~Lai's admission of guilt was published in a Canadian Chinese-language newspaper last week. He
said that he had committed tax evasion and that he had circumvented Chinese customs by waiting until
customs officers left for the day, and then unloading shipments of goods including gasoline and
cigarettes.

Mr.~Lai, who fled to Canada in 1999, has been described in Chinese news reports as the country's
No.1 fugitive. The Chinese authorities accuse him of masterminding a network that smuggled goods
worth as much as \$10 billion, under the protection of corrupt government officials.

Reports of Mr.~Lai's admission came to light after it was revealed that his former wife, Tsang
Mingna, had returned to China with their eldest daughter. Mr.~Matas said that Ms.~Tsang had made a
deal to return to China in exchange for all charges against her being dropped.

That approach by China could be an effort to convince Mr.~Lai that he will receive lenient treatment
should he return, Mr.~Matas suggested. ``It's pretty clear they're using the wife to get at him,''
he said.

China has given assurances that Mr.~Lai will not face execution if he returns, but Mr.~Matas said
the assurances could not be believed. The two countries do not have an extradition treaty.

Mr.~Lai, his former wife, and their three children applied for asylum\cite{asylum} after arriving in
Canada. They said that the refugee board that turned down their asylum requests in 2002 failed to
see that they were at risk of political persecution in China.

The asylum case has proceeded through court challenges and appeals since then. Meanwhile, Mr.~Lai
was granted a Canadian work permit in February, and he reportedly works in real estate.

Mr.~Matas said Mr.~Lai was still awaiting a decision by Canadian immigration officials on whether to
evaluate the risk he might face if sent back to China.

In 2001, Jiang Zemin, the Chinese president at the time, sent Prime Minister Jean Chr\'etien a note
with assurances that Mr.~Lai would not be executed. But Mr.~Matas dismissed those assurances.
``There's still no system in place to monitor assurances,'' Mr.~Matas said. ``They won't even let
the Red Cross into the prisons.''

China has repeatedly said that the Lai case is the country's biggest scandal. Exhibitions and
television shows have detailed his alleged crimes and excesses. The authorities say that he plied
officials with prostitutes and alcohol in a seven-story mansion built for that purpose.

Eight people connected to the case have already been executed in China. Several others have been
jailed for sending funds to help Mr.~Lai pay for his defense.

\section{China Charges 4 Rio Tinto Employees; Spying Allegations Are in Question }

\lettrine{C}{hina} formally charged an Australian citizen and three other employees of the
British-Australian mining giant Rio Tinto late Tuesday with commercial bribery and trade secrets
infringement in a case that has rocked the global steel industry and created a diplomatic uproar.

But a brief statement on the case, reported by the official Xinhua news agency, did not mention the
much more serious charges of espionage that Beijing had earlier alleged and that had frayed
relations between China and Australia in recent weeks. It was unclear if the Chinese government
included those charges in its more detailed filing on the case or had backed away from them under
international pressure.

The case has generated significant anxiety among foreign corporations and investors worried that
China could stifle business by pressing state secrets charges against executives.

In the statement released through Xinhua, prosecutors said the four Rio Tinto employees had used
``improper means'' to obtain commercial secrets from China's government-controlled iron and steel
industry. Although commercial bribery is a criminal charge, the sentence for anyone convicted of the
crime is significantly shorter than those for states secrets violations.

The announcement came more than a month after the four employees were detained in Shanghai and
accused of stealing state secrets from China's steel industry and harming the nation's economic
interests. It followed weeks of diplomatic pressure from the Australian government, which had
complained about a lack of transparency in the case.

After the detentions, Prime Minister Kevin Rudd of Australia had cautioned China to deal fairly,
openly and judiciously with the employees because the world would be watching how it handled the
case. Australian government officials did not reply early Wednesday to requests for comment on the
statement of the charges.

If China indeed chose not to charge the employees with violating the country's state secrets law, it
would drastically alter the case because legal scholars say that prosecutions under that law are not
transparent and are notoriously difficult to defend against.

Jerome A.~Cohen, a professor of law at New York University and a expert on China's legal system,
said the language used in the prosecutor's announcement was unclear but suggested that Beijing might
have chosen to downgrade the charges.

``It doesn't sound like they are doing this under the state secrets law,'' Mr.~Cohen said after
being read the Xinhua statement. ``They may be trying to lower the temperature in the world by
getting it out of the state secrets context.''

The detentions of the four employees, including three Chinese citizens, had stunned many
multinational companies doing business in the country because the government had invoked its murky
espionage law in a case that many legal experts said involved allegations that sounded like
commercial bribery rather than spying.

The accusations in the case have also caused turmoil\cite{turmoil} in China's steel industry, the
world's largest, and disrupted the country's already contentious annual iron-ore negotiations with
some of the world's biggest producers, including Rio Tinto.

A spokesman at Rio Tinto's headquarters in London said late Tuesday that the company was not yet
aware that the charges had been filed.

In recent weeks Rio Tinto has strongly denied any wrongdoing in the case and insisted that the
company had no evidence that its employees, including Stern Hu, a Chinese-born Australian citizen,
had ever engaged in bribery or wrongdoing.

Mr.~Hu, the general manager of Rio Tinto's Shanghai office and the company's top iron ore salesman
in China, is believed to be one of the highest ranking foreign executives ever charged with
commercial bribery and trade secrets infringement involving a state-controlled industry.

\section{China Scales Back Software Filter Plan }

\lettrine{C}{hinese} officials retreated on Thursday from a plan to install so-called
anti-pornography software on every computer sold here, saying instead that Internet cafes, schools
and other public places must use the program, but that individual consumers will be spared.

The industry and information technology minister, Li Yizhong, said the notion that the program,
called Green Dam/Youth Escort, would be required on every new computer was ``a misunderstanding''
spawned by poorly written regulations.

The ministry order, first issued last May 19, had stirred an outcry from Chinese Internet users and
foreign computer manufacturers alike, arguing that the software ran counter to China's proclaimed
goal of creating an information-based society.

The United States warned China that the installation requirement could be seen as a violation of
world trade regulations.

Although the government insists that the program is meant to shield children from online
pornography, its filter -- automatically updated by the government -- targeted many topics with
political overtones. Free-speech advocates said that the program was a government attempt to extend
its control of political opinions into people's living rooms.

The information ministry previously had suspended the Green Dam pre-installation mandate on June 30,
one day before it was to take effect, saying that computer makers needed more time to accommodate it
in their manufacturing.

The Thursday statement by Mr.~Li appeared to make that suspension permanent. Mr.~Li said the
government would neither require the program to come pre-installed on new computers or force
computer makers to include the program on a CD with optional software.

A few Asian computer manufacturers, led by China-based Lenovo and Taiwan's Acer, nevertheless
include the software on computers sold in China.

Although Mr.~Li's concession is a step backward for the Green Dam program, the software remains
mandatory in schools, Internet cafes and other sites used by scores of millions of people. The
government already takes extraordinary steps to monitor computer use in Internet cafes, which remain
common in a nation where owning a computer remains a comparative luxury.

China has sought to increase government control over ordinary people's use of computers in recent
months. The government has systematically blocked ordinary citizens from viewing foreign-based
websites like Facebook, Flickr and YouTube that sometimes include comment critical of the
government.

Domestic websites with political content also have increasingly been censored or blocked. Experts
are divided over whether the increased censorship is a temporary measure in a year filled with
sensitive events, including the coming 60th anniversary of modern China's founding, or is a
permanent attempt to clamp down on unapproved speech.

The government recently proposed a requirement that all users of online chat rooms and bulletin
boards use their real names when posting comments, a move that would stifle the
sometimes-freewheeling debate on many sites. Until now, government censors have played a
cat-and-mouse game with anonymous Internet users who posted comments that flout\cite{flout} approved
positions.

\section{New Army Handbook Teaches Afghanistan Lessons }

\lettrine{M}{ore} than a year has passed since an Afghan police commander turned on coalition forces
and helped insurgents carry out a surprise attack that killed nine Americans, wounded more than 30
United States and Afghan troops and nearly resulted in the loss of an allied outpost in one of the
deadliest engagements of the war.

Within days of the attack, Army historians and tactical analysts arrived in eastern Afghanistan to
review the debacle near Wanat, interviewing soldiers who survived the intense battle, in which
outnumbered Americans exchanged gunfire for more than four hours with insurgents\cite{insurgent},
often at distances closer than 50 feet.

Now, that effort to harvest lessons from the firefight of July 13, 2008, has contributed to a new
battlefield manual that will be delivered over coming days to Army units joining the fight in
Afghanistan with the troop increase ordered by President Obama.

The handbook, ``Small-Unit Operations in Afghanistan,'' strikes a tone of respect for the Taliban
and other insurgent groups, which are acknowledged to be extremely experienced fighters; even more,
American soldiers are warned that the insurgents rapidly adapt to shifts in tactics.

In page after page, the handbook draws on lessons from Wanat and other missions, some successful and
some that resulted in death and injury for American and allied forces. The manual can be read as an
effort to push the nuances\cite{nuance} of the complex counterinsurgency fight now under way in
Afghanistan down from the generals and colonels to newly minted privates as well as to the sergeants
and junior officers who lead small units into combat.

Copies of the 123-page handbook, produced by the Center for Army Lessons Learned, are being
distributed throughout the service and are available to NATO allies and other nations with troops in
Afghanistan. A copy was provided in advance to The New York Times by an official involved in the
distribution, who said consideration was being given to a broader public release.

The manual includes a chapter titled ``Cultural Engagements,'' offering guidance to small-unit
leaders on building relationships with wavering village elders and trust among distrustful village
residents -- a process that cannot be left to senior officers who may be back at headquarters.

Implicit in the instructions is a warning that troops are at risk if they are aloof\cite{aloof} from
the locals and uncaring of their needs -- and of the certain dangers if intelligence sources are
used incorrectly.

One lesson of Wanat was that American troops, who had set up the firebase five days before the
attack, were caught unaware of collusion between the district police chief and the Taliban.

The manual describes how to train better for the defense of remote forward operating bases in harsh
Afghan terrain, especially in contested areas where the loyalties of local people are uncertain. The
detailed ``how to'' lists include instructions on such battlefield techniques as deploying
mortars\cite{mortar} more effectively than soldiers did at Wanat, where they did not take into
account terrain that provided cover for attackers.

In the fight now under way in Afghanistan, even small platoons may be expected to patrol areas and
conduct both combat operations and civilian reconstruction missions traditionally assigned to much
larger combat units.

``Every soldier or leader involved in command post operations is one less soldier or leader
available to send on patrol, provide security, or staff a quick-reaction force,'' the handbook says.
One lesson of Wanat was that the primitive forward firebase was understaffed.

The handbook's publication days after the first anniversary of the Wanat battle was first noted by a
blog operated by the Combined Arms Center, the Army's headquarters for advanced learning and leader
development at Fort Leavenworth, Kan.

The blogger, Frontier 6, said the handbook drew together lessons of successes in Afghanistan, as
well as what has been gleaned from operations when Americans left the battlefield badly bloodied.

``Although the losses at Wanat were tragic, a close scrutiny\cite{scrutiny} of the action with an
eye to lessons learned can save lives in the future,'' the blogger wrote, noting that the handbook
also built on analysis from an insurgent ambush of American troops this past April in the Korangal
Valley, also in eastern Afghanistan.

It is perhaps one of the worst-kept secrets in the Army, but Frontier 6 is the Internet alias of
Lt.~Gen. William B.~Caldwell IV, the Fort Leavenworth commander and occasional blogger, who
previously served as the top military spokesman in Iraq.

Combat commanders acknowledge how much they rely on the analysis and lessons-learned manuals sent
from headquarters back in the United States. ``The education of our force is the best weapon we
have,'' said Gen. Stanley A.~McChrystal, the new commander in Afghanistan. ``Counterinsurgency is
complex, nuanced and ever-changing, and success is dependent on a fighting force that can recognize
these changes and adapt to them.''

In distilling\cite{distil} lessons into practical advice for the troops, General Caldwell is
building on an effort brought to popular attention by a predecessor, Gen. David H.~Petraeus, now
commander of American forces in the Middle East.

Under General Petraeus's leadership at Fort Leavenworth, the military released a counterinsurgency
field manual credited with helping turn around the war in Iraq and ending the armed services' focus
on heavily armored conventional warfare.

With that manual's release, the American military was forced to embrace the messy irregular warfare
that had been the core competency solely of a small specialty branch in the armed services -- the
Army Special Forces, known as Green Berets.

\section{W.T.O. Rules Against China's Limits on Imports }

\lettrine{T}{he} World Trade Organization gave the United States a victory on Wednesday in its trade
battle with China, ruling that Beijing had violated international rules by limiting imports of
books, songs and movies.

The W.T.O. panel decision in Geneva buttresses growing complaints from the United States and Europe
that China is becoming increasingly nationalistic in its trade policies. It also offers some hope
that China will remove its restrictions on media and reduce rampant piracy\cite{piracy} of
intellectual property, though the country can appeal.

But even if China changes its policy in light of the decision, Western companies could struggle to
increase their sales anytime soon. The ruling does not affect a quota that caps at 20 the number of
foreign films that can be released in Chinese movie theaters each year.

Also, because of piracy, Chinese consumers are so accustomed to paying very little for DVDs, or
downloading movies or songs free on the Internet, that American movie companies already sell
authorized DVDs of their movies for much less in China than in the United States -- and still
struggle to find buyers.

Still, Ron Kirk, the United States trade representative, praised the panel's legal finding. ``This
decision promises to level the playing field for American companies working to distribute
high-quality entertainment products in China,'' Mr.~Kirk said, ``so that legitimate American
products can get to market and beat out the pirates.''

For the American media industry, the ruling essentially means that the W.T.O. supports demands by
United States movie studios, book and newspaper publishers, and record labels that they be allowed
to sell more directly to the Chinese consumer, rather than first going through a middleman, often a
state-owned enterprise, as China has required. It does not necessarily mean the Chinese consumer
will have access to a broader array of American films, books and music -- although those industries
hope that may eventually occur.

``American companies now have the right to trade without going through a Chinese intermediary at the
border,'' said James Bacchus, a lawyer at Greenberg Traurig in Washington who represented the China
Copyright Alliance, a consortium of media companies, in the case.

Dan Glickman, chairman of the Motion Picture Association of America, acknowledged that the decision
might not immediately result in a wider array of American movies available to the Chinese consumer.

``I wouldn't say it will have a night-and-day, revolutionary impact right away,'' he said. However,
he added, ``It's hard for me to believe that the import quota, which has been in effect for 10
years, will be there in perpetuity with this decision.''

Either side may appeal the panel's ruling. It is difficult, although not impossible, for a panel
decision like this one to clear the way for the petitioning country to impose trade sanctions on the
country that broke the rules.

The ruling goes to the heart of one of the biggest trade issues pending between China and the West:
whether intellectual property, like copyrighted songs, books and movies, should be granted the same
kind of protection from discriminatory trade practices as manufactured goods.

China has enjoyed double-digit economic growth through most of the last three decades in part
because of rapid expansion of exports, virtually all of which have been manufactured goods. But
Chinese imports have grown much more slowly, particularly if imports of goods for export are
excluded, like computer chips from Japan that are assembled in China into consumer electronics for
shipment to the United States.

One reason for the slow growth in imports has been China's restrictions on imported books, movies
and other content. Demand is met by pirated copies made in China; the latest Hollywood movies are on
DVDs on street corners across China within days of their release, at a cost of \$1 or less -- much
less in inland cities and for the buyer who bargains aggressively.

The Chinese government had no immediate reaction to the decision, which was released late at night
Beijing time. Chinese state media also initially ignored the decision. Officials sometimes wait a
day or two to respond to adverse trade developments.

In its petition to the W.T.O., the United States criticized China's requirement that most
copyrighted material be imported through a few government-designated companies, which tend to be
wholly owned or majority-owned by the government. The restrictions also required foreign financial
news services to operate through a government-designated distributor.

The panel condemned this, saying in its report that ``it also appears that foreign individuals and
enterprises, including those not invested or registered in China, are accorded treatment less
favorable than that accorded to enterprises in China with respect to the right to trade.''

Many of these restrictions, like limiting the number of foreign movies that can be shown each year
in Chinese theaters, have been aimed partly at limiting foreign influence in China but also at
sheltering domestic industries.

The panel stopped short of endorsing an American requests for a ruling on whether Chinese censorship
had unfairly restricted imports. The panel said that this question was outside its purview; for the
same reason, the panel also declined to rule on whether China's approval processes were too onerous
for would-be distributors of imported entertainment.

Like the United Nations, the W.T.O. has limited power to enforce decisions. But criticism from the
W.T.O. can shame countries, and panel rulings against other countries have frequently become the
basis for bilateral or multilateral negotiations that result in policy changes.

The Bush administration filed the original complaint in 2007, partly to head off possible
legislation requiring a more confrontational trade policy toward China. The Obama administration now
faces pressure from the Democratic majority in Congress to take more assertive action in response to
China's trade surplus during the current recession, and could use the ruling as evidence that the
issue is already being addressed. It may also use the victory as a precedent to take more cases
against China to the trade organization, said Gary Clyde Hufbauer, a trade expert at the Peterson
Institute for International Economics.

But while China has lost two other W.T.O. panel rulings in the last 13 months, regarding high taxes
on imported auto parts and lax enforcement of counterfeiting laws, China has not changed its
policies in either case.

``They've got a poor record of compliance. They keep filing appeals,'' said Lyle Vander Schaaf, a
partner in Washington at the law firm Bryan Cave who specializes in W.T.O. dispute panels and has
not advised either side in any of the three panel decisions against China.

\section{In Beijing, a Brand-New Old Quarter }

\lettrine{T}{his} is a city conspicuously\cite{conspicuous} unfriendly to pedestrians. Six-lane ring
highways cut through old neighborhoods, making places that are rather close as the crow flies very
far apart if you walk, since you've often got to navigate long cloverleaf intersections and widely
spaced overpasses simply to cross the street.

And that's one of the things, in a city far more oriented toward the car than shoe leather, that
makes the recent opening of a big pedestrian-only network of streets in the historic heart of
Beijing so noteworthy. It's in the area called Dashalar, which, during the Manchu-ruled Qing
Dynasty, was the ethnic Chinese part of this city, the warren of streets and avenues that lay just
outside the imposing city wall.

Now, the only parts left of the wall -- torn down at the order of Mao Zedong despite the entreaties
of China's leading architectural historian to preserve it -- are two immense gates. The more famous
of them is known as the Qianmen, for Front Gate, and beyond it, running in an arrow-straight line
due south, is Dashalar's main street, Qianmen Dajie, or Front Gate Big Street, where cars are now
banned.

Over the last decade or so, most of the old Dashalar neighborhood has been razed. Tens of thousands
of residents, whose single-story courtyard homes were often models of
dilapidation\cite{dilapidation} and obsolescence\cite{obsolescence}, were required, with payment of
various amounts of compensation, to move out. This process is vividly recounted in a book published
a year ago, ``The Last of Old Beijing'' by Michael Meyer.

But while most of the old residential areas of Beijing that have been demolished were transformed
into zones of high-rise hotels, shopping centers and office buildings, Dashalar has been converted,
perhaps ironically, back into what it once was -- or, more accurately, into an idealized, postcard
version of what it once was, a shopping street lined by three-story traditional Chinese buildings,
with balconies above the first floor, latticed balustrades, red columns, bright gold-leaf Chinese
signs against backgrounds of black wood.

``Western preservationists won't like this,'' said Chen Xiangming, a former Beijing resident who is
professor of sociology and international studies at Trinity College in Hartford, Connecticut, and
the editor of a new book, ``Shanghai Rising.'' We were on a recent walk through the brand-new
pedestrian mall.

They won't like it because creating a new copy of an old neighborhood fails to meet the standard of
what is often called authenticity\cite{authenticity}, which requires the preservation of the actual
antique buildings. The opening of the Dashalar pedestrian district in this sense represents the
completion of a strange sort of circle.

In the past few years, even as a lot of Beijing's old neighborhoods have been demolished, there's
been the realization that Old Beijing had value, if mostly for the city's growing tourism industry.
And so, a couple of weeks ago, after two years of construction that took place behind concealing
high fences themselves decorated with idealized images of Old Beijing, Qianmen Big Street formally
opened to pedestrians and shoppers.

``It's a sort of Disneyfication of Beijing,'' Mr.~Chen allowed, ``and Disneyfication is the opposite
of authenticity.'' But Mr.~Chen is less critical of this in the Chinese context than some others
might be. In part this is because he remembers, having grown up in Beijing in the 1960s, how poor
and drab the city was, and he appreciates the better and richer material life it has now.

``I had the Communist experience,'' he said, ``which some of the critics of Beijing's policy might
not have had, and that can color their views.''

But there's another reason not to judge the theme-park aspect of Beijing too harshly, and that's
because it's part of a global phenomenon.

``The Disneyfication of cities is not happening only in China,'' Professor Chen said. ``Many
deindustrialized cities are doing it to revitalize themselves through tourism.

``This is local to China, but it's also Western because it's been Disneyfied,'' he said. ``The more
nuanced counter-argument to the anti-Disneyfication argument is that you have to Disneyfy with a
different twist, or else you don't have anything special to attract tourists.''

What might clash with this particular restoration is the presence in these replicated antique
Chinese shops of the same global brands you find in the other shopping centers in this city. There's
a Swatch store about to open under a latticed terrace, the sign pointing down the street to an H\&M
(``Fashion heaven is straight ahead,'' the sign says, in English.) Other foreign brands are
reportedly on their way.

One of the characteristics of China is that, contrary to repeated entreaties by the propaganda
machinery not to ``worship foreign things,'' most of the prestige\cite{prestige} brands are foreign
imports (though they are ``imports'' frequently manufactured in China under foreign license).

Still, if the worship of foreign things is evidenced in Dashalar, there do seem to be more Chinese
traditional products there than you would find in other shopping arcades, including some that have
been in their locations in the old Chinese city for hundreds of years.

There's Duyiju, for example, making and selling dumplings since 1738, and several venerable silk
shops, including Ruifuxiang, which, according to a sign in Chinese and English, has been on its
current site ``Since the eighth year of the reign of Emperor Xienfeng in the Qing Dynasty'' -- in
other words, since 1862.

There are slipper shops, clothing shops and Chinese medicine shops. Mr.~Chen and I had a fine lunch
at Quan Ju De, perhaps Beijing's most famous roast duck restaurant, though we ate it off plastic
plates made to look like Qing Dynasty tableware, and that seemed an especially ersatz touch. And
then there's Liubiju, Beijing's most famous preserved-foods store, where porcelain vats contain
dozens of preserved vegetables, sauces, condiments and pastes, a few of them vaguely recognizable
for Westerners, most of them not.

Given recent history, Beijing's replica of itself certainly seems better than yet another grouping
of high-rises, though the question certainly needs to be raised whether anything that could have
been put in old Dashalar could be as fascinating to visitors as the neighborhood that was destroyed.

\section{Chinese Artist Says He Was Barred From Rights Advocate's Trial }

\lettrine{A}{ prominent} Chinese artist and frequent critic of the Communist Party said he was hit
by police officers and put under detention in his hotel room in western China on Wednesday when he
tried to go testify at the trial of a civil rights advocate.

The artist, Ai Weiwei, best known for helping to design the Bird's Nest Olympic Stadium in Beijing,
said dozens of police officers barged into his hotel room early Wednesday and the rooms of others
who had traveled to Chengdu, the capital of Sichuan Province, to testify on behalf of Tan Zuoren,
the rights advocate.

The opening session of Mr.~Tan's trial began Wednesday morning and ended hours later without a
verdict. The courtroom was closed to the public. Mr.~Tan, a well-known writer, has been charged with
subversion. He is believed to be on trial because of his role in pushing for an official
investigation into widespread school collapses during the Sichuan earthquake last year, and for
trying to organize a group event in June to commemorate the 20th anniversary of the massacre of
civilians by government forces during the Tiananmen Square protests of 1989.

Earlier this month, another rights advocate, Huang Qi, went on trial on a charge of possessing state
secrets. Like Mr.~Tan, he pushed for the government to investigate the school collapses.

During the earthquake last year, thousands of students were killed when school buildings in Sichuan
and other provinces crumbled, even as surrounding buildings remained standing. Grieving parents said
shoddy construction and corruption were responsible, and demanded that officials investigate. Local
governments went to great efforts to silence the parents, ordering the police to detain them, or
handing out cash payments in exchange for the parents' dropping their complaints.

Initial reports from the official news media said about 7,000 schoolrooms collapsed and as many as
10,000 children might have died. In May, the government released the first official toll of students
killed, saying 5,335 were dead or missing.

Earlier this year, Mr.~Ai sent volunteers to Sichuan to collect the names of students who had been
killed. He began posting the names on his blog and kept a running tally. Government censors then
blocked his blog, Mr.~Ai said, while police officers in Sichuan detained some of his volunteers and
beat a few of them.

Mr.~Ai said in a telephone interview on Wednesday that he had arrived in Chengdu the previous day to
testify on behalf of Mr.~Tan but that the judge did not allow him to. He said that he had planned to
appear at court anyway but that he and 10 or 11 of his volunteers were prevented from doing so by
the officers who barged into their hotel rooms and kept them under watch.

``They left a couple of hours later, but some stayed in the hallway and some in the lobby to keep an
eye on us, to make sure we failed to attend the trial,'' Mr.~Ai said from the hotel.

Mr.~Ai posted grainy digital photographs on Twitter of police officers in the hotel hallway.

Someone answering the phone at the police headquarters in Chengdu declined to comment.

The court did not allow anyone to testify on behalf of Mr.~Tan at the trial on Wednesday, said Pu
Zhiqiang, Mr.~Tan's lawyer.

He said he would submit a written defense statement to the court by Monday.

``However, unless this is an extremely rare case, Tan will be found guilty,'' he said.

\section{Sony Plans to Adopt Common Format for E-Books }

\lettrine{P}{aper} books may be low tech, but no one will tell you how and where you can read them.

For many people, the problem with electronic books is that they come loaded with just those kinds of
restrictions. Digital books bought today from Amazon.com, for example, can be read only on Amazon's
Kindle device or its iPhone software.

Some restrictions on the use of e-books are likely to remain a fact of life. But some publishers and
consumer electronics makers are aiming to give e-book buyers more flexibility by rallying around a
single technology standard for the books. That would also help them counter Amazon, which has taken
an early lead in the nascent market.

On Thursday, Sony Electronics, which sells e-book devices under the Reader brand, plans to announce
that by the end of the year it will sell digital books only in the ePub format, an open standard
created by a group including publishers like Random House and HarperCollins.

Sony will also scrap its proprietary anticopying software in favor of technology from the software
maker Adobe that restricts how often e-books can be shared or copied.

After the change, books bought from Sony's online store will be readable not just on its own device
but on the growing constellation of other readers that support ePub. Those include the Plastic Logic
eReader, a thin device that has been in development for nearly a decade and is expected to go on
sale early next year.

``There is going to be a proliferation of different reading devices, with different features and
capabilities and prices for a different set of consumer requirements,'' said Steve Haber, president
of Sony's digital reading unit. ``If people are going to this e-book shopping mall, they are going
to want to shop at all the stores, and not just be required to shop at one store.''

Sony's move comes amid mounting concern about Amazon's market power in the budding category of
electronic books. E-book sales in the United States hit a record \$14 million in June, a 136.2
percent increase from a year earlier, according to the Association of American Publishers.

Amazon does not divulge its e-book revenue, but analysts say it most likely accounted for a majority
of those sales.

Amazon inadvertently demonstrated one potential consequence of e-book restrictions last month when
it discovered that it had sold unauthorized copies of George Orwell's ``1984'' and ``Animal Farm''
and then removed the books from the Kindle libraries of people who had bought the novels.

Amazon's chief executive, Jeffrey P.~Bezos, later apologized for the move, but not before advocates
used the episode to rail against limitations on digital reading.

``People need to remember, when they buy books that come with digital rights management, they don't
have the freedoms they normally would have with a book,'' said Holmes Wilson, campaigns manger of
the Free Software Foundation, which obtained the signatures of nearly 4,000 authors and tech pundits
on a petition saying Amazon's anticopying software was a ``clear threat to the free exchange of
ideas.''

Companies like Sony and Adobe do not want to abandon anticopying measures, fearing that piracy of
books would run rampant. Rather, they want to push the e-book industry toward common standards to
avoid a replay of Apple's domination of the digital music business.

Early this decade, Apple sold music from its iTunes store that was protected by its own FairPlay
software and could be played only on the iPod.

The result was what is known as ``lock-in.'' Apple built up extraordinary market power and leverage
to dictate terms to the major music labels on matters like the price of digital songs. Then, as now,
second-tier players banded together to promote the increased flexibility and choice that open
standards gave to consumers.

``If the business terms and conditions end up being dictated to publishers by one bookseller who has
a chokehold over the value chain, publishers are going to have a hard time staying profitable,''
said Bill McCoy, general manager for Adobe's digital publishing business.

For Sony, which introduced its Reader devices more than a year before the Kindle arrived, the move
to open formats is part of a strategy to make up lost ground. Sony recently introduced two new, less
expensive devices and announced it was dropping its price for new releases and best sellers to
\$9.99. Later in the year, the company will begin selling a third Reader that will, like the Kindle,
allow users to buy e-books wirelessly.

Amazon, for its part, believes it can go it alone, without embracing industry standards. An Amazon
spokesman would not comment for this article, but Mr.~Bezos has said before that his goal was to
``make Kindle books available on as many hardware devices as possible.'' That suggests it will soon
introduce versions of its Kindle software for the Palm Pre and other reading devices.

Allen Weiner, an analyst at the technology research firm Gartner, says there is one more company
that must declare its allegiance to either an open or closed world for e-books: Apple.

If, as expected, Apple soon introduces a tablet computer that can function as a reading device, and
if it embraces an open standard like ePub, Amazon will have to reconsider its closed approach,
Mr.~Weiner said.

``If you see some Adobe executive up on stage with Steve Jobs when they announce the tablet, at that
point Amazon has a lot to worry about,'' he said.

\section{Obama Is Taking an Active Role in Talks on Health Care Plan }

\lettrine{I}{n} pursuing his proposed overhaul of the health care system, President Obama has
consistently presented himself as aloof from the legislative fray, merely offering broad principles.
Prominent among them is the creation of a strong, government-run insurance plan to compete with
private insurers and press for lower costs.

Behind the scenes, however, Mr.~Obama and his advisers have been quite active, sometimes negotiating
deals with a degree of cold-eyed political realism potentially at odds with the president's
rhetoric.

Early last month, for example, hospital officials were poised to appear at the White House to
announce a deal limiting their industry's share of the costs of the overhaul proposal when a wave of
jitters swept through the group. Senator Max Baucus, the Finance Committee chairman and a party to
the deal, had abruptly pulled out of the event. Was he backing away from his end of the deal?

Not to worry, Jim Messina, the deputy White House chief of staff, told the hospital lobbyists,
according to White House officials and lobbyists briefed on the call. The White House was standing
behind the deal, Mr.~Messina told them, capping the industry's costs at a maximum of \$155 billion
over 10 years in exchange for its political support.

Some Democrats and industry lobbyists now argue that, in negotiating deals through Mr.~Baucus's
committee with powerful health care interests, the White House was tacitly signaling as early as
last spring that it might end up accepting something more modest than the government insurer the
president has said he prefers.

The Finance Committee, for example, appears to be coalescing around the idea of nonprofit insurance
cooperatives instead of a government-run plan. It is a proposal the health care industry prefers,
but many liberal Democrats oppose, in both cases because cooperatives are likely to have less
leverage over health care prices.

Rahm Emanuel, the White House chief of staff, disputed that the administration had elevated the work
of the Senate finance panel above the four other committees that have all approved strong government
insurers.

``They are an important committee,'' Mr.~Emanuel said. ``They have a bipartisan process. The
president would like that to work, just as he is proud that the other committees have done their
work. They don't get an exalted status over everybody else.''

But he also acknowledged the political realities that have made the Finance Committee's
still-unfinished cooperative plan a center of attention.

``We have heard from both chambers that the House sees a public plan as essential for the final
product, and the Senate believes it cannot pass it as constructed and a co-op is what they can do,''
Mr.~Emanuel said. ``We are cognizant of that fact.''

Asked whether the president would accept the weaker co-op, Mr.~Emanuel declined to comment. ``I am
not going to fast-forward the process,'' he said.

Industry lobbyists and moderate Democrats in both chambers, though, argue that the White House's
actions behind the scenes show a recognition that the finance panel's anticipated compromise is the
most likely template for any final legislation.

``The House has largely been a sideshow,'' said Representative Jim Cooper of Tennessee, a member of
the so-called Blue Dog caucus of conservative Democrats. ``The Senate Finance Committee is where it
really matters. That's the bottleneck.''

Members and staff of the four other committees working on health care say the White House has
generally stayed on the sidelines of their work. ``They have been -- what is a good way to put it?
-- available for consultation,'' Mr.~Cooper said. ``They haven't been hands-off, but they haven't
been hands-on.''

Mr.~Obama and his top aides have immersed themselves in the Senate Finance Committee process. The
president talks to Mr.~Baucus several times a week, people briefed on their conversations say.
Mr.~Obama has also held a few calls with the panel's ranking Republican, Senator Charles E.~Grassley
of Iowa.

In addition, Mr.~Obama invited both senators to a private lunch at the White House early in the
summer and met with six panel members for another White House session last week. White House
advisers have held long evening and weekend meetings with Finance Committee staff members, even
poring over copies of the Tax Code together.

Nancy-Ann DeParle, charged with leading the White House health effort, has a standing biweekly
meeting with Mr.~Baucus, while Peter R.~Orszag, the White House budget director, has spent so much
time in the senator's office that he helps himself to the Coke Zeros tucked away in Mr.~Baucus's
personal refrigerator.

Lobbyists for both the drug and hospital industries say that, as early as June, White House
officials directed them to work out cost-saving deals with Mr.~Baucus's committee.

Drug industry lobbyists said they negotiated a deal to contribute \$80 billion over 10 years toward
the cost of an overhaul with Mr.~Baucus, under White House supervision, before taking it to the
president for final approval. House lawmakers have said they were caught by surprise when it was
announced.

Hospital industry lobbyists, speaking on condition of anonymity for fear of alienating the White
House, say they negotiated their \$155 billion in concessions with Mr.~Baucus and the administration
in tandem. House staff members were present, including for at least one White House meeting, but
their role was peripheral, the lobbyists said.

Several hospital lobbyists involved in the White House deals said it was understood as a condition
of their support that the final legislation would not include a government-run health plan paying
Medicare rates -- generally 80 percent of private sector rates -- or controlled by the secretary of
health and human services.

``We have an agreement with the White House that I'm very confident will be seen all the way through
conference,'' one of the industry lobbyists, Chip Kahn, director of the Federation of American
Hospitals, told a Capitol Hill newsletter.

Mr.~Emanuel and liberal Democrats argued that the White House had worked more closely with the
Senate Finance Committee because it was stepping in to break up legislative logjams. In the same
way, they said, Mr.~Obama and Mr.~Emanuel had personally interceded to resolve a last-minute revolt
by conservative House Democrats that threatened to derail a bill in the House Energy and Commerce
Committee at the end of July.

Representative Henry A.~Waxman, the California Democrat who is chairman of the Energy and Commerce
Committee, said the president had personally assured House members that he did not intend to let the
Senate Finance Committee determine the final legislation.

``This is going to be a genuine conference with give and take,'' Mr.~Waxman said. ``The president
has said that personally to the senators, and he has said it personally to us.''

He added: ``The president has said he wants a public option to keep everybody honest. He hasn't said
he wants a co-op as a public option.''

Still, industry lobbyists say they are not worried. ``We trust the White House,'' Mr.~Kahn said.
``We are confident that the Senate Finance Committee will produce a bill we fully can endorse.'' 

\section{After Protest, Chinese Officials Halt Steel Mill Sale }

\lettrine{A}{ Chinese} provincial government on Sunday halted the privatization of a state-owned
steel mill where thousands of workers protested last week and took an official hostage, in the
latest sign of increasing labor activism in the country's steel industry.

The apparent capitulation\cite{capitulation} of the government of Henan province in central China
came three weeks after rioting workers beat to death the executive overseeing the sale to a private
business of another state-owned steel company, Tonghua Iron \& Steel Works, in the northeastern
province of Jilin.

Local, provincial and national government agencies have been reluctant to use overwhelming force
against protesting workers. China Daily newspaper reported Saturday that police had tried to break
through the ranks of workers Friday in the latest incident, at Linzhou Iron \& Steel in Anyang City,
in Henan.

China Daily did not say if the police had been successful. The official Xinhua news agency said that
the workers decided Saturday to halt their protests, which had attracted up to 3,000 participants at
a time, after a government mediation team agreed to reconsider the takeover. Xinhua did not mention
what became of the official who had been held hostage.

The success of steel workers in blocking privatization could embolden\cite{embolden} workers in
other industries, experts on Chinese labor issues said Sunday.

``It is no longer possible to push through privatization regardless, without considering the
workers' interests,'' said Geoffrey Crothall, a spokesman for the China Labor Bulletin, a labor
rights advocacy group based in Hong Kong.

In a sign of high-level interest in the recent unrest, the government-sponsored All China Federation
of Trade Unions has posted a prominent series of commentaries at the top of its Web page under the
heading, ``Corporate restructuring: participation of the trade union is essential.''

Chinese law has long required that each privatization be approved by the workers' congress of the
affected company. But local government officials and company managers have frequently been able to
rig\cite{rig} the approval by running the congresses themselves, Mr.~Crothall said.

In this weekend's action, Xinhua said that the Henan provincial government and the province's
Communist Party committee issued a decision that ``issues regarding the future of Linzhou Iron \&
Steel Co.~Ltd.~and benefits of its workers should be decided by its workers' congress.''

The Chinese steel industry, the world's largest and a cornerstone of the country's
construction-dependent economy, is in turmoil\cite{turmoil} this year, which may have fed labor
unrest.

The global economic downturn has severely hurt the sector, with Chinese steel exports down 15.4
percent in the first half of this year. In addition, the Chinese government is locked in a series of
disputes with Australia over iron ore imports, and has arrested four Rio Tinto executives accused of
bribery and trade secret infringement during price negotiations with Chinese steel makers. Rio Tinto
has strongly denied any wrongdoing.

Faced with a glut\cite{glut} of steelmaking capacity and many small steel companies vying to buy
iron ore, Beijing officials on Thursday ordered a three-year moratorium\cite{moratorium} on the
construction of any new steel mills or the expansion of existing ones.

But the government has had less success in its efforts to force a consolidation of existing mills,
which might strengthen their bargaining power with the handful of multinationals that dominate the
global iron ore business. Local and provincial government agencies have been wary of losing control
of businesses that are often vital to their economies, and many workers are opposed to
consolidation.

Yet many older, less efficient steel mills are deeply troubled. Founded in 1969, the Linzhou mill
has not been operating since March because its sales fell, it ran short of cash and it failed to
meet environmental standards.

Xinhua said that Fengbao Iron \& Steel had already paid 180 million yuan, or \$26.5 million, of the
259 million yuan it bid at an auction to acquire Linzhou Steel. The 180 million yuan will be
refunded, Xinhua said, citing an unidentified local official.

Mary E.~Gallagher, the director of the Center for Chinese Studies at the University of Michigan in
Ann Arbor, predicted that China would be able to continue privatizing state-owned enterprises in the
years ahead, but that it would be costlier than during the big shift toward private ownership from
1997 to 2001. At least 30 million workers at state-owned enterprises were laid off then, but local
protests did not spread.

``I'm betting they can do it again, but that it will cost them with higher compensation packages to
laid-off workers,'' she said.

\section{American Held in Myanmar Is Released }

\lettrine{A}{n} American man who was sentenced in Myanmar last week to seven years of hard labor for
intruding at the home of Daw Aung San Suu Kyi, the pro-democracy leader, was released Sunday and
left the country together with a visiting American senator, Jim Webb.

The release of the American came one day after Mr.~Webb, a Democrat from Virginia, met with the
leader of Myanmar's ruling junta, Senior Gen. Than Shwe, and with Mrs.~Aung San Suu Kyi, who was
sentenced Tuesday to 18 months of house arrest.

At a news conference in Bangkok on Sunday, Mr.~Webb said the meetings and the release of the
American, John Yettaw, were gestures that could be helpful as part of the foundations for a
relationship ``of good will and confidence building so that we can have a better situation in the
future.'' Although he said he was not an emissary\cite{emissary} of the administration, Mr.~Webb's
visit came at a time when the United States was exploring the possibility of a more cooperative
relationship with the military junta. Mr.~Webb said he would discuss his trip with State Department
officials when he returned to Washington.

For decades, Washington's policy toward the junta has been confrontational and
punitive\cite{punitive}, with an ever more elaborate package of economic and diplomatic sanctions in
reaction to the generals' human rights abuses and political repression.

Secretary of State Hillary Rodham Clinton has said that the administration is exploring options for
a more fruitful\cite{fruitful} relationship, saying neither sanctions by Western nations nor
engagement by Myanmar's neighbors has affected the junta's behavior.

In an indication that diplomatic contacts have already been under way between the two governments,
Mr.~Webb said Mr.~Yettaw's arrest had ``set back'' the process toward more cooperative relations.
Mr.~Yettaw, 53, appeared tired but smiled as he and the senator arrived at a Bangkok airport.
Mr.~Webb said Mr.~Yettaw ``is not a well man'' and was receiving medical treatment in Bangkok. He
had suffered seizures before Tuesday's verdict, and Mr.~Webb said he had had a ``medical incident''
on Sunday when he was released.

``I believe what happened was regrettable,'' Mr.~Webb said. ``He was trying to help. He's not a
mean-spirited human being.''

Mr.~Yettaw, of Falcon, Mo., had been sentenced to seven years of hard labor after swimming across a
lake and spending two days at the home of Mrs.~Aung San Suu Kyi in early May.

Ms.~Aung San Suu Kyi, 64, was convicted and sentenced at the same trial along with two housemates
for violating the terms of her house arrest in connection with Mr.~Yettaw's intrusion.

An initial sentence of three years at hard labor for all three women was commuted to house arrest by
Gen. Than Shwe in what many analysts saw as a reaction to international criticism.

At the news conference, Mr.~Webb emphasized the importance of China, Myanmar's powerful
patron\cite{patron} in the region, in resolving the country's problems. Mr.~Webb has been a longtime
advocate of greater engagement with Myanmar. It was not clear from his remarks at the news
conference whether Mrs.~Aung San Suu Kyi agreed with him.

``We had a very long discussion about when they work and when they don't work,'' the senator said.
He said Mrs.~Aung San Suu Kyi did not make a statement regarding the continuation or softening of
sanctions. ``She would prefer to wait to have a unified statement from her party at large,''
Mr.~Webb said.

\section{China Warms to New Credo\cite{credo}: Business First }

\lettrine{S}{o} far this week, the World Trade Organization has rebuffed China in an important case
involving Chinese restrictions on imported books and movies. The Chinese government dropped
explosive espionage\cite{espionage} charges against executives of a foreign mining giant, the
Anglo-Australian Rio Tinto, after a global corporate outcry. And on Thursday, the government said it
had backed off another contentious\cite{contentious} plan to install censorship software on all new
computers sold here.

Throughout its long economic boom, China has usually managed to separate its aggressive push into
the global business arena from domestic politics, which remained tightly controlled by the Communist
Party. But events this week raise the question of just how long it will be before the two meet.

In each of those matters, politics and business collided, and business won. Business does not always
win, and when it does, as in these cases, the reasons are as often as not a matter of guesswork. But
in at least some high-profile matters, China appears to be facing the reality that the outside
business world can be freewheeling and defiant\cite{defiant} when its profits are threatened. And so
China's authoritarian system may also have to evolve in ways its top leaders may not readily
endorse.

Beijing has a global footprint now, a consequence of its booming domestic growth and breakneck
international expansion. And decisions that once were made on purely parochial\cite{parochial}
grounds -- like censoring Web sites, protecting the interests of its state-owned companies and
restricting the flow of foreign news and entertainment into China -- now have international
ramifications\cite{ramification}.

``This is a country in the middle of a big transition in its global role,'' said Kenneth Lieberthal,
a veteran\cite{veteran} China analyst now at the Brookings Institution. ``They've always looked in
the past to what's good for China, and they still do. But for the first time, added to that is the
consideration that they're in the position of being rule-makers, not just rule-takers.''

China's leaders, he said, ``are just beginning to learn how to handle that.''

Consider the following: Since late May, Beijing's Industry and Information Technology Ministry had
more or less insisted that so-called anti-pornography software, called Green Dam-Youth Escort, would
eventually be packaged with every newly purchased computer.

On Thursday, the ministry backed down, calling the requirement a ``misunderstanding'' spawned by
badly written rules. Officials offered no other explanation, but the retreat followed weeks of
protests by outsiders -- from foreign computer makers to foreign governments to foreign corporate
branch offices -- that said the software stifled free speech, compromised corporate security and
threatened computers' stability.

Computers are not the only example.

This week, the World Trade Organization told Beijing that it could no longer force providers of
American books, music and films to distribute their goods through a local partner. Foreign companies
saw that rule as an impediment\cite{impediment} to reaching a broad Chinese audience with their
products. The Chinese market is flooded with pirated CDs and DVDs whose contents' creators receive
no money.

The Chinese legally may appeal the decision, but the foreign minister, Yang Jiechi, indicated in a
Geneva speech that simply ignoring it was not an option. China worked for years to join the global
trading system and is bound, as much as other nations are, by its rules.

``China will never seek to advance its interests at the expense of others,'' Mr.~Yang said,
according to Reuters.

Similarly, Chinese prosecutors appeared this week to retreat from earlier statements that they would
prosecute employees of Rio Tinto as spies for stealing state secrets.

While the espionage\cite{espionage} allegations were not spelled out, they were apparently related
to delicate commercial negotiations over the price of China's imports of iron ore for its steel
mills.

Rio Tinto executives have strongly denied the accusations, and both the United States and Australia
said China's actions could have both business and diplomatic repercussions.

While the Rio Tinto employees still face lesser charges of bribery and theft of trade secrets, the
espionage threats stirred broad unease among foreign companies operating in China, which feared that
they could face persecution and closed-door trials for engaging in what much of the world would
regard as bare-knuckle business tactics.

Yet whether such instances represent trends or exceptions -- or neither -- remains a matter of some
debate.

Increasingly, many experts say, Chinese officials appear to be aware that their actions have far
broader ramifications than they might have had even a few years ago.

``Fifteen years ago, the mantra\cite{mantra} in China was, 'We're the victims of a system that's
stacked against us,' '' said James V.~Feinerman, an expert on Chinese law and policy at Georgetown
University in Washington.

China's entry into the world trading system, he said, is slowly helping to change the nation's view
of itself from that of an outsider to an insider with a stake in the global system's success.

Other experts note, however, that what outsiders see as carefully calculated policy changes may in
fact be nothing of the sort.

The government's decision to install censorship software on computers -- and its subsequent reversal
-- is but one example, they say; the original proposal was probably pushed by a government clique
that found itself outflanked once Internet users and foreign corporations began objecting to the
plan.

``Is China susceptible to international pressure? Of course it is,'' said Charles Freeman, a leading
China scholar at the Center for Strategic and International Studies in Washington.

``China does have international interests, and they are impacted by what it does domestically,'' he
said. ``There's a constant battle between agencies over how much political capital to expend on
international issues against domestic interests.''

In any case, few experts are willing to stake their reputations on a prediction that Beijing's
recent softening of some positions signifies a strong trend.

To the contrary, Mr.~Feinerman said, China had undergone ``a real pushback'' in the last five years
on some fronts, reasserting political dogma in some areas where commercial norms and the rule of law
had begun to have more sway.

And Jonathan Hecht, an expert on Chinese law at Yale University's China Law Center, said that
developments in China should be viewed against a history of great leaps forward on such matters,
followed by equally great retreats.

``I've given up predicting long-term trends,'' he said. 

\section{Nortel Sale to Ericsson Stirs Protest in Canada }

\lettrine{A}{t} the start of this decade, Nortel Networks was perhaps the most hated company in
Canada. The collapse of technology stocks combined with an accounting scandal destroyed the widely
held telecommunications company's share price and, with that, the savings of thousands of Canadians.

But now that Nortel is being broken up through bankruptcy proceedings, many Canadians have found, if
not a renewed love, at least a feeling of nostalgia\cite{nostalgia} for the 114-year-old company.

The prospect of foreign companies' picking the bones of what was once Canada's largest corporation
has prompted the hasty reassessment of Nortel.

James L.~Balsillie and Michael Lazaridis, the two men who brought the world the BlackBerry, are
among the prominent Canadians who are calling on the government to review the proposed \$1.13
billion sale of Nortel's wireless infrastructure business to Ericsson, the Swedish company that
dominates that business. Last week, the industry committee of the House of Commons in Canada held an
unusual emergency meeting to discuss Nortel's fate.

Nortel is bankrupt, its shares are delisted, its former top executives remain under criminal
indictment, its businesses are rapidly shrinking and its chief executive and most of the board
resigned this week. But several Canadian politicians argue that the company may be too valuable to
fall into Swedish hands.

``We need to put a magnifying glass on this,'' Brian Masse, a member of the industry committee, said
from his constituency\cite{constituency} in Windsor, Ontario, on Tuesday.

The renewed interest in Nortel among Canadians has several roots. Mr.~Balsillie and Mr.~Lazaridis
appear to have taken up the mantle because, in their version of events, Nortel reneged\cite{renege}
on a deal to sell the wireless assets to their company, Research in Motion.

Other Canadians are still smarting from sale of several prominent Canadian companies to foreign
competitors before the credit crisis, including the mining companies Inco and Alcan. And there has
been considerable publicity surrounding the decision by United States Steel to drastically scale
back the operations of Stelco, a Canadian steelmaker it acquired in 2007. The Canadian government is
arguing in court that the American steel company violated commitments it made under Canadian
investment laws to maintain employment and production.

The debate over Nortel has focused on exactly what remains of value at the company and, more
broadly, what Nortel owes Canadians.

``People are a bit miffed\cite{miff} that a lot of government money went into Nortel through tax
credits'' for research and development, said Mihkel M.~Tombak, a professor of strategic management
at the University of Toronto.

While that was certainly the case in the past, the Nortel executives who appeared at the committee
hearing repeatedly pointed out that the company's lack of profits had left it unable to claim tax
credits since at least 2001. They were unable to discuss the government's contribution before then,
and a series of financial restatements related to the bookkeeping fraud made it difficult to assess
that from Nortel's published statements.

Under its deal, Ericsson will obtain the right to use the patents derived from that research
although, at least for now, Nortel will retain ownership. Ericsson will also acquire the company's
wireless business, including a research-and-development unit that is largely based in Ottawa.

Nortel's greatest strength is a wireless technology known as CDMA Used almost exclusively in
North America, CDMA is now, as Nortel's executives at the hearing acknowledged, a technology
that is gradually becoming obsolete.

In his testimony, Mr.~Lazaridis, who directs Research in Motion's research, said that Nortel's most
valuable property is patents related to the next generation of wireless networks.

``Nortel's intellectual property and work force, devoted to next-generation wireless research known
in the industry as L.T.E., or long-term evolution, is nothing short of a national treasure that
Canada must not lose,'' Mr.~Lazaridis told the members of Parliament, adding that putting it in
foreign hands could endanger Canada's national security and make it more difficult for Canadian
carriers to restore future network failures. The Canadian government can turn down foreign purchases
on national security grounds.

While Research in Motion runs a worldwide data network to support its BlackBerry e-mail service, it
has never been in the business of producing the software and hardware used by carriers. But
Mr.~Lazaridis said that his company now wanted a stake in the development of L.T.E. and believed at
several points it had struck a deal with Nortel to buy its technology for about \$1.1 billion. He
told the committee of shaking hands on deals and exchanging e-mail messages with Nortel over the
wording of a press release to announce them.

In the end, however, Nortel surprised Research in Motion by filing in Canada and the United States
for bankruptcy protection. ``We felt like we were snookered,'' Mr.~Lazaridis said.

Unlike Ericsson and Nokia Siemens Networks, Research in Motion did not enter the auction for
Nortel's wireless assets, which was supervised by an American court. Mr.~Lazaridis said that was
because of conditions attached to the process, including the licensing rather than the sale of the
patents.

Testifying after Mr.~Lazaridis, Paul B.~Schabas, a Canadian lawyer representing Ericsson, suggested
that Research in Motion should have raised its complaints earlier and in court.

``It could have come and objected or made submissions,'' Mr.~Schabas said. ``It's never done that.
It's only come here today to complain about the process.''

Several Canadian politicians, echoing Mr.~Lazaridis, have emphasized the value of Nortel's L.T.E.
patents. But unlike its competitors, Nortel has found no carriers that are interested in buying its
technology, and some analysts question the importance of what it has developed to date.

``It may be the crown jewels of Canada, but it's no longer a jewel of the wireless industry,'' said
John Strand, the chief executive of Strand Consult, a wireless technology consulting company based
in Denmark. ``It's not like Nortel has the keys to L.T.E. Nortel is a like a sick old man who has
been lying in a hospital bed for years.''

Mr.~Masse, who is a member of the labor-backed New Democratic Party, is unclear about what he wants
the Canadian government to do with Nortel beyond reviewing the purchase under the Investment Canada
Act.

During a visit to Panama on Tuesday, Prime Minister Stephen Harper told reporters that the sale
would be reviewed under the act to see if it was ``in Canada's national interest,'' but he ruled out
any other intervention. Mr.~Masse, however, suggested that rather than approve the sale of the
Nortel wireless unit, the government might consider keeping Nortel intact through temporary
government ownership, the solution that saved Chrysler, the chief employer in his hometown.

Mr.~Strand, however, has a message from Denmark for anxious Canadians. Several Danish companies were
once leading makers and designers of handsets. Acquisitions and failures, however, either eliminated
them or put those companies under foreign control. But subsequent investments by foreign owners in
Denmark, Mr.~Strand said, have made the country's wireless industry bigger than ever and based
largely around high-paying engineering, research and development jobs.

``If you have the right people with the right skills, you have no problem,'' Mr.~Strand said.

\section{For Winter Games in Vancouver, Ice Isn't So Easy }

\lettrine{I}{ce} appears such a simple concoction\cite{concoction} -- water turned really cold. This
time of year, its favored role is usually the one it plays at the bottom of a glass.

But ice is a year-round, lifelong science for the people hired by the Vancouver Organizing Committee
to produce complex forms of it for the Winter Olympics, coming in February.

Beyond the usual challenges in constructing ice surfaces to meet the needs of different sports in
different arenas, Vancouver's location presents a twist, with its combination of sea-level elevation
and high humidity, unique among Winter Olympic host cities.

Toss\cite{toss} in the usual broadcast demands accompanying the Olympics -- extra banks of hot
lights and inflexible event schedules among them -- and the so-called icemeisters may want to make
some extra cubes for a stiff drink.

Already, six months before the Olympics, they oversee a busy calendar of events and training
sessions to test the limits of the refrigeration systems, and gauge the effects of spectators in the
building and athletes on the ice. They experiment with water filtration and air circulation and
train teams of workers in the art of resurfacing.

About half of the Olympics will take place on carefully crafted ice between one and two inches
thick. Varying by arena, it must be a specific temperature, texture, composition, even color (thank
you, television), whether spread across a vast surface inside an arena or down the side of a
mountain. It must hold its consistency for weeks despite the collaborative efforts of ice's enemies,
from the obvious (the sun, sharp objects and thundering 1,400-pound bobsleds) to the not-so-obvious
(open arena doors, spectators and spinning 90-pound figure skaters).

``You can't just go out there and make ice,'' said Hans Wuthrich, in charge of the surface at the
newly built curling arena, where the final step is a delicate spritz of scientifically configured
water droplets strong enough to alter the course of 44 pounds of sliding granite.

The five ice specialists, each with deep Canadian ties, have extensive experience from previous
Olympics. On behalf of ice, they helped design new locales and the upgrades to existing ones. They
toured Vancouver's water-treatment plants to study their product's key ingredient. They ponder every
ice-dooming possibility.

Consider the challenge facing Tracy Seitz, who will make ice for competitors in bobsled, luge and
skeleton at the Whistler Sliding Centre.

The serpentine track, nearly a mile long, starts at an elevation of 3,080 feet and drops to 2,582
feet. In February, it sometimes snows at the top and rains at the bottom. Occasionally, the opposite
occurs.

The track's straight parts are U-shaped. High-banked curves are C-shaped, arcing up about 15 feet to
a gravity-defying overhang to accommodate speeds of more than 95 miles an hour.

Parts of the track angle toward the midday sun. Parts are always in shadow. Some competitions will
be at night.

In mid-September, with the track chilled by refrigeration running through pipes in the concrete,
Seitz and his crew will spray water from hoses until there is an inch of ice. That is thick enough
to absorb abuse from sliders experiencing more than 5 Gs of force around a corner. It is thin enough
to let the refrigeration system below the ice combat threats above it.

The trick is to keep the ice, no matter the forces working against it, between 23 to 25 degrees
Fahrenheit. Warmer, and the outer layers can turn slushy. Colder, and the brittle ice can crack in
sheets. Inconsistencies can be dangerous, even fatal.

``We want to make sure the ice doesn't change considerably in the heat of a race,'' Seitz said.

Competitions can take several hours, he noted. But medals are often determined by hundredths of a
second.

Even indoors, most ice begins the same way -- atop concrete disguising a maze of refrigeration
pipes. Water is added in thin layers because that makes stronger ice than filling a rink with an
inch of water as if it were a giant ice-cube tray.

``The water in Vancouver is incredible,'' said Mark Messer, the icemeister for the Richmond Olympic
Oval, where speed skating will take place. ``It's very, very clean. We have a filtration unit that
we're using at Richmond, and when putting it through there it almost makes the water too clean. You
need to have a slight bit of impurities there to just kind of hold things together.''

Messer, whose full-time job is plant manager and ice specialist at the Olympic oval in Calgary,
Alberta, billed as the fastest ice in the world, experimented ``to give us a combination of a lot of
glide in the ice and still have a lot of grip so the ice doesn't break when you make it cold.''

Atop its thin base, the ice is painted, usually in a tone commonly called Broadcast Gray. It looks
white compared with concrete slabs and creates a stark, television-friendly background for the
kaleidoscope of colors that swirl and race atop it. Lines and logos are added with stencils or
applied like decals.

Sensors are placed in the ice before final layers are added. Other sensors measure air temperature
and humidity in the building. They warn of the slightest changes. Similar controls will be in place
at the two hockey arenas, overseen by Dan Craig, facilities operations manager for the National
Hockey League.

For Vancouver, the biggest concern is humidity, which averages about 80 percent during February. The
goal is to keep it below 50 percent inside each arena, and below 40 percent in most, even if it is
pouring rain outside. When humid air hits ice, frost is created.

``As soon as you have frost, you can't get that out of your ice,'' the icemeister Kameron Kiland
said, as if describing red wine on white carpet.

Vancouver's indoor arenas are packed with extra dehumidifiers, more than at other Olympics.
Spectators will enter through large, tented buffer zones between the outside elements and the
surface of the ice.

``The big thing for us is the load-in of the spectators,'' said Kiland, overseeing ice for Pacific
Coliseum, where both figure skaters (who want soft ice) and short-track speed skaters (who want hard
ice) will alternately compete. ``If you've got a coliseum with upwards of 13,000 fans loading in for
an event and they're all wearing fleeces, catching moisture and bringing it in, you've got to dry
that off and catch it before they bring it into the bowl.''

The last thing the damp fans will be thinking about is the ice, in all its complex transparency. 

\section{Living and Fighting Alongside Men, and Fitting In }

\lettrine{T}{here} is no mistaking that this dusty, gravel-strewn camp northeast of Baghdad is
anything other than a combat outpost in a still-hostile land. And there is no mistaking that women
in uniform have had a transformative effect on it.

They have their own quarters, boxy trailers called CHUs (the military's acronym for containerized
housing units, pronounced ``chews'').

There are women's bathrooms and showers, alongside the men's. Married couples live together. The
base's clinic treats gynecological problems and has, alongside the equipment needed to treat the
trauma\cite{trauma} of modern warfare, an ultrasound machine.

Opponents of integrating women in combat zones long feared that sex would mean the end of American
military prowess. But now birth control is available -- the PX at Warhorse even sold out of condoms
one day recently -- reflecting a widely accepted reality that soldiers have sex at outposts across
Iraq.

The wars in Iraq and Afghanistan are the first in which tens of thousands of American military women
have lived, worked and fought with men for prolonged periods. Wars without front lines, they have
done more than just muddle\cite{muddle} the rules meant to keep women out of direct enemy contact.

They have changed the way the United States military goes to war. They have reshaped life on bases
across Iraq and Afghanistan. They have cultivated a new generation of women with a warrior's
ethos\cite{ethos} -- and combat experience -- that for millennia was almost exclusively the preserve
of men.

And they have done so without the disruption\cite{disruption} of discipline and unit cohesion that
some feared would unfold at places like Warhorse.

``There was a lot of debate over where women should be,'' said Brig. Gen. Heidi V.~Brown, one of the
two highest ranking women in Iraq today, recalling the start of the war. ``Here we are six years
later, and you don't hear about it. You shouldn't hear about it.''

In many ways, General Brown's career trajectory\cite{trajectory} since the war began reflects the
expanded role for women at war.

In 2003, as a colonel\cite{colonel}, she commanded a Patriot air-defense brigade\cite{brigade} that
joined the push from Kuwait to Baghdad, losing nine soldiers in a maintenance battalion outside
Nasiriya three days after the invasion began. One of them, Pfc. Lori Ann Piestewa, was the first
woman killed in action in Iraq; Pfc. Jessica D.~Lynch was captured in the same attack. Now, as the
American role in the war declines, General Brown will oversee the logistics of withdrawing the vast
amounts of military hardware in Iraq over the next year.

``We've needed -- needed -- the contributions of both our men and women,'' said Brig. Gen. Mary
A.~Legere, the director of intelligence for the American war effort here and the other highest
ranking woman in Iraq.

The military, of course, is not gender blind, especially in a war zone.

Sexual harassment in a still-predominantly male institution remains a problem. So does sexual
assault. Both are underreported, soldiers and officers here say, because the rigidity of the
military chain of command can make accusations uncomfortable and even risky for victims living in
close quarters with the men they accuse.

As a precaution, women are advised to travel in pairs, particularly in smaller bases populated with
Iraqi troops and civilians. Capt.~Margaret D.~Taafe-McMenamy, commander of the intelligence analysis
cell at Warhorse, carries a folding knife and a heavy, ridged flashlight -- a Christmas gift from
her husband, whom she lives with here -- as a precaution when she is out at night on the base.

Staff Sgt. Patricia F.~Bradford, 27, a psychological operations soldier, said that
slights\cite{slight}, subtle and not, were common, and some were easier to brush off than others.
Women are still viewed derisively at times in the confined, occasionally tense space of an outpost
like Warhorse.

``You're a bitch, a slut or a dyke -- or you're married, but even if you're married, you're still
probably one of the three,'' Sergeant\cite{sergeant} Bradford said.

At the same time, she and other female soldiers cope with the slights, showing a disarming
brashness.

``I think being a staff sergeant -- and a bitch -- helps deflect those things,'' she added.

The issues that arise in having women in combat -- harassment, bias, hardship, even sexual relations
-- are, she and others said, a matter of discipline, maturity and professionalism rather than an
argument for separating the sexes.

Sergeant Bradford recalled the day during her first tour when her convoy moved south while a soldier
with whom she was then engaged to be married moved north on the same highway. She listened on the
radio as his convoy came under an attack that continued after she was out of range.

``For four days, I had no idea what happened to him,'' she said, ``but I still had to continue my
mission, because that's what you do when you're a soldier.'' (He emerged unscathed, she later
learned.)

Unforeseen Issues

Such issues were not foreseen when the war in Iraq began in 2003, even though the initial invasion
force included women in the vanguard\cite{vanguard}.

On a practical level, the military was not prepared to house and otherwise address the specific
needs of women in a war zone -- including issues like health and privacy.

Early on, bases were largely makeshift and far more dangerous. Few soldiers, male or female, had
more than rudimentary\cite{rudimentary} quarters or latrines\cite{latrine}. None had much privacy.

Sgt. Dawn M.~Cloukey, a communications specialist, spent her first tour in Iraq in 2005 and 2006 as
the only woman among 45 soldiers, operating a retransmission station in the mountains of northern
Iraq and then in the center of Baghdad. She lived out of a rucksack, with no toilet or room of her
own. She described the experience as isolating.

``I always felt like the plague\cite{plague},'' she said at Warhorse, on her second tour in Iraq,
where she handles communications for the commander of the First Stryker Brigade of the 25th
Infantry\cite{infantry} Division.

As the United States military settled into more permanent bases, many initial difficulties abated,
as the Army gradually adapted to the new reality of waging war with a mixed force. So have the
soldiers themselves.

Women have sought acceptance in a still-predominately male environment not by emphasizing their sex
but rather by displaying their toughness, their willingness to adjust to conditions that are less
than ideal.

``I've kicked my guys out of the truck to pee in a bottle like that,'' Sgt. Joelene M.~Lachance, a
soldier with the 172nd Military Intelligence Battalion, said at Warhorse, pointing to one of the
liter water bottles that are ubiquitous at bases in Iraq. ``Cut the bottle off and pee in the bottle
and then dispose of it. Sometimes it's an issue, but most of the time, I just make do.

``I don't try to, like, 'I can't sleep here,' '' she continued. ``If they're sleeping there, I'm
sleeping there. I spent five days out in the truck once -- with six of my guys, sleeping on the
floor.''

Warhorse still reverberates\cite{reverberate} with the rumble\cite{rumble} of armored
convoys\cite{convoy} and the thud\cite{thud} of helicopters ferrying troops and, at times, the
wounded. It is just north of Baquba, the regional capital of Diyala Province, one of the most
restive\cite{restive} provinces in Iraq. Here, the war is not over. Warhorse will very likely be
among the last bases to close in Iraq before American troops withdraw in full.

At the outset of the war, the introduction of women into outposts like Warhorse raised fears not
just of abuse or harassment, but also of sex and pregnancy. The worst of those fears, officers say,
have not materialized.

In fact, sex in America's war zones is fairly common, soldiers say, and has not generally proved
disruptive\cite{disruptive}.

In April, the latest iteration of General Order No.~1, the rules governing the behavior of soldiers
in Iraq broadly, quietly relaxed the explicit prohibition on sex in a war zone, though it still bars
sex with Iraqis and spending the night in someone else's CHU. Some commands, including Baghdad,
retain broader restrictions, for example, on being in CHUs belonging to members of the opposite sex.

``The chain of command already has to deal with enough,'' Captain Taafe-McMenamy said. ``They don't
really want to have to punish soldiers for dating.''

Women do become pregnant -- a condition that, intentional or not, in or out of wedlock, requires the
woman to be flown out within two weeks, causing personnel disruptions in individual units.

The Army and Marine Corps declined to say exactly how many women left Iraq and Afghanistan as a
result of pregnancies, but it appears to be relatively rare and has had little effect on overall
readiness, commanders say. At Warhorse, the First Stryker Brigade, which has thousands of soldiers,
has sent only three women home because of pregnancies in 10 months in Iraq, the brigade said.

``There was a fear if we integrate units, you will have a bunch of young people with raging
hormones, and it will end up in too many unwanted pregnancies, and it's more trouble than it's
worth,'' said Peter Mansoor, a former battalion commander in Iraq who, until retiring recently,
served as Gen. David H.~Petraeus's executive officer. ``With good leadership and mentorship, we have
been able to keep those problems to a minimum.''

Taking On New Roles

Roughly 1 in 20 of the 5,600 soldiers at Warhorse is female, a smaller ratio than in the military as
a whole. Nonetheless, they are fully integrated in the base's operations.

Many of the women at Warhorse serve in jobs that have traditionally accommodated women: the base
hospital, food service, supply and administration.

Others, though, serve on the brigade staff, in intelligence and psychological operations, which
until recently were part of the Special Forces and thus off limits to women.

``We have changed so much,'' Col.~Burt K.~Thompson, the commander at Warhorse, said of the Army,
noting that every time he leaves the base, his patrol includes two women, including Sergeant Cloukey
``on comms'' -- communications -- and a medic, Sgt. Evette T.~Lee-Stewart. ``To have a female on an
infantry brigade staff? Oh my God.''

Like many commanders who have served in Iraq or Afghanistan, he said that women have ended the
debate over their role by their performance.

``I've relieved males from command,'' he said. ``I've never relieved a female commander in two and a
half years as commander.''

The nature of the war has also done much to change the debate over combat roles. Any trip off the
heavily secured bases now effectively invites contact with the enemy.

Many women have also been pulled off their regular jobs and trained to search Iraqi women at
checkpoints because of local cultural sensitivities, putting them as much at risk as any male
counterpart.

When Specialist Jennifer M.~Hoeppner goes ``outside the wire'' at Warhorse, as going on patrol is
known, she clambers into what she calls ``the best seat in the truck,'' the turret atop the Army's
newest armored vehicle, the MRAP.

``I'm the gunner on all our missions,'' she said, having qualified for the M240B machine gun at an
expert level.

``I think some of the males are a little confused when I go up,'' Specialist Hoeppner said.
``They're like, 'Who's your gunner?' ''

Women are also increasingly ``attached'' to infantry and armored units that train and advise Iraq's
police and military forces. Now that almost all American combat forces have pulled back to bases
outside of Iraq's cities, that training has become the main mission in Iraq.

The involvement of women in it has been a cultural shock for Iraqi men far less accustomed to
dealing with women professionally, especially in the military.

Women spoke of inappropriate comments or uncomfortable flattery, and even gifts. ``It was everything
from candy to lingerie,'' said Capt.~Victoria Ferreira, 29, who spent a year with an 11-person squad
training Iraqi officers. ``How do you react to that? 'Thank you?' ''

For the most part, though, Iraqis seem to accept the role of women in the American military -- they
have even expanded their own ranks for tasks like searching women at checkpoints -- even if it seems
unlikely that women will be incorporated more widely into the Iraqi armed forces anytime soon.

``I think now, six years since the war started, they've learned to adapt or tolerate the fact that
in the American Army we have high ranking positions that are filled by women,'' said Capt.~Violeta
Z.~Sifuentes, who commands the 591st Military Police Company.

It was not always so, she recalled of her first tour in Samarra in 2006. ``They always thought my
platoon sergeant or my squad leader was the one in charge until I was like, 'Listen here. I'm in
charge whether you like it or not.' ''

The captain's remarks were typical. The women serving in today's military represent a generational
shift. They are confident young women who have not had to fight the same gender battles their
predecessors in uniform did.

``I never felt like I had to fight to succeed in the Army'' was how Captain Taafe-McMenamy, who is
27, put it.

Adapting to the Tasks

Women in today's military say they do not feel the same pressure to prove themselves. They adapt and
expect others to adapt. They preserve their femininity without making much of it.

Specialist Hoeppner and her roommate, Sergeant Bradford, belong to the 361st Tactical Psychological
Operations Company, which patrols the towns and villages of Diyala with infantry squads to spread
and collect information.

On a recent patrol in the small village of Shifta, they seemed more of a novelty to the Iraqis they
encountered than the soldiers they patrolled with, taking up defensive positions alongside their
male colleagues whenever they paused.

``I actually had this million-dollar idea my first deployment,'' Sergeant Bradford said of her tour
as a truck driver hauling supplies in 2004. ``I was like, I need something that's like a beer bong
that I can hold in place so I can pee standing up without pulling my pants down. Cause we were truck
drivers. We'd stop on the side of the road. There's no bushes. I was telling one of my soldiers
about this great idea, and he said they already make that.''

She produced from her bunk in her CHU a device sold by REI called a ``feminine urinary director.''
``It's even pink,'' Specialist Hoeppner interjected.

Warhorse's supply officer -- a woman -- acquired dozens of them.

``The first time one of them came around a truck and saw me peeing on a tire,'' she said of one of
her male colleagues, ``I thought he was going to have a heart attack.''

\section{H.P. Tries to Keep the Ink Flowing}

\lettrine{Y}{ou've} heard about the fears of a jobless economic recovery. For Hewlett-Packard, there
is an even darker threat looming as the world works its way out of a recession: an inkless recovery.

In the last year, H.P.'s sales of printers and related supplies like ink have tumbled. The major
cause for the decline remains the weak global economy. Businesses have spent less on printing
products, and with unemployment high, fewer workers are around to hit the ink-burning Print button
on their computer screens.

Analysts contend that the recession has created a culture of reluctant printer users. Companies have
urged workers to keep a close eye on printing costs. So they are using less paper and choosing
costly color ink as a last resort. Such policies push people toward digital documents they read on
mobile devices as well as their computer screens.

``People have learned some of these lessons more rapidly as a result of the economic conditions than
they have in the past,'' said Steve Reynolds, an analyst with Lyra Research, which follows the
printing and imaging industries. ``My thought is that some of these habits will become permanent.''

In addition, online videos, Web-based TV and social networking sites -- where consumers are spending
more time -- do not lend themselves to stirring the inky rivers. ``A Twitter stream is not
printable,'' Mr.~Reynolds said.

H.P.'s printing group has long been one of the company's star performers. It accounts for nearly a
quarter of overall revenue. Printer ink remains one of the most expensive liquids on the planet --
more valuable than expensive perfumes -- providing H.P. with far higher profit margins than PCs and
other types of computing hardware provide.

On Tuesday, H.P. showed how its printer business remained vulnerable to the recession when it
reported third-quarter financial results. H.P.'s printing and imaging revenue fell 20 percent, to
\$5.7 billion, as sales of supplies tumbled 13 percent and sales of printers fell 23 percent.

These latest results add to several months of sharper than expected declines from H.P.'s printing
group, and the company has been scrambling to raise prices and adjust its inventory levels to offset
the slump.

``What has been the most striking issue for investors has been the fact that printer supplies have
really fallen off,'' said A.~M. Sacconaghi, an analyst for the investment research firm Sanford
C.~Bernstein. ``There was always this belief that people keep printing like they keep eating even
during recessions.''

H.P.'s overall business has recovered as sales in China and the United States have added some
stability to the technology sector, Mark V.~Hurd, H.P.'s chief executive, said in an interview.

H.P., based in Palo Alto, Calif., reported net income of \$1.6 billion, or 67 cents a share, for the
quarter ended July 31, a 19 percent decline from the \$2 billion in net income, or 80 cents a share,
reported in the period last year. Excluding charges, H.P. earned 91 cents a share, beating a
forecast of 90 cents a share from analysts polled by Thomson Reuters.

H.P.'s revenue declined 2 percent, to \$27.5 billion, from \$28 billion a year ago.

H.P.'s other businesses struggled during the quarter as well, with PC sales tumbling 18 percent and
computer server and storage sales falling 23 percent. The company's main bright spot was its
services division, which recorded a 93 percent increase in revenue to \$8.5 billion on the back of
the Electronic Data Systems acquisition.

The company's shares declined 2 percent in extended trading after the announcement.

Although H.P. has been carefully watching printing habits change for years, Mr.~Hurd said H.P. was
not in danger.

He said people examining H.P.'s printing results needed to take into account currency fluctuations
that had hurt the company, as well as costs tied to inventory adjustments. For example, such costs,
Mr.~Hurd said, would account for most of the 13 percent slide in printer supplies revenue last
quarter. The sales of ink by companies like Best Buy, Wal-Mart and Tesco were flat to slightly down,
he said.

``People are printing just as much as they did last year,'' Mr.~Hurd said. Rather than curtailing
their printing behavior, businesses and people have simply opted not to buy spare ink cartridges at
the same rate as they did during better financial times, according to Mr.~Hurd.

Still, H.P's printing supply revenue had grown at close to 10 percent in the three previous years,
and has now fallen. And it is that sharp shift that has investors worried whether longer-term
changes in behavior are at work, Mr.~Sacconaghi said.

H.P. has spent years trying to carve out new printing businesses, while also tacitly acknowledging
changes in consumer behavior, by seeking ways to make more Web content like Twitter streams and
MySpace pages printable.

The company sells large machines like its Indigo digital presses for handling commercial printing
jobs, ranging from printing magazines to billboards to labels on consumer products.

Digital presses produce only about 12 percent of the more than 50 trillion printed objects made each
year, according to H.P. Each point of market share H.P. can gain in this large-scale printing
segment could translate into vast increases in revenue.

H.P. has been encouraging businesses to use digital presses as money-making tools, since they are
more flexible than cheaper, traditional analog presses.

For example, H.P. has pushed the idea that retailers like Wal-Mart could use in-store digital
presses to quickly make placards offering price cuts on certain products or to tweak the placards
for each store's demographics, making Spanish signs in certain regions for instance.

For consumers, H.P has developed software that can manipulate Web content into a more printable
format. The company's researchers have found that girls ages 14 to 24 who tend not to print much
will in fact turn things like collections of photos and comments from MySpace into a foldable cube
that can be shared with their friends.

\section{Appeals Court Orders New Trial in Brocade Case }

\lettrine{A}{ federal} appeals court has tossed\cite{toss} out the criminal conviction of the first
Silicon Valley executive to go to trial in a stock options scandal that triggered investigations at
hundreds of companies and charges against at least a dozen executives.

The 9th U.S.~Circuit Court of Appeals ordered a new trial Tuesday for Gregory Reyes, former chief
executive of Brocade Communications Systems Inc., citing prosecutorial misconduct. The court said a
prosecutor\cite{prosecutor} lied when he told the jury that Brocade's finance department was unaware
that Reyes was doling out backdated stock options to company employees.

A judge last year sentenced Reyes to one year and nine month in prison and fined him \$15 million
after a jury convicted him on all 10 felony\cite{felony} counts, including fraud, falsified
accounting, conspiracy and filing false financial statements.

``In representing the United States, a federal prosecutor has a special duty not to
impede\cite{impede} the truth,'' the appeals court wrote. ``The record demonstrates that the
prosecution argued to the jury material facts that the prosecution knew were false, or at the very
least had strong reason to doubt.''

The court ordered a new trial, noting that the government's original case appeared ``relatively
strong.''

The appeals court on Tuesday also upheld the conviction of Stephanie Jensen, the company's former
vice president of human resources. But the court ordered that she be given a new sentence for
falsifying corporate records. Jensen was sentenced last year to serve four months in prison and pay
\$1.25 million.

Jensen and Reyes have remained free pending their appeals.

The Reyes' case was the government's highest-profile backdating prosecution. A spokesman with the
U.S.~Attorney's office in San Francisco didn't return a telephone call Tuesday. The prosecutor
accused of misconduct, Timothy Crudo, has left the U.S.~Attorney's office for private practice.
Crudo didn't return a call seeking comment.

Reyes' primary defense during the six-week trial in 2007 was that the finance department knew of his
plan to retain talented employees by granting them options to buy the company's stock at
artificially deflated prices. Reyes' position at trial was that he relied on the finance department
to accurately record the stock options.

Reyes wasn't accused of personally profiting from the backdating. But he was accused of defrauding
investors by withholding information from Brocade's finance department that caused them to file
faulty financial statements.

Prosecutors said Reyes' scheme made Brocade appear more profitable than it was, which indirectly
helped him pump up the company's stock price and command higher pay for himself.

San Jose, Calif.-based Brocade, which makes switches that connect companies' servers to their data
storage systems, wiped out hundreds of millions of dollars in previously reported profits from 1999
to 2004 after its books were corrected.

\section{Spokesman for Taliban Is Captured, Pakistan Says}

\lettrine{P}{akistan} has captured the chief spokesman of the Pakistani Taliban, security officials
said Tuesday, dealing another blow to the militant network two weeks after its top commander was
apparently killed in an American airstrike.

The spokesman, Maulvi Said Muhammad, who also goes by the name Maulvi Omar, was captured by local
Pakistani tribesmen in Mohmand Agency, a Taliban stronghold in western Pakistan, and turned over to
Pakistani forces, security officials said.

The officials, who spoke on the condition of anonymity because they were not authorized to discuss
the arrest publicly, said Mr.~Omar confirmed the death of the top commander, Baitullah Mehsud,
during an interrogation.

It was impossible to verify their accounts.

One security official said Mr.~Omar, who was captured along with two bodyguards, had indicated that
there were ``problems'' within the movement over leadership.

Security officials said Mr.~Omar had told them that he wanted to go to the Orakzai tribal region for
consultations over a successor to Mr.~Mehsud, but that he was forced to try to return to his native
Bajaur tribal region because of fighting between militant factions along the route. He was stopped
by pro-government tribesmen, turned over to the paramilitary Frontier Corps and then airlifted to
Peshawar, the main city in western Pakistan.

Televised images showed a prisoner believed to be Mr.~Omar being bundled into a helicopter. One
senior military official said Mr.~Omar had apparently been beaten by the tribesmen.

Separately, a suicide bomber rammed his explosive-laden car into a security checkpoint in North
Waziristan on Tuesday, killing seven people, according to a Pakistani security official.

\section{Mental Stress Training Is Planned for U.S.~Soldiers}

\lettrine{T}{he} Army plans to require that all 1.1 million of its soldiers take intensive training
in emotional resiliency, military officials say.

The training, the first of its kind in the military, is meant to improve performance in combat and
head off the mental health problems, including depression, post-traumatic stress disorder and
suicide, that plague about one-fifth of troops returning from Afghanistan and Iraq.

Active-duty soldiers, reservists and members of the National Guard will receive the training, which
will also be available to their family members and to civilian employees.

The new program is to be introduced at two bases in October and phased in gradually throughout the
service, starting in basic training. It is modeled on techniques that have been tested mainly in
middle schools.

Usually taught in weekly 90-minute classes, the methods seek to defuse or expose common habits of
thinking and flawed beliefs that can lead to anger and frustration -- for example, the tendency to
assume the worst. (``My wife didn't answer the phone; she must be with someone else.'')

The Army wants to train 1,500 sergeants by next summer to teach the techniques.

In an interview, Gen. George W.~Casey Jr., the Army's chief of staff, said the \$117 million program
was an effort to transform a military culture that has generally considered talk of emotions to be
so much hand-holding, a sign of weakness.

``I'm still not sure that our culture is ready to accept this,'' General Casey said. ``That's what I
worry about most.''

In an open exchange at an early training session here last week, General Casey asked a group of
sergeants what they thought of the new training. Did it seem too touchy-feely?

``I believe so, sir,'' said one, standing to address the general. He said a formal class would be a
hard sell to a young private ``who all he wants to do is hang out with his buddies and drink beer.''

But others disagreed, saying the program was desperately needed. And in the interview, General Casey
said the mental effects of repeated deployments -- rising suicide rates in the Army, mild traumatic
brain injuries, post-traumatic stress -- had convinced commanders ``that we need a program that
gives soldiers and their families better ways to cope.''

The general agreed to the interview after The New York Times learned of the program from Dr.~Martin
E.~P. Seligman, chairman of the University of Pennsylvania Positive Psychology Center, who has been
consulting with the Pentagon.

In recent studies, psychologists at Penn and elsewhere have found that the techniques can reduce
mental distress in some children and teenagers. But outside experts cautioned that the Army program
was more an experiment than a proven solution.

``It's important to be clear that there's no evidence that any program makes soldiers more
resilient,'' said George A.~Bonanno, a psychologist at Columbia University. But he and others said
the program could settle one of the most important questions in psychology: whether mental toughness
can be taught in the classroom.

``These are skills that apply broadly, they're things people use throughout life, and what we've
done is adapt them for soldiers,'' said Karen Reivich, a psychologist at Penn, who is helping the
Army carry out the program.

At the training session, given at a hotel near the university, 48 sergeants in full fatigues and
boots sat at desks, took notes, play-acted, and wisecracked as psychologists taught them about
mental fitness. In one role-playing exercise, Sgt. First Class James Cole of Fort Riley, Kan., and a
classmate acted out Sergeant Cole's thinking in response to an order late in the day to have his
exhausted men do one last difficult assignment.

``Why is he tasking us again for this job?'' the classmate asked. ``It's not fair.''

``Well, maybe,'' Sergeant Cole responded. ``Or maybe he's hitting us because he knows we're more
reliable.''

In another session, Dr.~Reivich asked the sergeants to think of situations when such internal
debates were useful.

One, a veteran of several deployments to Iraq, said he was out at dinner the night before when a
customer at a nearby table said he and his friends were being obnoxious\cite{obnoxious}.

``At one time maybe I would have thrown the guy out the window and gone for the
jugular\cite{jugular},'' the sergeant said. But guided by the new techniques, he fought the
temptation and decided to buy the man a beer instead. ``The guy came over and apologized,'' he said.

The training is based in part on the ideas of Dr.~Aaron Beck and the late Albert Ellis, who found
that mentally disputing unexamined thoughts and assumptions often defuses them. It also draws on
recent research suggesting that people can manage stress by thinking in terms of their psychological
strengths.

``Psychology has given us this whole language of pathology\cite{pathology}, so that a soldier in
tears after seeing someone killed thinks, 'Something's wrong with me; I have post-traumatic stress,'
'' or P.T.S.D., Dr.~Seligman said. ``The idea here is to give people a new vocabulary, to speak in
terms of resilience\cite{resilience}. Most people who experience trauma don't end up with P.T.S.D.;
many experience post-traumatic growth.''

Many of the sergeants were at first leery\cite{leery} of the techniques. ``But I think maybe it
becomes like muscle memory -- with practice you start to use them automatically,'' said Sgt. First
Class Darlene Sanders of Fort Jackson, S.C.

To track the effects of the program, the Army will require troops at all levels, from new recruits
to officers, to regularly fill out a 170-item questionnaire to evaluate their mental health, along
with the strength of their social support, among other things.

The program is not intended to diagnose mental health problems. The results will be kept private,
General Casey said.

The Army will track average scores in units to see whether the training has any impact on mental
symptoms and performance, said Gen. Rhonda Cornum, the director of Comprehensive Soldier Fitness,
who is overseeing the carrying out of the new resilience program. General Cornum said that the Army
had contracted with researchers at the University of Michigan to determine whether the training was
working, and added that corrections could be made along the way ``if the program is not having the
intended effect.''

This being the Army, the sergeants at the training session last week had questions about logistics.
How would teachers be evaluated? How and when would Reserve and Guard units get the training?

Perhaps the biggest question -- can an organization that has long suppressed talk of emotions now
open up? -- is unlikely to have an answer until next year at the earliest. But the Army's leaders
are determined to ask.

``For years, the military has been saying, 'Oh, my God, a suicide, what do we do now?' '' said
Col.~Darryl Williams, the program's deputy director. ``It was reactive. It's time to change that.''

\section{Chinese Public-Interest Lawyer Charged Amid Crackdown}

\lettrine{P}{rosecutors} have charged one of China's leading public-interest lawyers, Xu Zhiyong,
with tax evasion, his lawyer said on Tuesday, continuing a government crackdown on this nation's
small band of activist lawyers and scholars that has lasted months.

Mr.~Xu, 36, is a founder of the Open Constitution Initiative, known in Chinese as Gongmeng, a
nonprofit group that often has taken on high-profile cases involving citizens' civil rights. The
government shut down the organization's Gongmeng legal center on July 17, three days after accusing
it of tax violations, and the police seized Mr.~Xu on July 29.

In an interview on Tuesday, his lawyer, Zhou Ze, said Mr.~Xu was formally charged on Aug.~12. Mr.~Xu
could face seven years in prison if he is tried and convicted. The prosecutors now must seek an
indictment\cite{indictment}, but that is widely considered a formality.

The government's main accusation is that Mr.~Xu's group failed to pay taxes on a \$100,000 grant
from Yale University that was earmarked for the legal center. But human rights advocates and foreign
political analysts agree that the charges are politically inspired, part of what seems to be a
growing effort by security officials to shut down independent activism and especially activism that
is supported with foreign funds.

The government has moved this year to block many foreign-based Web sites and social-networking
services used by Chinese activists and, often, by Chinese citizens. It also has taken action against
a host of activist scholars and lawyers, effectively disbarring about 50 lawyers earlier this
summer. Gao Zhisheng, whose aggressive legal campaigns earned him a reputation as a gadfly, has not
been heard from since being taken into custody more than six months ago.

A number of activist lawyers have been beaten or harassed by unidentified assailants while working
this year on cases. And one of China's most prominent political dissidents, Liu Xiaobo, has been
held virtually incommunicado in a suburban Beijing detention center since December.

Separately, the Beijing financial publication Economic Observer reported on Tuesday that the
government had begun a broad inquiry into the so-called resident representative offices of
foreign-based enterprises -- in essence, offices that many foreign groups, including many charities
and nonprofit organizations, establish on Chinese soil.

The newspaper quoted an unnamed source as saying that the government was drafting new regulations
governing the offices and that many of the existing offices were suspected of violating Chinese law.

The current rules exempt the offices from paying taxes but also bar them from conducting business
activities. The unnamed source was quoted as saying that many offices have flouted that prohibition,
while others are guilty of lesser violations like failing to report address changes or renew expired
registrations.

The report stated that investigators had already begun visiting some resident representative
offices. In at least some instances, the investigators have been accompanied by the police, Nicholas
Bequelin, a researcher with the Hong Kong office of Human Rights Watch, said in an interview on
Tuesday.

The government's move against Gongmeng and Mr.~Xu, he said, has sent a chill sweeping over China's
activist organizations, in large part because Gongmeng is widely seen as one of the most scrupulous
groups working to expand the rule of law. Indeed, Mr.~Xu, a professor at Beijing's University of
Posts and Telecommunications, has been an elected member of a local governing body, the People's
Congress in Beijing's Haidian district, since 2003.

``He was doing everything aboveboard,'' said Mr.~Bequelin, who called Mr.~Xu ``the voice of
moderation'' in public-interest legal circles. ``If he goes down, who is safe?''

\section{One Person's Boondoggle, Another's Necessity}

\lettrine{T}{he} \$14.7 million for a new airport on an Alaskan island that averages only 42 flights
a month. The half-million dollars for a new skateboard park in unemployment-ravaged Rhode Island.
The \$3.4 million for fencing and tunnels to keep Florida turtles from becoming roadkill.

Those proposals for spending federal stimulus money were all criticized by cable news commentators,
Republican officials and, in the case of the airport, the inspector general of the Transportation
Department. But they have something else in common, too. They are popular locally. And they
underscore a truth that has been evident since the New Deal: sometimes the boondoggle is in the eye
of the beholder.

Dan Clarion, the maintenance manager and harbor master in Ouzinkie, a town of about 170 people on an
island off Alaska, said the new airport was vital for safety reasons. The island's current runway is
too short, Mr.~Clarion said, and its location leaves it vulnerable to dangerous crosswinds that
sometimes shut down the airport, stranding islanders who rely on planes for their mail, their
groceries and their access to the mainland when the seas are too high.

``Many times it's the only way we can get in and out,'' Mr.~Clarion said by telephone.

The outcry that greeted Pawtucket, R.I., over its plans to spend stimulus money on a playground with
a skate park at a junior high school was so fierce that city officials soon scrapped the idea and
decided to spend their stimulus money on road work instead. But residents of Pawtucket had sought
the skate park for a decade to lure skaters off the steps of the city's library, the granite walls
by the old mill and even the steps of City Hall. So the city is building it anyway, using the
original source of federal financing that it had planned to use before the stimulus law passed.

And Florida's plan to build a wall and culverts to help turtles cross a deadly stretch of road by
Lake Jackson near Tallahassee was mocked by Senator Tom Coburn, Republican of Oklahoma, in a report
he issued in questioning many stimulus projects.

``Why did the turtle cross the road?'' the report asked. ``To get to the other side of a stimulus
project.''

But local advocates have sought the passage there for years, monitoring the steady stream of turtle
deaths, building the temporary fence that protects turtles there now and mounting a large
grass-roots campaign that eventually won state support for a more permanent ``eco-passage.''

Of course, none of these projects are likely to stimulate much in the way of new, long-term economic
activity, a goal of the federal stimulus plan, once the construction crews decamp. But neither will
many of the street repaving jobs that are getting the bulk of the highway money in the stimulus
bill, however badly they may be needed.

One of the main purposes of the stimulus is to use government spending to prop up the economy in the
short term by offsetting the decline in private-sector spending. But in doing so, the Obama
administration has given rise to a debate that goes back many decades.

``The divide hasn't really changed since the Depression,'' said Nick Taylor, an author who
chronicled the fierce debates of the New Deal era over boondoggles in ``American-Made: The Enduring
Legacy of the W.P.A.'' (Bantam, 2008). Mr.~Taylor said many of the projects that were mocked back
then as boondoggles went on to have long and productive lives. But he added that there were some
truly wasteful projects back then as well, like a lake dug at a spot in North Carolina that turned
out not to have a water source.

Before it became a bad word, ``boondoggle'' was an innocent, humble craft. It was the Boy Scouts of
America who claimed credit for coining the word, to refer to the plaited leather lanyards that they
made and wore around their necks.

That all changed on April 3, 1935, at a hearing in New York City on how New Deal relief money was
being spent. A Brooklyn crafts teacher reluctantly testified that he was paid to show the jobless
how to make ``boon doggles.'' The outcry was swift. ``\$3,187,000 Relief is Spent to Teach Jobless
to Play,'' trumpeted a front-page headline the next day in The New York Times. `` 'Boon Doggles'
Made.''

A new, more sinister meaning was born, and the word came to signify government make-work, later
referring to wasteful government projects in general. Critics used it to criticize scores of
projects, but President Franklin D.~Roosevelt took a longer view. ``If we can boondoggle ourselves
out of this Depression,'' Roosevelt said, ``that word is going to be enshrined in the hearts of the
American people for years to come.''

The Depression ended, but the word has not exactly been enshrined in the hearts of the people. Now,
as a new administration in a new century tries to pump \$787 billion into the economy with a
combination of spending and tax cuts, questions about the usefulness of projects are raging as
loudly as ever, even as some economists begin to credit the stimulus with averting an even worse
downturn.

Boondoggle-phobia was a guiding principle in the writing of the bill. After a water park in Florida
was mocked for seeking stimulus money while the bill was being drafted, a provision was added
barring states and local governments from using their money to build swimming pools. (Swimming pools
may sound frivolous, but they are also popular: the government built hundreds of them during the New
Deal, including 11 of New York City's biggest, grandest pools, many of which are still in use.)

Another provision barred state and local governments from spending their stimulus money on zoos,
which were another big recipient of New Deal money. (The National Zoological Park in Washington,
which is part of the Smithsonian Institution, is apparently exempt: it is getting \$11.4 million
from the stimulus for repairs and safety projects.)

Robert D.~Leighninger Jr., a sociologist who wrote ``Long-Range Public Investment: The Forgotten
Legacy of the New Deal'' (South Carolina University Press, 2007), recounted the story of a Works
Progress Administration official in Arizona who went off in search of boondoggles, and discovered
that the towns he visited seemed to like their own projects but questioned those of their neighbors.

``I've been hunting all over the state for one, but everywhere I go I'm told it's in the next
county,'' the official was quoted as saying in a 1936 newspaper article. ``So far I haven't been
able to catch up with a real, live one.''

\section{Supreme Court to Hear Case on Executive Pay}

\lettrine{L}{ast} summer, Richard A.~Posner, a federal appeals court judge, issued a surprising and
prescient dissent. Executive pay is out of control, he said, and the marketplace cannot be trusted
to rein it in.

Judge Posner is a conservative with libertarian leanings, and he is a leader of the law and
economics movement associated with the University of Chicago. He often relies on economic analysis
in his judicial decisions, and he believes that many questions are best sorted out by the
marketplace.

But corporate America has insulated pay decisions from market discipline, Judge Posner wrote.
``Executive compensation in large publicly traded firms often is excessive,'' he added, ``because of
the feeble incentives of boards of directors to police compensation.''

The Supreme Court will hear the case this fall, as anger over huge bonuses paid to the executives of
failing companies continues to grow. The case, Jones v. Harris Associates, may turn out to be the
court's first significant statement on the corporate culture that helped lead to the Great
Recession.

The case arose from the enormous fees mutual\cite{mutual} funds pay to their investment advisers. A
three-judge panel of Judge Posner's court, the United States Court of Appeals for the Seventh
Circuit, in Chicago, threw out a lawsuit brought by the investors in three Oakmark mutual funds who
said the funds had overpaid their investment adviser, Harris Associates.

The panel decision, written by Chief Judge Frank H.~Easterbrook, another leader of the law and
economics movement, said the marketplace could be trusted to regulate fees. Judge Posner, dissenting
from the full court's decision not to rehear the case, said competition had not been effective in
keeping the compensation under control.

Before last year's market collapse, the mutual fund industry held more than \$11 trillion in
retirement and personal savings, and it paid advisers perhaps \$100 billion in fees.

Mutual funds are odd enterprises. They are typically formed and run by their investment advisers,
which select the fund's board of directors. That board then negotiates the adviser's fees.

Here is how Warren Buffett analyzed the situation in his 2003 letter to shareholders: ``Year after
year, at literally thousands of funds, directors had routinely rehired the incumbent management
company, however pathetic its performance had been. Just as routinely, the directors had mindlessly
approved fees that in many cases far exceeded those that could have been negotiated.''

The plaintiffs in the case before the Supreme Court claimed that Harris Associates had charged their
funds twice as much as it charged its unaffiliated clients, like pension funds.

The Oakmark funds paid Harris Associates 1 percent of the first \$2 billion in assets; independent
clients were charged roughly one-half of 1 percent of the first \$500 million. One percent of a
billion dollars is nice work if you can get it.

``Mutual funds rarely fire their advisers,'' Judge Easterbrook acknowledged. But, he continued,
``investors can and do 'fire' advisers cheaply and easily by moving their money elsewhere.'' A 2007
study from John C.~Coates IV and R.~Glenn Hubbard supported this conclusion, finding that mutual
fund fees are kept in check by the movement of investors' money.

But a brief supporting the plaintiffs filed in the Supreme Court by three economists, Ian Ayres,
Robert E.~Litan and Joseph R.~Mason, questioned that study. New research in behavioral economics,
the brief said, showed that most investors had a very poor grasp of rudimentary truths about
probability and a disproportionate aversion to taking losses.

Mutual fund investors thus tend to look at past performance rather than fees. And they have a
tendency to sell winning investments too early and hold losing ones too long.

Even if mutual fund investors could be counted on to act rationally, the economists' brief said,
they do not have ready access to the information they need to make sensible choices.

Instead of counting on investor behavior to keep fees in check, the brief concluded, courts should
look to how much advisers charged independent clients like pension funds. A supporting brief from
the federal government made the same point.

There is academic research to support this view, too.

``In contrast to mutual fund investors,'' Diane Del Guercio and Paula A.~Tkac wrote in a 2002 study,
``pension clients punish poorly performing managers by withdrawing assets under management and do
not flock disproportionately to recent winners.''

But Judge Easterbrook questioned the value of such comparisons. The two kinds of clients, he said,
may have different needs. In its brief urging the Supreme Court not to hear the case, Harris
Associates added that the Oakmark funds had outperformed ``virtually every fund in their peer
groups.''

Still, the tide seems to be turning toward skepticism about outsize compensation. In April, a month
after the Supreme Court agreed to hear an appeal from Judge Easterbrook's decision, the federal
appeals court in St Louis allowed a suit against another investment adviser, Ameriprise Financial,
to go forward. It was the first ruling in favor of unhappy mutual fund investors suing over
advisers' fees since Congress imposed a fiduciary duty on advisers in 1970.

Judge Easterbrook said the law had only a minor role to play, requiring no more than making sure
that advisers ``make full disclosure and play no tricks.''

But when public sentiment, economic research and even Judge Posner argue for more vigorous judicial
examination of whether compensation is fair, the Supreme Court may just agree.

\section{House Ethics Inquiry Has Roots in Untidy Yard}

\lettrine{C}{ould} an untended lawn and delinquent\cite{delinquent} mortgage payments lead to a
Congressional ethics investigation?

That is the question apparently facing Representative Laura Richardson, a Democrat whose property in
Sacramento has been the subject of questions by investigators from the Office of Congressional
Ethics.

The nonpartisan board, which has the power to recommend a formal investigation to the standards
committee of the House of Representatives, has questioned a neighbor and a real estate broker about
their dealings with Ms.~Richardson.

At the heart of the review is a modest three-bedroom home in the Curtis Park district of Sacramento
that Ms.~Richardson bought shortly after being elected to the State Assembly in 2006, and shortly
before being elected to Congress the next year.

The back-to-back electoral wins apparently prevented Ms.~Richardson, a former city councilwoman from
Long Beach, from taking up residence permanently in Curtis Park, a leafy, manicured neighborhood on
Sacramento's south side. Her ascendant political career also distracted her from lawn care,
residents of the street said.

``The front yard grass started getting overwhelmingly high, and she'd come and leave, and nothing
would ever get done,'' said Sean Padovan, 62, a retired police sergeant who lives four houses down
from the Richardson house. ``I figured she was busy. But it got worse and worse.''

Some of her neighbors found the property so unsightly that they decided to take the gardening into
their own hands.

And that is, apparently, where some of the questions began.

Peter Thomsen, a retired banker across the street, said he was recently interviewed by investigators
from the ethics office as to whether his horticultural efforts -- including watering the property's
beleaguered ivy plant -- were meant to curry Ms.~Richardson's favor.

``They were trying to determine if I had a positive or negative relationship'' with the
congresswoman, Mr.~Thomsen said. ``They were very clearly focused as to whether we had done work on
her property to her benefit.''

Questions about lawn care are not the only ones surrounding Ms.~Richardson's house, which the
congresswoman at one point lost to foreclosure but later regained ownership.

According to papers filed in Sacramento County, the Curtis Park home was bought from the bank
trustee that held the mortgage in May 2008 by a local real estate broker, James York, for \$388,000
-- roughly \$147,000 less than what Ms.~Richardson had paid a year before.

Shortly after that sale, however, the bank, Washington Mutual, rescinded the sale and the property
was returned to Ms.~Richardson, according to Mr.~York. Mr.~York said he sued both the bank and
Ms.~Richardson, and later reached an agreement over the property.

In an e-mail message on Friday, Mr.~York said he was prevented by that agreement from commenting
about the changes in ownership. ``I can't discuss the terms,'' he wrote.

But Mr.~York confirmed that he had been sent a letter of inquiry from the ethics office earlier this
year, though he did not reveal its contents.

Calls to the media line for Washington Mutual were not returned, and Ms.~Richardson's spokesman,
Michael Eagle, did not respond to repeated questions last week about the circumstances surrounding
Ms.~Richardson's regaining ownership of the house, which has been the subject of considerable
speculation in the California news media.

In a written statement issued on Ms.~Richardson's behalf, Mr.~Eagle called reports about her Curtis
Park property ``personal, biased, partisan and in some cases deliberately inaccurate.''

Mr.~Eagle also said in another statement that the property was ``neither deteriorating nor a
nuisance'' and was under renovation. And by early this month, there were signs of work under way
inside the house, and cut, green grass in the front yard.

Ron O'Connor, operations manager with the City of Sacramento code enforcement department, said his
officers had been to Ms.~Richardson's home on several occasions over the last year, for complaints
about issues like rotting fruit, overgrown weeds and a police report of a squatter living in the
garage. But Mr.~O'Connor concurred with Mr.~Eagle that the house was not blighted or neglected.

``I could live it in now if it had a little more property where I could have a garden,''
Mr.~O'Connor said.

Leo Wise, staff director and chief counsel of the Office of Congressional Ethics, said it was the
office's policy not to confirm or deny a review.

Reviews by the ethics office, which was established last year, have two phases. Preliminary reviews
take 30 days. If the board members vote for a further investigation, the office has 59 days to make
its recommendations to the House standards committee.

\section{Diabetes Case Shows Pitfalls of Treatment Rules}

\lettrine{I}{t} sounds like a simple idea for improving health care: draw up guidelines on how best
to treat a particular illness and then pay doctors to follow them. That strategy, which some
insurers and health plans already employ, has been embraced during the health care debate by some
lawmakers in Congress who want to extend it more broadly.

The goal is to improve treatment and, at the same time, save money. But setting guidelines that are
good for every patient, it turns out, can get messy, with some experts warning that a big national
plan of this sort poses risks. A recent case involving treatment for diabetes, one of the nation's
most pervasive illnesses, illustrates the difficulties.

Last year, a national guideline-setting group abruptly withdrew a controversial diabetes standard it
adopted in 2006 that called for aggressive control of blood sugar, or glucose. The change came after
a large federal study indicated that lowering glucose too quickly or too much in some patients could
harm or even kill them.

In medical journal articles and elsewhere over the last year, some diabetes experts have lashed out
at the group's initial decision to approve the guideline, saying they warned back in 2006 that it
was medically ill-advised for some patients.

``This was a case in which the advocates of a disease got caught up in their disease rather than the
interests of patients,'' said Dr.~Rodney A.~Hayward, a diabetes expert at the University of Michigan
who had opposed the benchmark.

Critics like Dr.~Hayward have also suggested that pharmaceutical companies influenced the guideline
so they could sell more glucose-lowering drugs like insulin. The group that set the guideline, a
Washington organization called the National Committee for Quality Assurance, received about \$3
million, or 10 percent of its revenue, last year from drug and medical device makers.

The group's officials, and several outside experts who advise it, rejected such suggestions of
industry influence. Still, some experts question why the group, which sets standards that are widely
used by insurers and private health plans, should take even small amounts of drug-industry money.

``There should not be any industry funding of a group that is involved in working on national
guidelines,'' said Dr.~Jerome E.~Groopman, a professor of medicine at Harvard.

To many experts, the diabetes case shows how setting one guideline that works for all patients
suffering from the same disease can be tricky.

The main problem is that many guidelines are based not on rigorous studies like clinical trials but
on weaker types of medical evidence. And critics like Dr.~Groopman have argued that the
guideline-setting process is often influenced by industry or by medical ideologues looking to
advance their personal agendas.

``These guidelines often come out of specialty societies who tend to want to treat rather than not
to treat,'' said Arthur Levin, the executive director of the Center for Medical Consumers, a patient
advocacy group. Mr.~Levin also serves as an adviser to the National Committee for Quality Assurance
and took part in the diabetes decision.

The diabetes controversy began after several studies showed that patients with diabetes could
significantly benefit from tighter control of their blood sugar. As a result, in 2005, the American
Diabetes Association, along with others, began urging the National Committee on Quality Assurance to
adopt an aggressive glucose control standard. The group's guidelines are used by insurers and others
to assess the performance of doctors who work for them and, in some cases, to determine whether they
qualify for a bonus.

On its face, the glucose proposal made sense. Excess levels of glucose can cause some of the worst
complications of diabetes, including blindness and kidney failure. But lowering glucose too much can
also create problems for some patients, particularly older, sicker ones, causing them to black out
or to have seizures.

Pointing to such risks, some doctors like Dr.~Hayward, who also works for the Veterans
Administration, opposed extending the guideline too broadly. They said that one significant group of
diabetics, those who also had heart problems or cardiovascular disease, had never been studied
during the trials that showed the benefits of aggressive glucose control.

But supporters of the move, like Dr.~Richard Kahn, then a top official of the American Diabetes
Association, rejected such arguments. Dr.~Kahn argued that even if the standard were adopted,
doctors would still have plenty of discretion about how to treat their patients. ``All of us can no
longer sit silently thinking that our discipline is unified toward improving diabetes care,'' he
wrote in a 2006 e-mail message as the argument unfolded.

In mid-2006, a panel assembled by the National Committee for Quality Assurance that included
insurers, major employers, doctors and Mr.~Levin, the patient advocate, unanimously approved the new
standard.

Critics like Dr.~Hayward were furious and later said they believed that diabetes experts with
drug-industry ties had played a role in guiding the decision.

Both Mr.~Levin and another adviser to the quality committee -- Dr.~Paul Wallace of Kaiser
Permanente, the big health care provider -- said they saw no evidence of that. Also, officers of the
National Committee for Quality Assurance insist that medical products companies, who sponsor some of
the group's activities, do not play a role in setting its treatment standards.

The diabetes standard was not in place for long. It was challenged last year when federal
researchers abruptly halted a study of glucose control in older diabetics with cardiovascular
disease, the group not previously studied.

In that trial, researchers aggressively treated patients with a variety of drugs in an effort to
bring their glucose down to the level of a person without diabetes, but patients receiving the
treatment started dying at an unexpected rate.

In light of the findings, the National Committee for Quality Assurance withdrew its standard and
convened an expert panel that this time included critics like Dr.~Hayward. A year ago, the group
issued a more nuanced guideline that, among other things, exempted patients over 55 years old with
heart or cardiovascular problems.

Those who opposed the initial measure, like Dr.~Joseph Selby, a research director at Kaiser
Permanente, feel their position has been vindicated. They see a straightforward lesson from the
episode: faulty guidelines can pose risks to patients, particularly when linked to doctors' pay,
which is an idea under consideration in Washington.

If such treatment guidelines are created ``in the absence of good evidence,'' Dr.~Selby said, ``you
risk perversion or even harm.''

\section{Democrats Seem Set to Go Alone on a Health Bill}

\lettrine{G}{iven} hardening Republican opposition to Congressional health care proposals, Democrats
now say they see little chance of the minority's cooperation in approving any overhaul, and are
increasingly focused on drawing support for a final plan from within their own ranks.

Top Democrats said Tuesday that their go-it-alone view was being shaped by what they saw as
Republicans' purposely strident tone against health care legislation during this month's
Congressional recess, as well as remarks by leading Republicans that current proposals were flawed
beyond repair.

Rahm Emanuel, the White House chief of staff, said the heated opposition was evidence that
Republicans had made a political calculation to draw a line against any health care changes, the
latest in a string of major administration proposals that Republicans have opposed.

``The Republican leadership,'' Mr.~Emanuel said, ``has made a strategic decision that defeating
President Obama's health care proposal is more important for their political goals than solving the
health insurance problems that Americans face every day.''

The Democratic shift may not make producing a final bill much easier. The party must still reconcile
the views of moderate and conservative Democrats worried about the cost and scope of the legislation
with those of more liberal lawmakers determined to win a government-run insurance option to compete
with private insurers.

On the other hand, such a change could alter the dynamic of talks surrounding health care
legislation, and even change the substance of a final bill. With no need to negotiate with
Republicans, Democrats might be better able to move more quickly, relying on their large majorities
in both houses.

Democratic senators might feel more empowered, for example, to define the authority of the nonprofit
insurance cooperatives that are emerging as an alternative to a public insurance plan.

Republicans have used the Congressional break to dig in hard against the overhaul outline drawn by
Democrats. The Senate's No.~2 Republican, Jon Kyl of Arizona, is the latest to weigh in strongly,
saying Tuesday that the public response lawmakers were seeing over the summer break should persuade
Democrats to scrap their approach and start over.

``I think it is safe to say there are a huge number of big issues that people have,'' Mr.~Kyl told
reporters in a conference call from Arizona. ``There is no way that Republicans are going to support
a trillion-dollar-plus bill.''

The White House has also interpreted critical comments by Senator Charles E.~Grassley of Iowa, the
top Republican negotiator in a crucial Finance Committee effort to reach a bipartisan compromise, as
a sign that there is little hope of reaching a deal politically acceptable to both parties.

Mr.~Grassley, who is facing the possibility of a Republican primary challenge next year, has gotten
an earful in traveling around his home state. At one gathering last week, in a city park in the
central Iowa town of Adel, a man rose from the crowd and urged him to ``stand up and fight'' the
Democratic plans. If he does not, the man yelled, ``we will vote you out!''

The White House, carefully following Mr.~Grassley's activities, presumed he was no longer interested
in negotiating with Democrats after he initially made no effort to debunk misinformation that the
legislation could lead to ``death panels'' empowered to judge who would receive care.

Citing a packed schedule, Mr.~Grassley has also put off plans for the bipartisan group of Finance
Committee negotiators to meet in either Iowa or Maine, the home of another Republican member of the
group, Senator Olympia J.~Snowe, before Congress resumes.

Further, Mr.~Grassley said this week that he would vote against a bill unless it had wide support
from Republicans, even if it included all the provisions he wanted. ``I am negotiating for
Republicans,'' he told MSNBC.

In an interview on Tuesday, Mr.~Grassley said he had simply been repeating earlier comments that he
would not support a measure that did not have significant Republican support. He said that raucous
town-hall-style meetings might have made the job of reaching a compromise harder, but that he had
not given up.

``It may be more difficult than it was before,'' he said. ``I am intent on talking. I am intent on
seeing what we can do.''

Administration officials, who maintain that Republicans are badly mischaracterizing the legislation
that has emerged from three House committees and the Senate health committee, said they had hoped to
achieve some level of bipartisan support. But they are becoming increasingly convinced that they
will instead have to navigate the complicated politics among varying Democratic factions.

The officials said the White House hoped to make the case to the American people that it was
Republicans who had abandoned the effort at bipartisanship. Republicans countered by saying that
they simply opposed the legislation and that the public outcry had validated their view and
solidified their opposition.

This week's careful administration maneuvering on whether a public insurance option was an essential
element of any final bill was seemingly part of the new White House effort to find consensus among
Democrats, since the public plan has been resisted by moderate and conservative Democrats who could
be crucial to winning the votes for passage if no Republicans are on board.

For the second time in two days, Mr.~Obama did not mention health care on Tuesday, a marked
departure from the aggressive public relations campaign he mounted in July and early August. The
White House is striving to stay out of the fray, aides said, until the president can get away on
vacation this weekend.

Even as the administration showed some flexibility, angering liberal Democrats who consider a public
plan essential, Republicans turned their attacks from the public option to the health care
cooperative idea being promoted by some Senate Democrats.

In what Democrats regarded as further evidence that Republicans were not serious about negotiating,
Mr.~Kyl and Representative Eric Cantor of Virginia, the second-ranking House Republican, described a
co-op as a public option carrying another name.

The continuing opposition was noted Tuesday by Robert Gibbs, the White House spokesman, who said of
Republicans that at best ``only a handful seem interested in the type of comprehensive reform that
so many people believe is necessary to ensure the principles and the goals that the president has
laid out.''

\section{In North Korea, Clinton Helped Unveil a Mystery }

\lettrine{T}{he} visit was arranged under a veil of secrecy with the help of an unlikely broker: a
high-level American intelligence officer who spent much of his career trying to unlock the mysteries
of North Korea.

When former President Bill Clinton landed in Pyongyang on Aug.~4 to win the release of two
imprisoned American journalists, senior officials said he met an unexpectedly spry North Korean
leader, Kim Jong-il, who welcomed him with a long dinner that night, even proposing to stay up
afterward.

Mr.~Kim was flanked by two longtime aides, and he gave no hint that North Korea was in the throes of
a succession struggle, despite the widespread questions over how long he might live.

Mr.~Clinton and the Obama administration were determined not to extend a public-relations coup to
Mr.~Kim, who expressed a desire for better relations with the United States. But the visit is
already setting off ripples that could change the tenor of the relationship between the United
States and North Korea.

On Wednesday, diplomats from North Korea plan to visit Gov. Bill Richardson of New Mexico with an
undisclosed agenda, a senior administration official said Tuesday. Like Mr.~Clinton, Mr.~Richardson
has traveled to Pyongyang to negotiate the release of Americans held there, in his case in the
mid-1990s.

The White House approved the visit, which the official said did not signal any movement toward the
resumption of official talks with North Korea and the United States. But the meeting, which he said
the North Koreans requested, comes on the heels of conciliatory\cite{conciliatory} gestures toward
South Korea, and suggests a concerted effort on the part of the North.

Mr.~Clinton steered\cite{steer} clear of broader issues during his humanitarian mission, officials
said. Indeed, he did not even ask to see Mr.~Kim, requesting instead a meeting with ``an appropriate
official.'' To help the former president in case something went awry, the White House recommended
John Podesta, an adviser to both Mr.~Clinton and President Obama, join his delegation.

And to ensure he would not leave empty-handed, Mr.~Clinton asked that a member of his entourage meet
with the journalists, Laura Ling and Euna Lee, shortly after he landed to make sure they were safe,
said a senior administration official, who had been briefed on the visit.

For all the billions of dollars a year that the United States spends on intelligence gathering about
mysterious and unpredictable countries like North Korea, it took just 20 hours on the ground in
Pyongyang by a former president to give the Obama administration its first detailed look into a
nuclear-armed nation that looms as one of its greatest foreign threats.

The details about Mr.~Clinton's visit came from interviews with multiple government officials, who
spoke on the condition of anonymity because they were not authorized to speak publicly.

On Tuesday, Mr.~Clinton went to the White House to brief Mr.~Obama and his top aides about the trip.
Even before the 40-minute session in the Situation Room, Mr.~Clinton had spoken to the president by
phone and briefed his national security adviser, Gen. James L.~Jones.

But the meeting was rich in symbolism, and the president invited Mr.~Clinton to the Oval Office to
talk further. The White House said little about what the men discussed, beyond noting that Mr.~Obama
had wanted to thank Mr.~Clinton for winning the release of Ms.~Ling and Ms.~Lee.

The role of the intelligence officer, Joseph R.~DeTrani, in arranging the visit, has not previously
been reported. Mr.~DeTrani is the government's senior officer responsible for collecting and
analyzing intelligence on North Korea. His efforts to pave the way for Mr.~Clinton's visit offer a
glimpse into how the administration has been forced to use unorthodox methods to overcome the lack
of formal communications.

During the Bush administration, when the United States was in still in talks with North Korea, the
White House did not use intelligence officers for these purposes, an official familiar with the
talks said. Indeed, before taking the job of North Korea mission manager in the Office of the
Director of National Intelligence in 2006, Mr.~DeTrani served as the special envoy to the six-party
talks with North Korea.

More than anything else, senior officials said, Mr.~Clinton's visit served to clear up some of the
shadows surrounding Kim Jong-il's health. After suffering a stroke last year, he looked frail in
photos, spurring questions about who might replace him.

Those questions have not gone away, officials said, but they may recede\cite{recede} a bit after
Mr.~Clinton's visit. So, too, may the speculation about internal battles, given the apparent good
standing of Kim Kye-gwan, the chief nuclear negotiator, and another foreign policy official, Kang
Sok-ju, who also took part in the meetings.

The former president did not engage in a substantive discussion about North Korea's nuclear program.
Nor did the North Korean leader give Mr.~Clinton any indication that his nation would relinquish its
nuclear ambitions -- a condition the United States has set for resuming negotiations, officials
said.

During his one-hour meeting, officials said, Mr.~Clinton advised the North Korean that he could win
favor with South Korea and Japan by resolving cases of their citizens who had been abducted by North
Korea. The dinner, which lasted over two hours, was ``chitchat,'' the official said. ``It was not
substantive.''

North Korea has sent other conciliatory signals. Kim Jong-il met last week with the head of a South
Korean conglomerate and agreed to restart several tourism ventures, which allow people from the
South to visit the North. North Korea said it would also allow reunions of Korean families divided
by the border. Also, Yonhap, the South Korean news agency, said Wednesday the North would send a
delegation to the funeral of Kim Dae-jung, a former president of the South.

It is not clear whether the overtures represent a change of heart or a growing desperation for
money, as the North comes under increasingly strict United Nations sanctions. The White House says
it is determined not to ease the pressure. In a deliberate bit of timing, it dispatched a senior
diplomat, Philip S.~Goldberg, to Asia on Tuesday to discuss ways to enforce the sanctions.

Still, officials and analysts said, Mr.~Clinton's visit was valuable, largely because North Korea is
so opaque. Victor Cha, a top North Korea adviser in the Bush administration, said, ``The Clinton
trip has got a lot of people rethinking and reassessing.'' 

\section{As Japan's Mediums Die, Ancient Tradition Fades}

\lettrine{I}{ts} name means the Mountain of Horror, which seems an apt description for this sacred
Buddhist\cite{Buddhist} site inside the crater\cite{crater} of a dormant volcano. The weather-beaten
temple here is surrounded by a lifeless lake and a wasteland of naked rock reeking of sulfur that
conjures images of Buddhist hell.

But during the mountain's twice annual religious festivals, visitors come by the busload to line up
before a row of small tents in a corner of the temple. Within are the ``itako'' -- elderly, often
blind women who hold s\'eance-like ceremonies that customers hope will allow them to commune with
spirits of the dead.

These spiritual mediums seem out of place in a hyper-modern nation better known for bullet trains
and hybrid cars. Found only in peripheral areas like this volcano on the far northern tip of Japan's
main island, and only dimly known to most Japanese, the itako are among the last remaining adherents
to ancient shamanistic beliefs that predate Buddhism and modern forms of Shintoism, Japan's two main
religions, historians say.

They have survived government efforts to stamp them out, as well as the continuing disdain of many
Japanese, who look down on them as charlatans who trade in superstition. Even the deputy abbot at
Bodai-ji, Mount Osore's temple, said the itako were not connected to the temple, which he said only
tolerates their presence.

Now, however, even these last remaining itako are vanishing. Only four graying itako appeared at Mt.
Osore's weeklong summer festival this year, three having died of old age in the last year. Worse,
the only practicing medium younger than retirement age -- 40-year-old Keiko Himukai, known among
believers as the last itako -- stopped coming this year for health reasons.

``We can see a very ancient flame dying out before our eyes,'' Ms.~Himukai said in a separate
interview. ``But traditions have to change with the times.''

Junichi Tonosaki, a historian in the prefectural museum in Aomori, where Mount Osore is located,
said the number of itako had fallen from about 20 a decade ago. He said they began gathering at
Mount Osore in the last century as their numbers began to dwindle, to make it easier for customers
to find them. The volcano's 1,200-year-old temple is believed by many here to be a gathering point
for souls of the dead before Buddhist reincarnation.

Mr.~Tonosaki and other historians say itako and other shamanistic mediums were common across Japan
in medieval times, when this was often the only occupation available for the blind. But they were
suppressed in the late 19th century, as Japan built a modern nation. In recent times, they have
survived only on the geographic margins, in rural northern Japan and on the southern island of
Okinawa.

Shojiro Kurokawa, 82, can remember as a child in the 1930s when residents of his and other nearby
villages would trek to the temple to hold weeklong festivals of all-night dancing, singing and
s\'eances. In those days, he said, there were more than 100 itako.

``This is an era when children ignore their parents and forget about the dead,'' said Mr.~Kurokawa,
who runs an inn near the temple that caters to visitors of the spiritual mediums.

Now, there are so few itako that visitors routinely wait in line for several hours to see one. Itako
charge 3,000 yen, or about \$30, for each spirit called in a roughly 10-minute ceremony.

One family of three came from Tokyo, a day's drive away. Masako Toyama, 68, said she came to speak
with her husband, who died suddenly last summer of cancer. She, her son and his wife said it was
their first visit to Mount Osore. She said she knew of the itako from growing up in northern Japan.

``They were a scary but also soothing presence,'' Ms.~Toyama said. ``Japanese still need this sort
of emotional support.''

``I wanted to give the itako a try, to see if this is real,'' said the son, Shinji Toyama, a
41-year-old salaryman at a medical testing firm.

When the Toyamas' turn finally came after six and a half hours of waiting, they seemed almost taken
aback by the itako's modern appearance, in pink-tinted glasses and a flower-patterned shirt.

The itako, Setsu Aoyama, began by lifting a long strand of dark beads and began a short chant: ``I
call the spirit who died on July 11,'' the date Ms.~Toyama's husband, Shigeto, died.

In the same rhythmic cadence, swaying her head with eyes shut, she assumed the voice of Shigeto. ``I
didn't go to a doctor soon enough,'' she intoned. ``Men don't listen to things like that.''

Ms.~Himukai, the 40-year-old itako, says she enters a trance in which she feels the presence of the
spirit and its mood, which she expresses in her own words. She said she decided to begin the
three-year period of study to become a spiritual medium as a teenager, after an itako near her rural
village cured her of an ailment that doctors could not fix.

She said she felt guilty about not going to Mount Osore. However, she said she may no longer be able
to attend because of health problems, including a chronic stomach ailment. Instead, she said she
wanted to write a book or make a movie about the itako.

``The end can also be the beginning of something new,'' said Ms.~Himukai, who wore a plain gray suit
with pants and spoke in a whisper.

After the ceremony, the Toyamas had mixed reactions. The widow said her heart had been put at ease.
But the younger generation was less convinced.

``I didn't feel like it was really my father in front of me,'' said Shinji, the son. But he said he
wanted to come again next year, to try a different itako. ``Maybe we just had beginner's bad luck by
choosing the wrong one.''

\section{Bush Official, in Book, Tells of Pressure on '04 Vote}

\lettrine{T}{om} Ridge, the first secretary of homeland security, asserts in a new book that he was
pressured by top advisers to President George W.~Bush to raise the national threat level just before
the 2004 election in what he suspected was an effort to influence the vote.

After Osama bin Laden released a threatening videotape four days before the election, Attorney
General John Ashcroft and Defense Secretary Donald H.~Rumsfeld pushed Mr.~Ridge to elevate the
public threat posture but he refused, according to the book. Mr.~Ridge calls it a ``dramatic and
inconceivable'' event that ``proved most troublesome'' and reinforced his decision to resign.

The provocative allegation provides fresh ammunition for critics who have accused the Bush
administration of politicizing national security. Mr.~Bush and his Democratic challenger, Senator
John Kerry of Massachusetts, were locked in a tight race heading into that final weekend, and some
analysts concluded that even without a higher threat level, the bin Laden tape helped the president
win re-election by reminding voters of the danger of Al Qaeda.

Keith M.~Urbahn, a spokesman for Mr.~Rumsfeld, said the defense secretary supported letting the
public know if intelligence agencies believed there was a greater threat, and pointed to a variety
of chilling Qaeda warnings in those days, including one tape vowing that ``the streets of America
will run red with blood.''

``Given those facts,'' Mr.~Urbahn said, ``it would seem reasonable for senior administration
officials to discuss the threat level. Indeed, it would have been irresponsible had that discussion
not taken place.''

Mr.~Urbahn said ``the storyline advanced by his publisher seemingly to sell copies of the book is
nonsense.''

Mr.~Ashcroft could not be reached for comment. But Mark Corallo, who was his spokesman at the
Justice Department, dismissed Mr.~Ridge's account. ``Didn't happen,'' he said. ``Now would be a good
time for Mr.~Ridge to use his emergency duct tape.''

Frances Fragos Townsend, who was Mr.~Bush's homeland security adviser, said that ``there was a
fulsome debate'' about the threat level but that ``the politics of it were not ever a factor.''

Mr.~Ridge's book, called ``The Test of Our Times'' and due out Sept.~1 from Thomas Dunne Books, is
the latest by a Bush adviser to disclose internal disagreements and establish distance from an
unpopular administration. Mr.~Ridge complains that he was never invited to National Security Council
meetings, that Mr.~Rumsfeld would rarely meet with him and that the White House pressured him to
include a justification for the Iraq war in a speech.

He also writes that he lobbied unsuccessfully before Hurricane Katrina in 2005 to replace Michael
D.~Brown as head of the Federal Emergency Management Agency and that the White House killed his
proposal to open a homeland security regional office in New Orleans.

The most sensational assertion was the pre-election debate in 2004 about the threat level, first
reported by U.S.~News \& World Report. Mr.~Ridge writes that the bin Laden tape alone did not
justify a change in the nation's security posture but describes ``a vigorous, some might say
dramatic, discussion'' on Oct.~30 to do so.

``There was absolutely no support for that position within our department. None,'' he writes. ``I
wondered, 'Is this about security or politics?' Post-election analysis demonstrated a significant
increase in the president's approval rating in the days after the raising of the threat level.''

Mr.~Ridge provides no evidence that politics motivated the discussion. Until now, he has denied
politics played a role in threat levels. Asked by Eric Lichtblau of The New York Times if politics
ever influenced decisions on threat warnings, he volunteered to take a lie-detector test. ``Wire me
up,'' Mr.~Ridge said, according to Mr.~Lichtblau's book, ``Bush's Law.'' ``Not a chance. Politics
played no part.''

\section{Rise of the Super-Rich Hits a Sobering Wall}

\lettrine{T}{he} rich have been getting richer for so long that the trend has come to seem almost
permanent.

They began to pull away from everyone else in the 1970s. By 2006, income was more concentrated at
the top than it had been since the late 1920s. The recent news about resurgent Wall Street pay has
seemed to suggest that not even the Great Recession could reverse the rise in income inequality.

But economists say -- and data is beginning to show -- that a significant change may in fact be
under way. The rich, as a group, are no longer getting richer. Over the last two years, they have
become poorer. And many may not return to their old levels of wealth and income anytime soon.

For every investment banker whose pay has recovered to its prerecession levels, there are several
who have lost their jobs -- as well as many wealthy investors who have lost millions. As a result,
economists and other analysts say, a 30-year period in which the super-rich became both wealthier
and more numerous may now be ending.

The relative struggles of the rich may elicit little sympathy from less well-off families who are
dealing with the effects of the worst recession in a generation. But the change does raise several
broader economic questions. Among them is whether harder times for the rich will ultimately benefit
the middle class and the poor, given that the huge recent increase in top incomes coincided with
slow income growth for almost every other group. In blunter terms, the question is whether the
better metaphor for the economy is a rising tide that can lift all boats -- or a zero-sum game.

Just how much poorer the rich will become remains unclear. It will be determined by, among other
things, whether the stock market continues its recent rally and what new laws Congress passes in the
wake of the financial crisis. At the very least, though, the rich seem unlikely to return to the
trajectory they were on.

Last year, the number of Americans with a net worth of at least \$30 million dropped 24 percent,
according to CapGemini and Merrill Lynch Wealth Management. Monthly income from stock dividends,
which is concentrated among the affluent, has fallen more than 20 percent since last summer, the
biggest such decline since the government began keeping records in 1959.

Bill Gates, Warren E.~Buffett, the heirs to the Wal-Mart Stores fortune and the founders of Google
each lost billions last year, according to Forbes magazine. In one stark example, John McAfee, an
entrepreneur who founded the antivirus software company that bears his name, is now worth about \$4
million, from a peak of more than \$100 million. Mr.~McAfee will soon auction off his last big
property because he needs cash to pay his bills after having been caught off guard by the
simultaneous crash in real estate and stocks.

``I had no clue,'' he said, ``that there would be this tandem collapse.''

Some of the clearest signs of the reversal of fortunes can be found in data on spending by the
wealthy. An index that tracks the price of art, the Mei Moses index, has dropped 32 percent in the
last six months. The New York Yankees failed to sell many of the most expensive tickets in their new
stadium and had to drop the price. In one ZIP code in Vail, Colo., only five homes sold for more
than \$2 million in the first half of this year, down from 34 in the first half of 2007, according
to MDA Dataquick. In Bronxville, an affluent New York suburb, the decline was to two, from 17,
according to Coldwell Banker Residential Brokerage.

``We had a period of roughly 50 years, from 1929 to 1979, when the income distribution tended to
flatten,'' said Neal Soss, the chief economist at Credit Suisse. ``Since the early '80s, incomes
have tended to get less equal. And I think we've entered a phase now where society will move to a
more equal distribution.''

No More '50s and '60s

Few economists expect the country to return to the relatively flat income distribution of the 1950s
and 1960s. Indeed, they say that inequality is likely to remain significantly greater than it was
for most of the 20th century. The Obama administration has not proposed completely rewriting the
rules for Wall Street or raising the top income-tax rate to anywhere near 70 percent, its level as
recently as 1980. Market forces that have increased inequality, like globalization, are also not
going away.

But economists say that the rich will probably not recover their losses immediately, as they did in
the wake of the dot-com crash earlier this decade. That quick recovery came courtesy of a new bubble
in stocks, which in 2007 were more expensive by some measures than they had been at any other point
save the bull markets of the 1920s or 1990s. This time, analysts say, Wall Street seems unlikely to
return soon to the extreme levels of borrowing that made such a bubble possible.

Any major shift in the financial status of the rich could have big implications. A drop in their
income and wealth would complicate life for elite universities, museums and other institutions that
received lavish donations in recent decades. Governments -- federal and state -- could struggle,
too, because they rely heavily on the taxes paid by the affluent.

Perhaps the broadest question is what a hit to the wealthy would mean for the middle class and the
poor. The best-known data on the rich comes from an analysis of Internal Revenue Service returns by
Thomas Piketty and Emmanuel Saez, two economists. Their work shows that in the late 1970s, the
cutoff to qualify for the highest-earning one ten-thousandth of households was roughly \$2 million,
in inflation-adjusted, pretax terms. By 2007, it had jumped to \$11.5 million.

The gains for the merely affluent were also big, if not quite huge. The cutoff to be in the top 1
percent doubled since the late 1970s, to roughly \$400,000.

By contrast, pay at the median -- which was about \$50,000 in 2007 -- rose less than 20 percent,
Census data shows. Near the bottom of the income distribution, the increase was about 12 percent.

Some economists say they believe that the contrasting trends are unrelated. If anything, these
economists say, any problems the wealthy have will trickle down, in the form of less charitable
giving and less consumer spending. Over the last century, the worst years for the rich were the
early 1930s, the heart of the Great Depression.

Other economists say the recent explosion of incomes at the top did hurt everyone else, by
concentrating economic and political power among a relatively small group.

``I think incredibly high incomes can have a pernicious effect on the polity and the economy,'' said
Lawrence Katz, a Harvard economist. Much of the growth of high-end incomes stemmed from market
forces, like technological innovation, Mr.~Katz said. But a significant amount also stemmed from the
wealthy's newfound ability to win favorable government contracts, low tax rates and weak financial
regulation, he added.

The I.R.S. has not yet released its data for 2008 or 2009. But Mr.~Saez, a professor at the
University of California, Berkeley, said he believed that the rich had become poorer. Asked to
speculate where the cutoff for the top one ten-thousandth of households was now, he said from \$6
million to \$8 million.

For the number to return to \$11 million quickly, he said, would probably require a large financial
bubble.

Making More Money

The United States economy experienced two such bubbles in recent years -- one in stocks, the other
in real estate -- and both helped the rich become richer. Mr.~McAfee, whose tattoos and tinted hair
suggest an independent streak, is an extreme but telling example. For two decades, at almost every
step of his career, he figured out a way to make more money.

In the late 1980s, he founded McAfee Associates, the antivirus software company. It gave away its
software, unlike its rivals, but charged fees to those who wanted any kind of technical support.
That decision helped make it a huge success. The company went public in 1992, in the early years of
one of biggest stock market booms in history.

But Mr.~McAfee is, by his own description, an atypical businessman -- easily bored and given to
serial obsessions. As a young man, he traveled through Mexico, India and Nepal and, more recently,
he wrote a book called, ``Into the Heart of Truth: The Spirit of Relational Yoga.'' Two years after
McAfee Associates went public, he was bored again.

So he sold his remaining stake, bringing his gains to about \$100 million. In the coming years, he
started new projects and made more investments. Almost inevitably, they paid off.

``History told me that you just keep working, and it is easy to make more money,'' he said, sitting
in the kitchen of his adobe-style house in the southwest corner of New Mexico. With low tax rates,
he added, the rich could keep much of what they made.

One of the starkest patterns in the data on inequality is the extent to which the incomes of the
very rich are tied to the stock market. They have risen most rapidly during the biggest bull
markets: in the 1920s and the 20 years starting in 1987.

``We are coming from an abnormal period where a tremendous amount of wealth was created largely by
selling assets back and forth,'' said Mohamed A.~El-Erian, chief executive of Pimco, one of the
country's largest bond traders, and the former manager of Harvard's endowment.

Some of this wealth was based on real economic gains, like those from the computer revolution. But
much of it was not, Mr.~El-Erian said. ``You had wealth creation that could not be tied to the
underlying economy,'' he added, ``and the benefits were very skewed: they went to the assets of the
rich. It was financial engineering.''

But if the rich have done well in bubbles, they have taken enormous hits to their wealth during
busts. A recent study by two Northwestern University economists found that the incomes of the
affluent tend to fall more, in percentage terms, in recessions than the incomes of the middle class.
The incomes of the very affluent -- the top one ten-thousandth -- fall the most.

Over the last several years, Mr.~McAfee began to put a large chunk of his fortune into real estate,
often in remote locations. He bought the house in New Mexico as a playground for himself and fellow
aerotrekkers, people who fly unlicensed, open-cockpit planes. On a 157-acre spread, he built a
general store, a 35-seat movie theater and a cafe, and he bought vintage cars for his visitors to
use.

He continued to invest in financial markets, sometimes borrowing money to increase the potential
returns. He typically chose his investments based on suggestions from his financial advisers. One of
their recommendations was to put millions of dollars into bonds tied to Lehman Brothers.

For a while, Mr.~McAfee's good run, like that of many of the American wealthy, seemed to continue.
In the wake of the dot-com crash, stocks started rising again, while house prices just continued to
rise. Outside's Go magazine and National Geographic Adventure ran articles on his New Mexico
property, leading to him to believe that ``this was the hottest property on the planet,'' he said.

But then things began to change.

In 2007, Mr.~McAfee sold a 10,000-square-foot home in Colorado with a view of Pike's Peak. He had
spent \$25 million to buy the property and build the house. He received \$5.7 million for it. When
Lehman collapsed last fall, its bonds became virtually worthless. Mr.~McAfee's stock investments
cost him millions more.

One day, he realized, as he said, ``Whoa, my cash is gone.''

His remaining net worth of about \$4 million makes him vastly wealthier than most Americans, of
course. But he has nonetheless found himself needing cash and desperately trying to reduce his
monthly expenses.

He has sold a 10-passenger Cessna jet and now flies coach. This week his oceanfront estate in Hawaii
sold for \$1.5 million, with only a handful of bidders at the auction. He plans to spend much of his
time in Belize, in part because of more favorable taxes there.

Next week, his New Mexico property will be the subject of a no-floor auction, meaning that
Mr.~McAfee has promised to accept the top bid, no matter how low it is.

``I am trying to face up to the reality here that the auction may bring next to nothing,'' he said.

In the past, when his stock investments did poorly, he sold real estate and replenished his cash.
This time, that has not been an option.

Stock Market Mystery

The possibility that the stock market will quickly recover from its collapse, as it did earlier this
decade, is perhaps the biggest uncertainty about the financial condition of the wealthy. Since
March, the Standard \& Poor's 500-stock index has risen 49 percent.

Yet Wall Street still has a long way to go before reaching its previous peaks. The S.\& P.~500
remains 35 percent below its 2007 high. Aggregate compensation for the financial sector fell 14
percent from 2007 to 2008, according to the Securities Industry and Financial Markets Association --
far less than profits or revenue fell, but a decline nonetheless.

``The difference this time,'' predicted Byron R.~Wein, a former chief investment strategist at
Morgan Stanley, who started working on Wall Street in 1965, ``is that the high-water mark that
people reached in 2007 is not going to be exceeded for a very long time.''

Without a financial bubble, there will simply be less money available for Wall Street to pay itself
or for corporate chief executives to pay themselves. Some companies -- like Goldman Sachs and
JPMorgan Chase, which face less competition now and have been helped by the government's attempts to
prop up credit markets -- will still hand out enormous paychecks. Over all, though, there will be
fewer such checks, analysts say. Roger Freeman, an analyst at Barclays Capital, said he thought that
overall Wall Street compensation would, at most, increase moderately over the next couple of years.

Beyond the stock market, government policy may have the biggest effect on top incomes. Mr.~Katz, the
Harvard economist, argues that without policy changes, top incomes may indeed approach their old
highs in the coming years. Historically, government policy, like the New Deal, has had more lasting
effects on the rich than financial busts, he said.

One looming policy issue today is what steps Congress and the administration will take to
re-regulate financial markets. A second issue is taxes.

In the three decades after World War II, when the incomes of the rich grew more slowly than those of
the middle class, the top marginal rate ranged from 70 to 91 percent. Mr.~Piketty, one of the
economists who analyzed the I.R.S. data, argues that these high rates did not affect merely post-tax
income. They also helped hold down the pretax incomes of the wealthy, he says, by giving them less
incentive to make many millions of dollars.

Since 1980, tax rates on the affluent have fallen more than rates on any other group; this year, the
top marginal rate is 35 percent. President Obama has proposed raising it to 39 percent and has said
he would consider a surtax on families making more than \$1 million a year, which could push the top
rate above 40 percent.

What any policy changes will mean for the nonwealthy remains unclear. There have certainly been
periods when the rich, the middle class and the poor all have done well (like the late 1990s), as
well as periods when all have done poorly (like the last year). For much of the 1950s, '60s and
'70s, both the middle class and the wealthy received raises that outpaced inflation.

Yet there is also a reason to think that the incomes of the wealthy could potentially have a bigger
impact on others than in the past: as a share of the economy, they are vastly larger than they once
were.

In 2007, the top one ten-thousandth of households took home 6 percent of the nation's income, up
from 0.9 percent in 1977. It was the highest such level since at least 1913, the first year for
which the I.R.S. has data.

The top 1 percent of earners took home 23.5 percent of income, up from 9 percent three decades
earlier.

\section{Will the Hard-Core Starbucks Customer Pay More? The Chain Plans to Find Out}

\lettrine{A}{s} the recession wears on and fewer people are splurging at Starbucks, the coffee
chain's response is to raise prices. On Thursday, Starbucks stores in several cities started
charging up to 30 cents more for some specialty beverages, though the company is charging less for
some basic drinks.

The price adjustments will be made at stores nationwide in coming months.

The move comes as Starbucks toes a tricky line between remaining a premium brand while retaining
price-sensitive customers who can buy cheaper coffee at other shops. This summer, Starbucks has been
fending off competition from McDonald's, which has introduced a new line of lower-priced espresso
drinks that have proved popular.

The price of sugary Starbucks concoctions with several ingredients, like Frappuccinos and caramel
macchiatos, will increase an average 10 cents to 15 cents, but in some cases as much as 30 cents, or
about 8 percent, said Valerie O'Neil, a Starbucks spokeswoman.

The price of the most popular beverages -- 12-ounce lattes and brewed coffees -- will decrease an
average of 5 cents to 15 cents. This is the first time Starbucks has lowered prices, she said. It
last increased the price of drinks by 5 cents in 2007.

``The goal of it is really to continue to find ways to balance the value we're providing for our
customers with the business responsibilities,'' Ms.~O'Neil said. In the three months ending June 28,
sales at Starbucks stores open at least a year fell 5 percent from the same period the year before.
Starbucks is also experimenting with other ways to make more money, like selling alcoholic beverages
in a few stores in Seattle.

Starbucks is also trying to heighten its appeal with a redesign of some of its menus so they more
prominently feature less expensive brewed coffee, and it is running a summer promotion for iced
drinks.

Raising prices in the middle of a big downturn in consumer spending might not be such a crazy idea,
said J.~Miguel Villas-Boas, a marketing professor who studies pricing and consumer choice at the
Haas School of Business of the University of California, Berkeley. Though it is paradoxical, it is a
strategy premium companies often take, he said.

``Given that McDonald's is capturing some of the consumers less interested in the premium that
Starbucks offers, then the consumers left out for Starbucks are the consumers willing to pay more,
so Starbucks says, 'Let's charge them,' '' he said.

Connie Williams, 57, a saleswoman in Los Angeles, is one of those customers. She said the price
changes would not affect her daily habit of buying a 12-ounce nonfat cappuccino ``with a little
shake of chocolate'' at Starbucks. ``I'm hooked,'' she said. ``It's like asking a cocaine addict,
'If I raise my prices, are you going to buy less?' ''

But some customers said the price increase would make them turn away. Elise Walls, 24, a student in
San Francisco, visits Starbucks nearly every day for a 20-ounce green ginger tea. On Thursday, she
noticed it cost 5 cents more. ``I was like, 'Wait a minute,' because it's already expensive
anyway,'' she said. She plans to start buying boxes of tea bags and making it herself.

The sugary, creamy drinks that will now be more expensive at Starbucks are the same type of beverage
that McDonald's has been marketing in a huge advertising campaign.

Starbucks is safe raising the prices of specialty drinks because they are where the company best
differentiates itself, said Kenneth Davids, editor of Coffee Review.

Mr.~Davids recently published a taste test comparing espresso drinks at Starbucks and McDonald's.
The difference between standard lattes was subtle, he said, but ``what surprised me was how dramatic
the difference was when it got to the additives,'' like syrup and whipped cream.

Starbucks customers get what they pay for, he said. ``With McDonald's, it was just terrible stuff
that I don't think anybody except maybe a 5-year-old might like.''

\section{Obama Insists Health Plan Will Pass}

\lettrine{P}{resident} Obama said Thursday that he ``would love to have more Republicans engaged''
in the health care negotiations under way in Congress, but expressed doubt that a bipartisan
compromise could be reached because he suspected the party's leadership was intent on defeating his
signature domestic priority.

``I think early on, a decision was made by the Republican leadership that said, 'Look, let's not
give him a victory,' '' Mr.~Obama said, accusing his rivals of trying to reprise the efforts that
went into derailing the Clinton health care agenda.

``There are some folks,'' he added, ``who are taking a page out of that playbook.''

But one lawmaker who participated Thursday night in a conference call among six Republican and
Democratic senators trying to reach bipartisan agreement on a health care plan said the senators
agreed to keep working together and asked staff members to focus on making health coverage more
affordable while holding down medical costs.

``We are hanging in there,'' said Senator Kent Conrad, Democrat of North Dakota, after the telephone
conference of almost one hour and 40 minutes. ``Is it easy? No.~But I tell you, I have never seen a
group with a greater sense of purpose than this group.''

As the president prepared to leave for vacation on Friday, he sought to allay concerns about his
health care agenda from both sides of the political spectrum, hoping to clear up what he called ``a
lot of confusion.'' He appeared on a conservative-leaning radio talk show and later rallied his own
network of supporters in a coordinated White House effort to push forward on remaking the health
care system.

``You know,'' Mr.~Obama said, ``passing a big bill like this is always messy.''

The president did not entirely close the door to reaching a bipartisan deal, saying he was ``happy
to make sensible compromises.'' But he added that Democrats would not ``give up on the core
principle that Americans who don't have health insurance should get it.''

Mr.~Obama's comments came hours before six members of the Senate Finance Committee were to consult
in a conference call on the state of their negotiations toward a bipartisan health care agreement.
The leaders of the talks, Senator Max Baucus, Democrat of Montana and chairman of the panel, and the
senior Republican, Senator Charles E.~Grassley of Iowa, have professed some confidence this week
that a deal could be struck, but the tone from Mr.~Obama and other Democrats suggested skepticism.

As the bipartisan Finance Committee negotiators prepared to discuss a bill that almost certainly
would not include a public insurance plan, Speaker Nancy Pelosi told reporters in San Francisco on
Thursday that legislation without such a public plan could not win House approval.

Those involved in the Senate negotiations continued to express confidence that they could ultimately
reach an agreement when Congress returned next month, though they said any legislation produced
would have to be scaled back from measures that have cleared other committees in the House and the
Senate.

Mr.~Obama talked by telephone with two of the negotiators, Mr.~Conrad and Senator Olympia J.~Snowe,
Republican of Maine. Ms.~Snowe is still seen by administration officials and Senate Democrats as a
possible ally on the health care legislation, even if other Republicans are unwilling to back a
bill.

Republican Congressional leaders disputed the president's characterization that they had worked in a
coordinated effort to oppose the administration's health care agenda, saying they had legitimate
questions about the cost and scope of the plan.

``Americans are rightly skeptical about the administration's approach to overhauling everyone's
health care and about the more than \$1 trillion price tag,'' said Senator Mitch McConnell of
Kentucky, the Republican leader. ``And Americans are concerned about funding new government programs
through massive cuts to Medicare and taxes on small business, which is not the bipartisan approach
Americans expected.''

Republicans aside, Mr.~Obama had plenty to contend with in his own party.

The president took a handful of questions Thursday on the radio talk show, which originated from the
White House and was hosted by Michael Smerconish of Philadelphia. Mr.~Obama renewed his support for
a public insurance option, seeking to quell the uproar among some Democrats that he was abandoning
his position to make the plan more acceptable to moderates.

``The press got excited, and some folks on the left got a little excited,'' Mr.~Obama said,
referring to the news reports this week that the administration was leaning toward health care
cooperatives as an alternative to a government-run program to compete with private insurers.

Still, the president added that the public option was merely ``one choice of many choices,'' adding
that no one should be ``obligated to go into a public plan.''

Asked by one caller what elements needed to be included in a health care plan, Mr.~Obama listed four
points: reducing the cost of health care, protecting consumers from insurance abuses, providing
affordable coverage to uninsured Americans and not adding to the deficit.

Mr.~Smerconish, who attracts conservative-leaning listeners, endorsed Mr.~Obama's candidacy last
year and moderated over a session of polite questions that belied the intense anger that had erupted
over the airwaves this month. Mr.~Obama said he had gone on the show to dispel myths that he said
had overshadowed the debate, among them that the health plan would provide coverage to illegal
immigrants.

One caller, who said that he had supported Mr.~Obama and worked for his campaign, said he sensed
that the president's ``knees are buckling a little bit'' on his pledge to overhaul the nation's
health care system.

``I'm confident that we're going to get it done,'' Mr.~Obama said.

Later, Mr.~Obama spoke to a gathering of Organizing for America, the group of activists who
supported his candidacy. Many members of the organization, along with other liberal groups, have
accused him of wavering on his pledge to include a public insurance option as a central part of his
health care plan, but no such concerns were expressed on Thursday as Mr.~Obama delivered a pep talk
to the group in a meeting at the Democratic National Committee that was broadcast on its Web site.

``The best offense against lies is the truth,'' said Mr.~Obama, barely mentioning the discord in his
own party as he urged his supporters to focus on the true opposition. ``All we can do is just keep
on pushing the truth.''

\section{The Daughter Deficit}

\lettrine{I}{n} the late 1970s, a Ph.D. student named Monica Das Gupta was conducting
anthropological\cite{anthropological} fieldwork in Haryana, a state in the north of India. She
observed something striking about families there: parents had a fervent\cite{fervent} preference for
male offspring. Women who had given birth to only daughters were desperate for sons and would keep
having children until they had one or two. Midwives were even paid less when a girl was born. ``It's
something you notice coming from outside,'' says Das Gupta, who today studies population and public
health in the World Bank's development research group. ``It just leaps out at you.''

Das Gupta saw that educated, independent-minded women shared this prejudice in Haryana, a state that
was one of India's richest and most developed. In fact, the bias against girls was far more
pronounced there than in the poorer region in the east of India where Das Gupta was from. She
decided to study the issue in Punjab, then India's richest state, which had a high rate of female
literacy and a high average age of marriage. There too the prejudice for sons flourished. Along with
Haryana, Punjab had the country's highest percentage of so-called missing girls -- those aborted,
killed as newborns or dead in their first few years from neglect.

Here was a puzzle: Development seemed to have not only failed to help many Indian girls but to have
made things worse.

It is rarely good to be female anywhere in the developing world today, but in India and China the
situation is dire\cite{dire}: in those countries, more than 1.5 million fewer girls are born each
year than demographics would predict, and more girls die before they turn 5 than would be expected.
(In China in 2007, there were 1.73 million births -- and a million missing girls.) Millions more
grow up stunted, physically and intellectually, because they are denied the health care and the
education that their brothers receive.

Among policymakers, the conventional wisdom is that such selective brutality toward girls can be
mitigated by two factors. One is development: surely the wealthier the home, the more educated the
parents, the more plugged in to the modern economy, the more a family will invest in its girls. The
other is focusing aid on women. The idea is that a mother who has more money, knowledge and
authority in the family will direct her resources toward all her children's health and education.
She will fight for her girls.

Yet these strategies -- though invaluable -- underestimate the complexity of the situation in
certain countries. To be sure, China and India are poor. But in both nations, girls are actually
more likely to be missing in richer areas than in poorer ones, and in cities than in rural areas.
Having more money, a better education and (in India) belonging to a higher caste all raise the
probability that a family will discriminate against its daughters. The bias against girls applies in
some of the wealthiest and best-educated nations in the world, including, in recent years, South
Korea, Taiwan and Singapore. It also holds among Indian immigrants in Britain and among Chinese,
Indian and South Korean immigrants in the United States. In the last few years, the percentage of
missing girls has been among the highest in the middle-income, high-education nations of the
Caucasus: Armenia, Azerbaijan and Georgia.

Nor does a rise in a woman's autonomy or power in the family necessarily counteract prejudice
against girls. Researchers at the International Food Policy Research Institute have found that while
increasing women's decision-making power would reduce discrimination against girls in some parts of
South Asia, it would make things worse in the north and west of India. ``When women's power is
increased,'' wrote Lisa C.~Smith and Elizabeth M.~Byron, ``they use it to favor boys.''

Why should this be? A clue lies in what Das Gupta uncovered in her research in Punjab in the 1980s.
At the time, it was assumed that parents in certain societies simply did not value girls. And in
important ways, this was true. But Das Gupta complicated this picture. She found that it was not
true that all daughters were mistreated equally. A firstborn daughter was not typically subjected to
inferior treatment; she was treated like her brothers. But a subsequent daughter born to an educated
mother was 2.36 times as likely to die before her fifth birthday as her siblings were to die before
theirs -- mainly because she was less likely to see a doctor. It turned out that a kind of economic
logic was at work: with a firstborn girl, families still had plenty of chances to have a boy; but
with each additional girl, the pressure to have a son increased. The effect of birth order that Das
Gupta discovered has now been confirmed in subsequent studies of missing girls.

What unites communities with historically high rates of discrimination against girls is a rigid
patriarchal culture that makes having a son a financial and social necessity. When a daughter grows
up and marries, she essentially becomes chattel\cite{chattel} in her husband's parents' home and has
very limited contact with her natal\cite{natal} family. Even if she earns a good living, it will be
of no help to her own parents in their old age. So for parents, investing in a daughter is truly, in
the Hindi expression, planting a seed in the neighbor's garden. Sons, by contrast, provide a kind of
social security. A family with only daughters will also likely lose its land when the father dies:
although women can legally inherit property, in areas of north India and China, they risk
ostracism\cite{ostracism} or even murder if they claim what is theirs. And sons are particularly
important to mothers, who acquire power and authority when they have married sons. Sons, according
to Chinese custom, are also needed to care for the souls of dead ancestors.

What Das Gupta discovered is that wealthier and more educated women face this same
imperative\cite{imperative} to have boys as uneducated poor women -- but they have smaller families,
thus increasing the felt urgency of each birth. In a family that expects to have seven children, the
birth of a girl is a disappointment; in a family that anticipates only two or three children, it is
a tragedy.

Thus development can worsen, not improve, traditional discrimination. This can happen in other ways
too. With the access it brings to cutting-edge technology, development can also offer more
sophisticated and easier options for exercising old-fashioned prejudice. In China and in the north
and west of India, for instance, the spread of ultrasound technology, which can inform parents of
the sex of their fetus, has turned a pool of missing girls into an ocean. The birth of girls has
long been avoided through infanticide, which is still practiced often in China. But there are even
more couples who would abort a pregnancy than would kill a newborn. Ultrasound has been advertised
in India as ``pay 5,000 rupees today and save 500,000 rupees tomorrow.'' In both countries, it is
illegal to inform parents of the sex of their fetus, and sex-selective abortion is banned. But it is
practiced widely and rarely punished.

Finally, because higher education and income levels generate more resources, development offers new
opportunities to discriminate against living girls. After all, if people are very poor, boys and
girls are necessarily deprived equally -- there is little to dole out to anyone. But as parents gain
the tools to help their children survive and thrive (and indeed, all children do better as their
parents' education and income levels advance), they allocate advantages like doctor visits to boys
and firstborn girls, leaving subsequent daughters behind.

To be sure, development can eventually lead to more equal treatment for girls: South Korea's birth
ratios are now approaching normality. But policymakers need to realize that this type of development
works slowly and mainly indirectly, by softening a son-centered culture. The solution is not to
abandon development or to stop providing, say, microcredit to women. But these efforts should be
joined by an awareness of the unintended consequences of development and by efforts, aimed at
parents, to weaken the cultural preference for sons.

The lesson here is subtle but critical: Development brings about immense and valuable cultural
change -- much of it swiftly -- but it doesn't necessarily change all aspects of a culture at the
same rate. (India and China have myriad laws outlawing discrimination against girls that are widely
ignored. And how to explain the persistence of missing girls among Asian immigrants in America?) In
the short and medium terms, the resulting clashes between modern capabilities and old prejudices can
make some aspects of life worse before they make them better.

\section{Bernanke, a Hero to His Own, Can't Shake Critics}

\lettrine{B}{en} S.~Bernanke, chairman of the Federal Reserve, no longer looks sleep-deprived.

He still works seven days a week, but earlier this month he took two days off -- for the first time
in two years -- to attend his son's wedding. And he often gets home for dinner and even out to
baseball games every few weeks.

As central bankers and economists from around the world gather on Thursday for the Fed's annual
retreat in Jackson Hole, Wyo., most are likely to welcome Mr.~Bernanke as a conquering hero. In
Washington and on Wall Street, it would be a surprise if President Obama did not nominate
Mr.~Bernanke for a second term, even though he is a Republican and was appointed by President George
W.~Bush.

But the White House has remained silent. And despite Mr.~Bernanke's credibility in financial
circles, both he and the Fed as an institution have come under political fire from lawmakers in both
parties over the handling of particular bailouts and the scope of the Fed's power.

He has been frustrated that many in Congress do not give the Fed what he believes is enough credit
for what it has accomplished. Indeed, Mr.~Bernanke has met privately with hundreds of lawmakers in
recent months to explain the Fed's strategy.

Fellow economists, however, are heaping praise on Mr.~Bernanke for his bold actions and steady hand
in pulling the economy out of its worst crisis since the 1930s. Tossing\cite{toss} out the Fed's
standard playbook, Mr.~Bernanke orchestrated a long list of colossal\cite{colossal} rescue programs:
Wall Street bailouts, shotgun weddings, emergency loan programs, vast amounts of newly printed money
and the lowest interest rates in American history.

Even one of his harshest critics now praises him.

``He realized that the great recession could turn into the Great Depression 2.0, and he was very
aggressive about taking the actions that needed to be taken,'' said Nouriel Roubini, chairman of
Roubini Global Economics, who had long criticized Fed officials for ignoring the dangers of the
housing bubble.

But Mr.~Bernanke is hardly breathing easy. Unemployment is still at 9.4 percent, and the central
bank's own forecasts assume that it will remain that high through the end of next year. Even if all
goes according to plan, Fed officials said, Mr.~Bernanke's current popularity could sink if the
recovery proves slower than many people expect.

While the White House keeps mum about Mr.~Bernanke's future, the leading Democratic candidates to
replace him include Lawrence H.~Summers, director of the National Economic Council; Janet L.~Yellen,
president of the Federal Reserve Bank of San Francisco; Alan S.~Blinder, a Princeton economist and
former Fed vice chairman; and Roger Ferguson, another former Fed vice chairman.

Mr.~Bernanke faces two major challenges. On the economic front, the Fed has to decide when and how
it will reverse all its emergency measures and raise interest rates back to normal without either
stalling the economy or igniting inflation.

On the political front, Mr.~Bernanke is trying to defend the Fed's power and independence as the
White House and Congress debate plans to overhaul the system of financial regulation.

Democrats like Senator Christopher J.~Dodd of Connecticut, chairman of the Senate Banking Committee,
contend that the Fed was too cozy with banks and Wall Street firms as the mortgage crisis was
building. House Republicans, and some Democrats, complain that the Fed already has too much power.

``Why does the Fed deserve more authority when institutionally it seemed to have failed to prevent
the current crisis?'' asked Senator Dodd last month.

The political battle over President Obama's plan to overhaul financial regulation has put
Mr.~Bernanke in an awkward position.

Fed officials support the administration's proposals to put them in charge of systemic risk like the
growth of reckless mortgage lending or the misuse of financial derivatives\cite{derivative}. But
they chafe\cite{chafe} at the plan to shift the Fed's consumer-protection functions, which protect
people from deceptive and unfair lending practices, to a new agency.

Mr.~Bernanke has avoided publicly criticizing the White House's call for an independent consumer
regulatory agency. While acknowledging that the Federal Reserve did nothing to stop mortgage
practices during the housing bubble, Mr.~Bernanke has argued that the Fed has since written tough
new protections for both mortgage borrowers and credit card customers.

``We think the Fed can play a constructive role in protecting consumers,'' he told the House
Financial Services Committee last month.

Mr.~Bernanke and other Fed officials now concede they failed to anticipate the full danger posed by
the explosion of subprime mortgage lending. As recently as the spring of 2007, Mr.~Bernanke still
contended that the problems of the housing market were largely ``contained'' to subprime mortgages.
When panic over mortgage-backed securities began spreading through the broader credit markets in
late July 2007, Fed officials initially refused to cut interest rates.

By December 2007, Mr.~Bernanke became increasingly convinced that the economy itself was in trouble
but policy makers were unable to reach agreement and decided not to reduce interest rates.

At a meeting on Jan.~21, 2008, the Fed slashed the benchmark federal funds rate by 0.75 percent, to
3.5 percent, the biggest one-time reduction in decades. Nine days later, officials cut the rate
again, down to 3 percent.

As the credit crisis deepened, Mr.~Bernanke urged Fed officials to devise proposals that had never
been tried before. They responded with a kaleidoscope of emergency loan programs to a wide array of
industries.

``He has had tremendous courage throughout this episode,'' said Frederic S.~Mishkin, a professor at
Columbia University's business school and a former Fed governor.

Amid the chaos, Fed and Treasury officials made numerous mistakes. Their original idea for the \$700
billion to buy up bad mortgage assets held by banks has yet to get off the ground.

But economists say Mr.~Bernanke's most important accomplishment was to create
staggering\cite{stagger} amounts of money out of thin air.

All told, the Federal Reserve has expanded its balance sheet to \$1.9 trillion today, from about
\$900 billion a year ago. Analysts now caution that Mr.~Bernanke's job is only half complete. He
will eventually have to reel all that money back. He has already laid out elements of the Fed's
``exit strategy,'' but Fed officials have been careful to say it is still too early to pull back any
time soon.

\section{Dude, You Are So (Not) Obama }

\lettrine{L}{ast} week, if you wanted to use the latest slang to tell a friend he was cool, you
could have called him ``Obama,'' as in: ``Dude, you're rocking the new Pre phone? You are so
Obama.''

This week? Best not to risk it.

The sudden shift in meaning has nothing to do with the fortunes of the president, regardless of what
the health care debate may do to his cool factor. The fault rests entirely with what has happened to
the life span of slang, which seems to shorten with every click of the mouse.

``Obama'' was one of the most noteworthy new entries in ``U.C.L.A. Slang 6,'' a recently released
compendium of student colloquialisms.

And the word's very inclusion in the dictionary signifies that its street cred has
evaporated\cite{evaporate}.

``I think that word has completely left us,'' said Pamela Munro, the U.C.L.A. professor who edited
the latest edition of that dictionary, which is compiled once every four years.

What's a hipster (hepcat?) to do? Keeping up with the latest slang is at once easier and harder than
ever. The number of slang dictionaries is growing, both online and off, not to mention social
networking media that invent and discard words, phrases and memes at the speed of broadband. The
life of slang is now shorter than ever, say linguists, and what was once a reliable code for
identifying members of an in-group or subculture is losing some of its magic.

The Internet ``is robbing slang of a lot of its sociolinguistic exclusionary power,'' said Robert
A.~Leonard, a linguistics professor at Hofstra in Hempstead, N.Y., whose slang credentials include
being a founding member of the doo-wop group Sha Na Na, formed in the late 1960s. ``If you are in a
real inside group, you are manufacturing slang so that you can exclude the wannabes.''

And that becomes harder, he added, as the whole world has access to your language.

Part of the problem is that electronic media are making it too easy to compile dictionaries like
``U.C.L.A. Slang 6.'' While slang dictionaries have been around in one form or another since the
18th century, they now number more than a dozen in print, to say nothing of online resources like
UrbanDictionary.com and slangsite.com that are updated hundreds of times daily.

``It used to be that the guys who did all the slang dictionaries would take years just to track it
down,'' said David Crystal, a linguist who has written more than 30 books. ``Now, with the Internet,
you just put something up on Facebook and say 'send me in your slang terms,' and in a few days you
have hundreds of examples and the ability to check it out with other readers.''

Urban Dictionary, which is 10 years old, may be the ultimate example. It is a sprawling, chaotic
collection of street talk, all of it user-submitted, giving everyone access to the meaning behind
the coded lyrics of someone like Lil Wayne.

Slang is meant to be ``something that keeps groups together and keeps people out,'' said Aaron
Peckham, who began the site when he was a freshman at California Polytechnic State University. But
the slang on Urban Dictionary is ``from every group you can imagine, so it's helping people
understand each other.''

In July, the Urban Dictionary attracted 15 million unique users and 1,000 new words a day,
Mr.~Peckham said.

But widespread understanding is the opposite of what slang is about. Indeed, it is accepted wisdom
among linguists that once a word actually shows up in a slang dictionary, it effectively ceases to
be slang.

What does that mean for language in the age of Urban Dictionary? Will it be the end of slang?
Hardly, says Jonathon Green, author of ``Chambers Slang Dictionary'' and other books about
neologisms. ``I think slang is the salsa, the great hot sauce on our language,'' he said. ``I think
apart from losing its power, it keeps reproducing itself. There are now 2,500 words for 'drunk.'
Soon there will be 3,500.''

In the past, slang has proved resilient, if only because it recycles itself almost as quickly as it
wears out. Mr.~Green said that, when he was a hippie in England in the 1960s, ``the language that we
were using was in fact the language of 1930s black America, though very few of us were aware of
this.''

Back then, ``it took 20, 30 years to cross the Atlantic,`` he said. ``The difference now is it takes
20 to 30 hours.''

Tracking a word's arc from hip to lame is notoriously difficult, especially because different social
groups grab hold of different terms at different times. But Professor Leonard cites ``pwned''
(rhymes with owned) as a relatively recent invention that is already falling out of favor with the
gaming crowd that coined it. ``It comes from the mistyping of 'owned' on some computer game,'' as in
``I just owned you,'' he said.

And the site Gawker recently tried without success to ban the phrase ``I'm just sayin,' '' which has
become ubiquitous on blogs and Twitter as a way of defanging -- disingenuously, perhaps -- a
potentially confrontational statement. In a sign of just how quickly such phrases are now being
co-opted, CNN recently unveiled a segment called ``Just Sayin,' '' in which an anchorwoman offers a
common-sense opinion about some curious cultural phenomenon, followed by the passive-aggressive
catch phrase.

``These words that once might have marked you as a member of a certain group or social set, as soon
as they get into circulation, they spread very quickly and lose their specialness,'' said Gabriel
Snyder, Gawker's editor in chief.

Indeed, banning tired slang has become something of a regular feature on the site, a reaction to the
rapid pace at which insider terms become worn out.

So is slang in danger of losing its cool? Not exactly, says Ms.~Munro, the U.C.L.A. linguistics
professor.

People who learn slang secondhand, she says, will tend to use it incorrectly.

``I feel that your grandmother would have a real hard time sounding like Lil Wayne,'' she said. 

\section{Tension Increases as China and Australia Grow Closer }

\lettrine{C}{hina}'s diplomatic relationship with Australia, so recently flourishing despite
occasional spats, this month has taken a severe turn toward the governmental
equivalent\cite{equivalent} of thrown dinner plates.

Public exchanges between the nations, already testy after China's detention of four employees of the
British-Australian mining giant Rio Tinto, grew sharper when Australia granted a visa to Rebiya
Kadeer, the American-based rights advocate for China's Muslim Uighur minority. Ms.~Kadeer was
accused by Chinese officials of plotting riots last month in China's Xinjiang region.

The Australians recalled their Chinese ambassador to the capital, Canberra, for talks on Wednesday,
after a week in which Beijing's state-controlled news media excoriated\cite{excoriate} Australia's
``Sinophobic politicians'' and suggested that China's billions were better spent trading with
friendlier nations.

The Chinese also canceled planned visits by Vice Premier Li Keqiang, the heir apparent to Prime
Minister Wen Jiabao, and the vice foreign minister, He Yafei, who was supposed to attend a meeting
of Asian nations. Columnists in the Chinese press have also advocated limiting Chinese tourism in
Australia and curbing the number of Chinese students studying there.

Australia's prime minister, Kevin Rudd, countered that the nations' relationship is always ``full of
challenges'' and that their broader ties will endure. ``We share enormous common interests with our
friends in China, but we have continuous differences,'' he was quoted as saying.

Hardly all Australians are persuaded. ``I really don't think there's anything that Australia can
do,'' J.~Bruce Jacobs, a China specialist at Monash University in Australia, said of the tiff. ``The
Chinese seem to have various people they like to pick on -- the French, because of the Dalai Lama,
and us, because of Kadeer. I think all of this is driven by political imperatives\cite{imperative}
within China.''

Mr.~Jacobs was referring to China's decision to boycott a European Union summit meeting last
December because the union's leader then, President Nicolas Sarkozy of France, planned to meet the
Dalai Lama, the spiritual leader of Tibetans. The Chinese accuse the Dalai Lama of plotting to split
Tibet from China.

In the latest case, China sought this month to keep Australia from granting Ms.~Kadeer a visa to
attend the screening of a film about her life, then tried to prevent her from making a speech to the
National Press Club.

They were further examples of how Australians and Chinese have chafed at their inexorably growing
ties. Trade between China and Australia has grown sevenfold this decade, making China Australia's
largest trading partner. Chinese investment in Australia, while still small compared with its
investment in the United States, is mushrooming.

But Australians worry that Chinese investment is directed at their vast natural resources, turning
them into a sort of open-pit mine for Chinese interests.

Mr.~Rudd, who is fluent in Mandarin and was once an Australian diplomat in Beijing, has advocated
deeper cooperation with China in global economic forums and with President Obama.

Despite that, the relationship has foundered lately on two issues that frequently divide Beijing and
the West: Chinese industrial policy and human rights.

Many Chinese expressed frustration this summer after the collapse of a deal for a state-owned
company to acquire a 19 percent stake in Rio Tinto, a crucial supplier of iron ore to China's steel
mills. Although economic factors stopped the deal, domestic suspicion of China's intentions toward
an Australian corporate icon was an undercurrent in the talks.

Relations deteriorated further in July after China arrested four Rio Tinto employees involved in
iron-ore sales on espionage charges, accusing them of stealing state secrets about ore pricing. The
Chinese later decided to pursue only lesser charges, but outraged Australians -- including Mr.~Rudd
-- warned that the prosecution would threaten China's commercial relations with the outside world.

In Beijing, one political analyst said Friday that the rift was unlikely to cause lasting damage to
Chinese-Australian relations.

``The mainstream of the two countries' relationship remains stable and friendly, even though there
are some kinds of problems,'' said the analyst, Su Hao, director of the Center for Strategic and
Conflict Management at China Foreign Affairs University in Beijing. The current spats, he said, are
``technical'' issues.

\section{Workers Return to Restive China Region }

\lettrine{T}{he} first wave of workers has arrived in the annual migration to China's
restive\cite{restive} western region of Xinjiang this year to pick cotton, according to a report on
Friday by Xinhua, the state news agency.

The workers are mostly ethnic Han and are the first large batch of migrant workers to make the
journey to Xinjiang since deadly ethnic rioting broke out this summer.

The Xinhua report appeared intended to support the government's contention that harmony had been
restored to the region and that migrant workers felt safeenough to travel there in large numbers.

On July 5, ethnic Uighur mobs went on a rampage through the streets of Urumqi, the capital of
Xinjiang, after clashes erupted between Uighur protesters and riot police officers. At least 197
people were killed and 1,721 injured, most of them Han civilians, according to state news
organizations.

The rioting was the deadliest outburst of ethnic violence in China in decades.

Uighurs interviewed in Urumqi said the government had underestimated the number of Uighur
casualties\cite{casualty}, especially those victimized by vengeful Han mobs in the days afterward.
The Uighurs are a Turkic-speaking people who mostly follow Sunni Islam. The Han are the dominant
group across China. In Xinjiang, the Uighurs are still the largest ethnic group, but an influx of
Han migrants in recent decades is changing the demographic.

The Chinese government has blamed Rebiya Kadeer, a Uighur exile who lives in the Washington area,
for the initial rioting. She has denied involvement. Uighurs in Urumqi said anger against the Han
and the Han-dominated government had been simmering\cite{simmer}. The Uighurs say they feel intense
discrimination in Xinjiang, and they often complain about the Han migrants who compete with Uighurs
for jobs.

Many Uighurs question why companies in Xinjiang, including large government-run farms, employ so
many Han migrants when there is a high unemployment rate among Uighurs.

The Xinhua report on Friday said most of the cotton pickers were women from the northwestern
provinces of Gansu and Qinghai, the Ningxia Hui Autonomous Region and the central province of Henan.

A total of about 1,600 workers from Henan arrived at their work sites late Wednesday and early
Thursday, Xinhua reported. In all, 100 trains full of workers will arrive in Xinjiang through
mid-September.

Xinjiang, a vast territory that makes up one-sixth of China, is the largest cotton-producing region
in the country.

The Xinjiang Production and Construction Corps, the state enterprise that runs the bingtuan, the
vast farms set up in the 1950s and afterward by the People's Liberation Army to help populate the
frontier region, reported a drop in the number of migrants going to Xinjiang after the July riots.

``In half a month after the violence, there was a huge shortage of tomato pickers,'' Wang Wei, a
bingtuan official, told Xinhua. ``Many migrant workers didn't dare to come here.''

Another bingtuan official, Sun Zhiqiang, said the migrants earned the equivalent of nearly \$15 a
day on average.

\section{A Privileged World Begins to Give Up Its Secrets}

\lettrine{A}{bout} 10 years ago, when I was working in Frankfurt, Germany's banking capital, I was
invited to the top floor of the glittering skyscraper headquarters of one of the country's most
venerable\cite{venerable} banks. There, I was treated to something that, it was made clear to me,
few eyes usually had the privilege of seeing -- a tour of its private art collection, an impressive
spattering of modern and ancient European and American masters.

The point was, those pictures reflected the bank's wealth. And the fact the secretive treasures were
kept forever behind closed doors for the enjoyment of the privileged few reflected its power.

If that seems like a different era, it is. Banks around the world are reeling\cite{reel}, as we
know; the European banks' losses are among the most ruinous\cite{ruinous}. And their prestige and
putative\cite{putative} secrecy and independence received a further blow last week, when the
government of Switzerland agreed to release to the United States the names of 4,450 American
citizens suspected of using secret Swiss accounts at UBS, the country's biggest bank, for tax
evasion\cite{evasion}.

The victory for the United States was made possible by evidence from an American-born whistleblower
-- code name Tarantula -- a disgruntled former UBS employee from the Boston area who was working in
Switzerland. Until he left the bank, he was part of a UBS team that made frequent trips across the
Atlantic to aggressively market investment strategies to rich Americans to elude\cite{elude} the
scrutiny\cite{scrutiny} of the Internal Revenue Service.

But it would be wrong to see the settlement as a one-off strike against just one bank by a single
government. It is in fact the result of a broader political moment created in the wake of the global
financial crisis when disenchantment with financial globalization is causing governments to
repatriate wealth back to within national borders, especially at a time when countries badly need to
balance their books.

Just a few years ago, in the pre-crisis era, the shadowy workings of cross-border banking -- and
what may or may not have been happening there -- were generally overlooked.

And, while some of the alleged tax evaders may be the war criminals, gunrunners or despots usually
linked with secret foreign bank accounts, the target of the latest efforts are much more likely to
include rich businessmen and high-net-worth individuals. ``There is a political movement because of
the financial debacle,'' said one veteran European banker who insisted on speaking anonymously
because he has retired. ``They are turning toward the so-called rich and want to hurt them.''

Of course, the United States looks at it a bit differently. Prosecutors have contended that in the
UBS case alone, wealthy Americans hid billions of dollars, thereby evading taxes of hundreds of
millions of dollars a year.

While Switzerland is arguably the largest off-shore center, it is not the only one. Supporters of
its banking secrecy code point out that the code is wrapped up in the country's claims to
neutrality\cite{neutrality} and being above the global political fray. But secrecy has also turned
out to be immensely lucrative; according to some estimates one-quarter of the world's offshore money
now resides in Switzerland.

Other countries or territories have copied the model -- Liechtenstein, Bermuda, the Cayman Islands,
Macao and Hong Kong among them. And while Switzerland is probably seen as the most conservative,
blue chip, upstanding offshore haven, the others are measured by a sliding scale of probity and
association with dubious\cite{dubious} business practices, if not crime. The European banker said
that in the early 1990s, following the fall of the Soviet Union, he worked in Switzerland where he
said agents of Russian expats would show up with ``boxes of cash'' from Cyprus, a popular haven for
capital fleeing the Russian authorities and the country's post-collapse chaos.

The backlash against this illicit\cite{illicit} world has not been confined to the United States; it
is apparent across Europe, too.

France will become of one of the first European countries to put in place a new tax treaty with
Switzerland to improve transparency and access to banking information. Germany is in discussions
with Liechtenstein over issues related to tax evasion by German companies and individuals.
Liechtenstein has also struck a disclosure agreement with Britain, encouraging British clients of
Liechtenstein banks to volunteer information to British tax authorities in return for reduced
penalties. In Italy, tax officials have started an investigation into whether the estate of the late
Gianni Agnelli, the former chairman of Fiat, has money hidden away in Switzerland. In Britain, the
government has become particularly exercised by tax competition -- the offering of low tax rates and
other advantages like tax secrecy to lure capital away.

In the Swiss settlement last week, the American authorities got the information they needed after
they saw an opportunity in the weakness of UBS, a bank that once enjoyed a sterling\cite{sterling}
global reputation but has suffered billions of dollars in losses linked to United States subprime
securities and had to be saved by a big government bailout last October. For the Swiss government,
the deal lifts the immediate threat of heftier legal action and frees the bank -- one of the
mainstays\cite{mainstay} of the Swiss economy -- to concentrate on recovery.

But will anything really change? Although the United States is supposed to learn the identities of a
few thousand tax evaders, those names will go first to an intermediate tax administration in
Switzerland for review. The actual process of recovering the names may become lost in bureaucracy
and foot-dragging.

Moreover, as The Times reported last week, smaller Swiss banks say they are confident that they can
continue to profit by finding new, more elaborate ways to protect the privacy of their clients.
Those banks continue to help clients hide billions of dollars through complex structures in offshore
havens.

But the I.R.S. commissioner, Doug Schulman, said the agreement with UBS was a ``major step forward''
in the government's efforts to pierce bank secrecy, and he warned that ``wealthy Americans who have
hidden their money offshore will find themselves in a jam.''

In the new political climate, expect to see a few rich Americans shifting uncomfortably.

\section{Asia's Recovery Highlights China's Ascendance }

\lettrine{I}{n} past global slowdowns, the United States invariably led the way out, followed by
Europe and the rest of the world. But for the first time, the catalyst is coming from China and the
rest of Asia, where resurgent\cite{resurgent} economies are helping the still-shaky West recover
from the deepest recession since World War II.

Economists have long predicted that an increasingly powerful China would come to rival and
eventually surpass the United States in economic influence. While the American economy is still more
than three times the size of China's, the nascent\cite{nascent} global recovery suggests that this
long-anticipated change could arrive sooner than had been expected.

Such a shift would have significant ramifications\cite{ramification} for the United States and the
rest of the West, even after the global economic recovery takes hold.

``The economic center of gravity has been shifting for some time, but this recession marks a turning
point,'' said Neal Soss, chief economist for Credit Suisse in New York. ``It's Asia that's lifting
the world, rather than the U.S., and that's never happened before.''

China's government-dominated, top-down economy is surging after Chinese banks doled out more than
\$1 trillion in loans in the first half of the year, in addition to a nearly \$600 billion
government stimulus program.

Though the benefits are manifest, some economists wonder whether China is laying the groundwork for
sustainable growth or just increasing its export capacity despite more frugal\cite{frugal} spending
habits on the part of Western consumers.

``The big question is what happens next,'' said Kenneth S.~Rogoff, a professor of economics at
Harvard. ``If the consumer in the United States and Europe doesn't come back, I'm not sure Asia has
a Plan B.''

But robust demand among Chinese consumers and businesses is one reason oil prices have doubled to
more than \$70 a barrel since bottoming out early this year, and China is likely to keep buying
American debt as Washington borrows heavily to finance its myriad\cite{myriad} stimulus and bailout
plans.

The United States is also being shoved aside as the make-or-break customer for export-driven nations
like Germany and Japan. China overtook the United States as Japan's leading trading partner in the
first half of 2009, while in Europe manufacturers are looking east instead of west.

``What we're losing in the trans-Atlantic trade with the U.S., we are gaining in China,'' said Jens
Nagel, head of the international department of the German Exporters Association.

In the near term, however, the United States should benefit from a resurgent Asia, as the American
economy finally begins growing again, as expected in the second half of 2009.

``Vigorous rebounds overseas, particularly in East Asia, suggest that U.S.~imports and exports will
soon improve,'' Mr.~Soss said.

Last week, Hewlett-Packard pointed to double-digit revenue growth in China as a rare bright spot in
an otherwise lackluster earnings report. Meanwhile, overall American exports to China have already
been picking up, rising to \$5.5 billion in June from \$4.1 billion in January.

``The numbers are volatile\cite{volatile}, but the trend is clear,'' said Robert Brusca of FAO
Economics in New York. ``It's a big contrast with Japan, where U.S.~exports are still dropping, but
China is different.''

Of course, other factors have played a significant role in helping the global economy begin to
stabilize, including trillions of dollars in support from central banks for frozen credit markets,
as well as bailouts and rescues of major financial institutions, insurers and automobile companies.

But as the engine for future demand growth shifts from the government back to the private sector,
and Americans remain wary of returning to their free-spending ways, Asian consumption is expected to
pick up at least some of the slack. And if China does slow, as some experts fear it could in the
second half of 2009, the United States' effort to climb out of recession could be that much harder.

After the recession of 2001-2 and the slowdown in the early 1990s, the American economy served as
the global locomotive, said Michael Saunders, head of European economics research for Citigroup.

Back then, he said, China and other Asian countries lacked huge cash reserves that could
buttress\cite{buttress} them in the event of recession. But in the last decade, China has enjoyed
huge trade surpluses with the West, and it holds \$2.13 trillion in foreign reserves, solidifying
its position as a rapidly emerging economic power.

Citigroup recently increased its estimate for annual Chinese economic growth to 8.7 percent in 2009
from 8.2 percent, and to 9.8 percent next year from 8.8 percent.

While economists like Mr.~Soss expect that growth to spill over to the United States shortly, the
effect is already visible in Europe.

Indeed, after the French and German economies shocked most economists this month by turning in
positive performances for the second quarter, the normally conservative Deutsche Bank released a
report titled, ``Eurozone Q2 GDP: Made in China?''

For now, the answer seems to be yes. ``It's quite amazing, because usually Asia doesn't play such a
big role in European exports or output,'' said Gilles Moec, senior European economist with Deutsche
Bank in London.

French exports to China and other East Asian economies rose 18.7 percent in the second quarter,
according to customs data, a sharp turnaround from the 16.2 percent drop recorded in the previous
quarter. Overall exports to the region from the 16 countries that use the euro currency increased
6.3 percent in the second quarter, reversing a 6.2 percent drop in the first quarter, Mr.~Moec said.

While Western European countries have been more timid about embarking on big spending programs
because of their already mounting deficits, and European banks took huge hits on their holdings of
subprime American debt, Beijing does not face either obstacle.

In the first half of 2009, Chinese banks lent a record \$1.1 trillion in new loans, setting off
fears that the lending binge might create a bubble over the long term.

China's moves have also helped its neighbors increase industrial production sharply from recession
lows. Since hitting a trough in late 2008 and early 2009, industrial production has jumped 28
percent in Korea and 26 percent in Taiwan. In July, American industrial production rose for the
first time since December 2007, but it remains just half a percentage point above the bottom in
June.

``Asia is still relatively small in the world, but it reflects how the world is changing, and
economic power does translate, of course, into political power,'' said Simon Johnson, a former chief
economist for the International Monetary Fund and now a senior fellow at the Peterson Institute for
International Economics. ``You can use it to win friends and influence people, as the Chinese are
already doing in Africa and Latin America.''

\section{Nokia Plans to Start Making Netbooks }

\lettrine{N}{okia}, the world's largest maker of mobile phones, said on Monday that it would make
and sell its own netbook computers, expanding into the fastest-growing segment of the PC market.

The company, which is struggling to keep up with the iPhone from Apple and the BlackBerry from
Research In Motion in the smartphone market, has not made a computer since the 1980s.

Like smartphones, netbooks are increasingly being used to access the Internet from outside the home.

``It is fair to say that Nokia had no choice but to do this,'' said Neil Mawston, an analyst at
Strategy Analytics in London. ``It's a high-growth, high-value segment that Nokia couldn't afford to
ignore.''

In the only details it has revealed, the company said that the Nokia Booklet 3G would run Microsoft
Windows software and use an Intel Atom processor. The company said it would have a 10-inch screen
and work on third-generation mobile and Wi-Fi networks.

Kai Oistamo, the executive vice president for devices at Nokia, said the company's decision was
driven by consumer demand for practical, lightweight devices designed for mobile broadband.

``The Nokia Booklet 3G is a natural evolution for us,'' Mr.~Oistamo said in a statement. ``We will
create something compelling.''

Nokia said it would provide more details, like when it would begin selling netbooks and for how
much, at an event in Stuttgart next week.

Netbooks, which first appeared on the market about 18 months ago, are expected to make up a third of
the 46 million portable computers that will be sold this year in Western Europe, according to the
International Data Corporation.

Western Europe is the biggest market for the devices, accounting for 70 percent of global sales,
I.D.C. said.

The market in Western Europe is dominated by two Taiwanese manufacturers, Acer and Asus, which
combine for 49 percent of sales. Samsung of South Korea and two American computer makers,
Hewlett-Packard and Dell, combine to sell 32 percent.

Nokia, which has a 38 percent share of the global cellphone market, is hoping to exploit its
distribution channels and longstanding relationships with mobile operators to sell netbooks.

But the telecommunications industry has had mixed results so far in trying to sell netbooks, mostly
as a lure to win long-term customers for mobile data plans, said Eszter Morvay, an analyst for
I.D.C. in London.

During the 2008 Christmas holiday season, operators in Britain, France and Italy bought far too many
netbooks for use in their promotional efforts, Ms.~Morvay said, and were unable to sell them during
the holiday rush.

In Britain, netbooks that are not subsidized as part of any promotion are selling for as little as
\textsterling150, or \$246, Ms.~Morvay said. In Germany, they sell for less than 300 euros, or
\$429.

``The netbook market is incredibly competitive,'' Ms.~Morvay said. ``The handset and netbook markets
are very different in all aspects. So far, the telco industry is still learning how to sell these.''

\section{In Health Care Debate, Bioethicist Becomes a Lightning Rod for Criticism }

\lettrine{F}{ew} people hold a more uncomfortable place at the health care debate's intersection
between nuanced policy and cable-ready political rhetoric than President Obama's special health care
adviser, Dr.~Ezekiel J.~Emanuel.

Largely quoting his past writings out of context this summer, Betsy McCaughey, a former lieutenant
governor of New York, labeled Dr.~Emanuel a ``deadly doctor'' who believes health care should be
``reserved for the nondisabled'' -- a false assertion that Representative Michele Bachmann,
Republican of Minnesota, repeated on the House floor.

Former Gov. Sarah Palin of Alaska has asserted that Dr.~Emanuel's ``Orwellian'' approach to health
care would ``refuse to allocate medical resources to the elderly, the infirm and the disabled who
have less economic potential,'' accusations similarly made by the political provocateur Lyndon
H.~LaRouche Jr.

In fact, Dr.~Emanuel has written more than a million words on health care, some of which form the
philosophical underpinnings of the Obama administration plan and some of which have enough
free-market elements to win grudging respect from some conservative opponents.

The debate over Dr.~Emanuel shows how subtle philosophical arguments that have long bedeviled
bioethicists are being condensed, oversimplified and distorted in the griddle-hot health care
debate. His writings grapple\cite{grapple} with some of the most complex issues of medical ethics,
like who should get the kidney transplant, the younger patient or the one who is older and sicker?

Perhaps it should come as no surprise that Dr.~Emanuel, an oncologist, has come to personify the
most intense attacks on the president's plan.

He is the older brother of the White House chief of staff, Rahm Emanuel, and the Hollywood
superagent Ari Emanuel. As a leading bioethicist at Harvard and at the National Institutes of
Health, Dr.~Emanuel had a reputation for pushing limits while exploring uncomfortable life-and-death
issues in starkly academic terms.

The level of vitriol against him has led even some conservative opponents to defend Dr.~Emanuel
while expressing concern that it is overtaking what they say are more vital real-world critiques.

``He is a serious oncologist and bioethicist, so the kinds of charges that have been raised against
him are particularly inappropriate,'' said Gail R.~Wilensky, a Republican and senior White House
health care adviser under the first President George Bush who criticizes Mr.~Obama's plan as being
too reliant on the federal government.

Given Dr.~Emanuel's well-publicized repudiations of doctor-assisted suicide and voluntary
euthanasia, and his calls for a national health insurance voucher system that would eventually
eliminate Medicare, Medicaid and employer-provided insurance -- nonstarters at the White House --
Dr.~Emanuel says he is perplexed by depictions of him as a socialist euthanasia proponent.

``You can only call me someone who's interested in euthanizing patients and denying care to patients
by willful distortion of my record,'' he said in an interview.

Dr.~Emanuel rose to prominence in the late 1980s with a popular standardized medical directive that
made it easier for terminally ill patients to share their wishes with doctors before becoming too
sick to speak for themselves.

Concerned with the hard questions that arise without such directives, Dr.~Emanuel included in his
1991 book, ``The Ends of Human Life'' (Harvard University Press), a critique of a court ruling
upholding a family's request to end treatment for a dying, mentally incapacitated daughter. He
argued that the ruling, in the case of Karen Ann Quinlan, did not provide an adequate ethical
framework for such a weighty decision in the absence of a patient's stated wishes.

In a 1997 article in The Atlantic, he argued against doctor-assisted suicide and euthanasia, warning
it would ``become the rule in the context of demographic and budgetary pressures,'' and ``would make
us want to extend the option to others who, in society's view, are suffering and leading purposeless
lives'' -- concerns reflecting the exact opposite of the views his critics now ascribe to him.

Peter R.~Orszag, the president's budget director, said in an interview that he had hired Dr.~Emanuel
on his own merits, as opposed to his brother's advice, after he offered to help with health care
policy. Mr.~Orszag said he was not surprised that Dr.~Emanuel's writings had drawn scrutiny.

``You can look at anyone who has written tons of stuff and play the same game,'' he said.

Ms.~McCaughey seemed to have evidence for her conclusion that ``he explicitly defends discrimination
against older patients'' in a recent New York Post opinion article. She quoted from a paper he
co-wrote for Lancet in January: ``Even if 25-year-olds receive priority over 65-year-olds, everyone
who is 65 years now was previously 25.''

But she did not report that the paper was addressing the allocation of ``very scarce resources''
like kidneys or vaccines, not the system in general.

Dr.~Emanuel's argument -- that young adults should take priority in vying for limited health
resources because they will get more years of life from them -- is a fairly mainstream if unpleasant
approach to a problem with only bad choices, ethicists and doctors of varying persuasions say.

``These kinds of dilemmas go on every day in clinical practice,'' said Dr.~Scott Gottlieb, a
physician and scholar at the American Enterprise Institute, a conservative research group. ``There's
a very big leap to say his contemplations about how doctors contend with these issues extends to
saying he believes government should take on these issues.''

Dr.~Gottlieb opposes the administration's proposals, calling them too prescriptive, too expensive,
and too open to eventual increased rationing.

In a brief interview, Ms.~McCaughey said that either way, because of its Medicare cost cuts, ``the
president's proposal will force hospitals to operate with scarce resources.''

The administration disputes that assertion.

Ms.~McCaughey, Ms.~Palin and others have based accusations that Dr.~Emanuel would direct treatment
away from the disabled on a 1996 paper he wrote for the Hastings Center bioethics institute.

In it, Dr.~Emanuel did not assert that ``medical care should be reserved for the nondisabled, `` as
the critics have said.

The paper laid out what he called a growing consensus among competing political philosophies about
how a society should allocate health care services. In clinical terms, he said that consensus held
that those who ``are irreversibly prevented from being or becoming participating citizens'' should
not be guaranteed the same level of treatment as others.

He cited as an example, ``not guaranteeing health services to patients with dementia.''

Dr.~Emanuel said he was simply describing a consensus held by others, not himself.

But even some colleagues said in interviews that the paper did not go far enough in repudiating the
view.

``He doesn't ever endorse it, nor does he explicitly distance himself from it,'' said Thomas
H.~Murray, president of the Hastings Center. But, Mr.~Murray added, ``anyone who would attribute
this isolated sentence to his convictions, it's just unfair.''

Dr.~Emanuel said he understood some of the criticisms.

``Maybe if I had been a smarter, more careful thinker about how people could interpret it, I would
have qualified it and condemned it more robustly,'' he said. ``In my 1.2, 1.3 million written words,
you can't find another sentence that even comes close to advocating that in my voice. When I
advocate, I'm not shy.''

\section{Airlines Are Sweetening Frequent-Flier Programs }

\lettrine{A}{fter} years of shrinking benefits and increasing fees, frequent-flier programs are
courting their most loyal customers again.

In recent weeks, American Airlines introduced new awards for one-way flights; United dropped its
\$75 to \$100 fees for booking an award within three weeks of travel; and Delta Air Lines made it
easier for elite frequent fliers to retain their V.I.P. status.

Carriers have also been dangling mileage offers, including double-mile promotions, and have even
resurrected bonuses for booking online.

Those who follow the programs describe the changes as a slight shift in a more consumer-friendly
direction, giving customers more ways to earn and redeemmiles.

``What we're seeing is a bit of a resurgence\cite{resurgence} on the airlines' part in focusing on
their frequent-flier programs,'' said Tim Winship, who tracks loyalty programs for
SmarterTravel.com. But he added that the issue members care about most -- being able to use their
miles for free tickets -- remains a challenge, especially with planes flying nearly full.

``I think award availability remains the elephant in the room,'' he said. ``The travel slump has not
resulted in any increase in award availability, and in fact, if anything my sense is that people are
having a more difficult time now.''

That assessment, while hotly debated on travel Web sites, is difficult to confirm.

Randy Petersen, founder of the frequent-flier hubs InsideFlyer.com and FlyerTalk.com, said 2008 was
a record year for award redemption, which he attributed to travelers using miles for more mundane
trips, like visits to relatives in less exotic locales.

``In a good economy, everyone wants to use their miles to go to Hawaii, London or San Francisco,''
he said. ``But when the recession hit, award redemption really started to spread out because people
were using their miles not as much for vacation as necessary travel.''

Another factor spurring award redemption, Mr.~Petersen said, has been the shift to ``dynamic
pricing.'' If a seat that requires 25,000 miles is available for the outbound leg of a trip, but
there is only a 50,000-mile seat on the return, many airlines now offer the itinerary for an amount
in between -- say, 37,500 miles -- rather than charging the higher price for the entire trip.

In a similar vein, the new one-way award from American Airlines lets members book a flight after as
little as 12,500 miles, a good option for students returning to college or snowbirds heading south
for the winter.

There have also been changes in financial regulations that require airlines to account for unused
frequent-flier miles as a liability, which has motivated carriers to make awards easier to redeem.

``There really is an incentive now for the airlines to get all of their miles redeemed because it's
expensive for them to carry them as a liability,'' Mr.~Petersen said.

Even so, he acknowledged that as travelers prioritize low-priced tickets over flying with a
preferred airline, carriers need to work harder to earn loyalty, especially among their elite
fliers.

Delta Airlines, for example, recently announced new benefits for its elite SkyMiles members, which
take effect in 2010.

Among the more significant changes, Delta added a ``diamond'' tier to its existing levels of
Medallion membership (silver, gold and platinum) -- all referred to as having elite status, which
comes with benefits like free upgrades and fee waivers.

Paul R.~Skrbec, a Delta spokesman, said adding a fourth tier encouraged Medallion members to stay
with Delta, rather than switch to another carrier after obtaining platinum status. The higher the
tier, the better the benefits. For example, a diamond member will get a 125 percent mileage bonus
versus a 100 percent bonus for platinum elites.

Seth Miller, a freelance technology consultant in Manhattan, is one of the elite frequent fliers
airlines have been courting, having maintained top-tier platinum status in Continental's OnePass
program for 10 years.

``The legacy airlines especially are doing as much as they can to generate loyalty, and these days
that means crazy bonuses -- double miles for this, triple miles for that, miles are being given out
like candy on Halloween,'' Mr.~Miller said.

``It's great now,'' he added, but he worries that the airlines' largess could have a downside. ``The
question is, what happens when everyone tries to take advantage of all these benefits that are being
handed out now? I've got elite status, but am I going to be able to get upgrades or will everybody
else have status, too?''

Another question is whether nonelite frequent fliers have dropped off the radar as airlines focus on
their biggest spenders -- a message that average consumers may be responding to by shifting their
loyalty to other rewards programs, like those offered on credit cards.

A research company that tracks loyalty programs found that membership in credit card reward programs
surpassed membership in frequent-flier programs for the first time in 2009.

According to Colloquy, the company that conducted the research, the average household in the United
States is signed up for 14 loyalty programs, ranging from grocery stores and gas stations to
airlines and hotels, but actively participates in only six.

The recession has diminished participation in multiple travel programs, said Kelly Hlavinka, a
partner at Colloquy. She said this could bring about a return to the original premise of loyalty
rewards: to cement a relationship with just one airline or hotel.

``Savvy travelers may be saying, 'I may not be able to spread my business out to two or three
airlines, but I can consolidate my travel with one company,' '' Ms.~Hlavinka said. ``The real
opportunity for airlines is to try to keep that business with their airline.''


\section{Tucked Away in Shanghai, Hidden Lives}

\lettrine{F}{or} the last couple of months I have spent the first part of each day either teaching
at a Chinese university or writing.

Nearly every afternoon, though, in what has distinctly felt like the start of a new day, I have set
off with camera in hand by motorcycle and subway to some of the fast-disappearing old neighborhoods
of this city, to knock on the doors of hundreds of ordinary, working-class people.

These encounters with strangers have plunged me deep into a world experienced by few foreigners, and
indeed, one might venture, few Chinese -- particularly those of the middle class.

Through the time spent in the cramped, dimly lit homes of my subjects -- people whose portraits I've
taken for a long-term photographic project about the city's oldest neighborhoods -- I may have
learned as much about Shanghai and about China as I did in five busy years as a correspondent here.

Typically, I enter their world by climbing up a rickety, twisting wooden staircase, ducking to avoid
bumping my head in the near-total darkness. This experience, eerie at first, but now utterly
familiar, has come to feel appropriate for a photographic adventure, like the adjustment of one's
eyes, and perspective, upon entry into a darkroom.

My subjects come fresh to the experience, so it has been unexpected and unquestionably strange for
each of them, at least initially. Once they have overcome their surprise at the sight of a very
tall, camera-bearing, Chinese-speaking foreigner in the sanctum\cite{sanctum} of their tiny homes,
the most common question has been: What could possibly be interesting about a place like this?

The answer is: plenty.

The demographics of this city, said to have the oldest population in a country that has begun to age
rapidly, has come to life before my eyes. I had not expected to find so much evidence of China's
thriving quasi-underground religious culture here. In house after house, I found people worshiping
privately as Christians or Buddhists. Asked how she had come to the church, a woman who had been
sent to the countryside as a youth in the Cultural Revolution told me she had been converted by her
neighbors. ``Everyone in this building believes in Christ,'' she said.

We are ever more accustomed to dazzling images of China, the fast-rising nation that may soon
surpass the United States and lead the world, according to one increasingly widespread trope. Those
who know a bit about the country will be aware that there are still many hundreds of millions of
people in the countryside who have not yet found a spot on China's economic
escalator\cite{escalator}.

Even in China's richest city, huge numbers of people eke out a very modest existence. To be sure,
these are very often migrants from provinces like Anhui or Jiangsu, or even further afield. But more
than most Chinese would suspect -- particularly the proud, newly affluent generations of Shanghai
people who look at my photographs and sniff ``wai di ren,'' or ``outsiders'' -- a great many of the
denizens of the city's dilapidated but character-rich old quarters are natives.

Much has been written lately about growing social inequality in China. The country's social
divisions, however, are much more complicated than statistics suggest, involving lots of
fine-grained, identity-based prejudices.

I think, for example, of the poor and jobless Shanghainese parents in the old garment district who
told me of their eagerness to be relocated across the river to Pudong, where the environment would
be better, in part they said, because there would be fewer of the ``wai di'' people, whom they
dismissed as having ``no culture.''

Others pessimistically dismissed the likelihood that China's increasing prosperity would continue to
lift all boats. ``I'm frightened for my son's future,'' said a migrant from Henan. ``China's biggest
problem is the population. There are just too many of us, and the competition for opportunity is
murderous.''

Inevitably, the theme of relocation comes up often in encounters like these, given the frantic pace
of redevelopment. Some people are pleased with the take-it-or-leave-it buyout arrangements the
government has offered to pave the way for the construction of high-rises; others respond with
fatalism. ``If the country needs this land, what can I do?'' said one elderly man.

A great many people spoke bitterly and with surprising candor\cite{candor}, though, about what they
see as a crisis of social justice. Here, I think of the 75-year-old owner of a tiny barbershop whose
neighborhood came down before my eyes this summer.

``What they are doing here is simply unfair,'' he said, telling me how thugs had been dispatched to
beat up residents who refused to quietly make way for the demolition. ``There is no rule of law. The
'lao bai xing' have no rights at all.'' That old phrase, meaning the nameless masses, never seemed
more appropriate.

Others told me the stories of corrupt local officials, whom they said offered higher compensation
for relocated people who were willing to pay bribes. These anecdotes took on special potency in a
summer where a nearly completed apartment building fell on its side, killing a worker and setting
off lurid rumors of government corruption.

I learned that large numbers of Chinese understand and value democratic ideals and yearn for them to
be applied here. ``We may have gotten richer, but our politics have not really evolved since
imperial times,'' said one elderly man. ``Chinese people want democracy as much as anyone else, and
one day we will have it.''

\section{China Approves Law Governing Armed Police Force}

\lettrine{S}{enior} members of China's legislature approved a law on Thursday detailing the
authority of the People's Armed Police, a large paramilitary force that was criticized in some
quarters as slow to respond to the riots in the western Xinjiang region last month, in which nearly
200 people died.

The new law appears to address those concerns by clarifying how and when the troops may be deployed.

The law is the first to explicitly govern the force, whose members serve as border guards, security
guards for government officials, firefighters and relief workers during disasters but who are best
known outside China for their role in suppressing political and social unrest.

The troops will have authority in ``handling rebellion, riots, large-scale serious criminal
violence, terror attacks and other social safety incidents,'' according to a summary of the law
published by the state-run Xinhua news agency before the measure was approved.

The legislation also apparently removes the authority of county-level local officials to summon the
force to handle disorders. Chinese citizens have complained that the armed police are sometimes
enlisted by low-level government officials to abusively bolster their powers, sometimes with
excessive force.

Estimates of the size of the People's Armed Police have ranged as high as 1.5 million troops, but
the government said the number in 2006 was 660,000. Once seen as corrupt and ill trained, the armed
police have become more professional in recent years. But they sometimes are regarded as
unaccountable, and human rights advocates have accused the troops of brutality and insensitivity to
the law in handling civil disorders.

The troops' contact with average Chinese citizens has grown in recent years as China's economic
transformation -- and dislocation -- has led to an increasing number of spontaneous street protests.

President Hu Jintao, in Xinjiang on Tuesday, told armed police troops there that ensuring social
stability was ``the most urgent task'' they faced.

Wang Yukai, a professor at China's National School of Administration, said in an interview on
Thursday that stripping local officials of the authority to summon the police was perhaps the most
important provision of the new legislation.

``This is to prevent the misuse of the armed police by local governments, to prevent the deaths of
innocent people,'' he said.

Conflict is inevitable given China's rapid changes, Mr.~Wang said, but ``sometimes it is
counterproductive to use force.''

``I think the government understands this and does not want to see it happen,'' he added.

\section{China's Press: Still Not Free, but More Freewheeling}

\lettrine{M}{odern} Weekly is about as hip as it gets in China, and about as successful as it gets
in this country's small universe of independent magazines. It comes in a glossy, large format with
slick graphics, lots of short quotes and pictures and some sophisticated reporting in separate
sections -- news, business, lifestyle and culture -- including a good deal of news, business,
lifestyle and culture in countries other than China.

The flagship publication of Shao Zhong, the People's Republic of China's first private media
entrepreneur, the magazine has a circulation of 700,000, relatively still small given the country's
population of 1.3 billion, but it is nonetheless a trend-setter in this country, where the media in
general are still heavily censored and controlled.

``We give our readers an amazing amount of information,'' Mr.~Shao said during a recent conversation
in his Beijing office. ``We tell them a lot about what's going on in the world, and in these areas
there's a lot of space.''

What Mr.~Shao, who uses the English name Thomas, meant by a lot of space is political space, room
for reporters to roam. He was responding to the inevitable question from a foreign journalist about
what restrictions he has to accept in order to run a successful magazine empire in China. He
wouldn't, he was asked, be doing much reporting on Tibet or on violations of human rights in China,
would he?

``I tend to stay away from topics like that,'' Mr.~Shao said.

Of course he does, and he has to. It's the bargain that Chinese journalism has to make to survive in
this country, where everybody knows what can and can't be reported. When a couple of weeks ago, for
example, China's public security bureau arrested Xu Zhiyong, an activist lawyer and founder of the
Open Constitution Initiative, which advocates the rule of law in China, the news was widely reported
abroad, but blacked out in China itself. (Although Mr.~Xu was released on bail on Sunday, he is
still expected to face charges of tax evasion.)

Similarly, there's very little reporting on the situation in the western region of Xinjiang, scene
of last month's ethnic riots, except for dutiful repetition of the government's position and
denunciations of Rebiya Kadeer, the exiled Uighur leader whom China blames for the disturbances.
Chinese readers are likely to find in their own press more about the private life of Barack Obama
than about their own president, Hu Jintao, because reporting is pretty free on Mr.~Obama yet tightly
regulated on Mr.~Hu.

In a recent issue of Mr.~Shao's Modern Weekly, there was a two-page spread on what the magazine
called ``India's Secret Submarine,'' all by way of reporting on India's naval development over the
past decade or so. It is out of the question of course that Modern Weekly would carry out an
unauthorized examination of China's naval development, which would be deemed a state secret.

As everybody knows, the press is not free in China. Yet there's no question that things are
changing, and Modern Weekly illustrates the emergence of a press that, certainly, is at least more
freewheeling than ever before in China, and separate from the government.

``Rarely will you find a successful magazine these days that's completely run by the government,''
Mr.~Shao said.

Mr.~Shao, 48, is tall, lean, and casual. He started out in his native town of Guangzhou in China's
south, where he put in a few years in city government, as a member of the Guangzhou City Development
and Planning Commission.

Around 1990, he said, he began, thanks to China's opening to the outside world, to see a lot of
foreign magazines, and he contrasted them unfavorably to those being published in China, which, he
said, were ``ugly.''

He got some very specific inspiration from a book, ``Seventy Years of Time Magazine'' that
chronicled how Time's founder, Henry Luce, built his magazine publishing empire, a story that
inspired Mr.~Shao to strive to be a sort of Chinese Luce -- ironic, given Luce's famous
animosity\cite{animosity} to China's Communists.

Mr.~Shao liked the way Luce's magazines combined pictures and texts and divided coverage into
sections -- similar to the sections that Mr.~Shao publishes in Modern Weekly.

But in China, anyone who wants to publish a magazine has to find a partner with a publishing
license, and that means a state-owned newspaper or magazine, which, essentially furnishes legal
cover, making no investment, taking no risk, and doing no work.

``There have been a lot of difficulties,'' Mr.~Shao said. ``I've failed in some startups and lost
some partners, that sort of thing. You have to collaborate with a license holder, and when you're
successful, they just take you over.''

Mr.~Shao's first partner for his first weekly startup was the Guangdong Daily, which, after a while,
took control of that magazine, showing Mr.~Shao to the door.

``When that happens, there's nothing you can do,'' he said. ``You can't sue the guy. You just have
to start over.''

That's what Mr.~Shao has done, clearly with success. His company, Modern Media Group, now publishes
10 magazines, with 10 different official license partners; it has a staff of 700 employees and
offices all around China. His Beijing office, where I met him, is sleekly modern, like his
magazines, and is adorned with large contemporary paintings of the sort being sold for millions of
dollars at Sotheby's these days.

If Mr.~Shao feels frustrated by the restrictions still placed on journalists in China, he doesn't
show it.

He seems happy to be able to be private entrepreneur with enough scope to publish magazines that he
feels are useful, and perhaps that's the broader compromise made in China today -- you can operate
freely in lots of areas, as long as you agree to stay out of the areas that the authorities have
fenced off.

``Nobody tells you what to do,'' Mr.~Shao said. ``There are no rules. You have to feel your way
around with content and not send up any red flags, so we concentrate on international news and also
lifestyle and trends, which have much less to do with government.''

It's a useful thing to do, Mr.~Shao feels as well as a profitable one.

``We put our readers in step with what's going on in the world,'' he said.

\section{Hey, PC, Who Taught You to Fight Back?}

\lettrine{S}{ean} siler would never be mistaken for a movie star. A former Navy officer who wears
glasses and is a tad on the heavy side, Mr.~Siler works at Microsoft, where he oversees the Windows
division's adoption of new Internet connectivity software called IPv6.

But there were audible gasps\cite{gasp} last summer when Mr.~Siler, 39, auditioned for Microsoft's
new ad campaign for Windows, created by Crispin Porter \& Bogusky, the Miami agency best known for
its cheeky\cite{cheeky} work for Mini Cooper and Burger King.

``I was like, 'Are you kidding?' '' recalls Rob Reilly, one of the agency's executive creative
directors. ``It couldn't have been more perfect.''

Everybody agreed that Mr.~Siler looked exactly like PC, the character played by the comedian John
Hodgman in Apple's popular ``Get a Mac'' ads that lampoon\cite{lampoon} Windows-based computers and
those who love them. Two weeks later, Mr.~Siler reported to a nearby television studio. The agency
dressed him in PC's dorky uniform -- white shirt, baggy khakis, brown sport coat and matching brown
tie -- and handed him a script with the lines: ``I'm a PC. And I've been made into a stereotype.''

Mr.~Siler joined a parade of environmentalists, budget-conscious laptop shoppers, mixed martial arts
fighters, mash-up DJs and remarkably tech-savvy preschoolers who appear in Microsoft's new campaign,
which is intended to show that real Windows users aren't all clueless drones\cite{drone}.

For Mr.~Siler, the experience was almost like being a geeky incarnation\cite{incarnation} of Brad
Pitt. His e-mail address was on the screen, and he received 4,000 messages from viewers -- some from
grateful parents whose children had wanted expensive Macs over PCs and now had second thoughts.

Crispin put up a video on YouTube in which Mr.~Siler discussed his role in the campaign; it was
viewed more than 702,000 times. At work, he was constantly interrupted by his fellow Microsoft
employees. ``For a couple of weeks,'' Mr.~Siler recalls, ``I had people coming by my office and
saying: 'Hey, you are the PC guy, aren't you? That's so cool!' ''

His mother wasn't so sure. ``You look so horrible,'' she told him. ``You don't look anything like
that man. Why did they make you look so bad?''

Somebody better explain to Mr.~Siler's mother that this isn't a beauty contest; it's an ad war, one
destined to go down in history with the cola wars of the 1980s and '90s and the Hertz-Avis feud of
the 1960s. According to TNS Media Intelligence, Apple spent \$264 million on television ads last
year, 71 percent more than Microsoft. In the first six months of 2009, however, Microsoft responded
with \$163 million worth of commercials, more than twice Apple's spending.

Surprisingly, Microsoft, which has never been known for running cool ads, has landed some punches.
Shortly after the Microsoft campaign started, Apple unleashed commercials that mocked its competitor
as spending money on advertising when it should have been fixing Vista, its much-maligned operating
system.

``It got Apple's attention, didn't it?'' says Robert X.~Cringely, host of PBS's NerdTV.

FOR years, Microsoft was the stodgy market leader. It sold 90 percent of the world's operating
system software, and generally left the advertising to Dell, H.P. and other hardware makers who
licensed Windows. The only time Microsoft hawked its most recognizable brand on television was when
the latest version of the software hit the shelves. Then the company flooded the airwaves with
commercials full of loud music and swirling imagery saying that the new version of Windows is out --
and that it's awesome!

Apple is the classic smaller insurgent. Its share for desktops and laptops in the United States is
just over 8 percent. Every time Apple grabs another point of market share from Microsoft's partners,
its stock price climbs. And one way that Apple has tried to gain share is by running clever ads that
ridicule everything Microsoft stands for.

There's no better example than ``Get a Mac,'' unveiled three years ago by Apple's longtime ad
agency, TBWA/Chiat/Day. No technology company would choose Mr.~Hodgman's character, PC, to personify
its brand. He reeks of the past. He boasts of using his desktop to make spreadsheets and ridicules
his more youthful friend, Mac, played by the actor Justin Long, for using his desktop for
``juvenile'' pursuits like blogging and movie making -- even through it's clear that PC would like
to be in on the fun. He just can't get his Windows computer to do his bidding.

Like a classic sitcom character -- think Ralph Kramden of ``The Honeymooners'' -- PC is always
dreaming up ill-advised schemes intended to show his superiority. He's thwarted by viruses, system
crashes and other problems more associated with Windows-based computers than Apple's products --
and, recently, he has become a hapless apologist for Vista. Mr.~Long's character smugly watches his
friend's pratfalls, glancing at the audience with raised eyebrows as if to say, ``If only this poor
guy would buy a Mac\ldots''

PC will never learn. Not as long as he keeps driving sales for Apple. Since 2006, the year that he
first appeared in all his pasty-faced glory, Apple's share of the computer desktop market in the
United States has more than doubled, according to IDC, the technology industry research firm. Its
stock price, meanwhile, has risen 142 percent since May 2006, while Microsoft's has barely budged.
Yes, the astonishing success of newer Apple products like the iPod and the iPhone has helped. But
the PC character should also take a bow. (Representatives of Apple and TBWA/Chiat/Day declined to be
interviewed for this article.)

Apple's ads put Microsoft in a bind. One of Madison Avenue's rules is that a market leader never
acknowledges a smaller competitor in its advertising. What's more, if Microsoft responded with ads
that backfired, it would look just like Mr.~Hodgman's character. Maybe it was better to grin and
bear it.

Then, last year, Microsoft hired Crispin Porter and struck back with uncharacteristic wit. There was
Mr.~Siler's star turn. The agency also handed bunches of cash to shoppers and asked them to choose
between a PC and a Mac. Lauren, a 20-something in one of the ``Laptop Hunter'' spots, is
giddy\cite{giddy} about the money she has left over when she selects a \$699 H.P. with a 17-inch
screen, rather than a \$1,000 Mac with a 13-inch screen. ``I guess I'm just not cool enough to be a
Mac person,'' she sighs. This time, the joke was on Apple. In a recession, it's pretty hip to save
\$300.

Microsoft's effort to inspire PC pride seemed to resonate\cite{resonate} after its debut last
September. According to IDC, Mac shipments in the United States plummeted 20 percent in the fourth
quarter of 2008 versus the previous quarter, as the economy went into a tailspin, while those of PCs
manufactured by Dell and H.P. fell only 13 percent and 3 percent, respectively.

Microsoft was quick to declare victory -- maybe too quick. In the second quarter this year, Mac
sales in the United States rebounded 34 percent, IDC said, while Dell and H.P. had more modest
gains. Even more humbling for Microsoft was the company's announcement in late July that its
year-over-year operating income for the quarter declined 29 percent.

As a result, some analysts have argued that the Microsoft campaign has failed. But they, too, may be
too hasty. We are only weeks away from the Oct.~22 release of Windows 7, which may undo much of the
company's self-inflicted damage from Vista. PC users, many of whom skipped buying Vista machines,
could be holding off until then to buy. And the introduction of Windows 7 will be accompanied by yet
another Crispin Porter ad blitz.

``You are not so embarrassed to take your PC out of the bag on a plane anymore,'' said Mr.~Reilly at
the ad agency. ``It's actually kind of cool that you do. I know this is working.''

EVERY Wednesday, Lee Clow, the creative director of TBWA/Chiat/Day, travels from Los Angeles to
Cupertino, Calif., for his weekly meeting with Steven P.~Jobs, the Apple chief executive. They
started doing this years ago and have created ads that are as stylish and cool as anything on
television. Usually, the subtext of these ads is that Microsoft is the Evil Empire.

Mr.~Jobs started working with Mr.~Clow, a laid-back former surfer dude, in the early 1980s when
Mr.~Clow helped to create Apple's path-breaking ``1984'' television commercial introducing the
Macintosh. The ad's unsubtle message was that buyers of the new machine would be striking against
I.B.M., portrayed as Apple's Orwellian foe.

Mr.~Jobs struggled to persuade Apple's board to run the ad, which was directed by Ridley Scott.
Mr.~Clow was similarly adamant\cite{adamant} when his boss, the late Jay Chiat, tried to shelve it.
The ad ran only once, during the 1984 Super Bowl, but it has never been forgotten.

Apple forced out Mr.~Jobs the next year and hired a new ad agency, BBDO. But when Mr.~Jobs returned
triumphantly to the company in 1997, he reunited with TBWA/Chiat/Day. Mr.~Clow brought him the idea
for ``Think Different,'' a campaign that identified Apple with figures like Bob Dylan, Albert
Einstein and Martin Luther King Jr.~Mr.~Jobs used it to introduce the iMac and to re-establish Apple
as an iconoclast.

TBWA/Chiat/Day went on to create the 2002 ``Switchers'' campaign, in which the director Errol Morris
filmed real computer users describing why they ditched their PCs for a Mac. Who can forget Ellen
Feiss, the slow-talking teenager who made the hearts of young geeks flutter when she explained how
her PC ate her homework? ``It was, like, beep beep beep beep beep beep beep,'' Ms.~Feiss says. ``And
then, like, half of my paper was gone.''

Then came the iPod ads from TWBA/Chiat/Day that not only helped drive sales of Apple's breakout
product, but also made stars of little-known indie rock acts like Feist. Such is the power of
Apple's marketing wizardry.

Many of Apple's new customers were plugging their iPods into PCs. Mr.~Clow proposed ``Get a Mac'' to
get them thinking about springing for an Apple machine. Mr.~Jobs was intrigued. But he wanted the
ads to be perfect.

``The discussion within Apple was: 'Is this the right tone? How young a guy should Mac be? How dorky
do we make PC look?' '' recalls Ken Segall, a former TBWA/Chiat/Day creative director who worked
early on as a consultant for Apple on the campaign. ``It went many rounds before Steve was
comfortable with the idea. Then he loved it.''

IN spring 2007, a year after Apple introduced the ``Get a Mac'' ads, Steve Ballmer, the Microsoft
C.E.O., barged into the office of Mich Mathews, head of the company's central marketing group. The
two had talked about a campaign that would repair the damage from the Apple ads.

Ms.~Mathews recalls Mr.~Ballmer enthusiastically asking her, ``When are we going to move?''

Advertising has never seemed to be part of Microsoft's DNA. The chairman, Bill Gates, ``never really
seemed to get marketing,'' says Rob Enderle, a longtime technology industry analyst. And for many
years, Mr.~Enderle says, Mr.~Ballmer ``just didn't think it was worth spending the money on it.''

The company's Windows campaigns seemed to reflect executives' lack of interest. Perhaps the best
example was the push for Microsoft Vista in 2007, created by McCann Erickson with the slogan ``The
Wow Starts Now.'' It showed people gaping in childlike wonder at the newest version of Windows. But
Vista, to put it mildly, didn't live up to the ads.

``The operating system was visually beautiful,'' said Jeff Musser, a former McCann Erickson creative
director who worked on the campaign. ``But it was a bad product. I didn't really hear anybody
saying, 'Wow.' ''

There were also cultural issues at Microsoft when it came to advertising. On Madison Avenue, they
say that the more hands that touch an advertisement, the worse it becomes. Microsoft felt
differently. ``They thought the more people saw it and gave an opinion, the better it would be,''
Mr.~Musser said. ``That's how you develop software. It's not how you develop great creative.''

So Ms.~Mathews tried to change things. She set up a nine-member task force to figure out a marketing
strategy and keep meddlers at arm's length.

In February 2008, Microsoft picked Crispin Porter. At the agency, Mr.~Reilly was initially
apprehensive. He didn't even own a PC; he had an ultraslim MacBook Air. (He has since bought himself
two PCs -- a Sony Vaio and a Lenovo ThinkPad.)

The adman also wondered whether Microsoft was ready for a Crispin campaign. Mr.~Reilly himself
oversees the agency's irreverent work for Burger King, aimed at young men hungering for menu items
like the Triple Whopper.

He wanted to come up with a campaign that would redefine Windows, and he counseled against ads that
attacked Apple. Then he changed his tune. Last summer in Apple ads, Mr.~Hodgman's PC character
morphed into a personification of Microsoft itself. PC was haunted by problems with Vista. He took
up yoga to calm his nerves, only to discover that his teacher was on edge because Vista wreaked
havoc on her billing system. PC tried to find peace by creating a line of herbal teas with names
like ``Crashy-Time Camomile'' and ``Raspberry Restart.''

``As the tone of their campaign became more and more negative, we were like, 'We gotta do
something,' '' Mr.~Reilly said. ``That's where the whole notion of 'I'm a PC' and putting a face on
our users came about. We have a billion users. That's who our cast is, whereas Apple is just two
fictitious characters.''

Microsoft recruited influential Windows fans like the ``Desperate Housewives'' star Eva Longoria.
``I feel bad about the little PC guy,'' she said this month. ``He is always getting beaten up.'' It
also brought in some who would appeal to niche audiences, like the Pittsburgh mash-up D.J.~Gregg
Gillis, who is better known as Girl Talk.

When Mr.~Ballmer finally saw the ads in September, he congratulated Ms.~Mathews and gave her a
high-five. Then, Ms.~Mathews says, he started shouting, ``I'm a PC!''

THE new Windows campaign got off to an inauspicious\cite{inauspicious} start. Puzzling ads featuring
Mr.~Gates kidding around with the comedian Jerry Seinfeld left a lot of people scratching their
heads. The ads quickly disappeared.

As the ``I'm a PC'' ads with Mr.~Siler replaced them two weeks later, Apple's ``Get a Mac'' spots
disappeared. Microsoft doesn't think that was a coincidence. When PC and Mac reappeared, it was in
the advertising that criticized Microsoft as spending on ads rather than on Vista.

Microsoft thought that it had scored a point. ``You've got to look at that and say, 'You are not
advertising to consumers; you're advertising to the Microsoft marketing department,' '' Ms.~Mathews
says. ``I just admit that did bring a smile to my face.''

Emboldened, Microsoft continued its barrages. In February, it unveiled its ``Rookies'' ads, arguing
that PCs are so easy to use that even Kylie, an adorable 4 1/2-year old, could upload a picture of
her goldfish, Dorothy, onto her PC and e-mail it to her relatives. You want to make fun of Kylie,
Apple? Microsoft and Crispin dare you to try it.

The next month, Microsoft deployed its ``Laptop Hunters'' ads. They clearly moved the needle in
Microsoft's favor. Ted Marzilli, a managing director of BrandIndex, a company that tracks consumer
perceptions, said that at the beginning of the year, adults thought Apple offered more value than
Microsoft. In May, however, Microsoft closed the gap in the firm's surveys. ``Apple took a hit,''
Mr.~Marzilli said. ``Since then, they have been neck and neck.''

In June, Microsoft felt that it had more reason to gloat\cite{gloat}. The chief operating officer,
B.~Kevin Turner, says he got a call from an Apple lawyer who asked him to change the ads because
Apple was lowering its prices by \$100. ``I did cartwheels down the hallway,'' Mr.~Turner
subsequently boasted in speech at a New Orleans conference.

Then Apple announced its second-quarter rebound. And for some analysts, it seemed like game over.
``The reality is that Apple's business has been impacted by the overall economy, not by Microsoft's
campaign,'' said Gene Munster, senior research analyst at Piper Jaffray. ``Those 'What can I get for
1,000 bucks' ads? That was a clever campaign. But it never really caught on. If you compare it to
'Get a Mac,' it didn't even register.''

And yet Apple keeps responding. On Friday, it released its Snow Leopard operating system a month
ahead of schedule, accompanied by a new round of ``Get a Mac'' ads. One involves a red-headed woman
who is clearly intended to resemble Microsoft's Lauren. PC introduces her to his suave friend, a
top-of-the-line model played by Patrick Warburton, who was David Puddy on ``Seinfeld.'' She declines
to buy a Windows machine when they can't promise that she won't have virus woes.

Microsoft, however, has found it enjoys mixing it up with Apple on the airwaves. In July,
Mr.~Ballmer told analysts that Crispin's work had been ``quite effective.'' He promised that
Microsoft would continue investing heavily in Windows marketing. ``We didn't do that three, four,
five, six years ago,'' he added.

For Mr.~Siler, this is a welcome change. ``I've never seen more pride at Microsoft,'' he says. ``You
walk through the campus, and you see people's laptops that have 'I'm a PC' stickers on them. I walk
in the company store, and there are these huge banners that say, 'I'm a PC' and shirts and ties and
mugs. I think I made a difference. My God, that's so cool!''

\section{Apple Reaches Agreement to Sell the iPhone in China}

\lettrine{A}{pple} has reached a deal with a Chinese wireless carrier, China Unicom, to sell the
iPhone in China, the world's biggest cellphone market.

Under a multiyear deal announced Friday, the iPhone is expected to go on sale in the fourth quarter,
China Unicom executives said in Hong Kong. They declined to give financial details or reveal how
much the iPhone would cost, saying only that the price would be ``competitive.''

Unicom, one of three major state-owned carriers, would be the first Chinese phone company to
formally support the iPhone, though unlocked iPhones brought in from other markets are in wide use
in China.

The mobile market in China trails the United States, Japan and some others in financial size but it
has 650 million mobile phone accounts and is seen as a major prize for foreign companies. However,
Edward Yu, chief executive at Analysys International, a technology research firm in Beijing, said it
remained to be seen whether the iPhone would appeal to ``the mainstream China population.''

The chief executive of Unicom, Chang Xiaobin, said the companies would not share revenue, with
Unicom instead buying the phones in batches from Apple and offering them with subsidies.

Many Chinese already own iPhones bought on the country's thriving gray market. There are well over
1.5 million of the devices in China, according to estimates by BDA China Ltd., a technology research
firm in Beijing.

Analysts said that could work in favor of Apple and China Unicom by providing a ready pool of users
who might be willing to upgrade to an authorized iPhone that offers a warranty, fewer glitches, more
applications and faster speeds.

An Apple spokeswoman, Jill Tan in Hong Kong, did not respond to an e-mail message seeking comment.

The iPhone, which also functions as a music player, camera and Web browser, could help Unicom
compete against China Mobile, which dominates the mainland's mobile market.

\section{Who's Driving Twitter's Popularity? Not Teens}

\lettrine{K}{risten} Nagy, an 18-year-old from Sparta, N.J., sends and receives 500 text messages a
day. But she never uses Twitter, even though it publishes similar snippets of conversations and
observations.

``I just think it's weird and I don't feel like everyone needs to know what I'm doing every second
of my life,'' she said.

Her reluctance to use Twitter, a feeling shared by others in her age group, has not doomed the
microblogging service. Just 11 percent of its users are aged 12 to 17, according to comScore.
Instead, Twitter's unparalleled explosion in popularity has been driven by a decidedly older group.
That success has shattered a widely held belief that young people lead the way to popularizing
innovations.

``The traditional early-adopter model would say that teenagers or college students are really
important to adoption,'' said Andrew Lipsman, director of industry analysis at comScore. Teenagers,
after all, drove the early growth of the social networks Facebook, MySpace and Friendster.

Twitter, however, has proved that ``a site can take off in a different demographic than you expect
and become very popular,'' he said. ``Twitter is defying the traditional model.''

In fact, though teenagers fueled the early growth of social networks, today they account for 14
percent of MySpace's users and only 9 percent of Facebook's. As the Web grows up, so do its users,
and for many analysts, Twitter's success represents a new model for Internet success. The notion
that children are essential to a new technology's success has proved to be largely a myth.

Adults have driven the growth of many perennially popular Web services. YouTube attracted young
adults and then senior citizens before teenagers piled on. Blogger's early user base was adults and
LinkedIn has built a successful social network with professionals as its target.

The same goes for gadgets. Though video games were originally marketed for children, Nintendo Wiis
quickly found their way into nursing homes. Kindle from Amazon caught on first with adults and many
gadgets, like iPhones and GPS devices, are largely adult-only.

Similarly, Twitter did not attract the young trendsetters at the outset. Its growth has instead come
from adults who might not have used other social sites before Twitter, said Jeremiah Owyang, an
industry analyst studying social media. ``Adults are just catching up to what teens have been doing
for years,'' he said.

Many young people, who have used Facebook since they began using the Internet and for whom text
messaging is their primary method of communication, say they simply do not have a need for Twitter.

Almost everyone under 35 uses social networks, but the growth of these networks over the last year
has come from older adults, according to a report from Forrester Research issued Tuesday. Use of
social networking by people aged 35 to 54 grew 60 percent in the last year.

Another reason that teenagers do not use Twitter may be that their lives tend to revolve around
their friends. Though Twitter's founders originally conceived of the site as a way to stay in touch
with acquaintances, it turns out that it is better for broadcasting ideas or questions and answers
to the outside world or for marketing a product. It is also useful for marketing the person doing
the tweeting, a need few teenagers are attuned\cite{attune} to.

``Many people use it for professional purposes -- keeping connected with industry contacts and
following news,'' said Evan Williams, Twitter's co-founder and chief executive. ``Because it's a
one-to-many network and most of the content is public, it works for this better than a social
network that's optimized for friend communication.''

Wendy Grazier, a mother in Arkansas, said her two teenaged daughters thought Twitter was ``lame,''
yet they asked her to follow teenage pop stars like Miley Cyrus and Taylor Swift on Twitter so she
could report back on what the celebrities wrote. Why won't they deign to do it themselves? ``It
seems more, like, professional, and not something that a teenager would do,'' said 16-year-old
Miranda Grazier. ``I think I might join when I'm older.''

The public nature of Twitter is particularly sensitive for the under-18 set, whether because they
want to hide what they are doing from their parents or, more often, because their parents restrict
their interaction with strangers on the Web.

Georgia Marentis, a 14-year-old in Great Falls, Va., uses Facebook instead of Twitter because she
can choose who sees her updates. ``My parents wouldn't want me to have everything going on in my
life displayed for the entire world,'' she said. (Of course, because of the public nature of social
networks and the ease of creating a fake identity on the Web, even sites with more privacy settings
have proved dangerous for young people in some cases.)

Many young people use the Web not to keep up with the issues of the day but to form and express
their identities, said Andrea Forte, who studied how high school students use social media for her
dissertation. (She will be an assistant professor at Drexel University in the spring.)

``Your identity on Twitter is more your ability to take an interesting conversational turn, throw an
interesting bit of conversation out there. Your identity isn't so much identified by the music you
listen to and the quizzes you take,'' as it is on Facebook, she said. She called Twitter ``a
comparatively adult kind of interaction.''

For Twitter's future, young people's ambivalence\cite{ambivalence} could be a good thing. Teenagers
may be more comfortable using new technologies, but they are also notoriously fickle\cite{fickle}.
Although they drove the growth of Friendster and MySpace, they then moved on from those sites to
Facebook.

Perhaps Twitter's experience will encourage Web start-ups to take a more realistic view of who uses
the Web and go after a broader audience, Ms.~Forte said. ``Older populations are a smart thing to be
thinking about, as opposed to eternally going after the 15- through 19-year-olds,'' she said.

\vocabulary
\end{document}
